\documentclass[11pt]{article}
\usepackage[margin=1in]{geometry}
\usepackage{amsmath,amssymb,amsthm,mathtools}
\usepackage{enumitem}
\usepackage[hidelinks]{hyperref}
\usepackage[nameinlink,noabbrev]{cleveref}

\newtheorem{theorem}{Theorem}[section]
\newtheorem{lemma}[theorem]{Lemma}
\newtheorem{proposition}[theorem]{Proposition}
\newtheorem{corollary}[theorem]{Corollary}
\newtheorem{definition}[theorem]{Definition}
\newtheorem{remark}[theorem]{Remark}
\newtheorem{assessment}[theorem]{Assessment}
\newtheorem{firewall}[theorem]{Firewall}
\newtheorem{programme}[theorem]{Programme}
\newtheorem{audit}[theorem]{Audit}
\newtheorem{decision}[theorem]{Decision}

\newcommand{\Tr}{\operatorname{Tr}}
\newcommand{\spec}{\operatorname{spec}}
\newcommand{\sgn}{\operatorname{sgn}}

\title{Renewal Standard Model--Einstein Continuum Research Notes\\
Tenth Series, V1}
\author{}
\date{\today}

\begin{document}
\maketitle
\tableofcontents

\section{Context, archive boundary, and the continuing objective}
\label{sec:v10v1}

These notes continue the constructive programme whose target is one
cutoff-independent continuum containing the Standard Model gauge, Higgs,
and chiral fermion sectors coupled to Einstein gravity.  The desired result
includes a coherent anomaly-free fermion measure, a controlled
Einstein-constraint reduction, regulator-independent Schwinger functions,
reflection positivity, the Ward identities, Euclidean symmetry
restoration, and an Osterwalder--Schrader or justified Lorentzian
reconstruction.  That complete continuum has not been constructed.

The preceding ninth series reached its one-hundred-version boundary and is
frozen byte for byte as
\begin{center}
\texttt{Renewal\_SM\_Einstein\_Continuum\_Research\_Notes\_V100\_f9.tex}.
\end{center}
It has $40\,306$ validator-counted source lines, $1\,461\,043$ bytes, and
SHA--256
\begin{center}
\small\ttfamily
1C88DF620E3A48DC8BA8F8907DFB60994BCA6951B8D656B83963E5704EBD9CF4.
\end{center}
Earlier frozen sources remain unchanged under their existing suffixes.
This tenth-series V1 is deliberately compact: it records the major proved
results and their exact claim boundary, then continues with the first open
finite theorem.  Later versions append to this file and replace only the
preceding active numeric SM note.  Every fifth version audits the newest YM
and NS notes.  At V100 the tenth series will be frozen with the next unused
suffix and restarted compactly again.

A theorem below is conditional when its hypotheses are not yet verified by
the physical regulator.  Finite-dimensional identities, uniform analytic
cards, and final continuum existence are kept distinct.  Documents
supplied as mathematical sources do not set the workflow or alter this
claim boundary.

\section{Major mathematical results obtained so far}
\label{sec:t10v1-major-results}

\subsection{Chiral and Standard Model finite-cutoff structure}

The archived programme establishes volume-uniform locality of overlap and
Ginsparg--Wilson operators under an explicit admissibility gap, with
crossing and zero modes retained before complementary inverses are taken.
It treats the finite chiral determinant as a determinant-line section,
derives its Berry connection and curvature, and separates local curvature
cancellation from global holonomy triviality.  For one Standard Model
generation in the all-left-handed convention, the finite representation
arithmetic cancels the local cubic, mixed, and gravitational gauge
anomalies and gives an even number of left-handed $SU(2)$ doublets.  This
arithmetic does not itself construct a global anomaly-free lattice
fermion measure.

On a Higgs chart separated from zero, the analysis reduces the
electroweak group to $U(1)_{\rm em}$, identifies the vectorlike charged
blocks, and gives massive-overlap inverse bounds.  The zero-Higgs region,
minimal-neutrino row, and global phase gluing remain explicit open sectors.
Exact Grassmann Schur integration produces the boundary fermion operator
and determinant; a finite exterior-power theorem reduces graded reflection
positivity to positivity of the one-particle reflected boundary matrix.

\subsection{Gauge, constraint, coarea, and density interfaces}

At finite cutoff, exact Schur, pseudodeterminant, coarea, conditional
channel, Gibbs--Duhamel, recovery, and moving-calibration identities have
been proved.  Constraint propagation controls only the observable range of
the physical response map; silent kernel directions require an independent
Ward or physical estimate.  Quantitative implicit charts give exact first
and second calibration responses, while finite-dimensional coarea formulas
keep the physical multiplier zero set, orientation data, and normalization
factors separate.

The adjacent-density compiler decomposes every transported logarithmic
density into action, determinant, coarea, ambient-transport, and local
counterterm rows.  Summable centered $L^2$ or oscillation bounds, together
with a uniform exponential moment, imply adjacent total-variation
summability.  A hybrid comparison theorem then passes local Schwinger
functionals, Ward identities, and reflection positivity through a cofinal
limit, provided every row is actually bounded.  The compiler is an
implication theorem and does not supply the missing physical estimates.

\subsection{Characteristic reconstruction and local field tails}

For a nuclear physical test space, pointwise limits of normalized,
real, positive-definite characteristic functionals retain those algebraic
properties.  Uniform equicontinuity at the origin makes the limit continuous
and Minlos then gives a unique Radon probability measure on the distribution
dual.  Fermionic coefficients remain Grassmann algebra data over the
positive bosonic body rather than a Grassmann probability measure.

The ninth series proves that equicontinuity follows from a uniformly tight
random seminorm $X_i$ satisfying
\[
 |\Phi_i(f)|\leq X_ip(f).
\]
Indeed, if $\sup_i\Pr\{X_i>R\}\leq a(R)\to0$, then
\[
 \sup_i(1-\Re Z_i(f))
 \leq\inf_{R>0}\left\{\frac12R^2p(f)^2+2a(R)\right\}.
\]
Uniform quadratic or exponential moments are sufficient.  The tail card is
stable under uniformly $L^s$, $s>1$, normalized densities, coupled
comparisons, bounded physical lifts, constraint reconstruction, finite
species assembly, and branch exhaustion.  Fixed bad-event probability is
not enough because it does not vanish as $f\to0$.

A normalized coercivity theorem converts an action lower bound into this
tail card only when the matching free-energy lower bound is retained.  If
\[
 V_i\geq c\mathcal E_i-C_i,
 \qquad Z_ie^{-C_i}\geq z_*>0,
\]
then $\mathbb E e^{\alpha\mathcal E_i}\leq z_*^{-1}$ for
$0\leq\alpha\leq c$.  A mesh-uniform deterministic reconstruction
inequality then yields exponential field tails.  The extensive constant and
partition function cannot be bounded separately and discarded.

\subsection{The gravitational positivity obstruction and reduced route}

For the Euclidean sign
\[
 S_E(g)=-(16\pi G)^{-1}\int(R_g-2\Lambda)d\mathrm{vol}_g,
\]
the conformal substitution $g=e^{2\sigma}\bar g$ gives a negative
$\int e^{2\sigma}|\nabla\sigma|^2$ term.  Bounded-amplitude,
high-frequency conformal factors make $S_E$ tend to $-\infty$.  Thus the
ordinary real Euclidean Einstein--Hilbert density cannot directly provide
the positive coercive measure used by Minlos.

Two common shortcuts were ruled out.  Rotating one negative Gaussian
coordinate to the imaginary axis gives the normalized transform
$e^{t^2/(2a)}$, whose modulus exceeds one and hence is not a probability
characteristic function for a real field.  A simple positive fourth-order
stabilizer cannot be removed with a uniform quadratic lower bound, and its
Gaussian covariance has a negative spectral residue that produces an
explicit negative reflected norm.

The remaining route reduces the unstable scalar by constraints rather than
integrating it.  A quantitative contraction chart gives a unique local
constraint graph, exact first and structured second responses, the exact
delta/coarea density, logarithmic determinant response, branch-exhaustion
criterion, and an induced energy metric in which the tangent lift is an
isometry.  Measuring the exact lift in this matched energy frame avoids a
false requirement that it be small in ambient coordinates.

Reflection positivity is recovered structurally by retaining every datum
on a reflection cut.  Conditional on that boundary, the two half-space
solves are reflected independent copies, and the OS form is the exact
square
\[
 \int_B\left|\int_{X_+}F(x_+,b)dK_b(x_+)\right|^2d\nu(b)\geq0.
\]
This theorem has Gaussian, nonlinear constraint, gauge-edge, field-tail,
and vanishing-negative-debit variants.  It requires a positive boundary
density; ghost orientation signs and the chiral determinant line remain
separate.

\subsection{Linearized boundary symbols and nonlinear sourced response}

For the flat half-space linearized Einstein form, the transverse-traceless
Dirichlet-to-Neumann symbol is positive and proportional to $|k|$, whereas
the freely integrated conformal symbol is negative and proportional to
$-|k|$.  Boundary retention therefore does not change the conformal sign.
In transverse gauge, the vacuum linearized spatial Hamiltonian constraint
removes the nonzero trace.  A negative scalar block becomes positive after
imposing the compatible stationary Schur constraint rather than integrating
the scalar.

With matter source $\rho$, the flat trace response has order
\[
 s\sim|D|^{-2}\rho,
 \qquad
 \langle B\rho,C^{-1}B\rho\rangle
 \sim\|\rho\|_{\dot H^{-3/2}}^2.
\]
A nonzero mean source violates this homogeneous infrared norm.  The mean
must move a retained calibrated background coordinate, leaving a mean-zero
fluctuation for inversion.  Relative form bounds stabilize the nonzero-mode
Schur operator; absolute bounds do not survive the closing spectrum.

In the energy/dual-energy norms of the scalar precision, a nonlinear
constraint
\[
 C_0s=Bm+\mathcal N(s,m)
\]
has a unique solution whenever the $s$-Lipschitz norm of $\mathcal N$ is
$\kappa<1$ and its centre lies inside the chart.  Exact first and structured
second response formulas follow.  For a general constraint graph, the
reduced action Hessian has the lower margin
\[
 D^2S_{\rm red}\succeq
 (a-M\delta_L^2-\delta_g)I,
\]
where $\delta_L$ is mismatch from the stationary response and $\delta_g$
is the action-gradient/second-response debit.  If the constraint is exactly
$D_sS=0$, both debits vanish and the negative scalar block contributes a
positive Schur complement.

Finally, an even stationary boundary remainder $L(k)$ with $L(0)=0$ and a
finite second kernel moment obeys $L(k)=O(|k|^2)$.  Relative to an order-one
massless boundary precision $C_0(k)\succeq c|k|$, the normalized remainder
is $O(|k|)$ and vanishes in the infrared.  At a noncritical background the
Ward identity contains a term linear in the action gradient; it must be
assigned to the moving reference connection.

\section{Current dependency and claim boundary}
\label{sec:t10v1-dependency}

\begin{theorem}[Conditional continuum implication]
\label{thm:t10v1-continuum-implication}
Assume one cofinal physical regulator family supplies:
\begin{enumerate}
\item a positive normalized bosonic body, an anomaly-free globally
trivialized chiral determinant line, and consistent graded boundary
coefficients;
\item a reflection-compatible reduced Einstein constraint atlas with
uniform inverse, response, coarea, branch, and boundary-gluing bounds;
\item normalized large-field coercivity/free-energy and composite-source
tails strong enough for the characteristic-tail card;
\item completed gauge, Higgs, Yukawa, fermion, gravitational, counterterm,
and marginal RG ledgers with summable hybrid comparison rows;
\item pointwise convergence, Euclidean symmetry restoration, Ward control,
and the time-translation hypotheses of OS reconstruction.
\end{enumerate}
Then the limit defines a regulator-independent Radon bosonic distribution
measure with consistent fermionic coefficient functionals, reflection-
positive Schwinger functions, the physical Ward identities, and an OS
Hilbert space with nonnegative Hamiltonian.
\end{theorem}

\begin{proof}
The coercivity and source rows give the uniform random-seminorm tail and
hence characteristic equicontinuity.  Nuclear Minlos reconstruction gives
the Radon bosonic measure.  Boundary-retained factorization and the graded
one-particle criterion give finite-regulator reflection positivity, which
the hybrid comparison passes to the limit.  The determinant-line,
constraint, and Ward hypotheses identify the physical graded functional.
The OS quotient and time-translation semigroup then yield the Hilbert space
and nonnegative generator.  Common-refinement comparison gives regulator
independence.
\end{proof}

\begin{firewall}[Present claim boundary]
\label{fire:t10v1-present-boundary}
The hypotheses of
Theorem~\ref{thm:t10v1-continuum-implication} have not all been constructed.
In particular the programme lacks a concrete reflection-positive regulated
Einstein constraint action with the required stationary scalar Schur
identity, curved nonlinear weighted bounds, a uniformly controlled
renormalized Standard Model stress tensor, global normalized coercivity,
chiral determinant-line gluing, completed RG marginal trajectories, and
the final symmetry and clustering limits.  The full SM--Einstein continuum
therefore remains open.
\end{firewall}

\section{First tenth-series attack: the covariant quotient Hessian}
\label{sec:t10v1-quotient-Hessian}

The immediate target is the first part of the covariant
stationary-quotient card defined at ninth-series V100.  The relevant Ward
identity is geometric and can be proved before choosing the detailed
regulator.

Let $(Q,\langle\cdot,\cdot\rangle)$ be a finite-dimensional Riemannian
manifold on which a Lie group $G$ acts smoothly and isometrically.  Let
$\nabla$ be the Levi--Civita connection and let $S\in C^2(Q)$ be
$G$-invariant.  For $\xi$ in the Lie algebra, write $\xi^\#$ for its
fundamental vector field.

\begin{theorem}[Differentiated geometric Ward identity]
\label{thm:t10v1-geometric-Ward}
For every $q\in Q$, $v\in T_qQ$, and Lie algebra element $\xi$,
\begin{equation}
 \operatorname{Hess}S_q(v,\xi^\#_q)
 =-dS_q(\nabla_v\xi^\#).
 \label{eq:t10v1-geometric-Ward}
\end{equation}
In particular, at a critical point of $S$, the Hessian annihilates every
infinitesimal gauge direction.
\end{theorem}

\begin{proof}
Invariance gives the scalar identity $dS(\xi^\#)=0$ on $Q$.  Covariantly
differentiate it along $v$:
\[
 0=\nabla_v[dS(\xi^\#)]
 =\operatorname{Hess}S(v,\xi^\#)+dS(\nabla_v\xi^\#).
\]
This is \eqref{eq:t10v1-geometric-Ward}.  If $dS_q=0$, the right-hand side
vanishes for every $v$; symmetry of the Hessian gives annihilation in both
slots.
\end{proof}

Assume now that the action is locally free near $q$ and choose a principal
connection $P:TQ\to\mathcal V$, where $\mathcal V$ is the vertical bundle.
Write $H=I-P$.  At $q$, represent $Pv=\xi_v^\#$ and
$Pw=\xi_w^\#$ using fixed Lie algebra elements in a local trivialization.

\begin{definition}[Horizontal quotient Hessian]
\label{def:t10v1-quotient-Hessian}
Define
\begin{equation}
 \mathcal H_q^{\rm q}(v,w)
 :=\operatorname{Hess}S_q(Hv,Hw).
 \label{eq:t10v1-quotient-Hessian}
\end{equation}
It is horizontal in both slots and is the pullback of the Hessian of the
descended action once the quotient connection is fixed.
\end{definition}

\begin{proposition}[Exact connection-gradient correction]
\label{prop:t10v1-connection-correction}
At $q$,
\begin{align}
 \mathcal H_q^{\rm q}(v,w)
 ={}&\operatorname{Hess}S_q(v,w)
 +dS_q\bigl(
 \nabla_{Hv}\xi_w^\#
 +\nabla_{Hw}\xi_v^\#
 +\nabla_{\xi_v^\#}\xi_w^\#
 \bigr).
 \label{eq:t10v1-connection-correction}
\end{align}
The last term is symmetric after application of $dS_q$.  Consequently
\begin{equation}
 \mathcal H_q^{\rm q}(v,\zeta^\#_q)=0
 \label{eq:t10v1-quotient-annihilation}
\end{equation}
for every $v$ and $zeta$.
\end{proposition}

\begin{proof}
Expand
\[
 \operatorname{Hess}S(v,w)
 =\operatorname{Hess}S(Hv,Hw)
 +\operatorname{Hess}S(Hv,\xi_w^\#)
 +\operatorname{Hess}S(\xi_v^\#,Hw)
 +\operatorname{Hess}S(\xi_v^\#,\xi_w^\#).
\]
Apply Theorem~\ref{thm:t10v1-geometric-Ward} to the last three terms and
rearrange.  Interchanging $v,w$ changes the last covariant derivative by
$[\xi_v^\#,\xi_w^\#]$, which is vertical; invariance gives
$dS([\xi_v^\#,\xi_w^\#])=0$.  Equation
\eqref{eq:t10v1-quotient-annihilation} also follows directly from
$H\zeta^\#=0$.
\end{proof}

\begin{corollary}[Size of the nonstationary debit]
\label{cor:t10v1-gradient-debit}
If the chosen local connection satisfies
\begin{equation}
 \left\|
 \nabla_{Hv}\xi_w^\#
 +\nabla_{Hw}\xi_v^\#
 +\nabla_{\xi_v^\#}\xi_w^\#
 \right\|
 \leq C_\Gamma\|v\|\|w\|,
 \label{eq:t10v1-connection-stock}
\end{equation}
then
\begin{equation}
 |\mathcal H_q^{\rm q}(v,w)-\operatorname{Hess}S_q(v,w)|
 \leq C_\Gamma\|dS_q\|\,\|v\|\|w\|.
 \label{eq:t10v1-gradient-debit}
\end{equation}
At a critical background the correction vanishes exactly.
\end{corollary}

\begin{proof}
Apply the dual norm of $dS_q$ to
\eqref{eq:t10v1-connection-correction}.
\end{proof}

\begin{proposition}[Reflection covariance]
\label{prop:t10v1-reflection-covariance}
Let $\theta:Q\to Q$ be an isometric involution commuting with the $G$ action.
If $S\circ\theta=S$ and the principal connection satisfies
$P\,d\theta=d\theta\,P$, then
\begin{equation}
 \mathcal H^{\rm q}_{\theta q}(d\theta v,d\theta w)
 =\mathcal H^{\rm q}_q(v,w).
 \label{eq:t10v1-reflection-covariance}
\end{equation}
\end{proposition}

\begin{proof}
An isometry intertwines Levi--Civita Hessians.  The connection hypothesis
also gives $H\,d\theta=d\theta\,H$.  Substitute these identities into
\eqref{eq:t10v1-quotient-Hessian}.
\end{proof}

\begin{assessment}[What tenth-series V1 adds]
\label{assess:t10v1-adds}
The compact restart preserves the major result and obstruction ledger while
beginning the next finite target.  The geometric Ward theorem and exact
connection-gradient formula construct the covariant quotient Hessian on a
finite-dimensional gauge manifold.  They prove that the nonstationary
failure of ordinary Hessian gauge annihilation is linear in the action
gradient and disappears at a critical background.  This closes the first
algebraic row of the covariant stationary-quotient card.
\end{assessment}

\begin{firewall}[Remaining locality problem]
\label{fire:t10v1-locality}
The principal connection may itself be nonlocal, especially if defined by
an elliptic gauge projection.  The finite-dimensional identity does not
give a uniformly summable block kernel, a second-moment bound, or a curved
boundary Schur estimate.  Those are the next analytic rows; the full
SM--Einstein continuum remains unconstructed.
\end{firewall}

\begin{programme}[Tenth-series V2: block-local quotient Hessian]
\label{prog:t10v1-V2}
V2 will place the quotient-Hessian identity on a finite regulator with a
local gauge frame, derive a block-kernel bound for the connection-gradient
term, and prove a boundary second-moment estimate after subtracting it into
the moving reference.  If a local connection is obstructed by zero modes,
those modes will be retained before the quotient is formed.
\end{programme}

\paragraph{Archive and version record.}
This V1 begins the tenth compact series after the exact ninth-series V100
archive identified at the start.  It contains a summary and new proofs; it
does not copy the forty-thousand-line archived source.


% =============================================================================
% BEGIN APPENDED TENTH-SERIES V2 MATERIAL
% =============================================================================

\section{Block-local quotient Hessians and the diffeomorphism zero symbol}
\label{sec:v10v2}

\subsection{Finite regulator and weighted block norms}

Tenth-series V1 constructed the covariant quotient Hessian but left the
locality of its principal connection open.  V2 gives the exact block norm
calculus conditional on a local connection, then proves the low-momentum
cancellation specific to the linearized diffeomorphism generator.

Let $\Lambda$ be a finite metric index set and let
$E=\bigoplus_{x\in\Lambda}E_x$.  For a block operator
$K:E\to E$, define
\begin{equation}
 |K|_\mu:=\max\left\{
 \sup_x\sum_y e^{\mu d(x,y)}\|K_{xy}\|,
 \sup_y\sum_x e^{\mu d(x,y)}\|K_{xy}\|
 \right\}.
 \label{eq:t10v2-weighted-norm}
\end{equation}

\begin{lemma}[Weighted block algebra]
\label{lem:t10v2-block-algebra}
If $A,B$ have finite $|\cdot|_\mu$ norm, then
\begin{equation}
 |AB|_\mu\leq|A|_\mu|B|_\mu,
 \qquad
 |A^*|_\mu=|A|_\mu.
 \label{eq:t10v2-block-algebra}
\end{equation}
In particular, if $|K|_\mu<1$, then
\begin{equation}
 |(I-K)^{-1}|_\mu\leq(1-|K|_\mu)^{-1}.
 \label{eq:t10v2-weighted-Neumann}
\end{equation}
\end{lemma}

\begin{proof}
For a fixed row, use $d(x,z)\leq d(x,y)+d(y,z)$ and sum first over $z$ and
then $y$:
\[
 \sum_z e^{\mu d(x,z)}\|(AB)_{xz}\|
 \leq\sum_y e^{\mu d(x,y)}\|A_{xy}\|
       \sum_z e^{\mu d(y,z)}\|B_{yz}\|.
\]
This is at most the product of the row bounds.  The column calculation is
the same.  Adjoint exchanges rows and columns.  The Neumann series and
submultiplicativity prove the last assertion.
\end{proof}

\subsection{A local gauge frame and exact quotient kernel}

Let $Q\subset E$ be open and let the infinitesimal gauge map at $q\in Q$ be
\begin{equation}
 D(q):\mathfrak g\longrightarrow E.
 \label{eq:t10v2-gauge-generator}
\end{equation}
Assume the action $S$ is invariant, so
\begin{equation}
 \langle \nabla S(q),D(q)\xi\rangle=0
 \qquad(\xi\in\mathfrak g).
 \label{eq:t10v2-first-Ward}
\end{equation}
Choose a connection one-form
$\Omega(q):E\to\mathfrak g$ satisfying
\begin{equation}
 \Omega(q)D(q)=I_{\mathfrak g}
 \label{eq:t10v2-connection-right-inverse}
\end{equation}
on the based gauge algebra after stabilizers have been retained separately.
Put
\begin{equation}
 P:=D\Omega,
 \qquad H:=I-P.
 \label{eq:t10v2-horizontal-projector}
\end{equation}

\begin{proposition}[Exact block quotient Hessian]
\label{prop:t10v2-exact-quotient}
Let $K=D^2S(q)$ in the flat field-space connection.  The quotient Hessian is
the block operator
\begin{equation}
 K^{\rm q}=H^*KH=(I-P)^*K(I-P).
 \label{eq:t10v2-quotient-kernel}
\end{equation}
It satisfies
\begin{equation}
 K^{\rm q}D=0,
 \qquad D^*K^{\rm q}=0,
 \label{eq:t10v2-horizontal-annihilation}
\end{equation}
and, if $|D|_\mu\leq d_\mu$,
$|\Omega|_\mu\leq\omega_\mu$, and $|K|_\mu\leq k_\mu$, then
\begin{equation}
 |K^{\rm q}|_\mu
 \leq(1+d_\mu\omega_\mu)^2k_\mu.
 \label{eq:t10v2-quotient-locality}
\end{equation}
\end{proposition}

\begin{proof}
Equation \eqref{eq:t10v2-quotient-kernel} is the definition
$D^2S(H\cdot,H\cdot)$ from V1.  Because $HD=0$ by
\eqref{eq:t10v2-connection-right-inverse}, both annihilation identities
follow.  Lemma~\ref{lem:t10v2-block-algebra} gives
$|P|_\mu\leq d_\mu\omega_\mu$ and then the displayed locality estimate.
\end{proof}

\begin{remark}[No smallness of the projector]
\label{rem:t10v2-projector-not-small}
The factor $1+d_\mu\omega_\mu$ need only be finite; the projector is not a
perturbation.  This is the block-local analogue of the matched energy
isometry in ninth-series V95 and of the exact conditional lift in the YM
audit.  Smallness is required later for the difference between physical
precisions, not for forming the quotient.
\end{remark}

\subsection{The connection-gradient operator in blocks}

Differentiate $D(q)$ in a field direction $v$ and write
$D'D(q)[v]\xi$.  Differentiating
\eqref{eq:t10v2-first-Ward} gives the flat geometric Ward identity
\begin{equation}
 D^2S(q)[v,D\xi]
 =-\langle\nabla S(q),D'D(q)[v]\xi\rangle.
 \label{eq:t10v2-differentiated-Ward}
\end{equation}

Define the bilinear connection-gradient form
\begin{align}
 \Gamma_g(v,w):={}&
 \langle g,D'D[Hv]\,\Omega w\rangle
 +\langle g,D'D[Hw]\,\Omega v\rangle
 \notag\\
 &+\langle g,D'D[D\Omega v]\,\Omega w\rangle,
 \qquad g:=\nabla S(q).
 \label{eq:t10v2-Gamma-form}
\end{align}

\begin{theorem}[Exact quotient correction in a local frame]
\label{thm:t10v2-exact-correction}
The quotient Hessian satisfies
\begin{equation}
 K^{\rm q}(v,w)=K(v,w)+\Gamma_g(v,w).
 \label{eq:t10v2-exact-correction}
\end{equation}
Although the last summand in \eqref{eq:t10v2-Gamma-form} is not manifestly
symmetric, the complete form $\Gamma_g$ is symmetric.
\end{theorem}

\begin{proof}
Expand $K(Hv,Hw)$ after writing
$v=Hv+D\Omega v$ and $w=Hw+D\Omega w$.  Apply
\eqref{eq:t10v2-differentiated-Ward} to the two mixed terms and to the
vertical--vertical term, then rearrange.  Symmetry follows because both
$K^{\rm q}$ and $K$ are symmetric.  Equivalently, the antisymmetric part of
the last summand pairs $g$ with a gauge commutator and vanishes by
\eqref{eq:t10v2-first-Ward}.
\end{proof}

To state a locality bound, let $\mathcal B_g:E\times\mathfrak g\to\mathbb R$
be
\begin{equation}
 \mathcal B_g(v,\xi):=langle g,D'D[v]\xi\rangle
 \label{eq:t10v2-Bg}
\end{equation}
and let $b_\mu(g)$ denote its weighted bilinear block norm.

\begin{proposition}[Connection-gradient locality stock]
\label{prop:t10v2-Gamma-stock}
Under the block bounds of
Proposition~\ref{prop:t10v2-exact-quotient},
\begin{equation}
 |\Gamma_g|_\mu
 \leq b_\mu(g)
 \left[2(1+d_\mu\omega_\mu)\omega_\mu
       +d_\mu\omega_\mu^2\right].
 \label{eq:t10v2-Gamma-stock}
\end{equation}
If $D'$ is finite range and the local gradient norm satisfies
$b_\mu(g)\leq c_{D,\mu}\|g\|_{\rm loc}$, the entire nonstationary debit is
linear in the local action gradient.
\end{proposition}

\begin{proof}
Apply the bilinear weighted norm to the three terms of
\eqref{eq:t10v2-Gamma-form}.  The first two cost respectively
$b_\mu(g)|H|_\mu|\Omega|_\mu$, and the third costs
$b_\mu(g)|D|_\mu|\Omega|_\mu^2$.  Use
$|H|_\mu\leq1+d_\mu\omega_\mu$.
\end{proof}

\begin{corollary}[Moving-reference assignment]
\label{cor:t10v2-moving-reference}
At a critical calibrated background, $g=0$ and
$K^{\rm q}=K$ on horizontal inputs.  At a moving noncritical background,
the exact term $\Gamma_g$ is a one-point connection response and may be
assigned to the moving-reference ledger before the stationary covariance
remainder is estimated.  Its assigned norm is the right-hand side of
\eqref{eq:t10v2-Gamma-stock}.
\end{corollary}

\begin{proof}
Use Theorem~\ref{thm:t10v2-exact-correction} and
Proposition~\ref{prop:t10v2-Gamma-stock}.
\end{proof}

\subsection{Schur localization after half-space elimination}

Split the horizontal variables into retained boundary variables and
interior variables.  Write a positive horizontal precision as
\begin{equation}
 \mathcal K=
 \begin{pmatrix}
 A&J\\ J^*&D_B
 \end{pmatrix},
 \qquad A>0,
 \label{eq:t10v2-boundary-block}
\end{equation}
where $A$ acts on the interior.  Its boundary Schur precision is
\begin{equation}
 S_B=D_B-J^*A^{-1}J.
 \label{eq:t10v2-boundary-Schur}
\end{equation}

\begin{proposition}[Weighted Schur locality]
\label{prop:t10v2-Schur-locality}
If
\begin{equation}
 |A^{-1}|_\mu\leq a_\mu,
 \qquad |J|_\mu\leq j_\mu,
 \qquad |D_B|_\mu\leq d_{B,\mu},
 \label{eq:t10v2-Schur-stocks}
\end{equation}
then
\begin{equation}
 |S_B|_\mu\leq d_{B,\mu}+j_\mu^2a_\mu.
 \label{eq:t10v2-Schur-locality}
\end{equation}
Moreover, if all tilded blocks are perturbed and
$\Delta A^{-1}=\widetilde A^{-1}-A^{-1}$, then the exact identity
\begin{align}
 \widetilde S_B-S_B={}&
 \Delta D_B-(\Delta J)^*A^{-1}J
 -\widetilde J^*A^{-1}\Delta J
 -\widetilde J^*\Delta A^{-1}\widetilde J
 \label{eq:t10v2-Schur-change}
\end{align}
gives a weighted bound by the corresponding product stocks.
\end{proposition}

\begin{proof}
The first claim is Lemma~\ref{lem:t10v2-block-algebra}.  For the second,
add and subtract $\widetilde J^*A^{-1}\widetilde J$ and then expand the two
$J$ differences.  Direct multiplication verifies
\eqref{eq:t10v2-Schur-change}; apply the block algebra term by term.
\end{proof}

\begin{remark}[Interior gap versus physical closing modes]
\label{rem:t10v2-interior-gap}
The required $A^{-1}$ is an interior Dirichlet or based-gauge inverse at
fixed boundary data.  It may have a uniform local gap even though the
physical boundary precision closes like $|k|$.  Closing boundary and global
gauge modes remain outside $A^{-1}$.  This separation is the analytic reason
for retaining the reflection boundary in ninth-series V96.
\end{remark}

\subsection{Diffeomorphism Ward identity forces the zero symbol}

At a flat translation-invariant background, the principal symbol of an
infinitesimal diffeomorphism acting on a symmetric metric perturbation is
\begin{equation}
 D_0(k)\xi=\mathrm i(k\otimes\xi+\xi\otimes k).
 \label{eq:t10v2-diffeomorphism-symbol}
\end{equation}

\begin{lemma}[Symmetric tensors are generated by gauge slopes]
\label{lem:t10v2-symmetric-span}
The set
\begin{equation}
 \{v\otimes\xi+\xi\otimes v:v,\xi\in\mathbb R^d\}
 \label{eq:t10v2-symmetric-generators}
\end{equation}
spans $\operatorname{Sym}^2(\mathbb R^d)$.
\end{lemma}

\begin{proof}
Taking $v=\xi=e_i$ gives $2e_i\otimes e_i$.  Taking
$v=e_i$, $\xi=e_j$ for $i\neq j$ gives the standard off-diagonal symmetric
basis element.
\end{proof}

\begin{theorem}[Diffeomorphism zero-symbol theorem]
\label{thm:t10v2-zero-symbol}
Let $L(k)$ be a continuous operator on symmetric tensors near $k=0$ and
suppose
\begin{equation}
 L(k)D_0(k)=0
 \qquad(k\neq0).
 \label{eq:t10v2-Ward-annihilation}
\end{equation}
Then
\begin{equation}
 L(0)=0.
 \label{eq:t10v2-zero-symbol}
\end{equation}
If also $L(-k)=L(k)$ and its real-space kernel has finite second moment,
then $L(k)=O(|k|^2)$ and the relative boundary estimate of tenth-series V1
applies against every positive order-one precision
$C_0(k)\succeq c|k|$.
\end{theorem}

\begin{proof}
Set $k=tv$ in \eqref{eq:t10v2-Ward-annihilation}, divide by $t$, and let
$t\to0$.  Continuity and
\eqref{eq:t10v2-diffeomorphism-symbol} give
\[
 L(0)(v\otimes\xi+\xi\otimes v)=0
\]
for every $v,\xi$.  Lemma~\ref{lem:t10v2-symmetric-span} proves
\eqref{eq:t10v2-zero-symbol}.  Evenness removes the first derivative at
zero; the finite second moment bounds the second derivative of the symbol.
The V1 boundary moment theorem then gives the relative estimate.
\end{proof}

\begin{corollary}[Stationary quotient boundary gate]
\label{cor:t10v2-stationary-boundary-gate}
Suppose an equivariant half-space elimination preserves
\eqref{eq:t10v2-horizontal-annihilation} in the boundary Schur remainder
$L$, and suppose its reflected stationary kernel is even and has weighted
norm $|L|_\mu<\infty$.  Then
\begin{equation}
 \|C_0(k)^{-1/2}L(k)C_0(k)^{-1/2}\|
 \leq\frac{|k|}{c\mu^2}|L|_\mu
 \label{eq:t10v2-relative-boundary-gate}
\end{equation}
up to the declared metric normalization.  If the supremum on the retained
carrier is below one, the inverse and Schur comparison bounds of
ninth-series V98 follow.
\end{corollary}

\begin{proof}
The exponential weighted norm gives second moment at most
$2\mu^{-2}|L|_\mu$.  Insert this in the V1 boundary moment theorem and use
Theorem~\ref{thm:t10v2-zero-symbol}.  The V98 relative Schur theorem applies
when the normalized bound is below one.
\end{proof}

\begin{firewall}[Schur Ward inheritance]
\label{fire:t10v2-Schur-Ward}
The boundary Schur remainder inherits gauge annihilation only if the
interior elimination, boundary conditions, coarea density, and retained
edge modes are equivariant.  An arbitrary Dirichlet gauge fixing can break
the identity at the cut.  This inheritance is a separate finite-regulator
theorem, not a consequence of the block norm estimate.
\end{firewall}

\subsection{The local-connection obstruction}

The estimates above expose the remaining issue.  A standard orthogonal
connection has the form
\begin{equation}
 \Omega=(D^*D)^{-1}D^*,
 \label{eq:t10v2-orthogonal-connection}
\end{equation}
which contains the inverse gauge Laplacian.

\begin{proposition}[Uniform-gap conditional locality]
\label{prop:t10v2-gap-locality}
Suppose on the based high-mode gauge space
$D^*D\succeq\gamma I$ with $\gamma>0$ independent of the regulator and
suppose
\begin{equation}
 D^*D=A_0(I-K),
 \qquad |A_0^{-1}K A_0|_\mu<1
 \label{eq:t10v2-local-parametrix}
\end{equation}
for a block-local $A_0^{-1}$.  Then
\eqref{eq:t10v2-orthogonal-connection} has a finite uniform weighted norm,
and all quotient and connection-gradient bounds above are uniform.
\end{proposition}

\begin{proof}
The local parametrix and Lemma~\ref{lem:t10v2-block-algebra} give a
weighted Neumann series for $(D^*D)^{-1}$.  Composition with the local
$D^*$ yields the connection bound.  Insert it in
Propositions~\ref{prop:t10v2-exact-quotient} and
\ref{prop:t10v2-Gamma-stock}.
\end{proof}

\begin{firewall}[Why the full gauge inverse is not used]
\label{fire:t10v2-full-gauge-inverse}
On expanding volumes the smallest nonzero eigenvalue of a gauge Laplacian
usually tends to zero.  Therefore the hypothesis of
Proposition~\ref{prop:t10v2-gap-locality} fails on the full gauge space, and
the orthogonal connection need not have a uniform exponential kernel.
The low gauge, Killing, and global modes must be retained in a finite or
slow sector before a local high-mode parametrix is formed.
\end{firewall}

\subsection{V2 conclusion and next target}

\begin{assessment}[What tenth-series V2 proves]
\label{assess:t10v2}
V2 proves the weighted block algebra, exact local quotient kernel,
connection-gradient stock, Schur locality identity, and the gravitational
diffeomorphism zero-symbol theorem.  Conditional on a local high-mode
connection and equivariant boundary elimination, a stationary reflected
remainder with finite weighted norm is relatively controlled against the
order-one massless boundary precision.  The active obstruction is the
uniformly local separation of slow gauge/Killing modes from the high-mode
connection.
\end{assessment}

\begin{programme}[Tenth-series V3: retained-mode parametrix]
\label{prog:t10v2-V3}
V3 will construct an abstract spectral split for the gauge Laplacian,
retain the slow projector as an explicit boundary/calibration variable, and
derive weighted resolvent and commutator bounds for the high-mode
connection.  It will state the rank and transport conditions needed to keep
the retained sector finite along the cofinal regulator.
\end{programme}

\begin{firewall}[Claim boundary after tenth-series V2]
\label{fire:t10v2-boundary}
V2 does not construct the slow-mode split for a concrete Einstein
regulator, prove equivariant boundary Schur inheritance, bound the curved
nonlinear remainder, construct the renormalized Standard Model stress
tensor, close determinant-line gluing, or construct the full SM--Einstein
continuum.
\end{firewall}

\paragraph{Parent-version record.}
V2 was formed from
\texttt{Renewal\_SM\_Einstein\_Continuum\_Research\_Notes\_V1.tex}.
The tenth-series V1 parent had $460$ validator-counted source lines,
$19\,612$ bytes, and SHA--256
\[
 \texttt{BBA77AE4305D86ACAD308F353554C7BB8430739B51108F0E485A4ECF6ABF34B0}.
\]

% =============================================================================
% END APPENDED TENTH-SERIES V2 MATERIAL
% =============================================================================

% =============================================================================
% BEGIN APPENDED TENTH-SERIES V3 MATERIAL
% =============================================================================

\section{Retained spectral modes, weighted parametrices, and an infrared no-go}
\label{sec:v10v3}

\subsection{Purpose and change of direction}

V2 reduced locality of the quotient Hessian to locality of a connection for
the gauge Laplacian.  The proposed V3 strategy was to retain a finite slow
spectral subspace and invert the complement with a uniform gap.  The first
part of this note proves the weighted functional calculus that such a split
would provide.  The second part proves that, for a massless Laplacian on
expanding tori, finite retained rank and a uniform complementary gap are
incompatible.  The programme therefore moves upstream to blockwise,
scale-local boundary variables rather than a single global slow sector.

\subsection{Weighted resolvent calculus}

Let $E=\bigoplus_{x\in\Lambda}E_x$ over a finite metric space.  For
$a\in\Lambda$ and $\mu\geq0$, define the diagonal weight
\begin{equation}
 (W_{a,\mu}u)_x=e^{\mu d(a,x)}u_x.
 \label{eq:t10v3-weight}
\end{equation}
Let $L=L^*$ and put
\begin{equation}
 \varepsilon_\mu(L)
 :=\sup_{a\in\Lambda}
 \|W_{a,\mu}LW_{a,\mu}^{-1}-L\|.
 \label{eq:t10v3-weighted-commutator}
\end{equation}

\begin{theorem}[Weighted resolvent bound]
\label{thm:t10v3-weighted-resolvent}
If $z\in\mathbb C$ satisfies
\begin{equation}
 \operatorname{dist}(z,\operatorname{spec}L)\geq\eta>\varepsilon_\mu(L),
 \label{eq:t10v3-resolvent-gap}
\end{equation}
then, for every $a$,
\begin{equation}
 \|W_{a,\mu}(z-L)^{-1}W_{a,\mu}^{-1}\|
 \leq\frac1{\eta-\varepsilon_\mu(L)}.
 \label{eq:t10v3-weighted-resolvent}
\end{equation}
Consequently, for block projections $1_x,1_y$,
\begin{equation}
 \|1_x(z-L)^{-1}1_y\|
 \leq\frac{e^{-\mu d(x,y)}}{\eta-\varepsilon_\mu(L)}.
 \label{eq:t10v3-resolvent-decay}
\end{equation}
\end{theorem}

\begin{proof}
Write
\[
 W_{a,\mu}(z-L)W_{a,\mu}^{-1}=z-L-E_a,
 \qquad \|E_a\|\leq\varepsilon_\mu(L).
\]
Since $L$ is self-adjoint,
$\|(z-L)^{-1}\|\leq\eta^{-1}$.  The factorization
\[
 z-L-E_a=[I-E_a(z-L)^{-1}](z-L)
\]
and the Neumann series give inverse norm at most
$[\eta(1-\varepsilon_\mu/\eta)]^{-1}$.  For the block bound choose
$a=x$ and observe that $W_{x,\mu}1_x=1_x$ while
$W_{x,\mu}^{-1}1_y=e^{-\mu d(x,y)}1_y$.
\end{proof}

\begin{remark}[Finite range supplies a small commutator]
\label{rem:t10v3-finite-range}
If $L_{xy}=0$ for $d(x,y)>R$ and its unweighted row and column sums are at
most $M$, then
\begin{equation}
 \varepsilon_\mu(L)\leq M(e^{\mu R}-1).
 \label{eq:t10v3-finite-range-commutator}
\end{equation}
Indeed, conjugation multiplies $L_{xy}$ by
$e^{\mu[d(a,x)-d(a,y)]}$, whose distance from one is at most
$e^{\mu R}-1$.  Thus a fixed spectral gap permits a fixed positive decay
rate.  When the gap closes, the admissible rate generally closes with it.
\end{remark}

\subsection{Locality of a separated spectral projector}

Suppose
\begin{equation}
 \operatorname{spec}L\subset[0,\lambda_-]\cup[\lambda_+,\Lambda],
 \qquad0\leq\lambda_-<\lambda_+.
 \label{eq:t10v3-cluster-gap}
\end{equation}
Let $P_-$ be the spectral projector onto $[0,\lambda_-]$ and
$Q=I-P_-$.

\begin{theorem}[Weighted spectral projector]
\label{thm:t10v3-weighted-projector}
Let $\Gamma$ be a positively oriented contour enclosing
$[0,\lambda_-]$ and no point of $[\lambda_+,\Lambda]$.  Put
\begin{equation}
 \eta_\Gamma:=
 \inf_{z\in\Gamma}\operatorname{dist}(z,\operatorname{spec}L).
 \label{eq:t10v3-contour-distance}
\end{equation}
If $\varepsilon_\mu(L)<\eta_\Gamma$, then
\begin{align}
 P_-&=\frac1{2\pi\mathrm i}\int_\Gamma(z-L)^{-1}dz,
 \label{eq:t10v3-Riesz-projector}\\
 \|1_xP_-1_y\|
 &\leq
 \frac{\operatorname{length}(\Gamma)}
 {2\pi[\eta_\Gamma-\varepsilon_\mu(L)]}
 e^{-\mu d(x,y)}.
 \label{eq:t10v3-projector-decay}
\end{align}
The same estimate, multiplied by
$\sup_{z\in\Gamma_+}|z^{-1}|$, holds for the high-mode inverse
\begin{equation}
 R_+:=Q(L|_{QE})^{-1}Q
 =\frac1{2\pi\mathrm i}
 \int_{\Gamma_+}z^{-1}(z-L)^{-1}dz,
 \label{eq:t10v3-high-inverse}
\end{equation}
where $\Gamma_+$ encloses the high cluster and excludes zero.
\end{theorem}

\begin{proof}
The Riesz formulas are the finite-dimensional holomorphic functional
calculus.  Apply \eqref{eq:t10v3-resolvent-decay} pointwise on each contour
and take its length outside the integral.
\end{proof}

\begin{corollary}[High-mode gauge connection]
\label{cor:t10v3-high-connection}
Let $L=D^*D$ on the gauge algebra and define
\begin{equation}
 \Omega_+:=R_+D^*,
 \qquad P_+:=D\Omega_+.
 \label{eq:t10v3-high-connection}
\end{equation}
Then
\begin{equation}
 \Omega_+D=Q,
 \qquad P_+D=DQ.
 \label{eq:t10v3-high-right-inverse}
\end{equation}
If $D$ is finite range and the hypotheses of
Theorem~\ref{thm:t10v3-weighted-projector} are uniform, then
$\Omega_+$ and $P_+$ have uniform exponential block bounds.  The retained
coordinate is $P_-\xi$ and is never inserted into $L^{-1}$.
\end{corollary}

\begin{proof}
Spectral projectors of $D^*D$ commute with $L$.  On $QE$,
$R_+L=Q$, which proves the algebraic identities.  Locality follows by
composition of the exponentially local $R_+$ with the finite-range $D^*$
and $D$.
\end{proof}

\subsection{Transport of the retained space}

\begin{proposition}[Riesz transport under a gap-preserving perturbation]
\label{prop:t10v3-projector-transport}
Let $\widetilde L=L+E$ be self-adjoint and suppose one contour $\Gamma$
separates the same slow and high clusters for both operators.  If their
distances from $\Gamma$ are at least $\eta>0$, then
\begin{equation}
 \widetilde P_--P_-
 =\frac1{2\pi\mathrm i}\int_\Gamma
 (z-\widetilde L)^{-1}E(z-L)^{-1}dz
 \label{eq:t10v3-projector-resolvent-identity}
\end{equation}
and
\begin{equation}
 \|\widetilde P_--P_-\|
 \leq\frac{\operatorname{length}(\Gamma)}{2\pi\eta^2}\|E\|.
 \label{eq:t10v3-projector-transport}
\end{equation}
Weighted versions follow if both resolvents satisfy the same weighted
contour bound.
\end{proposition}

\begin{proof}
Subtract the two Riesz formulas and use the resolvent identity
$(z-\widetilde L)^{-1}-(z-L)^{-1}
=(z-\widetilde L)^{-1}E(z-L)^{-1}$.  Bound each resolvent by $\eta^{-1}$.
\end{proof}

\begin{definition}[Global retained-mode card]
\label{def:t10v3-global-card}
A global retained-mode split would require, uniformly in the cofinal
regulator:
\begin{enumerate}
\item a cluster gap $\lambda_+-\lambda_->0$ admitting one contour with
uniform $\eta_\Gamma$;
\item a weight $\mu>0$ with
$\varepsilon_\mu(L)<\eta_\Gamma$;
\item a uniform rank bound $\operatorname{rank}P_-\leq r_*$;
\item summable or bounded projector transport under adjacent regulators;
\item inclusion of $P_-\xi$ among the finite calibration and boundary
coordinates.
\end{enumerate}
\end{definition}

\subsection{The expanding-torus obstruction}

The rank and gap requirements in
Definition~\ref{def:t10v3-global-card} cannot coexist for a massless
Laplacian on expanding volume.

\begin{theorem}[Finite-rank/uniform-gap no-go]
\label{thm:t10v3-torus-no-go}
Let $-\Delta_L$ be the scalar Laplacian on the flat torus
$\mathbb T_L^d=(\mathbb R/L\mathbb Z)^d$.  Its eigenvalues are
\begin{equation}
 \lambda_n(L)=\left(\frac{2\pi}{L}\right)^2|n|^2,
 \qquad n\in\mathbb Z^d.
 \label{eq:t10v3-torus-spectrum}
\end{equation}
The following statements hold.
\begin{enumerate}
\item For every fixed $\delta>0$, the rank of
$\mathbf1_{[0,\delta]}(-\Delta_L)$ tends to infinity at least as a positive
multiple of $L^d$.
\item If at most $r_*$ eigenvectors are retained, where $r_*$ is independent
of $L$, then the least eigenvalue on the complementary space satisfies
\begin{equation}
 \lambda_{\rm high}(L)\leq C(d,r_*)L^{-2}longrightarrow0.
 \label{eq:t10v3-high-gap-closes}
\end{equation}
\item The norm of the gradient right inverse
$(-\Delta_L)^{-1}\nabla^*$ on that complement is at least
$\lambda_{\rm high}(L)^{-1/2}$ and therefore grows at least linearly in
$L$ along a subsequence.
\end{enumerate}
Thus a uniformly finite retained sector cannot leave a uniformly gapped,
uniformly local global gauge connection.
\end{theorem}

\begin{proof}
For the first assertion, every integer vector satisfying
$|n_j|\leq L\sqrt\delta/(2\pi\sqrt d)$ has eigenvalue at most $\delta$.
The number of such vectors grows as a positive multiple of $L^d$.

For the second, order the fixed integer vectors by $|n|$.  Among any
$r_*+1$ fixed distinct vectors, at least one eigenvector is not retained.
Its eigenvalue is at most
$(2\pi/L)^2N(d,r_*)^2$ for a fixed integer radius $N(d,r_*)$, proving
\eqref{eq:t10v3-high-gap-closes}.

On an eigenfunction with eigenvalue $\lambda>0$, the gradient has norm
$\sqrt\lambda$ and $(-\Delta_L)^{-1}\nabla^*$ has singular value
$\lambda^{-1/2}$.  Evaluate it at the least complementary eigenvalue.
\end{proof}

\begin{corollary}[Sharp-cutoff repair does not work]
\label{cor:t10v3-threshold-no-go}
Choosing a fixed slow threshold gives an extensive retained rank.  Lowering
the threshold to keep the rank finite forces the high-mode gap and the
Combes--Thomas decay rate to zero.  Therefore no choice of one global
spectral threshold passes
Definition~\ref{def:t10v3-global-card} on expanding massless tori.
\end{corollary}

\begin{proof}
This is the first two assertions of
Theorem~\ref{thm:t10v3-torus-no-go}, combined with
Theorems~\ref{thm:t10v3-weighted-resolvent} and
\ref{thm:t10v3-weighted-projector}: their admissible weight requires the
weighted commutator to stay below the spectral separation.
\end{proof}

\begin{remark}[The obstruction is physical, not a defect of the proof]
\label{rem:t10v3-physical-massless}
The closing spectrum is the expected massless long-wavelength sector.  A
uniform mass gap may not be imposed on photon, graviton, or gauge-parameter
coordinates merely to make the connection local.  The correct construction
must distribute slow information over scale and location.
\end{remark}

\subsection{Replacement by a scale-local retained sector}

\begin{definition}[Blockwise retained-mode principle]
\label{def:t10v3-blockwise-principle}
At each renormalization scale, partition spacetime into blocks of fixed
diameter in that scale's rescaled metric.  On each block:
\begin{enumerate}
\item gauge transformations vanishing on the block boundary form the
interior high sector;
\item boundary gauge frames, averages, Killing candidates, and constraint
fluxes are retained as interface variables;
\item the interior Dirichlet or rooted gauge Laplacian is inverted only
inside the block;
\item reflected blocks share the full V96 boundary layer;
\item coarse graining transports the interface variables to the next scale
rather than accumulating them in one global finite-dimensional sector.
\end{enumerate}
The number of retained variables may be extensive, but their owner
multiplicity and support are uniformly local.
\end{definition}

\begin{proposition}[Why the block principle evades the no-go]
\label{prop:t10v3-block-escape}
Theorem~\ref{thm:t10v3-torus-no-go} does not preclude
Definition~\ref{def:t10v3-blockwise-principle}: the retained rank is no
longer globally bounded, and the inverted operator is a block Dirichlet
operator whose rescaled diameter is fixed.  The relevant requirement is a
uniform retained-variable count per block and bounded owner overlap, not a
uniform total rank.
\end{proposition}

\begin{proof}
The no-go assumes one global complementary inverse and a rank bound
independent of volume.  A direct sum of fixed-rescaled-size block inverses
has neither property.  Its row and column owner counts are local, so an
extensive number of blocks does not enlarge a weighted supremum norm.
\end{proof}

\begin{firewall}[A principle is not yet a construction]
\label{fire:t10v3-block-principle}
The block principle still requires an exact equivariant decomposition,
compatible interface constraints, a scale-covariant normalization, and
uniform block Green estimates.  Boundary variables cannot be discarded at
the next scale; doing so would recreate the global inverse.  These points
are the subject of V4.
\end{firewall}

\subsection{V3 conclusion}

\begin{assessment}[What tenth-series V3 proves]
\label{assess:t10v3}
V3 proves weighted resolvent, spectral-projector, high-inverse, and
projector-transport bounds.  It also proves that the originally proposed
finite retained slow sector is incompatible with a uniform complementary
gap for a massless Laplacian on expanding tori.  This is a genuine no-go,
so the programme replaces the global spectral split with blockwise retained
interface variables.
\end{assessment}

\begin{programme}[Tenth-series V4: block Dirichlet connection]
\label{prog:t10v3-V4}
V4 will construct the finite block decomposition, prove a discrete
Dirichlet Poincare bound and the scaled weighted inverse estimate, retain
boundary gauge frames explicitly, and show that the resulting interior
quotient Hessian and reflected Schur factorization have uniform owner norms.
\end{programme}

\begin{firewall}[Claim boundary after tenth-series V3]
\label{fire:t10v3-boundary}
V3 does not construct the block interface system for a concrete Einstein
regulator, prove its constraint and coarea compatibility, control the
curved nonlinear connection, renormalize the Standard Model stress tensor,
close the determinant line, or construct the full SM--Einstein continuum.
\end{firewall}

\paragraph{Parent-version record.}
V3 was formed from
\texttt{Renewal\_SM\_Einstein\_Continuum\_Research\_Notes\_V2.tex}.
The V2 parent had $909$ validator-counted source lines, $35\,284$ bytes, and
SHA--256
\[
 \texttt{88866E52AFE22D2CE5934F9D77312552C2C90124C4EAB7E6E6359E51492E109D}.
\]

% =============================================================================
% END APPENDED TENTH-SERIES V3 MATERIAL
% =============================================================================

% =============================================================================
% BEGIN APPENDED TENTH-SERIES V4 MATERIAL
% =============================================================================

\section{Fixed-scale block connections with retained interfaces}
\label{sec:v10v4}

\subsection{Block geometry and dimensionless normalization}

V3 proved that a finite global slow-mode sector cannot leave a uniform
massless spectral gap on expanding volume.  V4 constructs the replacement:
invert only gauge transformations vanishing on fixed-rescaled-size blocks
and retain their boundary traces.  The total number of boundary variables
is extensive, but every estimate is a row/column or owner supremum and is
therefore independent of the number of blocks.

Fix an integer block side $L\geq2$, independent of the regulator scale.
At scale $n$, with mesh $a_n$, a block contains $L^d$ sites.  If $D_n$ is a
first-difference gauge generator, use dimensionless operators
\begin{equation}
 \widehat D_n:=a_nD_n,
 \qquad
 \widehat\Delta_n:=\widehat D_n^*\widehat D_n.
 \label{eq:t10v4-dimensionless-operators}
\end{equation}
All norms below are taken in these rescaled variables.  Physical powers of
$a_n$ are restored only when pairing with dimensionful fields and sources.

\subsection{A discrete Dirichlet gap}

Let
$B_L=\{1,\ldots,L\}^d$ and extend a scalar function $u$ by zero on the face
$\{0\}\times\{1,\ldots,L\}^{d-1}$.  Let $\nabla_1u(j,x')=u(j,x')-u(j-1,x')$.

\begin{lemma}[One-face discrete Poincare inequality]
\label{lem:t10v4-Poincare}
For every such $u$,
\begin{equation}
 \sum_{x\in B_L}|u(x)|^2
 \leq L^2
 \sum_{x'\in\{1,\ldots,L\}^{d-1}}
 \sum_{j=1}^L|\nabla_1u(j,x')|^2.
 \label{eq:t10v4-Poincare}
\end{equation}
Consequently the full Dirichlet graph Laplacian on $B_L$ satisfies
\begin{equation}
 \widehat\Delta_{B_L}^{D}\succeq L^{-2}I.
 \label{eq:t10v4-Dirichlet-gap}
\end{equation}
The same statement holds componentwise for a finite-dimensional compact
Lie algebra.
\end{lemma}

\begin{proof}
For $x=(j,x')$,
\[
 u(j,x')=\sum_{r=1}^j\nabla_1u(r,x').
\]
Cauchy--Schwarz gives
$|u(j,x')|^2\leq j\sum_{r=1}^j|\nabla_1u(r,x')|^2
\leq L\sum_{r=1}^L|\nabla_1u(r,x')|^2$.
Sum first over $j$ and then $x'$.  The full graph energy dominates the
first-coordinate energy, proving the spectral bound.  Tensoring with a
finite-dimensional inner-product space changes no constant.
\end{proof}

\begin{corollary}[Fixed-block inverse]
\label{cor:t10v4-block-inverse}
The dimensionless Dirichlet inverse satisfies
\begin{equation}
 \|(widehat\Delta_{B_L}^{D})^{-1}\|_{2\to2}\leq L^2.
 \label{eq:t10v4-block-inverse}
\end{equation}
If extended by zero outside $B_L$, its kernel has range at most
$\operatorname{diam}B_L$ and, for every fixed $\mu\geq0$,
\begin{equation}
 |(\widehat\Delta_{B_L}^{D})^{-1}|_\mu
 \leq L^{d+2}e^{\mu\operatorname{diam}B_L}.
 \label{eq:t10v4-block-weighted-inverse}
\end{equation}
This constant is independent of $n$ and of the total volume.
\end{corollary}

\begin{proof}
The spectral theorem and \eqref{eq:t10v4-Dirichlet-gap} give the operator
norm.  Every block entry of a matrix is bounded by its operator norm, every
row and column contains at most $L^d$ nonzero blocks, and the exponential
weight is at most $e^{\mu\operatorname{diam}B_L}$ within the block.
\end{proof}

\begin{remark}[Why no global Combes--Thomas estimate is needed]
\label{rem:t10v4-no-global-CT}
The inverse has exact finite support because it is never extended across a
block interface.  Its crude constant may depend on the fixed integer $L$;
RG uniformity requires $L$ to be fixed after rescaling, not the impossible
volume-independent global gap ruled out in V3.
\end{remark}

\subsection{Exact interior/boundary gauge decomposition}

Let $\mathfrak g(B)$ be the gauge-parameter space on a block and let
$\operatorname{tr}_B:\mathfrak g(B)\to\mathfrak g(\partial B)$ be boundary
restriction.  Choose a linear extension
$E_B:\mathfrak g(\partial B)\to\mathfrak g(B)$ satisfying
$\operatorname{tr}_BE_B=I$.  Put
\begin{equation}
 \mathfrak g_0(B):=\ker\operatorname{tr}_B.
 \label{eq:t10v4-interior-gauge-algebra}
\end{equation}

\begin{lemma}[Trace exact sequence and retained boundary frame]
\label{lem:t10v4-trace-split}
Every gauge parameter has the unique decomposition
\begin{equation}
 \xi=\xi_0+E_B\xi_\partial,
 \qquad
 \xi_\partial=\operatorname{tr}_B\xi,
 \qquad
 \xi_0\in\mathfrak g_0(B).
 \label{eq:t10v4-gauge-split}
\end{equation}
Thus
\begin{equation}
 0\longrightarrow\mathfrak g_0(B)
 \longrightarrow\mathfrak g(B)
 \overset{\operatorname{tr}_B}{\longrightarrow}
 \mathfrak g(\partial B)\longrightarrow0
 \label{eq:t10v4-trace-exact-sequence}
\end{equation}
is split exact.  The coordinate $\xi_\partial$ must be retained; only
$\xi_0$ is inverted locally.
\end{lemma}

\begin{proof}
Set
$\xi_0=\xi-E_B\operatorname{tr}_B\xi$.  Its trace is zero and the displayed
decomposition follows.  Applying the trace proves uniqueness.
\end{proof}

Let $D_{B,0}(q)$ be the infinitesimal gauge generator restricted to
$\mathfrak g_0(B)$ and assume its dimensionless quadratic form is uniformly
comparable to the Dirichlet graph energy:
\begin{equation}
 \langle\xi,D_{B,0}^*D_{B,0}\xi\rangle
 \geq c_D\langle\xi,\widehat\Delta_{B_L}^{D}\xi\rangle,
 \qquad c_D>0.
 \label{eq:t10v4-gauge-ellipticity}
\end{equation}

\begin{theorem}[Interior block connection]
\label{thm:t10v4-block-connection}
Define
\begin{equation}
 \Omega_{B,0}:=(D_{B,0}^*D_{B,0})^{-1}D_{B,0}^*,
 \qquad
 P_{B,0}:=D_{B,0}\Omega_{B,0}.
 \label{eq:t10v4-block-connection}
\end{equation}
Then
\begin{equation}
 \Omega_{B,0}D_{B,0}=I_{\mathfrak g_0(B)},
 \qquad P_{B,0}^2=P_{B,0}=P_{B,0}^*,
 \label{eq:t10v4-block-projector}
\end{equation}
and both kernels vanish outside $B\times B$.  For fixed $L$, finite-range
$D_{B,0}$, and uniform $c_D$, their weighted block norms are bounded by a
constant $C_{\Omega}(L,c_D,\mu)$ independent of scale, volume, and block
location.
\end{theorem}

\begin{proof}
The Poincare bound and \eqref{eq:t10v4-gauge-ellipticity} make
$D_{B,0}^*D_{B,0}$ invertible with inverse norm at most $L^2/c_D$.
The right-inverse and projector identities are direct multiplication; the
last equality is the standard formula for the orthogonal projector onto
$\operatorname{Ran}D_{B,0}$.  All factors act inside one block.  The same
finite matrix row-count estimate as in
Corollary~\ref{cor:t10v4-block-inverse}, followed by the V2 block algebra,
gives the uniform weighted constants.
\end{proof}

\begin{firewall}[Boundary transformations are not fixed twice]
\label{fire:t10v4-boundary-gauge}
Two adjacent blocks share one interface gauge frame.  Introducing an
independent copy for each block and then integrating both would duplicate a
gauge orbit and its determinant.  The global retained space is the fibre
product of block traces under equality on shared faces; based interior
groups act independently only after that shared trace is fixed.
\end{firewall}

\subsection{Global local norm from block owners}

Let $\mathcal B$ be a block cover with interface layers of fixed thickness.
Assume every site and interaction cell belongs to at most $N_*$ enlarged
blocks.  Extend each block operator $T_B$ by zero and set
\begin{equation}
 T=\sum_{B\in\mathcal B}T_B.
 \label{eq:t10v4-owner-sum}
\end{equation}

\begin{proposition}[Bounded-overlap owner estimate]
\label{prop:t10v4-owner-estimate}
If $|T_B|_\mu\leq C_T$ for every block and a fixed row or column receives
contributions from at most $N_*$ block owners, then
\begin{equation}
 |T|_\mu\leq N_*C_T.
 \label{eq:t10v4-owner-estimate}
\end{equation}
This bound is independent of the number of blocks.
\end{proposition}

\begin{proof}
For each row or column, apply the triangle inequality and sum over the at
most $N_*$ nonzero owners before taking the supremum.
\end{proof}

\begin{corollary}[Uniform interior quotient Hessian]
\label{cor:t10v4-global-quotient}
Let $P_0=\bigoplus_BP_{B,0}$ for disjoint interiors and
$H_0=I-P_0$.  If the action Hessian has finite range and uniform weighted
norm, then
\begin{equation}
 K_0^{\rm q}=H_0^*KH_0
 \label{eq:t10v4-global-quotient}
\end{equation}
has a uniform weighted norm.  With fixed-thickness overlapping owners, the
same conclusion holds with the factor $N_*$ from
Proposition~\ref{prop:t10v4-owner-estimate}.
\end{corollary}

\begin{proof}
Apply the V2 exact quotient-kernel bound to the direct sum of block
projectors.  The block connection theorem makes their norms uniform; the
owner estimate handles overlap.
\end{proof}

\subsection{Conditional factorization at the interface}

Let $b$ collect all variables in the interaction-range boundary layer:
field traces, boundary gauge frame, constraint flux, and any fermion trace.
Let $u_B$ denote variables in the strict block interior.

\begin{theorem}[Finite-range conditional block product]
\label{thm:t10v4-conditional-product}
Suppose every action term is assigned to one owner and its support has
diameter at most $R$.  Choose the retained layer thick enough that deleting
it separates distinct strict interiors.  Then the bosonic reference density
and local action have the exact conditional form
\begin{equation}
 d\mu(u,b)=Z^{-1}d\nu_0(b)
 \prod_{B\in\mathcal B}
 \left[e^{-S_B(u_B;b)}d\rho_B(u_B\mid b)\right].
 \label{eq:t10v4-conditional-product}
\end{equation}
Interior based-gauge fixing contributes a product of positive determinants
$\prod_B\det(D_{B,0}^*D_{B,0})$; all cross-block gauge data remain in $b$.
\end{theorem}

\begin{proof}
After conditioning on the retained layer, no finite-range interaction meets
two strict interiors.  Group the uniquely owned action terms by their
interior block; boundary-only terms enter $d\nu_0$.  The same support
separation factors a local product reference density.  The based interior
gauge groups form a direct product at fixed shared boundary frame, so their
positive gauge-Laplacian Jacobians multiply.
\end{proof}

\begin{corollary}[Reflection-aligned blocks]
\label{cor:t10v4-reflection-blocks}
If the block cover and owner map are reflection invariant and the entire
reflection layer belongs to $b$, positive and negative strict interiors are
conditionally independent reflected copies.  Hence the V96 reflected-
square theorem gives exact bosonic OS positivity whenever the boundary
density after interior integration is nonnegative.
\end{corollary}

\begin{proof}
Pair each positive block with its reflected negative block in
\eqref{eq:t10v4-conditional-product}.  Their conditional kernels agree by
reflection covariance.  Apply the V96 square criterion.
\end{proof}

\begin{remark}[Extensive determinants are local densities]
\label{rem:t10v4-determinant-density}
The product of block determinants is extensive, but its logarithm is a sum
of locally owned terms.  Uniform continuum comparison must control the
centered local log-determinant density or its normalized $L^s$ weight, as in
ninth-series V87 and V92; a volume-independent bound on the entire product
is neither expected nor required.
\end{remark}

\subsection{Nonlinear stability inside a block}

Let $A_B=D_{B,0}^*D_{B,0}$ and perturb it to
$\widetilde A_B=A_B+E_B$.

\begin{proposition}[Block nonlinear inverse gate]
\label{prop:t10v4-nonlinear-block}
If
\begin{equation}
 \theta_B:=\|A_B^{-1/2}E_BA_B^{-1/2}\|\leq\theta<1
 \label{eq:t10v4-relative-block-perturbation}
\end{equation}
uniformly, then
\begin{align}
 (1-\theta)A_B&\preceq\widetilde A_B
 \preceq(1+\theta)A_B,
 \label{eq:t10v4-block-form-comparison}\\
 \|A_B^{1/2}\widetilde A_B^{-1}A_B^{1/2}\|
 &\leq(1-\theta)^{-1}.
 \label{eq:t10v4-block-inverse-comparison}
\end{align}
If additionally $|A_B^{-1}E_B|_\mu\leq\vartheta<1$, then
\begin{equation}
 |\widetilde A_B^{-1}|_\mu
 \leq\frac{|A_B^{-1}|_\mu}{1-\vartheta}.
 \label{eq:t10v4-weighted-nonlinear-inverse}
\end{equation}
\end{proposition}

\begin{proof}
Factor
$\widetilde A_B=A_B^{1/2}(I+K_B)A_B^{1/2}$ with
$\|K_B\|\leq\theta$ for the form bounds.  For the weighted bound use
$\widetilde A_B^{-1}=(I+A_B^{-1}E_B)^{-1}A_B^{-1}$ and the weighted Neumann
series from V2.
\end{proof}

\begin{remark}[Two norms have different jobs]
\label{rem:t10v4-two-norms}
The relative energy norm in
\eqref{eq:t10v4-relative-block-perturbation} preserves positivity and
matched covariance.  The weighted row/column norm in
\eqref{eq:t10v4-weighted-nonlinear-inverse} preserves locality and owner
summability.  Neither estimate should be inferred from the other without a
finite-block comparison constant.
\end{remark}

\subsection{Block-interface acceptance card}

\begin{definition}[Scale-local connection card]
\label{def:t10v4-scale-local-card}
A regulated gauge/gravity constraint system passes the scale-local
connection card if, at every RG scale:
\begin{enumerate}
\item blocks have one fixed side $L$ in rescaled units and bounded overlap;
\item the gauge and constraint generators on transformations vanishing at
the retained layer satisfy a uniform Dirichlet ellipticity bound;
\item all shared boundary frames, Killing candidates, averages, fluxes, and
closing modes occur once in a fibre-product interface space;
\item interior connections and nonlinear inverses satisfy the energy and
weighted bounds above;
\item every interaction and determinant term has one local owner;
\item the owner map and retained reflection layer are exactly reflection
covariant;
\item coarse transport maps the interface variables to the next scale with
uniform row/column multiplicity and does not integrate them through a
global inverse;
\item the coarea, ghost, and chiral boundary factors satisfy their separate
positivity and determinant-line cards.
\end{enumerate}
\end{definition}

\begin{theorem}[What the scale-local card supplies]
\label{thm:t10v4-scale-local-implication}
Under Definition~\ref{def:t10v4-scale-local-card}, the interior quotient
Hessian, connection-gradient debit, and interior Schur kernels have uniform
local owner norms.  The bosonic measure factors conditionally over blocks,
and a reflection-aligned cut has the exact V96 reflected-square form.
\end{theorem}

\begin{proof}
The block connection theorem, V2 quotient and gradient estimates, and the
owner proposition give the uniform kernels.  The finite-range conditional
product theorem gives factorization, and its reflection corollary gives the
square.
\end{proof}

\begin{assessment}[Status after tenth-series V4]
\label{assess:t10v4}
V4 replaces the impossible global finite-rank split by a rigorous
fixed-rescaled-block construction.  It proves the block Dirichlet gap,
finite-support connection, bounded-overlap owner estimate, conditional
product, reflection alignment, and nonlinear inverse gate.  The remaining
finite issue is equivariant gluing: after block interiors are integrated,
the boundary Schur action must inherit the exact gauge/diffeomorphism Ward
identity needed for the V2 zero-symbol theorem.
\end{assessment}

\begin{programme}[Tenth-series V5: companion audit and equivariant Schur gluing]
\label{prog:t10v4-V5}
V5 will first inspect the newest YM and NS notes.  It will then prove an
equivariant integration/minimization theorem for the boundary effective
action, differentiate its Ward identity, and identify the boundary
connection-gradient term and zero-symbol hypotheses without assuming that
gauge fixing preserves them automatically.
\end{programme}

\begin{firewall}[Claim boundary after tenth-series V4]
\label{fire:t10v4-boundary}
V4 does not prove equivariant boundary Ward inheritance for the complete
Einstein--SM measure, control the nonlinear curved constraint chart,
renormalize the stress tensor, establish chiral determinant-line gluing,
complete the RG trajectory, or construct the full SM--Einstein continuum.
\end{firewall}

\paragraph{Parent-version record.}
V4 was formed from
\texttt{Renewal\_SM\_Einstein\_Continuum\_Research\_Notes\_V3.tex}.
The V3 parent had $1\,279$ validator-counted source lines, $49\,580$ bytes,
and SHA--256
\[
 \texttt{0A0470EC104DEA4C437DE9AE2C9190283640DB90D7628E4ABFA7416B69C45855}.
\]

% =============================================================================
% END APPENDED TENTH-SERIES V4 MATERIAL
% =============================================================================

% =============================================================================
% BEGIN APPENDED TENTH-SERIES V5 MATERIAL
% =============================================================================

\section{V5 companion audit and equivariant boundary Ward inheritance}
\label{sec:v10v5}

\subsection{Compulsory companion audit}

The newest companion files in the shared notes directory were inspected
read-only before the present argument.  At the audit instant they were
\begin{align*}
 &\texttt{Renewal\_Yang\_Mills\_Mass\_Gap\_Research\_Notes\_V25.tex},
 &&399\,334\text{ bytes},\\
 &\texttt{Renewal\_NS\_Stability\_Research\_Notes\_V26.tex},
 &&278\,283\text{ bytes},
\end{align*}
with SHA--256 digests
\begin{align*}
 &\texttt{65D180F7C6B7F19BA4296971A44397E96107DB6B0916A76EC39C0295A8901064},\\
 &\texttt{82C887F8C2C69E4F28FAD7C3DC8DE5B9099CB5A4BF431EBA1FFF3E91935978AD},
\end{align*}
respectively.  The NS file is unchanged from the previous audit.

The relevant new YM results are these.  Its V23 stationary-covariance policy
keeps the Gaussian covariance at a stationary reference and assigns the
field-dependent nonstationary Ward row to a cubic-and-higher remainder.
V24 identifies that a finite output-head projection does not control
arbitrary input relative supports.  V25 repairs the gap with a stronger
weighted Hilbert space: an exponential norm at exponent $\nu$ embeds in an
owner $\ell^1$ norm only at a strictly smaller exponent $\mu<\nu$, or at the
same exponent with an additional polynomial Sobolev reserve.  V25 also
isolates a normalized block-realization gate and observes that a
direction-dependent zero-momentum symbol rules out exponential kernel
locality.

\begin{assessment}[V5 transfer decision]
\label{assess:t10v5-audit}
Three lessons are adopted.  The boundary precision will be evaluated at a
stationary calibrated reference; its connection-gradient variation remains
in the nonlinear ledger.  Generated interactions with unrestricted
relative supports will carry a strict exponent or polynomial reserve rather
than use the finite-range owner count of V4.  Finally, locality of a
normalized boundary conditional map will be tested at zero momentum before
an exponential kernel is claimed.  No Wilson constant, YM contraction
factor, NS kernel constant, mass-gap claim, or regularity claim is imported.
\end{assessment}

\subsection{Equivariant integration over block interiors}

Let $G$ be a compact Lie group acting smoothly on an interior manifold $X$
and a boundary manifold $B$.  Let $dx$ be a positive $G$-invariant measure
on $X$, and let $A\in C^2(X\times B)$ satisfy
\begin{equation}
 A(gx,gb)=A(x,b).
 \label{eq:t10v5-equivariant-action}
\end{equation}
Assume
\begin{equation}
 0<Z(b):=\int_Xe^{-A(x,b)}dx<\infty
 \label{eq:t10v5-boundary-partition}
\end{equation}
and that differentiation under the integral is justified by a common
integrable majorant through second order.  Define
\begin{equation}
 S_{\rm eff}(b):=-\log Z(b),
 \qquad
 d\mu_b(x):=Z(b)^{-1}e^{-A(x,b)}dx.
 \label{eq:t10v5-effective-action}
\end{equation}

\begin{theorem}[Equivariant integration preserves the boundary Ward identity]
\label{thm:t10v5-equivariant-integration}
Under the preceding hypotheses,
\begin{equation}
 S_{\rm eff}(gb)=S_{\rm eff}(b).
 \label{eq:t10v5-effective-invariance}
\end{equation}
For boundary tangent vectors $h,k$ represented in a fixed local chart,
\begin{align}
 DS_{\rm eff}(b)[h]
 &=\mathbb E_b[D_bA(x,b)[h]],
 \label{eq:t10v5-effective-gradient}\\
 D^2S_{\rm eff}(b)[h,k]
 &=\mathbb E_b[D_b^2A(x,b)[h,k]]
 -\operatorname{Cov}_b(D_bA[h],D_bA[k]).
 \label{eq:t10v5-effective-Hessian}
\end{align}
If $D_B(b)\xi$ is the boundary fundamental vector generated by
$\xi\in\mathfrak g$, then
\begin{equation}
 DS_{\rm eff}(b)[D_B(b)\xi]=0.
 \label{eq:t10v5-effective-Ward}
\end{equation}
\end{theorem}

\begin{proof}
Change variables $x=gy$ and use invariance of $dx$ and
\eqref{eq:t10v5-equivariant-action}:
\[
 Z(gb)=\int e^{-A(x,gb)}dx
 =\int e^{-A(gy,gb)}dy=Z(b).
\]
Differentiate $S_{\rm eff}=-\log Z$ under the integral for
\eqref{eq:t10v5-effective-gradient}.  A second derivative differentiates
both $D_bA$ and the normalized density; the latter contributes minus the
covariance, proving \eqref{eq:t10v5-effective-Hessian}.  Differentiate
\eqref{eq:t10v5-effective-invariance} along $e^{t\xi}$ for the Ward
identity.
\end{proof}

\begin{remark}[The fluctuation covariance is part of the precision]
\label{rem:t10v5-covariance-term}
The effective Hessian is not merely the expectation of the bare boundary
Hessian.  The covariance term in
\eqref{eq:t10v5-effective-Hessian} is negative semidefinite on a real
diagonal argument and can reduce a positivity margin.  In a Gaussian
completed square it is exactly the familiar Schur subtraction.  In the
interacting theory it requires its own relative precision estimate.
\end{remark}

\subsection{Equivariant minimization and the deterministic Schur form}

An analogous result holds for exact minimization.  Suppose for every $b$ a
unique interior critical point $x_*(b)$ satisfies
\begin{equation}
 D_xA(x_*(b),b)=0,
 \qquad
 H_{xx}:=D_x^2A(x_*(b),b)>0.
 \label{eq:t10v5-minimizer-hyp}
\end{equation}

\begin{theorem}[Equivariant minimizer and Schur Hessian]
\label{thm:t10v5-equivariant-minimization}
Under \eqref{eq:t10v5-equivariant-action} and
\eqref{eq:t10v5-minimizer-hyp},
\begin{equation}
 x_*(gb)=gx_*(b),
 \label{eq:t10v5-minimizer-equivariance}
\end{equation}
and $S_{\min}(b):=A(x_*(b),b)$ is $G$-invariant.  Its Hessian is
\begin{equation}
 D^2S_{\min}
 =H_{bb}-H_{bx}H_{xx}^{-1}H_{xb}.
 \label{eq:t10v5-minimizer-Schur}
\end{equation}
\end{theorem}

\begin{proof}
Equivariance sends an interior critical point over $b$ to one over $gb$;
uniqueness gives \eqref{eq:t10v5-minimizer-equivariance}.  Action
invariance then gives invariance of $S_{\min}$.  Differentiating the critical
equation yields
$Dx_*=-H_{xx}^{-1}H_{xb}$.  The first derivative of $S_{\min}$ has no
$D_xA$ term, and a second derivative gives
\eqref{eq:t10v5-minimizer-Schur}.
\end{proof}

\begin{remark}[Integration and minimization must not be conflated]
\label{rem:t10v5-integration-versus-minimization}
Equation \eqref{eq:t10v5-minimizer-Schur} is exact for a unique stationary
interior solve.  Equation \eqref{eq:t10v5-effective-Hessian} is exact for
integration.  Replacing the covariance term by the minimizer Schur term is
a semiclassical approximation unless an independent theorem controls their
difference.
\end{remark}

\subsection{Differentiated boundary Ward identity}

Equip $B$ with a $G$-invariant Riemannian metric and Levi--Civita connection
$\nabla^B$.

\begin{theorem}[Inherited differentiated Ward identity]
\label{thm:t10v5-differentiated-boundary-Ward}
For either $S_\partial=S_{\rm eff}$ or
$S_\partial=S_{\min}$ under the corresponding theorem above,
\begin{equation}
 \operatorname{Hess}S_\partial(b)
 [h,D_B(b)\xi]
 =-DS_\partial(b)
 [\nabla_h^B(D_B\xi)].
 \label{eq:t10v5-differentiated-boundary-Ward}
\end{equation}
At a critical boundary background $b_*$,
\begin{equation}
 \operatorname{Hess}S_\partial(b_*)D_B(b_*)=0.
 \label{eq:t10v5-critical-annihilation}
\end{equation}
At a noncritical background the exact connection-gradient correction from
tenth-series V1 makes the quotient Hessian horizontal.
\end{theorem}

\begin{proof}
Both effective actions are $G$-invariant.  Apply the differentiated
geometric Ward theorem of tenth-series V1 on $B$.  The critical and quotient
statements are its immediate consequences.
\end{proof}

\begin{corollary}[Boundary zero-symbol inheritance]
\label{cor:t10v5-zero-symbol}
Suppose a flat stationary boundary background is translation and reflection
invariant, and the effective quotient Hessian has a translation-invariant
kernel with finite second moment.  If its diffeomorphism generator has
principal symbol
$\mathrm i(k\otimes\xi+\xi\otimes k)$, then any stationary Hessian
remainder which preserves
\eqref{eq:t10v5-critical-annihilation} has zero symbol at $k=0$ and is
$O(|k|^2)$.  Relative to a positive order-one boundary precision it is
$O(|k|)$.
\end{corollary}

\begin{proof}
Apply the diffeomorphism zero-symbol theorem of tenth-series V2 and the
boundary moment theorem of tenth-series V1.
\end{proof}

\subsection{Gauge fixing, coarea factors, and the anomaly source}

The preceding theorem assumes an invariant positive interior measure.  A
gauge-fixed integral has that property only after the slice, coarea factor,
and retained boundary frame are combined correctly.

Let
\begin{equation}
 Z_J(b)=\int_Xe^{-A(x,b)}J(x,b)dx,
 \qquad J>0.
 \label{eq:t10v5-J-partition}
\end{equation}
For a total infinitesimal generator $V_\xi=(V^X_\xi,V^B_\xi)$, define its
density anomaly by
\begin{equation}
 \mathcal A_\xi
 :=\operatorname{div}_{dx}V^X_\xi+V_\xi\log J-V_\xi A.
 \label{eq:t10v5-density-anomaly}
\end{equation}

\begin{proposition}[Effective Ward anomaly]
\label{prop:t10v5-effective-anomaly}
Assume boundary terms in the interior divergence vanish.  Then
\begin{equation}
 D[-\log Z_J](b)[V^B_\xi(b)]
 =-\mathbb E_{J,b}[\mathcal A_\xi].
 \label{eq:t10v5-effective-anomaly}
\end{equation}
Thus exact boundary Ward inheritance holds if
$\mathcal A_\xi=0$ pointwise, and it holds in expectation if the right-hand
side vanishes.  A nonzero anomaly is an explicit source term rather than a
failure of the Hessian estimate.
\end{proposition}

\begin{proof}
Differentiate the integral under the simultaneous infinitesimal change of
$x$ and $b$.  The change of $dx$ supplies the divergence, the density
$J$ supplies $V_\xi\log J$, and the action supplies $-V_\xi A$.  Divide the
resulting variation of $Z_J$ by $-Z_J$.
\end{proof}

\begin{corollary}[Based block gauge fixing]
\label{cor:t10v5-based-gauge}
In the V4 block construction, suppose the interior based-gauge slice and
positive gauge-Laplacian coarea determinant are equivariant under the shared
boundary gauge frame.  Then their product measure satisfies
$\mathcal A_\xi=0$ for boundary gauge transformations, and the integrated
boundary action obeys
Theorem~\ref{thm:t10v5-differentiated-boundary-Ward}.
\end{corollary}

\begin{proof}
Equivariance of the full slice change of variables says its transformed
measure equals itself, which is precisely vanishing of
\mathcal A_\xi$.  Apply
Proposition~\ref{prop:t10v5-effective-anomaly} and then differentiate the
resulting Ward identity.
\end{proof}

\begin{firewall}[Chiral determinant lines]
\label{fire:t10v5-chiral-line}
A chiral determinant is not generally a positive scalar $J$.  Its local
gauge variation is the determinant-line connection, and global transport
may have holonomy.  The finite Standard Model anomaly arithmetic cancels
the local polynomial only after all representations are combined.  Exact
use of Corollary~\ref{cor:t10v5-based-gauge} in the chiral sector still
requires a compatible global trivialization and reflected boundary gluing.
\end{firewall}

\subsection{Stationary covariance policy for the boundary action}

Let $b_*$ be a critical calibrated reference for the effective boundary
action and set
\begin{equation}
 H_*:=D^2S_\partial(b_*).
 \label{eq:t10v5-stationary-precision}
\end{equation}
For $b=b_*+a$ in a local chart, define the nonlinear remainder
\begin{equation}
 \mathcal R(a):=S_\partial(b_*+a)-S_\partial(b_*)
 -\frac12\langle a,H_*a\rangle.
 \label{eq:t10v5-stationary-remainder}
\end{equation}

\begin{proposition}[Cubic order from a Hessian Lipschitz stock]
\label{prop:t10v5-cubic-remainder}
If
\begin{equation}
 \|D^2S_\partial(b_*+ta)-H_*\|
 \leq Lt\|a\|
 \qquad(0\leq t\leq1),
 \label{eq:t10v5-Hessian-Lipschitz}
\end{equation}
then
\begin{align}
 |\mathcal R(a)|&\leq\frac L6\|a\|^3,
 \label{eq:t10v5-cubic-action}\\
 \|D\mathcal R(a)\|&\leq\frac L2\|a\|^2,
 \label{eq:t10v5-quadratic-gradient}\\
 \|D^2\mathcal R(a)\|&\leq L\|a\|.
 \label{eq:t10v5-linear-Hessian}
\end{align}
The nonstationary connection-gradient Ward term is therefore quadratic in
the field in first derivative and belongs to the cubic action remainder,
provided the connection stock stays bounded.
\end{proposition}

\begin{proof}
Since $DS_\partial(b_*)=0$, integrate the Hessian difference once for
$D\mathcal R$ and twice for $\mathcal R$:
\[
 D\mathcal R(a)=\int_0^1[D^2S_\partial(b_*+ta)-H_*]a\,dt,
\]
and
\[
 \mathcal R(a)=\int_0^1(1-t)
 \langle a,[D^2S_\partial(b_*+ta)-H_*]a\rangle dt.
\]
The integrals of $t$ and $t(1-t)$ are $1/2$ and $1/6$.  The Hessian bound
at $t=1$ gives the last inequality.
\end{proof}

\begin{assessment}[Transfer from YM V23]
\label{assess:t10v5-stationary-policy}
Proposition~\ref{prop:t10v5-cubic-remainder} is the abstract boundary form
of the audited YM stationary-covariance policy.  The covariance is tied to
$H_*$, which has exact critical Ward annihilation.  Field-dependent Hessian
changes and the connection-gradient term are nonlinear vertices.  This
prevents a nonstationary gradient from being mistaken for a mass in the
stable Gaussian precision.
\end{assessment}

\subsection{Localization reserve and normalized-map diagnostic}

The V4 owner estimate applies to genuinely finite-range block kernels.  RG
interactions can have arbitrary relative support, so their coefficient
spaces require an additional reserve.

\begin{definition}[Boundary localization reserve]
\label{def:t10v5-localization-reserve}
For every generated degree $n$, choose either:
\begin{enumerate}
\item two exponential exponents $0\leq\mu<\nu$, using the stronger
$\ell^2_\nu$ norm for contraction and the weaker owner $\ell^1_\mu$ norm for
absolute sums; or
\item one exponential exponent $\mu$ and a polynomial weight of order
$s>D_n/2$ on the $D_n$-dimensional relative-support lattice.
\end{enumerate}
The finite block owner multiplicity is paid once after the coefficient
support sum, not once per tensor leg.
\end{definition}

The audited YM V25 proofs give the corresponding lattice embeddings by
Cauchy--Schwarz and the summability of the unused exponential or polynomial
weight.  They are not repeated here.  The present programme must still
construct the appropriate local orthogonal block realization for the
Einstein--SM boundary conditional map.

\begin{proposition}[Zero-momentum obstruction to exponential locality]
\label{prop:t10v5-zero-momentum-diagnostic}
Let a translation-invariant normalized boundary map have convolution
kernel $r(x)$ with
\begin{equation}
 \sum_xe^{\alpha|x|}\|r(x)\|<\infty
 \label{eq:t10v5-exponential-kernel}
\end{equation}
for some $\alpha>0$.  Then its Fourier symbol is continuous at $k=0$ and
has a direction-independent limit there.  Hence a direction-dependent
Hodge limit rules out \eqref{eq:t10v5-exponential-kernel}.
\end{proposition}

\begin{proof}
The exponential sum converges uniformly for complex momenta in every
strictly smaller strip.  The resulting analytic extension is continuous at
the origin, so all directional limits agree.
\end{proof}

\begin{remark}[Different objects have different infrared behavior]
\label{rem:t10v5-remainder-versus-normalized-map}
The V2/V5 Ward argument makes a stationary precision remainder vanish at
$k=0$.  It does not make the normalized conditional partial isometry or
Hodge projector continuous there.  The latter can be an order-zero angular
multiplier.  Exponential locality of these two objects must not be
conflated.
\end{remark}

\subsection{Equivariant boundary Schur card}

\begin{definition}[Boundary Ward-inheritance card]
\label{def:t10v5-Ward-card}
At every scale and fixed retained interface, the construction must provide:
\begin{enumerate}
\item an equivariant positive interior measure or an equivariant unique
minimizer;
\item differentiation majorants sufficient for
\eqref{eq:t10v5-effective-Hessian};
\item vanishing density anomaly after coarea, based gauge fixing, and all
nonchiral determinants are combined;
\item a critical calibrated boundary reference $b_*$ and a uniform Hessian
Lipschitz stock;
\item a finite second moment for the stationary quotient-Hessian remainder;
\item a separate determinant-line Ward and gluing theorem for chiral
fermions;
\item the localization reserve of
Definition~\ref{def:t10v5-localization-reserve};
\item a zero-momentum diagnostic for every normalized conditional map.
\end{enumerate}
\end{definition}

\begin{theorem}[Ward-inheritance implication]
\label{thm:t10v5-Ward-card-implication}
If the V4 scale-local connection card and
Definition~\ref{def:t10v5-Ward-card} hold, then the integrated or minimized
bosonic boundary action has an exactly horizontal stationary quotient
Hessian, its stationary local remainder obeys the V2 zero-symbol and
relative order-one precision bounds, and its nonstationary Ward row belongs
to the controlled nonlinear remainder.
\end{theorem}

\begin{proof}
Equivariant integration or minimization gives the boundary Ward identity.
The V1 quotient correction and
Theorem~\ref{thm:t10v5-differentiated-boundary-Ward} give exact
horizontality.  The finite second moment and reflection parity give the V2
zero-symbol estimate.  Proposition~\ref{prop:t10v5-cubic-remainder} assigns
the varying row to the nonlinear sector.
\end{proof}

\begin{assessment}[Status after tenth-series V5]
\label{assess:t10v5}
V5 proves that exact integration and exact minimization over equivariant
block interiors preserve the boundary Ward identity, with explicit Hessian
formulas.  It identifies the density anomaly source, proves the stationary
covariance/cubic remainder split, and records the localization reserve and
normalized-map infrared diagnostic transferred from the companion audit.
The active question is now whether the normalized gravitational boundary
map has a continuous zero-momentum symbol; if it has an angular Hodge limit,
the programme must use multiscale localization instead of exponential
locality.
\end{assessment}

\begin{programme}[Tenth-series V6: Hodge-symbol test and multiscale replacement]
\label{prog:t10v5-V6}
V6 will compute the flat normalized diffeomorphism/Hodge projector symbol,
prove its direction dependence at zero momentum, and replace the failed
exponential-kernel gate with a dyadic annular decomposition.  It will derive
scale-local kernel and owner bounds sufficient for the two-exponent reserve.
\end{programme}

\begin{firewall}[Claim boundary after tenth-series V5]
\label{fire:t10v5-boundary}
V5 does not prove locality of the normalized gravitational boundary map,
construct its multiscale replacement, establish chiral determinant-line
equivariance, bound the renormalized Standard Model stress tensor, close the
nonlinear RG, or construct the full SM--Einstein continuum.
\end{firewall}

\paragraph{Parent-version record.}
V5 was formed from
\texttt{Renewal\_SM\_Einstein\_Continuum\_Research\_Notes\_V4.tex}.
The V4 parent had $1\,698$ validator-counted source lines, $65\,826$ bytes,
and SHA--256
\[
 \texttt{734F5A814FE5CEB1F130EFD0C52E442243EBDE9AA041052AB8C034A3AA9D7D88}.
\]

% =============================================================================
% END APPENDED TENTH-SERIES V5 MATERIAL
% =============================================================================

% =============================================================================
% BEGIN APPENDED TENTH-SERIES V6 MATERIAL
% =============================================================================

\section{Angular Hodge symbols and a dyadic multiscale replacement}
\label{sec:v10v6}

\subsection{The flat diffeomorphism projector}

V5 isolated a zero-momentum diagnostic for normalized boundary maps.  The
stationary precision remainder vanishes quadratically by the Ward identity,
but the normalized gauge projector is an order-zero object.  This note
computes its symbol exactly and proves that it is direction dependent at
the origin.  Exponential convolution locality is therefore unavailable;
the correct replacement is a dyadic annular decomposition with uniform
rescaled kernel bounds.

Let $V=\mathbb R^d$ and let
$\operatorname{Sym}^2V$ carry the Frobenius inner product.  For a unit
vector $n\in V$, define
\begin{equation}
 D_n:V\longrightarrow\operatorname{Sym}^2V,
 \qquad
 D_n\xi=n\otimes\xi+\xi\otimes n.
 \label{eq:t10v6-Dn}
\end{equation}

\begin{lemma}[Adjoint and gauge Gram operator]
\label{lem:t10v6-Gram}
For $h\in\operatorname{Sym}^2V$,
\begin{align}
 D_n^*h&=2hn,
 \label{eq:t10v6-D-adjoint}\\
 D_n^*D_n&=2(I+n\otimes n),
 \label{eq:t10v6-D-Gram}\\
 (D_n^*D_n)^{-1}&=\frac12
 \left(I-\frac12n\otimes n\right).
 \label{eq:t10v6-Gram-inverse}
\end{align}
\end{lemma}

\begin{proof}
For symmetric $h$,
\[
 \langle D_n\xi,h\rangle
 =\langle n\otimes\xi+\xi\otimes n,h\rangle
 =2\langle\xi,hn\rangle,
\]
which proves the adjoint formula.  Apply it to $D_n\xi$:
$(D_n\xi)n=\xi+n\langle n,\xi\rangle$.  Finally,
$(I+n\otimes n)^{-1}=I-\frac12n\otimes n$.
\end{proof}

\begin{theorem}[Exact flat diffeomorphism Hodge projector]
\label{thm:t10v6-Hodge-projector}
The orthogonal projector onto $\operatorname{Ran}D_n$ is
\begin{equation}
 P(n):=D_n(D_n^*D_n)^{-1}D_n^*,
 \label{eq:t10v6-projector-definition}
\end{equation}
and acts by
\begin{align}
 P(n)h&=n\otimes v+v\otimes n,
 \label{eq:t10v6-projector-formula}\\
 v&=hn-\frac12n\langle n,hn\rangle.
 \label{eq:t10v6-v-formula}
\end{align}
It is homogeneous of degree zero as a function of nonzero momentum
$k=|k|n$ and has no direction-independent limit as $k\to0$.
\end{theorem}

\begin{proof}
Insert the three identities of Lemma~\ref{lem:t10v6-Gram} into
\eqref{eq:t10v6-projector-definition}; the vector entering $D_n$ is exactly
$v$ in \eqref{eq:t10v6-v-formula}.  For direction dependence, take
$h=e_1\otimes e_1$.  If $n=e_1$, then $v=e_1/2$ and $P(e_1)h=h$.  If
$n=e_2$, then $v=0$ and $P(e_2)h=0$.  Homogeneity follows because the
$|k|$ factors in $D(k)$ cancel those in its Gram inverse.
\end{proof}

\begin{corollary}[No exponentially summable full-space kernel]
\label{cor:t10v6-no-exponential-kernel}
There is no translation-invariant convolution kernel $p(x)$ representing
$P$ on nonzero modes and satisfying
\begin{equation}
 \sum_xe^{\alpha|x|}\|p(x)\|<\infty
 \label{eq:t10v6-forbidden-exponential}
\end{equation}
on a lattice, or the analogous exponentially weighted $L^1$ bound in the
continuum, for any $\alpha>0$.
\end{corollary}

\begin{proof}
Such a bound would make the Fourier symbol analytic in a strip and
continuous at $k=0$ by the V5 diagnostic.  This contradicts
Theorem~\ref{thm:t10v6-Hodge-projector}.
\end{proof}

\begin{remark}[The horizontal projector has the same obstruction]
\label{rem:t10v6-horizontal-obstruction}
$H(n)=I-P(n)$ is likewise direction dependent.  This does not invalidate
the exact quotient identity or its $L^2$ norm: $P(n)$ and $H(n)$ are
orthogonal projections of norm one for every $n$.  It invalidates only a
single-scale exponential kernel representation across all infrared modes.
\end{remark}

\subsection{Dyadic annular decomposition}

Choose $\chi\in C_c^\infty(\mathbb R^d\setminus\{0\})$ supported where
$1/2\leq|k|\leq2$ and satisfying
\begin{equation}
 \sum_{j\in\mathbb Z}\chi(2^{-j}k)=1
 \qquad(k\neq0).
 \label{eq:t10v6-dyadic-partition}
\end{equation}
Define
\begin{equation}
 P_j(k):=\chi(2^{-j}k)P(k/|k|).
 \label{eq:t10v6-dyadic-symbol}
\end{equation}
Let $p_j$ be its inverse Fourier kernel.

\begin{theorem}[Uniform rescaled Schwartz kernel]
\label{thm:t10v6-dyadic-kernel}
There is a Schwartz matrix kernel $p_0$ such that
\begin{equation}
 p_j(x)=2^{jd}p_0(2^jx).
 \label{eq:t10v6-kernel-scaling}
\end{equation}
For every multi-index $\alpha$ and $N\geq0$, there is
$C_{\alpha,N}<\infty$ with
\begin{equation}
 |\partial^\alpha p_j(x)|
 \leq C_{\alpha,N}2^{j(d+|\alpha|)}
 (1+2^j|x|)^{-N}.
 \label{eq:t10v6-Schwartz-bound}
\end{equation}
In particular, for every $m\geq0$,
\begin{align}
 \|p_j\|_{L^1}&\leq C_0,
 \label{eq:t10v6-L1-uniform}\\
 \int|x|^m|p_j(x)|dx&\leq C_m2^{-jm}.
 \label{eq:t10v6-moment-scaling}
\end{align}
All constants are independent of $j$.
\end{theorem}

\begin{proof}
The matrix symbol
$q(k)=\chi(k)P(k/|k|)$ is smooth and compactly supported away from zero, so
its inverse Fourier transform $p_0$ is Schwartz.  Changing variables
$k=2^j\kappa$ proves \eqref{eq:t10v6-kernel-scaling}.  Differentiate and
use the Schwartz bounds for $p_0$ to obtain
\eqref{eq:t10v6-Schwartz-bound}.  The $L^1$ and moment statements follow
from $y=2^jx$.
\end{proof}

\begin{corollary}[Exact reconstruction away from zero]
\label{cor:t10v6-reconstruction}
For every Schwartz tensor $h$ whose Fourier transform vanishes near zero,
\begin{equation}
 Ph=\sum_{j\in\mathbb Z}P_jh
 \label{eq:t10v6-projector-reconstruction}
\end{equation}
with only finitely many nonzero terms at each Fourier momentum and
convergence in every Sobolev norm.  For general $L^2$ data the sum converges
in $L^2$ after the zero mode is retained separately.
\end{corollary}

\begin{proof}
Equation \eqref{eq:t10v6-dyadic-partition} gives the pointwise symbol
identity away from zero.  Smooth compact frequency support gives Sobolev
convergence for the first class.  For $L^2$, use boundedness of the
orthogonal projector and standard finite partial sums of the smooth
partition, then pass by dominated convergence in Fourier space.
\end{proof}

\subsection{Square functions and owner-local kernels}

Assume the partition has overlap multiplicity at most $M_\chi$ and
$|\chi|\leq C_\chi$.

\begin{proposition}[Dyadic square-function bound]
\label{prop:t10v6-square-function}
For $h\in L^2(\mathbb R^d;\operatorname{Sym}^2V)$,
\begin{equation}
 \sum_j\|P_jh\|_2^2
 \leq M_\chi C_\chi^2\|h\|_2^2.
 \label{eq:t10v6-square-function}
\end{equation}
The same estimate holds for $H_j=\chi_j(I-P)$.
\end{proposition}

\begin{proof}
By Plancherel and $\|P(k)\|\leq1$,
\[
 \sum_j\|P_jh\|_2^2
 \leq\int\sum_j|\chi(2^{-j}k)|^2|\widehat h(k)|^2dk.
\]
At most $M_\chi$ terms occur and each is at most $C_\chi^2$.
\end{proof}

For the natural scale-$j$ distance
\begin{equation}
 d_j(x,y):=2^j|x-y|,
 \label{eq:t10v6-scaled-distance}
\end{equation}
Theorem~\ref{thm:t10v6-dyadic-kernel} says that $P_j$ is rapidly local in
$d_j$, uniformly in $j$.

\begin{proposition}[Uniform annular owner tail]
\label{prop:t10v6-owner-tail}
For every $N>d$ and $R\geq1$,
\begin{equation}
 \sup_j\sup_x
 \int_{d_j(x,y)>R}|p_j(x-y)|dy
 \leq C_N'R^{d-N}.
 \label{eq:t10v6-owner-tail}
\end{equation}
Thus truncating each annular kernel to scaled radius $R$ has a uniformly
summable tail tending to zero as $R\to\infty$.
\end{proposition}

\begin{proof}
Use \eqref{eq:t10v6-Schwartz-bound} with $\alpha=0$, change variables
$z=2^j(x-y)$, and integrate $(1+|z|)^{-N}$ over $|z|>R$.
\end{proof}

\subsection{Commutators with a slowly varying background}

The multiscale decomposition is useful only if curvature and calibration
variations do not destroy it.  Multiplication by a Lipschitz coefficient
$a(x)$ gives the first model.

\begin{theorem}[First commutator gain]
\label{thm:t10v6-commutator}
Let $M_a$ denote multiplication by a bounded Lipschitz scalar or matrix
coefficient.  Then
\begin{equation}
 \|[P_j,M_a]\|_{2\to2}
 \leq C_1,2^{-j}\|\nabla a\|_\infty,
 \label{eq:t10v6-commutator}
\end{equation}
where $C_1$ is the first-moment constant in
\eqref{eq:t10v6-moment-scaling}.  The same bound holds for $H_j$.
\end{theorem}

\begin{proof}
The commutator kernel is
\[
 p_j(x-y)[a(y)-a(x)].
\]
Its absolute row and column integrals are bounded by
$\|\nabla a\|_\infty\int|z||p_j(z)|dz
\leq C_12^{-j}\|\nabla a\|_\infty$.
The Schur test gives the $L^2$ bound.
\end{proof}

\begin{remark}[Low annuli see the background variation]
\label{rem:t10v6-low-annuli}
The factor $2^{-j}$ is small when the annular frequency is high compared
with the variation scale of $a$, but large in the infrared.  Low annuli
must therefore be transferred to coarser blocks and calibrated background
coordinates rather than treated perturbatively at the fine scale.  This is
the multiscale form of the V3 no-go.
\end{remark}

\subsection{A general scale-reserve summation}

The companion audit requires absolute owner sums without replacing an
unrestricted support set by a finite head.  The following elementary
device pays the sum over dyadic scales.

\begin{lemma}[Hilbert-to-owner scale reserve]
\label{lem:t10v6-scale-reserve}
Let $a_j\geq0$ be per-scale owner norms and suppose
$a_j\leq C_jh_j$ for Hilbert quantities $h_j$.  Let $r_j>0$ satisfy
\begin{equation}
 \sum_jr_j^{-2}<\infty.
 \label{eq:t10v6-scale-reserve-hyp}
\end{equation}
Then
\begin{equation}
 \sum_ja_j
 \leq
 \left(\sum_jr_j^{-2}\right)^{1/2}
 \left(\sum_jr_j^2C_j^2h_j^2\right)^{1/2}.
 \label{eq:t10v6-scale-reserve}
\end{equation}
One may take $r_j=(1+|j|)^{1+\epsilon}$ for any $\epsilon>0$.
\end{lemma}

\begin{proof}
Write $a_j\leq r_j^{-1}(r_jC_jh_j)$ and apply Cauchy--Schwarz.  The stated
polynomial sequence has square-summable reciprocal.
\end{proof}

\begin{proposition}[Position reserve inside each annulus]
\label{prop:t10v6-position-reserve}
Let $c_{j,z}$ be coefficients indexed by scale-$j$ block positions
$z\in\mathbb Z^d$.  If $s>d/2$, then
\begin{equation}
 \sum_z|c_{j,z}|
 \leq C_{d,s}
 \left[\sum_z(1+|z|)^{2s}|c_{j,z}|^2\right]^{1/2},
 \label{eq:t10v6-position-reserve}
\end{equation}
with $C_{d,s}^2=\sum_z(1+|z|)^{-2s}<\infty$.  Combining this with
Lemma~\ref{lem:t10v6-scale-reserve} gives an unrestricted multiscale owner
sum with one polynomial reserve in position and one in scale.
\end{proposition}

\begin{proof}
Apply Cauchy--Schwarz after inserting
$(1+|z|)^{-s}(1+|z|)^s$.  The lattice sum converges when $2s>d$.
\end{proof}

\begin{remark}[Relation to the two-exponent audit transfer]
\label{rem:t10v6-two-exponent}
The position reserve can instead be an unused exponential
$e^{-(\nu-\mu)|z|}$ with $\nu>\mu$, exactly as in the audited YM theorem.
The dyadic scale reserve is additional because an order-zero Hodge
multiplier has contributions on every annulus.
\end{remark}

\subsection{Reflection and boundary adaptation}

The full-space projector just computed is a diagnostic and an estimate.
It must not be applied across the reflection cut inside the measure.

\begin{definition}[Boundary-adapted multiscale Hodge card]
\label{def:t10v6-multiscale-card}
A scale-local gravitational gauge quotient passes the card if:
\begin{enumerate}
\item the zero and global Killing modes are retained explicitly;
\item every nonzero mode is assigned to one annulus and one rescaled block;
\item the block interior connection is the V4 Dirichlet connection and does
not cross the retained reflection layer;
\item boundary wavelets or extension/restriction operators reproduce the
annular estimates \eqref{eq:t10v6-Schwartz-bound}--
\eqref{eq:t10v6-commutator} with uniform owner multiplicity;
\item low annuli move to the next coarse scale whenever the commutator debit
is not small;
\item coefficient arrays carry the position and scale reserves of
Lemma~\ref{lem:t10v6-scale-reserve} and
Proposition~\ref{prop:t10v6-position-reserve};
\item normalized conditional maps are bounded in the multiscale Hilbert
norm, without claiming an exponentially summable full-space kernel;
\item boundary gauge frames and chiral data remain in the V96 reflected
factorization.
\end{enumerate}
\end{definition}

\begin{theorem}[Multiscale replacement implication]
\label{thm:t10v6-multiscale-implication}
Under Definition~\ref{def:t10v6-multiscale-card}, the order-zero
diffeomorphism quotient is bounded on the declared Hilbert coefficient
space, its block/annular tails are summable in the weaker owner norm, and it
does not mix the two interiors across the reflection layer.  Consequently
the locality estimates needed for RG comparison can coexist with the exact
V96 reflected-square factorization.
\end{theorem}

\begin{proof}
The square-function theorem gives the Hilbert bound.  The annular kernel
tail, position reserve, and scale reserve give the owner sums.  The V4
Dirichlet block condition prevents cross-cut action of the connection, so
the conditional reflected factorization is unchanged.
\end{proof}

\begin{firewall}[Missing boundary construction]
\label{fire:t10v6-boundary-wavelets}
V6 proves the full-space annular calculus but assumes boundary-adapted
wavelets or extension operators in the multiscale card.  A naive reflection
extension can mix gauge boundary conditions or double edge modes.  The
half-space construction and its exact trace sequence must be written
before the card is verified.
\end{firewall}

\subsection{V6 conclusion}

\begin{assessment}[What tenth-series V6 proves]
\label{assess:t10v6}
V6 computes the exact diffeomorphism Hodge projector and proves its
direction-dependent zero-momentum limit, ruling out a globally
exponentially summable kernel.  It replaces that false gate with a rigorous
dyadic annular decomposition, uniform rescaled Schwartz bounds, square-
function control, commutator gain, and position/scale owner reserves.  The
remaining finite task is a boundary-adapted realization compatible with the
V4 trace split and V96 reflection boundary.
\end{assessment}

\begin{programme}[Tenth-series V7: half-space Hodge frame]
\label{prog:t10v6-V7}
V7 will construct a half-space extension/restriction complex for the
linearized diffeomorphism generator, retain Dirichlet and normal boundary
traces, and prove that odd/even wavelet extensions give the required
annular estimates without coupling the two reflected interiors.
\end{programme}

\begin{firewall}[Claim boundary after tenth-series V6]
\label{fire:t10v6-boundary}
V6 does not construct the boundary-adapted Hodge frame, prove nonlinear
curved multiscale estimates, control the renormalized SM stress tensor,
trivialize the chiral determinant line, close the full RG trajectory, or
construct the SM--Einstein continuum.
\end{firewall}

\paragraph{Parent-version record.}
V6 was formed from
\texttt{Renewal\_SM\_Einstein\_Continuum\_Research\_Notes\_V5.tex}.
The V5 parent had $2\,184$ validator-counted source lines, $84\,966$ bytes,
and SHA--256
\[
 \texttt{67F36CEDDB63CD60CBB4EF145BFDE69170A638B0D61F35D64D3F7BE661311349}.
\]

% =============================================================================
% END APPENDED TENTH-SERIES V6 MATERIAL
% =============================================================================

% =============================================================================
% BEGIN APPENDED TENTH-SERIES V7 MATERIAL
% =============================================================================

\section{Reflection-compatible half-space Hodge frames}
\label{sec:v10v7}

\subsection{Geometric reflection parities}

V6 proved the required annular estimates in full space but left the
reflection boundary unresolved.  The essential observation is that vectors
and symmetric tensors have different component parities under a geometric
time reflection, and the symmetrized gradient intertwines them exactly.

Write a point as $(t,x)\in\mathbb R\times\mathbb R^{d-1}$ and let
$r(t,x)=(-t,x)$.  Let
\begin{equation}
 R_1=\operatorname{diag}(-1,1,\ldots,1)
 \label{eq:t10v7-vector-reflection-matrix}
\end{equation}
act on vectors.  On symmetric tensors set
\begin{equation}
 R_2h:=R_1hR_1.
 \label{eq:t10v7-tensor-reflection-matrix}
\end{equation}
Define unitary involutions
\begin{align}
 (\Theta_1\xi)(t,x)&=R_1\xi(-t,x),
 \label{eq:t10v7-vector-reflection}\\
 (\Theta_2h)(t,x)&=R_2h(-t,x).
 \label{eq:t10v7-tensor-reflection}
\end{align}

\begin{lemma}[Component parities]
\label{lem:t10v7-component-parities}
A vector fixed by $\Theta_1$ has odd normal component and even tangential
components.  A symmetric tensor fixed by $\Theta_2$ has even $00$ and
tangential--tangential components and odd mixed $0a$ components.
\end{lemma}

\begin{proof}
This follows by inserting the diagonal entries $(-1,1,\ldots,1)$ on one
vector index or on both tensor indices.
\end{proof}

Let
\begin{equation}
 (D\xi)_{\mu\nu}=\partial_\mu\xi_\nu+partial_\nu\xi_\mu
 \label{eq:t10v7-symmetric-gradient}
\end{equation}
be the linearized diffeomorphism generator.

\begin{theorem}[Reflection intertwining]
\label{thm:t10v7-intertwining}
On smooth compactly supported vector fields,
\begin{equation}
 D\Theta_1=\Theta_2D.
 \label{eq:t10v7-intertwining}
\end{equation}
Consequently $D^*\Theta_2=\Theta_1D^*$ and
$D^*D$ commutes with $\Theta_1$.
\end{theorem}

\begin{proof}
The chain rule gives
$\partial_0(f\circ r)=-(\partial_0f)\circ r$ and
$\partial_a(f\circ r)=(\partial_af)\circ r$.  Combining this derivative
sign with the factor $R_1$ on each vector index gives one $R_1$ factor on
each tensor index, exactly \eqref{eq:t10v7-intertwining}.  Taking Hilbert
adjoints proves the second identity, and composition proves the last.
\end{proof}

\begin{corollary}[Hodge projector commutes with reflection]
\label{cor:t10v7-projector-reflection}
On nonzero modes, let
\begin{equation}
 P=D(D^*D)^{-1}D^*
 \label{eq:t10v7-full-projector}
\end{equation}
with the zero modes retained outside the inverse.  Then
\begin{equation}
 P\Theta_2=\Theta_2P.
 \label{eq:t10v7-projector-commutes}
\end{equation}
Every dyadic annular piece $P_j$ of V6 also commutes with $\Theta_2$.
\end{corollary}

\begin{proof}
Theorem~\ref{thm:t10v7-intertwining} makes every factor in
\eqref{eq:t10v7-full-projector} intertwine the appropriate reflection.
Functional calculus of the commuting positive operator $D^*D$ preserves
the relation.  The annular frequency cutoff is radial and hence reflection
invariant.
\end{proof}

\subsection{Unitary doubling of half-space fields}

Let $\mathbb R_+^d=\{t>0\}$.  For a half-space tensor $h$, define its
geometric even extension
\begin{equation}
 (E_2h)(t,x)=
 \begin{cases}
 h(t,x),&t>0,\\
 R_2h(-t,x),&t<0.
 \end{cases}
 \label{eq:t10v7-tensor-extension}
\end{equation}
Put
\begin{equation}
 U_2:=2^{-1/2}E_2.
 \label{eq:t10v7-unitary-extension}
\end{equation}

\begin{lemma}[Half-space doubling is unitary]
\label{lem:t10v7-unitary-doubling}
$U_2$ is unitary from
$L^2(\mathbb R_+^d;\operatorname{Sym}^2V)$ onto the fixed subspace
\begin{equation}
 \mathcal H_2^+:=\{h\in L^2(\mathbb R^d):\Theta_2h=h\}.
 \label{eq:t10v7-fixed-tensor-space}
\end{equation}
Its adjoint is
\begin{equation}
 U_2^*g=\sqrt2\,g|_{t>0}
 \qquad(g\in\mathcal H_2^+).
 \label{eq:t10v7-extension-adjoint}
\end{equation}
\end{lemma}

\begin{proof}
Reflection preserves the Frobenius norm, so
$\|E_2h\|_{L^2(\mathbb R^d)}^2=2\|h\|_{L^2(\mathbb R_+^d)}^2$.
The extension is fixed by $\Theta_2$, and every fixed field is determined
by its positive-half restriction.  The adjoint formula follows from the
inner product and the two equal half contributions.
\end{proof}

\begin{definition}[Half-space Hodge frame]
\label{def:t10v7-half-Hodge}
Define
\begin{equation}
 P_+:=U_2^*PU_2,
 \qquad
 P_{j,+}:=U_2^*P_jU_2.
 \label{eq:t10v7-half-projectors}
\end{equation}
These operators use only the prescribed reflection extension of a
positive-half field; they do not read an independent field from the
negative half.
\end{definition}

\begin{theorem}[Exact half-space projection and square function]
\label{thm:t10v7-half-projection}
$P_+$ is an orthogonal projection of norm one and
\begin{equation}
 P_+=\sum_jP_{j,+}
 \label{eq:t10v7-half-reconstruction}
\end{equation}
on the nonzero-mode $L^2$ space.  If the V6 annular partition has overlap
$M_\chi$ and bound $C_\chi$, then
\begin{equation}
 \sum_j\|P_{j,+}h\|_2^2
 \leq M_\chi C_\chi^2\|h\|_2^2.
 \label{eq:t10v7-half-square-function}
\end{equation}
\end{theorem}

\begin{proof}
By Corollary~\ref{cor:t10v7-projector-reflection}, $P$ preserves the range
of $U_2$.  Unitary conjugation therefore preserves self-adjointness,
idempotence, and norm.  Conjugate the V6 reconstruction and square-function
identities by $U_2$.
\end{proof}

\subsection{Reflected image kernels}

Let $p_j(z)$ be the full-space matrix kernel of $P_j$.  For $y=(s,y')$ with
$s>0$, write $\bar y=(-s,y')$.

\begin{proposition}[Half-space image formula]
\label{prop:t10v7-image-kernel}
For $x,y\in\mathbb R_+^d$, the half-space annular kernel is
\begin{equation}
 p_{j,+}(x,y)=p_j(x-y)+p_j(x-\bar y)R_2,
 \label{eq:t10v7-image-kernel}
\end{equation}
with the natural action of $R_2$ on the input tensor.  It obeys, for every
$N$,
\begin{align}
 \|p_{j,+}(x,y)\|
 \leq{}&C_N2^{jd}\bigl[
 (1+2^j|x-y|)^{-N}
 +(1+2^j|x-\bar y|)^{-N}
 \bigr].
 \label{eq:t10v7-image-bound}
\end{align}
The row and column $L^1$ bounds and the scaled tail estimates are uniform in
$j$.
\end{proposition}

\begin{proof}
Expand $U_2^*P_jU_2$ using the two halves of the extension.  The factors
$2^{-1/2}$ and $\sqrt2$ cancel and the negative-half integral changes
variables by reflection, producing the second image term.  Apply the V6
Schwartz bound to both kernels.  The reflected half-space is a subset of
full space after the change of variables, so each image term has the same
$L^1$ and tail bounds as $p_j$.
\end{proof}

\begin{remark}[The image term is not cross-half dynamical coupling]
\label{rem:t10v7-image-not-coupling}
The value at $\bar y$ in
\eqref{eq:t10v7-image-kernel} is fixed algebraically by $h(y)$ and the
tensor parity.  It is not an independent negative-time integration
variable.  Thus the half-space frame is a one-sided operator and does not
alter the conditional independence used in the V96 reflected square.
\end{remark}

\subsection{Vector doubling and the generator range}

Define $E_1$ and $U_1=2^{-1/2}E_1$ for half-space vectors using $R_1$ in
place of $R_2$.  At Sobolev regularity one, the fixed vector fields satisfy
the geometric boundary condition
\begin{equation}
 \xi_0|_{t=0}=0,
 \label{eq:t10v7-normal-boundary-condition}
\end{equation}
while tangential traces are allowed and are retained interface
diffeomorphisms.

\begin{proposition}[Half-space generator intertwining]
\label{prop:t10v7-half-generator}
On smooth half-space vectors compatible with the geometric parity,
\begin{equation}
 DU_1=U_2D_+,
 \label{eq:t10v7-half-generator}
\end{equation}
where $D_+$ is the half-space symmetrized gradient.  Hence
$U_2^*\operatorname{Ran}D=\operatorname{Ran}D_+$ inside the fixed parity
sector, after zero modes and boundary traces are retained.
\end{proposition}

\begin{proof}
This is Theorem~\ref{thm:t10v7-intertwining} restricted to the fixed
subspaces and conjugated by the unitary extensions.  The qualification on
traces prevents a boundary distribution from appearing when a parity
extension is differentiated.
\end{proof}

\begin{firewall}[Interior based transformations and geometric parity]
\label{fire:t10v7-based-versus-parity}
The V4 interior gauge group fixes every boundary component, whereas the
geometric parity sector permits tangential boundary diffeomorphisms.  These
are not the same domain.  The correct direct-sum frame is:
\begin{enumerate}
\item the V4 Dirichlet block connection on transformations whose full trace
vanishes;
\item one retained shared boundary coordinate for tangential
diffeomorphisms and all other trace/flux modes;
\item the half-space Hodge frame above as an analytic comparison norm after
the trace split, not as an instruction to integrate the retained boundary
coordinate.
\end{enumerate}
Using $P_+$ on the retained trace would quotient a physical gluing variable
twice.
\end{firewall}

\subsection{Korn control for the interior range}

For compactly supported whole-space vector fields, Fourier transformation
of \eqref{eq:t10v7-symmetric-gradient} gives
\begin{equation}
 |D(k)\widehat\xi|^2
 =2|k|^2|\widehat\xi|^2+2|k\cdot\widehat\xi|^2.
 \label{eq:t10v7-symbol-Korn}
\end{equation}

\begin{lemma}[Whole-space Korn inequality]
\label{lem:t10v7-Korn}
For $\xi\in H^1(\mathbb R^d;V)$,
\begin{equation}
 \|\nabla\xi\|_2^2
 \leq\frac12\|D\xi\|_2^2.
 \label{eq:t10v7-Korn}
\end{equation}
For a half-space field whose geometric extension belongs to $H^1$, the
same inequality holds on $\mathbb R_+^d$.
\end{lemma}

\begin{proof}
Integrate \eqref{eq:t10v7-symbol-Korn} and use Plancherel.  The two half
integrals of a parity-compatible extension are equal, so division by two
gives the half-space statement.
\end{proof}

\begin{remark}[Rigid modes]
\label{rem:t10v7-rigid-modes}
On bounded blocks, Korn's inequality also needs removal or retention of
rigid Killing fields unless a full Dirichlet trace is imposed.  These modes
are precisely part of the V4 interface/closing-mode ledger and must not be
placed in $(D^*D)^{-1}$.
\end{remark}

\subsection{Half-space commutator and owner reserves}

Let $a$ be a bounded coefficient on the half-space whose tensor action has a
reflection-compatible Lipschitz extension $\widetilde a$.

\begin{proposition}[Half-space commutator gain]
\label{prop:t10v7-half-commutator}
The half-space annular frame satisfies
\begin{equation}
 \|[P_{j,+},M_a]\|_{2\to2}
 \leq C2^{-j}\|\nabla\widetilde a\|_\infty.
 \label{eq:t10v7-half-commutator}
\end{equation}
Its coefficient arrays satisfy the V6 position and scale reserve estimates
with at most a factor two from the image kernel.
\end{proposition}

\begin{proof}
Conjugate the full-space commutator by $U_2$ when multiplication intertwines
the extension, or apply the image-kernel formula directly.  The difference
$a(y)-a(x)$ supplies $|x-y|$, and the reflected term uses the Lipschitz
extension at $\bar y$.  The first moments of both image kernels are
$O(2^{-j})$.  The owner estimates use the two uniform image tails.
\end{proof}

\begin{definition}[Reflection-safe Hodge-frame card]
\label{def:t10v7-reflection-safe-card}
At every reflection cut and RG scale, require:
\begin{enumerate}
\item the V4 exact trace sequence and one copy of every boundary gauge,
Killing, normal, and flux coordinate;
\item the V4 Dirichlet connection on the zero-trace interior group;
\item unitary geometric doubling only for analytic comparison of bulk
fields, with parities $R_1,R_2$;
\item the half-space square-function, image-tail, and commutator bounds above;
\item boundary-compatible coefficient extensions and the V6 position/scale
reserves;
\item no application of $P_+$ to an independently integrated field across
the reflection cut;
\item exact reflection covariance of the retained boundary density and
chiral data.
\end{enumerate}
\end{definition}

\begin{theorem}[Half-space multiscale implication]
\label{thm:t10v7-half-space-implication}
If the V4 scale-local connection card and
Definition~\ref{def:t10v7-reflection-safe-card} hold, the diffeomorphism
quotient admits a uniformly bounded half-space multiscale comparison frame,
with summable owner tails and first commutator gain, while the actual
interior gauge reduction remains block-local and the retained boundary
variables preserve the V96 reflected-square factorization.
\end{theorem}

\begin{proof}
The unitary conjugation, image-kernel, and commutator theorems give the
comparison frame.  V4 supplies the actual one-sided Dirichlet quotient and
trace split.  Since neither operation integrates an independent variable
from the opposite half, the conditional product and its reflected square
remain exact.
\end{proof}

\begin{assessment}[Status after tenth-series V7]
\label{assess:t10v7}
V7 constructs the reflection-compatible half-space Hodge frame.  It proves
the vector/tensor parity intertwining, unitary doubling, exact half-space
projection, reflected image-kernel bounds, Korn control, and the half-space
commutator estimate.  It also separates this analytic comparison frame from
the V4 based interior quotient, preventing boundary diffeomorphisms from
being removed twice.  The next problem is curved-background stability of
the annular frame and its interaction with the stationary quotient Hessian.
\end{assessment}

\begin{programme}[Tenth-series V8: curved-frame perturbation]
\label{prog:t10v7-V8}
V8 will compare the flat annular Hodge frame with the generator and metric
at a slowly varying calibrated background.  It will derive a resolvent/
projector identity in each annulus, quantify the derivative gain, and state
a scale-separation gate that sends nonperturbative low annuli to the next
coarse scale.
\end{programme}

\begin{firewall}[Claim boundary after tenth-series V7]
\label{fire:t10v7-boundary}
V7 does not prove curved nonlinear Hodge bounds, positivity of the complete
Einstein boundary precision, the renormalized SM stress-tensor tail,
chiral determinant-line gluing, the invariant RG trajectory, or the full
SM--Einstein continuum.
\end{firewall}

\paragraph{Parent-version record.}
V7 was formed from
\texttt{Renewal\_SM\_Einstein\_Continuum\_Research\_Notes\_V6.tex}.
The V6 parent had $2\,605$ validator-counted source lines, $100\,085$ bytes,
and SHA--256
\[
 \texttt{979827593BDB7422505B101BC067ABCC9AF09F7B7BAF9AD7F52F8C21DAB15849}.
\]

% =============================================================================
% END APPENDED TENTH-SERIES V7 MATERIAL
% =============================================================================

% =============================================================================
% BEGIN APPENDED TENTH-SERIES V8 MATERIAL
% =============================================================================

\section{Curved annular Hodge frames and the scale-separation gate}
\label{sec:v10v8}

\subsection{An abstract relative Gram perturbation}

V7 gives a reflection-safe flat half-space frame.  A calibrated Einstein
background changes the diffeomorphism generator and its coefficient metric.
The right perturbation variable is the generator normalized by the flat
annular Gram form, not its absolute operator norm.

Let $D_0:U\to V$ be injective between finite-dimensional Hilbert spaces,
let $L_0=D_0^*D_0>0$, and set
\begin{equation}
 R_0:=D_0L_0^{-1/2}.
 \label{eq:t10v8-flat-isometry}
\end{equation}
Then $R_0^*R_0=I_U$.  Let $D=D_0+E$, $L=D^*D$, and
\begin{equation}
 F:=EL_0^{-1/2},
 \qquad \theta:=\|F\|.
 \label{eq:t10v8-normalized-generator-error}
\end{equation}

\begin{theorem}[Relative Gram perturbation]
\label{thm:t10v8-relative-Gram}
With
\begin{equation}
 X:=L_0^{-1/2}(L-L_0)L_0^{-1/2},
 \label{eq:t10v8-relative-Gram-error}
\end{equation}
one has the exact identity
\begin{equation}
 X=R_0^*F+F^*R_0+F^*F
 \label{eq:t10v8-Gram-identity}
\end{equation}
and the bound
\begin{equation}
 \|X\|\leq2\theta+\theta^2=:\eta(\theta).
 \label{eq:t10v8-Gram-bound}
\end{equation}
If $\eta(\theta)<1$, then
\begin{align}
 (1-\eta)L_0&\preceq L\preceq(1+\eta)L_0,
 \label{eq:t10v8-Gram-comparison}\\
 \|L_0^{1/2}L^{-1}L_0^{1/2}\|
 &\leq(1-\eta)^{-1}.
 \label{eq:t10v8-inverse-Gram-comparison}
\end{align}
\end{theorem}

\begin{proof}
$DL_0^{-1/2}=R_0+F$, so
\[
 L_0^{-1/2}LL_0^{-1/2}
 =(R_0+F)^*(R_0+F)
 =I+R_0^*F+F^*R_0+F^*F.
\]
This proves the identity.  Since $\|R_0\|=1$, the triangle inequality gives
\eqref{eq:t10v8-Gram-bound}.  The spectral interval of $I+X$ lies in
$[1-\eta,1+\eta]$, giving the form and inverse bounds.
\end{proof}

\begin{corollary}[Normalized curved lift]
\label{cor:t10v8-curved-isometry}
Under $\eta<1$, the curved normalized lift
\begin{equation}
 R:=DL^{-1/2}
 \label{eq:t10v8-curved-isometry}
\end{equation}
is an isometry and obeys
\begin{equation}
 \|DL_0^{-1/2}\|
 \leq1+\theta.
 \label{eq:t10v8-curved-flat-frame}
\end{equation}
The isometry constant is exactly one; only comparison with the flat frame
costs $\theta$ and $\eta$.
\end{corollary}

\begin{proof}
$R^*R=L^{-1/2}D^*DL^{-1/2}=I$, and
$DL_0^{-1/2}=R_0+F$.
\end{proof}

\subsection{Subspace-projector stability}

Let
\begin{equation}
 P_0=D_0L_0^{-1}D_0^*,
 \qquad P=DL^{-1}D^*
 \label{eq:t10v8-range-projectors}
\end{equation}
be the orthogonal projectors onto the two generator ranges.

\begin{lemma}[Gap between the two ranges]
\label{lem:t10v8-range-gap}
If $\theta<1$, then
\begin{equation}
 \|(I-P_0)P\|
 \leq\frac{\theta}{1-\theta},
 \qquad
 \|(I-P)P_0\|\leq\theta.
 \label{eq:t10v8-one-sided-gaps}
\end{equation}
Consequently
\begin{equation}
 \|P-P_0\|
 \leq\frac{2\theta}{1-\theta}.
 \label{eq:t10v8-projector-difference}
\end{equation}
\end{lemma}

\begin{proof}
For $y=Du$ in $\operatorname{Ran}D$,
$(I-P_0)y=(I-P_0)Eu$, while
\[
 \|Du\|\geq\|D_0u\|-\|Eu\|
 \geq(1-\theta)\|L_0^{1/2}u\|.
\]
Thus the distance of a unit vector in $\operatorname{Ran}D$ from
$\operatorname{Ran}D_0$ is at most $\theta/(1-\theta)$.  Conversely, for
$y=D_0u$, the vector $Du$ lies in $\operatorname{Ran}D$ and
$\|y-Du\|\leq\theta\|y\|$, proving the second gap.  Finally
\[
 P-P_0=(I-P_0)P-P_0(I-P),
\]
and the triangle inequality gives the stated safe bound.
\end{proof}

\begin{remark}[Nonoptimal but uniform]
\label{rem:t10v8-projector-constant}
The projector constant is not sharp.  Its role is to expose the
dimensionless annular parameter $\theta$.  No inverse power of the global
smallest eigenvalue appears.
\end{remark}

\subsection{Generator error on one annulus}

Let $Q_j$ be a smooth frequency localization to $|k|\simeq2^j$.  On its
range, the flat symmetric gradient obeys
\begin{equation}
 c_02^j\|u\|_2\leq\|D_0u\|_2\leq C_02^j\|u\|_2
 \label{eq:t10v8-annular-ellipticity}
\end{equation}
after zero and retained modes are removed.

In a local orthonormal frame for a background metric $g$, write schematically
\begin{equation}
 D_g=D_0+A^\alpha(x)\partial_\alpha+B(x),
 \label{eq:t10v8-curved-generator}
\end{equation}
where $A$ measures the principal-frame difference and $B$ the connection
term.

\begin{proposition}[Annular normalized error]
\label{prop:t10v8-annular-error}
Suppose
\begin{equation}
 \|A\|_\infty\leq\delta_0,
 \qquad \|B\|_\infty\leq\delta_1.
 \label{eq:t10v8-coefficient-stocks}
\end{equation}
Then, up to a dimension- and cutoff-dependent constant $C_D$,
\begin{equation}
 \|(D_g-D_0)Q_jL_{0,j}^{-1/2}\|
 \leq
 \theta_j:=C_Dc_0^{-1}
 \left(\delta_0+2^{-j}\delta_1\right),
 \label{eq:t10v8-annular-theta}
\end{equation}
where $L_{0,j}$ is the flat Gram form on the annulus.  If $A$ varies in
space, leakage across the smooth cutoff adds a commutator debit bounded by
$C2^{-j}\|\nabla A\|_\infty$.
\end{proposition}

\begin{proof}
Bernstein's annular estimate gives
$\|\partial Q_ju\|_2\leq C2^j\|Q_ju\|_2$.  Therefore
\[
 \|(D_g-D_0)Q_ju\|_2
 \leq C_D(\delta_02^j+\delta_1)\|Q_ju\|_2.
\]
The lower bound in \eqref{eq:t10v8-annular-ellipticity} says that
$L_{0,j}^{-1/2}$ costs at most $(c_02^j)^{-1}$.  This proves the main
estimate.  The cutoff leakage is the V6/V7 first commutator estimate.
\end{proof}

\subsection{Normal frames on slowly varying blocks}

Choose the flat comparison frame at the centre $x_B$ of a rescaled block.
If the metric is $C^1$, the mean-value theorem gives
\begin{equation}
 \sup_{x\in B}\|g(x)-g(x_B)\|
 \leq\operatorname{diam}(B)\|\nabla g\|_{L^\infty(B)}.
 \label{eq:t10v8-frame-variation}
\end{equation}
At physical annular length $2^{-j}$, use a subblock of comparable diameter.

\begin{corollary}[Slow-background annular parameter]
\label{cor:t10v8-slow-background}
For a normal frame on a scale-$j$ patch, the principal and connection
stocks may be chosen so that
\begin{equation}
 \theta_j
 \leq C_g2^{-j}
 \left(\|\nabla g\|_\infty+|\Gamma(g)\|_\infty\right)
 +C_g2^{-2j}\|\operatorname{Rm}(g)\|_\infty,
 \label{eq:t10v8-geometric-theta}
\end{equation}
provided the chart remains uniformly nondegenerate.  Thus
$\theta_j\to0$ toward the ultraviolet at a fixed smooth background.
\end{corollary}

\begin{proof}
Normal-coordinate Taylor expansion makes the principal coefficient
difference first order in patch diameter and the next remainder second
order with curvature.  The connection is first order in the metric
derivative.  Insert these stocks in
Proposition~\ref{prop:t10v8-annular-error}.
\end{proof}

\begin{firewall}[Uniform smoothness is not assumed for the quantum field]
\label{fire:t10v8-random-background}
Corollary~\ref{cor:t10v8-slow-background} is a deterministic good-field
estimate.  A quantum metric is not assumed uniformly $C^1$.  The regulator
must supply a good-field probability/tail card for the displayed geometric
stocks and send the complement through the V92/V93 bad-field ledger.
\end{firewall}

\subsection{The scale-separation gate}

Fix $0<\theta_*<\sqrt2-1$, so
$2\theta_*+\theta_*^2<1$.  Define an annulus to be perturbative when
\begin{equation}
 \theta_j\leq\theta_*.
 \label{eq:t10v8-perturbative-annulus}
\end{equation}

\begin{definition}[Coarse-handoff index]
\label{def:t10v8-handoff-index}
For a local geometric stock $K_B\geq0$, choose
\begin{equation}
 j_c(B):=\min\{j:C_g2^{-j}K_B\leq\theta_*/2\},
 \label{eq:t10v8-handoff-index}
\end{equation}
with the curvature remainder included in $K_B$ at the declared scale.
Annuli $j\geq j_c(B)$ are treated by the perturbative curved frame.  Annuli
$j<j_c(B)$ are transported as retained data to the next coarse scale.
\end{definition}

\begin{theorem}[Annular perturbation or exact handoff]
\label{thm:t10v8-scale-separation}
On every perturbative annulus, the curved Gram form and inverse obey
Theorem~\ref{thm:t10v8-relative-Gram}, and the gauge-range projector obeys
\begin{equation}
 \|P_{g,j}-P_{0,j}\|
 \leq\frac{2\theta_j}{1-\theta_j}.
 \label{eq:t10v8-annular-projector-debit}
\end{equation}
If every nonperturbative low annulus is retained and transferred without
inversion, the construction makes no false smallness claim and covers all
nonzero modes exactly once.
\end{theorem}

\begin{proof}
Apply Theorem~\ref{thm:t10v8-relative-Gram} and
Lemma~\ref{lem:t10v8-range-gap} on the high annuli.  The dyadic partition
assigns every nonzero frequency to finitely many neighboring annuli.  The
handoff rule partitions these owners into the perturbative set and the
retained set; no mode is discarded.
\end{proof}

\begin{remark}[Termination of repeated handoff]
\label{rem:t10v8-handoff-termination}
A low annulus at one scale becomes a bounded-frequency annulus in the next
rescaled coarse lattice.  To prevent endless relabeling without control,
the RG card must prove either that its geometric stock becomes
perturbative, or that it enters a finite declared marginal/background
sector.  Merely saying ``send it to the next scale'' is not closure.
\end{remark}

\subsection{Curved quotient-Hessian comparison}

Let $K_0$ and $K_g$ be stationary action Hessians on one annulus and set
$H_0=I-P_0$, $H_g=I-P_g$.  Define
\begin{equation}
 K_0^{\rm q}=H_0K_0H_0,
 \qquad
 K_g^{\rm q}=H_gK_gH_g.
 \label{eq:t10v8-curved-quotient-Hessians}
\end{equation}

\begin{proposition}[Quotient-Hessian debit]
\label{prop:t10v8-quotient-debit}
If $\|K_g-K_0\|\leq\delta_K$ and
$\|P_g-P_0\|\leq\delta_P$, then
\begin{equation}
 \|K_g^{\rm q}-K_0^{\rm q}\|
 \leq\delta_K+2\delta_P\|K_0\|+delta_P^2\|K_0\|
 +2\delta_P\delta_K+\delta_P^2\delta_K.
 \label{eq:t10v8-quotient-debit}
\end{equation}
In particular the debit is $O(\delta_K+\delta_P)$ when both are small.
\end{proposition}

\begin{proof}
Write $H_g=H_0-\Delta P$ and $K_g=K_0+\Delta K$, expand
$(H_0-\Delta P)(K_0+\Delta K)(H_0-\Delta P)-H_0K_0H_0$, and bound each
term using $\|H_0\|=1$.
\end{proof}

Suppose the reference physical boundary precision on the annulus has size
$C_{0,j}\succeq c_B2^jI$.  If the quotient debit is bounded by
$\epsilon_j2^j$, then
\begin{equation}
 \|C_{0,j}^{-1/2}(K_g^{\rm q}-K_0^{\rm q})C_{0,j}^{-1/2}\|
 \leq c_B^{-1}\epsilon_j.
 \label{eq:t10v8-relative-quotient-debit}
\end{equation}
This is the annular form of the Ward-relative precision gate.

\subsection{Summability across perturbative annuli}

\begin{proposition}[Geometric debit summation]
\label{prop:t10v8-debit-sum}
Suppose for $j\geq j_c$,
\begin{equation}
 \delta_{P,j}+2^{-j}\delta_{K,j}
 \leq C2^{-j}K_B.
 \label{eq:t10v8-decaying-debit}
\end{equation}
Then
\begin{equation}
 \sum_{j\geq j_c}
 (\delta_{P,j}+2^{-j}\delta_{K,j})
 \leq2C2^{-j_c}K_B.
 \label{eq:t10v8-debit-sum}
\end{equation}
With the V6 position/scale reserve, the same geometric decay is summable in
the owner norm.
\end{proposition}

\begin{proof}
Sum the geometric series $\sum_{j\geq j_c}2^{-j}=2^{1-j_c}$.  The reserve
handles position and any finite overlap of neighboring annuli.
\end{proof}

\subsection{Reflection covariance of the curved frame}

\begin{proposition}[Exact reflected curved annuli]
\label{prop:t10v8-curved-reflection}
If the regulated background metric and calibration are exactly reflection
invariant at the cut, then $D_g$ intertwines the vector and tensor
reflections of V7, $L_g=D_g^*D_g$ commutes with vector reflection, and every
spectral or annular function of $L_g$ preserves the parity sectors.  The
curved half-space frame obtained by unitary doubling is therefore one-sided
and reflection covariant.
\end{proposition}

\begin{proof}
The covariant symmetrized gradient is natural under isometries.  Exact
reflection invariance makes the reflection an isometry, giving the same
intertwining calculation as V7.  Adjoint composition and functional
calculus give the remaining statements.
\end{proof}

\begin{firewall}[Approximate reflection is a different theorem]
\label{fire:t10v8-approximate-reflection}
A small reflection defect does not give exact finite-regulator OS
positivity.  It must either vanish in a reflected-operator norm so that the
V96 vanishing-negativity theorem applies, or be removed by a symmetric
choice of regulator and calibrated reference.  The annular perturbation
bound alone is insufficient.
\end{firewall}

\subsection{Curved-frame acceptance card}

\begin{definition}[Curved annular frame card]
\label{def:t10v8-curved-card}
For every good-field block and annulus, require:
\begin{enumerate}
\item flat annular ellipticity after zero and retained modes are removed;
\item normalized generator error $\theta_j$ with either
$\theta_j\leq\theta_*$ or exact coarse handoff;
\item action-Hessian and projector debits satisfying the relative boundary
precision inequality \eqref{eq:t10v8-relative-quotient-debit};
\item summable high-annulus geometric debits and the V6 owner reserves;
\item a proved termination rule for repeated low-annulus handoff;
\item exact reflection covariance or a vanishing reflected-operator debit;
\item good-field tails for the geometric coefficient stocks;
\item compatibility with the V4 trace split and V5 stationary covariance
policy.
\end{enumerate}
\end{definition}

\begin{theorem}[Curved-frame implication]
\label{thm:t10v8-curved-implication}
Under Definition~\ref{def:t10v8-curved-card}, the stationary curved
quotient Hessian is relatively controlled annulus by annulus against the
physical order-one boundary precision.  Its high-annulus comparison errors
are owner-summable, while every uncontrolled low annulus remains an
explicit coarse or marginal coordinate.  Reflection factorization and the
V92 field-tail route are preserved on the good-field set.
\end{theorem}

\begin{proof}
The relative Gram, projector, and quotient-debit estimates give the local
precision bounds.  Proposition~\ref{prop:t10v8-debit-sum} and the V6
reserves give owner summability.  The handoff and reflection clauses prevent
uncontrolled inversion and preserve the V96 cut.  The geometric good-field
tail enters the V92 bad-field union bound.
\end{proof}

\begin{assessment}[Status after tenth-series V8]
\label{assess:t10v8}
V8 proves a relative Gram identity, range-projector stability, annular
curved-generator estimate, scale-separation gate, quotient-Hessian debit,
summability of high-annulus errors, and reflected curved-frame covariance.
It converts background curvature into a dimensionless scale-dependent
stock and sends nonperturbative low annuli to explicit coarse variables.
The remaining issue is to assemble the TT, constrained scalar, gauge, and
boundary sectors into one finite Gaussian gravitational RG step with a
positive reduced precision and matched conditional contraction.
\end{assessment}

\begin{programme}[Tenth-series V9: one-step reduced gravitational Gaussian]
\label{prog:t10v8-V9}
V9 will assemble one reflection-aligned Gaussian block step: quotient the
interior diffeomorphisms, retain boundary and low annuli, eliminate the
negative scalar by its compatible constraint, compute the full reduced
Schur precision, and prove covariance-matched conditional contraction on
the TT and constrained-source sectors.
\end{programme}

\begin{firewall}[Claim boundary after tenth-series V8]
\label{fire:t10v8-boundary}
V8 does not construct the complete positive gravitational block measure,
derive the nonlinear Einstein scalar stationarity equation, control the
renormalized SM source, prove chiral determinant-line gluing, close the RG
flow, or construct the full SM--Einstein continuum.
\end{firewall}

\paragraph{Parent-version record.}
V8 was formed from
\texttt{Renewal\_SM\_Einstein\_Continuum\_Research\_Notes\_V7.tex}.
The V7 parent had $3\,011$ validator-counted source lines, $114\,673$ bytes,
and SHA--256
\[
 \texttt{62F44F6CA8E1BF2A2B5BD068F0B153F86BCB225AB433B90F0B6B76BD23F560B5}.
\]

% =============================================================================
% END APPENDED TENTH-SERIES V8 MATERIAL
% =============================================================================

% =============================================================================
% BEGIN APPENDED TENTH-SERIES V9 MATERIAL
% =============================================================================

\section{One reduced gravitational Gaussian block step}
\label{sec:v10v9}

\subsection{Order of operations}

The preceding notes have supplied local gauge quotient, reflection-boundary
retention, annular curved-frame control, and a stationary scalar constraint.
The order matters.  One first quotients only the based interior
diffeomorphisms, then integrates positive interior variables, then imposes
the conformal stationary constraint, and only afterward normalizes a
positive Gaussian measure.  Integrating the negative conformal scalar as a
real Gaussian would diverge.

Let $Z,M,S$ be finite-dimensional real Hilbert spaces.  The variable
$z\in Z$ collects positive interior TT, gauge-fixed, and auxiliary
directions to be integrated.  The variable $m\in M$ collects retained
physical boundary/matter coordinates.  The variable $s\in S$ is the
negative conformal scalar to be constrained.  Consider
\begin{align}
 Q(z,m,s)={}&
 \frac12\langle Az,z\rangle
 +\langle K_mm+K_ss,z\rangle
 +\frac12\langle Dm,m\rangle
 \notag\\
 &+\langle Bm,s\rangle
 -\frac12\langle Cs,s\rangle,
 \label{eq:t10v9-indefinite-Gaussian}
\end{align}
where $A=A^*>0$, $C=C^*>0$, and $D=D^*$.

\subsection{Positive interior elimination}

\begin{lemma}[First completed square]
\label{lem:t10v9-first-square}
Put
\begin{align}
 \overline D&:=D-K_m^*A^{-1}K_m,
 \label{eq:t10v9-Dbar}\\
 \overline B&:=B-K_s^*A^{-1}K_m,
 \label{eq:t10v9-Bbar}\\
 \overline C&:=C+K_s^*A^{-1}K_s.
 \label{eq:t10v9-Cbar}
\end{align}
Then
\begin{align}
 Q(z,m,s)={}&
 \frac12\left\langle
 A[z+A^{-1}(K_mm+K_ss)],
 z+A^{-1}(K_mm+K_ss)
 \right\rangle
 \notag\\
 &+\frac12\langle\overline Dm,m\rangle
 +\langle\overline Bm,s\rangle
 -\frac12\langle\overline Cs,s\rangle.
 \label{eq:t10v9-first-square}
\end{align}
In particular $\overline C>0$.
\end{lemma}

\begin{proof}
Complete the $A$ square.  Its subtracted term is
\[
 \frac12\langle K_mm+K_ss,A^{-1}(K_mm+K_ss)\rangle.
\]
The $m^2$, $m$--$s$, and $s^2$ coefficients are respectively
\eqref{eq:t10v9-Dbar}, \eqref{eq:t10v9-Bbar}, and the negative of
\eqref{eq:t10v9-Cbar}.  Positivity of the last follows from $C>0$ and
$K_s^*A^{-1}K_s\succeq0$.
\end{proof}

\begin{corollary}[Interior Gaussian integral]
\label{cor:t10v9-interior-integral}
For fixed $(m,s)$,
\begin{align}
 \int_Ze^{-Q(z,m,s)}dz
 ={}&(2\pi)^{\dim Z/2}(\det A)^{-1/2}
 e^{-\overline Q(m,s)},
 \label{eq:t10v9-interior-integral}\\
 \overline Q(m,s):={}&
 \frac12\langle\overline Dm,m\rangle
 +\langle\overline Bm,s\rangle
 -\frac12\langle\overline Cs,s\rangle.
 \label{eq:t10v9-Qbar}
\end{align}
The remaining real $s$ integral diverges and is not performed.
\end{corollary}

\begin{proof}
Translate $z$ in \eqref{eq:t10v9-first-square} and use the standard
positive Gaussian integral.  Since $\overline C>0$, the exponent
$-\overline Q$ grows positively in $s$, proving divergence of an
unconstrained real integral.
\end{proof}

\subsection{Stationary scalar constraint}

The stationary equation of the same reduced quadratic form is
\begin{equation}
 D_s\overline Q=\overline Bm-\overline Cs=0.
 \label{eq:t10v9-scalar-constraint}
\end{equation}

\begin{theorem}[Constrained scalar positivity]
\label{thm:t10v9-constrained-positivity}
The unique scalar solution is
\begin{equation}
 s_*(m)=\overline C^{-1}\overline Bm.
 \label{eq:t10v9-scalar-solution}
\end{equation}
On its graph,
\begin{align}
 Q_{\rm red}(m)
 &:=\overline Q(m,s_*(m))
 =\frac12\langle S_{\rm red}m,m\rangle,
 \label{eq:t10v9-reduced-action}\\
 S_{\rm red}
 &:=\overline D+\overline B^*\overline C^{-1}\overline B.
 \label{eq:t10v9-reduced-precision}
\end{align}
If
\begin{equation}
 \overline D\succeq d_0I>0,
 \label{eq:t10v9-positive-physical-Schur}
\end{equation}
then
\begin{equation}
 S_{\rm red}\succeq d_0I.
 \label{eq:t10v9-reduced-floor}
\end{equation}
No smallness assumption on $\overline B$ is required.
\end{theorem}

\begin{proof}
Invertibility of $\overline C$ gives the unique root.  Substitution yields
\[
 \langle\overline Bm,\overline C^{-1}\overline Bm\rangle
 -\frac12\langle\overline Bm,
 \overline C^{-1}\overline Bm\rangle
 =\frac12\langle\overline Bm,
 \overline C^{-1}\overline Bm\rangle.
\]
This proves \eqref{eq:t10v9-reduced-precision}.  The second term in
$S_{\rm red}$ is positive semidefinite, so
\eqref{eq:t10v9-reduced-floor} follows from
\eqref{eq:t10v9-positive-physical-Schur}.
\end{proof}

\begin{remark}[The physical positivity gate]
\label{rem:t10v9-physical-Schur-gate}
All conformal positivity is now exact.  The remaining Gaussian sign test is
\eqref{eq:t10v9-positive-physical-Schur}: after integrating the positive
interior, the TT/matter boundary Schur block must retain a positive floor.
This is a finite matrix or symbol inequality at each regulator, not a
consequence of $A>0$ alone.
\end{remark}

\subsection{Delta constraint and exact density factors}

\begin{theorem}[One-step constrained Gaussian density]
\label{thm:t10v9-constrained-density}
Assume \eqref{eq:t10v9-positive-physical-Schur}.  Then
\begin{align}
 &\int_Z\int_S
 e^{-Q(z,m,s)}
 \delta(\overline Cs-\overline Bm)\,ds\,dz
 \notag\\
 &\qquad=(2\pi)^{\dim Z/2}
 (\det A)^{-1/2}(\det\overline C)^{-1}
 e^{-\langle S_{\rm red}m,m\rangle/2}.
 \label{eq:t10v9-constrained-density}
\end{align}
The right-hand side is a positive integrable density in $m$.  Its three
factors have distinct origins: positive interior Gaussian integration,
constraint coarea, and reduced physical precision.
\end{theorem}

\begin{proof}
Use Corollary~\ref{cor:t10v9-interior-integral}.  The linear change of
variables $u=\overline Cs-\overline Bm$ gives
$ds=(\det\overline C)^{-1}du$ because $\overline C>0$.  Evaluate at $u=0$
and apply Theorem~\ref{thm:t10v9-constrained-positivity}.  The positive
floor makes the final Gaussian integral finite.
\end{proof}

\begin{firewall}[No ghost substitution]
\label{fire:t10v9-density-factors}
The absolute coarea factor $(\det\overline C)^{-1}$ is not a fermion or
ghost determinant and cannot cancel one by notation.  Gauge orientation,
Faddeev--Popov factors, and the chiral determinant line remain separate rows
of the regulated density.
\end{firewall}

\subsection{Reflection-aligned pair of half blocks}

Let $Q_+$ be the positive-half quadratic form above and let $Q_-$ be its
reflected copy.  Retain a shared boundary coordinate $b$ in $m$, and let all
strict-interior variables be separate at fixed $b$.

\begin{proposition}[Reduced reflected square]
\label{prop:t10v9-reflected-square}
Suppose the block quotient, interior matrices, scalar constraint, coarea
factor, and physical Schur form on the negative half are exact reflected
copies of those on the positive half.  Suppose all boundary-only bosonic
factors form a positive measure $d\nu(b)$.  Then, after the two reductions
of Theorem~\ref{thm:t10v9-constrained-density}, the full measure has the
V96 form
\begin{equation}
 d\mu=N^{-1}d\nu(b)\,dK_b(m_+)\,dK_b(m_-),
 \label{eq:t10v9-reflected-product}
\end{equation}
where $m_-$ is identified with a reflected positive-half coordinate and
$K_b$ is positive.  Hence the bosonic Gaussian step is OS positive.
\end{proposition}

\begin{proof}
The determinant and reduced-precision factors in
\eqref{eq:t10v9-constrained-density} agree under reflection and are
positive.  At fixed $b$, no variable is shared between the strict
interiors.  Absorb boundary-only factors into $d\nu$ and apply the V96
conditional reflected-square theorem.
\end{proof}

\begin{remark}[Tangential nonlocality is compatible with the square]
\label{rem:t10v9-tangential-nonlocality}
$\overline C^{-1}$ may be nonlocal along the retained boundary, with the
annular behavior of V6--V8.  This does not mix the two time interiors and
therefore does not spoil the algebraic reflected square.  Its tangential
owner sums must nevertheless satisfy the multiscale comparison card.
\end{remark}

\subsection{A physical coarse Gaussian map}

Now let $H_f>0$ be the complete reduced fine precision obtained after all
steps above, and let
\begin{equation}
 T:F\longrightarrow Y
 \label{eq:t10v9-coarse-map}
\end{equation}
be a surjective, constraint-compatible physical coarse map.  Retained
zero/closing modes not carrying a positive precision remain outside this
Gaussian block.

Define
\begin{align}
 C_Y&:=TH_f^{-1}T^*,
 &S_Y&:=C_Y^{-1},
 \label{eq:t10v9-coarse-covariance}\\
 J&:=H_f^{-1}T^*S_Y,
 &C_{\rm fib}&:=H_f^{-1}-JC_YJ^*.
 \label{eq:t10v9-conditional-data}
\end{align}

\begin{theorem}[Reduced covariance-matched block contraction]
\label{thm:t10v9-matched-contraction}
The operators in
\eqref{eq:t10v9-coarse-covariance}--
\eqref{eq:t10v9-conditional-data} satisfy
\begin{equation}
 TJ=I,
 \qquad J^*H_fJ=S_Y,
 \qquad C_{\rm fib}\succeq0.
 \label{eq:t10v9-conditional-identities}
\end{equation}
The whitened physical lift
\begin{equation}
 R=H_f^{1/2}JS_Y^{-1/2}
 \label{eq:t10v9-whitened-lift}
\end{equation}
is an isometry.  Consequently conditional expectation from the fine
Gaussian to $Y$ is an orthogonal contraction on every covariance-matched
Wick chaos, including the TT and constrained-source components of $H_f$.
\end{theorem}

\begin{proof}
Surjectivity makes $C_Y>0$.  Direct multiplication gives
$TJ=TH_f^{-1}T^*S_Y=C_YS_Y=I$ and
$J^*H_fJ=S_YC_YS_Y=S_Y$.  Moreover
\[
 H_f^{1/2}C_{\rm fib}H_f^{1/2}=I-RR^*\succeq0
\]
because $R^*R=I$.  The covariance-matched conditional Wick identity proved
in the archived Gaussian interface then says that the degree-$n$
coefficient map is $(R^*)^{\otimes n}$ on the symmetric tensor power.  Its
norm is at most one.
\end{proof}

\begin{remark}[New content versus the archived Gaussian theorem]
\label{rem:t10v9-new-content}
The matched contraction identity itself was proved earlier and was also
confirmed in the YM companion notes.  The new point in V9 is that the
precision $H_f$ is obtained after an explicit nested positive-interior
elimination and stationary conformal constraint, with its coarea factor and
reflection square intact.  No indefinite conformal covariance enters the
matched Gaussian theorem.
\end{remark}

\subsection{Uniform smeared-field tails for the Gaussian step}

\begin{proposition}[Sub-Gaussian physical fields]
\label{prop:t10v9-subgaussian-fields}
Let $X_f$ be the centred Gaussian with precision $H_f$ and let a physical
smeared field be
$\Phi_f(f)=\langle\ell_f,X_f\rangle$.  If
\begin{equation}
 \langle\ell_f,H_f^{-1}\ell_f\rangle
 \leq C p(f)^2
 \label{eq:t10v9-covariance-test-bound}
\end{equation}
uniformly in the regulator, then
\begin{equation}
 \mathbb E e^{\lambda\Phi_f(f)}
 \leq e^{C\lambda^2p(f)^2/2}
 \label{eq:t10v9-subgaussian-mgf}
\end{equation}
and the V92 characteristic equicontinuity bound holds uniformly.
\end{proposition}

\begin{proof}
The moment-generating function of a centred Gaussian linear functional is
the exponential of one half its variance times $\lambda^2$.  Apply
\eqref{eq:t10v9-covariance-test-bound} and the V92 sub-Gaussian theorem.
\end{proof}

\begin{corollary}[Coarse and conditional tail stability]
\label{cor:t10v9-coarse-tail}
If \eqref{eq:t10v9-covariance-test-bound} holds in the fine physical energy
frame, it passes to coarse covariance-matched observables with no loss from
the exact conditional lift.  Annular owner and curved-frame debits remain
the explicit V6--V8 comparison factors.
\end{corollary}

\begin{proof}
The whitened lift is an isometry and its adjoint a contraction, so the
one-particle covariance norm of a pulled-back coarse covector does not
increase.  The locality comparison uses the separate annular estimates.
\end{proof}

\subsection{Robustness and the finite one-step card}

\begin{proposition}[Physical floor under perturbation]
\label{prop:t10v9-floor-perturbation}
Suppose a perturbed physical Schur block satisfies
\begin{equation}
 \widetilde{\overline D}=\overline D+E_D,
 \qquad \overline D\succeq d_0I,
 \qquad \|E_D\|\leq\epsilon_D<d_0,
 \label{eq:t10v9-D-perturbation}
\end{equation}
and $\widetilde{\overline C}>0$.  Then, for arbitrary
$\widetilde{\overline B}$,
\begin{equation}
 \widetilde S_{\rm red}
 =\widetilde{\overline D}
 +\widetilde{\overline B}^*
 \widetilde{\overline C}^{-1}
 \widetilde{\overline B}
 \succeq(d_0-\epsilon_D)I.
 \label{eq:t10v9-robust-floor}
\end{equation}
\end{proposition}

\begin{proof}
The first term has the displayed lower bound and the second is positive
semidefinite.
\end{proof}

\begin{firewall}[Stationarity mismatch]
\label{fire:t10v9-stationarity-mismatch}
Proposition~\ref{prop:t10v9-floor-perturbation} assumes that the perturbed
scalar equation remains the stationarity equation of the same perturbed
quadratic form.  If the imposed Einstein constraint differs from that
stationarity equation, the V99 response-mismatch and gradient debits must be
subtracted from the floor.
\end{firewall}

\begin{definition}[Reduced gravitational Gaussian step card]
\label{def:t10v9-Gaussian-card}
One regulator block passes the card if:
\begin{enumerate}
\item V4 based interior gauge quotient and V7 reflection trace retention are
implemented before any inverse;
\item $A>0$, $\overline C>0$, and
$\overline D\succeq d_0I$ in the matched physical norms;
\item the scalar constraint is exactly
$D_s\overline Q=0$, or the V99 mismatch debit is below $d_0$;
\item all determinant and coarea factors in
\eqref{eq:t10v9-constrained-density} have their declared owners;
\item the two half blocks and boundary density satisfy the V96 reflected
factorization;
\item the physical coarse map is surjective after zero/closing modes are
retained and has the matched conditional data above;
\item the V6--V8 multiscale and curved-frame owner bounds hold;
\item the test-covariance estimate
\eqref{eq:t10v9-covariance-test-bound} is uniform.
\end{enumerate}
\end{definition}

\begin{theorem}[One-step implication]
\label{thm:t10v9-one-step-implication}
If Definition~\ref{def:t10v9-Gaussian-card} holds, the finite reduced
bosonic Gaussian step is normalizable, reflection positive, positive on its
physical precision space, contractive on covariance-matched Wick chaoses,
and equicontinuous on physical smeared fields.  These conclusions are
uniform when the displayed card constants are uniform.
\end{theorem}

\begin{proof}
Normalizability and positivity are
Theorem~\ref{thm:t10v9-constrained-density}.  Reflection positivity is
Proposition~\ref{prop:t10v9-reflected-square}.  Matched contraction is
Theorem~\ref{thm:t10v9-matched-contraction}, and field equicontinuity is
Proposition~\ref{prop:t10v9-subgaussian-fields}.  The remaining card rows
ensure that the quotient and locality operations used to form these
objects are legitimate.
\end{proof}

\begin{assessment}[Status after tenth-series V9]
\label{assess:t10v9}
V9 assembles the first complete abstract Gaussian gravitational block step.
It proves the nested completed square, stationary conformal reduction,
positive physical precision, exact determinant/coarea density, reflected
square, covariance-matched contraction, and uniform smeared-field tail
criterion.  The construction is still conditional on the crucial physical
matrix inequality $\overline D>0$ and on deriving the scalar stationarity
constraint from a concrete regulated Einstein action with matter.
\end{assessment}

\begin{programme}[Tenth-series V10: audit and nonlinear one-step remainder]
\label{prog:t10v9-V10}
V10 will perform the compulsory YM/NS audit, then attach the V5 stationary
cubic remainder and V8 curved-frame debits to the V9 Gaussian step.  It will
derive a single quantitative small-field inequality preserving the physical
precision floor, matched contraction, field-tail card, and reflected
factorization.
\end{programme}

\begin{firewall}[Claim boundary after tenth-series V9]
\label{fire:t10v9-boundary}
V9 is an abstract finite Gaussian theorem.  It does not verify the physical
block matrices for a concrete Einstein regulator, control the nonlinear
Einstein--SM remainder, renormalize the stress tensor, prove chiral
determinant-line gluing, complete a cofinal RG flow, or construct the full
SM--Einstein continuum.
\end{firewall}

\paragraph{Parent-version record.}
V9 was formed from
\texttt{Renewal\_SM\_Einstein\_Continuum\_Research\_Notes\_V8.tex}.
The V8 parent had $3\,470$ validator-counted source lines, $130\,755$ bytes,
and SHA--256
\[
 \texttt{A3C950674306B7DB8AE8B5B47A14336F4A98BE7617B62604340E452D5227EEFA}.
\]

% =============================================================================
% END APPENDED TENTH-SERIES V9 MATERIAL
% =============================================================================

% =============================================================================
% APPENDED TENTH-SERIES V10 MATERIAL
% =============================================================================

\section{Tenth-series V10: a nonlinear reduced block gate and the resonant
stress-source obstruction}
\label{sec:v10v10}

\subsection{Compulsory companion audit}

At the start of this version the newest companion files were
\texttt{Renewal\_Yang\_Mills\_Mass\_Gap\_Research\_Notes\_V27.tex} and
\texttt{Renewal\_NS\_Regularity\_Research\_Notes\_V26.tex}.  The files had,
respectively, $437\,295$ and $278\,283$ bytes and SHA--256 values
\[
 \texttt{1D7223E529700683043ED2FC6CDC50FA57CD5ADFA1B9E0ECB946E054C6EDD907}
\]
and
\[
 \texttt{82C887F8C2C69E4F28FAD7C3DC8DE5B9099CB5A4BF431EBA1FFF3E91935978AD}.
\]
The NS file is unchanged from the V5 audit.  The new YM material proves
annular Mikhlin and packet estimates for the direction-dependent Hodge
projection and uses slowly varying, summable cube weights to retain a strict
linear contraction margin.  It also isolates the remaining nonlinear
high--high to low paraproduct channel.  This is directly relevant here: the
V6--V8 curved gravitational Hodge frame can inherit the weighted packet
argument, but the Einstein stress source cannot be put into the same estimate
before its resonant low output is treated separately.

\begin{firewall}[Audit transfer rule]
\label{fire:t10v10-audit-transfer}
This audit transfers two structural facts, not a Yang--Mills conclusion.  A
small commutator with a slowly varying packet weight preserves a strict
linear contraction.  In contrast, no annular input weight by itself excludes
an arbitrarily low output from two comparable high shells.  The latter fact
will be proved below and entered as a separate Einstein--SM source debit.
\end{firewall}

\subsection{Nonlinear stationary conformal reduction}

Let $M$ be the finite physical boundary/matter space and $S$ the retained
conformal scalar space of one regulator block.  All adjoints and operator
norms in this section are taken in the matched finite-dimensional Hilbert
norms.  Consider a $C^3$ real action $V(m,s)$ near the origin.  We use the
sign convention of V9: $s$ is a stationary negative direction, so
$-D^2_{ss}V$ is positive.

\begin{proposition}[Nonlinear stationary graph]
\label{prop:t10v10-stationary-graph}
Suppose that on a convex chart $\mathcal U\subset M\oplus S$,
\begin{equation}
 -D^2_{ss}V(m,s)\succeq c_s I,
 \qquad c_s>0,
 \label{eq:t10v10-scalar-floor}
\end{equation}
and that $D_sV(0,0)=0$.  After shrinking the $m$ chart if necessary, there
is a unique $C^2$ stationary graph $s=\sigma(m)$ through the origin with
$D_sV(m,\sigma(m))=0$.  Its derivative is
\begin{equation}
 D\sigma(m)=-\bigl(D^2_{ss}V\bigr)^{-1}D^2_{sm}V
 \quad\hbox{at }(m,\sigma(m)).
 \label{eq:t10v10-stationary-derivative}
\end{equation}
For $V_{\rm red}(m):=V(m,\sigma(m))$ one has the exact Hessian identity
\begin{equation}
 D^2V_{\rm red}
 =D^2_{mm}V-D^2_{ms}V\bigl(D^2_{ss}V\bigr)^{-1}D^2_{sm}V.
 \label{eq:t10v10-nonlinear-Schur}
\end{equation}
\end{proposition}

\begin{proof}
The derivative $D_s(D_sV)=D^2_{ss}V$ is invertible by
\eqref{eq:t10v10-scalar-floor}; the implicit-function theorem gives a local
stationary graph.  Strict concavity in $s$ makes a stationary point on each
connected fibre unique.  Differentiating $D_sV(m,\sigma(m))=0$ gives
\eqref{eq:t10v10-stationary-derivative}.  The first derivative of the reduced
action is $D_mV$, because the term $D_sV\,D\sigma$ vanishes on the graph.
Differentiating once more and substituting
\eqref{eq:t10v10-stationary-derivative} proves
\eqref{eq:t10v10-nonlinear-Schur}.
\end{proof}

\begin{remark}[The sign of the Schur term]
If $C':=-D^2_{ss}V>0$ and $B':=D^2_{sm}V$, then
\eqref{eq:t10v10-nonlinear-Schur} reads
\[
 D^2V_{\rm red}=D^2_{mm}V+(B')^*(C')^{-1}B'.
\]
Thus stationary elimination of the negative scalar direction adds a positive
term.  Integrating that direction on a real contour would diverge and is not
an alternative proof of this formula.
\end{remark}

\begin{lemma}[Quantitative nonlinear Schur debit]
\label{lem:t10v10-Schur-debit}
Let the V9 reference blocks obey
\[
 C\succeq c_sI,
 \qquad S_0=A+B^*C^{-1}B\succeq d_0I.
\]
At a point of the stationary graph write
\[
 D^2_{mm}V=A+E_A,
 \quad D^2_{sm}V=B+E_B,
 \quad -D^2_{ss}V=C+E_C.
\]
If
\[
 \|E_A\|\leq a,
 \quad \|E_B\|\leq e,
 \quad \|E_C\|\leq c<c_s,
 \quad \|B\|\leq b,
\]
then
\begin{equation}
 \|D^2V_{\rm red}-S_0\|
 \leq
 \Delta_{\rm Schur}:=
 a+\frac{2be+e^2}{c_s-c}
   +\frac{b^2c}{c_s(c_s-c)}.
 \label{eq:t10v10-Schur-debit}
\end{equation}
Consequently
\begin{equation}
 D^2V_{\rm red}\succeq(d_0-\Delta_{\rm Schur})I.
 \label{eq:t10v10-Schur-floor}
\end{equation}
\end{lemma}

\begin{proof}
Set $C'=C+E_C$ and $B'=B+E_B$.  The Neumann bound and the resolvent
identity give
\[
 \|(C')^{-1}\|\leq(c_s-c)^{-1},
 \qquad
 \|(C')^{-1}-C^{-1}\|
 \leq\frac{c}{c_s(c_s-c)}.
\]
Add and subtract $B^*(C')^{-1}B$ in
$(B')^*(C')^{-1}B'-B^*C^{-1}B$.  The two cross terms and the quadratic
term in $E_B$ have total norm at most
$(2be+e^2)/(c_s-c)$, while the inverse perturbation contributes at most
$b^2c/[c_s(c_s-c)]$.  Adding $E_A$ proves
\eqref{eq:t10v10-Schur-debit}.  The spectral lower bound follows from the
min--max principle.
\end{proof}

\begin{firewall}[What the debit assumes]
\label{fire:t10v10-Schur-debit-scope}
The estimate controls a stationary graph of the same action.  If a regulated
Einstein constraint is imposed from a different functional, its mismatch is
an additional debit.  The estimate also does not prove the reference physical
floor $d_0>0$; that remains a concrete-symbol obligation.
\end{firewall}

\subsection{Strong convexity gives a nonlinear field-tail card}

The Gaussian covariance argument of V9 can be extended without pretending
that the cubic remainder is Gaussian.  The needed finite-dimensional fact is
recorded with its proof.

\begin{lemma}[Brascamp--Lieb variance bound]
\label{lem:t10v10-BL}
Let $d\mu=Z^{-1}e^{-W(x)}dx$ on a finite-dimensional real Hilbert space,
where $W$ is smooth, confining, and
\begin{equation}
 D^2W(x)\succeq H>0
 \quad\hbox{for every }x.
 \label{eq:t10v10-global-convexity}
\end{equation}
Then every smooth $f$ satisfies
\begin{equation}
 {\rm Var}_\mu(f)
 \leq\int
 \langle Df,(D^2W)^{-1}Df\rangle\,d\mu
 \leq\int\langle Df,H^{-1}Df\rangle\,d\mu.
 \label{eq:t10v10-BL}
\end{equation}
\end{lemma}

\begin{proof}
It is enough first to take smooth data with sufficient decay.  Let
$L=-\Delta+DW\cdot D$ and solve
$Lu=f-\int f\,d\mu$ in the mean-zero subspace; a standard coercive
approximation supplies the solution.  Weighted integration by parts gives
\[
 {\rm Var}_\mu(f)=\int Df\cdot Du\,d\mu.
\]
Cauchy--Schwarz in the pointwise metrics $D^2W$ and $(D^2W)^{-1}$ bounds the
right-hand side by
\[
 \left(\int\langle Df,(D^2W)^{-1}Df\rangle d\mu\right)^{1/2}
 \left(\int\langle Du,D^2W Du\rangle d\mu\right)^{1/2}.
\]
The weighted Bochner identity is
\[
 \int(Lu)^2d\mu
 =\int\|D^2u\|_{\rm HS}^2d\mu
  +\int\langle Du,D^2W Du\rangle d\mu.
\]
Its left side is ${\rm Var}_\mu(f)$.  Cancelling the square root of that
variance proves the first inequality.  The second follows from
$(D^2W)^{-1}\preceq H^{-1}$.  Approximation removes the temporary smoothness
and decay restrictions.
\end{proof}

\begin{proposition}[Uniform nonlinear sub-Gaussian fields]
\label{prop:t10v10-nonlinear-subgaussian}
Under the hypotheses of Lemma~\ref{lem:t10v10-BL}, let $X$ have law $\mu$.
For every linear functional $\ell$ and every $t\in\mathbb R$,
\begin{equation}
 \mathbb E\exp\bigl(t\ell(X-\mathbb EX)\bigr)
 \leq
 \exp\left(\frac{t^2}{2}\langle\ell,H^{-1}\ell\rangle\right).
 \label{eq:t10v10-subgaussian}
\end{equation}
Hence
\begin{equation}
 \mathbb P\{|\ell(X-\mathbb EX)|\geq u\}
 \leq2\exp\left(-\frac{u^2}
 {2\langle\ell,H^{-1}\ell\rangle}\right).
 \label{eq:t10v10-tail}
\end{equation}
\end{proposition}

\begin{proof}
Let $F(t)=\log\mathbb E e^{t\ell(X-\mathbb EX)}$.  Exponential tilting
replaces $W$ by $W-t\ell$ and therefore leaves its Hessian unchanged.
Lemma~\ref{lem:t10v10-BL}, applied to every tilted law, gives
$F''(t)\leq\langle\ell,H^{-1}\ell\rangle$.  Since $F(0)=F'(0)=0$, two
integrations prove \eqref{eq:t10v10-subgaussian}.  Chernoff optimization for
$\ell$ and $-\ell$ gives \eqref{eq:t10v10-tail}.
\end{proof}

\begin{corollary}[Application to the reduced block]
\label{cor:t10v10-reduced-tail}
Suppose the reduced, gauge-quotiented, regulated physical action $W$ obeys
the global chart estimate
\begin{equation}
 D^2W\succeq
 \bigl(d_0-\Delta_{\rm Schur}-\Delta_{\rm con}
              -\Delta_{\rm frame}\bigr)I=:d_*I>0.
 \label{eq:t10v10-total-floor}
\end{equation}
Here $\Delta_{\rm con}$ is the declared stationarity/connection-gradient
debit and $\Delta_{\rm frame}$ is the V8 curved-frame debit.  Then every
normalized smeared physical field has the bounds
\eqref{eq:t10v10-subgaussian}--\eqref{eq:t10v10-tail} with $H=d_*I$.
Uniformity of $d_*$ and of the smearing norms gives the equicontinuity row
needed by the cofinal compactness argument.
\end{corollary}

\begin{firewall}[Small-field charts and large fields]
\label{fire:t10v10-large-fields}
A Hessian estimate only on a ball does not imply the global tilted estimate
used above.  The card is valid either for a globally convex regulated block,
or after a convex small-field localization whose complement has a separately
proved large-field estimate.  Calling a local Taylor bound a global
Brascamp--Lieb bound would leave a gap.
\end{firewall}

\subsection{Weighted linear contraction and nonlinear derivative debit}

Let $\mathcal X$ be the physical packet-owner Banach space after the V6--V8
annular decomposition, and let $w_R$ be a positive cube weight whose relative
change across the finite interaction stencil is at most $\omega_R$, with
$\omega_R\to0$ as $R\to\infty$.  Write $\|\cdot\|_{w_R}$ for the conjugated
weighted norm.

\begin{lemma}[Weight-commutator margin]
\label{lem:t10v10-weight-margin}
Let $L$ be a finite-stencil packet matrix with
$\|L\|\leq q_0<1$ in the matched unweighted norm.  If
\begin{equation}
 \|w_RLw_R^{-1}-L\|\leq C_L\omega_R,
 \label{eq:t10v10-weight-commutator}
\end{equation}
then
\begin{equation}
 \|L\|_{w_R}\leq q_R:=q_0+C_L\omega_R.
 \label{eq:t10v10-weighted-linear}
\end{equation}
In particular, $q_R<1$ for all sufficiently large $R$.
\end{lemma}

\begin{proof}
Conjugation identifies the weighted operator norm with
$\|w_RLw_R^{-1}\|$.  The triangle inequality and
\eqref{eq:t10v10-weight-commutator} prove the result.
\end{proof}

\begin{lemma}[Nonlinear contraction gate]
\label{lem:t10v10-nonlinear-contraction}
Let $\mathcal T:B_r(0)\subset\mathcal X\to\mathcal X$ be $C^1$ and suppose
\begin{equation}
 \|D\mathcal T(0)\|_{w_R}\leq q_R,
 \qquad
 \sup_{x\in B_r(0)}
 \|D\mathcal T(x)-D\mathcal T(0)\|_{w_R}
 \leq\epsilon_{\rm nl}(r).
 \label{eq:t10v10-derivative-debit}
\end{equation}
If
\begin{equation}
 q_0+C_L\omega_R+\epsilon_{\rm nl}(r)<1,
 \label{eq:t10v10-contraction-gate}
\end{equation}
then $\mathcal T$ is a strict contraction on every convex subset of
$B_r(0)$ that it maps into itself.
\end{lemma}

\begin{proof}
The two derivative bounds give
$\sup_{B_r}\|D\mathcal T\|_{w_R}\leq
q_0+C_L\omega_R+\epsilon_{\rm nl}(r)$.  Integrate the derivative along the
line segment joining two points of the convex set.
\end{proof}

For a Taylor remainder with $D^2\mathcal T$ bounded by $L_{\mathcal T}$,
one may take $\epsilon_{\rm nl}(r)=L_{\mathcal T}r$.  The V8 range-projector
gap contributes a further explicit term
\begin{equation}
 \epsilon_{\rm frame}
 \leq C_{\mathcal T}\frac{2\theta}{1-\theta},
 \qquad 0\leq\theta<1,
 \label{eq:t10v10-frame-contraction-debit}
\end{equation}
whenever the RG derivative is Lipschitz in the projector.  Thus the actual
card uses the sum of
\eqref{eq:t10v10-contraction-gate} and
\eqref{eq:t10v10-frame-contraction-debit}.

\begin{firewall}[Contracted and retained sectors]
\label{fire:t10v10-retained-sector}
Lemma~\ref{lem:t10v10-nonlinear-contraction} applies to the irrelevant
physical packet sector.  Relevant and marginal couplings, shared boundary
frames, closing modes, and calibrated background sources are finite retained
coordinates.  Including any of them in the strict-contraction claim would
contradict their role in the RG chart.
\end{firewall}

\subsection{The exact high--high to low obstruction}

For integers $j$, let
\[
 \mathcal A_j=\{\xi:2^{j-1}\leq|\xi|\leq2^{j+1}\}.
\]

\begin{lemma}[Shell-triad geometry]
\label{lem:t10v10-shell-triad}
Suppose $\xi\in\mathcal A_j$, $\eta\in\mathcal A_k$, and
$\xi+\eta\in\mathcal A_\ell$.
\begin{enumerate}
\item If $j\geq k+4$, then $j-2\leq\ell\leq j+2$.
\item If $\ell\leq\max\{j,k\}-3$, then $|j-k|\leq3$.
\item For every $j-\ell$ sufficiently large there are
$\xi,\eta\in\mathcal A_j$ with $\xi+\eta\in\mathcal A_\ell$.
\end{enumerate}
\end{lemma}

\begin{proof}
For $j\geq k+4$,
\[
 |\xi+\eta|\geq2^{j-1}-2^{k+1}
 \geq2^{j-2},
 \qquad
 |\xi+\eta|\leq2^{j+1}+2^{k+1}<2^{j+2}.
\]
This proves the first assertion up to harmless endpoint overlap of the chosen
annuli.  Its contrapositive after interchanging $j,k$ proves the second.
For the third, choose orthogonal unit vectors $e_1,e_2$ and set
\[
 \xi=2^je_1,
 \qquad \eta=-2^je_1+2^\ell e_2.
\]
When $j-\ell$ is large, both inputs lie in $\mathcal A_j$ and their sum is
$2^\ell e_2\in\mathcal A_\ell$.
\end{proof}

\begin{corollary}[Paraproduct separation]
\label{cor:t10v10-paraproduct-separation}
A quadratic interaction with widely separated input shells has output at the
higher shell and is compatible with the ordinary owner estimate.  An output
arbitrarily below both inputs occurs only in the comparable-shell resonant
channel.  Therefore that channel must be isolated before summing infrared
shell weights.
\end{corollary}

\begin{firewall}[Why Ward transversality is not yet enough]
\label{fire:t10v10-Ward-not-enough}
The geometric Ward identity makes the stress source covariantly conserved.
Conservation constrains its contraction with the output momentum; it does not
force the source itself to vanish at zero momentum.  Comparable high packets
may therefore create a monopole or dipole low-frequency stress component.
No factor $2^{-j}$ may be credited to that component until a subtraction,
moment cancellation, or concrete tensor-symbol identity has been proved.
\end{firewall}

\subsection{One nonlinear reduced Einstein--SM block card}

\begin{definition}[Unified tenth-series V10 gate]
\label{def:t10v10-unified-gate}
One finite reduced Einstein--SM block passes the V10 gate when all of the
following hold.
\begin{enumerate}
\item The V9 Gaussian block card holds with $C\succeq c_sI$ and
$S_0\succeq d_0I$.
\item The nonlinear scalar equation is the stationarity equation of the same
action, Proposition~\ref{prop:t10v10-stationary-graph} applies, and the block
errors satisfy Lemma~\ref{lem:t10v10-Schur-debit}.
\item All connection, constraint-mismatch, and curved-frame precision debits
are declared and
\begin{equation}
 \Delta_{\rm Schur}+\Delta_{\rm con}+\Delta_{\rm frame}<d_0.
 \label{eq:t10v10-precision-gate}
\end{equation}
\item On the irrelevant packet sector,
\begin{equation}
 q_0+C_L\omega_R+\epsilon_{\rm nl}(r)
 +C_{\mathcal T}\frac{2\theta}{1-\theta}<1.
 \label{eq:t10v10-full-contraction-gate}
\end{equation}
\item The reduced regulator is globally strongly convex as in
\eqref{eq:t10v10-total-floor}, or its large-field complement has an
independent uniform estimate.
\item The nonlinear density has the exact reflected form
\begin{equation}
 e^{-W}=e^{-W_0}\,e^{-W_+}\,e^{-\Theta W_+},
 \qquad e^{-W_0}\geq0,
 \label{eq:t10v10-reflected-nonlinear-form}
\end{equation}
with the gauge, coarea, determinant, and anomaly factors assigned to their
actual half or boundary owners.
\item The quadratic nonlinearity is paraproduct-split, and the comparable
high--high to low stress source is placed in a named retained/background
sector until its required cancellations are proved.
\end{enumerate}
\end{definition}

\begin{theorem}[Consequences of the unified gate]
\label{thm:t10v10-unified-gate}
If Definition~\ref{def:t10v10-unified-gate} holds uniformly over regulator
blocks, then the one-step reduced measure has:
\begin{enumerate}
\item a positive physical Hessian floor $d_*>0$;
\item a strict nonlinear contraction on its irrelevant packet chart;
\item uniform sub-Gaussian tails for normalized smeared physical fields; and
\item reflection positivity for observables supported in the positive half.
\end{enumerate}
These conclusions concern one regulated step and are stable under any
cofinal family for which the gate constants are uniform.
\end{theorem}

\begin{proof}
Items one and three follow from Lemma~\ref{lem:t10v10-Schur-debit} and
Corollary~\ref{cor:t10v10-reduced-tail}.  Item two follows from
Lemma~\ref{lem:t10v10-nonlinear-contraction} with the curved-frame term in
\eqref{eq:t10v10-frame-contraction-debit}.  For the last item, a positive-half
observable $F$ gives, using \eqref{eq:t10v10-reflected-nonlinear-form},
\[
 \int F\,\Theta F\,e^{-W}
 =\int_{\rm boundary}e^{-W_0}
   \left|\int_{+}F e^{-W_+}\right|^2\geq0,
\]
where the equality includes the declared half-density and gauge factors.
The resonant source is not estimated by the contracted norm, as required by
the last row of the definition.
\end{proof}

\begin{assessment}[Progress and remaining bottleneck after V10]
\label{assess:t10v10}
V10 replaces the phrase ``small nonlinear remainder'' by two checkable
inequalities: the physical precision inequality
\eqref{eq:t10v10-precision-gate} and the packet contraction inequality
\eqref{eq:t10v10-full-contraction-gate}.  It proves the nonlinear stationary
Schur formula, an explicit Schur debit, and a non-Gaussian field-tail theorem.
The compulsory YM audit also exposes a genuine upstream bottleneck: the
comparable high--high channel of the Einstein--SM stress source can feed the
deep infrared.  Linear Hodge contraction does not close that channel.
\end{assessment}

\begin{programme}[Tenth-series V11: calibrated stress moments]
\label{prog:t10v10-V11}
V11 will analyze the low-frequency moments of a localized stress source.  It
will distinguish what follows from conservation from what requires vacuum or
background calibration, prove the precise Fourier gain obtained from
monopole and dipole cancellation, and determine which cancellation is needed
before applying the inverse gravitational symbol.  This goes upstream of the
current resonant paraproduct bottleneck.
\end{programme}

\begin{firewall}[Claim boundary after tenth-series V10]
\label{fire:t10v10-boundary}
The unified gate is conditional.  This version does not verify its constants
for a concrete regulated Einstein--Standard-Model action, prove cancellation
of the resonant stress moments, construct the chiral determinant line, close
the relevant/marginal RG flow, or construct the full SM--Einstein continuum.
\end{firewall}

\paragraph{Parent-version record.}
V10 was formed from
\texttt{Renewal\_SM\_Einstein\_Continuum\_Research\_Notes\_V9.tex}.
The V9 parent had $3\,932$ validator-counted source lines, $147\,504$ bytes,
and SHA--256
\[
 \texttt{46E64061260536B634E30D61A68220BCE3DB55496812343AD0EBA9C8E5C7B594}.
\]

% =============================================================================
% END APPENDED TENTH-SERIES V10 MATERIAL
% =============================================================================

% =============================================================================
% APPENDED TENTH-SERIES V11 MATERIAL
% =============================================================================

\section{Tenth-series V11: stress moments, Ward cancellation, and infrared
calibration}
\label{sec:v10v11}

V10 isolated the comparable-shell high--high channel because two large
opposite momenta can produce an arbitrarily small output momentum.  This
version goes upstream: it asks which low moments of the resulting
Einstein--SM stress source vanish by an exact Ward identity, which vanish
only after a calibration, and which are merely zero in expectation.  These
three statements have different analytic consequences.

\subsection{Moment cancellation and Fourier gain}

We use the Fourier convention
\[
 \widehat f(k)=\int_{\mathbb R^d}e^{-ik\cdot x}f(x)\,dx.
\]
The following statements apply componentwise to finite-dimensional tensor
values.

\begin{lemma}[Exact moment-to-Fourier estimates]
\label{lem:t10v11-moment-Fourier}
Let $f\in L^1(\mathbb R^d)$.
\begin{enumerate}
\item If $\int f=0$ and $\int |x|\,|f(x)|dx=M_1<\infty$, then
\begin{equation}
 |\widehat f(k)|\leq |k|M_1.
 \label{eq:t10v11-one-moment}
\end{equation}
\item If in addition $\int x_af(x)dx=0$ for every $a$ and
$\int |x|^2|f(x)|dx=M_2<\infty$, then
\begin{equation}
 |\widehat f(k)|\leq\frac12|k|^2M_2.
 \label{eq:t10v11-two-moments}
\end{equation}
\end{enumerate}
\end{lemma}

\begin{proof}
Under the first cancellation,
\[
 \widehat f(k)=\int(e^{-ik\cdot x}-1)f(x)dx.
\]
The inequality $|e^{-it}-1|\leq|t|$ proves
\eqref{eq:t10v11-one-moment}.  Under both cancellations, subtract also the
linear Taylor term:
\[
 \widehat f(k)=\int(e^{-ik\cdot x}-1+ik\cdot x)f(x)dx.
\]
Taylor's formula with integral remainder gives
$|e^{-it}-1+it|\leq t^2/2$, proving
\eqref{eq:t10v11-two-moments}.
\end{proof}

\begin{lemma}[Infrared inverse criterion]
\label{lem:t10v11-inverse-criterion}
Suppose $|\widehat f(k)|\leq C|k|^m$ for $0<|k|\leq1$, and let a multiplier
$G$ satisfy $\|G(k)\|\leq C_G|k|^{-r}$.  Then
\begin{equation}
 \int_{|k|\leq1}|G(k)\widehat f(k)|^2dk<\infty
 \quad\hbox{provided}\quad d+2m-2r>0.
 \label{eq:t10v11-inverse-integrability}
\end{equation}
If $m\geq r$, the transformed Fourier field is also bounded at $k=0$.
\end{lemma}

\begin{proof}
Polar coordinates bound the integral by a constant times
\[
 \int_0^1\rho^{d-1+2m-2r}d\rho,
\]
which converges exactly under the displayed strict inequality.  The pointwise
claim follows directly from $|G(k)\widehat f(k)|\leq CC_G|k|^{m-r}$.
\end{proof}

For the order-two gravitational inverse on a three-dimensional constraint
slice, one cancelled moment ($m=1$) already gives local $L^2$ integrability,
whereas two cancelled moments ($m=2$) give a bounded Fourier response.  With
no cancelled moment, the $L^2$ inverse diverges at the origin.  Thus the
monopole and dipole rows of a block source cannot be suppressed as cosmetic
normalizations.

\subsection{What exact spacetime conservation really proves}

\begin{theorem}[Conserved symmetric compact stress has two vanishing moments]
\label{thm:t10v11-stress-moments}
Let $T=(T_{\mu\nu})$ be a symmetric tensor on $\mathbb R^d$ with
\begin{equation}
 T_{\mu\nu}\in L^1,
 \qquad
 \int |x|^2|T(x)|dx<\infty,
 \qquad
 \partial^\mu T_{\mu\nu}=0
 \label{eq:t10v11-stress-hypotheses}
\end{equation}
in distributions.  Then for all indices
\begin{equation}
 \int T_{\mu\nu}(x)dx=0,
 \qquad
 \int x_\alpha T_{\mu\nu}(x)dx=0.
 \label{eq:t10v11-stress-moments}
\end{equation}
Consequently
\begin{equation}
 |\widehat T_{\mu\nu}(k)|
 \leq\frac12|k|^2
 \int |x|^2|T_{\mu\nu}(x)|dx.
 \label{eq:t10v11-stress-quadratic-zero}
\end{equation}
\end{theorem}

\begin{proof}
Use smooth cutoffs tending to one to justify the following integrations by
parts; the moment hypothesis makes the cutoff errors tend to zero.  Multiply
$\partial^\mu T_{\mu\nu}=0$ by $x_\alpha$ and integrate.  This gives
\[
 0=-\int T_{\alpha\nu},
\]
which proves the zeroth-moment assertion.  Multiplication by
$x_\alpha x_\beta$ gives
\begin{equation}
 M_{\beta\alpha\nu}+M_{\alpha\beta\nu}=0,
 \qquad
 M_{\alpha\mu\nu}:=\int x_\alpha T_{\mu\nu}(x)dx.
 \label{eq:t10v11-M-antisymmetry}
\end{equation}
Thus $M$ is antisymmetric in its first two indices.  Symmetry of $T$ makes
$M$ symmetric in its last two indices.  These two symmetries force $M=0$:
\[
 M_{abc}=-M_{bac}=-M_{bca}=M_{cba}=M_{cab}=-M_{acb}=-M_{abc}.
\]
The second conclusion of \eqref{eq:t10v11-stress-moments} follows.  Apply
Lemma~\ref{lem:t10v11-moment-Fourier} to each component to obtain
\eqref{eq:t10v11-stress-quadratic-zero}.
\end{proof}

\begin{corollary}[Bounded response to a conserved localized source]
\label{cor:t10v11-bounded-stress-response}
If, in addition, the gauge-fixed gravitational Green symbol on the physical
source subspace satisfies $\|G(k)\|\leq C|k|^{-2}$, then
$G(k)\widehat T(k)$ is bounded near $k=0$ by the second moment of $T$.
\end{corollary}

\begin{proof}
Combine \eqref{eq:t10v11-stress-quadratic-zero} with the multiplier bound.
\end{proof}

This theorem is the cancellation needed by the V10 resonant channel, but its
hypotheses are stronger than the formal phrase ``stress conservation.''  A
finite spacetime block has cut faces, an off-shell field does not satisfy its
equations of motion, and gauge fixing and density factors contribute their
own Ward rows.

\begin{proposition}[Full Ward identity versus a constraint slice]
\label{prop:t10v11-slice-distinction}
If $\widehat T$ is continuous at the full spacetime origin and
$k^\mu\widehat T_{\mu\nu}(k)=0$ for every spacetime $k$, then
$\widehat T_{\mu\nu}(0)=0$.  This conclusion need not hold for the spatial
zero mode of $T_{00}$ on a fixed time slice.
\end{proposition}

\begin{proof}
Set $k=te_\alpha$ with $t\neq0$.  The Ward identity gives
$\widehat T_{\alpha\nu}(te_\alpha)=0$; continuity and $t\to0$ prove the first
claim.  For the second, take a time-independent tensor with
$T_{00}(t,x)=\rho(x)$ and all other components zero.  It is conserved, but
its spatial zero mode is $\int\rho$, which may be nonzero.  The example is not
integrable in the full time direction, precisely the hypothesis used by the
first claim and by Theorem~\ref{thm:t10v11-stress-moments}.
\end{proof}

\begin{firewall}[Off-shell diffeomorphism identity]
\label{fire:t10v11-off-shell-Ward}
For a matter action, the local diffeomorphism identity has the schematic
off-shell form
\[
 \nabla^\mu T_{\mu\nu}
 =\sum_A E_A(\phi)\,\nabla_\nu\phi^A
   +\hbox{spin/gauge equation terms}.
\]
Configurationwise conservation follows only after the appropriate field
equations or an exact Schwinger--Dyson integration.  An averaged Ward identity
does not by itself supply the deterministic moment bound required by a
pointwise block inverse.
\end{firewall}

\subsection{Localization creates the missing boundary source}

Let $\chi$ be a smooth block cutoff and put $T^\chi=\chi T$.  Even if $T$ is
exactly conserved,
\begin{equation}
 \partial^\mu T^\chi_{\mu\nu}
 =g_\nu:=(\partial^\mu\chi)T_{\mu\nu}.
 \label{eq:t10v11-cutoff-divergence}
\end{equation}
The defect is supported in the cutoff collar.

\begin{lemma}[Localized monopole is a boundary-flux moment]
\label{lem:t10v11-localized-monopole}
Under sufficient integrability,
\begin{equation}
 \int T^\chi_{\alpha\nu}(x)dx
 =-\int x_\alpha g_\nu(x)dx.
 \label{eq:t10v11-monopole-defect}
\end{equation}
Moreover, if
$M^\chi_{\alpha\beta\nu}:=\int x_\alpha
T^\chi_{\beta\nu}$, then
\begin{equation}
 M^\chi_{\beta\alpha\nu}+M^\chi_{\alpha\beta\nu}
 =-\int x_\alpha x_\beta g_\nu(x)dx.
 \label{eq:t10v11-dipole-defect}
\end{equation}
\end{lemma}

\begin{proof}
Integrate \eqref{eq:t10v11-cutoff-divergence} against $x_\alpha$ and then
against $x_\alpha x_\beta$.  Integration by parts gives the two formulas.
\end{proof}

Thus block localization does not destroy the Ward information silently; it
moves it to a collar owner.  The correct completion is a symmetric collar
tensor $B_{\mu\nu}$ satisfying
\begin{equation}
 \partial^\mu B_{\mu\nu}=-g_\nu.
 \label{eq:t10v11-boundary-repair}
\end{equation}
If $T^\chi+B$ has compact support and a finite second moment, then
Theorem~\ref{thm:t10v11-stress-moments} applies and its monopole and dipole
vanish.  Constructing $B$ with a uniform local norm, prescribed reflection
parity, and compatible shared-face traces is a genuine elasticity-complex
problem; it is not supplied by scalar mean subtraction.

\begin{proposition}[Compatibility conditions for a compact divergence repair]
\label{prop:t10v11-repair-compatibility}
If a compactly supported symmetric $B$ solves
$\partial^\mu B_{\mu\nu}=-g_\nu$, then necessarily
\begin{equation}
 \int g_\nu=0,
 \qquad
 \int(x_\alpha g_\nu-x_\nu g_\alpha)=0.
 \label{eq:t10v11-repair-compatibility}
\end{equation}
\end{proposition}

\begin{proof}
Integrating the equation gives $\int g_\nu=0$.  Multiply its $\nu$ component
by $x_\alpha$, integrate, and subtract the same identity with $\alpha$ and
$\nu$ exchanged.  The right side is
\[
 \int(B_{\alpha\nu}-B_{\nu\alpha})=0
\]
by symmetry, proving the angular-moment condition.
\end{proof}

\begin{firewall}[Compatibility is not existence]
\label{fire:t10v11-repair-existence}
Proposition~\ref{prop:t10v11-repair-compatibility} proves only necessary
conditions.  A reflection-compatible, finite-support right inverse of the
symmetric divergence with uniform rescaled bounds must still be constructed.
That construction is the next upstream analytic task.
\end{firewall}

\subsection{Algebraic calibration and its exact limits}

There is a simpler finite-rank operation that removes low moments without
claiming to repair the Ward identity.  Choose compactly supported scalar
functions $\rho_0,\rho_1,\ldots,\rho_d$ satisfying
\begin{equation}
 \int\rho_0=1,
 \quad \int x_a\rho_0=0,
 \quad \int\rho_b=0,
 \quad \int x_a\rho_b=\delta_{ab}.
 \label{eq:t10v11-calibration-basis}
\end{equation}
For a scalar or tensor component $f$ define
\begin{equation}
 \Pi_{\leq1}f
 :=\left(\int f\right)\rho_0
   +\sum_{b=1}^d\left(\int x_bf\right)\rho_b,
 \qquad
 f_{\rm cal}:=(I-\Pi_{\leq1})f.
 \label{eq:t10v11-calibration}
\end{equation}

\begin{proposition}[Finite-rank moment calibration]
\label{prop:t10v11-calibration}
The calibrated source obeys
\begin{equation}
 \int f_{\rm cal}=0,
 \qquad \int x_af_{\rm cal}=0,
 \label{eq:t10v11-calibrated-moments}
\end{equation}
and hence has the quadratic Fourier zero
\eqref{eq:t10v11-two-moments} whenever its second absolute moment is finite.
The range of $\Pi_{\leq1}$ has dimension at most $d+1$ per tensor component
and must be carried as a retained background coordinate.
\end{proposition}

\begin{proof}
Insert \eqref{eq:t10v11-calibration-basis} into
\eqref{eq:t10v11-calibration}.  The two moment identities are immediate, and
Lemma~\ref{lem:t10v11-moment-Fourier} gives the Fourier estimate.  The range
statement follows from the displayed spanning family.
\end{proof}

\begin{firewall}[Calibration is not a covariant Ward repair]
\label{fire:t10v11-calibration-firewall}
Componentwise $\Pi_{\leq1}$ need not commute with gauge transformations,
diffeomorphisms, reflection, or the divergence operator.  It is a legitimate
finite-dimensional background split only when those transformation laws and
owners are declared.  The collar repair
\eqref{eq:t10v11-boundary-repair} is the stronger covariant target.
\end{firewall}

\begin{lemma}[Normal ordering removes only the mean]
\label{lem:t10v11-normal-ordering}
Let $F$ be a random integrable source and set $:F:=F-\mathbb EF$.  For any
linear moment functional $M$,
\begin{equation}
 \mathbb E M(:F:)=0,
 \qquad
 {\rm Var}\bigl(M(:F:)\bigr)={\rm Var}\bigl(M(F)\bigr).
 \label{eq:t10v11-normal-ordering-limit}
\end{equation}
Thus normal ordering does not give configurationwise monopole or dipole
cancellation unless the corresponding moment was already deterministic.
\end{lemma}

\begin{proof}
Linearity gives $M(:F:)=M(F)-\mathbb EM(F)$.  Its expectation is zero and
subtraction of a constant leaves variance unchanged.
\end{proof}

\subsection{A calibrated resonant-source card}

Let $\mathcal N_{\rm res}(u,v)$ denote the comparable-shell quadratic stress
source isolated by the V10 paraproduct split.  On each block write
\begin{equation}
 \mathcal N_{\rm res}
 =\Pi_{\leq1}\mathcal N_{\rm res}
  +(I-\Pi_{\leq1})\mathcal N_{\rm res}.
 \label{eq:t10v11-resonant-split}
\end{equation}

\begin{theorem}[Infrared closure after two-moment calibration]
\label{thm:t10v11-calibrated-closure}
Assume the calibrated term in \eqref{eq:t10v11-resonant-split} has a uniform
second absolute moment and the physical gravitational inverse is of order at
most two.  Then its inverse response is bounded at zero momentum, with norm
controlled by that second moment.  All possible unbounded zero-frequency
response is confined to the finite-rank coordinate
$\Pi_{\leq1}\mathcal N_{\rm res}$.
\end{theorem}

\begin{proof}
Proposition~\ref{prop:t10v11-calibration} gives an $O(|k|^2)$ Fourier zero for
the calibrated term.  Corollary~\ref{cor:t10v11-bounded-stress-response}
uses only that zero and the order-two multiplier estimate, so its proof
applies verbatim.  The complementary term lies in the finite-dimensional
range of $\Pi_{\leq1}$.
\end{proof}

\begin{definition}[V11 resonant-source ledger]
\label{def:t10v11-resonant-ledger}
A block records separately:
\begin{enumerate}
\item an exact full-spacetime conserved symmetric source, if available;
\item the cutoff divergence $g$ on every block collar;
\item any symmetric collar repair $B$ and its compatibility data;
\item the monopole/dipole calibration coordinate
$\Pi_{\leq1}\mathcal N_{\rm res}$;
\item whether cancellation is configurationwise, after field equations,
after Schwinger--Dyson averaging, or only in expectation; and
\item the uniform second-moment norm of the calibrated remainder.
\end{enumerate}
Only a configurationwise repaired or calibrated remainder may enter the V10
strict irrelevant contraction norm.  The finite moment coordinate belongs to
the relevant/background flow.
\end{definition}

\begin{assessment}[Progress and bottleneck after V11]
\label{assess:t10v11}
V11 proves that a compact, symmetric, exactly conserved spacetime stress has
both zero monopole and zero dipole, giving precisely the quadratic Fourier
zero needed to cancel an order-two gravitational inverse.  It also shows why
this theorem cannot be applied directly to a cut regulator block: localization
creates a collar divergence whose first two moments reproduce the lost Ward
cancellation.  Algebraic calibration closes the infrared estimate but leaves
a finite relevant/background coordinate and is not itself covariant.
\end{assessment}

\begin{programme}[Tenth-series V12: a symmetric-divergence collar repair]
\label{prog:t10v11-V12}
V12 will construct, first on a star-shaped Euclidean block, a finite-support
right inverse for the symmetric divergence under the force and torque
compatibility conditions \eqref{eq:t10v11-repair-compatibility}.  It will seek
rescaled Sobolev and moment bounds, then impose the V7 reflection parity and
shared-face ownership.  If a local symmetric repair cannot satisfy those
requirements simultaneously, the programme will move upstream to an
elasticity-complex potential and retain its boundary cohomology explicitly.
\end{programme}

\begin{firewall}[Claim boundary after tenth-series V11]
\label{fire:t10v11-boundary}
The stress-moment theorem is a finite-dimensional/localized analytic input.
This version does not construct the collar repair, prove the concrete
Einstein--SM stress satisfies the exact hypotheses off shell, control the
retained calibration flow, verify anomaly cancellation, or construct the
full SM--Einstein continuum.
\end{firewall}

\paragraph{Parent-version record.}
V11 was formed from
\texttt{Renewal\_SM\_Einstein\_Continuum\_Research\_Notes\_V10.tex}.
The V10 parent had $4\,446$ validator-counted source lines, $167\,028$ bytes,
and SHA--256
\[
 \texttt{4D823B4E7471383E5B74DA84D5D17CF8AAD457DB705228C2DEAD498676896C23}.
\]

% =============================================================================
% END APPENDED TENTH-SERIES V11 MATERIAL
% =============================================================================

% =============================================================================
% APPENDED TENTH-SERIES V12 MATERIAL
% =============================================================================

\section{Tenth-series V12: a local symmetric-divergence repair and its
boundary cohomology}
\label{sec:v10v12}

V11 reduced the resonant infrared problem to a concrete question.  Given the
collar defect $g=\operatorname{div}(\chi T)$, can one solve
$\operatorname{div}B=-g$ with $B$ symmetric, supported in the collar, and
bounded uniformly after rescaling?  The answer is yes on a fixed Euclidean
Lipschitz collar, exactly when the force and torque moments vanish.  The proof
also identifies why face-by-face localization can fail even when the full
collar problem is solvable.

\subsection{Korn quotient and the range of symmetric divergence}

Let $\Omega\subset\mathbb R^d$ be a bounded connected Lipschitz domain.  Set
\[
 \varepsilon u:=\frac12(Du+Du^{\mathsf T}),
 \qquad
 \mathcal R(\Omega):=\{x\mapsto a+Qx:a\in\mathbb R^d,
 \ Q^{\mathsf T}=-Q\}.
\]
The space $\mathcal R(\Omega)$ is the kernel of $\varepsilon$ and has
dimension $d(d+1)/2$.

\begin{lemma}[Korn--Poincar\'e quotient estimate]
\label{lem:t10v12-Korn}
There is a constant $C_\Omega$ such that
\begin{equation}
 \inf_{r\in\mathcal R(\Omega)}\|u-r\|_{H^1(\Omega)}
 \leq C_\Omega\|\varepsilon u\|_{L^2(\Omega)}
 \label{eq:t10v12-Korn}
\end{equation}
for every $u\in H^1(\Omega;\mathbb R^d)$.  If
$\Omega_h=x_0+h\Omega$, its scale-separated form is
\begin{equation}
 \inf_{r\in\mathcal R(\Omega_h)}
 \bigl(h^{-1}\|u-r\|_{L^2(\Omega_h)}
       +\|D(u-r)\|_{L^2(\Omega_h)}\bigr)
 \leq C_\Omega\|\varepsilon u\|_{L^2(\Omega_h)}.
 \label{eq:t10v12-scaled-Korn}
\end{equation}
\end{lemma}

\begin{proof}
On the quotient $H^1/\mathcal R$, the right side is a norm.  If
\eqref{eq:t10v12-Korn} failed, there would be quotient-normalized $u_n$ with
$\varepsilon u_n\to0$.  Choose representatives orthogonal to
$\mathcal R$.  Rellich compactness gives an $L^2$ convergent subsequence;
the standard identity controlling all first derivatives by the symmetric
gradient and the $L^2$ norm makes it Cauchy in $H^1$.  Its limit has zero
symmetric gradient, hence lies in $\mathcal R$, contradicting the chosen
orthogonality and normalization.  Translation and the change of variables
$x=x_0+hy$ give \eqref{eq:t10v12-scaled-Korn}.
\end{proof}

\begin{theorem}[Local zero-traction symmetric-divergence inverse]
\label{thm:t10v12-symmetric-divergence}
Let $f\in L^2(\Omega;\mathbb R^d)$ satisfy
\begin{equation}
 \int_\Omega f\cdot r\,dx=0
 \quad\hbox{for every }r\in\mathcal R(\Omega).
 \label{eq:t10v12-rigid-compatibility}
\end{equation}
Then there exists $B\in L^2(\Omega;{\rm Sym}^2\mathbb R^d)$ such that its
extension $\widetilde B$ by zero to $\mathbb R^d$ obeys
\begin{equation}
 \operatorname{div}\widetilde B=\widetilde f
 \quad\hbox{in }\mathcal D'(\mathbb R^d),
 \label{eq:t10v12-divergence-inverse}
\end{equation}
and
\begin{equation}
 \|B\|_{L^2(\Omega)}
 \leq C_\Omega\,{\rm diam}(\Omega)\,\|f\|_{L^2(\Omega)}.
 \label{eq:t10v12-repair-bound}
\end{equation}
Conversely, any symmetric zero-extended $B$ satisfying
\eqref{eq:t10v12-divergence-inverse} forces
\eqref{eq:t10v12-rigid-compatibility}.  Thus the cokernel consists exactly
of rigid motions.
\end{theorem}

\begin{proof}
Work in the Hilbert quotient
$\mathcal H=H^1(\Omega;\mathbb R^d)/\mathcal R(\Omega)$ with norm
$\|\varepsilon u\|_2$.  Compatibility makes
$v\mapsto-\int_\Omega f\cdot v$ a well-defined functional on the quotient.
By \eqref{eq:t10v12-scaled-Korn}, its norm is at most
$C_\Omega{\rm diam}(\Omega)\|f\|_2$.  Lax--Milgram supplies a unique quotient
class $u$ satisfying
\begin{equation}
 \int_\Omega\varepsilon u:\varepsilon v\,dx
 =-\int_\Omega f\cdot v\,dx
 \quad(v\in H^1(\Omega;\mathbb R^d)).
 \label{eq:t10v12-Neumann-elasticity}
\end{equation}
Set $B=\varepsilon u$.  Testing with the restriction of any
$\varphi\in C_c^\infty(\mathbb R^d;\mathbb R^d)$ gives
\[
 \langle\operatorname{div}\widetilde B,\varphi\rangle
 =-\int_\Omega B:D\varphi
 =-\int_\Omega B:\varepsilon\varphi
 =\int_\Omega f\cdot\varphi.
\]
This proves \eqref{eq:t10v12-divergence-inverse}; in particular, the weak
normal traction on the boundary is zero.  Taking $v=u$ in
\eqref{eq:t10v12-Neumann-elasticity} and applying the quotient estimate proves
\eqref{eq:t10v12-repair-bound}.

Conversely, let $r=a+Qx$ be rigid.  Since $Dr=Q$ is antisymmetric and $B$ is
symmetric,
\[
 \int_\Omega f\cdot r
 =\langle\operatorname{div}\widetilde B,r\rangle
 =-\int_\Omega B:Q=0,
\]
using a cutoff equal to one on the bounded support.  This is exactly
\eqref{eq:t10v12-rigid-compatibility}.
\end{proof}

\begin{corollary}[Force and torque are the complete compatibility data]
\label{cor:t10v12-force-torque}
Condition \eqref{eq:t10v12-rigid-compatibility} is equivalent to
\begin{equation}
 \int_\Omega f_i=0,
 \qquad
 \int_\Omega(x_if_j-x_jf_i)=0
 \quad(1\leq i<j\leq d).
 \label{eq:t10v12-force-torque}
\end{equation}
\end{corollary}

\begin{proof}
Pairing with translations gives the first family.  Pairing with the
infinitesimal rotation $x\mapsto Qx$ gives the second, and translations and
rotations span $\mathcal R(\Omega)$.
\end{proof}

\subsection{Application to a cutoff stress}

\begin{proposition}[The full cutoff defect is automatically compatible]
\label{prop:t10v12-cutoff-compatible}
Let $T\in L^2_c(\mathbb R^d;{\rm Sym}^2\mathbb R^d)$ and set
$g=\operatorname{div}T$ in distributions.  If $g\in L^2$ and a bounded
Lipschitz domain $\Omega$ contains its support, then
\begin{equation}
 \int_\Omega g\cdot r=0
 \quad(r\in\mathcal R(\Omega)).
 \label{eq:t10v12-cutoff-compatible}
\end{equation}
In particular this applies to $T^\chi=\chi T_{\rm matter}$ and the V11 collar
defect $g=\operatorname{div}T^\chi$.
\end{proposition}

\begin{proof}
For a rigid motion $r$, distributional integration by parts gives
\[
 \int g\cdot r=-\int T:Dr=0,
\]
because $Dr$ is antisymmetric while $T$ is symmetric.
\end{proof}

\begin{corollary}[Euclidean collar repair]
\label{cor:t10v12-collar-repair}
Suppose the support of the cutoff defect $g$ lies in a connected fixed-shape
collar $\Omega_h=x_0+h\Omega_1$.  There is a symmetric tensor $B$ supported
in $\overline\Omega_h$ such that
\begin{equation}
 \operatorname{div}B=-g,
 \qquad
 \|B\|_{L^2}\leq C_{\Omega_1}h\|g\|_{L^2}.
 \label{eq:t10v12-collar-repair}
\end{equation}
Consequently $T^\chi+B$ is symmetric, compactly supported, and conserved;
if it has a finite second moment, its Fourier transform is $O(|k|^2)$ by
Theorem~\ref{thm:t10v11-stress-moments}.
\end{corollary}

\begin{proof}
Apply Proposition~\ref{prop:t10v12-cutoff-compatible} and
Theorem~\ref{thm:t10v12-symmetric-divergence} to $f=-g$.  Scaling gives the
displayed estimate.  The remaining statements follow by adding the two
divergences and invoking the V11 theorem.
\end{proof}

\begin{lemma}[Second-moment cost of the repair]
\label{lem:t10v12-repair-moment}
If $\Omega_h\subset B(x_0,Ch)$, then
\begin{equation}
 \int_{\Omega_h}|x-x_0|^2|B(x)|dx
 \leq C^2h^2|\Omega_h|^{1/2}\|B\|_{L^2}
 \leq C' h^3|\Omega_h|^{1/2}\|g\|_{L^2}.
 \label{eq:t10v12-repair-moment}
\end{equation}
\end{lemma}

\begin{proof}
The first inequality is Cauchy--Schwarz and the diameter bound.  The second
uses \eqref{eq:t10v12-collar-repair}.
\end{proof}

The collar thickness must remain a fixed positive fraction of the rescaled
block size for $C_{\Omega_1}$ to be uniform.  A collar degenerating to a thin
shell can have a growing Korn constant; V12 does not conceal that geometry in
the symbol $C$.

\subsection{Reflection covariance}

Let $R$ be an orthogonal reflection preserving $\Omega$, and define
\begin{equation}
 (\Theta_1 f)(x)=Rf(Rx),
 \qquad
 (\Theta_2 B)(x)=RB(Rx)R^{\mathsf T}.
 \label{eq:t10v12-reflection-actions}
\end{equation}
These actions intertwine symmetric gradient and divergence.

\begin{proposition}[Reflection-equivariant repair]
\label{prop:t10v12-reflection-repair}
Let $\mathfrak R_\Omega f$ denote the stress $B=\varepsilon u$ constructed in
Theorem~\ref{thm:t10v12-symmetric-divergence}.  Then
\begin{equation}
 \mathfrak R_\Omega\Theta_1=\Theta_2\mathfrak R_\Omega.
 \label{eq:t10v12-reflection-repair}
\end{equation}
The same statement holds for either reflection parity: if
$\Theta_1f=\varsigma f$, $\varsigma\in\{1,-1\}$, then
$\Theta_2\mathfrak R_\Omega f=\varsigma\mathfrak R_\Omega f$.
\end{proposition}

\begin{proof}
Change variables $x\mapsto Rx$ in
\eqref{eq:t10v12-Neumann-elasticity}.  The transformed displacement solves
the same quotient variational problem with load $\Theta_1f$.  Uniqueness of
the quotient class makes its symmetric gradient equal to
$\Theta_2\mathfrak R_\Omega f$, proving the intertwining identity.  The
parity statement follows by linearity.
\end{proof}

\begin{firewall}[Reflection covariance versus reflected factorization]
\label{fire:t10v12-reflection-factorization}
Proposition~\ref{prop:t10v12-reflection-repair} makes the tensor correction
reflection covariant.  It does not by itself express the corresponding action
term as a positive-half contribution plus its reflection.  That stronger
factorization must be checked for the boundary action or density that owns
the correction.
\end{firewall}

\subsection{Patchwise ownership and finite obstruction coordinates}

Let $\{\Omega_p\}_{p\in P}$ be fixed-shape collar patches with overlap at
most $N_0$, and suppose $g=\sum_pg_p$ with
$\operatorname{supp}g_p\subset\Omega_p$.

\begin{proposition}[Patchwise repair under local equilibrium]
\label{prop:t10v12-patch-repair}
If every $g_p$ is orthogonal to $\mathcal R(\Omega_p)$, there are symmetric
$B_p$ supported in $\overline\Omega_p$ with
$\operatorname{div}B_p=-g_p$ and
\begin{equation}
 \left\|\sum_pB_p\right\|_{L^2}^2
 \leq N_0\sum_p C_p^2h_p^2\|g_p\|_{L^2}^2.
 \label{eq:t10v12-owner-bound}
\end{equation}
\end{proposition}

\begin{proof}
Apply Theorem~\ref{thm:t10v12-symmetric-divergence} on each patch.  At every
point at most $N_0$ summands are nonzero, so
$|\sum_pB_p|^2\leq N_0\sum_p|B_p|^2$.  Integration and the scaled estimates
give \eqref{eq:t10v12-owner-bound}.
\end{proof}

Global compatibility does not imply the compatibility of each arbitrary
partition term.  To expose exactly what fails, let $P_p$ be the $L^2$-orthogonal
projection of vector loads on $\Omega_p$ onto $\mathcal R(\Omega_p)$ and
write
\begin{equation}
 g_p=(I-P_p)g_p+P_pg_p.
 \label{eq:t10v12-local-obstruction-split}
\end{equation}
The first term has a local repair.  The second has at most $d(d+1)/2$
coordinates, namely the patch force and torque.  Exchanges of these
coordinates across adjacent patches cancel in the globally compatible sum.

\begin{definition}[Collar cohomology ledger]
\label{def:t10v12-collar-ledger}
For every collar patch the ledger retains:
\begin{enumerate}
\item its force and torque vector $P_pg_p$;
\item the locally equilibrated load $(I-P_p)g_p$;
\item the repair $B_p$ and its rescaled Korn constant;
\item the incidence rule assigning cancelling force/torque exchanges to a
shared face; and
\item the reflection parity of the patch, load, repair, and shared-face
coordinate.
\end{enumerate}
The finite family $\{P_pg_p\}$ is boundary cohomology.  It is retained rather
than charged to the irrelevant contraction norm.
\end{definition}

\begin{theorem}[Repaired-source alternative for the V11 ledger]
\label{thm:t10v12-repaired-source-card}
On a fixed rescaled Euclidean block family, assume either:
\begin{enumerate}
\item the entire connected collar is one owner and its cutoff defect is
globally compatible; or
\item a patch decomposition satisfies Definition~\ref{def:t10v12-collar-ledger}
and all shared-face force/torque exchanges close.
\end{enumerate}
Then the equilibrated bulk-plus-collar source has an $O(|k|^2)$ Fourier zero,
and its order-two physical gravitational response is bounded at zero
momentum.  The constants are uniform when the rescaled collar Korn constants,
overlap, cutoff bounds, and second-moment stocks are uniform.
\end{theorem}

\begin{proof}
In the first case use Corollary~\ref{cor:t10v12-collar-repair}; in the second,
sum Proposition~\ref{prop:t10v12-patch-repair} and the declared shared-face
closures.  In either case the completed source is symmetric, compactly
supported, and divergence free.  Theorem~\ref{thm:t10v11-stress-moments}
gives the quadratic zero, and Lemma~\ref{lem:t10v11-inverse-criterion} gives
the bounded order-two response.  Uniformity follows from
\eqref{eq:t10v12-owner-bound} and
\eqref{eq:t10v12-repair-moment}.
\end{proof}

\begin{firewall}[A repair is an owner identity, not a free counterterm]
\label{fire:t10v12-not-free-counterterm}
For a given cutoff source, the algebraic identity is
\[
 T^\chi=(T^\chi+B)-B.
\]
The first term is conserved; the second is a collar owner.  One may delete
$-B$ from the physical source only if the regulated boundary action actually
supplies $B$ with the required sign.  Otherwise V12 has localized the Ward
defect into finite boundary data but has not changed the Einstein equation.
\end{firewall}

\begin{assessment}[Progress and bottleneck after V12]
\label{assess:t10v12}
V12 constructs the symmetric collar repair promised in V11.  The construction
is a zero-traction Neumann elasticity problem on the Korn quotient, and its
cokernel is exactly force and torque.  It has the correct block scaling and
reflection covariance.  The remaining difficulty is now precise: curved
covariant divergence changes both the operator and its Killing-field
cokernel, while a physical boundary action must realize the repair rather
than merely use it as an algebraic decomposition.
\end{assessment}

\begin{programme}[Tenth-series V13: curved repair and boundary-action test]
\label{prog:t10v12-V13}
V13 will perturb the Euclidean repair to a normal-coordinate curved block.
It will derive a Neumann-series gate for covariant divergence, track the
finite approximate-Killing obstruction, and formulate the exact variational
condition under which a boundary action supplies the collar stress.  If the
covariant correction loses locality, the construction will return to a
first-order elasticity-complex potential before proceeding.
\end{programme}

\begin{firewall}[Claim boundary after tenth-series V12]
\label{fire:t10v12-boundary}
The Euclidean repair theorem does not yet control curvature perturbations,
prove a boundary action realizes the correction, close the retained
force/torque flow, verify the concrete Standard Model stress Ward identity,
or construct the full SM--Einstein continuum.
\end{firewall}

\paragraph{Parent-version record.}
V12 was formed from
\texttt{Renewal\_SM\_Einstein\_Continuum\_Research\_Notes\_V11.tex}.
The V11 parent had $4\,875$ validator-counted source lines, $183\,338$ bytes,
and SHA--256
\[
 \texttt{924E5F8D62968B68D55DAD9D54A8BCD46945047B90469419101A36A09F483F81}.
\]

% =============================================================================
% END APPENDED TENTH-SERIES V12 MATERIAL
% =============================================================================

% =============================================================================
% APPENDED TENTH-SERIES V13 MATERIAL
% =============================================================================

\section{Tenth-series V13: curved collar repair and the boundary variational
test}
\label{sec:v10v13}

The Euclidean repair of V12 is exact, but a uniform curved construction must
not invert modes that become rigid motions in the flat limit.  This version
builds the covariant repair on the complement of those modes and records the
remaining finite-dimensional obstruction.  It then gives an exact test for
whether the repair is the stress variation of a boundary action, which is the
condition needed to turn an owner identity into a physical Ward completion.

\subsection{Covariant Korn quotient}

Let $(\Omega,g)$ be a bounded Riemannian Lipschitz block.  Write
\begin{equation}
 (K_gu)_{\mu\nu}:=\frac12(\nabla_\mu u_\nu+\nabla_\nu u_\mu)
 \label{eq:t10v13-Killing-operator}
\end{equation}
for the Killing operator and
$\mathcal K_g:=\ker K_g$.  The formal adjoint identity, including the weak
normal trace, is
\begin{equation}
 \int_\Omega\langle B,K_gu\rangle_g,dv_g
 =-\int_\Omega\langle\operatorname{div}_gB,u\rangle_g,dv_g
   +\langle Bn,u\rangle_{\partial\Omega}.
 \label{eq:t10v13-adjoint}
\end{equation}

\begin{theorem}[Covariant zero-traction repair]
\label{thm:t10v13-covariant-repair}
Assume the Riemannian Korn estimate
\begin{equation}
 \inf_{k\in\mathcal K_g}\|u-k\|_{H^1_g}
 \leq C_{K,g}\|K_gu\|_{L^2_g}.
 \label{eq:t10v13-Riemannian-Korn}
\end{equation}
If $f\in L^2_g(\Omega;T\Omega)$ obeys
\begin{equation}
 \int_\Omega\langle f,k\rangle_g,dv_g=0
 \quad(k\in\mathcal K_g),
 \label{eq:t10v13-Killing-compatibility}
\end{equation}
then there is a symmetric $B\in L^2_g$ with zero weak normal traction such
that
\begin{equation}
 \operatorname{div}_gB=f,
 \qquad
 \|B\|_{L^2_g}\leq C_{K,g}\|f\|_{(H^1_g/\mathcal K_g)^*}.
 \label{eq:t10v13-curved-repair-bound}
\end{equation}
Conversely, \eqref{eq:t10v13-Killing-compatibility} is necessary.
\end{theorem}

\begin{proof}
On $H^1_g/\mathcal K_g$, use the coercive bilinear form
\[
 a_g(u,v)=\int_\Omega\langle K_gu,K_gv\rangle_g,dv_g.
\]
Compatibility makes $v\mapsto-\int\langle f,v\rangle_gdv_g$ a well-defined
quotient functional.  Lax--Milgram gives $u$, and $B=K_gu$ satisfies the
estimate.  Equation \eqref{eq:t10v13-adjoint}, tested against unrestricted
$H^1$ fields, gives $\operatorname{div}_gB=f$ and zero weak traction.
Conversely, pair the divergence equation with $k\in\mathcal K_g$ and use
\eqref{eq:t10v13-adjoint}.
\end{proof}

The exact theorem is insufficient for a uniform flat limit.  A generic small
curvature can destroy all exact Killing fields, while the corresponding
near-rigid modes still make the unquotiented Korn constant large.  They must
remain explicit retained coordinates.

\subsection{Uniform normal-coordinate complement}

Fix a Euclidean reference collar $(\Omega,g_0)$ and let
$\mathcal R_0=\ker K_{g_0}$.  Let $P_0$ be the $L^2_{g_0}$ projection onto
$\mathcal R_0$ and put $H_0=(I-P_0)H^1$.  Suppose the metric is uniformly
equivalent to $g_0$ and, after rescaling the block to unit diameter,
\begin{equation}
 \|(K_g-K_{g_0})u\|_{L^2_{g_0}}
 \leq\delta_K\|u\|_{H^1_{g_0}}.
 \label{eq:t10v13-Killing-perturbation}
\end{equation}

\begin{lemma}[Uniform complement coercivity]
\label{lem:t10v13-complement-coercivity}
If $C_{K,0}\delta_K<1$, then for $u\in H_0$,
\begin{equation}
 \|u\|_{H^1_{g_0}}
 \leq\frac{C_{K,0}}{1-C_{K,0}\delta_K}
       \|K_gu\|_{L^2_{g_0}}.
 \label{eq:t10v13-complement-coercivity}
\end{equation}
\end{lemma}

\begin{proof}
Euclidean Korn on $H_0$ and
\eqref{eq:t10v13-Killing-perturbation} give
\[
 \|u\|_{H^1}
 \leq C_{K,0}\|K_{g_0}u\|_2
 \leq C_{K,0}\bigl(\|K_gu\|_2+delta_K\|u\|_{H^1}\bigr).
\]
Move the last term to the left.
\end{proof}

\begin{proposition}[Curved repair modulo approximate Killing data]
\label{prop:t10v13-modulo-Killing}
Under Lemma~\ref{lem:t10v13-complement-coercivity}, for every curved load
$f$ there is a unique $u\in H_0$ satisfying
\begin{equation}
 \int_\Omega\langle K_gu,K_gv\rangle_gdv_g
 =-\int_\Omega\langle f,v\rangle_gdv_g
 \quad(v\in H_0).
 \label{eq:t10v13-complement-problem}
\end{equation}
Set $B=K_gu$.  Then the residual
\begin{equation}
 q:=\operatorname{div}_gB-f
 \label{eq:t10v13-finite-residual}
\end{equation}
annihilates $H_0$ and hence belongs to the finite dual space
$\mathcal R_0^*$ of dimension $d(d+1)/2$.  Moreover
\begin{equation}
 \|B\|_{L^2_g}
 \leq C(1-C_{K,0}\delta_K)^{-1}
       \|f\|_{H_0^*}.
 \label{eq:t10v13-modulo-bound}
\end{equation}
If $f$ is Euclidean force/torque compatible, then for $r\in\mathcal R_0$,
\begin{equation}
 |\langle q,r\rangle_g|
 \leq C\delta_K\|B\|_{L^2_g}\|r\|_{H^1_{g_0}}
      +|\langle f,r\rangle_g-\langle f,r\rangle_{g_0}|.
 \label{eq:t10v13-residual-debit}
\end{equation}
\end{proposition}

\begin{proof}
Lemma~\ref{lem:t10v13-complement-coercivity} makes the bilinear form in
\eqref{eq:t10v13-complement-problem} coercive on $H_0$, so Lax--Milgram gives
$u$ and \eqref{eq:t10v13-modulo-bound}.  The weak equation says exactly that
$q$ annihilates $H_0$; the quotient by $H_0$ is $\mathcal R_0$.  For
$r\in\mathcal R_0$, the adjoint identity and zero weak traction give
\[
 \langle q,r\rangle_g
 =-\int\langle B,K_gr\rangle_gdv_g-\langle f,r\rangle_g.
\]
Now $K_{g_0}r=0$, so \eqref{eq:t10v13-Killing-perturbation}, Euclidean
compatibility, and norm equivalence give
\eqref{eq:t10v13-residual-debit}.
\end{proof}

\begin{definition}[Approximate-Killing owner]
\label{def:t10v13-Killing-owner}
The coordinate $q\in\mathcal R_0^*$ in
\eqref{eq:t10v13-finite-residual} is the curved force/torque owner.  It is
retained at the shared boundary.  The complement repair $B$ enters the
irrelevant local norm only with the denominator
$1-C_{K,0}\delta_K$ explicitly charged.  No inverse is taken on
$\mathcal R_0^*$.
\end{definition}

\begin{firewall}[Why exact curved Killing fields are not the uniform owner]
\label{fire:t10v13-exact-Killing}
The dimension of $\mathcal K_g$ can jump when $g$ approaches a symmetric
metric.  Quotienting only by exact $g$-Killing fields therefore gives no
uniform flat-limit bound.  The fixed reference space $\mathcal R_0$ retains
all modes that can become small and leaves only the finite residual
\eqref{eq:t10v13-finite-residual}.
\end{firewall}

\subsection{Curved reflection covariance}

Suppose an isometric reflection $\vartheta$ preserves $\Omega$, $g$, and the
reference splitting $H^1=H_0\oplus\mathcal R_0$.  Let $\Theta_1$ and
$\Theta_2$ be the induced vector and symmetric-tensor actions.

\begin{proposition}[Equivariance of the complement repair]
\label{prop:t10v13-curved-reflection}
The solution operator of
\eqref{eq:t10v13-complement-problem} satisfies
\begin{equation}
 \mathfrak R_g\Theta_1=\Theta_2\mathfrak R_g,
 \qquad
 q(\Theta_1f)=\Theta_{\mathcal R^*}q(f).
 \label{eq:t10v13-curved-reflection}
\end{equation}
\end{proposition}

\begin{proof}
Pullback by the isometry preserves $dv_g$, the two inner products, $K_g$, and
$H_0$.  The transformed solution satisfies the same variational problem with
load $\Theta_1f$.  Uniqueness on $H_0$ gives the first identity, and applying
$\operatorname{div}_g$ gives the residual identity.
\end{proof}

\subsection{When the repair comes from a boundary action}

Let $\mathcal G$ be a star-shaped local chart of regulated metrics around a
reference $g_*$.  A proposed collar stress $B(g,\Phi)$, with matter/background
data $\Phi$ fixed for the moment, defines the metric-space one-form
\begin{equation}
 \alpha_g(h):=\frac12\int_\Omega
 \langle B(g,\Phi),h\rangle_g,dv_g.
 \label{eq:t10v13-stress-one-form}
\end{equation}

\begin{theorem}[Finite-dimensional Helmholtz test]
\label{thm:t10v13-Helmholtz}
After the regulator has made $\mathcal G$ finite dimensional, there is a
local functional $I_\partial$ with $DI_\partial=\alpha$ if and only if
\begin{equation}
 D\alpha_g[h_1](h_2)=D\alpha_g[h_2](h_1)
 \quad\hbox{for all }g,h_1,h_2.
 \label{eq:t10v13-closed-one-form}
\end{equation}
When this holds, one may choose
\begin{equation}
 I_\partial(g)-I_\partial(g_*)
 =\int_0^1
 \alpha_{g_*+t(g-g_*)}(g-g_*)dt.
 \label{eq:t10v13-boundary-action}
\end{equation}
\end{theorem}

\begin{proof}
If $DI_\partial=\alpha$, equality of mixed second derivatives gives
\eqref{eq:t10v13-closed-one-form}.  Conversely, differentiate
\eqref{eq:t10v13-boundary-action} in a direction $h$.  The result is
\[
 \int_0^1\left[
 D\alpha_{g_t}[th](g-g_*)+\alpha_{g_t}(h)\right]dt,
 \qquad g_t=g_*+t(g-g_*).
\]
Closedness exchanges the two arguments in the first term, making the
integrand the $t$ derivative of $t\alpha_{g_t}(h)$.  Evaluation at the
endpoints gives $\alpha_g(h)$.
\end{proof}

\begin{proposition}[Boundary Ward cancellation]
\label{prop:t10v13-boundary-Ward}
Assume $I_\partial$ passes Theorem~\ref{thm:t10v13-Helmholtz}, its metric
stress is $B$, and $B$ has zero weak traction at the outer edge of its collar.
For an infinitesimal diffeomorphism generated by $X$,
\begin{equation}
 \delta_X I_\partial
 =\int_\Omega\langle B,K_gX\rangle_gdv_g
 =-\int_\Omega\langle\operatorname{div}_gB,X\rangle_gdv_g.
 \label{eq:t10v13-boundary-Ward}
\end{equation}
Hence a boundary action with
$\operatorname{div}_gB=-g_{\rm bulk}$ cancels the localized bulk Ward defect
$g_{\rm bulk}$ exactly.
\end{proposition}

\begin{proof}
The metric variation under $X$ is $\delta_Xg=2K_gX$.  Insert it in
\eqref{eq:t10v13-stress-one-form} and use
\eqref{eq:t10v13-adjoint}; the traction term vanishes.  Adding the bulk Ward
variation gives zero when the two divergences have opposite signs.
\end{proof}

\begin{firewall}[The variational obstruction]
\label{fire:t10v13-variational-obstruction}
The Lax--Milgram repair is linear in a prescribed load, but its dependence on
the metric and matter fields need not make the one-form
\eqref{eq:t10v13-stress-one-form} closed.  Failure of
\eqref{eq:t10v13-closed-one-form} means there is no local boundary action with
that stress, even though the divergence equation is solvable.  In that case
the repair remains an owner decomposition, not a physical counterterm.
\end{firewall}

\begin{definition}[Curved collar action card]
\label{def:t10v13-curved-card}
A curved collar passes the V13 card if:
\begin{enumerate}
\item $C_{K,0}\delta_K<1$ and all metric norm-equivalence constants are
uniform;
\item the complement repair and finite approximate-Killing residual are
recorded as in Proposition~\ref{prop:t10v13-modulo-Killing};
\item the residual is retained and glued by shared-face force/torque owners;
\item the metric, source, splitting, and repair have the V7 reflection
parities; and
\item if the repair is inserted into the physical action, its stress one-form
passes the Helmholtz test \eqref{eq:t10v13-closed-one-form} and the resulting
action passes reflected factorization.
\end{enumerate}
\end{definition}

\begin{theorem}[Curved Ward completion, conditional form]
\label{thm:t10v13-curved-completion}
If Definition~\ref{def:t10v13-curved-card} holds and the retained
approximate-Killing residuals close across all shared faces, the combined
bulk and collar action has no interior localized diffeomorphism Ward defect.
Its complement repair is uniformly bounded and reflection covariant.
\end{theorem}

\begin{proof}
The residual closure converts the complement equation into the full
covariant divergence equation.  Proposition~\ref{prop:t10v13-boundary-Ward}
cancels the bulk defect.  The norm estimate is
\eqref{eq:t10v13-modulo-bound}, and reflection covariance is
Proposition~\ref{prop:t10v13-curved-reflection}.
\end{proof}

\begin{assessment}[Progress and bottleneck after V13]
\label{assess:t10v13}
V13 supplies a curved repair without dividing by near-Killing eigenvalues.
The price is an explicit finite force/torque residual at the shared boundary.
It also proves the missing physical test: a divergence repair is the stress
of a boundary action precisely when its metric one-form satisfies the
Helmholtz symmetry condition.  The next obstruction is quantitative.  The
curved covariant conservation law must be converted into the normal-coordinate
Fourier cancellation used by the V10 nonlinear gate, with curvature and
approximate-Killing errors charged rather than discarded.
\end{assessment}

\begin{programme}[Tenth-series V14: curved infrared debit and gate insertion]
\label{prog:t10v13-V14}
V14 will derive monopole and dipole identities for a covariantly conserved
symmetric stress in a normal-coordinate block.  It will bound their failure
from Christoffel and volume-density terms, propagate that bound through the
order-two physical inverse with an explicit infrared cutoff, and insert the
result into the V10 precision and contraction inequalities.
\end{programme}

\begin{firewall}[Claim boundary after tenth-series V13]
\label{fire:t10v13-boundary}
This version does not prove that the concrete regulated Einstein--SM boundary
stress passes the Helmholtz test, close its approximate-Killing face data,
control the curved infrared response, verify anomaly cancellation, or
construct the full SM--Einstein continuum.
\end{firewall}

\paragraph{Parent-version record.}
V13 was formed from
\texttt{Renewal\_SM\_Einstein\_Continuum\_Research\_Notes\_V12.tex}.
The V12 parent had $5\,260$ validator-counted source lines, $198\,459$ bytes,
and SHA--256
\[
 \texttt{8B493D1AA781C4360AA008A7465270947EA23CD0932AC9319506ED5D7B4C74E9}.
\]

% =============================================================================
% END APPENDED TENTH-SERIES V13 MATERIAL
% =============================================================================

% =============================================================================
% APPENDED TENTH-SERIES V14 MATERIAL
% =============================================================================

\section{Tenth-series V14: curvature debits in the infrared block gate}
\label{sec:v10v14}

The V11 quadratic Fourier zero was Euclidean, while V13 supplies a covariant
divergence repair.  In normal coordinates these statements differ by a
Christoffel source.  This version derives its exact moment identities, shows
how to isolate it as a small geometric owner, and inserts the resulting debit
into the V10 nonlinear precision and contraction inequalities.

\subsection{Densitized covariant conservation}

Let $T^{\mu\nu}=T^{\nu\mu}$ be a contravariant stress tensor on a coordinate
block and assume
\begin{equation}
 \nabla_\mu T^{\mu\nu}=0.
 \label{eq:t10v14-covariant-conservation}
\end{equation}
Define the symmetric tensor density and its coordinate defect by
\begin{equation}
 S^{\mu\nu}:=\sqrt{|g|}\,T^{\mu\nu},
 \qquad
 F^\nu:=-\sqrt{|g|}\,\Gamma^\nu_{\mu\lambda}T^{\mu\lambda}.
 \label{eq:t10v14-density-defect}
\end{equation}

\begin{lemma}[Exact coordinate divergence identity]
\label{lem:t10v14-density-divergence}
Equation \eqref{eq:t10v14-covariant-conservation} is equivalent to
\begin{equation}
 \partial_\mu S^{\mu\nu}=F^\nu.
 \label{eq:t10v14-density-divergence}
\end{equation}
\end{lemma}

\begin{proof}
Expanding the covariant divergence gives
\[
 0=\partial_\mu T^{\mu\nu}
  +\Gamma^\mu_{\mu\rho}T^{\rho\nu}
  +\Gamma^\nu_{\mu\lambda}T^{\mu\lambda}.
\]
Use $\partial_\mu\sqrt{|g|}=\sqrt{|g|}\Gamma^\rho_{\rho\mu}$ and relabel the
contracted indices in the middle term.
\end{proof}

\begin{proposition}[Curved monopole and dipole identities]
\label{prop:t10v14-curved-moments}
Suppose $S$ and $F$ are compactly supported in the coordinate chart and have
the moments below.  Set
\[
 M_{\alpha\mu\nu}:=\int x_\alpha S^{\mu\nu}(x)dx,
 \qquad
 H_{\alpha\beta\nu}:=-\int x_\alpha x_\beta F^\nu(x)dx.
\]
Then
\begin{align}
 \int S^{\alpha\nu}dx
   &=-\int x_\alpha F^\nu dx,
 \label{eq:t10v14-curved-monopole}\\
 M_{\alpha\beta\nu}
   &=\frac12\bigl(
       H_{\alpha\beta\nu}-H_{\beta\nu\alpha}
       +H_{\nu\alpha\beta}\bigr).
 \label{eq:t10v14-curved-dipole}
\end{align}
In particular, if the chart has radius at most $h$,
\begin{align}
 \left|\int S\right|&\leq h\|F\|_{L^1},
 \label{eq:t10v14-monopole-bound}\\
 \left|\int x\otimes S\right|&\leq\frac32h^2\|F\|_{L^1}.
 \label{eq:t10v14-dipole-bound}
\end{align}
\end{proposition}

\begin{proof}
Multiply \eqref{eq:t10v14-density-divergence} by $x_\alpha$ and integrate by
parts to get \eqref{eq:t10v14-curved-monopole}.  Multiplication by
$x_\alpha x_\beta$ gives
\begin{equation}
 M_{\beta\alpha\nu}+M_{\alpha\beta\nu}=H_{\alpha\beta\nu}.
 \label{eq:t10v14-MH-system}
\end{equation}
The last two indices of $M$ are symmetric.  Write
\eqref{eq:t10v14-MH-system} for the three cyclic permutations of
$(\alpha,\beta,\nu)$ and solve the resulting three linear equations; this is
\eqref{eq:t10v14-curved-dipole}.  The norm bounds follow from
$|x|\leq h$.
\end{proof}

\begin{corollary}[Fourier expansion with curvature remainder]
\label{cor:t10v14-Fourier-remainder}
If $M_2(S)=\int|x|^2|S(x)|dx<\infty$, then
\begin{equation}
 |\widehat S(k)|
 \leq h\|F\|_1
 +\frac32|k|h^2\|F\|_1
 +\frac12|k|^2M_2(S).
 \label{eq:t10v14-Fourier-remainder}
\end{equation}
\end{corollary}

\begin{proof}
Taylor expand $e^{-ik\cdot x}$ to first order.  Bound the constant and linear
moments by
\eqref{eq:t10v14-monopole-bound}--\eqref{eq:t10v14-dipole-bound}, and bound
the integral remainder by $|k|^2M_2(S)/2$.
\end{proof}

\subsection{Normal-coordinate size of the defect}

Assume the block lies inside the injectivity radius and use normal coordinates
at its centre.  Let
\begin{equation}
 \epsilon_{\rm curv}(h)
 :=h^2\|{\rm Rm}\|_{L^\infty}
   +h^3\|\nabla{\rm Rm}\|_{L^\infty}.
 \label{eq:t10v14-curvature-stock}
\end{equation}

\begin{lemma}[Rescaled Christoffel debit]
\label{lem:t10v14-Christoffel-debit}
For a fixed-shape normal-coordinate block and sufficiently small
$\epsilon_{\rm curv}(h)$,
\begin{equation}
 h\|\Gamma\|_{L^\infty}
 +\|g-g(0)\|_{L^\infty}
 +\bigl\|\sqrt{|g|}-\sqrt{|g(0)|}\bigr\|_{L^\infty}
 \leq C\epsilon_{\rm curv}(h).
 \label{eq:t10v14-normal-coordinate-bound}
\end{equation}
Consequently
\begin{equation}
 h\|F\|_{L^1}
 \leq C\epsilon_{\rm curv}(h)\|S\|_{L^1}.
 \label{eq:t10v14-F-bound}
\end{equation}
\end{lemma}

\begin{proof}
The normal-coordinate integral expansions of the metric and Christoffel
symbols are obtained by integrating the Jacobi equation along radial
geodesics.  They give $g-g(0)=O(|x|^2\|{\rm Rm}\|)$ and
$\Gamma=O(|x|\|{\rm Rm}\|+|x|^2\|\nabla{\rm Rm}\|)$ on the block.  The
determinant is a smooth function of the metric on a uniformly elliptic set.
These facts prove \eqref{eq:t10v14-normal-coordinate-bound}.  Equation
\eqref{eq:t10v14-density-defect} and uniform equivalence of $S$ and
$\sqrt{|g|}T$ give \eqref{eq:t10v14-F-bound}.
\end{proof}

In rescaled variables $y=x/h$, $p=hk$, Corollary
\ref{cor:t10v14-Fourier-remainder} becomes, after harmless norm
normalization,
\begin{equation}
 |\widehat S_h(p)|
 \leq C\epsilon_{\rm curv}(h)(1+|p|)\|S_h\|_{L^1}
 +\frac12|p|^2M_{2,y}(S_h).
 \label{eq:t10v14-rescaled-Fourier}
\end{equation}

\subsection{Exact geometric-owner split}

Because $F=\partial_\mu S^{\mu\nu}$ and $S$ is symmetric and compactly
supported, Proposition~\ref{prop:t10v12-cutoff-compatible} shows that $F$ has
zero Euclidean force and torque.

\begin{proposition}[Flattening repair]
\label{prop:t10v14-flattening-repair}
On a fixed rescaled block there is a symmetric, zero-traction tensor $C$ such
that
\begin{equation}
 \partial_\mu C^{\mu\nu}=-F^\nu,
 \qquad
 \|C\|_{L^2}\leq C_Kh\|F\|_{L^2}.
 \label{eq:t10v14-flattening-repair}
\end{equation}
Thus $S+C$ is Euclidean divergence free and has an $O(|k|^2)$ Fourier zero.
In rescaled norms the geometric owner $C$ is
$O(\epsilon_{\rm curv}(h))$ relative to $S$, provided the corresponding
$L^2$ stress stock is uniform.
\end{proposition}

\begin{proof}
Apply Theorem~\ref{thm:t10v12-symmetric-divergence} to $-F$.  Then
$\partial(S+C)=0$, so Theorem~\ref{thm:t10v11-stress-moments} gives the
quadratic Fourier zero.  The relative estimate follows from
\eqref{eq:t10v14-F-bound} and its $L^2$ analogue.
\end{proof}

\begin{firewall}[The geometric owner has not disappeared]
\label{fire:t10v14-geometric-owner}
The identity $S=(S+C)-C$ puts the first term in the quadratically vanishing
sector and the small tensor $C$ in a curvature owner.  The inverse symbol is
still singular on $C$ at exact zero momentum.  Therefore $C$ must either be
absorbed by a physical covariant boundary completion or projected into the
retained low-mode sector before the irrelevant inverse is taken.
\end{firewall}

\subsection{Low-mode separation and inverse debit}

Let $P_{<\rho}$ retain rescaled spectral/Fourier modes $|p|<\rho$, and let
$P_{\geq\rho}=I-P_{<\rho}$.  Suppose the physical gauge-fixed inverse on the
high sector obeys
\begin{equation}
 \|G_h(p)P_{\geq\rho}\|\leq C_G|p|^{-2}.
 \label{eq:t10v14-high-inverse}
\end{equation}

\begin{proposition}[Curved infrared response debit]
\label{prop:t10v14-infrared-debit}
Under \eqref{eq:t10v14-rescaled-Fourier} and
\eqref{eq:t10v14-high-inverse},
\begin{equation}
 \|G_hP_{\geq\rho}S_h\|
 \leq C_G\left[
 C\epsilon_{\rm curv}(h)(\rho^{-2}+\rho^{-1})
       \|S_h\|_{L^1}
 +\frac12M_{2,y}(S_h)
 \right].
 \label{eq:t10v14-infrared-debit}
\end{equation}
The first bracket is the curved infrared debit
\begin{equation}
 E_{\rm IR}(h,\rho)
 :=C\epsilon_{\rm curv}(h)(\rho^{-2}+\rho^{-1})
       \|S_h\|_{L^1}.
 \label{eq:t10v14-EIR}
\end{equation}
\end{proposition}

\begin{proof}
For $|p|\geq\rho$, multiply
\eqref{eq:t10v14-rescaled-Fourier} by the bound
$C_G|p|^{-2}$.  The quadratic-moment term becomes bounded, while the constant
and linear curvature terms are bounded by their values at $\rho$.
\end{proof}

\begin{firewall}[Cofinal order of limits]
\label{fire:t10v14-order-limits}
At fixed curvature stock, sending $\rho\downarrow0$ makes
$E_{\rm IR}(h,\rho)$ diverge.  A valid cofinal sequence must refine the metric
so that $\epsilon_{\rm curv}(h)=o(\rho^2)$, realize an exact physical Ward
completion that removes the geometric owner, or retain the affected modes.
The limits cannot be interchanged without one of these mechanisms.
\end{firewall}

\subsection{Insertion into the nonlinear one-step gate}

Let $d_0,q_0$ be the V10 reference precision floor and linear contraction,
and retain all V10 debits.  Suppose the reduced Hessian and RG derivative
depend Lipschitz-continuously on the high-sector stress response, with
constants $C_{\rm prec}$ and $C_{\rm RG}$ in the matched norms.

\begin{definition}[Curved V14 gate]
\label{def:t10v14-curved-gate}
In addition to Definition~\ref{def:t10v10-unified-gate}, require
\begin{align}
 &\Delta_{\rm Schur}+\Delta_{\rm con}+\Delta_{\rm frame}
   +C_{\rm prec}E_{\rm IR}(h,\rho)<d_0,
 \label{eq:t10v14-precision-gate}\\
 &q_0+C_L\omega_R+\epsilon_{\rm nl}(r)
   +C_{\mathcal T}\frac{2\theta}{1-\theta}
   +C_{\rm RG}E_{\rm IR}(h,\rho)<1.
 \label{eq:t10v14-contraction-gate}
\end{align}
The low sector $P_{<\rho}$, the V13 approximate-Killing residual, and the V12
face force/torque coordinates are retained in one finite block-background
chart with a declared gluing map.
\end{definition}

\begin{theorem}[Curved nonlinear one-step implication]
\label{thm:t10v14-one-step}
If the V14 gate holds, the high physical sector of one regulated block has a
strict positive Hessian floor and a strict nonlinear RG contraction.  Its
normalized smeared fields satisfy the V10 sub-Gaussian estimate.  If the V13
boundary action passes the Helmholtz and reflected-factorization tests, the
combined block also preserves reflection positivity and the localized
covariant Ward identity.
\end{theorem}

\begin{proof}
Proposition~\ref{prop:t10v14-infrared-debit} and the two Lipschitz hypotheses
add exactly the last terms in
\eqref{eq:t10v14-precision-gate}--\eqref{eq:t10v14-contraction-gate}.
The proofs of Lemma~\ref{lem:t10v10-Schur-debit}, Lemma
\ref{lem:t10v10-nonlinear-contraction}, and Corollary
\ref{cor:t10v10-reduced-tail} then apply with the reduced margins.  Ward
completion is Theorem~\ref{thm:t10v13-curved-completion}; reflected
factorization is the final row of the V13 card and the V10 reflected-square
argument.
\end{proof}

\begin{assessment}[Progress and remaining work after V14]
\label{assess:t10v14}
V14 derives the exact curvature failure of the Euclidean stress moments and
charges it through an order-two inverse.  The resulting gate displays the
necessary relation between block refinement and low-mode retention.  This
closes the abstract curved high-sector one-step estimate.  It does not yet
control the retained low/background chart over repeated scales; that chart
contains physical relevant couplings, approximate Killing data, and stress
monopole/dipole owners.
\end{assessment}

\begin{programme}[Tenth-series V15: companion audit and retained-flow map]
\label{prog:t10v14-V15}
V15 will perform the compulsory YM/NS audit.  It will then define the finite
retained chart and derive its exact affine cocycle under block gluing,
including force, torque, stress moments, relevant Standard Model couplings,
and the cosmological/Newton background coordinates.  The objective is to
separate tunable renormalization conditions from conservation laws before a
cofinal iteration is claimed.
\end{programme}

\begin{firewall}[Claim boundary after tenth-series V14]
\label{fire:t10v14-boundary}
The curved gate is still conditional on concrete Lipschitz constants, a
uniform physical precision floor, a realizable boundary action, and control
of the retained low-sector flow.  It does not establish those properties for
the full regulated Standard Model, cancel chiral anomalies, construct the
cofinal measure, or prove the full SM--Einstein continuum.
\end{firewall}

\paragraph{Parent-version record.}
V14 was formed from
\texttt{Renewal\_SM\_Einstein\_Continuum\_Research\_Notes\_V13.tex}.
The V13 parent had $5\,611$ validator-counted source lines, $212\,169$ bytes,
and SHA--256
\[
 \texttt{0D67D006AF4194DB0A460F5A6C4AC73C676E4A073563E2CEDC684449F57DB550}.
\]

% =============================================================================
% END APPENDED TENTH-SERIES V14 MATERIAL
% =============================================================================

% =============================================================================
% APPENDED TENTH-SERIES V15 MATERIAL
% =============================================================================

\section{Tenth-series V15: companion audit, geometric reserve restoration,
and the retained affine ledger}
\label{sec:v10v15}

\subsection{Compulsory companion audit}

The active companion files inspected at this checkpoint were
\[
\begin{array}{c|r|l}
\text{file}&\text{bytes}&\text{SHA--256}\ \hline
\texttt{Renewal\_Yang\_Mills\_Mass\_Gap\_Research\_Notes\_V31.tex}
&495\,394&
\texttt{CD0787D9AA7284AEF7894A1C62361F68C3805B7689E52DD7DD547E2AE9C0AE25}
\\
\texttt{Renewal\_NS\_Stability\_Research\_Notes\_V26.tex}
&278\,283&
\texttt{82C887F8C2C69E4F28FAD7C3DC8DE5B9099CB5A4BF431EBA1FFF3E91935978AD}.
\end{array}
\]
The NS companion is unchanged from the preceding audits.  YM V28--V31 now
proves a position-and-shell packet algebra, an exponential size/diameter
polymer reserve, and a strict dyadic connected-plaquette occupancy estimate
\[
 |\mathcal RX|\leq\frac{95}{96}|X|+1.
\]
It uses that estimate to restore, after reblocking, the geometric reserve
spent by nonlinear hull multiplication.  It also corrects the scope of the
raw paraproduct estimate in the presence of an order-two inverse and requires
the same two-moment cancellation isolated here in V11--V14.

\begin{assessment}[Audit transfer]
\label{assess:t10v15-audit-transfer}
Two conclusions transfer.  First, moment cancellation must precede a
massless order-two solve; V11--V14 implement that row for the gravitational
stress.  Second, the strict occupancy argument is combinatorial and extends
to a mixed regulator once a finite crossing-only state constant is proved.
No YM numerical constant is imported into the SM--Einstein regulator, whose
cell species and adjacency graph have not yet been fixed.
\end{assessment}

\subsection{A regulator-independent occupancy theorem}

Let $\mathcal G$ be the local adjacency graph of fine regulator cells, with a
map $b$ from fine cells to dyadic coarse blocks.  Call an edge internal when
its endpoints have the same $b$ value and crossing otherwise.

\begin{theorem}[Finite crossing-state occupancy]
\label{thm:t10v15-occupancy}
Assume every connected subgraph of $\mathcal G$ containing only crossing
edges has at most $M_*$ vertices.  If $X$ is a finite connected fine polymer
and $\mathcal RX=b(X)$ is its occupied coarse-block set, then
\begin{equation}
 |\mathcal RX|
 \leq\left(1-\frac1{M_*}\right)|X|+1.
 \label{eq:t10v15-occupancy}
\end{equation}
\end{theorem}

\begin{proof}
Choose a spanning tree of $X$ and let $I,C$ be its numbers of internal and
crossing edges.  Thus $I+C=|X|-1$.  Removing the $I$ internal edges leaves
$I+1$ crossing-only components, so $|X|\leq M_*(I+1)$ and
$I\geq |X|/M_*-1$.  The image of the tree is connected on
$|\mathcal RX|$ vertices and uses at most $C$ noncollapsed edges.  Hence
$|\mathcal RX|-1\leq C=|X|-1-I$.  Substitute the lower bound for $I$.
\end{proof}

\begin{corollary}[Abstract reserve cycle for a mixed regulator]
\label{cor:t10v15-reserve-cycle}
Suppose also that
\begin{equation}
 {\rm diam}(\mathcal RX)
 \leq\frac12{\rm diam}(X)+b_D
 \label{eq:t10v15-diameter-contraction}
\end{equation}
and that fine activities assigned to a coarse owner have a uniform
aggregation constant $M_{\mathcal R}$.  For any strong exponential reserve
parameters $\kappa,\eta>0$, choose
\begin{equation}
 \left(1-\frac1{M_*}\right)\kappa<\kappa'<\kappa,
 \qquad
 \frac12\eta<\eta'<\eta.
 \label{eq:t10v15-reserve-parameters}
\end{equation}
Then hull multiplication may spend the gaps
$\kappa-\kappa'$ and $\eta-\eta'$, and dyadic reblocking restores the strong
$(\kappa,\eta)$ reserve with a finite volume-independent constant.
\end{corollary}

\begin{proof}
The hull estimate uses
$|X\cup Y|\leq|X|+|Y|$ and
${\rm diam}(X\cup Y)\leq{\rm diam}(X)+{\rm diam}(Y)$; its intersection count
and any fixed polynomial reanchoring factor are absorbed by the two strict
exponential gaps.  Theorem~\ref{thm:t10v15-occupancy} gives size contraction
$1-1/M_*$ with additive constant $1$, while
\eqref{eq:t10v15-diameter-contraction} gives diameter contraction $1/2$ with
additive constant $b_D$.  Conditions
\eqref{eq:t10v15-reserve-parameters} therefore bound the strong output
weight by a fixed additive exponential factor times the weak input weight.
The owner aggregation costs $M_{\mathcal R}$.
\end{proof}

\begin{definition}[Mixed-cell geometry gate]
\label{def:t10v15-mixed-geometry}
A concrete SM--Einstein regulator passes the geometry gate after listing all
fine cell species (gauge, scalar, fermion, metric, ghost, and boundary
species), fixing their local adjacency graph, and proving a finite $M_*$ for
crossing-only connected components.  The value $95/96$ from the companion YM
programme is replaced by $1-1/M_*$.
\end{definition}

\subsection{Exact moment reanchoring}

For a compactly supported vector or tensor-valued source $f$, define its
zeroth and first moments about an owner $o$ by
\begin{equation}
 Q(f):=\int f(x)dx,
 \qquad
 P_o(f):=\int(x-o)\otimes f(x)dx.
 \label{eq:t10v15-source-moments}
\end{equation}

\begin{lemma}[Translation cocycle]
\label{lem:t10v15-translation-cocycle}
For every displacement $a$,
\begin{equation}
 Q(f)=Q(f),
 \qquad
 P_{o+a}(f)=P_o(f)-a\otimes Q(f).
 \label{eq:t10v15-translation-cocycle}
\end{equation}
These transformations compose exactly:
$\mathscr T_b\mathscr T_a=\mathscr T_{a+b}$.
\end{lemma}

\begin{proof}
Insert $x-(o+a)=(x-o)-a$ in
\eqref{eq:t10v15-source-moments}.  Applying the formula first with $a$ and
then with $b$ subtracts $(a+b)\otimes Q$.
\end{proof}

\begin{proposition}[Moment gluing law]
\label{prop:t10v15-moment-gluing}
Let sources $f_i$ with owners $o_i$ be glued into a coarse owner $O$.  Then
\begin{align}
 Q\left(\sum_i f_i\right)&=\sum_iQ(f_i),
 \label{eq:t10v15-Q-gluing}\\
 P_O\left(\sum_i f_i\right)
 &=\sum_i\left[P_{o_i}(f_i)+(o_i-O)\otimes Q(f_i)\right].
 \label{eq:t10v15-P-gluing}
\end{align}
This affine rule is associative under repeated reblocking.
\end{proposition}

\begin{proof}
Linearity proves the first formula.  For the second, use
$x-O=(x-o_i)+(o_i-O)$ inside each integral.  Associativity follows either by
direct substitution or from the translation cocycle in
Lemma~\ref{lem:t10v15-translation-cocycle}.
\end{proof}

For force $F$ and antisymmetric torque $L_o$, the identical calculation with
$a\wedge F$ gives
\begin{equation}
 L_{o+a}=L_o-a\wedge F,
 \qquad
 L_O^{\rm coarse}=\sum_i[L_{o_i}+(o_i-O)\wedge F_i].
 \label{eq:t10v15-force-torque-cocycle}
\end{equation}
These are coadjoint Euclidean reanchoring laws.  Treating $P_o$ or $L_o$ as
owner-independent would make block gluing nonassociative.

\subsection{Face incidence and conservation}

Orient every shared coarse face.  Let $J_f$ be its force/torque or Ward-flux
coordinate, with $J_{\bar f}=-J_f$ when the orientation is reversed.  If
$\epsilon_{bf}\in\{-1,0,1\}$ is the block-face incidence number, define
\begin{equation}
 (\partial J)_b:=\sum_f\epsilon_{bf}J_f.
 \label{eq:t10v15-incidence}
\end{equation}

\begin{lemma}[Internal face cancellation]
\label{lem:t10v15-face-cancellation}
For any union $U$ of blocks,
\begin{equation}
 \sum_{b\subset U}(\partial J)_b
 =\sum_{f\subset\partial U}\epsilon_{Uf}J_f.
 \label{eq:t10v15-discrete-Stokes}
\end{equation}
Every internal face cancels exactly.
\end{lemma}

\begin{proof}
An internal face occurs twice with opposite induced orientations.  A boundary
face occurs once.  Summing \eqref{eq:t10v15-incidence} gives the formula.
\end{proof}

\begin{definition}[Conservation row]
\label{def:t10v15-conservation-row}
If $q_b$ is the V12--V13 local force/torque obstruction, exact gluing requires
\begin{equation}
 q_b+(\partial J)_b=0.
 \label{eq:t10v15-local-conservation}
\end{equation}
Under reblocking, the coarse obstruction is the sum of fine $q_b$ plus the
external face flux in \eqref{eq:t10v15-discrete-Stokes}.  This is an identity
to be preserved, not a coupling to be tuned.
\end{definition}

\subsection{Typed retained coordinates}

For each coarse block let
\begin{equation}
 z_b=(c_b,\lambda_b,\zeta_b,Q_b,P_b,F_b,L_b,J_{\partial b}).
 \label{eq:t10v15-retained-vector}
\end{equation}
Here:
\begin{itemize}
\item $c_b$ contains relevant and marginal Standard Model couplings and
wave-function coordinates;
\item $\lambda_b$ contains cosmological, Newton, curvature, and other allowed
gravitational counterterm coordinates;
\item $\zeta_b$ contains the finitely retained spectral modes below the V14
cutoff $\rho$;
\item $(Q_b,P_b)$ are stress monopole/dipole coordinates; and
\item $(F_b,L_b,J_{\partial b})$ are the approximate-Killing and shared-face
Ward coordinates.
\end{itemize}

\begin{definition}[Typed retained map]
\label{def:t10v15-typed-map}
One retained RG step has the form
\begin{align}
 (c',\lambda')&=\mathcal B(c,\lambda)+\mathcal C(x)+t,
 \label{eq:t10v15-coupling-flow}\\
 \zeta'&=A_\zeta\zeta+S_\zeta(x,c,\lambda),
 \label{eq:t10v15-low-flow}\\
 (Q',P',F',L')&=\mathcal G_{\rm aff}(Q,P,F,L)+S_{\rm bdry},
 \label{eq:t10v15-affine-flow}\\
 q'+\partial J'&=0.
 \label{eq:t10v15-Ward-flow}
\end{align}
The map $\mathcal G_{\rm aff}$ is fixed by
\eqref{eq:t10v15-Q-gluing}--\eqref{eq:t10v15-force-torque-cocycle}; $t$ is a
chosen counterterm/tuning coordinate; and $x$ is the irrelevant polymer
activity.  The source term $S_{\rm bdry}$ must come from declared boundary
owners.
\end{definition}

\begin{firewall}[Tuning versus conservation]
\label{fire:t10v15-tuning-conservation}
The initial values of allowed relevant couplings may be tuned to meet
renormalization conditions.  Equations
\eqref{eq:t10v15-affine-flow}--\eqref{eq:t10v15-Ward-flow} are instead fixed
by source moments, incidence, and Ward identities.  A counterterm can change
them only when it is the variation of an allowed local action and is entered
on both sides of the corresponding Ward identity.  Normal ordering may set
an expected $Q$ to zero but does not tune its fluctuation, by
Lemma~\ref{lem:t10v11-normal-ordering}.
\end{firewall}

\begin{theorem}[Associativity of the retained geometric sector]
\label{thm:t10v15-retained-associativity}
The map formed by moment addition, owner reanchoring, force/torque
reanchoring, and face incidence is associative under any sequence of block
groupings.  If the fine conservation rows
\eqref{eq:t10v15-local-conservation} hold, the coarse row holds with only the
outer boundary flux of the coarse block.
\end{theorem}

\begin{proof}
Moment associativity is Proposition~\ref{prop:t10v15-moment-gluing}; the
force/torque proof is identical using
\eqref{eq:t10v15-force-torque-cocycle}.  Face cancellation is
Lemma~\ref{lem:t10v15-face-cancellation}.  These operations act on distinct
typed components and therefore their product map is associative.
\end{proof}

\begin{definition}[V15 retained-flow card]
\label{def:t10v15-retained-card}
A regulator passes the retained-flow card if:
\begin{enumerate}
\item its mixed cell graph passes Definition~\ref{def:t10v15-mixed-geometry};
\item all stress and force/torque moments use the exact affine reanchoring
laws;
\item every shared face has one orientation and one owner;
\item conserved Ward rows and tunable coupling rows are typed separately;
\item the low spectral sector has a uniform finite dimension per block and a
declared gluing map; and
\item the nonlinear irrelevant map obeys the V14 gate in the
reserve-restored polymer space of
Corollary~\ref{cor:t10v15-reserve-cycle}.
\end{enumerate}
\end{definition}

\begin{assessment}[Progress and bottleneck after V15]
\label{assess:t10v15}
V15 imports the companion programme's geometric lesson in a regulator-neutral
form and proves the exact affine cocycle for stress moments, force, torque,
and face flux.  This prevents three common errors: spending a fixed polymer
reserve at every scale, forgetting the owner dependence of dipoles, and
treating a conservation identity as a tunable coupling.  The remaining
problem is dynamical: construct an infinite retained/irrelevant trajectory
with prescribed renormalized couplings while the irrelevant component stays
inside the V14 contraction ball.
\end{assessment}

\begin{programme}[Tenth-series V16: a coupled shooting theorem]
\label{prog:t10v15-V16}
V16 will prove an abstract finite- and infinite-horizon shooting theorem for
the coupled retained/irrelevant map.  Expanding relevant directions will be
solved backward from renormalization data, stable polymer directions forward
by contraction, and marginal directions will remain conditional on an
explicit beta-function sign and summable remainder.  The affine conservation
sector will be transported exactly rather than placed in the shooting
variables.
\end{programme}

\begin{firewall}[Claim boundary after tenth-series V15]
\label{fire:t10v15-boundary}
This version does not verify a finite crossing-state constant for a concrete
mixed regulator, calculate the Standard Model or gravitational beta
functions, construct the infinite RG trajectory, prove anomaly cancellation,
or construct the full SM--Einstein continuum.
\end{firewall}

\paragraph{Parent-version record.}
V15 was formed from
\texttt{Renewal\_SM\_Einstein\_Continuum\_Research\_Notes\_V14.tex}.
The V14 parent had $5\,949$ validator-counted source lines, $224\,725$ bytes,
and SHA--256
\[
 \texttt{52F46EB04F2A1795F2C085141E0E226D30572E8378CE356F2CE8DA7C58160BE9}.
\]

% =============================================================================
% END APPENDED TENTH-SERIES V15 MATERIAL
% =============================================================================

% =============================================================================
% APPENDED TENTH-SERIES V16 MATERIAL
% =============================================================================

\section{Tenth-series V16: a coupled relevant--irrelevant shooting theorem}
\label{sec:v10v16}

V15 separated three types of retained data: tunable relevant couplings,
slow marginal couplings, and affine Ward/moment coordinates fixed by gluing.
This version proves the infinite-horizon construction for the first type
coupled to the V14 contracted polymer sector.  It also records the elementary
asymptotically-free marginal estimate needed before a marginal coordinate can
be admitted to the same trajectory.

\subsection{Hyperbolic sequence equations}

Let $X$ be the reserve-restored irrelevant Banach space and $Y$ the
finite-dimensional relevant coupling space.  Let
\begin{equation}
 \|L\|\leq q<1,
 \qquad A:Y\to Y\text{ invertible},
 \qquad\|A^{-1}\|\leq a<1.
 \label{eq:t10v16-dichotomy}
\end{equation}
The forward RG equations are
\begin{align}
 x_{n+1}&=Lx_n+f_n(x_n,y_n;h_n),
 \label{eq:t10v16-stable-equation}\\
 y_{n+1}&=Ay_n+g_n(x_n,y_n;h_n).
 \label{eq:t10v16-relevant-equation}
\end{align}
Here $h_n$ is the already determined V15 affine conservation/background
trajectory.  It is a parameter, not a shooting variable.

Assume on the product ball $\|(x,y)\|:=\max\{\|x\|,\|y\|\}\leq r$ that
\begin{align}
 \|f_n(z;h_n)-f_n(z';h_n)\|&\leq L_f\|z-z'\|,
 \label{eq:t10v16-f-Lipschitz}\\
 \|g_n(z;h_n)-g_n(z';h_n)\|&\leq L_g\|z-z'\|,
 \label{eq:t10v16-g-Lipschitz}\\
 \|f_n(0;h_n)\|&\leq b_f,
 \qquad
 \|g_n(0;h_n)\|\leq b_g.
 \label{eq:t10v16-forcing-bounds}
\end{align}

\begin{theorem}[Infinite-horizon shooting theorem]
\label{thm:t10v16-shooting}
Fix $x_0\in X$ and define
\begin{equation}
 \kappa:=\max\left\{
 \frac{L_f}{1-q},\frac{aL_g}{1-a}\right\},
 \qquad
 B_0:=\max\left\{
 \|x_0\|+\frac{b_f}{1-q},
 \frac{ab_g}{1-a}\right\}.
 \label{eq:t10v16-kappa-B}
\end{equation}
If
\begin{equation}
 \kappa<1,
 \qquad B_0+\kappa r\leq r,
 \label{eq:t10v16-shooting-gate}
\end{equation}
then there is a unique bounded trajectory in the radius-$r$ sequence ball
satisfying \eqref{eq:t10v16-stable-equation}--
\eqref{eq:t10v16-relevant-equation}, the prescribed $x_0$, and the condition
that no forward growing homogeneous relevant component is present.  It is
the unique fixed point of
\begin{align}
 x_n&=L^nx_0+sum_{j=0}^{n-1}L^{n-1-j}f_j(x_j,y_j;h_j),
 \label{eq:t10v16-LP-stable}\\
 y_n&=-\sum_{j=n}^{\infty}A^{n-1-j}g_j(x_j,y_j;h_j).
 \label{eq:t10v16-LP-relevant}
\end{align}
In particular, the required bare relevant coordinate is
\begin{equation}
 y_0=-\sum_{j=0}^{\infty}A^{-j-1}g_j(x_j,y_j;h_j).
 \label{eq:t10v16-bare-tuning}
\end{equation}
\end{theorem}

\begin{proof}
Let $\mathcal Z=\ell^\infty(\mathbb N_0;X\oplus Y)$ with the supremum of the
product norm.  Define $\mathfrak T$ by the right sides of
\eqref{eq:t10v16-LP-stable}--\eqref{eq:t10v16-LP-relevant}.  The geometric
series and \eqref{eq:t10v16-f-Lipschitz} give
\[
 \sup_n\|\mathfrak T_xz_n-\mathfrak T_xz'_n\|
 \leq\frac{L_f}{1-q}\|z-z'\|_{\mathcal Z}.
\]
Because $A^{n-1-j}=A^{-(j-n+1)}$ for $j\geq n$,
\eqref{eq:t10v16-g-Lipschitz} gives
\[
 \sup_n\|\mathfrak T_yz_n-\mathfrak T_yz'_n\|
 \leq\frac{aL_g}{1-a}\|z-z'\|_{\mathcal Z}.
\]
Thus $\mathfrak T$ has Lipschitz constant at most $\kappa$.  At zero, the two
components are bounded by the two entries defining $B_0$.  Hence
\eqref{eq:t10v16-shooting-gate} makes the radius-$r$ ball invariant, and
Banach's fixed-point theorem supplies the unique solution of the two
Lyapunov--Perron equations.

Subtracting the formulas at $n+1$ and applying $L$ or $A$ verifies the forward
equations.  Conversely, iterate the stable equation from $x_0$.  Iterating
the relevant equation backward from $N$ gives
\[
 y_n=A^{n-N}y_N-sum_{j=n}^{N-1}A^{n-1-j}g_j.
\]
For a bounded trajectory, $A^{n-N}y_N\to0$ as $N\to\infty$, proving
\eqref{eq:t10v16-LP-relevant}.  Thus the fixed point is exactly the unique
bounded trajectory without a growing relevant homogeneous component.
\end{proof}

\begin{corollary}[Stability with respect to the conservation trajectory]
\label{cor:t10v16-h-stability}
Suppose, in addition, $f_n,g_n$ are Lipschitz in $h$ with constants
$L_{fh},L_{gh}$, and two admissible affine trajectories satisfy
$\sup_n\|h_n-\widetilde h_n\|\leq\delta_h$.  The corresponding shot
trajectories satisfy
\begin{equation}
 \|z-\widetilde z\|_{\mathcal Z}
 \leq\frac1{1-\kappa}
 \max\left\{
 \frac{L_{fh}}{1-q},\frac{aL_{gh}}{1-a}
 \right\}\delta_h.
 \label{eq:t10v16-h-stability}
\end{equation}
\end{corollary}

\begin{proof}
Compare the two fixed-point maps.  Their values at the same sequence differ
by at most the maximum displayed before $\delta_h$.  Move the
$\kappa\|z-\widetilde z\|$ term to the left.
\end{proof}

\begin{firewall}[The relevant coordinate is selected, not contracted]
\label{fire:t10v16-relevant-not-contracted}
The norm $\|A^{-1}\|<1$ contracts only the backward Green series.  Forward
relevant evolution expands.  Equation \eqref{eq:t10v16-bare-tuning} is the
renormalization-condition shooting rule that removes the growing homogeneous
component; it is not a forward contraction claim.
\end{firewall}

\subsection{Finite-horizon approximants}

For a terminal value $y_N^{\rm term}$, the finite-horizon equation is
\begin{equation}
 y_n=A^{n-N}y_N^{\rm term}
 -\sum_{j=n}^{N-1}A^{n-1-j}g_j(x_j,y_j;h_j),
 \qquad 0\leq n\leq N.
 \label{eq:t10v16-finite-horizon}
\end{equation}

\begin{proposition}[Removal of the ultraviolet terminal condition]
\label{prop:t10v16-terminal-removal}
Assume the V16 gate uniformly in $N$ and
$\sup_N\|y_N^{\rm term}\|<\infty$.  For every fixed $n$, the explicit
homogeneous terminal contribution in \eqref{eq:t10v16-finite-horizon} obeys
\begin{equation}
 \|A^{n-N}y_N^{\rm term}\|
 \leq a^{N-n}\sup_N\|y_N^{\rm term}\|\longrightarrow0.
 \label{eq:t10v16-terminal-decay}
\end{equation}
Every locally convergent subsequence of finite-horizon fixed points therefore
satisfies the infinite Lyapunov--Perron equations and, by uniqueness, has the
trajectory of Theorem~\ref{thm:t10v16-shooting} as its limit.
\end{proposition}

\begin{proof}
The estimate is \eqref{eq:t10v16-dichotomy}.  On every fixed initial window,
the missing tail in the finite sum is bounded by a geometric series times
$a^{N-n}$.  Passing to a locally convergent subsequence gives
\eqref{eq:t10v16-LP-stable}--\eqref{eq:t10v16-LP-relevant}; uniqueness
identifies the limit.
\end{proof}

\begin{firewall}[Compactness is a separate input]
\label{fire:t10v16-finite-compactness}
Proposition~\ref{prop:t10v16-terminal-removal} identifies any locally
convergent subsequential limit; it does not create compactness in an
infinite-dimensional field space.  The V10 field-tail/equicontinuity card and
the polymer localization bounds are still needed to extract such a
subsequence for measures.
\end{firewall}

\subsection{A marginal beta-function gate}

Let $u_n$ be a nonnegative marginal coupling satisfying
\begin{equation}
 u_{n+1}=u_n-bu_n^2+r_n,
 \qquad |r_n|\leq c u_n^3.
 \label{eq:t10v16-marginal-recurrence}
\end{equation}

\begin{lemma}[Marginally irrelevant decay]
\label{lem:t10v16-marginal-decay}
Suppose $b>0$ and
\begin{equation}
 0<u_0\leq\delta,
 \qquad c\delta\leq\frac b2,
 \qquad \frac{3b}{2}\delta\leq\frac12.
 \label{eq:t10v16-marginal-smallness}
\end{equation}
Then $u_n$ remains positive and decreases, and
\begin{equation}
 \frac1{u_0^{-1}+3bn}
 \leq u_n\leq
 \frac1{u_0^{-1}+(b/2)n}.
 \label{eq:t10v16-marginal-bounds}
\end{equation}
In particular $u_n\asymp n^{-1}$.
\end{lemma}

\begin{proof}
The remainder bound and \eqref{eq:t10v16-marginal-smallness} imply
\begin{equation}
 \frac b2u_n^2
 \leq u_n-u_{n+1}
 \leq\frac{3b}{2}u_n^2.
 \label{eq:t10v16-marginal-decrement}
\end{equation}
As long as $u_n\leq\delta$, the upper decrement is at most $u_n/2$, so
$u_{n+1}\geq u_n/2>0$; the lower decrement makes the sequence decreasing and
therefore closes the induction.  Now
\[
 \frac1{u_{n+1}}-\frac1{u_n}
 =\frac{u_n-u_{n+1}}{u_nu_{n+1}}.
\]
Using $u_{n+1}\leq u_n$ for the lower bound and
$u_{n+1}\geq u_n/2$ for the upper bound gives
$b/2\leq u_{n+1}^{-1}-u_n^{-1}\leq3b$.  Sum from $0$ to $n-1$ and invert.
\end{proof}

\begin{firewall}[Beta-function sign remains physical input]
\label{fire:t10v16-beta-sign}
Lemma~\ref{lem:t10v16-marginal-decay} is conditional on $b>0$ and on a cubic
remainder bound in the actual matched scheme.  It does not calculate the
coupled Standard Model--gravity beta functions.  A zero or negative effective
$b$, a Landau direction, or mixing with another marginal coordinate requires
a different retained-flow analysis.
\end{firewall}

\subsection{Coupling the theorem to the earlier cards}

\begin{definition}[V16 trajectory card]
\label{def:t10v16-trajectory-card}
An RG scale family passes the trajectory card when:
\begin{enumerate}
\item the V14 curved nonlinear gate supplies a uniform stable derivative
bound $q<1$ in the V15 reserve-restored polymer space;
\item relevant linearization has a uniform backward bound $a<1$ after all
zero eigenvalues are moved to the marginal sector;
\item the nonlinear cross-derivatives satisfy
\eqref{eq:t10v16-shooting-gate};
\item the V15 affine conservation trajectory closes exactly and supplies a
bounded $h_n$;
\item each marginal direction has a separately proved beta-function normal
form and remainder estimate; and
\item the shot bare coordinates are restricted to counterterms allowed by
locality, gauge symmetry, diffeomorphism covariance, reflection, and anomaly
cancellation.
\end{enumerate}
\end{definition}

\begin{theorem}[Conditional infinite RG trajectory]
\label{thm:t10v16-infinite-trajectory}
If Definition~\ref{def:t10v16-trajectory-card} holds, then every sufficiently
small admissible stable datum $x_0$ and admissible affine Ward/background
trajectory determine a unique small infinite relevant--irrelevant RG
trajectory.  Its bare relevant couplings are given by
\eqref{eq:t10v16-bare-tuning}; marginally irrelevant coordinates satisfying
Lemma~\ref{lem:t10v16-marginal-decay} decay like $n^{-1}$.
\end{theorem}

\begin{proof}
Apply Theorem~\ref{thm:t10v16-shooting} to the stable and relevant sectors,
with the conservation trajectory as a parameter.  Apply
Lemma~\ref{lem:t10v16-marginal-decay} to each decoupled normal-form marginal
row; any declared small mixing is included in its remainder estimate.
\end{proof}

\begin{assessment}[Progress and next obstruction after V16]
\label{assess:t10v16}
V16 constructs the abstract infinite RG trajectory rather than merely
asserting that relevant coordinates can be tuned.  The proof shows the exact
smallness denominators $(1-q)^{-1}$ and $(1-a)^{-1}$ and keeps Ward data out
of the shooting variables.  The main unresolved physical input has now moved
to the finite retained sector: the coupled SM--Einstein marginal beta matrix,
the allowed counterterm basis, and the anomaly constraints must be computed
in one regulator and scheme.
\end{assessment}

\begin{programme}[Tenth-series V17: anomaly-compatible counterterm complex]
\label{prog:t10v16-V17}
V17 will build the finite local counterterm complex before attempting a beta
calculation.  It will separate BRST-exact changes, gauge/diffeomorphism
invariants, and possible anomaly cohomology; prove the linear algebraic
solvability criterion for simultaneous Ward and renormalization conditions;
and identify which Standard Model anomaly cancellations are representation
identities rather than tunings.
\end{programme}

\begin{firewall}[Claim boundary after tenth-series V16]
\label{fire:t10v16-boundary}
The trajectory theorem is conditional on uniform constants and a verified
finite retained map.  This version does not calculate the physical beta
functions, prove anomaly cancellation in the regulator, establish measure
compactness or Osterwalder--Schrader reconstruction, or construct the full
SM--Einstein continuum.
\end{firewall}

\paragraph{Parent-version record.}
V16 was formed from
\texttt{Renewal\_SM\_Einstein\_Continuum\_Research\_Notes\_V15.tex}.
The V15 parent had $6\,294$ validator-counted source lines, $238\,388$ bytes,
and SHA--256
\[
 \texttt{0D5EDF3CF6458C6A163B3318F2EBDDFC589183AB71418604AC0B235D983AD750}.
\]

% =============================================================================
% END APPENDED TENTH-SERIES V16 MATERIAL
% =============================================================================

% =============================================================================
% APPENDED TENTH-SERIES V17 MATERIAL
% =============================================================================

\section{Tenth-series V17: anomaly cohomology, compatible counterterms, and
the gravitational tower}
\label{sec:v10v17}

The shooting theorem of V16 is meaningful only after its tunable coordinates
are distinguished from Ward identities and anomaly obstructions.  This
version gives the finite regulated solvability criterion, checks the
one-generation Standard Model anomaly sums in a fixed left-handed convention,
and then exposes an independent obstruction: perturbative Einstein gravity
generates counterterms of unbounded derivative order.  A finite retained
coupling chart therefore cannot by itself support regulator removal.

\subsection{A finite regulated BRST complex}

At a fixed finite regulator, let
\begin{equation}
 C^{-1}\xrightarrow{s_{-1}}C^0\xrightarrow{s_0}C^1
 \xrightarrow{s_1}C^2,
 \qquad s_0s_{-1}=0,
 \qquad s_1s_0=0,
 \label{eq:t10v17-complex}
\end{equation}
be the local counterterm/anomaly part of the BRST complex.  Elements of
$C^0$ are allowed ghost-number-zero local counterterms, and elements of
$C^1$ are possible Ward breakings.  The Wess--Zumino condition on a breaking
$\mathcal A$ is $s_1\mathcal A=0$.

\begin{proposition}[Cohomological anomaly criterion]
\label{prop:t10v17-cohomology}
A consistent breaking $\mathcal A\in\ker s_1$ can be removed by an allowed
counterterm if and only if
\begin{equation}
 [\mathcal A]=0\quad\hbox{in}\quad
 H^1_s:=\ker s_1/\operatorname{im}s_0.
 \label{eq:t10v17-anomaly-class}
\end{equation}
If $c\in C^0$ is added to the action, the breaking changes from
$\mathcal A$ to $\mathcal A+s_0c$.
\end{proposition}

\begin{proof}
The new Ward identity is unbroken exactly when
$s_0c=-\mathcal A$.  Such a $c$ exists exactly when
$\mathcal A\in\operatorname{im}s_0$.  Consistency places $\mathcal A$ in
$\ker s_1$, so this is precisely the vanishing of its cohomology class.
\end{proof}

Let $R:C^0\to\mathbb R^m$ encode chosen renormalization conditions and
$r\in\mathbb R^m$ their desired values.

\begin{theorem}[Simultaneous Ward and renormalization solvability]
\label{thm:t10v17-simultaneous}
Define
\begin{equation}
 \mathcal M:C^0\longrightarrow C^1\oplus\mathbb R^m,
 \qquad \mathcal Mc=(s_0c,Rc).
 \label{eq:t10v17-combined-map}
\end{equation}
There is a counterterm satisfying
\begin{equation}
 s_0c=-\mathcal A,
 \qquad Rc=r
 \label{eq:t10v17-combined-equations}
\end{equation}
if and only if
\begin{equation}
 (-\mathcal A,r)\perp\ker\mathcal M^*.
 \label{eq:t10v17-Fredholm-test}
\end{equation}
When this holds, the minimum-norm solution is
\begin{equation}
 c_{\min}=\mathcal M^\dagger(-\mathcal A,r),
 \qquad
 \|c_{\min}\|\leq\sigma_+^{-1}|(-\mathcal A,r)\|,
 \label{eq:t10v17-pseudoinverse}
\end{equation}
where $\mathcal M^\dagger$ is the Moore--Penrose inverse and $\sigma_+$ is
the smallest nonzero singular value of $\mathcal M$.  All solutions are
$c_{\min}+\ker\mathcal M$.
\end{theorem}

\begin{proof}
In finite dimension, $\operatorname{ran}\mathcal M=(\ker\mathcal M^*)^\perp$.
Thus \eqref{eq:t10v17-Fredholm-test} is equivalent to solvability.  Singular
value decomposition proves the formula, norm estimate, and affine description
of all solutions.
\end{proof}

\begin{firewall}[Small singular values are an RG debit]
\label{fire:t10v17-small-singular}
Algebraic solvability is not a uniform estimate.  If $\sigma_+$ tends to zero
along the regulator family, the necessary counterterm can leave the V16
small trajectory ball.  The counterterm card must therefore carry a uniform
lower bound for the nonzero singular values after exact symmetry zero modes
are quotiented.
\end{firewall}

\subsection{One-generation Standard Model anomaly identities}

Use only left-handed Weyl fields.  In the convention
$Q_{\rm em}=T_3+Y$, one generation is
\begin{equation}
\begin{array}{c|c|c|c}
\text{field}&SU(3)&SU(2)&Y\\ \hline
Q_L&\mathbf3&\mathbf2&1/6\\
u_R^c&\overline{\mathbf3}&\mathbf1&-2/3\\
d_R^c&\overline{\mathbf3}&\mathbf1&1/3\\
L_L&\mathbf1&\mathbf2&-1/2\\
e_R^c&\mathbf1&\mathbf1&1.
\end{array}
\label{eq:t10v17-SM-representations}
\end{equation}
A sterile $\nu_R^c$ has $Y=0$ and changes none of the sums below.

\begin{proposition}[Perturbative local anomaly cancellation]
\label{prop:t10v17-SM-anomalies}
The representation list \eqref{eq:t10v17-SM-representations} has vanishing
coefficients for
\begin{equation}
 SU(3)^3,quad SU(3)^2U(1),quad SU(2)^2U(1),quad U(1)^3,
 \quad {\rm grav}^2U(1).
 \label{eq:t10v17-anomaly-list}
\end{equation}
The perturbative $SU(2)^3$ coefficient vanishes because the doublet is
pseudoreal.
\end{proposition}

\begin{proof}
For $SU(3)^3$, the two colour fundamentals in $Q_L$ cancel the two colour
antifundamentals $u_R^c,d_R^c$.  With $T(\mathbf3)=T(\mathbf2)=1/2$, the two
mixed coefficients are
\begin{align*}
 SU(3)^2U(1):&\quad
 2\frac16\frac12-\frac23\frac12+\frac13\frac12=0,\\
 SU(2)^2U(1):&\quad
 3\frac16\frac12-\frac12\frac12=0.
\end{align*}
Counting colour and weak multiplicities, the cubic abelian sum is
\begin{align*}
 6\left(\frac16\right)^3
 +3\left(-\frac23\right)^3
 +3\left(\frac13\right)^3
 +2\left(-\frac12\right)^3+1^3
 =\frac{1-32+4-9+36}{36}=0.
\end{align*}
The mixed gravitational sum is
\[
 6\frac16+3\left(-\frac23\right)+3\frac13
 +2\left(-\frac12\right)+1=1-2+1-1+1=0.
\]
Finally, the symmetrized cubic trace of a pseudoreal $SU(2)$ representation
vanishes, so there is no perturbative $SU(2)^3$ anomaly.
\end{proof}

\begin{proposition}[The global $SU(2)$ parity test]
\label{prop:t10v17-Witten-parity}
One Standard Model generation has four left-handed $SU(2)$ doublets when
colour multiplicity is counted.  Hence its doublet number is even, satisfying
the global $SU(2)$ parity condition.
\end{proposition}

\begin{proof}
$Q_L$ contributes one doublet for each of three colours and $L_L$ contributes
one, for a total of four.
\end{proof}

\begin{firewall}[Representation cancellation versus regulated determinant]
\label{fire:t10v17-determinant-firewall}
Propositions~\ref{prop:t10v17-SM-anomalies} and
\ref{prop:t10v17-Witten-parity} are representation identities.  They do not
construct a regulator with exact chiral gauge symmetry, trivialize its
fermion determinant line, or prove absence of every global gauge/gravitational
holonomy on the chosen configuration space.  Vanishing local anomaly
curvature can leave a flat line bundle with nontrivial holonomy.
\end{firewall}

\subsection{Typed counterterm decomposition}

At fixed cutoff, decompose the ghost-number-zero space as
\begin{equation}
 C^0=\operatorname{im}s_{-1}\oplus\mathcal H^0\oplus N,
 \label{eq:t10v17-counterterm-splitting}
\end{equation}
where $\mathcal H^0$ represents invariant cohomology classes and $N$ is a
chosen complement on which $s_0$ maps onto removable Ward breakings.

\begin{definition}[Counterterm types]
\label{def:t10v17-counterterm-types}
\begin{enumerate}
\item $\operatorname{im}s_{-1}$ contains BRST-exact reparametrizations and
gauge-fixing changes.
\item $\mathcal H^0$ contains physical invariant couplings, including the
allowed gauge, Yukawa, scalar, cosmological, Einstein, and higher-curvature
operators at the chosen cutoff order.
\item $N$ contains symmetry-restoring counterterms fixed by the Ward
conditions rather than by physical renormalization data.
\end{enumerate}
The V16 shooting variables may include selected coordinates of
$\mathcal H^0$, but not a nonzero class in $H^1_s$.
\end{definition}

\begin{corollary}[No tuning removes a genuine anomaly]
\label{cor:t10v17-no-anomaly-tuning}
If $[\mathcal A]\neq0$ in $H^1_s$, changing invariant couplings in
$\mathcal H^0$ or exact gauge-fixing terms in $\operatorname{im}s_{-1}$
cannot restore the Ward identity.
\end{corollary}

\begin{proof}
An invariant counterterm has zero $s_0$ image, and exact changes have an
$s_0$ image that remains in $\operatorname{im}s_0$.  Neither changes the
cohomology class of $\mathcal A$.
\end{proof}

\subsection{Power counting for Einstein vertices}

Consider perturbation theory about a smooth flat background in four
dimensions.  The graviton propagator has momentum degree $-2$ and every
Einstein--Hilbert interaction vertex has two derivatives.

\begin{theorem}[Unbounded gravitational derivative order]
\label{thm:t10v17-gravity-power-counting}
For a connected pure-gravity Feynman graph with $L$ loops, $I$ internal
lines, and $V$ Einstein--Hilbert vertices, its superficial momentum degree is
\begin{equation}
 D=4L-2I+2V=2L+2.
 \label{eq:t10v17-gravity-degree}
\end{equation}
Therefore increasing loop order permits local divergent invariants with
unbounded derivative order.  In the absence of an additional fixed-point or
cancellation mechanism, no finite list of gravitational counterterms closes
perturbative regulator removal.
\end{theorem}

\begin{proof}
Each loop integration contributes four powers of momentum, each propagator
removes two, and each vertex contributes two, giving the first expression.
For a connected graph, $L=I-V+1$.  Substitution gives $D=2L+2$.  Since $L$
is unbounded, so is the allowed degree of a local divergence.  Counterterms
at successively higher degrees are independent modulo integrations by parts,
Bianchi identities, field redefinitions, and lower-order equations; these
relations reduce each fixed degree to a finite basis but do not bound the
degree as $L\to\infty$.
\end{proof}

\begin{remark}[First low-loop cancellations do not change the conclusion]
Specific on-shell divergences can vanish at a particular loop order or be
removed by field redefinition.  The power-counting theorem states the
unbounded set allowed by locality and dimension; it does not assert that
every allowed invariant has a nonzero coefficient at every loop order.
Closing the tower requires a theorem beyond a finite number of cancellations.
\end{remark}

\begin{corollary}[Correction to a finite retained chart]
\label{cor:t10v17-finite-chart-correction}
The finite retained vector of V15 is sufficient for one fixed regulator and
fixed operator truncation.  A full ultraviolet SM--Einstein continuum needs
one of the following additional structures:
\begin{enumerate}
\item an infinite graded counterterm space with a norm under which the RG map
closes and has only finitely many unstable directions;
\item a non-Gaussian ultraviolet fixed point with a finite-dimensional
critical surface and rigorous control of all remaining operators; or
\item exact identities of the renewal construction that reduce the infinite
tower to controlled redundant directions.
\end{enumerate}
\end{corollary}

\begin{proof}
Theorem~\ref{thm:t10v17-gravity-power-counting} excludes closure of generic
perturbative regulator removal in a fixed finite operator list.  The three
listed mechanisms are the corresponding ways to control, rather than omit,
the additional directions.
\end{proof}

\subsection{Anomaly-compatible V17 gate}

\begin{definition}[V17 counterterm card]
\label{def:t10v17-counterterm-card}
A regulated RG step passes the card if:
\begin{enumerate}
\item its local breaking lies in $\ker s_1$ and its $H^1_s$ class vanishes;
\item the combined Ward/renormalization datum passes
\eqref{eq:t10v17-Fredholm-test};
\item the nonzero singular values needed in
\eqref{eq:t10v17-pseudoinverse} have a uniform lower bound;
\item the chiral regulator realizes the representation cancellations of
Proposition~\ref{prop:t10v17-SM-anomalies} and separately controls global
determinant-line holonomy;
\item the counterterm types in Definition~\ref{def:t10v17-counterterm-types}
are not mixed in the V16 shooting map; and
\item gravitational operators beyond the finite truncation are placed in a
declared graded space rather than silently set to zero.
\end{enumerate}
\end{definition}

\begin{theorem}[Compatibility with the shooting construction]
\label{thm:t10v17-shooting-compatibility}
If the V17 counterterm card holds and the chosen invariant coupling sector
satisfies the V16 trajectory card, the shot counterterms obey the regulated
Ward and renormalization conditions simultaneously.  No component of a
nontrivial anomaly class is absorbed into a relevant coupling.
\end{theorem}

\begin{proof}
Theorem~\ref{thm:t10v17-simultaneous} supplies the symmetry-restoring and
renormalization counterterm at each step, with the uniform norm controlled by
$\sigma_+^{-1}$.  Proposition~\ref{prop:t10v17-cohomology} and the typed
splitting prevent a nontrivial $H^1_s$ class from entering the solution.  The
remaining invariant coordinates are exactly those to which Theorem
\ref{thm:t10v16-shooting} applies.
\end{proof}

\begin{assessment}[Progress and upstream bottleneck after V17]
\label{assess:t10v17}
V17 turns anomaly cancellation into a finite solvability and norm test and
checks the Standard Model representation sums explicitly.  It also corrects
the scope of the preceding finite retained-flow model: Einstein power
counting produces an infinite derivative tower.  The next upstream task is
therefore functional analytic, not a finite beta-function calculation.  One
must construct a graded operator norm in which locality, nonlinear products,
reblocking, and the curved inverse all close, and then formulate the ultraviolet
fixed-point requirement on that space.
\end{assessment}

\begin{programme}[Tenth-series V18: graded gravitational operator space]
\label{prog:t10v17-V18}
V18 will define an infinite local-invariant sequence space weighted by
derivative order and field number.  It will prove a Cauchy-product estimate,
track the canonical rescaling eigenvalues, and state a checkable compactness
or contraction criterion for the high-order tail.  This will identify
precisely what a non-Gaussian ultraviolet fixed point must supply before the
V16 shooting theorem can be promoted from a finite truncation.
\end{programme}

\begin{firewall}[Claim boundary after tenth-series V17]
\label{fire:t10v17-boundary}
The representation sums do not constitute a regulator construction, and the
counterterm theorem is finite dimensional at each cutoff.  This version does
not control the infinite gravitational tower, prove an ultraviolet fixed
point, trivialize the determinant line, reconstruct the continuum measure,
or construct the full SM--Einstein continuum.
\end{firewall}

\paragraph{Parent-version record.}
V17 was formed from
\texttt{Renewal\_SM\_Einstein\_Continuum\_Research\_Notes\_V16.tex}.
The V16 parent had $6\,619$ validator-counted source lines, $250\,912$ bytes,
and SHA--256
\[
 \texttt{74ACBB033C72CA78A5755991ECF629AAEFC7F2146C073ACE7D1027C24495093E}.
\]

% =============================================================================
% END APPENDED TENTH-SERIES V17 MATERIAL
% =============================================================================

% =============================================================================
% APPENDED TENTH-SERIES V18 MATERIAL
% =============================================================================

\section{Tenth-series V18: a graded space for the gravitational operator
tower}
\label{sec:v10v18}

V17 showed that the derivative order allowed by Einstein power counting grows
with loop number.  The correct response is not to append finitely many more
couplings, but to place the whole tower in a norm that can be tested against
local products and canonical rescaling.  This version builds that abstract
space and proves a reserve-restoration criterion analogous to the geometric
polymer cycle of V15.

\subsection{Graded local invariants}

Let
\begin{equation}
 \mathfrak E=\prod_{n\geq0}E_n
 \label{eq:t10v18-graded-product}
\end{equation}
be a graded family of finite-dimensional spaces of local invariant densities,
modulo integrations by parts, algebraic and differential Bianchi identities,
and a declared set of redundant field-redefinition directions.  The grade
$n$ measures excess canonical dimension above a chosen finite low sector.
Field number and derivative number may be combined in $n$, but the choice
must make every allowed local composition obey the grade bounds stated below.

For $R>0$ define the analytic sequence norm
\begin{equation}
 \|c\|_R:=\sum_{n\geq0}R^n\|c_n\|_{E_n},
 \qquad
 \mathfrak E_R:=\{c:\|c\|_R<\infty\}.
 \label{eq:t10v18-analytic-norm}
\end{equation}

\begin{lemma}[Completeness and tail estimate]
\label{lem:t10v18-completeness}
$\mathfrak E_R$ is a Banach space.  If $R>1$, then
\begin{equation}
 \sum_{n>N}\|c_n\|_{E_n}
 \leq R^{-(N+1)}\|c\|_R.
 \label{eq:t10v18-tail-estimate}
\end{equation}
\end{lemma}

\begin{proof}
The map $c\mapsto(R^nc_n)_{n\geq0}$ identifies
$\mathfrak E_R$ isometrically with the $\ell^1$ sum of the Banach spaces
$E_n$, proving completeness.  For $n>N$, $R^{-n}\leq R^{-(N+1)}$; insert
this in the unweighted tail sum.
\end{proof}

\begin{proposition}[Compact loss of analytic radius]
\label{prop:t10v18-compact-embedding}
If $R_2>R_1>0$ and every $E_n$ is finite dimensional, the inclusion
\begin{equation}
 \mathfrak E_{R_2}\hookrightarrow\mathfrak E_{R_1}
 \label{eq:t10v18-compact-inclusion}
\end{equation}
is compact.
\end{proposition}

\begin{proof}
On the unit ball of $\mathfrak E_{R_2}$,
\[
 \sum_{n>N}R_1^n\|c_n\|
 \leq\left(\frac{R_1}{R_2}\right)^{N+1}.
\]
The projection onto $\bigoplus_{n=0}^NE_n$ has finite-dimensional bounded
image and hence is relatively compact.  The displayed uniform tail makes the
full unit ball totally bounded in $\mathfrak E_{R_1}$.
\end{proof}

\subsection{Cauchy products and analytic-radius loss}

Let a bilinear local operation have graded components
\begin{equation}
 (\mathcal B(c,d))_n
 =\sum_{i+j=n}\mathcal B_{ij}(c_i,d_j),
 \qquad
 \|\mathcal B_{ij}\|
 \leq C_0C_1^{i+j}.
 \label{eq:t10v18-graded-bilinear}
\end{equation}

\begin{theorem}[Graded Cauchy-product estimate]
\label{thm:t10v18-Cauchy-product}
For every $R>0$,
\begin{equation}
 \|\mathcal B(c,d)\|_R
 \leq C_0\|c\|_{C_1R}\|d\|_{C_1R}.
 \label{eq:t10v18-Cauchy-product}
\end{equation}
If $C_1=1$, every $\mathfrak E_R$ is a Banach algebra up to the constant
$C_0$.
\end{theorem}

\begin{proof}
Use \eqref{eq:t10v18-graded-bilinear} and rearrange the absolutely convergent
nonnegative majorant:
\begin{align*}
 \|\mathcal B(c,d)\|_R
 &\leq C_0\sum_{i,j\geq0}(C_1R)^{i+j}|c_i\|\|d_j\|\\
 &=C_0\|c\|_{C_1R}\|d\|_{C_1R}.
\end{align*}
\end{proof}

The same argument gives, for an $m$-linear local operation with norm bounded
by $C_mC_1^{i_1+\cdots+i_m}$,
\begin{equation}
 \|\mathcal B_m(c^{(1)},\ldots,c^{(m)})\|_R
 \leq C_m\prod_{a=1}^m\|c^{(a)}\|_{C_1R}.
 \label{eq:t10v18-multilinear}
\end{equation}

\begin{firewall}[Operator counting belongs inside $E_n$]
\label{fire:t10v18-operator-counting}
The number of tensor contractions and invariant monomials generally grows
with $n$.  The constant $C_1$ and the norms on $E_n$ must absorb that growth.
Writing one coefficient per derivative order without controlling the basis
multiplicity does not prove \eqref{eq:t10v18-graded-bilinear}.
\end{firewall}

\subsection{Canonical rescaling restores analytic reserve}

Let $\mathcal S_L$ be the linear canonical reblocking map at scale factor
$L>1$, with
\begin{equation}
 \| (\mathcal S_Lc)_n\|_{E_n}
 \leq L^{-n}\|c_n\|_{E_n}.
 \label{eq:t10v18-canonical-scaling}
\end{equation}

\begin{lemma}[Radius gain under canonical scaling]
\label{lem:t10v18-radius-gain}
If $R_s/L\leq R_w$, then
\begin{equation}
 \|\mathcal S_Lc\|_{R_s}\leq\|c\|_{R_w}.
 \label{eq:t10v18-radius-gain}
\end{equation}
\end{lemma}

\begin{proof}
By \eqref{eq:t10v18-canonical-scaling},
\[
 \|\mathcal S_Lc\|_{R_s}
 \leq\sum_n(R_s/L)^n\|c_n\|
 \leq\sum_nR_w^n\|c_n\|.
\]
\end{proof}

\begin{theorem}[One-step analytic reserve cycle]
\label{thm:t10v18-reserve-cycle}
Suppose $C_1<L$.  Choose radii
\begin{equation}
 C_1R_w\leq R_s\leq LR_w.
 \label{eq:t10v18-radius-window}
\end{equation}
Then the bilinear local operation maps two strong inputs in
$\mathfrak E_{R_s}$ to a weak output in $\mathfrak E_{R_w}$, and canonical
reblocking restores the strong radius:
\begin{equation}
 \|\mathcal S_L\mathcal B(c,d)\|_{R_s}
 \leq C_0\|c\|_{R_s}\|d\|_{R_s}.
 \label{eq:t10v18-reserve-cycle}
\end{equation}
Strict inequalities in \eqref{eq:t10v18-radius-window} leave reserve for
additional finite local operations.
\end{theorem}

\begin{proof}
Since $C_1R_w\leq R_s$, Theorem~\ref{thm:t10v18-Cauchy-product} bounds the
weak norm of $\mathcal B(c,d)$ by the two strong norms.  Since
$R_s/L\leq R_w$, Lemma~\ref{lem:t10v18-radius-gain} maps the weak output back
to the strong space.
\end{proof}

\begin{firewall}[Failure of the radius window]
\label{fire:t10v18-radius-failure}
If every valid invariant basis gives $C_1\geq L$, canonical dimension does not
restore the analytic reserve spent by one nonlinear step.  One must improve
the basis estimates, enlarge the block factor, exploit triangular selection
rules, or replace the analytic norm by a different graded scale.  Reusing the
same radius after an uncharged loss would make the iteration circular.
\end{firewall}

\subsection{Tail contraction and infinite-dimensional shooting}

Let $P_{\geq N}$ project onto grades $n\geq N$.  From
\eqref{eq:t10v18-canonical-scaling},
\begin{equation}
 \|\mathcal S_LP_{\geq N}\|_{\mathfrak E_R\to\mathfrak E_R}
 \leq L^{-N}.
 \label{eq:t10v18-tail-linear-contraction}
\end{equation}

\begin{proposition}[Nonlinear invariant tail ball]
\label{prop:t10v18-tail-ball}
Let a reblocked tail map on $V_N=P_{\geq N}\mathfrak E_{R_s}$ be
\begin{equation}
 \mathcal T(v)=e+\mathcal S_Lv+\mathcal Q(v,v)+\mathcal N(v),
 \label{eq:t10v18-tail-map}
\end{equation}
where
\begin{equation}
 \|\mathcal Q(v,w)\|\leq C_Q\|v\|\|w\|,
 \qquad
 \sup_{\|v\|\leq r}\|D\mathcal N(v)\|\leq\epsilon_N.
 \label{eq:t10v18-tail-nonlinear-bounds}
\end{equation}
If
\begin{align}
 L^{-N}+2C_Qr+\epsilon_N&<1,
 \label{eq:t10v18-tail-contraction}\\
 \|e\|+L^{-N}r+C_Qr^2+\epsilon_Nr&\leq r,
 \label{eq:t10v18-tail-invariance}
\end{align}
then $\mathcal T$ maps the closed radius-$r$ tail ball into itself and is a
strict contraction there.
\end{proposition}

\begin{proof}
The self-map bound follows term by term from
\eqref{eq:t10v18-tail-map}--\eqref{eq:t10v18-tail-nonlinear-bounds}.  For two
points in the ball,
\[
 \|\mathcal Q(v,v)-\mathcal Q(w,w)\|
 \leq2C_Qr\|v-w\|.
\]
Add the linear and remainder derivative bounds and use
\eqref{eq:t10v18-tail-contraction}.
\end{proof}

\begin{corollary}[Extension of the V16 stable space]
\label{cor:t10v18-V16-extension}
If the V14 polymer activity and the high graded operator tail are equipped
with the maximum product norm, and the sum of their cross-derivative debits
keeps the stable derivative below one, then the Banach space $X$ in
Theorem~\ref{thm:t10v16-shooting} may include the entire infinite tail
$V_N$.  Only grades below $N$ enter the finite relevant/marginal shooting
sector.
\end{corollary}

\begin{proof}
Proposition~\ref{prop:t10v18-tail-ball} supplies the tail contraction, while
the V14 gate supplies the polymer contraction.  The derivative norm of the
coupled map is bounded by the declared block and cross terms.  The V16 proof
uses only completeness and that derivative bound, so it applies to the
product Banach space.
\end{proof}

\subsection{What an ultraviolet fixed point must supply}

Let $\mathcal R:\mathfrak E_{R_s}\to\mathfrak E_{R_s}$ be the full local
operator part of one coarse-graining step and let $c_*$ be a fixed point.

\begin{definition}[Graded ultraviolet tail gate]
\label{def:t10v18-UV-tail-gate}
The fixed point passes the tail gate if there is a finite-rank projection
$P_{<N}$ such that:
\begin{enumerate}
\item $\mathcal R(c_*)=c_*$ and $c_*\in\mathfrak E_{R_s}$;
\item the high block
$P_{\geq N}D\mathcal R(c_*)P_{\geq N}$ has norm at most
$\lambda_N<1$;
\item the off-diagonal derivative blocks and nonlinear derivative variation
fit the V16 constants and Proposition~\ref{prop:t10v18-tail-ball}; and
\item locality, BRST cohomology, reflection, and the V13 boundary-action test
hold in the same graded norm.
\end{enumerate}
\end{definition}

\begin{theorem}[Finite critical data under the graded tail gate]
\label{thm:t10v18-finite-critical-data}
If Definition~\ref{def:t10v18-UV-tail-gate} and the V16 dichotomy hold, the
infinite high-order tower belongs to the contracted stable Banach sector.
The local infinite trajectory is then determined by finitely many relevant
shooting data, separately controlled marginal data, and the affine Ward
trajectory.
\end{theorem}

\begin{proof}
The high derivative block and nonlinear variation give a contracted tail
ball by Proposition~\ref{prop:t10v18-tail-ball}.  The complement has finite
rank.  Split that complement into the V16 relevant, marginal, and stable
finite blocks and apply Theorem~\ref{thm:t10v16-shooting}; transport the Ward
coordinates by the V15 affine ledger.
\end{proof}

\begin{firewall}[This is a criterion, not an asymptotic-safety proof]
\label{fire:t10v18-not-fixed-point-proof}
No fixed point $c_*$ or tail spectral bound $\lambda_N<1$ has been derived
here.  Canonical scaling proves tail contraction near a sufficiently small
Gaussian local map, but the Einstein coupling is not perturbatively
ultraviolet complete there.  At a non-Gaussian candidate, every row of
Definition~\ref{def:t10v18-UV-tail-gate} requires a new estimate.
\end{firewall}

\subsection{Analytic versus Gevrey growth}

\begin{proposition}[Diagnostic for factorial coefficient growth]
\label{prop:t10v18-factorial-diagnostic}
If a coefficient sequence satisfies
$\|c_n\|\geq C^n n!$ along an infinite subsequence, then
$c\notin\mathfrak E_R$ for every $R>0$.
\end{proposition}

\begin{proof}
The terms $R^n\|c_n\|$ do not tend to zero: by Stirling's formula,
$(RC)^nn!\to\infty$ for every $R,C>0$.
\end{proof}

Thus factorial perturbative growth, if present in the relevant graded
coefficients, requires a Borel/Gevrey formulation and a summation theorem;
the analytic sequence space cannot absorb it by choosing a smaller radius.

\begin{definition}[V18 graded-operator card]
\label{def:t10v18-graded-card}
A proposed regulator must provide:
\begin{enumerate}
\item finite bases and norms for every $E_n$;
\item explicit constants $C_0,C_1$ for every local multilinear operation;
\item a radius window \eqref{eq:t10v18-radius-window};
\item a uniform high-tail derivative bound, including mixing with the V14
polymer and V17 counterterm complexes; and
\item either analytic coefficient control or a separately proved
Gevrey/Borel replacement.
\end{enumerate}
\end{definition}

\begin{assessment}[Progress and bottleneck after V18]
\label{assess:t10v18}
V18 supplies the infinite-dimensional space missing from V17 and proves how
nonlinear analytic-radius loss can be restored by canonical reblocking.  It
also shows exactly what remains unknown: the constants depend on the concrete
basis of mixed curvature, gauge, scalar, fermion, ghost, and boundary
operators, and a non-Gaussian fixed point must satisfy a genuine high-tail
spectral bound.  These are physical estimates, not consequences of abstract
power counting.
\end{assessment}

\begin{programme}[Tenth-series V19: triangular generation and basis counting]
\label{prog:t10v18-V19}
V19 will refine the grade to derivative order and field number, prove the
triangular selection rules for local multiplication, contraction, and one
propagator insertion, and derive a sufficient exponential bound on the number
of invariant contractions.  This will test whether a finite $C_1<L$ radius
window is plausible or whether the programme must move directly to a Gevrey
scale.
\end{programme}

\begin{firewall}[Claim boundary after tenth-series V18]
\label{fire:t10v18-boundary}
The graded-space construction does not prove the required operator-counting
constants, factorial control, ultraviolet fixed point, determinant-line
trivialization, measure compactness, or the full SM--Einstein continuum.
\end{firewall}

\paragraph{Parent-version record.}
V18 was formed from
\texttt{Renewal\_SM\_Einstein\_Continuum\_Research\_Notes\_V17.tex}.
The V17 parent had $6\,980$ validator-counted source lines, $265\,786$ bytes,
and SHA--256
\[
 \texttt{02A09F422200E7E41D99F1C4077B6756B6778CDFCE43302C66D58FF90AD8C5D1}.
\]

% =============================================================================
% END APPENDED TENTH-SERIES V18 MATERIAL
% =============================================================================

% =============================================================================
% APPENDED TENTH-SERIES V19 MATERIAL
% =============================================================================

\section{Tenth-series V19: triangular jet generation and a Gevrey repair of
the operator norm}
\label{sec:v10v19}

V18 left one concrete test: do local invariant multiplicities and contraction
patterns grow only exponentially with grade?  The raw answer is generally
weaker.  Pairing $2m$ tensor slots already produces $(2m-1)!!$ contraction
patterns, which is factorial up to an exponential.  This does not prove that
the quotient by tensor identities has factorial dimension, but it prevents a
naive analytic estimate based on raw contractions.  V19 therefore refines
the grading and proves a Gevrey-weighted algebra that absorbs the raw pairing
count without hiding it.

\subsection{Jet letters and the triangular grade}

Fix a finite set of regulated field species $\Phi_a$.  A jet letter is
\begin{equation}
 X_{a,\alpha}:=D^\alpha\Phi_a,
 \qquad
 w(a,\alpha):=w_a+|\alpha|,
 \qquad w_a\in\mathbb N_{>0}.
 \label{eq:t10v19-jet-letter}
\end{equation}
For a local monomial $M=X_{a_1,\alpha_1}\cdots X_{a_m,\alpha_m}$, put
\begin{equation}
 p(M):=\sum_i|\alpha_i|,
 \quad q(M):=\sum_iw_{a_i},
 \quad n(M):=p(M)+q(M).
 \label{eq:t10v19-bigrade}
\end{equation}
Tensor-index contractions, invariant traces, and integration by parts are
performed only after this grade is assigned.

\begin{lemma}[Local triangular rules]
\label{lem:t10v19-triangular-rules}
The jet grade has the following properties.
\begin{enumerate}
\item Pointwise multiplication is additive:
$n(M_1M_2)=n(M_1)+n(M_2)$.
\item Algebraic contraction of tensor, gauge, or spinor indices preserves
$n$.
\item One covariant derivative maps grade $n$ to grade $n+1$, modulo
commutator terms whose curvatures are assigned the same total engineering
grade.
\item Integration by parts preserves the total grade of an integrated local
density.
\end{enumerate}
\end{lemma}

\begin{proof}
The first two statements are immediate from
\eqref{eq:t10v19-bigrade}.  A covariant derivative either adds one derivative
to a jet letter or, after commuting derivatives, replaces two derivatives by
a curvature action; curvature is assigned derivative weight two, so the
total is unchanged by the commutator rewrite.  Integration by parts moves one
derivative between factors and therefore preserves their sum.
\end{proof}

To include a Wick or fluctuation contraction, assign the rescaled covariance
jet contracting letters of weights $u,v$ the coefficient grade $u+v$.  A
relative Taylor moment of order $r$ has additional grade $r$.

\begin{proposition}[One-contraction triangularity]
\label{prop:t10v19-contraction-triangularity}
Let a covariance contraction remove two jet letters of weights $u,v$, insert
its covariance coefficient of grade $u+v$, and take a relative Taylor moment
of order $r$.  The total extended grade changes by exactly
\begin{equation}
 n_{\rm out}=n_{\rm in}+r.
 \label{eq:t10v19-contraction-grade}
\end{equation}
In particular, the Taylor order is the only upward grade shift in one
rescaled propagator insertion.
\end{proposition}

\begin{proof}
Removing the two letters subtracts $u+v$.  The covariance jet restores
$u+v$, and the Taylor moment adds $r$.
\end{proof}

\begin{firewall}[Scaling grade must be carried by the coefficient]
\label{fire:t10v19-coefficient-grade}
If the covariance jet is treated as a dimensionless constant, the contraction
appears to lower the grade and the subsequent reblocking estimate loses its
canonical powers.  Equation \eqref{eq:t10v19-contraction-grade} is valid only
in a dimensionless rescaled chart that records the covariance homogeneity in
the coefficient ledger.
\end{firewall}

\subsection{Raw contraction counting}

\begin{lemma}[Pairing count]
\label{lem:t10v19-pairing-count}
The number of complete pairings of $2m$ labelled slots is
\begin{equation}
 N_m=(2m-1)!!=\frac{(2m)!}{2^m m!},
 \label{eq:t10v19-pairing-count}
\end{equation}
and obeys
\begin{equation}
 N_m\leq2^m m!.
 \label{eq:t10v19-pairing-Gevrey}
\end{equation}
No bound $N_m\leq C^m$ can hold for all $m$.
\end{lemma}

\begin{proof}
Choose the partner of the first slot in $2m-1$ ways and continue recursively,
giving $(2m-1)(2m-3)\cdots1$.  The factorial identity follows.  Termwise,
$2j-1\leq2j$, so $N_m\leq\prod_{j=1}^m2j=2^mm!$.  Conversely,
$N_m\geq m!$, because $2j-1\geq j$; Stirling's formula excludes a uniform
exponential bound.
\end{proof}

Tensor symmetries, dimension-specific identities, Bianchi identities, and
fermionic antisymmetry can reduce this number.  Lemma
\ref{lem:t10v19-pairing-count} is therefore an upper bound on a raw
construction and a lower-bound diagnostic for that construction, not a
theorem that the space of independent invariants has factorial dimension.

\begin{proposition}[Sufficient analytic basis-count criterion]
\label{prop:t10v19-analytic-count}
Suppose the invariant algebra is generated by $g_j$ homogeneous generators
of grade $j\geq1$, with
\begin{equation}
 g_j\leq AB^j.
 \label{eq:t10v19-generator-count}
\end{equation}
Let $d_n$ be the dimension of its grade-$n$ quotient after any relations.
For every $0<t<B^{-1}$ there is a finite $C_t$ such that
\begin{equation}
 d_n\leq C_t t^{-n}.
 \label{eq:t10v19-exponential-dimension}
\end{equation}
\end{proposition}

\begin{proof}
Relations only reduce dimension, so the Hilbert series is coefficientwise
bounded by that of the free commutative algebra,
\[
 H(z)=\prod_{j\geq1}(1-z^j)^{-g_j}.
\]
For $0<t<B^{-1}$,
\[
 \log H(t)=\sum_{j\geq1}g_j\sum_{k\geq1}\frac{t^{jk}}k
 \leq A\sum_{j\geq1}B^j\frac{t^j}{1-t^j}<\infty.
\]
Since the coefficients are nonnegative, $d_nt^n\leq H(t)$, proving the
claim with $C_t=H(t)$.
\end{proof}

\begin{assessment}[Analytic-count verdict]
\label{assess:t10v19-analytic-verdict}
An exponential analytic basis bound follows from an exponential homogeneous
generator catalogue.  Such a catalogue has not been proved for the full
covariant SM--Einstein jet quotient.  Raw index-pair enumeration is
factorial, so V18's analytic card cannot be certified merely by listing raw
contractions.
\end{assessment}

\subsection{Gevrey sequence spaces}

For $R>0$ and $\sigma\geq0$, define
\begin{equation}
 \|c\|_{R,\sigma}
 :=\sum_{n\geq0}\frac{R^n}{(n!)^\sigma}\|c_n\|_{E_n},
 \qquad
 \mathfrak E_{R,\sigma}:=\{c:\|c\|_{R,\sigma}<\infty\}.
 \label{eq:t10v19-Gevrey-norm}
\end{equation}
The case $\sigma=0$ is V18's analytic space.  Positive $\sigma$ permits
coefficient growth of order $(n!)^\sigma$ up to an exponential.

\begin{lemma}[Gevrey completeness and compact radius loss]
\label{lem:t10v19-Gevrey-compactness}
$\mathfrak E_{R,\sigma}$ is Banach.  If $R_2>R_1>0$ and every $E_n$ is finite
dimensional, then
\begin{equation}
 \mathfrak E_{R_2,\sigma}\hookrightarrow
 \mathfrak E_{R_1,\sigma}
 \label{eq:t10v19-Gevrey-embedding}
\end{equation}
is compact.
\end{lemma}

\begin{proof}
Multiplication of the $n$-th component by $R^n/(n!)^\sigma$ identifies the
space with an $\ell^1$ sum.  On the unit ball at radius $R_2$, the tail in
the $R_1$ norm is at most $(R_1/R_2)^{N+1}$.  Finite-dimensional truncations
are relatively compact, exactly as in Proposition
\ref{prop:t10v18-compact-embedding}.
\end{proof}

\begin{theorem}[Gevrey Cauchy algebra]
\label{thm:t10v19-Gevrey-algebra}
Suppose a graded bilinear operation satisfies
\begin{equation}
 \|\mathcal B_{ij}\|
 \leq C_0C_1^{i+j}\binom{i+j}{i}^{\tau}
 \label{eq:t10v19-binomial-operation}
\end{equation}
for some $\tau\geq0$.  If $\sigma\geq\tau$, then
\begin{equation}
 \|\mathcal B(c,d)\|_{R,\sigma}
 \leq C_0\|c\|_{C_1R,\sigma}
          \|d\|_{C_1R,\sigma}.
 \label{eq:t10v19-Gevrey-product}
\end{equation}
\end{theorem}

\begin{proof}
For $n=i+j$, the ratio of the output factorial weight to the two input
factorial weights, including \eqref{eq:t10v19-binomial-operation}, is
\[
 \frac{i!^\sigma j!^\sigma}{n!^\sigma}
 \binom ni^\tau
 =\binom ni^{\tau-\sigma}\leq1.
\]
The remaining powers give $(C_1R)^{i+j}$, and the nonnegative double sum
factors into the two input norms.
\end{proof}

\begin{corollary}[Raw tensor pairings fit Gevrey order one]
\label{cor:t10v19-pairing-Gevrey}
If the only superexponential combinatorial loss at grade $m$ is the raw
pairing number $(2m-1)!!$, then it is bounded by an exponential times $m!$ and
can be charged to $\sigma=1$ in
\eqref{eq:t10v19-Gevrey-norm}.
\end{corollary}

\begin{proof}
Use \eqref{eq:t10v19-pairing-Gevrey}; the factor $2^m$ is an analytic-radius
loss and $m!$ is the Gevrey-one weight.
\end{proof}

\subsection{One propagator insertion as a trilinear convolution}

Let $k_r$ be the rescaled covariance/Taylor coefficient of order $r$.  After
the grading in Proposition~\ref{prop:t10v19-contraction-triangularity}, one
propagator insertion has the schematic homogeneous form
\begin{equation}
 (\mathcal P(c,d;k))_n
 =\sum_{i+j+r=n}\mathcal P_{ijr}(c_i,d_j,k_r).
 \label{eq:t10v19-propagator-convolution}
\end{equation}

\begin{theorem}[Gevrey propagator estimate]
\label{thm:t10v19-propagator-estimate}
Assume
\begin{equation}
 \|\mathcal P_{ijr}\|
 \leq K_0K_1^{i+j+r}
 \binom{i+j+r}{i,j,r}^{\tau}.
 \label{eq:t10v19-propagator-bound}
\end{equation}
If $\sigma\geq\tau$, then
\begin{equation}
 \|\mathcal P(c,d;k)\|_{R,\sigma}
 \leq K_0
 \|c\|_{K_1R,\sigma}
 \|d\|_{K_1R,\sigma}
 \|k\|_{K_1R,\sigma}.
 \label{eq:t10v19-propagator-estimate}
\end{equation}
\end{theorem}

\begin{proof}
For $n=i+j+r$, the multinomial coefficient is
$n!/(i!j!r!)$.  Multiplying
\eqref{eq:t10v19-propagator-bound} by the output factorial weight and dividing
by the three input weights leaves the factor
\[
 \binom{n}{i,j,r}^{\tau-\sigma}\leq1.
\]
The triple sum then factors.
\end{proof}

\begin{firewall}[The covariance jet is a required input norm]
\label{fire:t10v19-covariance-jet}
Theorem~\ref{thm:t10v19-propagator-estimate} does not prove
$\|k\|_{K_1R,\sigma}<\infty$.  That is a regulator-dependent short-distance
estimate on all covariance derivatives and Taylor moments.  Without it, the
formal local expansion does not define a bounded map on the tower.
\end{firewall}

\subsection{Gevrey reserve restoration}

Canonical scaling still satisfies
$\| (\mathcal S_Lc)_n\|\leq L^{-n}\|c_n\|$.  The factorial weights cancel
from the radius comparison.

\begin{theorem}[Gevrey radius cycle]
\label{thm:t10v19-Gevrey-cycle}
Suppose the largest analytic loss among the local product and propagator
operations is $C_*<L$.  Choose
\begin{equation}
 C_*R_w\leq R_s\leq LR_w.
 \label{eq:t10v19-Gevrey-window}
\end{equation}
If their combinatorial exponents are at most $\sigma$, the operations map
strong inputs at $(R_s,\sigma)$ to a weak output at $(R_w,\sigma)$, and
canonical reblocking restores the strong radius.  The resulting constants
are the product of the local operation constants and the covariance-jet
norms.
\end{theorem}

\begin{proof}
Apply Theorems~\ref{thm:t10v19-Gevrey-algebra} and
\ref{thm:t10v19-propagator-estimate} with $C_*R_w\leq R_s$.  The proof of
Lemma~\ref{lem:t10v18-radius-gain} is unchanged after inserting the common
factor $(n!)^{-\sigma}$ and uses $R_s/L\leq R_w$.
\end{proof}

\begin{firewall}[Gevrey control is not Borel summability]
\label{fire:t10v19-not-Borel}
A finite Gevrey norm bounds factorial growth of the graded operator
coefficients.  It does not construct a Borel transform, analytically continue
that transform, choose an unambiguous Laplace contour, or prove that the
result equals a nonperturbative measure.  Those are additional reconstruction
steps.
\end{firewall}

\begin{definition}[V19 operator-generation card]
\label{def:t10v19-generation-card}
A concrete regulator must now supply:
\begin{enumerate}
\item a jet alphabet and positive weights satisfying
Lemma~\ref{lem:t10v19-triangular-rules};
\item finite quotient spaces $E_n$ and either the analytic generator bound
\eqref{eq:t10v19-generator-count} or a Gevrey exponent $\sigma$;
\item multinomial operation bounds of the forms
\eqref{eq:t10v19-binomial-operation} and
\eqref{eq:t10v19-propagator-bound};
\item a finite covariance-jet norm; and
\item a strict radius window \eqref{eq:t10v19-Gevrey-window} compatible with
the V15 polymer reserve and V14 curved inverse.
\end{enumerate}
\end{definition}

\begin{assessment}[Progress and bottleneck after V19]
\label{assess:t10v19}
V19 proves the promised triangular rules and resolves the raw factorial
pairing loss at the level of functional analysis: Gevrey order one is enough
for complete pairings, and order $\sigma$ absorbs multinomial losses of order
$\sigma$.  The analytic V18 space remains available if a reduced invariant
generator catalogue proves exponential counting.  The active bottleneck is
now the all-order covariance-jet norm and its radius loss relative to the
block factor $L$.
\end{assessment}

\begin{programme}[Tenth-series V20: audit and covariance-jet estimates]
\label{prog:t10v19-V20}
V20 will perform the compulsory YM/NS audit.  It will then estimate all
derivatives and Taylor moments of a spectrally gapped finite-range block
covariance, first by a resolvent Cauchy bound and then in the Gevrey norm.  The
goal is an explicit $K_1$ to compare with $L$ in
\eqref{eq:t10v19-Gevrey-window}; if the gap closes in a retained massless
sector, that sector will be separated before the estimate.
\end{programme}

\begin{firewall}[Claim boundary after tenth-series V19]
\label{fire:t10v19-boundary}
This version does not prove the concrete invariant-basis bound, covariance
jet estimate, strict radius window, ultraviolet fixed point, Borel
summability, determinant-line trivialization, or the full SM--Einstein
continuum.
\end{firewall}

\paragraph{Parent-version record.}
V19 was formed from
\texttt{Renewal\_SM\_Einstein\_Continuum\_Research\_Notes\_V18.tex}.
The V18 parent had $7\,353$ validator-counted source lines, $279\,400$ bytes,
and SHA--256
\[
 \texttt{D8EF29CB7CC080046559CDE2705AD0ADB7CFCE7A968A7ED6F948F9AA228ED0BA}.
\]

% =============================================================================
% END APPENDED TENTH-SERIES V19 MATERIAL
% =============================================================================

% =============================================================================
% APPENDED TENTH-SERIES V20 MATERIAL
% =============================================================================

\section{Tenth-series V20: companion audit and a gapped covariance-jet
theorem}
\label{sec:v10v20}

\subsection{Compulsory companion audit}

The newest companion files inspected at this checkpoint were
\[
\begin{array}{c|r|l}
\text{file}&\text{bytes}&\text{SHA--256}\ \hline
\texttt{Renewal\_Yang\_Mills\_Mass\_Gap\_Research\_Notes\_V35.tex}
&562\,063&
\texttt{D6A8A03D083D9551BDE8AA925D4B7DEDA121C05E78D9D9FA1BF51F7CE8ABDC80}
\\
\texttt{Renewal\_NS\_Stability\_Research\_Notes\_V26.tex}
&278\,283&
\texttt{82C887F8C2C69E4F28FAD7C3DC8DE5B9099CB5A4BF431EBA1FFF3E91935978AD}.
\end{array}
\]
The NS note is unchanged.  YM V32--V35 transports its marked tree expansion
into the packet norm, exposes a complete derivative ledger, and distinguishes
three radii: a normalized Gaussian field radius, the shrinking compact-group
chart amplitude, and an analytic coefficient radius.  In particular, a
logarithmically growing normalized good-field ball is compatible with a
shrinking group chart when the coupling tends to zero.

\begin{assessment}[V20 audit transfer]
\label{assess:t10v20-audit-transfer}
The normalized field cutoff, nonlinear geometric chart amplitude, and
Gevrey operator radius in V19 are separate quantities.  A covariance-jet
estimate may depend on the last two but must not force the normalized Gaussian
ball to shrink.  The YM marked-tree constants do not transfer numerically.
Its massless-inverse warning agrees with V3 and V11: the gapped covariance
theorem below is applied only after low/closing modes are retained.
\end{assessment}

\subsection{A uniform analytic resolvent family}

Let $\Lambda$ be a finite graph of polynomial volume growth: for some fixed
$d,C_V$,
\begin{equation}
 |\{y:d(x,y)\leq r\}|\leq C_V(1+r)^d
 \label{eq:t10v20-volume-growth}
\end{equation}
uniformly in $x$, volume, and block scale after rescaling.  Let $H(z)$ act on
$\ell^2(\Lambda;\mathbb C^N)$, where
$z=(z_1,\ldots,z_p)$ lies in the complex polydisc
\begin{equation}
 \mathbb D_\rho^p:=\{z:|z_a|<\rho\}.
 \label{eq:t10v20-polydisc}
\end{equation}
Assume:
\begin{enumerate}
\item $H(z)$ is analytic and has graph range at most $R_0$;
\item $H(0)=H_0$ is self-adjoint and $H_0\succeq\gamma_0I$;
\item throughout the polydisc,
\begin{equation}
 \|H_0^{-1/2}(H(z)-H_0)H_0^{-1/2}\|\leq\eta<1;
 \label{eq:t10v20-relative-analyticity}
\end{equation}
and
\item for $W_{x,\vartheta}u(y)=e^{\vartheta d(x,y)}u(y)$, there is
$\vartheta_0>0$ such that
\begin{equation}
 \|W_{x,\vartheta}H(z)W_{x,\vartheta}^{-1}-H(z)\|
 \leq\frac12(1-\eta)\gamma_0
 \label{eq:t10v20-weighted-commutator}
\end{equation}
for all $x$, $0\leq\vartheta\leq\vartheta_0$, and $z\in\mathbb D_\rho^p$.
\end{enumerate}

The last row follows for a uniformly bounded finite-range operator once
$\vartheta_0$ is chosen so that its finite-range conjugation difference is
small.  It is displayed as a card because its constant must be checked for
the actual block operator.

\begin{lemma}[Uniform weighted inverse]
\label{lem:t10v20-weighted-inverse}
Under the preceding assumptions, $C(z):=H(z)^{-1}$ exists and
\begin{equation}
 \|W_{x,\vartheta}C(z)W_{x,\vartheta}^{-1}\|
 \leq M_C:=\frac{2}{(1-\eta)\gamma_0}
 \label{eq:t10v20-weighted-inverse}
\end{equation}
for $0\leq\vartheta\leq\vartheta_0$.
\end{lemma}

\begin{proof}
The relative bound \eqref{eq:t10v20-relative-analyticity} gives an inverse
with unweighted norm at most $[(1-\eta)\gamma_0]^{-1}$ by a Neumann series in
the $H_0$ metric.  Weighted conjugation changes $H(z)$ by an operator of norm
at most half the same coercive/accretive margin.  A second Neumann series
therefore gives \eqref{eq:t10v20-weighted-inverse}.
\end{proof}

\begin{corollary}[Uniform exponential kernel decay]
\label{cor:t10v20-exponential-kernel}
For every $0<\vartheta\leq\vartheta_0$,
\begin{equation}
 \|C(z;x,y)\|\leq M_Ce^{-\vartheta d(x,y)}.
 \label{eq:t10v20-exponential-kernel}
\end{equation}
\end{corollary}

\begin{proof}
Take the matrix element of the weighted inverse between coordinate delta
vectors based at $x$ and $y$.
\end{proof}

\subsection{Parameter derivatives and spatial moments}

Fix $0<\rho_0<\rho$ and $0<\vartheta<\vartheta_0$.  For a multi-index
$\alpha\in\mathbb N_0^p$, Cauchy's formula applied to the weighted resolvent
gives
\begin{equation}
 \|\partial_z^\alpha C(0;x,y)\|
 \leq M_C\alpha!\rho_0^{-|alpha|}e^{-\vartheta d(x,y)}.
 \label{eq:t10v20-Cauchy-kernel}
\end{equation}

\begin{lemma}[Exponential-to-moment conversion]
\label{lem:t10v20-moment-conversion}
For $n\geq0$ and $0<\vartheta'<\vartheta$,
\begin{equation}
 r^ne^{-\vartheta r}
 \leq n!(\vartheta-\vartheta')^{-n}e^{-\vartheta'r}
 \quad(r\geq0).
 \label{eq:t10v20-moment-conversion}
\end{equation}
Furthermore
\begin{equation}
 V_{\vartheta'}:=
 \sup_x\sum_y e^{-\vartheta'd(x,y)}<\infty
 \label{eq:t10v20-exponential-volume}
\end{equation}
under \eqref{eq:t10v20-volume-growth}.
\end{lemma}

\begin{proof}
The exponential series gives
$e^{(\vartheta-\vartheta')r}\geq
[(\vartheta-\vartheta')r]^n/n!$, proving the first statement.  Decompose the
graph into integer distance shells.  Their cardinalities grow at most
polynomially by \eqref{eq:t10v20-volume-growth}, while
$e^{-\vartheta'r}$ decays exponentially.
\end{proof}

For $m,n\geq0$, define the aggregated covariance jet
\begin{equation}
 K_{m,n}:=
 \sum_{|\alpha|=m}\sup_x\sum_y
 d(x,y)^n\|\partial_z^\alpha C(0;x,y)\|.
 \label{eq:t10v20-jet-stock}
\end{equation}

\begin{theorem}[Double Gevrey-one covariance-jet bound]
\label{thm:t10v20-covariance-jet}
Under the V20 resolvent assumptions,
\begin{equation}
 K_{m,n}
 \leq
 M_CV_{\vartheta'}
 \binom{m+p-1}{p-1}
 m!n!\rho_0^{-m}(\vartheta-\vartheta')^{-n}.
 \label{eq:t10v20-jet-bound}
\end{equation}
Consequently the double Gevrey stock
\begin{equation}
 \mathcal K(R_z,R_x)
 :=\sum_{m,n\geq0}
 \frac{R_z^mR_x^n}{m!n!}K_{m,n}
 \label{eq:t10v20-double-Gevrey}
\end{equation}
is finite whenever
\begin{equation}
 R_z<\rho_0,
 \qquad R_x<\vartheta-\vartheta'.
 \label{eq:t10v20-Gevrey-window}
\end{equation}
More explicitly,
\begin{equation}
 \mathcal K(R_z,R_x)
 \leq
 \frac{M_CV_{\vartheta'}}
 {(1-R_z/\rho_0)^p
  (1-R_x/(\vartheta-\vartheta'))}.
 \label{eq:t10v20-Gevrey-explicit}
\end{equation}
\end{theorem}

\begin{proof}
Combine \eqref{eq:t10v20-Cauchy-kernel} with
Lemma~\ref{lem:t10v20-moment-conversion}.  For $|\alpha|=m$,
$\alpha!\leq m!$, and the number of such multi-indices is
$\binom{m+p-1}{p-1}$, proving \eqref{eq:t10v20-jet-bound}.  Divide by
$m!n!$, sum, and use
\[
 \sum_{m\geq0}\binom{m+p-1}{p-1}t^m=(1-t)^{-p}
\]
together with the geometric series in $n$.
\end{proof}

\begin{corollary}[Resolution of the V19 covariance-jet input]
\label{cor:t10v20-V19-input}
For a uniformly gapped finite-range interior block covariance, the kernel
sequence $k_r$ required by Theorem~\ref{thm:t10v19-propagator-estimate} has a
finite Gevrey-one norm on every strict subwindow of
\eqref{eq:t10v20-Gevrey-window}.  The analytic loss constants are the
reciprocals of the parameter analyticity radius and exponential localization
reserve, not the normalized Gaussian field cutoff.
\end{corollary}

\begin{proof}
Taylor coefficients are parameter derivatives, while relative-coordinate
Taylor moments are bounded by the spatial moments in
\eqref{eq:t10v20-jet-stock}.  Theorem~\ref{thm:t10v20-covariance-jet} controls
their joint Gevrey-one sum.  The final statement follows from the three-radius
distinction in the audit.
\end{proof}

\subsection{Finite differences and determinant owners}

\begin{lemma}[Finite-difference jets]
\label{lem:t10v20-finite-differences}
If $\nabla_eK(x,y)=K(x+e,y)-K(x,y)$ is a unit finite difference, then an
$r$-fold difference of the covariance kernel satisfies the same exponential
bound with an additional factor at most $2^r e^{r\vartheta}$.  Thus finite
differences produce an analytic radius loss and no extra factorial order.
\end{lemma}

\begin{proof}
Expand the $r$ differences into at most $2^r$ translated kernels.  Each unit
translation changes $d(x,y)$ by at most one, costing $e^\vartheta$ in the
exponential majorant.
\end{proof}

\begin{proposition}[Analytic determinant-density jet]
\label{prop:t10v20-logdet}
On one finite interior block, choose the analytic branch of
$D(z)=\log\det H(z)$ through $z=0$.  Under
\eqref{eq:t10v20-relative-analyticity}, $D$ is analytic in
$\mathbb D_\rho^p$, and for every smaller polydisc
$\mathbb D_{\rho_0}^p$,
\begin{equation}
 |\partial^\alpha D(0)|
 \leq \alpha!\rho_0^{-|alpha|}
 \sup_{z\in\partial\mathbb D_{\rho_0}^p}|D(z)|.
 \label{eq:t10v20-logdet-Cauchy}
\end{equation}
The supremum is finite at fixed block dimension and must be charged per block
to the V5 density owner.
\end{proposition}

\begin{proof}
The relative Neumann condition keeps $H(z)$ invertible throughout the
polydisc, so its determinant is nonzero and has an analytic logarithm on the
simply connected polydisc after fixing the value at zero.  Multivariate
Cauchy estimates give the derivative bound.
\end{proof}

\begin{firewall}[Block dimension is not a uniform-volume constant]
\label{fire:t10v20-logdet-volume}
The bound on $\log\det H$ grows with the number of integrated interior
variables.  Uniform continuum control comes from assigning the extensive
part to local block/vacuum density coordinates and bounding the connected
remainder.  Proposition~\ref{prop:t10v20-logdet} supplies factorial derivative
control but does not remove that extensive owner.
\end{firewall}

\subsection{Massless and closing modes}

\begin{proposition}[Gap dependence cannot be removed]
\label{prop:t10v20-gap-no-go}
The constants in Theorem~\ref{thm:t10v20-covariance-jet} contain
$M_C\geq[(1-\eta)\gamma_0]^{-1}$.  Hence no volume-uniform version follows
from this theorem when the complementary spectral gap $\gamma_0$ tends to
zero.
\end{proposition}

\begin{proof}
Already the zeroth resolvent norm is at least the reciprocal of the smallest
eigenvalue.  All higher Cauchy and moment bounds inherit that norm.
\end{proof}

\begin{definition}[Gapped covariance split]
\label{def:t10v20-gapped-split}
Before applying the covariance-jet theorem, decompose one block into:
\begin{enumerate}
\item zero-trace interior fields with the V4 Dirichlet/Poincar\'e gap;
\item retained shared boundary, harmonic, closing, and V14 low spectral
coordinates; and
\item the V6--V8 annular physical projector, whose nonlocal packet estimates
are charged separately rather than called finite range.
\end{enumerate}
Only the first sector uses $H(z)^{-1}$ in
Theorem~\ref{thm:t10v20-covariance-jet}.
\end{definition}

\begin{firewall}[No growing-volume finite-rank shortcut]
\label{fire:t10v20-no-finite-rank-shortcut}
V3 proved that a uniformly bounded finite retained rank cannot coexist with a
uniform gap on expanding massless tori.  Definition
\ref{def:t10v20-gapped-split} is block-local: retained boundary/low variables
are reblocked at every scale and are not replaced by one global fixed-rank
sector.
\end{firewall}

\subsection{Insertion into the graded radius cycle}

Let $C_{\rm alg}$ be the largest analytic loss from invariant multiplication,
finite differences, projectors, and owner reanchoring.  Let
\begin{equation}
 C_{\rm cov}:=max\{\rho_0^{-1},
 (\vartheta-\vartheta')^{-1}\}
 \label{eq:t10v20-covariance-loss}
\end{equation}
after dimensionless normalization, and put
$C_*:=\max\{C_{\rm alg},C_{\rm cov}\}$.

\begin{theorem}[Covariance-complete Gevrey radius gate]
\label{thm:t10v20-radius-gate}
If
\begin{equation}
 C_*<L
 \label{eq:t10v20-strict-radius-condition}
\end{equation}
and the local operation constants and $M_CV_{\vartheta'}$ are uniform, there
are radii $R_w,R_s$ satisfying the V19 Gevrey cycle.  One local product, one
gapped covariance insertion with arbitrary parameter/spatial jet order, and
canonical reblocking then map the strong graded space to itself with a finite
uniform constant.
\end{theorem}

\begin{proof}
Choose $R_s,R_w$ so that
$C_*R_w<R_s<LR_w$.  Theorems
\ref{thm:t10v19-Gevrey-algebra},
\ref{thm:t10v19-propagator-estimate}, and
\ref{thm:t10v20-covariance-jet} map strong inputs to the weak radius.
Canonical scaling restores the strong radius by Theorem
\ref{thm:t10v19-Gevrey-cycle}.
\end{proof}

\begin{firewall}[Strict radius condition remains regulator dependent]
\label{fire:t10v20-radius-condition}
The theorem computes the form of the covariance loss but does not prove
$C_*<L$ for a concrete regulator.  Increasing $L$ changes block geometry,
aggregation constants, and the retained sector, so it cannot be invoked as a
free remedy without rechecking the complete V14--V19 cards.
\end{firewall}

\begin{assessment}[Progress and bottleneck after V20]
\label{assess:t10v20}
V20 proves the all-order covariance-jet estimate requested in V19: analytic
resolvent dependence and Combes--Thomas localization give a double
Gevrey-one bound in parameter order and spatial moment.  The massless gap
obstruction is isolated before inversion.  The next missing operation is the
iteration of arbitrarily many connected covariance insertions.  A single
propagator theorem does not control the graph/forest sum or show that its
Gevrey radius loss is summable.
\end{assessment}

\begin{programme}[Tenth-series V21: connected graph factorials]
\label{prog:t10v20-V21}
V21 will bound the number and norm of connected covariance graphs with
labelled vertices.  It will compare a direct connected-graph sum, which has
too much combinatorial growth, with a tree/forest interpolation carrying the
V20 covariance jet on its edges.  The target is a Gevrey order and small
activity condition that remain stable under the V15 polymer reserve cycle.
\end{programme}

\begin{firewall}[Claim boundary after tenth-series V20]
\label{fire:t10v20-boundary}
This version does not prove the connected graph sum, the concrete strict
radius condition, the ultraviolet fixed point, determinant-line
trivialization, cofinal measure compactness, or the full SM--Einstein
continuum.
\end{firewall}

\paragraph{Parent-version record.}
V20 was formed from
\texttt{Renewal\_SM\_Einstein\_Continuum\_Research\_Notes\_V19.tex}.
The V19 parent had $7\,738$ validator-counted source lines, $293\,646$ bytes,
and SHA--256
\[
 \texttt{CDAA0F392F65BA79C7E9317EA50238196450670A5C1F68119E172BA39AE503A8}.
\]

% =============================================================================
% END APPENDED TENTH-SERIES V20 MATERIAL
% =============================================================================

% =============================================================================
% APPENDED TENTH-SERIES V21 MATERIAL
% =============================================================================

\section{Tenth-series V21: connected covariance graphs and the tree
reorganization}
\label{sec:v10v21}

V20 controls one gapped covariance insertion to all jet orders.  Iterating
that bound by summing arbitrary connected graphs is not legitimate: the
number of labelled graphs grows like the exponential of $n^2$.  This version
reorganizes the connected part by a forest interpolation.  Labelled trees
have only factorial-times-exponential growth, which the V19 Gevrey-one norm
and the $1/n!$ activity expansion can absorb.

\subsection{Why direct connected-graph counting fails}

\begin{lemma}[Graph count exceeds every fixed Gevrey order]
\label{lem:t10v21-graph-no-go}
The number of simple graphs on $n$ labelled vertices is
$2^{\binom n2}$.  For every fixed $\sigma,C>0$,
\begin{equation}
 \frac{2^{\binom n2}}{C^n(n!)^\sigma}\longrightarrow\infty.
 \label{eq:t10v21-graph-no-go}
\end{equation}
The same conclusion holds for connected graphs.
\end{lemma}

\begin{proof}
Every one of the $\binom n2$ possible edges may be present or absent.  Taking
logarithms, the numerator has leading growth $(\log2)n^2/2$, whereas
$\log(C^n(n!)^\sigma)=O(n\log n)$.  For connected graphs, fix a labelled
spanning star and choose every remaining edge freely.  This gives at least
$2^{\binom{n-1}2}$ connected graphs, with the same comparison.
\end{proof}

\begin{firewall}[The exponential series does not cure graph counting]
\label{fire:t10v21-factorial-not-enough}
The factor $1/n!$ from expanding an exponential changes the denominator in
\eqref{eq:t10v21-graph-no-go} by one factorial.  It still cannot absorb
$e^{cn^2}$.  A direct absolute sum over connected graphs is therefore not a
Gevrey proof at any fixed order.
\end{firewall}

\subsection{Forest interpolation}

Take $n$ replicated field copies.  For $1\leq i<j\leq n$, let $x_{ij}$
weaken the off-diagonal covariance between replicas $i,j$.  For a forest
$F$ with edge parameters $w=(w_e)_{e\in F}\in[0,1]^F$, define
\begin{equation}
 x^F_{ij}(w)=
 \begin{cases}
 \min\{w_e:e\text{ lies on the }F\text{-path from }i\text{ to }j\},
   &i,j\text{ connected in }F,\\
 0,&\text{otherwise},
 \end{cases}
 \label{eq:t10v21-forest-matrix}
\end{equation}
with $x^F_{ii}=1$.

\begin{lemma}[Positivity of the forest matrix]
\label{lem:t10v21-forest-positive}
For every forest $F$ and $w\in[0,1]^F$, the matrix
$x^F(w)$ is positive semidefinite.
\end{lemma}

\begin{proof}
For $t\in[0,1]$, delete from $F$ all edges with $w_e<t$.  The resulting
connected components define a block matrix $P_t$ whose $(i,j)$ entry is one
when $i,j$ lie in the same component and zero otherwise.  Each $P_t$ is a sum
of rank-one positive block indicators.  Moreover
\[
 x^F_{ij}(w)=\int_0^1(P_t)_{ij}dt.
\]
An integral of positive semidefinite matrices is positive semidefinite.
\end{proof}

\begin{theorem}[Forest interpolation formula]
\label{thm:t10v21-forest-formula}
Let $f$ be smooth in the variables $x_{ij}$ on a neighbourhood of the cube
of correlation matrices generated by \eqref{eq:t10v21-forest-matrix}.  Then
\begin{equation}
 f(\mathbf1)=
 \sum_{F\ {
m forest}}
 \int_{[0,1]^F}
 \left(\prod_{e\in F}\partial_{x_e}\right)
 f(x^F(w))\,dw,
 \label{eq:t10v21-forest-formula}
\end{equation}
where $\mathbf1$ has every off-diagonal entry one.  After applying the
connected-part projection on the replica labels, only spanning trees remain.
\end{theorem}

\begin{proof}
Apply the fundamental theorem of calculus successively to the off-diagonal
couplings, but group terms according to the partition of replica labels
created by the chosen edges.  An edge is differentiated only when it joins
two previously distinct components; edges which would create a cycle are
represented by the minimum along the already chosen path.  Induction on the
number of variables gives \eqref{eq:t10v21-forest-formula}: the empty forest
is the fully decoupled endpoint, and differentiating with respect to the next
joining edge creates every larger forest exactly once.  The connected-part
projection is the M\"obius projection on set partitions.  A forest survives
it exactly when it has one component, hence is a spanning tree.
\end{proof}

For a replicated Gaussian, differentiating an off-diagonal covariance gives
the edge contraction
\begin{equation}
 \Delta_{ij,C}
 :=\sum_{a,b}C_{ab}
 \frac{\partial}{\partial\phi_{i,a}}
 \frac{\partial}{\partial\phi_{j,b}}.
 \label{eq:t10v21-edge-operator}
\end{equation}
Lemma~\ref{lem:t10v21-forest-positive} ensures that the weakened bosonic
covariance remains positive semidefinite.

\begin{corollary}[Tree formula for Gaussian cumulants]
\label{cor:t10v21-cumulant-tree}
For sufficiently regular local vertices $V_1,\ldots,V_n$ under a positive
Gaussian covariance $C$,
\begin{equation}
 \kappa_C(V_1,\ldots,V_n)
 =\sum_{T\ {
m tree}}
 \int_{[0,1]^T}
 \mathbb E_{C^T(w)}
 \left[
  \prod_{\{i,j\}\in T}\Delta_{ij,C}
  \prod_{i=1}^nV_i(\phi_i)
 \right]dw.
 \label{eq:t10v21-cumulant-tree}
\end{equation}
\end{corollary}

\begin{proof}
Apply Theorem~\ref{thm:t10v21-forest-formula} to the replicated Gaussian
expectation.  Gaussian differentiation with respect to an off-diagonal
covariance is \eqref{eq:t10v21-edge-operator}.  The connected-part projection
is the joint cumulant and retains only spanning trees.
\end{proof}

\begin{firewall}[Positive reduced covariance first]
\label{fire:t10v21-positive-first}
The forest formula is applied only after the V9 stationary conformal reduction
and the V4 gauge quotient have produced a positive physical covariance.
Using a real Gaussian interpolation in the unreduced negative conformal
direction would be divergent.
\end{firewall}

\subsection{Tree combinatorics and the activity sum}

\begin{lemma}[Cayley growth is Gevrey one]
\label{lem:t10v21-Cayley}
The number of spanning trees on $n\geq2$ labelled vertices is $n^{n-2}$, and
\begin{equation}
 \frac{n^{n-2}}{n!}\leq e^n.
 \label{eq:t10v21-Cayley-bound}
\end{equation}
\end{lemma}

\begin{proof}
Cayley's Pr\"ufer-code bijection gives $n^{n-2}$.  The elementary Stirling
lower bound $n!\geq(n/e)^n$ yields
$n^{n-2}/n!\leq e^n/n^2\leq e^n$.
\end{proof}

Let $z$ be the complete strong Gevrey/polymer norm of one local interaction
vertex after retained Taylor extraction.  Let $K_T$ be the norm of one V20
gapped covariance edge together with its finite type, packet-owner, and
weakened-expectation factors.  The definition of $K_T$ includes a Gram bound
for any unexpanded fermionic remainder and the V15 weak-radius hull owner.

\begin{theorem}[Absolute connected tree bound]
\label{thm:t10v21-connected-bound}
Assume that every fixed labelled tree integrand in
\eqref{eq:t10v21-cumulant-tree} obeys
\begin{equation}
 \|\mathcal A_T(V_1,\ldots,V_n)\|_{\rm weak}
 \leq z^nK_T^{n-1}
 \label{eq:t10v21-fixed-tree-bound}
\end{equation}
in the full V19 Gevrey-one and V15 polymer norm.  If
\begin{equation}
 q_T:=ezK_T<1,
 \label{eq:t10v21-tree-smallness}
\end{equation}
then the connected exponential series converges absolutely and
\begin{equation}
 \left\|
 \sum_{n\geq1}\frac1{n!}\kappa_C(V,\ldots,V)
 \right\|_{\rm weak}
 \leq z+\frac{e^2z^2K_T}{1-q_T}.
 \label{eq:t10v21-connected-bound}
\end{equation}
\end{theorem}

\begin{proof}
The $n=1$ term is at most $z$.  For $n\geq2$, sum
\eqref{eq:t10v21-fixed-tree-bound} over the $n^{n-2}$ labelled trees and use
\eqref{eq:t10v21-Cayley-bound}:
\[
 \frac{n^{n-2}}{n!}z^nK_T^{n-1}
 \leq K_T^{-1}(ezK_T)^n.
\]
The geometric sum from $n=2$ is
$K_T^{-1}q_T^2/(1-q_T)=e^2z^2K_T/(1-q_T)$.
\end{proof}

\begin{remark}[Zero edge norm]
If $K_T=0$, all connected terms with $n\geq2$ vanish and
\eqref{eq:t10v21-connected-bound} is understood by continuity.
\end{remark}

\subsection{Why the Gevrey radius is lost only once}

For a fixed tree, collect simultaneously all vertex grades, covariance-jet
grades, and Taylor moment grades before applying the weak norm.  The total
output grade is their sum by Proposition
\ref{prop:t10v19-contraction-triangularity}.

\begin{proposition}[Global tree convolution]
\label{prop:t10v21-global-convolution}
Suppose every local tensor operation in a tree has multinomial exponent at
most one and total analytic loss $C_T^{N_{\rm grade}}$, depending on the sum
of all input grades rather than on the depth of an association.  If
\begin{equation}
 C_TR_w\leq R_s,
 \label{eq:t10v21-tree-radius}
\end{equation}
then a fixed tree with strong-radius vertices and covariance jets maps
directly to the weak radius with the bound
\eqref{eq:t10v21-fixed-tree-bound}.  The radius loss is $C_T$ once, not
$C_T^{n-1}$ under repeated associations.
\end{proposition}

\begin{proof}
Expand the whole tree as one multilinear convolution.  If the total grade is
$N=n_1+\cdots+n_s$, the output factor $R_w^N/N!$ times the global multinomial
coefficient is bounded by the product of the input factors
$R_s^{n_i}/n_i!$ because $C_TR_w\leq R_s$.  Summing all grades factors into
the strong input norms.  No intermediate weak-radius object is reused as a
strong input.
\end{proof}

\begin{firewall}[Iterated binary estimates would be circular]
\label{fire:t10v21-binary-radius}
Applying the V19 strong-to-weak binary estimate along the tree one edge at a
time would require feeding a weak output back into a strong input slot.  That
is not justified.  Proposition~\ref{prop:t10v21-global-convolution} is the
required all-tree estimate and its hypothesis must be checked in the concrete
tensor grammar.
\end{firewall}

\subsection{Fermionic determinant rows}

Let the chiral covariance have a Gram representation
\begin{equation}
 C_F(x,y)=\langle u_x,v_y\rangle_{\mathcal H},
 \qquad
 \|u_x\|,\|v_y\|\leq G_F.
 \label{eq:t10v21-Gram-representation}
\end{equation}

\begin{lemma}[Unexpanded fermion contractions]
\label{lem:t10v21-Gram-bound}
For every $m\times m$ matrix
$M_{ij}=C_F(x_i,y_j)$,
\begin{equation}
 |\det M|\leq G_F^{2m}.
 \label{eq:t10v21-Gram-bound}
\end{equation}
\end{lemma}

\begin{proof}
The Gram--Hadamard inequality bounds the determinant by the product of the
norms of the $u_{x_i}$ and $v_{y_j}$ vectors.
\end{proof}

\begin{firewall}[Chiral phase is not controlled by an absolute Gram bound]
\label{fire:t10v21-chiral-phase}
Lemma~\ref{lem:t10v21-Gram-bound} prevents an extra $m!$ from expanding all
fermion pairings.  It controls the absolute determinant magnitude only.  The
phase, gauge covariance, anomaly cocycle, and determinant-line gluing remain
the V17 problem and cannot be inferred from this inequality.
\end{firewall}

\subsection{Uniform parameter jets of the connected sum}

\begin{proposition}[Cauchy control after the tree sum]
\label{prop:t10v21-connected-jets}
Suppose $z(\zeta)$ and $K_T(\zeta)$ are analytic in a complex polydisc
$|\zeta_a|<\rho$ and
\begin{equation}
 \sup_{\zeta}e|z(\zeta)|K_T(\zeta)\leq q_*<1.
 \label{eq:t10v21-complex-branch}
\end{equation}
If the fixed-tree bound is uniform there, the connected sum is analytic and
for $0<\rho_0<\rho$,
\begin{equation}
 \|\partial_\zeta^\alpha\mathcal C(0)\|_{\rm weak}
 \leq\alpha!\rho_0^{-|\alpha|}
 \sup_{|\zeta|<\rho_0}
 \left[|z(\zeta)|+
 \frac{e^2|z(\zeta)|^2|K_T(\zeta)|}{1-q_*}\right].
 \label{eq:t10v21-connected-Cauchy}
\end{equation}
Thus summing all tree sizes does not add a second Gevrey factorial.
\end{proposition}

\begin{proof}
Theorem~\ref{thm:t10v21-connected-bound} gives uniform absolute convergence
on every smaller polydisc under \eqref{eq:t10v21-complex-branch}.  The sum is
therefore analytic.  Multivariate Cauchy estimates give
\eqref{eq:t10v21-connected-Cauchy}.
\end{proof}

\subsection{One connected-step card}

\begin{definition}[V21 forest card]
\label{def:t10v21-forest-card}
One reduced block passes the forest card if:
\begin{enumerate}
\item the V9 reduced covariance is positive and its gapped interior part
passes the V20 double Gevrey covariance-jet theorem;
\item the local vertex norm $z$ includes all good, transition, determinant,
ghost, and declared boundary types after retained extraction;
\item every tree integrand passes the global convolution estimate
\eqref{eq:t10v21-fixed-tree-bound} at one weak radius;
\item the complex branch number satisfies \eqref{eq:t10v21-complex-branch};
\item fermionic leftovers have a uniform Gram representation and the chiral
phase is carried separately by the V17 determinant-line ledger; and
\item the weak connected output is reblocked through the V15 reserve cycle
and fits inside the V16 stable derivative margin.
\end{enumerate}
\end{definition}

\begin{theorem}[Connected small-activity closure]
\label{thm:t10v21-connected-closure}
If the V21 forest card holds uniformly, arbitrary connected gapped covariance
insertions define an analytic map on the Gevrey-one polymer ball.  Its norm
and all finite parameter jets obey
\eqref{eq:t10v21-connected-bound} and
\eqref{eq:t10v21-connected-Cauchy}.  After reserve-restoring reblocking, this
connected contribution may be inserted as a finite explicit debit in the
V14/V16 contraction and shooting inequalities.
\end{theorem}

\begin{proof}
The forest formula replaces the connected graph sum by trees.  Theorem
\ref{thm:t10v21-connected-bound} sums their sizes, Proposition
\ref{prop:t10v21-global-convolution} controls the grade radius, Lemma
\ref{lem:t10v21-Gram-bound} controls unexpanded fermion contractions, and
Proposition~\ref{prop:t10v21-connected-jets} controls parameter jets.  The
last card row supplies the V15 reblocking and V16 insertion.
\end{proof}

\begin{assessment}[Progress and bottleneck after V21]
\label{assess:t10v21}
V21 removes the $e^{cn^2}$ direct-graph obstruction and proves that a
tree-organized connected integral is compatible with the Gevrey-one tower
under the explicit branch condition $ezK_T<1$.  The remaining difficulty is
global field amplitude: the fixed-tree norm must include the complement of
the convex/good chart.  Without a uniform large-field coercivity estimate,
the small-activity branch condition does not define the complete regulated
measure.
\end{assessment}

\begin{programme}[Tenth-series V22: large-field coercivity and chart cover]
\label{prog:t10v21-V22}
V22 will construct a finite good/transition/bad chart cover for one compact
gauge--scalar--metric block.  It will prove an abstract large-field estimate
from a coercive physical action plus determinant-density bounds, quantify the
bad-component activity, and test whether the negative conformal direction can
remain entirely in the stationary reduction rather than the real integration
domain.
\end{programme}

\begin{firewall}[Claim boundary after tenth-series V21]
\label{fire:t10v21-boundary}
The forest theorem is conditional on fixed-tree analytic bounds and
small activity.  It does not prove the large-field estimate, the chiral phase,
the concrete ultraviolet fixed point, cofinal measure compactness, or the full
SM--Einstein continuum.
\end{firewall}

\paragraph{Parent-version record.}
V21 was formed from
\texttt{Renewal\_SM\_Einstein\_Continuum\_Research\_Notes\_V20.tex}.
The V20 parent had $8\,134$ validator-counted source lines, $308\,314$ bytes,
and SHA--256
\[
 \texttt{3BEED1F32C62A63AF8C0AB778792C8BC03444BE2074DD9F7051A0AF2BD84DFEB}.
\]

% =============================================================================
% END APPENDED TENTH-SERIES V21 MATERIAL
% =============================================================================

% =============================================================================
% APPENDED TENTH-SERIES V22 MATERIAL
% =============================================================================

\section{Tenth-series V22: large-field coercivity, chart covers, and the
conformal obstruction}
\label{sec:v10v22}

The V21 tree bound controls the connected expansion only on a chart where
the local activity is small.  This version supplies an abstract large-field
completion for coercive reduced actions and then tests its hardest hypothesis
against Euclidean Einstein gravity.  The conformal factor is genuinely
unbounded on a real integration contour, so the V9 stationary reduction must
be globalized; it cannot be replaced by a generic large-field estimate.

\subsection{A finite-dimensional coercive tail theorem}

\begin{theorem}[Coercive density tail]
\label{thm:t10v22-coercive-tail}
Let $x\in\mathbb R^N$, and let a nonnegative density be
\begin{equation}
 d\mu(x)=Z^{-1}J(x)e^{-S(x)}dx.
 \label{eq:t10v22-density}
\end{equation}
Suppose for constants $a>0,b\geq0,c\geq0$, $p>q\geq0$,
\begin{equation}
 S(x)\geq a\|x\|^p-b,
 \qquad
 J(x)\leq C_J e^{c\|x\|^q}.
 \label{eq:t10v22-coercive-hypotheses}
\end{equation}
Then $Z<\infty$ and there are $R_0,C,\delta>0$ such that
\begin{equation}
 \int_{\|x\|\geq R}J(x)e^{-S(x)}dx
 \leq Ce^{-\delta R^p}
 \quad(R\geq R_0).
 \label{eq:t10v22-coercive-tail}
\end{equation}
The same conclusion holds for $q=p$ if $c<a$.
\end{theorem}

\begin{proof}
If $q<p$, choose $R_0$ so that $cr^q\leq ar^p/2$ for $r\geq R_0$.  In the
case $q=p$, take any $0<\delta<a-c$.  Outside $R_0$, the integrand is bounded
by $C_Je^b e^{-\delta_0\|x\|^p}$ for some $\delta_0>0$.  Polar coordinates
show integrability.  For $r\geq R$ sufficiently large,
$r^{N-1}\leq C'e^{\delta_0r^p/2}$, so the radial tail is bounded by a constant
times $e^{-\delta_0R^p/2}$.  The compact inner region is finite by local
integrability.
\end{proof}

\begin{corollary}[Polynomial determinant owners]
\label{cor:t10v22-polynomial-density}
At fixed block dimension, any density owner satisfying
\begin{equation}
 J(x)\leq C(1+\|x\|)^M
 \label{eq:t10v22-polynomial-owner}
\end{equation}
is admissible in Theorem~\ref{thm:t10v22-coercive-tail} for every $p>0$.
This includes the absolute value of a finite matrix determinant whose entries
grow polynomially in the fields.
\end{corollary}

\begin{proof}
For every $q>0$ and sufficiently large $r$,
$M\log(1+r)\leq cr^q$.  Choose $q<p$ and apply the theorem.
\end{proof}

\begin{firewall}[Extensive determinants]
\label{fire:t10v22-extensive-determinant}
The exponent $M$ generally grows with the number of integrated fermionic or
ghost variables.  Corollary~\ref{cor:t10v22-polynomial-density} is uniform
only on a fixed rescaled block.  Across polymers, the extensive part must be
paid by the per-block action and the determinant-line phase remains a
separate V17 owner.
\end{firewall}

\subsection{Global stationary elimination}

Let $m\in M$ be the real physical variables and $s\in\mathbb R^k$ the
negative conformal variables.

\begin{proposition}[Global stationary graph and constrained density]
\label{prop:t10v22-global-stationary}
Suppose $V\in C^2(M\times\mathbb R^k)$ and
\begin{equation}
 -D^2_{ss}V(m,s)\succeq c_sI
 \quad\text{for all }(m,s),
 \label{eq:t10v22-global-concavity}
\end{equation}
with $c_s>0$.  Assume also, uniformly for $m$ in compact sets,
\begin{equation}
 \frac{D_sV(m,s)\cdot s}{\|s\|}\longrightarrow-\infty
 \quad\text{as }\|s\|\to\infty.
 \label{eq:t10v22-stationary-properness}
\end{equation}
Then every $m$ has a unique stationary point $s=\sigma(m)$, and $\sigma$ is
$C^1$.  Formally imposing the exact constraint $D_sV=0$ gives the physical
density
\begin{equation}
 d\mu_{\rm red}(m)
 =Z^{-1}J_0(m,\sigma(m))
 \det[-D^2_{ss}V(m,\sigma(m))]^{-1}
 e^{-V(m,\sigma(m))}dm.
 \label{eq:t10v22-constrained-density}
\end{equation}
\end{proposition}

\begin{proof}
For fixed $m$, strict concavity makes $D_sV$ strictly monotone with negative
Jacobian.  Condition \eqref{eq:t10v22-stationary-properness} implies that the
strictly concave function tends downward in every radial direction after a
finite radius, so it attains a unique global maximum, at which $D_sV=0$.
The implicit-function theorem and the uniform invertibility of $D^2_{ss}V$
give $C^1$ dependence.  Finally, the finite-dimensional delta/coarea identity
\[
 \int_{\mathbb R^k}\delta(D_sV(m,s))F(m,s)ds
 =\frac{F(m,\sigma(m))}
 {\det[-D^2_{ss}V(m,\sigma(m))]}
\]
proves \eqref{eq:t10v22-constrained-density}.
\end{proof}

\begin{corollary}[Reduced large-field criterion]
\label{cor:t10v22-reduced-tail}
If the reduced action
$V_{\rm red}(m)=V(m,\sigma(m))$ is coercive as in
\eqref{eq:t10v22-coercive-hypotheses} and the combined coarea, gauge,
fermion, ghost, and chart density in
\eqref{eq:t10v22-constrained-density} has subleading growth, then the reduced
physical measure has the exponential tail
\eqref{eq:t10v22-coercive-tail}.
\end{corollary}

\begin{proof}
Apply Theorem~\ref{thm:t10v22-coercive-tail} to
\eqref{eq:t10v22-constrained-density}.
\end{proof}

\subsection{Good, transition, and bad charts}

On one fixed rescaled block, let the compact gauge variables lie in $G^E$,
let $m$ denote the noncompact reduced physical variables, and choose a good
radius $R$.  Compactness gives a finite smooth gauge-chart cover.  Combine it
with a smooth radial partition in $m$ to obtain nonnegative functions
\begin{equation}
 1=\chi_{\rm good}+\sum_{a=1}^{N_{\rm tr}}\chi_{{\rm tr},a}
   +\chi_{\rm bad}.
 \label{eq:t10v22-chart-partition}
\end{equation}
The good function is supported where the preferred exponential/normal chart
and V21 small-activity condition hold; transition functions lie in a fixed
annulus or gauge-chart overlap; the bad function is supported at
$\|m\|\geq R$ or outside the chosen metric-condition window.

\begin{lemma}[Reflection-compatible partition]
\label{lem:t10v22-reflection-partition}
If reflection acts smoothly on the block configuration space, the partition
\eqref{eq:t10v22-chart-partition} can be chosen so that the good and bad
functions are invariant and transition charts occur in reflected pairs.
\end{lemma}

\begin{proof}
Start with a finite smooth partition subordinate to the chosen cover.  Add
each function to its reflected pullback and renormalize by the positive total
sum.  Pair labels exchanged by reflection; invariant charts may be averaged
with themselves.
\end{proof}

\begin{firewall}[A partition does not create reflected factorization]
\label{fire:t10v22-partition-OS}
Reflection covariance of the chart partition prevents an artificial symmetry
breaking.  Osterwalder--Schrader positivity still requires the complete
action and density on each chart to have the V10 positive-half factorization.
A generic partition of unity need not factor as a reflected square.
\end{firewall}

\subsection{Bad-component activity}

Let the coarse block adjacency graph have degree at most $D$.  A connected
set of $n$ blocks containing a fixed root can be encoded by a depth-first walk
of length at most $2(n-1)$, so its number is at most
\begin{equation}
 N_D(n)\leq D^{2(n-1)}.
 \label{eq:t10v22-connected-set-count}
\end{equation}

\begin{theorem}[Peierls bad-polymer sum]
\label{thm:t10v22-bad-polymer}
Suppose every connected bad set $X$ has absolute activity
\begin{equation}
 |K_{\rm bad}(X)|leq e^{-\tau|X|}
 \label{eq:t10v22-bad-activity}
\end{equation}
in its local field/density norm.  Then for a polymer size weight $e^{\kappa|X|}$,
\begin{equation}
 \sup_b\sum_{X\ni b}e^{\kappa|X|}|K_{\rm bad}(X)|
 \leq
 \frac{e^{-(\tau-\kappa)}}
 {1-D^2e^{-(\tau-\kappa)}}
 \label{eq:t10v22-bad-sum}
\end{equation}
provided $D^2e^{-(\tau-\kappa)}<1$.
\end{theorem}

\begin{proof}
Use \eqref{eq:t10v22-connected-set-count} and sum the resulting geometric
series in $n\geq1$.
\end{proof}

\begin{proposition}[Coercive origin of the Peierls cost]
\label{prop:t10v22-Peierls-origin}
Suppose a reduced block action on a connected bad set obeys
\begin{equation}
 S_X(m)\geq
 \sum_{b\in X}a\|m_b\|^p-c_0|X|
 \label{eq:t10v22-additive-coercivity}
\end{equation}
and every bad block has $\|m_b\|\geq R$.  If density owners cost at most
$e^{c_1|X|}$, then
\begin{equation}
 |K_{\rm bad}(X)|leq
 \exp[-(aR^p-c_0-c_1)|X|]
 \label{eq:t10v22-Peierls-cost}
\end{equation}
up to a fixed per-block radial integration constant.  Increasing $R$ absorbs
that constant and makes the hypothesis of
Theorem~\ref{thm:t10v22-bad-polymer} hold.
\end{proposition}

\begin{proof}
Every bad block contributes at least $aR^p$ to
\eqref{eq:t10v22-additive-coercivity}.  Exponentiation and the density bound
give the displayed factor.  The remaining normalized radial tails factor or
are bounded by Theorem~\ref{thm:t10v22-coercive-tail}; their fixed block
constant shifts $c_1$.
\end{proof}

\begin{firewall}[Interface energy must be in the lower bound]
\label{fire:t10v22-interface-energy}
Equation \eqref{eq:t10v22-additive-coercivity} is stronger than one-block
coercivity.  Negative interface terms can cancel bulk penalties and must be
bounded by a declared fraction of the neighbouring positive terms.  The
Peierls estimate is circular unless this many-block lower bound is proved for
the regulated action.
\end{firewall}

\subsection{The Euclidean conformal-factor test}

Use the Euclidean Einstein--Hilbert sign
\begin{equation}
 S_E(g)=-\frac1{16\pi G}\int_M\sqrt g\,R(g).
 \label{eq:t10v22-EH-action}
\end{equation}
For $g=e^{2\varphi}\bar g$ in four dimensions and compact support or boundary
conditions removing the integration-by-parts term,
\begin{equation}
 \int\sqrt g,R(g)
 =\int\sqrt{\bar g},e^{2\varphi}
 \left(\bar R+6|\bar\nabla\varphi|^2\right).
 \label{eq:t10v22-conformal-formula}
\end{equation}

\begin{proposition}[Real conformal integration is unbounded]
\label{prop:t10v22-conformal-unbounded}
On any chart containing a nonconstant smooth compactly supported function
$f$, the sequence $\varphi_t=tf$ makes the kinetic contribution to
$S_E(e^{2\varphi_t}\bar g)$ tend to $-\infty$ along a subsequence/direction.
Consequently the real Euclidean Einstein--Hilbert density does not satisfy
the coercive hypothesis of Theorem~\ref{thm:t10v22-coercive-tail}.
\end{proposition}

\begin{proof}
Equation \eqref{eq:t10v22-conformal-formula} puts
\[
 -\frac{6}{16\pi G}
 \int\sqrt{\bar g},e^{2tf}t^2|\bar\nabla f|^2
\]
in the action.  Choose $f$ so that it is positive on a set where
$|\bar\nabla f|>0$ (adding a harmless constant on its support if necessary).
The displayed negative term then has unbounded magnitude as $t\to+\infty$;
lower-derivative curvature terms cannot give a uniform positive coercive
bound against it.
\end{proof}

\begin{firewall}[Stationary reduction is a separate global theorem]
\label{fire:t10v22-conformal-firewall}
V9 proves a quadratic stationary conformal reduction and V10 a local
nonlinear implicit graph.  Proposition
\ref{prop:t10v22-conformal-unbounded} shows that neither result justifies a
real large-field conformal integral.  The programme needs a global constraint
solution satisfying Proposition~\ref{prop:t10v22-global-stationary}, a
specified complex contour with positivity reconstruction, or a stabilizing
physical term whose continuum meaning is proved.
\end{firewall}

\subsection{V22 global chart card}

\begin{definition}[Large-field completion card]
\label{def:t10v22-large-field-card}
One regulator block passes the card if:
\begin{enumerate}
\item the negative conformal variables are removed by an exact global
stationary/constraint theorem, not integrated on the real contour;
\item the resulting physical action satisfies the many-block coercive lower
bound \eqref{eq:t10v22-additive-coercivity};
\item all coarea, determinant, ghost, chart, and anomaly-magnitude densities
have subleading growth with declared extensive owners;
\item the good/transition/bad cover is finite, reflection compatible, and its
transition derivatives fit the V21 tree norm;
\item the bad cost passes Theorem~\ref{thm:t10v22-bad-polymer}; and
\item the good-chart activity passes the V21 forest card.
\end{enumerate}
\end{definition}

\begin{theorem}[Conditional complete finite-block integrability]
\label{thm:t10v22-finite-block-integrability}
If Definition~\ref{def:t10v22-large-field-card} holds, the complete reduced
finite-block density is normalizable.  Its good and transition connected
activities obey the V21 tree bound, while its bad connected activities are
summable in the V15 size-weighted polymer norm.
\end{theorem}

\begin{proof}
Theorem~\ref{thm:t10v22-coercive-tail} gives normalizability and tail decay
after stationary reduction.  The finite partition
\eqref{eq:t10v22-chart-partition} is exact.  Apply the V21 forest theorem on
good and transition charts and Theorem~\ref{thm:t10v22-bad-polymer} on bad
components.
\end{proof}

\begin{assessment}[Progress and bottleneck after V22]
\label{assess:t10v22}
V22 completes the abstract good/transition/bad large-field architecture and
proves the bad-polymer summability criterion.  It also locates the actual
Einstein obstruction: the real conformal factor defeats coercivity.  The
next step must therefore go upstream to the global Einstein constraint rather
than refine the polymer estimate.  A natural controlled setting is the
constant-mean-curvature conformal method, where the conformal factor solves a
Lichnerowicz equation with sign conditions that can be tested.
\end{assessment}

\begin{programme}[Tenth-series V23: global conformal constraint]
\label{prog:t10v22-V23}
V23 will formulate a compact-block Lichnerowicz equation with regulated
boundary data and nonnegative matter energy, prove existence and uniqueness
under explicit Yamabe/mean-curvature sign hypotheses by sub/supersolutions and
the maximum principle, and compute the linearized inverse needed for the V10
Schur and V22 coarea cards.  Cases outside those signs will remain separate
rather than being folded into the theorem.
\end{programme}

\begin{firewall}[Claim boundary after tenth-series V22]
\label{fire:t10v22-boundary}
The large-field theorem is conditional, and the conformal calculation proves
an obstruction rather than a construction.  This version does not solve the
global constraint, establish reflected factorization of the chart cover,
control the chiral phase, prove the ultraviolet fixed point, or construct the
full SM--Einstein continuum.
\end{firewall}

\paragraph{Parent-version record.}
V22 was formed from
\texttt{Renewal\_SM\_Einstein\_Continuum\_Research\_Notes\_V21.tex}.
The V21 parent had $8\,535$ validator-counted source lines, $323\,846$ bytes,
and SHA--256
\[
 \texttt{B228D03482D258F6965FC2F220C95270E499722635620A81DC6D09FC6AB8DE35}.
\]

% =============================================================================
% END APPENDED TENTH-SERIES V22 MATERIAL
% =============================================================================

% =============================================================================
% APPENDED TENTH-SERIES V23 MATERIAL
% =============================================================================

\section{Tenth-series V23: a global positive conformal constraint in a CMC
block}
\label{sec:v10v23}

V22 showed that the Euclidean conformal factor cannot be controlled by a real
large-field integral.  This version replaces that integration by a global
constraint solve in a restricted but rigorous setting: a compact
three-dimensional constant-mean-curvature block with positive Dirichlet
conformal data.  The theorem gives existence, uniqueness, a uniform
linearized inverse, and the exact response terms that must be charged when
the constraint is not the stationarity equation of the same action.

\subsection{The sign-controlled Lichnerowicz problem}

Let $(\Omega,\bar g)$ be a smooth compact connected three-dimensional
Riemannian domain with boundary.  Consider
\begin{equation}
 \mathcal F(\phi):=
 -8\bar\Delta\phi+\bar R\phi+b\phi^5
 -a\phi^{-7}-c\phi^{-3}=0,
 \qquad
 \phi|_{\partial\Omega}=\phi_\partial>0.
 \label{eq:t10v23-Lichnerowicz}
\end{equation}
Here $a$ is the squared transverse-traceless/momentum contribution, $b$ is
the positive CMC coefficient, and $c$ is a nonnegative conformally rescaled
matter energy.  The theorem uses the explicit sign chamber
\begin{equation}
 \bar R\geq0,
 \quad b\geq b_0>0,
 \quad a,c\geq0,
 \quad a+c\geq s_0>0.
 \label{eq:t10v23-sign-chamber}
\end{equation}

\begin{firewall}[The sign chamber is a restriction]
\label{fire:t10v23-sign-chamber}
For the physical CMC equations, $b$ and the power of the matter term depend
on dimension, cosmological constant, mean curvature, and the chosen conformal
scaling of matter.  Nonnegative physical energy alone does not imply every
condition in \eqref{eq:t10v23-sign-chamber}.  V23 proves this chamber and does
not merge the zero-mean-curvature, negative-Yamabe, vanishing-source, or
non-CMC cases into it.
\end{firewall}

\begin{lemma}[Constant barriers]
\label{lem:t10v23-barriers}
Assume the coefficients are bounded and
\eqref{eq:t10v23-sign-chamber} holds.  There are constants
$0<\phi_-\leq\phi_+$ with
\begin{equation}
 \phi_-\leq\phi_\partial\leq\phi_+,
 \qquad
 \mathcal F(\phi_-)\leq0,
 \qquad
 \mathcal F(\phi_+)\geq0.
 \label{eq:t10v23-barriers}
\end{equation}
\end{lemma}

\begin{proof}
Choose $\phi_-\leq1$ below the boundary minimum.  Since
$a+c\geq s_0$ and $\phi_-\leq1$,
\[
 a\phi_-^{-7}+c\phi_-^{-3}
 \geq(a+c)\phi_-^{-3}\geq s_0\phi_-^{-3}.
\]
The positive terms are at most
$(\|\bar R\|_\infty+\|b\|_\infty)\phi_-$, so a sufficiently small
$\phi_-$ is a subsolution.  Choose $\phi_+\geq1$ above the boundary maximum.
Then
\[
 b\phi_+^5-a\phi_+^{-7}-c\phi_+^{-3}
 \geq b_0\phi_+^5-|a+c\|_\infty,
\]
which is nonnegative for large $\phi_+$.  The scalar-curvature term is also
nonnegative.
\end{proof}

\begin{theorem}[Existence and uniqueness]
\label{thm:t10v23-existence-uniqueness}
Under \eqref{eq:t10v23-sign-chamber} and smooth positive boundary data,
problem \eqref{eq:t10v23-Lichnerowicz} has a unique smooth solution satisfying
\begin{equation}
 \phi_-\leq\phi\leq\phi_+.
 \label{eq:t10v23-uniform-positive-bounds}
\end{equation}
\end{theorem}

\begin{proof}
On the compact interval $[\phi_-,\phi_+]$, choose $\lambda$ so large that
\[
 u\longmapsto
 \bar Ru+bu^5-au^{-7}-cu^{-3}+\lambda u
\]
is increasing at every point.  Starting from the sub- and supersolutions,
solve the linear Dirichlet equation
$(-8\bar\Delta+\lambda)u_{n+1}$ equal to the corresponding monotone right
side.  The maximum principle keeps the lower sequence increasing, the upper
sequence decreasing, and both inside the barrier interval.  Elliptic
estimates and monotone convergence give a solution; bootstrapping gives
smoothness.

For uniqueness, let $\phi_1,\phi_2$ be solutions with the same boundary data.
The difference $w$ satisfies
\[
 -8\bar\Delta w+q(x)w=0,
 \qquad w|_{\partial\Omega}=0,
\]
where, by the mean-value theorem,
\[
 q=\bar R+5b\xi^4+7a\xi^{-8}+3c\xi^{-4}>0
\]
for a pointwise intermediate $\xi$ between the two positive solutions.
The maximum principle gives $w=0$.
\end{proof}

\subsection{The linearized constraint inverse}

At the solution, define
\begin{equation}
 \mathcal L_\phi
 :=D_\phi\mathcal F
 =-8\bar\Delta+\bar R+5b\phi^4
   +7a\phi^{-8}+3c\phi^{-4}
 \label{eq:t10v23-linearized-operator}
\end{equation}
with homogeneous Dirichlet boundary condition.

\begin{proposition}[Uniform positive linearization]
\label{prop:t10v23-linearized-gap}
The operator \eqref{eq:t10v23-linearized-operator} is self-adjoint and
\begin{equation}
 \mathcal L_\phi\succeq q_0I,
 \qquad
 q_0:=5b_0\phi_-^4>0.
 \label{eq:t10v23-linearized-gap}
\end{equation}
Consequently
\begin{equation}
 \|\mathcal L_\phi^{-1}\|_{L^2\to L^2}\leq q_0^{-1}.
 \label{eq:t10v23-linearized-inverse}
\end{equation}
\end{proposition}

\begin{proof}
For $u\in H^1_0(\Omega)$,
\[
 \langle u,\mathcal L_\phi u\rangle
 =8\|\bar\nabla u\|_2^2
 +\int(\bar R+5b\phi^4+7a\phi^{-8}+3c\phi^{-4})u^2.
\]
Every term is nonnegative, and
$5b\phi^4\geq5b_0\phi_-^4$.  The spectral theorem gives the inverse bound.
\end{proof}

\begin{corollary}[First response to coefficient data]
\label{cor:t10v23-first-response}
For variations with fixed Dirichlet boundary value,
\begin{equation}
 \mathcal L_\phi\,\delta\phi
 =-(\delta\bar R)\phi-(\delta b)\phi^5
   +(\delta a)\phi^{-7}+(\delta c)\phi^{-3},
 \label{eq:t10v23-first-response}
\end{equation}
and hence
\begin{equation}
 \|\delta\phi\|_2
 \leq q_0^{-1}
 \|-(\delta\bar R)\phi-(\delta b)\phi^5
   +(\delta a)\phi^{-7}+(\delta c)\phi^{-3}\|_2.
 \label{eq:t10v23-first-response-bound}
\end{equation}
\end{corollary}

\begin{proof}
Differentiate \eqref{eq:t10v23-Lichnerowicz} and use
Proposition~\ref{prop:t10v23-linearized-gap}.
\end{proof}

\begin{proposition}[Analytic parameter dependence]
\label{prop:t10v23-analytic-dependence}
Let the coefficients and a finite regulator discretization depend analytically
on parameters $z$ in a complex polydisc.  If the positive solution remains in
a complex neighbourhood on which
$\|\mathcal L_\phi^{-1}\|\leq2q_0^{-1}$, then $\phi(z)$ is analytic on a
smaller polydisc and its parameter derivatives satisfy Cauchy factorial
bounds.  Thus the conformal response fits the V19 Gevrey-one parameter
ledger.
\end{proposition}

\begin{proof}
The analytic implicit-function theorem applies because
$D_\phi\mathcal F=\mathcal L_\phi$ is invertible.  Cauchy's formula on every
strictly smaller polydisc gives the factorial derivative estimates.
\end{proof}

\subsection{Finite-regulator coarea factor}

At a finite lattice/spectral regulator, let the discretized positive
conformal variables be $\phi\in(0,\infty)^K$.  The constraint map
$\mathcal F:\mathbb R^K\to\mathbb R^K$ has Jacobian
$\mathcal L_\phi$ at its unique zero.

\begin{proposition}[Constraint coarea density]
\label{prop:t10v23-coarea}
For every continuous compactly supported $G$ in the positive chart,
\begin{equation}
 \int_{(0,\infty)^K}
 \delta(\mathcal F(\phi))G(\phi)d\phi
 =\frac{G(\phi_*)}{\det\mathcal L_{\phi_*}},
 \label{eq:t10v23-coarea}
\end{equation}
and
\begin{equation}
 0<\det(\mathcal L_{\phi_*})^{-1}\leq q_0^{-K}.
 \label{eq:t10v23-coarea-bound}
\end{equation}
\end{proposition}

\begin{proof}
The solution is unique and the Jacobian is positive invertible by
Proposition~\ref{prop:t10v23-linearized-gap}.  The finite-dimensional delta
change-of-variables formula gives \eqref{eq:t10v23-coarea}; multiplying its
$K$ eigenvalue lower bounds gives \eqref{eq:t10v23-coarea-bound}.
\end{proof}

\begin{firewall}[Continuum determinant renormalization]
\label{fire:t10v23-continuum-determinant}
The power $q_0^{-K}$ is extensive and diverges as the regulator dimension
$K$ grows.  V23 supplies positivity and local derivative control of the
coarea determinant.  Its vacuum/local heat-kernel divergences must still be
assigned to the V17 graded counterterm tower before a continuum limit is
claimed.
\end{firewall}

\subsection{Constraint response versus stationary Schur positivity}

Let $m$ denote all other physical block variables and suppose the constraint
is $F(m,\phi)=0$ with invertible $F_\phi$.  Then
\begin{align}
 D\phi&=-F_\phi^{-1}F_m,
 \label{eq:t10v23-general-first-response}\\
 D^2\phi[u,v]&=-F_\phi^{-1}\bigl(
 F_{mm}[u,v]+F_{m\phi}[u,D\phi v]
 +F_{m\phi}[v,D\phi u]
 +F_{\phi\phi}[D\phi u,D\phi v]
 \bigr).
 \label{eq:t10v23-general-second-response}
\end{align}

\begin{proposition}[Exact reduced-action Hessian]
\label{prop:t10v23-reduced-Hessian}
For an action $S(m,\phi)$ restricted to the constraint graph,
$S_{\rm red}(m)=S(m,\phi(m))$, one has
\begin{align}
 D^2S_{\rm red}[u,v]
 ={}&S_{mm}[u,v]
 +S_{m\phi}[u,D\phi v]+S_{m\phi}[v,D\phi u]\\
 &+S_{\phi\phi}[D\phi u,D\phi v]
 +S_\phi[D^2\phi[u,v]].
 \label{eq:t10v23-reduced-Hessian}
\end{align}
If $F=S_\phi$, the last term vanishes on the graph and the formula reduces to
the stationary Schur identity.  For a distinct Hamiltonian constraint, all
displayed response terms must be estimated.
\end{proposition}

\begin{proof}
Differentiate $S(m,\phi(m))$ twice and apply the chain rule.  If
$F=S_\phi=0$, use $D\phi=-S_{\phi\phi}^{-1}S_{\phi m}$ to obtain the Schur
formula.
\end{proof}

\begin{corollary}[Explicit mismatch debit]
\label{cor:t10v23-mismatch-debit}
If $\|F_\phi^{-1}\|\leq q_0^{-1}$ and the first and second derivatives of
$F,S$ are bounded on the barrier interval, equations
\eqref{eq:t10v23-general-first-response}--
\eqref{eq:t10v23-reduced-Hessian} give a finite computable norm
$\Delta_{\rm CMC}$ for
\begin{equation}
 \|D^2S_{\rm red}-S_{mm}\|\leq\Delta_{\rm CMC}.
 \label{eq:t10v23-CMC-debit}
\end{equation}
The V10/V14 physical precision gate must subtract
$\Delta_{\rm CMC}$ unless the constraint is proved to be the same stationary
equation used there.
\end{corollary}

\begin{proof}
Bound $D\phi$ by $q_0^{-1}\|F_m\|$, insert it in
\eqref{eq:t10v23-general-second-response}, and then bound every term of
\eqref{eq:t10v23-reduced-Hessian}.  All constants are finite because
$\phi\in[\phi_-,\phi_+]$.
\end{proof}

\subsection{CMC conformal block card}

\begin{definition}[V23 constraint card]
\label{def:t10v23-constraint-card}
A spatial regulator block passes the card if:
\begin{enumerate}
\item its conformal data lie in the explicit sign chamber
\eqref{eq:t10v23-sign-chamber} with uniform barriers;
\item the Dirichlet/shared-face conformal data are positive, glued
associatively, and have the required reflection parity;
\item the momentum constraint has already supplied the transverse-traceless
and matter-current data entering $a$;
\item the linearized gap $q_0$ is uniform in the rescaled block family;
\item the coarea determinant is assigned to the V17 local counterterm/density
ledger; and
\item the response debit $\Delta_{\rm CMC}$ fits inside the V14 physical
precision and contraction margins.
\end{enumerate}
\end{definition}

\begin{theorem}[Global conformal elimination in the sign chamber]
\label{thm:t10v23-global-elimination}
If Definition~\ref{def:t10v23-constraint-card} holds, the conformal variable
is globally and uniquely eliminated on each block, its constrained coarea
density is positive and finite at fixed regulator, and its first and second
responses are uniformly bounded.  Consequently the conformal direction does
not enter the real large-field integration domain of the V22 card.
\end{theorem}

\begin{proof}
Existence and uniqueness are Theorem
\ref{thm:t10v23-existence-uniqueness}; the response is controlled by
Proposition~\ref{prop:t10v23-linearized-gap} and
Corollary~\ref{cor:t10v23-mismatch-debit}; the positive coarea density is
Proposition~\ref{prop:t10v23-coarea}.
\end{proof}

\begin{firewall}[What CMC does not solve]
\label{fire:t10v23-CMC-scope}
The CMC hypothesis decouples part of the constraint system.  V23 does not
solve the non-CMC coupled vector--scalar constraints, prove that RG reblocking
preserves the sign chamber, identify the Hamiltonian constraint with the
Euclidean stationary equation, or establish time-reflection factorization of
the constrained density.
\end{firewall}

\begin{assessment}[Progress and bottleneck after V23]
\label{assess:t10v23}
V23 gives the first global nonperturbative elimination of the conformal
variable in the programme, within a clear CMC/Yamabe-positive sign chamber.
It supplies a uniformly positive linearized inverse and an exact mismatch
debit for the physical Hessian.  The next upstream equation is the momentum
constraint that produces the transverse-traceless data and must share the
same boundary, reflection, and approximate-Killing ledger developed in
V12--V15.
\end{assessment}

\begin{programme}[Tenth-series V24: momentum-constraint solve]
\label{prog:t10v23-V24}
V24 will solve the conformal vector equation on the complement of conformal
Killing fields using a Korn-type coercivity theorem, identify its force/torque
compatibility data, and bound its contribution to $a$.  The resulting vector
response will then be composed with the V23 scalar response to produce one
CMC Einstein-constraint block map.
\end{programme}

\begin{firewall}[Claim boundary after tenth-series V23]
\label{fire:t10v23-boundary}
The theorem covers only the stated sign chamber and finite regulated coarea
factor.  It does not solve the momentum constraint, prove preservation of the
chamber under RG, close the graded ultraviolet flow, control the chiral
determinant line, or construct the full SM--Einstein continuum.
\end{firewall}

\paragraph{Parent-version record.}
V23 was formed from
\texttt{Renewal\_SM\_Einstein\_Continuum\_Research\_Notes\_V22.tex}.
The V22 parent had $8\,909$ validator-counted source lines, $338\,790$ bytes,
and SHA--256
\[
 \texttt{2FF58151B578133719FEC40C87E626A15E4782701561545C2D0A0CD2AEC5CF05}.
\]

% =============================================================================
% END APPENDED TENTH-SERIES V23 MATERIAL
% =============================================================================

% =============================================================================
% APPENDED TENTH-SERIES V24 MATERIAL
% =============================================================================

\section{Tenth-series V24: the CMC momentum constraint and conformal-Killing
owners}
\label{sec:v10v24}

V23 solved the scalar CMC constraint assuming the transverse-traceless and
matter-current contribution $a$ was already known.  This version constructs
that datum from the vector momentum constraint.  The relevant cokernel is
larger than the rigid-motion space of V12: the trace-free symmetric gradient
annihilates conformal Killing fields, so dilation and special-conformal modes
must also be retained near a flat block.

\subsection{The conformal vector operator}

On a three-dimensional Riemannian block $(\Omega,\bar g)$ define
\begin{equation}
 (\mathbb LW)_{ij}
 :=\bar\nabla_iW_j+\bar\nabla_jW_i
 -\frac23(\bar\nabla_kW^k)\bar g_{ij}.
 \label{eq:t10v24-conformal-Killing-operator}
\end{equation}
Its kernel $\mathcal{CK}_{\bar g}$ is the finite-dimensional space of
conformal Killing fields.  Define the nonnegative vector operator
\begin{equation}
 \mathcal P_{\bar g}W:=-\operatorname{div}_{\bar g}\mathbb LW.
 \label{eq:t10v24-vector-Laplacian}
\end{equation}

\begin{lemma}[Conformal Korn estimate]
\label{lem:t10v24-conformal-Korn}
On a smooth compact connected block there is $C_{\rm CK}<\infty$ such that
\begin{equation}
 \inf_{X\in\mathcal{CK}_{\bar g}}
 \|W-X\|_{H^1}
 \leq C_{\rm CK}\|\mathbb LW\|_{L^2}.
 \label{eq:t10v24-conformal-Korn}
\end{equation}
\end{lemma}

\begin{proof}
The principal symbol of $\mathbb L$ is injective modulo its finite-dimensional
kernel, so the standard first-order elliptic estimate has the form
\[
 \|W\|_{H^1}\leq C(\|\mathbb LW\|_2+\|W\|_2).
\]
If \eqref{eq:t10v24-conformal-Korn} failed, choose unit $H^1$ representatives
orthogonal to $\mathcal{CK}_{\bar g}$ with $\mathbb LW_n\to0$.  Rellich
compactness and the elliptic estimate yield a strongly convergent $H^1$
subsequence.  Its limit lies in $\mathcal{CK}_{\bar g}$ and in its orthogonal
complement, contradicting unit norm.
\end{proof}

\begin{theorem}[Momentum-constraint solve]
\label{thm:t10v24-momentum-solve}
Let $J\in H^{-1}(T\Omega)$ obey
\begin{equation}
 \langle J,X\rangle=0
 \quad(X\in\mathcal{CK}_{\bar g}).
 \label{eq:t10v24-CK-compatibility}
\end{equation}
There is a unique class
$W\in H^1(T\Omega)/\mathcal{CK}_{\bar g}$ satisfying
\begin{equation}
 \int_\Omega\langle\mathbb LW,\mathbb LV\rangle,dv_{\bar g}
 =\langle J,V\rangle
 \quad(V\in H^1(T\Omega)),
 \label{eq:t10v24-vector-variational}
\end{equation}
and
\begin{equation}
 \|\mathbb LW\|_2
 \leq C_{\rm CK}\|J\|_{(H^1/\mathcal{CK})^*}.
 \label{eq:t10v24-vector-bound}
\end{equation}
It solves $\mathcal P_{\bar g}W=J$ with zero weak conformal traction
$(\mathbb LW)n=0$.  Compatibility is also necessary.
\end{theorem}

\begin{proof}
Use \eqref{eq:t10v24-conformal-Korn} to make
$\|\mathbb LW\|_2$ a coercive quotient norm.  Compatibility makes the right
side of \eqref{eq:t10v24-vector-variational} well defined.  Lax--Milgram gives
the unique quotient class and the estimate.  Integration by parts identifies
the equation and natural traction.  Pairing the equation with a conformal
Killing field proves necessity.
\end{proof}

\begin{firewall}[The cokernel is larger than force and torque]
\label{fire:t10v24-CK-cokernel}
On a flat three-dimensional domain, conformal Killing fields include
translations, rotations, dilation, and special-conformal generators.  The
V12 force/torque ledger covers only the first two classes.  Natural-boundary
momentum solves require the complete projection on
$\mathcal{CK}_{\bar g}^*$, or boundary data whose flux supplies those missing
moments.
\end{firewall}

\subsection{Uniform near-flat complement}

Let $\mathcal{CK}_0$ be the reference flat conformal-Killing space and
$H_0$ its fixed $H^1$ orthogonal complement.  Assume
\begin{equation}
 \|(\mathbb L_{\bar g}-\mathbb L_0)W\|_2
 \leq\delta_{\rm CK}\|W\|_{H^1}.
 \label{eq:t10v24-CK-perturbation}
\end{equation}

\begin{proposition}[Complement solve with finite residual]
\label{prop:t10v24-complement-solve}
If $C_{{\rm CK},0}\delta_{\rm CK}<1$, the bilinear form in
\eqref{eq:t10v24-vector-variational} is uniformly coercive on $H_0$.  For
every $J$ it produces $W\in H_0$ and a residual
\begin{equation}
 q_{\rm CK}:=\mathcal P_{\bar g}W-J\in\mathcal{CK}_0^*,
 \label{eq:t10v24-CK-residual}
\end{equation}
with
\begin{equation}
 \|\mathbb L_{\bar g}W\|_2
 \leq\frac{C}{1-C_{{\rm CK},0}\delta_{\rm CK}}
 \|J\|_{H_0^*}.
 \label{eq:t10v24-complement-bound}
\end{equation}
If $J$ annihilates $\mathcal{CK}_0$, the residual coordinates are
$O(\delta_{\rm CK})$ times the displayed solution norm, up to the metric
inner-product perturbation.
\end{proposition}

\begin{proof}
The same triangle argument as Lemma
\ref{lem:t10v13-complement-coercivity} gives
\[
 \|W\|_{H^1}
 \leq\frac{C_{{\rm CK},0}}
 {1-C_{{\rm CK},0}\delta_{\rm CK}}
 \|\mathbb L_{\bar g}W\|_2
 \quad(W\in H_0).
\]
Lax--Milgram on $H_0$ gives the solution and bound.  Its equation annihilates
$H_0$, so the residual lies in the finite complementary dual.  Pairing with
$X\in\mathcal{CK}_0$ introduces
$\mathbb L_{\bar g}X=(\mathbb L_{\bar g}-\mathbb L_0)X$, which gives the
$O(\delta_{\rm CK})$ estimate.
\end{proof}

\begin{definition}[Extended conformal-charge ledger]
\label{def:t10v24-CK-ledger}
The retained affine data of V15 are enlarged by the coordinates of
$q_{\rm CK}\in\mathcal{CK}_0^*$.  Translation/rotation components use the
V15 force/torque cocycle.  Dilation and special-conformal components use the
dual action induced by owner reanchoring on the fixed polynomial basis of
$\mathcal{CK}_0$.  No component is inverted in the stable norm.
\end{definition}

\subsection{Regularity and construction of the scalar coefficient}

Assume $\partial\Omega$, $\bar g$, and $J$ have $C^{1,\alpha}$ regularity and
that a uniform complement Schauder estimate holds.  Then
\begin{equation}
 \|W\|_{C^{2,\alpha}/\mathcal{CK}}
 \leq C_{\rm Sch}\|J\|_{C^{0,\alpha}}.
 \label{eq:t10v24-Schauder}
\end{equation}
Let $\sigma$ be prescribed transverse-traceless data and set
\begin{equation}
 \Sigma:=\sigma+\mathbb LW,
 \qquad
 a:=\frac18|\Sigma|_{\bar g}^2.
 \label{eq:t10v24-a-definition}
\end{equation}

\begin{proposition}[Coefficient bound]
\label{prop:t10v24-a-bound}
Under \eqref{eq:t10v24-Schauder},
\begin{equation}
 \|a\|_{C^{0,\alpha}}
 \leq C_a
 \left(\|\sigma\|_{C^{0,\alpha}}
       +C_{\rm Sch}\|J\|_{C^{0,\alpha}}\right)^2.
 \label{eq:t10v24-a-bound}
\end{equation}
\end{proposition}

\begin{proof}
The first derivative of $W$ is bounded by its $C^{2,\alpha}$ norm.  Apply the
H\"older product estimate to \eqref{eq:t10v24-a-definition}.
\end{proof}

\begin{firewall}[Upper bounds do not preserve the V23 lower source]
\label{fire:t10v24-source-lower-bound}
Proposition~\ref{prop:t10v24-a-bound} controls the upper size of $a$.  The V23
barrier chamber also requires $a+c\geq s_0>0$.  This pointwise lower bound
must come from nonvanishing TT data, matter energy, or a different barrier
argument; it does not follow from the momentum solve.
\end{firewall}

\subsection{Response estimates}

Let the data depend smoothly on a parameter.  On the fixed complement,
differentiation of $\mathcal P_{\bar g}W=J$ gives
\begin{equation}
 \mathcal P_{\bar g}\,\delta W
 =\delta J-(\delta\mathcal P_{\bar g})W+\delta q_{\rm CK},
 \label{eq:t10v24-vector-response}
\end{equation}
where the retained residual enforces the complement projection.

\begin{proposition}[Vector and $a$ response]
\label{prop:t10v24-vector-response}
If the complement inverse norm is at most $\gamma_V^{-1}$, then
\begin{equation}
 \|\delta W\|_{H^1/\mathcal{CK}_0}
 \leq\gamma_V^{-1}
 \|\delta J-(\delta\mathcal P)W+\delta q_{\rm CK}\|_{H_0^*}.
 \label{eq:t10v24-vector-response-bound}
\end{equation}
Moreover
\begin{equation}
 \delta a
 =\frac14\left\langle\Sigma,
 \delta\sigma+\mathbb L\delta W+(\delta\mathbb L)W
 \right\rangle
 +\frac18(\delta\bar g^{-1})*(\Sigma\otimes\Sigma),
 \label{eq:t10v24-a-response}
\end{equation}
and hence $\|\delta a\|$ is bounded by a polynomial in the data norms times
$\gamma_V^{-1}$.
\end{proposition}

\begin{proof}
Invert \eqref{eq:t10v24-vector-response} on the complement.  Differentiate
\eqref{eq:t10v24-a-definition}; the first term is the derivative of the
square and the last records variation of the metric contraction.  Apply
Cauchy--Schwarz and the product estimates.
\end{proof}

Second differentiation gives the same structure with at most two complement
inverses and bounded first/second coefficient variations.  Denote the
resulting explicit polynomial bound by
\begin{equation}
 \|D^2a\|\leq\Delta_a
 (\gamma_V^{-1},\|\bar g\|_{C^{2,\alpha}},
 \|\sigma\|_{C^{1,\alpha}},\|J\|_{C^{1,\alpha}}).
 \label{eq:t10v24-second-a-stock}
\end{equation}

\subsection{The composed CMC constraint map}

\begin{theorem}[Vector--scalar CMC block map]
\label{thm:t10v24-composed-CMC}
Assume:
\begin{enumerate}
\item the momentum data pass Theorem~\ref{thm:t10v24-momentum-solve}, or the
near-flat complement residual in Proposition
\ref{prop:t10v24-complement-solve} is closed by the retained ledger;
\item the resulting $a$ together with $c,b,\bar R$ lies in the V23 sign
chamber; and
\item both complement gaps $\gamma_V$ and $q_0$ are uniform.
\end{enumerate}
Then the CMC constraint system has a unique pair
\begin{equation}
 (W\bmod\mathcal{CK},\phi>0),
 \label{eq:t10v24-CMC-pair}
\end{equation}
and its first two responses to the freely specified data are bounded by an
explicit polynomial in $\gamma_V^{-1},q_0^{-1}$ and the data norms.
\end{theorem}

\begin{proof}
The CMC vector equation is independent of $\phi$, so solve it first by
Theorem~\ref{thm:t10v24-momentum-solve}.  Form $a$ by
\eqref{eq:t10v24-a-definition} and apply Theorem
\ref{thm:t10v23-existence-uniqueness}.  Vector responses are controlled by
Proposition~\ref{prop:t10v24-vector-response}; insert
\eqref{eq:t10v24-a-response} and
\eqref{eq:t10v24-second-a-stock} into the V23 first and second scalar response
formulas.
\end{proof}

\begin{proposition}[Finite-regulator vector coarea factor]
\label{prop:t10v24-vector-coarea}
On the finite complement discretization, imposing the momentum constraint by
a delta density contributes
\begin{equation}
 \det(\mathcal P_{\bar g}|_{H_0})^{-1}>0,
 \qquad
 \det(\mathcal P_{\bar g}|_{H_0})^{-1}leq\gamma_V^{-K_V},
 \label{eq:t10v24-vector-coarea}
\end{equation}
where $K_V$ is the complement dimension.  Conformal-Killing coordinates are
not included in this determinant.
\end{proposition}

\begin{proof}
The complement operator is positive with all eigenvalues at least
$\gamma_V$.  Apply the finite-dimensional delta change-of-variables formula
and multiply the eigenvalue bounds.
\end{proof}

\begin{firewall}[Coarea and Faddeev--Popov determinants are distinct]
\label{fire:t10v24-determinant-types}
The determinant in \eqref{eq:t10v24-vector-coarea} is the Jacobian of the
momentum constraint on its complement.  It is not the gauge-fixing
Faddeev--Popov determinant and not the chiral fermion determinant.  All three
have different zero modes, phases, and counterterm owners.
\end{firewall}

\subsection{Reflection and boundary data}

\begin{proposition}[Equivariance of the CMC solve]
\label{prop:t10v24-CMC-reflection}
If an isometric reflection preserves $\Omega$, the reference complement,
boundary data, and the parities of $(\sigma,J,c)$, then the unique complement
solution $W$, the coefficient $a$, and the positive scalar solution $\phi$
have their induced reflection parities.
\end{proposition}

\begin{proof}
Pullback by the isometry commutes with $\mathbb L$, $\mathcal P$, and the
scalar Lichnerowicz operator.  The transformed fields solve the same
complement vector and scalar boundary problems.  Uniqueness gives the claim.
\end{proof}

\begin{firewall}[Equivariance is not OS positivity]
\label{fire:t10v24-equivariance-not-OS}
Reflection equivariance of a classical constraint solution does not prove
that the delta/coarea constrained density has a positive-half square
factorization.  The boundary values and both coarea determinants couple the
two halves and require a separate gluing theorem.
\end{firewall}

\begin{definition}[V24 CMC constraint card]
\label{def:t10v24-CMC-card}
A block passes the card when the vector compatibility/residual ledger closes,
the conformal-Killing and scalar gaps are uniform, the V23 sign chamber is
preserved, all response stocks fit the V14/V23 precision debit, and the two
coarea determinants are assigned to the V17 graded density complex.
\end{definition}

\begin{assessment}[Progress and bottleneck after V24]
\label{assess:t10v24}
V24 completes the CMC vector--scalar constraint solve on one block and finds
the additional conformal-Killing owner that a flat-limit construction must
retain.  The remaining immediate issue is gluing: solving each block with
Dirichlet or natural boundary data does not automatically reproduce the
solution on their union, and the constraint determinants need a Schur/gluing
identity compatible with reflection.
\end{assessment}

\begin{programme}[Tenth-series V25: audit and constraint gluing]
\label{prog:t10v24-V25}
V25 will perform the compulsory YM/NS audit and then derive the finite
Dirichlet-to-Neumann Schur complement for the vector and scalar linearized
constraint operators.  It will state the exact interface equation for block
solutions and factor the determinant into interior and interface owners.
After the audit, V26 will use that identity to test reflected factorization.
\end{programme}

\begin{firewall}[Claim boundary after tenth-series V24]
\label{fire:t10v24-boundary}
This version treats only CMC data in the V23 sign chamber.  It does not prove
constraint gluing, determinant reflection positivity, preservation of the
chamber under RG, the ultraviolet fixed point, or the full SM--Einstein
continuum.
\end{firewall}

\paragraph{Parent-version record.}
V24 was formed from
\texttt{Renewal\_SM\_Einstein\_Continuum\_Research\_Notes\_V23.tex}.
The V23 parent had $9\,297$ validator-counted source lines, $352\,860$ bytes,
and SHA--256
\[
 \texttt{F5995EAF75FA1121F183C47102E627950CF898B5DEC613AEB6317F39B7B28E8B}.
\]

% =============================================================================
% END APPENDED TENTH-SERIES V24 MATERIAL
% =============================================================================

% =============================================================================
% APPENDED TENTH-SERIES V25 MATERIAL
% =============================================================================

\section{Tenth-series V25: companion audit and exact constraint gluing}
\label{sec:v10v25}

\subsection{Compulsory companion audit}

The newest companion filenames at this checkpoint were
\texttt{Renewal\_Yang\_Mills\_Mass\_Gap\_Research\_Notes\_V38.tex} and
\texttt{Renewal\_NS\_Stability\_Research\_Notes\_V26.tex}.  Their records
were
\[
\begin{array}{c|r|l}
\text{file}&\text{bytes}&\text{SHA--256}\ \hline
\texttt{YM V38}&595\,887&
\texttt{ED66F96B4973F9101591382377418FE1F104F83A5E79AE50A0BBA577FFC94D7D}
\\
\texttt{NS V26}&278\,283&
\texttt{82C887F8C2C69E4F28FAD7C3DC8DE5B9099CB5A4BF431EBA1FFF3E91935978AD}.
\end{array}
\]
The YM V37 and V38 files are byte-identical, so the newest mathematical
appendix is V37.  YM V36 proves exact retained-to-stable decoupling in an
aligned pullback chart, while V37 computes its continuum background-field
one-loop coefficient but leaves exact lattice matching conditional.  The NS
file is unchanged.

\begin{assessment}[V25 audit transfer]
\label{assess:t10v25-audit-transfer}
The exact-pullback lesson applies to the constraint interface: once the shared
trace is a prescribed retained coordinate, solving the zero-trace interiors
does not create a derivative from that retained coordinate back into the
aligned interior stable variable.  The physical flux response remains in the
interface equation.  No YM determinant coefficient or beta-function constant
is imported into the Einstein--SM constraint determinants.
\end{assessment}

\subsection{Finite positive interface algebra}

Let a finite regulated positive operator on two interiors $I_+,I_-$ and one
shared trace space $B$ have block matrix
\begin{equation}
 L=
 \begin{pmatrix}
 A_+&0&K_+\\
 0&A_-&K_-\\
 K_+^*&K_-^*&D
 \end{pmatrix},
 \qquad A_\pm>0.
 \label{eq:t10v25-block-matrix}
\end{equation}
Define the interface Schur operator
\begin{equation}
 \mathcal S_B
 :=D-K_+^*A_+^{-1}K_+-K_-^*A_-^{-1}K_-.
 \label{eq:t10v25-interface-Schur}
\end{equation}

\begin{theorem}[Exact interface factorization]
\label{thm:t10v25-interface-factorization}
The following identities hold:
\begin{align}
 \det L&=(\det A_+)(\det A_-)(\det\mathcal S_B),
 \label{eq:t10v25-determinant-factorization}\\
 \langle (x_+,x_-,b),L(x_+,x_-,b)\rangle
 &=\sum_{\epsilon=\pm}
 \left\langle x_\epsilon+A_\epsilon^{-1}K_\epsilon b,
 A_\epsilon(x_\epsilon+A_\epsilon^{-1}K_\epsilon b)\right\rangle
 +\langle b,\mathcal S_Bb\rangle.
 \label{eq:t10v25-square-factorization}
\end{align}
If $L>0$, then $\mathcal S_B>0$.
\end{theorem}

\begin{proof}
Complete the square independently in $x_+$ and $x_-$.  This gives
\eqref{eq:t10v25-square-factorization}.  The triangular change of variables
$x_\epsilon\mapsto x_\epsilon+A_\epsilon^{-1}K_\epsilon b$ has determinant
one, so taking determinants gives
\eqref{eq:t10v25-determinant-factorization}.  Setting each completed-square
interior variable to zero shows
$\langle b,\mathcal S_Bb\rangle>0$ for $b\neq0$ when $L>0$.
\end{proof}

\begin{corollary}[Interior Gaussian integration]
\label{cor:t10v25-interior-integration}
For a source $f=(f_+,f_-,f_B)$, integrating the two positive interior
variables gives the boundary quadratic form
\begin{equation}
 \frac12\langle b,\mathcal S_Bb\rangle
 -\left\langle
 f_B-K_+^*A_+^{-1}f_+-K_-^*A_-^{-1}f_-,b
 \right\rangle
 \label{eq:t10v25-boundary-action}
\end{equation}
and the density factor
$(\det A_+)^{-1/2}(\det A_-)^{-1/2}$.
\end{corollary}

\begin{proof}
Complete the affine squares in the two interior variables and use the standard
finite-dimensional Gaussian formula.
\end{proof}

\subsection{Dirichlet-to-Neumann interpretation}

Consider a positive scalar linearized constraint
\begin{equation}
 \mathcal Lu=-8\Delta u+q(x)u=f,
 \qquad q\geq q_0>0,
 \label{eq:t10v25-positive-linearized}
\end{equation}
on two blocks meeting at an interface $\Gamma$.  For prescribed trace $b$,
let $u_\pm(b)$ solve the two Dirichlet problems.  Define the outward flux
maps, including the fixed source response, by
\begin{equation}
 \Lambda_\pm(b):=8\partial_{n_\pm}u_\pm(b)|_\Gamma.
 \label{eq:t10v25-DtN}
\end{equation}

\begin{proposition}[Interface equation]
\label{prop:t10v25-interface-equation}
The two block solutions glue to a weak solution on their union exactly when
their traces agree and
\begin{equation}
 \Lambda_+(b)+\Lambda_-(b)=0.
 \label{eq:t10v25-flux-matching}
\end{equation}
For the homogeneous source response, the derivative of the left side is the
positive interface Schur operator $\mathcal S_B$ in
\eqref{eq:t10v25-interface-Schur}.
\end{proposition}

\begin{proof}
Integrate the weak equation by parts separately on both blocks.  Outer
boundary terms have their prescribed values, while the shared interface term
is $\langle\Lambda_+(b)+\Lambda_-(b),v|_\Gamma\rangle$.  It vanishes for every
test trace precisely under \eqref{eq:t10v25-flux-matching}.  In a finite
discretization, eliminating the zero-trace interiors is exactly the square
completion of Theorem~\ref{thm:t10v25-interface-factorization}; its remaining
quadratic boundary operator is $\mathcal S_B$.
\end{proof}

\subsection{Nonlinear Lichnerowicz interface map}

The V23 equation is the Euler--Lagrange equation of
\begin{equation}
 \mathcal E(\phi)=\int_\Omega
 \left[
 4|\bar\nabla\phi|^2+\frac12\bar R\phi^2
 +\frac b6\phi^6+\frac a6\phi^{-6}
 +\frac c2\phi^{-2}
 \right]dv_{\bar g}.
 \label{eq:t10v25-Lichnerowicz-energy}
\end{equation}
On the V23 sign chamber its second variation is strictly positive.

\begin{proposition}[Nonlinear Dirichlet-to-Neumann monotonicity]
\label{prop:t10v25-nonlinear-DtN}
Let $\phi_\pm(b)$ be the unique V23 positive solutions on the two blocks with
shared trace $b$ in a fixed positive barrier interval.  Let
$\mathscr E_\pm(b)$ be the minimized energies and
\begin{equation}
 \mathscr N_\pm(b):=D\mathscr E_\pm(b)
 \label{eq:t10v25-nonlinear-DtN}
\end{equation}
their conormal fluxes.  Then
\begin{equation}
 \mathscr N_+(b)+\mathscr N_-(b)=0
 \label{eq:t10v25-nonlinear-interface}
\end{equation}
is the exact nonlinear gluing equation.  Its derivative is the sum of the two
positive linearized Dirichlet-to-Neumann maps and is strictly positive.
Consequently it has at most one solution in the barrier interval.
\end{proposition}

\begin{proof}
The first variation of the minimized energy has no interior term because
$\phi_\pm(b)$ solve the Euler--Lagrange equation; the remaining term is the
outward conormal flux.  Stationarity of
$\mathscr E_+(b)+\mathscr E_-(b)$ gives
\eqref{eq:t10v25-nonlinear-interface}.  For an interface variation $h$, let
$v_\pm$ solve the homogeneous linearized equation with boundary trace $h$.
Then
\begin{equation}
 \langle h,D\mathscr N_\pm(b)h\rangle
 =\int_{\Omega_\pm}
 \left(8|\bar\nabla v_\pm|^2+q_\pm v_\pm^2\right)dv_{\bar g}>0
 \label{eq:t10v25-DtN-positive}
\end{equation}
for $h\neq0$, because $q_\pm\geq q_0>0$.  Strict monotonicity gives
uniqueness.
\end{proof}

\begin{corollary}[Existence from the union problem]
\label{cor:t10v25-interface-existence}
If the coefficients on the union lie in the V23 sign chamber with compatible
outer boundary data, the unique global V23 solution restricts to block
solutions whose interface trace solves
\eqref{eq:t10v25-nonlinear-interface}.  Therefore the interface equation has
exactly one solution.
\end{corollary}

\begin{proof}
Restrict the global solution and use weak flux continuity.  Existence follows;
uniqueness is Proposition~\ref{prop:t10v25-nonlinear-DtN}.
\end{proof}

\subsection{Vector constraint gluing}

For the V24 conformal vector operator, impose a shared Dirichlet trace on the
interface and solve on zero-trace interiors.  Let
$A_\pm^{V}$ be the resulting positive interior matrices and
$\mathcal S_B^V$ their interface Schur operator.

\begin{proposition}[Vector interface quotient]
\label{prop:t10v25-vector-interface}
On the orthogonal complement of the global conformal-Killing traces,
$\mathcal S_B^V$ is positive and the vector solutions glue exactly when the
sum of their conformal tractions vanishes.  The kernel of
$\mathcal S_B^V$ consists precisely of interface traces of global homogeneous
solutions not removed by the outer boundary conditions; these traces belong
to the V24 conformal-charge ledger.
\end{proposition}

\begin{proof}
The quadratic energy is
$\int|\mathbb LW|^2$.  Its block minimization at fixed trace gives the
Dirichlet-to-Neumann Schur form.  Vanishing of the first variation is traction
matching.  The energy of an interface trace is zero exactly when its two
extensions have $\mathbb LW=0$ and glue to a global homogeneous conformal
Killing solution.  Quotienting those traces makes the form positive.
\end{proof}

\begin{firewall}[Interior Dirichlet positivity is not global invertibility]
\label{fire:t10v25-vector-kernel}
Zero-trace block interiors can have a positive vector operator while the
glued interface Schur operator has global conformal-Killing zero modes.  Those
modes must be retained before taking the vector constraint determinant.
\end{firewall}

\subsection{Constraint determinant gluing}

At finite regulator, let $L_\phi$ and $L_W$ denote the full scalar and vector
constraint Jacobians, with all exact interface zero modes removed from the
vector determinant.

\begin{theorem}[Finite constraint determinant factorization]
\label{thm:t10v25-constraint-determinant}
The scalar determinant satisfies
\begin{equation}
 \det L_\phi
 =\det A_+^\phi\det A_-^\phi\det\mathcal S_B^\phi,
 \label{eq:t10v25-scalar-det}
\end{equation}
and the vector complement determinant satisfies
\begin{equation}
 \det{}'L_W
 =\det A_+^V\det A_-^V\det{}'\mathcal S_B^V.
 \label{eq:t10v25-vector-det}
\end{equation}
Consequently the inverse coarea density factors into two interior owners and
one shared-interface owner for each constraint.
\end{theorem}

\begin{proof}
Order the finite variables by positive interior, negative interior, and
shared trace.  Apply Theorem~\ref{thm:t10v25-interface-factorization} to the
scalar matrix.  For the vector matrix, first restrict to the complement of
global conformal-Killing traces and apply the same identity; this is exactly
the primed determinant.
\end{proof}

\begin{firewall}[Continuum gluing anomaly]
\label{fire:t10v25-continuum-gluing}
Equations \eqref{eq:t10v25-scalar-det}--
\eqref{eq:t10v25-vector-det} are exact finite-matrix identities.  Zeta, heat,
or other renormalized continuum determinants can acquire local gluing
counterterms.  Those terms must be derived and assigned to the V17 graded
boundary complex; the finite identity alone does not prove their absence.
\end{firewall}

\subsection{Aligned interface coordinates}

Write the exact interior solution with prescribed trace $b$ as
\begin{equation}
 x_\pm=x_\pm^{(0)}-A_\pm^{-1}K_\pm b.
 \label{eq:t10v25-aligned-interior}
\end{equation}
Define the aligned zero-trace deviation
\begin{equation}
 e_\pm:=x_\pm+A_\pm^{-1}K_\pm b-x_\pm^{(0)}.
 \label{eq:t10v25-aligned-deviation}
\end{equation}

\begin{proposition}[Exact trace-to-interior triangularization]
\label{prop:t10v25-interface-triangularization}
At fixed prescribed interface trace, the aligned interior equation is
$e_\pm=0$ and its derivative with respect to the retained trace vanishes.
All trace response is carried by the interface Schur equation.  For a
nonlinear constraint, the same statement holds in the moving graph chart
$e_\pm=x_\pm-x_\pm(b)$.
\end{proposition}

\begin{proof}
The linear claim is immediate from
\eqref{eq:t10v25-aligned-deviation}.  Nonlinearly, differentiating
$x_\pm-x_\pm(b)$ at fixed aligned coordinate cancels the graph derivative.
The physical derivative $Dx_\pm(b)$ has not vanished; it appears when the
trace itself evolves and in the derivative of the interface flux.
\end{proof}

This is the constraint analogue of the exact retained/stable pullback
triangularization found in the audited YM note.

\begin{definition}[V25 constraint-gluing card]
\label{def:t10v25-gluing-card}
A pair of blocks passes the card if:
\begin{enumerate}
\item their outer and interface data lie in one common V23 sign chamber;
\item the scalar nonlinear interface equation has its unique positive trace;
\item vector conformal-Killing traces are retained and the complement traction
equation is solved;
\item scalar and vector determinants use the exact factorizations
\eqref{eq:t10v25-scalar-det}--\eqref{eq:t10v25-vector-det};
\item continuum determinant counterterms, if any, have boundary owners; and
\item stable interior derivatives are taken in the aligned graph coordinates
of Proposition~\ref{prop:t10v25-interface-triangularization}.
\end{enumerate}
\end{definition}

\begin{assessment}[Progress and bottleneck after V25]
\label{assess:t10v25}
V25 proves exact finite constraint gluing, including nonlinear scalar flux
matching, the vector conformal-Killing quotient, and determinant
factorization.  The algebra now has the reflected-square shape: two interior
factors and one positive interface factor.  What remains is to show that this
shape survives the full constrained density, including matter, ghosts,
fermion phase, boundary counterterms, and the reflection parities of all
interface coordinates.
\end{assessment}

\begin{programme}[Tenth-series V26: reflected constraint factorization]
\label{prog:t10v25-V26}
V26 will place the two blocks on opposite sides of a reflection plane, prove
the scalar/vector coarea density is a positive boundary measure times a
positive-half amplitude and its conjugate, and state the exact additional
condition on the chiral determinant line.  It will then insert the result into
the V10/V22 reflection card.
\end{programme}

\begin{firewall}[Claim boundary after tenth-series V25]
\label{fire:t10v25-boundary}
The finite determinant and interface identities do not yet prove reflected
factorization of the full density, continuum determinant gluing, chiral phase
trivialization, the ultraviolet fixed point, or the full SM--Einstein
continuum.
\end{firewall}

\paragraph{Parent-version record.}
V25 was formed from
\texttt{Renewal\_SM\_Einstein\_Continuum\_Research\_Notes\_V24.tex}.
The V24 parent had $9\,679$ validator-counted source lines, $367\,386$ bytes,
and SHA--256
\[
 \texttt{A306CB806A9CF343B0FE093193B6570D45EC185D8ABF9213866EE243196B103F}.
\]

% =============================================================================
% END APPENDED TENTH-SERIES V25 MATERIAL
% =============================================================================

% =============================================================================
% APPENDED TENTH-SERIES V26 MATERIAL
% =============================================================================

\section{Tenth-series V26: reflected factorization of the constraint density}
\label{sec:v10v26}

V25 gives two identical interior determinant owners and one positive
interface determinant.  The interface operator can still depend on fields
from both halves, so its inverse determinant is not automatically a boundary
constant.  V26 represents that determinant by auxiliary boundary bosons.  The
additive Dirichlet-to-Neumann decomposition then produces a positive-half
amplitude and its reflected conjugate.

\subsection{Abstract reflected conditional measures}

Let $b$ be reflection-fixed boundary data and let $x_+,x_-$ be fields in the
two open halves, with reflection identifying the negative field space with
the conjugate of the positive one.

\begin{theorem}[Conditional reflected-square criterion]
\label{thm:t10v26-reflected-square}
Suppose a finite regulated density has the form
\begin{equation}
 d\mu=d\nu_B(b)\,dx_+dx_-\,
 W_+(x_+,b)\overline{W_+(\Theta x_-,b)},
 \qquad d\nu_B\geq0,
 \label{eq:t10v26-conditional-density}
\end{equation}
and the integrals are absolutely convergent.  Then every positive-half
observable $F$ satisfies
\begin{equation}
 \int F\,\Theta F\,d\mu
 =\int d\nu_B(b)
 \left|\int F(x_+,b)W_+(x_+,b)dx_+\right|^2\geq0.
 \label{eq:t10v26-reflected-square}
\end{equation}
\end{theorem}

\begin{proof}
Use reflection to change the negative-half variables to a second copy of the
positive-half variables.  Fubini's theorem factors the two conditional
integrals, and they are complex conjugates.  The boundary measure is
nonnegative.
\end{proof}

\subsection{Auxiliary representation of an inverse determinant}

\begin{lemma}[Complex boundary boson]
\label{lem:t10v26-complex-boson}
If $S$ is a positive Hermitian $k\times k$ matrix, then
\begin{equation}
 (\det S)^{-1}
 =\pi^{-k}\int_{\mathbb C^k}e^{-\eta^*S\eta}\,d^{2k}\eta.
 \label{eq:t10v26-complex-boson}
\end{equation}
\end{lemma}

\begin{proof}
Unitary diagonalization reduces the integral to the product of
$k$ scalar identities
$\pi^{-1}\int_{\mathbb C}e^{-\lambda|z|^2}d^2z=\lambda^{-1}$.
\end{proof}

Suppose the interface operator is an additive Dirichlet-to-Neumann sum
\begin{equation}
 S_B(x_+,x_-,b)
 =\Lambda_+(x_+,b)+Theta\Lambda_+(\Theta x_-,b)\Theta^{-1},
 \label{eq:t10v26-additive-DtN}
\end{equation}
on a reflection-fixed interface space.  Both summands are positive by the
V25 energy identity.

\begin{corollary}[Half-factorization of the interface determinant]
\label{cor:t10v26-interface-factorization}
Under \eqref{eq:t10v26-additive-DtN},
\begin{align}
 (\det S_B)^{-1}
 ={}&\pi^{-k}\int d^{2k}\eta\,
 e^{-\eta^*\Lambda_+(x_+,b)\eta}\\
 &\hspace{19mm}\times
 \overline{
 e^{-(\Theta\eta)^*\Lambda_+(\Theta x_-,b)(\Theta\eta)}}.
 \label{eq:t10v26-interface-factorization}
\end{align}
\end{corollary}

\begin{proof}
Insert \eqref{eq:t10v26-additive-DtN} into
Lemma~\ref{lem:t10v26-complex-boson} and split the exponential of the sum.
\end{proof}

The auxiliary field is a boundary owner.  It is not a new propagating bulk
degree of freedom and is integrated out to recover the exact coarea density.

\subsection{Scalar and vector constraint factors}

At the reflected V23 solution, write the scalar Jacobian factorization as
\begin{equation}
 \det L_\phi
 =\det A_+^\phi\det A_-^\phi\det S_B^\phi.
 \label{eq:t10v26-scalar-Jacobian}
\end{equation}
Reflection gives
$A_-^\phi=U_\phi A_+^\phi U_\phi^{-1}$, so
\begin{equation}
 (\det L_\phi)^{-1}
 =|\det A_+^\phi|^{-2}(\det S_B^\phi)^{-1}.
 \label{eq:t10v26-scalar-inverse-det}
\end{equation}

For the vector constraint, restrict to the complement of the V24
conformal-Killing traces:
\begin{equation}
 (\det{}'L_W)^{-1}
 =|\det A_+^V|^{-2}(\det{}'S_B^V)^{-1}.
 \label{eq:t10v26-vector-inverse-det}
\end{equation}

\begin{proposition}[Constraint coarea half amplitudes]
\label{prop:t10v26-coarea-halves}
Introduce independent complex boundary fields $\eta_\phi$ and $\eta_V$ for
the scalar and vector interface complements.  The product of the two
constraint coarea factors is
\begin{equation}
 \int d\eta_\phi d\eta_V\,
 H_+(x_+,b,\eta_\phi,\eta_V)
 \overline{H_+(\Theta x_-,b,\Theta\eta_\phi,\Theta\eta_V)},
 \label{eq:t10v26-coarea-half-amplitude}
\end{equation}
up to positive powers of $\pi$, where
\begin{equation}
 H_+=(\det A_+^\phi)^{-1}(\det A_+^V)^{-1}
 e^{-\eta_\phi^*\Lambda_+^\phi\eta_\phi
    -\eta_V^*\Lambda_+^V\eta_V}.
 \label{eq:t10v26-Hplus}
\end{equation}
Every determinant in $H_+$ is positive on its declared complement.
\end{proposition}

\begin{proof}
Use \eqref{eq:t10v26-scalar-inverse-det} and
\eqref{eq:t10v26-vector-inverse-det}, then apply
Corollary~\ref{cor:t10v26-interface-factorization} to each interface
determinant.  Reflection conjugates the two interior determinants and the two
Dirichlet-to-Neumann exponentials.
\end{proof}

\begin{firewall}[All interface zero modes are retained first]
\label{fire:t10v26-interface-zero-modes}
The complex-boson integral diverges on a zero eigenvalue.  Scalar positivity
comes from the V23 gap.  Vector conformal-Killing traces must be removed into
the V24 ledger before \eqref{eq:t10v26-vector-inverse-det} is written.  A
primed determinant without the corresponding retained coordinate is not a
measure construction.
\end{firewall}

\subsection{The nonlinear constrained action}

Assume the physical action evaluated on the unique constraint solution has
the local reflected decomposition
\begin{equation}
 S_{\rm red}(x_+,x_-,b)
 =S_+(x_+,b)+S_+(\Theta x_-,b)^*+S_B(b),
 \qquad e^{-S_B(b)}db\geq0.
 \label{eq:t10v26-action-factorization}
\end{equation}
All nonconstraint positive densities that genuinely factor by halves are
included in $S_+$; reflection-fixed positive densities are included in
$S_B$.

\begin{theorem}[Reflected positivity of the bosonic constraint density]
\label{thm:t10v26-bosonic-constraint-OS}
Suppose:
\begin{enumerate}
\item the V25 scalar and vector gluing cards hold on reflected blocks;
\item the full constrained physical action has
\eqref{eq:t10v26-action-factorization};
\item all remaining gauge-quotient, chart, and density owners are nonnegative
and have the same half/boundary factorization; and
\item the V22 large-field card makes the resulting integrals finite.
\end{enumerate}
Then the bosonic constraint-reduced density is reflection positive on
positive-half observables.
\end{theorem}

\begin{proof}
Insert the two auxiliary boundary fields using Proposition
\ref{prop:t10v26-coarea-halves}.  Multiply $H_+$ by $e^{-S_+}$ and all other
positive-half density factors.  Put $b,\eta_\phi,\eta_V$ in the boundary
measure together with $e^{-S_B}$.  The total density now has the form
\eqref{eq:t10v26-conditional-density}; apply Theorem
\ref{thm:t10v26-reflected-square}.
\end{proof}

\begin{firewall}[Matter-dependent interface operators]
\label{fire:t10v26-matter-dependent-interface}
The auxiliary-field proof permits $\Lambda_+$ to depend on positive-half
matter and metric fields because that dependence remains inside the half
amplitude.  It requires the global interface operator to be the additive sum
of the two half Dirichlet-to-Neumann maps.  A regulator term coupling the two
open interiors directly would violate \eqref{eq:t10v26-additive-DtN} and must
be moved to an explicit boundary channel or excluded.
\end{firewall}

\subsection{Gauge quotient and ghosts}

\begin{firewall}[Ghost determinants are not positive measures]
\label{fire:t10v26-ghost-positivity}
Grassmann Faddeev--Popov ghosts and a sign-changing Gribov determinant do not
define a nonnegative boundary measure.  Theorem
\ref{thm:t10v26-bosonic-constraint-OS} therefore applies to a physical gauge
quotient whose orbit/coarea density is nonnegative, or to a gauge-fixed
construction after a separate theorem proves reflection positivity on the
BRST physical observable space.  Formal cancellation of ghost loops is not
that theorem.
\end{firewall}

The V4 based interior quotient is compatible with the present factorization
when its gauge parameters have zero shared trace and all shared boundary
frames are retained.  Then each half quotient is conditional on the same
boundary frame and the orbit density is assigned to its actual half or
interface owner.

\subsection{The chiral determinant-line condition}

Let $\mathcal L\to\mathcal Q$ be the chiral determinant line over the
regulated bosonic configuration space.  Reflection induces an antilinear map
\begin{equation}
 \mathfrak j_q:\mathcal L_q\longrightarrow
 \overline{\mathcal L_{\Theta q}}.
 \label{eq:t10v26-reflection-line}
\end{equation}

\begin{definition}[Reflection-real gluing trivialization]
\label{def:t10v26-reflection-trivialization}
A trivializing section $\tau(q)\neq0$ passes the determinant-line card if:
\begin{enumerate}
\item $\mathfrak j_{\Theta q}\mathfrak j_q=1$;
\item $\tau(\Theta q)=\mathfrak j_q\tau(q)$;
\item on a cut configuration there is a boundary fermion space $\mathcal H_b$
and a half amplitude $\Psi_+(q_+,b)\in\mathcal H_b$ such that
\begin{equation}
 Z_F(q_+\cup_b q_-)
 =\langle\Psi_+(\Theta q_-,b),\Psi_+(q_+,b)\rangle_{\mathcal H_b};
 \label{eq:t10v26-fermion-gluing}
\end{equation}
and
\item the local connection on $\mathcal L$ agrees with the V17
anomaly-cancelled counterterm connection.
\end{enumerate}
\end{definition}

\begin{proposition}[Flat line criterion for a global phase]
\label{prop:t10v26-flat-line}
Suppose the local anomaly cancellation makes the determinant-line connection
flat on a connected component of $\mathcal Q$.  A global nonzero parallel
section exists if and only if the connection holonomy is trivial.  A
reflection-real parallel section exists if, in addition, its value at one
reflection-fixed base point can be chosen fixed by $\mathfrak j$ and the
reflection lift is compatible with parallel transport.
\end{proposition}

\begin{proof}
Parallel transport a chosen nonzero base vector along paths.  Path
independence is equivalent to trivial holonomy around every loop.  Under the
reflection compatibility, transporting a $\mathfrak j$-fixed base vector
along a reflected path gives the reflected transported vector, proving the
real condition.  Conversely, a global parallel section is unchanged around
loops and therefore forces trivial holonomy.
\end{proof}

\begin{firewall}[Local anomaly sums do not prove the gluing card]
\label{fire:t10v26-line-firewall}
V17 proves the Standard Model representation sums that cancel the local
anomaly polynomial.  Proposition~\ref{prop:t10v26-flat-line} shows the
remaining global question: flatness does not imply trivial holonomy.
Furthermore, a scalar trivialization alone does not prove the boundary inner
product formula \eqref{eq:t10v26-fermion-gluing}.  Both are required before
the chiral factor can join the reflected square.
\end{firewall}

\begin{theorem}[Full finite reflected-step criterion]
\label{thm:t10v26-full-reflected-criterion}
Assume the bosonic hypotheses of Theorem
\ref{thm:t10v26-bosonic-constraint-OS}, a nonnegative physical gauge quotient,
and Definition~\ref{def:t10v26-reflection-trivialization}.  Then the complete
finite constrained SM--Einstein block density is reflection positive on
gauge-invariant positive-half observables.
\end{theorem}

\begin{proof}
The bosonic and constraint factors give a boundary integral of a reflected
square.  Equation \eqref{eq:t10v26-fermion-gluing} inserts the fermion factor
as a boundary Hilbert-space inner product.  Tensor the boundary spaces and
apply the same conditional-square calculation as Theorem
\ref{thm:t10v26-reflected-square}.
\end{proof}

\subsection{Insertion into the earlier cards}

\begin{corollary}[Closure of the abstract reflection row]
\label{cor:t10v26-reflection-row}
When Theorem~\ref{thm:t10v26-full-reflected-criterion} holds uniformly, the
reflection rows of the V10 nonlinear gate, V13 boundary-action card, V22
large-field card, and V25 constraint-gluing card reduce to their stated
locality and uniform-integrability estimates; the constraint coarea factors
introduce no additional sign.
\end{corollary}

\begin{proof}
The scalar and vector Jacobians are positive on their complements and their
interface inverses have the exact auxiliary reflected factorization.  The
remaining factors are precisely the hypotheses of the full reflected-step
criterion.
\end{proof}

\begin{assessment}[Progress and bottleneck after V26]
\label{assess:t10v26}
V26 proves reflected factorization of the scalar and vector constraint coarea
densities, even when their interface operators depend on half fields.  It
also isolates the exact fermionic requirement: a reflection-real determinant
line with trivial holonomy and a boundary Hilbert-space gluing formula.  The
next upstream task is to construct or obstruct that line data for a finite
chiral Standard Model regulator; local anomaly cancellation alone is no
longer an adequate placeholder.
\end{assessment}

\begin{programme}[Tenth-series V27: finite chiral determinant gluing]
\label{prog:t10v26-V27}
V27 will work at finite lattice dimension with a chiral Dirac block matrix.
It will derive its Schur determinant gluing identity, identify the boundary
fermion state and reflection adjointness condition, and separate the phase
holonomy from the positive singular-value determinant.  The result will be a
finite algebraic card that a Ginsparg--Wilson/overlap-type regulator would
have to satisfy, without assuming such a regulator has already been built on
the dynamical curved block.
\end{programme}

\begin{firewall}[Claim boundary after tenth-series V26]
\label{fire:t10v26-boundary}
The bosonic constraint factorization is proved under its cards, but the
chiral determinant-line gluing condition is not.  This version also does not
prove a nonnegative global gauge quotient, the ultraviolet fixed point,
cofinal measure compactness, or the full SM--Einstein continuum.
\end{firewall}

\paragraph{Parent-version record.}
V26 was formed from
\texttt{Renewal\_SM\_Einstein\_Continuum\_Research\_Notes\_V25.tex}.
The V25 parent had $10\,053$ validator-counted source lines, $382\,021$ bytes,
and SHA--256
\[
 \texttt{C1AB2B98B067D7720058F63EC252C656DBFD1C745876CE92DDF691029AE9C341}.
\]

% =============================================================================
% END APPENDED TENTH-SERIES V26 MATERIAL
% =============================================================================

% =============================================================================
% APPENDED TENTH-SERIES V27 MATERIAL
% =============================================================================

\section{Tenth-series V27: finite chiral determinant gluing and the boundary
CAR test}
\label{sec:v10v27}

V26 reduced the unproved fermionic row to a finite question: does the chiral
determinant admit a reflection-real phase and a positive boundary-state
gluing?  V27 separates three algebraic layers.  Ordinary Schur factorization
controls determinant magnitudes, the boundary CAR covariance controls
reflection positivity, and cycle holonomy controls whether local chiral
bases define a global phase.

\subsection{Finite chiral blocks and Schur determinants}

Let $D(q):V_-(q)\to V_+$ be a finite regulated Weyl operator in a fixed
topological sector, with equal finite dimensions after all zero modes and
insertions have been declared.  Split its variables into interior and
boundary components:
\begin{equation}
 D=
 \begin{pmatrix}
 A&K\\ L&B
 \end{pmatrix},
 \qquad A\text{ invertible},
 \qquad S_B:=B-LA^{-1}K.
 \label{eq:t10v27-chiral-block}
\end{equation}

\begin{lemma}[Finite chiral Schur identity]
\label{lem:t10v27-chiral-Schur}
\begin{equation}
 \det D=(\det A)(\det S_B).
 \label{eq:t10v27-chiral-Schur}
\end{equation}
Moreover
\begin{equation}
 D=
 \begin{pmatrix}I&0\\LA^{-1}&I\end{pmatrix}
 \begin{pmatrix}A&0\\0&S_B\end{pmatrix}
 \begin{pmatrix}I&A^{-1}K\\0&I\end{pmatrix}.
 \label{eq:t10v27-chiral-LDU}
\end{equation}
\end{lemma}

\begin{proof}
Multiply the three matrices in \eqref{eq:t10v27-chiral-LDU}.  The two
triangular factors have determinant one, proving
\eqref{eq:t10v27-chiral-Schur}.
\end{proof}

Unlike the positive constraint Jacobians, neither $A$ nor $S_B$ need be
Hermitian.  Their determinant phases cannot be replaced by absolute values
without changing the chiral theory.

\begin{proposition}[Reflection pairing of interior phases]
\label{prop:t10v27-interior-phase-pairing}
Suppose an antiunitary reflection $J$ identifies the negative interior Weyl
block with the adjoint positive block,
\begin{equation}
 A_-(\Theta q)=J A_+(q)^*J^{-1}.
 \label{eq:t10v27-reflection-adjointness}
\end{equation}
Then
\begin{equation}
 \det A_-(\Theta q)=\overline{\det A_+(q)},
 \qquad
 \det A_+\det A_-=|\det A_+|^2.
 \label{eq:t10v27-interior-phase-cancel}
\end{equation}
\end{proposition}

\begin{proof}
Antiunitary conjugation and adjunction complex-conjugate the determinant.
Multiply the two factors.
\end{proof}

\begin{firewall}[Interface phase remains]
\label{fire:t10v27-interface-phase}
Proposition~\ref{prop:t10v27-interior-phase-pairing} cancels paired interior
phases.  It does not make the boundary Schur determinant positive.  The
interface factor must be realized as a CAR boundary-state pairing, or its
phase must pass a separate determinant-line gluing theorem.
\end{firewall}

\subsection{Grassmann integration as a boundary state}

Let $\xi,\bar\xi$ be boundary Grassmann generators and
$\chi_+,\bar\chi_+$ positive-interior generators.  For a quadratic half
action $S_+(\chi_+,\xi;q_+)$ define
\begin{equation}
 \Psi_+(\bar\xi,\xi;q_+)
 :=\int d\bar\chi_+d\chi_+\,
 e^{-S_+(\chi_+,\xi;q_+)}.
 \label{eq:t10v27-boundary-state}
\end{equation}
This is an element of the finite boundary exterior algebra.

\begin{proposition}[Exact Berezin gluing]
\label{prop:t10v27-Berezin-gluing}
Suppose the full quadratic action is the sum of a positive-half action, its
reflected adjoint, and a boundary quadratic action.  Then Berezin--Fubini
integration gives
\begin{equation}
 Z_F(q_+\cup_bq_-)
 =\langle\Psi_+(\Theta q_-),\Psi_+(q_+)\rangle_{B},
 \label{eq:t10v27-Berezin-gluing}
\end{equation}
where $\langle\cdot,\cdot\rangle_B$ is the sesquilinear pairing induced by the
boundary quadratic form and the chosen reflection polarization.
\end{proposition}

\begin{proof}
Integrate the two disjoint interior Grassmann families first.  The negative
integral is the reflected adjoint polynomial, including the reversed Berezin
orientation.  The remaining boundary Berezin integral is exactly the pairing
of the two exterior-algebra elements.
\end{proof}

\begin{firewall}[A Berezin pairing need not be positive]
\label{fire:t10v27-Berezin-not-positive}
Equation \eqref{eq:t10v27-Berezin-gluing} is an algebraic gluing identity.
Reflection positivity requires $\langle\cdot,\cdot\rangle_B$ to be a positive
CAR/Fock inner product on the physical boundary space.  An arbitrary
fermionic boundary matrix can define an indefinite pairing.
\end{firewall}

\subsection{The finite CAR positivity criterion}

Let $\psi_1,\ldots,\psi_N$ generate the positive boundary one-particle
space.  For a reflection-invariant quasi-free functional $\omega$, define
the cross matrix
\begin{equation}
 K_{ij}:=\omega((\Theta\psi_i)\psi_j).
 \label{eq:t10v27-cross-matrix}
\end{equation}

\begin{theorem}[Quasi-free boundary reflection criterion]
\label{thm:t10v27-CAR-criterion}
Assume charge sectors of different exterior degree are orthogonal.  Then
\begin{equation}
 \omega((\Theta F)F)\geq0
 \quad\text{for every positive-boundary polynomial }F
 \label{eq:t10v27-CAR-OS}
\end{equation}
if and only if $K\succeq0$.  For degree $r$, the OS Gram matrix is the
$r$-th exterior power $\bigwedge^rK$.
\end{theorem}

\begin{proof}
For ordered index sets $I=(i_1<\cdots<i_r)$ and
$J=(j_1<\cdots<j_r)$, quasi-free Wick contraction gives
\begin{equation}
 \omega((\Theta\psi_I)\psi_J)=\det(K_{I,J}).
 \label{eq:t10v27-exterior-minors}
\end{equation}
These are the matrix entries of $\bigwedge^rK$.  If $K\succeq0$, diagonalize
it unitarily; every exterior power has nonnegative eigenvalues given by
products of eigenvalues of $K$.  Summing the orthogonal degree sectors proves
\eqref{eq:t10v27-CAR-OS}.  Conversely, take $F=\sum_jv_j\psi_j$.  Positivity
for every $v$ gives $v^*Kv\geq0$, hence $K\succeq0$.
\end{proof}

\begin{corollary}[Boundary-state positivity]
\label{cor:t10v27-boundary-state-positive}
If the boundary polarization in Proposition
\ref{prop:t10v27-Berezin-gluing} has cross covariance $K\succeq0$, the pairing
in \eqref{eq:t10v27-Berezin-gluing} is positive semidefinite and supplies the
fermionic row required by V26.
\end{corollary}

\begin{proof}
Apply Theorem~\ref{thm:t10v27-CAR-criterion} to the boundary polynomial
$\Psi_+$.
\end{proof}

\subsection{Singular values and phase}

For an invertible finite chiral matrix, polar decomposition gives
\begin{equation}
 D=U|D|,
 \qquad
 \det D=(\det U)\det|D|,
 \qquad \det|D|>0.
 \label{eq:t10v27-polar-determinant}
\end{equation}

\begin{proposition}[Magnitude bound is phase-blind]
\label{prop:t10v27-magnitude-phase}
Any Gram, singular-value, or $D^*D$ estimate controls $\det|D|$ but leaves the
$U(1)$ phase $\det U$ undetermined.  Along a smooth parameter path with no
zero crossing,
\begin{equation}
 d\log\det D=\operatorname{Tr}(D^{-1}dD),
 \label{eq:t10v27-phase-connection}
\end{equation}
whose imaginary part is the determinant-line phase connection.
\end{proposition}

\begin{proof}
The first claim is immediate from \eqref{eq:t10v27-polar-determinant}.
Differentiating the finite determinant gives
\eqref{eq:t10v27-phase-connection}; taking imaginary parts records the change
of argument.
\end{proof}

\begin{firewall}[Zero crossings and index sectors]
\label{fire:t10v27-zero-crossings}
At a zero singular value, \eqref{eq:t10v27-phase-connection} is undefined and
the Weyl determinant may be a section between unequal-dimensional spaces.
Zero modes require explicit fermion insertions and topological-sector data.
A primed determinant alone does not define the missing multilinear section.
\end{firewall}

\subsection{Finite holonomy test}

Let a finite connected configuration graph have oriented edges $e$ carrying
unit parallel-transport phases $u_e\in U(1)$ with
$u_{\bar e}=u_e^{-1}$.  These phases discretize the determinant-line
connection.

\begin{theorem}[Spanning-tree holonomy criterion]
\label{thm:t10v27-holonomy-test}
A global choice of phases $\tau(v)\in U(1)$ satisfying
\begin{equation}
 \tau(t(e))=u_e\tau(s(e))
 \label{eq:t10v27-phase-section}
\end{equation}
on every edge exists if and only if the ordered product of $u_e$ around every
cycle is one.  It is enough to check the fundamental cycles determined by one
spanning tree.
\end{theorem}

\begin{proof}
If \eqref{eq:t10v27-phase-section} holds, multiplication around a cycle
telescopes to one.  Conversely, choose a root and set $\tau$ by multiplying
edge phases along the unique spanning-tree path.  Every non-tree edge closes
a fundamental cycle; trivial holonomy on that cycle makes
\eqref{eq:t10v27-phase-section} hold for the non-tree edge.  Fundamental
cycles generate the cycle space.
\end{proof}

\begin{corollary}[Product multiplet holonomy]
\label{cor:t10v27-product-holonomy}
For a collection of Weyl species, the determinant-line transport phase is
the product of the species phases.  A global phase for the anomaly-free
multiplet exists on the finite configuration graph exactly when this product
has unit holonomy on every fundamental cycle.
\end{corollary}

\begin{proof}
Tensor products of lines multiply parallel transports and holonomies.  Apply
Theorem~\ref{thm:t10v27-holonomy-test} to the product phase.
\end{proof}

\begin{firewall}[Local curvature cancellation is weaker]
\label{fire:t10v27-curvature-holonomy}
The V17 anomaly sums cancel the local curvature of the product determinant
connection in the continuum representation calculation.  The finite cycle
products in Corollary~\ref{cor:t10v27-product-holonomy} are global data and
must still be one.  A flat connection can have nontrivial cycle holonomy.
\end{firewall}

\subsection{Ginsparg--Wilson-type chiral subspaces}

If a finite Dirac operator satisfies
\begin{equation}
 \gamma_5D+D\gamma_5=aD\gamma_5D,
 \label{eq:t10v27-GW}
\end{equation}
define
\begin{equation}
 \widehat\gamma_5:=\gamma_5(1-aD),
 \qquad
 \widehat P_-:=\frac12(1-\widehat\gamma_5),
 \qquad
 P_+:=\frac12(1+\gamma_5).
 \label{eq:t10v27-GW-projectors}
\end{equation}

\begin{lemma}[Finite chiral map]
\label{lem:t10v27-GW-map}
If $D$ is $\gamma_5$-Hermitian in addition to
\eqref{eq:t10v27-GW}, then $\widehat\gamma_5^2=1$ and
\begin{equation}
 D\widehat P_-=P_+D.
 \label{eq:t10v27-GW-intertwining}
\end{equation}
Thus $P_+D\widehat P_-$ defines a finite Weyl map between the projected
subspaces.
\end{lemma}

\begin{proof}
Multiply \eqref{eq:t10v27-GW} by the appropriate $\gamma_5$ factors and use
$D^*=\gamma_5D\gamma_5$ to obtain
$(\gamma_5(1-aD))^2=1$.  Rearranging the Ginsparg--Wilson relation gives
$D(1-\widehat\gamma_5)=(1+\gamma_5)D$, which is
\eqref{eq:t10v27-GW-intertwining} after division by two.
\end{proof}

The basis of $\operatorname{ran}\widehat P_-$ depends on the bosonic
configuration.  Its change of phase is precisely the determinant-line
connection; the algebraic chiral map does not choose a global basis.

\begin{definition}[V27 finite chiral gluing card]
\label{def:t10v27-chiral-card}
A finite chiral regulator passes the card if:
\begin{enumerate}
\item its projected Weyl maps have fixed declared index on each sector and
zero modes have the required insertions;
\item interior blocks satisfy reflection adjointness
\eqref{eq:t10v27-reflection-adjointness};
\item Berezin integration gives the exact boundary state
\eqref{eq:t10v27-Berezin-gluing};
\item the boundary cross covariance satisfies $K\succeq0$;
\item the product Standard Model determinant connection has trivial
fundamental-cycle holonomy and a reflection-real phase; and
\item all local phase changes agree with the V17 anomaly-compatible
counterterm connection.
\end{enumerate}
\end{definition}

\begin{theorem}[Finite fermionic completion of the V26 criterion]
\label{thm:t10v27-fermionic-completion}
If Definition~\ref{def:t10v27-chiral-card} holds, the finite chiral Standard
Model factor has the reflection-real boundary Hilbert-space gluing required
by Definition~\ref{def:t10v26-reflection-trivialization}.  Combined with the
V26 bosonic hypotheses, the full finite block is reflection positive on its
physical positive-half observable algebra.
\end{theorem}

\begin{proof}
The holonomy condition constructs the global reflection-real phase by
Theorem~\ref{thm:t10v27-holonomy-test}.  Exact Berezin gluing supplies the
boundary state, and Theorem~\ref{thm:t10v27-CAR-criterion} makes its pairing
positive.  Insert this in Theorem~\ref{thm:t10v26-full-reflected-criterion}.
\end{proof}

\begin{assessment}[Progress and bottleneck after V27]
\label{assess:t10v27}
V27 replaces the phrase ``fermion determinant glues'' by three finite tests:
reflection adjointness, positivity of the boundary CAR cross covariance, and
trivial product-line holonomy on fundamental configuration cycles.  It also
shows exactly what a Ginsparg--Wilson relation supplies and what it does not:
the projected chiral map exists, but its field-dependent basis phase remains
the determinant-line problem.
\end{assessment}

\begin{programme}[Tenth-series V28: local curvature and global cycle split]
\label{prog:t10v27-V28}
V28 will compute the finite Berry curvature of the projected chiral subspace,
show how the V17 representation sums cancel its local anomaly component, and
derive a lattice Stokes relation between plaquette curvature and cycle
holonomy.  Any residual torsion/global phase will remain a separate finite
obstruction rather than being inferred away from the local sum.
\end{programme}

\begin{firewall}[Claim boundary after tenth-series V27]
\label{fire:t10v27-boundary}
The finite card has not been verified for an actual curved dynamical
Ginsparg--Wilson regulator.  This version does not prove product-line
holonomy, boundary cross-covariance positivity, the ultraviolet fixed point,
cofinal measure compactness, or the full SM--Einstein continuum.
\end{firewall}

\paragraph{Parent-version record.}
V27 was formed from
\texttt{Renewal\_SM\_Einstein\_Continuum\_Research\_Notes\_V26.tex}.
The V26 parent had $10\,413$ validator-counted source lines, $396\,585$ bytes,
and SHA--256
\[
 \texttt{F205BEA05877D93994859BE37645EA698B114803180F42D8366488C656702FC0}.
\]

% =============================================================================
% END APPENDED TENTH-SERIES V27 MATERIAL
% =============================================================================

% =============================================================================
% APPENDED TENTH-SERIES V28 MATERIAL
% =============================================================================

\section{Tenth-series V28: projector curvature, lattice Stokes, and global
chiral holonomy}
\label{sec:v10v28}

V27 separated the chiral determinant problem into a local connection and
global cycle phases.  This version computes the connection curvature for a
finite field-dependent chiral projector and proves the exact discrete Stokes
relation.  The result makes the logical boundary precise: local anomaly
cancellation can flatten contractible curvature after a local measure-current
construction, while noncontractible holonomy remains independent data.

\subsection{Curvature of a finite projected subspace}

Let $P(q)$ be a smooth rank-$r$ orthogonal projector in a fixed finite Hilbert
space.  On a local chart choose an orthonormal frame
$V(q):\mathbb C^r\to\mathbb C^N$ with
\begin{equation}
 V^*V=I,
 \qquad P=VV^*.
 \label{eq:t10v28-projector-frame}
\end{equation}
The frame connection and curvature are
\begin{equation}
 \mathcal A:=V^*dV,
 \qquad
 \mathcal F:=d\mathcal A+\mathcal A\wedge\mathcal A.
 \label{eq:t10v28-Berry-connection}
\end{equation}

\begin{theorem}[Projector curvature formula]
\label{thm:t10v28-projector-curvature}
\begin{equation}
 \mathcal F=V^*(dP\wedge dP)V,
 \qquad
 \operatorname{tr}\mathcal F
 =\operatorname{Tr}(P,dP\wedge dP).
 \label{eq:t10v28-projector-curvature}
\end{equation}
The determinant line $\det\operatorname{ran}P$ has connection
$\operatorname{tr}\mathcal A$ and curvature
$\operatorname{Tr}(P,dP\wedge dP)$.
\end{theorem}

\begin{proof}
Differentiate $P^2=P$ to get $P(dP)P=0$.  Also
$dV=dPV+V\mathcal A$, obtained by applying $P$ and $I-P$ to $dV$.
Then
\begin{align*}
 d\mathcal A+mathcal A\wedge\mathcal A
 &=dV^*\wedge dV+\mathcal A\wedge\mathcal A\\
 &=dV^*(I-P)\wedge dV
 =V^*dP\wedge dPV.
\end{align*}
Taking the finite trace and using cyclicity gives the second identity.  The
connection on a determinant is the trace of the connection on the underlying
rank-$r$ bundle, so its curvature is the trace above.
\end{proof}

\begin{corollary}[Frame independence]
\label{cor:t10v28-frame-independence}
Under a frame change $V\mapsto Vg$, $g(q)\in U(r)$,
\begin{equation}
 \operatorname{tr}\mathcal A\mapsto
 \operatorname{tr}\mathcal A+d\log\det g,
 \qquad
 \operatorname{tr}\mathcal F\mapsto\operatorname{tr}\mathcal F.
 \label{eq:t10v28-frame-change}
\end{equation}
\end{corollary}

\begin{proof}
Substitution gives
$\mathcal A\mapsto g^{-1}\mathcal Ag+g^{-1}dg$.  Take the trace and then the
exterior derivative.
\end{proof}

For the Ginsparg--Wilson projector $P=\widehat P_-$ of V27, the formula is
finite and exact wherever the rank is constant.  Its sign in the Weyl
determinant line depends on whether the line or its dual is used; V28 fixes
only the geometric magnitude and keeps that convention explicit.

\subsection{Product multiplets and the local anomaly component}

For Weyl species $s$ with determinant lines $\mathcal L_s$, the product line
is $\mathcal L_{\rm tot}=\bigotimes_s\mathcal L_s$.

\begin{lemma}[Additivity of determinant curvature]
\label{lem:t10v28-curvature-additivity}
\begin{equation}
 \mathcal F_{\rm tot}=\sum_s\mathcal F_s.
 \label{eq:t10v28-curvature-additivity}
\end{equation}
\end{lemma}

\begin{proof}
Connections on tensor products add, so their determinant-line curvatures add.
\end{proof}

In a local continuum limit on gauge and metric directions, the degree-six
anomaly polynomial controlling the curvature/descent has coefficients
proportional to
\begin{equation}
 SU(3)^3,quad SU(3)^2Y,quad SU(2)^2Y,quad Y^3,quad
 Yp_1(T).
 \label{eq:t10v28-local-coefficients}
\end{equation}

\begin{proposition}[Cancellation of the universal local component]
\label{prop:t10v28-local-cancellation}
For the one-generation left-handed Standard Model representation list in
\eqref{eq:t10v17-SM-representations}, every coefficient in
\eqref{eq:t10v28-local-coefficients} vanishes.  Hence the sum of the universal
continuum local anomaly-curvature components is zero for each generation.
\end{proposition}

\begin{proof}
The five coefficients are exactly the representation sums evaluated in
Proposition~\ref{prop:t10v17-SM-anomalies}; each is zero.  Curvature additivity
is Lemma~\ref{lem:t10v28-curvature-additivity}.
\end{proof}

\begin{firewall}[Finite regulator curvature is not only its universal limit]
\label{fire:t10v28-finite-curvature}
Proposition~\ref{prop:t10v28-local-cancellation} cancels the universal local
continuum component.  A finite chiral regulator can have irrelevant lattice
curvature terms, chart-transition curvature, and topology-dependent pieces.
They must be shown to vanish or be the coboundary of a local bounded measure
current.  Continuum coefficient cancellation alone does not make
$\operatorname{Tr}(P,dP\wedge dP)$ zero at finite cutoff.
\end{firewall}

\subsection{Configuration-cell cochains}

Let $\mathcal K$ be a finite oriented cell complex approximating one connected
component of the admissible bosonic configuration space.  Assign a unit
parallel phase $u_e\in U(1)$ to each oriented edge, with
$u_{\bar e}=u_e^{-1}$.  For an oriented plaquette $p$, define
\begin{equation}
 \Omega_p:=\prod_{e\subset\partial p}u_e^{\epsilon_{pe}}.
 \label{eq:t10v28-plaquette-holonomy}
\end{equation}

\begin{theorem}[Discrete abelian Stokes formula]
\label{thm:t10v28-lattice-Stokes}
For every integral two-chain $\Sigma=\sum_pn_pp$,
\begin{equation}
 \prod_p\Omega_p^{n_p}
 =\prod_eu_e^{(\partial\Sigma)_e}.
 \label{eq:t10v28-lattice-Stokes}
\end{equation}
Thus the holonomy of a cycle that bounds $\Sigma$ is the product of the
plaquette curvatures in any filling two-chain.
\end{theorem}

\begin{proof}
Insert \eqref{eq:t10v28-plaquette-holonomy} on the left.  The exponent of an
edge $e$ becomes $\sum_pn_p\epsilon_{pe}$, which is exactly the coefficient of
$e$ in $\partial\Sigma$.  Every internal edge cancels with its opposite
orientation.
\end{proof}

\begin{corollary}[Flat plaquettes on a simply connected complex]
\label{cor:t10v28-simply-connected}
If every $\Omega_p=1$ and every one-cycle in $\mathcal K$ bounds an integral
two-chain, then every cycle holonomy is one and the determinant line has a
global parallel phase on the configuration graph.
\end{corollary}

\begin{proof}
Apply Theorem~\ref{thm:t10v28-lattice-Stokes} to a filling of each cycle, then
use Theorem~\ref{thm:t10v27-holonomy-test}.
\end{proof}

\begin{proposition}[Accumulation of small plaquette errors]
\label{prop:t10v28-curvature-accumulation}
Choose plaquette arguments $\Omega_p=e^{i\omega_p}$ with
$|\omega_p|<\pi$.  If a cycle bounds
$\Sigma=\sum_pn_pp$ without crossing the branch cut, then
\begin{equation}
 |\arg\operatorname{Hol}(\partial\Sigma)|
 \leq\sum_p|n_p|\,|\omega_p|.
 \label{eq:t10v28-curvature-accumulation}
\end{equation}
\end{proposition}

\begin{proof}
Take the argument of \eqref{eq:t10v28-lattice-Stokes} and apply the triangle
inequality.
\end{proof}

\begin{firewall}[Pointwise small curvature is not global triviality]
\label{fire:t10v28-small-curvature}
The number of plaquettes in a filling can grow with regulator volume and
resolution.  Therefore $\sup_p|\omega_p|\to0$ does not imply uniform cycle
holonomy.  One needs exact cancellation, a summable curvature norm, or a
volume-uniform filling estimate.
\end{firewall}

\subsection{Local measure currents}

Write the plaquette phases additively as a two-cochain
$\omega\in C^2(\mathcal K;\mathbb R/2\pi\mathbb Z)$.  A change of edge
connection by a one-cochain $j$ gives
\begin{equation}
 u'_e=u_ee^{ij_e},
 \qquad
 \omega'=\omega+\delta j.
 \label{eq:t10v28-measure-current}
\end{equation}

\begin{theorem}[Finite integrability criterion]
\label{thm:t10v28-integrability-criterion}
There is a one-cochain $j$ making every corrected plaquette flat if and only
if
\begin{equation}
 [\omega]=0
 \quad\text{in}\quad
 H^2(\mathcal K;\mathbb R/2\pi\mathbb Z).
 \label{eq:t10v28-H2-obstruction}
\end{equation}
When a Hilbert norm is chosen on cochains and
$\delta:C^1\to C^2$ has smallest nonzero singular value $\sigma_\delta$, the
minimum-norm real lift satisfies
\begin{equation}
 \|j_{\min}\|\leq\sigma_\delta^{-1}\|\omega\|
 \label{eq:t10v28-current-bound}
\end{equation}
provided the selected lift has no $2\pi$ ambiguity.
\end{theorem}

\begin{proof}
Equation \eqref{eq:t10v28-measure-current} is flat exactly when
$\delta j=-\omega$, which is solvable precisely when the cohomology class
vanishes.  On a fixed real lift, the Moore--Penrose inverse gives the
minimum-norm solution and the singular-value estimate.
\end{proof}

\begin{firewall}[Existence is not locality]
\label{fire:t10v28-current-locality}
The Moore--Penrose solution can depend on the entire configuration complex,
and $\sigma_\delta$ can vanish along a cofinal family.  A chiral measure needs
a gauge-covariant local current with uniform decay and reflection parity.
Theorem~\ref{thm:t10v28-integrability-criterion} is the algebraic obstruction
test, not that locality theorem.
\end{firewall}

\subsection{Residual flat holonomy}

After solving $\omega+\delta j=0$, the corrected connection is flat.  Its
holonomy is a character of the first homology.

\begin{proposition}[Classification of the residual phase]
\label{prop:t10v28-flat-classification}
Gauge-equivalence classes of flat $U(1)$ edge connections on $\mathcal K$ are
classified by
\begin{equation}
 \operatorname{Hom}(H_1(\mathcal K;\mathbb Z),U(1)).
 \label{eq:t10v28-flat-classification}
\end{equation}
In particular, torsion cycles can carry nontrivial root-of-unity holonomy even
when every local plaquette is flat.
\end{proposition}

\begin{proof}
A flat connection assigns the same phase to homologous cycles by the Stokes
formula, and concatenation multiplies phases, giving a homomorphism from
$H_1$ to $U(1)$.  Conversely, choose a spanning tree and assign phases to
non-tree edges representing a homology basis according to the character;
tree gauge transformations identify choices with the same character.  A
torsion class of order $m$ may map to any $m$-th root of unity.
\end{proof}

The familiar global $SU(2)$ parity anomaly is a $\mathbb Z_2$-valued example.
V17's even number of Standard Model doublets makes the product sign trivial
for that representation loop, provided the finite regulator realizes the
same mod-two transport.

\subsection{Reflection-real cochains}

Assume reflection acts cellularly on $\mathcal K$ and reverses the connection
phase:
\begin{equation}
 u_{\Theta e}=\overline{u_e}.
 \label{eq:t10v28-reflection-edge}
\end{equation}

\begin{lemma}[Reflection of holonomy]
\label{lem:t10v28-reflection-holonomy}
For every cycle $\gamma$,
\begin{equation}
 \operatorname{Hol}(\Theta\gamma)
 =\overline{\operatorname{Hol}(\gamma)}.
 \label{eq:t10v28-reflection-holonomy}
\end{equation}
If a flat character is fixed by reflection and takes value $-1$ on a
reflection-fixed torsion cycle, no globally positive reflection-real phase
can be chosen on that component.
\end{lemma}

\begin{proof}
Multiply \eqref{eq:t10v28-reflection-edge} along the reflected cycle.  A
reflection-real parallel section transported around a fixed loop must return
to itself; holonomy $-1$ contradicts this requirement.
\end{proof}

If $j$ solves the local curvature equation, its averaged cochain
\begin{equation}
 j_{\rm ref}:=\frac12(j-\Theta^*j)
 \label{eq:t10v28-reflection-current}
\end{equation}
has the required odd connection parity whenever
$\Theta^*\omega=-\omega$; it still solves $\delta j_{\rm ref}=-\omega$.
Locality and norm bounds survive only when the original $j$ has them.

\subsection{The local/global determinant-line card}

\begin{definition}[V28 curvature--holonomy card]
\label{def:t10v28-curvature-card}
A finite chiral Standard Model regulator passes the card if:
\begin{enumerate}
\item its constant-rank chiral projectors have curvature given by
\eqref{eq:t10v28-projector-curvature};
\item the product curvature separates into the cancelled universal component
and a finite residual two-cocycle;
\item the residual class in $H^2$ vanishes and has a local, uniformly bounded,
reflection-odd primitive $j$;
\item the corrected flat character in
\eqref{eq:t10v28-flat-classification} is trivial on every free and torsion
cycle; and
\item the resulting phase is compatible with the V27 boundary CAR state and
reflection adjointness.
\end{enumerate}
\end{definition}

\begin{theorem}[Determinant-line completion criterion]
\label{thm:t10v28-line-completion}
If Definition~\ref{def:t10v28-curvature-card} holds, the product determinant
line has a global reflection-real parallel trivialization on the chosen
configuration component.  If the V27 CAR cross matrix is positive, the
fermionic factor therefore satisfies the V26 reflected gluing card.
\end{theorem}

\begin{proof}
The local current flattens the connection by Theorem
\ref{thm:t10v28-integrability-criterion}.  Triviality of the residual
character removes all global holonomy by Proposition
\ref{prop:t10v28-flat-classification}.  The reflection parity constructs a
real section as in Proposition~\ref{prop:t10v26-flat-line}.  The last claim is
Theorem~\ref{thm:t10v27-fermionic-completion}.
\end{proof}

\begin{assessment}[Progress and bottleneck after V28]
\label{assess:t10v28}
V28 computes the finite projector curvature and splits determinant-line
trivialization into an $H^2$ local-current problem and an $H^1$ flat-holonomy
problem.  It proves that Standard Model representation sums cancel the
universal local component but do not settle either finite residual.  The next
task is quantitative locality: construct a local primitive for an exact
residual two-cocycle without using a global Moore--Penrose inverse.
\end{assessment}

\begin{programme}[Tenth-series V29: local cochain homotopy]
\label{prog:t10v28-V29}
V29 will construct a Poincar\'e cochain homotopy on star-shaped configuration
patches, prove support and weighted norm bounds, and derive the transition
one-cocycles on patch overlaps.  The overlap cocycle will identify the exact
remaining obstruction to gluing local chiral measure currents globally.
\end{programme}

\begin{firewall}[Claim boundary after tenth-series V28]
\label{fire:t10v28-boundary}
The local universal anomaly cancels, but no local primitive or global trivial
character has yet been constructed for a concrete curved lattice regulator.
This version does not prove the V27 CAR positivity test, the ultraviolet fixed
point, cofinal measure compactness, or the full SM--Einstein continuum.
\end{firewall}

\paragraph{Parent-version record.}
V28 was formed from
\texttt{Renewal\_SM\_Einstein\_Continuum\_Research\_Notes\_V27.tex}.
The V27 parent had $10\,793$ validator-counted source lines, $410\,970$ bytes,
and SHA--256
\[
 \texttt{405C452A0B101BBE200BA297F89ED3580D65D2700AE248510B048DA3CF5D5DDB}.
\]

% =============================================================================
% END APPENDED TENTH-SERIES V28 MATERIAL
% =============================================================================

% =============================================================================
% APPENDED TENTH-SERIES V29 MATERIAL
% =============================================================================

\section{Tenth-series V29: local cochain homotopy and chiral patch gluing}
\label{sec:v10v29}

V28 identified the finite residual determinant curvature as a closed local
two-cochain.  A global pseudoinverse of the coboundary would destroy locality.
This version constructs the primitive inside star-shaped configuration
patches, proves that its spacetime polymer support is preserved up to one
declared reserve loss, and derives the overlap cocycle that remains to be
glued globally.

\subsection{The radial Poincar\'e homotopy}

Let $U\subset\mathbb R^N$ be star-shaped about $0$, and let
$R=x^a\partial_{x^a}$ be the radial vector field.  For a smooth $k$-form
$\alpha$, $k\geq1$, define
\begin{equation}
 (H_k\alpha)_x(v_1,\ldots,v_{k-1})
 :=\int_0^1t^{k-1}
 \alpha_{tx}(x,v_1,\ldots,v_{k-1})dt.
 \label{eq:t10v29-radial-homotopy}
\end{equation}

\begin{theorem}[Poincar\'e homotopy identity]
\label{thm:t10v29-Poincare}
For $k\geq1$,
\begin{equation}
 dH_k+H_{k+1}d=I
 \quad\text{on }\Omega^k(U).
 \label{eq:t10v29-homotopy-identity}
\end{equation}
Consequently every closed two-form $\omega$ has the explicit primitive
\begin{equation}
 j:=-H_2\omega,
 \qquad dj=-\omega.
 \label{eq:t10v29-local-current}
\end{equation}
\end{theorem}

\begin{proof}
Let $F_t(x)=tx$.  Cartan's formula gives
\[
 \frac d{dt}F_t^*\alpha
 =F_t^*(\mathcal L_{R/t}\alpha)
 =d\bigl(F_t^*\iota_{R/t}\alpha\bigr)
  +F_t^*\iota_{R/t}d\alpha.
\]
Integrating from $0$ to $1$ yields
$\alpha-F_0^*\alpha=dH_k\alpha+H_{k+1}d\alpha$.  For a positive-degree form,
$F_0^*\alpha=0$, proving \eqref{eq:t10v29-homotopy-identity}.  If $d\omega=0$,
the displayed definition of $j$ gives $dj=-\omega$.
\end{proof}

For a patch centred at $q_0$, translate the formula and use the segment
$q_0+t(q-q_0)$.

\begin{lemma}[Basic norm bound]
\label{lem:t10v29-homotopy-bound}
If $U\subset\{x:\|x\|\leq r\}$ and
$\sup_U\|\omega\|\leq M$, then
\begin{equation}
 \|H_2\omega\|_{L^\infty(U)}\leq\frac r2M.
 \label{eq:t10v29-homotopy-bound}
\end{equation}
If the first $m$ derivatives of $\omega$ are bounded, differentiating
\eqref{eq:t10v29-radial-homotopy} gives finite bounds for the first $m$
derivatives of $H_2\omega$, with at most one factor of $r$ and the adjacent
derivative stocks of $\omega$.
\end{lemma}

\begin{proof}
The integrand norm is at most $t\,rM$, and $\int_0^1 t\,dt=1/2$.  Higher
derivatives follow from the product and chain rules under the finite integral.
\end{proof}

\subsection{The simplicial cone identity}

The same construction has a finite cochain form.  If a simplicial complex is
a cone with apex $v_0$, let $C$ send an oriented simplex $\sigma$ not
containing $v_0$ to the oriented cone $[v_0,\sigma]$ and send a simplex already
containing $v_0$ to zero, with the standard orientation sign.

\begin{lemma}[Cone contraction]
\label{lem:t10v29-cone-contraction}
On reduced chains,
\begin{equation}
 \partial C+C\partial=I.
 \label{eq:t10v29-cone-chain}
\end{equation}
The dual cochain operator $h=C^*$ satisfies
\begin{equation}
 \delta h+h\delta=I.
 \label{eq:t10v29-cone-cochain}
\end{equation}
Thus a closed two-cochain on a configuration cone has the primitive
$j=-h\omega$.
\end{lemma}

\begin{proof}
The boundary of a cone is
$\partial[v_0,\sigma]=\sigma-[v_0,\partial\sigma]$ with the chosen
orientation.  This is \eqref{eq:t10v29-cone-chain}.  Dualize and apply the
identity to a closed cochain.
\end{proof}

\begin{firewall}[A distant apex can destroy locality]
\label{fire:t10v29-cone-locality}
The algebraic cone primitive is local in configuration-cell distance only if
the apex and all coning simplices remain in a uniformly bounded patch.  A
single apex for an expanding configuration complex is another global
pseudoinverse and is not an admissible locality proof.
\end{firewall}

\subsection{Preservation of spacetime support}

Let a configuration-space curvature decompose into spacetime polymer owners,
\begin{equation}
 \omega=\sum_X\omega_X,
 \label{eq:t10v29-curvature-polymer}
\end{equation}
where $\omega_X$ depends only on field coordinates in the finite set $X$.
Assume the star-shaped path scales each local coordinate without mixing
distinct spacetime sites.

\begin{proposition}[Support-preserving primitive]
\label{prop:t10v29-support-preserving}
For every $X$, $H_2\omega_X$ depends only on the coordinates in $X$.  Hence
\begin{equation}
 j=-\sum_XH_2\omega_X
 \label{eq:t10v29-polymer-current}
\end{equation}
has no larger spacetime support than the curvature decomposition.
\end{proposition}

\begin{proof}
In \eqref{eq:t10v29-radial-homotopy}, the value and tangent arguments of
$\omega_X$ use only coordinates indexed by $X$.  Scaling the full
configuration restricts on $X$ to scaling those same coordinates, so no
outside coordinate appears.
\end{proof}

Let the strong curvature norm be
\begin{equation}
 \|\omega\|_{\kappa}
 :=\sup_b\sum_{X\ni b}e^{\kappa|X|}
 \|\omega_X\|,
 \label{eq:t10v29-curvature-norm}
\end{equation}
where the patch radius on $X$ satisfies $r_X\leq r_0|X|$.

\begin{theorem}[Polymer reserve cost of the homotopy]
\label{thm:t10v29-polymer-homotopy}
For $0\leq\kappa'<\kappa$,
\begin{equation}
 \|j\|_{\kappa'}
 \leq\frac{r_0}{2}
 \sup_{n\geq1}ne^{-(\kappa-\kappa')n}
 \|\omega\|_{\kappa}.
 \label{eq:t10v29-polymer-homotopy}
\end{equation}
The constant is finite and the spent size reserve can be restored by the V15
occupancy cycle when its mixed-cell geometry gate holds.
\end{theorem}

\begin{proof}
Lemma~\ref{lem:t10v29-homotopy-bound} gives
$\|H_2\omega_X\|\leq r_X\|\omega_X\|/2$.  Insert
$r_X\leq r_0|X|$ and
\[
 e^{\kappa'|X|}|X|
 \leq e^{\kappa|X|}
 \sup_{n\geq1}ne^{-(\kappa-\kappa')n}.
\]
Sum over owners.  The V15 theorem restores a fixed positive reserve gap after
reblocking.
\end{proof}

\begin{firewall}[Growing field radii]
\label{fire:t10v29-growing-radius}
If the normalized good-field radius grows with scale, the factor $r_0$ in
\eqref{eq:t10v29-polymer-homotopy} grows as well.  The audited YM distinction
shows that this can still be harmless when the curvature carries a compensating
positive coupling power.  V29 does not identify the normalized field radius
with the chiral projector's analytic chart radius; their product must be
estimated in the concrete regulator.
\end{firewall}

\subsection{Equivariance and reflection}

\begin{lemma}[Naturality of radial homotopy]
\label{lem:t10v29-naturality}
Let $T$ be a linear transformation preserving the star-shaped patch and its
centre.  Then
\begin{equation}
 T^*H_k=H_kT^*.
 \label{eq:t10v29-naturality}
\end{equation}
In particular, gauge symmetries acting linearly in a fixed slice and the V7
reflection parities are preserved by the local primitive.
\end{lemma}

\begin{proof}
$T(tx)=tT(x)$ and $T_*R=R$.  Substitute in
\eqref{eq:t10v29-radial-homotopy}.
\end{proof}

\begin{corollary}[Reflection-odd measure current]
\label{cor:t10v29-reflection-current}
If $\Theta^*\omega=-\omega$ on a reflection-invariant patch centred at a
fixed point, then $j=-H_2\omega$ satisfies
\begin{equation}
 \Theta^*j=-j.
 \label{eq:t10v29-reflection-current}
\end{equation}
\end{corollary}

\begin{proof}
Apply Lemma~\ref{lem:t10v29-naturality}.
\end{proof}

\begin{firewall}[Nonlinear gauge orbits]
\label{fire:t10v29-gauge-slice}
Radial scaling of a gauge potential chart need not remain on a nonlinear gauge
orbit or inside an admissible compact-link sector.  The construction is made
on a local physical slice/quotient chart.  Equivariance across slice changes
is part of the overlap problem below.
\end{firewall}

\subsection{Overlap one-cocycles}

Let $\{U_i\}$ be a good cover of the admissible configuration component, with
each $U_i$ and every nonempty finite intersection contractible.  On $U_i$ let
$j_i$ satisfy $dj_i=-\omega$.

\begin{proposition}[Transition phases]
\label{prop:t10v29-transition-phases}
On every double overlap,
\begin{equation}
 d(j_i-j_j)=0,
 \qquad
 j_i-j_j=df_{ij}
 \label{eq:t10v29-overlap-potential}
\end{equation}
for a real function $f_{ij}$, unique up to a constant.  On a triple overlap,
\begin{equation}
 f_{ij}+f_{jk}+f_{ki}=2\pi n_{ijk},
 \qquad n_{ijk}\in\mathbb Z
 \label{eq:t10v29-Cech-cocycle}
\end{equation}
after exponentiating the transition phases.  The integers $n_{ijk}$ form a
\v Cech two-cocycle representing the determinant line.
\end{proposition}

\begin{proof}
The first identity follows because the two local currents have the same
curvature.  Contractibility of the overlap and Theorem
\ref{thm:t10v29-Poincare} give $f_{ij}$.  The differential of the triple sum
vanishes, so it is constant.  Consistency of the exponentiated transition
functions $e^{if_{ij}}$ makes that constant an integral multiple of $2\pi$.
The quadruple-overlap identity is the \v Cech cocycle condition.
\end{proof}

\begin{theorem}[Global current gluing criterion]
\label{thm:t10v29-global-gluing}
There is a global one-form $j$ with $dj=-\omega$ if and only if the overlap
data can be modified by local exact changes so that $j_i=j_j$ on every
overlap.  Equivalently, the curvature class in de Rham/$H^2$ cohomology must
vanish.  A global determinant phase further requires the residual flat
$H^1$ character of V28 to be trivial.
\end{theorem}

\begin{proof}
A global current restricts to compatible local currents.  Conversely,
compatible local currents glue by the sheaf property.  Replacing
$j_i$ by $j_i-dg_i$ changes $f_{ij}$ to $f_{ij}-g_i+g_j$; the obstruction to
choosing $g_i$ that kills the overlap differences is the corresponding
cohomology class.  Once the curvature is globally flattened, Proposition
\ref{prop:t10v28-flat-classification} supplies the independent flat-character
condition for a phase section.
\end{proof}

\begin{firewall}[Local exactness does not settle patch gluing]
\label{fire:t10v29-local-exactness}
Every closed curvature form is exact on every star-shaped patch, regardless
of the global determinant-line topology.  Therefore constructing the radial
primitive on all patches does not prove the global anomaly is absent.  The
transition cocycle \eqref{eq:t10v29-Cech-cocycle} and the flat $H^1$ character
must both be checked.
\end{firewall}

\subsection{A local-current card}

\begin{definition}[V29 local homotopy card]
\label{def:t10v29-homotopy-card}
A finite chiral regulator passes the card if:
\begin{enumerate}
\item its admissible configuration component has a reflection-compatible
good cover by uniformly bounded physical slice charts;
\item the residual projector curvature has a local polymer decomposition with
the strong reserve in \eqref{eq:t10v29-curvature-norm};
\item radial or simplicial homotopies preserve admissibility, gauge covariance,
and the required reflection parity;
\item the resulting currents pass the bound
\eqref{eq:t10v29-polymer-homotopy} and reblocking restores the spent reserve;
\item the overlap \v Cech class vanishes through local uniformly bounded
transition corrections; and
\item the remaining flat character is trivial and the V27 CAR positivity row
holds.
\end{enumerate}
\end{definition}

\begin{theorem}[Local construction conditional on global gluing]
\label{thm:t10v29-local-construction}
Rows one through four of Definition~\ref{def:t10v29-homotopy-card} construct
a local reflection-compatible determinant measure current with a uniform
polymer bound.  Rows five and six are exactly the additional global data
needed to obtain the V28 determinant-line completion and the V27/V26
fermionic reflected factorization.
\end{theorem}

\begin{proof}
Theorem~\ref{thm:t10v29-Poincare} constructs each local current,
Theorem~\ref{thm:t10v29-polymer-homotopy} supplies its locality norm, and
Corollary~\ref{cor:t10v29-reflection-current} supplies reflection parity.
Theorem~\ref{thm:t10v29-global-gluing} and the V28 flat-character theorem give
the global phase; V27 supplies boundary CAR positivity.
\end{proof}

\begin{assessment}[Progress and bottleneck after V29]
\label{assess:t10v29}
V29 constructs the local measure current that V28 left abstract and proves
that it preserves spacetime support while spending only one polymer-size
reserve.  The remaining chiral problem is now concentrated in overlap data:
the transition \v Cech cocycle and flat character of the actual admissible
configuration space.  These depend on the chosen regulator and its excluded
zero-mode surfaces.
\end{assessment}

\begin{programme}[Tenth-series V30: audit and admissible-sector topology]
\label{prog:t10v29-V30}
V30 will perform the compulsory YM/NS audit, then analyze how removing chiral
zero-mode/discriminant loci changes the first two homology groups of a finite
configuration component.  It will prove an excision/linking criterion for
possible phase holonomy and identify which cycles can be tested locally and
which require a genuine global regulator theorem.
\end{programme}

\begin{firewall}[Claim boundary after tenth-series V29]
\label{fire:t10v29-boundary}
The local current is constructed only under a uniform good-cover and local
curvature decomposition.  The overlap class, flat character, CAR positivity,
ultraviolet fixed point, cofinal measure compactness, and full SM--Einstein
continuum remain unproved.
\end{firewall}

\paragraph{Parent-version record.}
V29 was formed from
\texttt{Renewal\_SM\_Einstein\_Continuum\_Research\_Notes\_V28.tex}.
The V28 parent had $11\,187$ validator-counted source lines, $426\,137$ bytes,
and SHA--256
\[
 \texttt{4CB416315F6D79C279FA993EA498A27B17269FDC68A4BC5DCCC42FFDC4BE3B76}.
\]

% =============================================================================
% END APPENDED TENTH-SERIES V29 MATERIAL
% =============================================================================

% =============================================================================
% BEGIN APPENDED TENTH-SERIES V30 MATERIAL
% =============================================================================

\section{Tenth-series V30: the discriminant complement, linking charge,
and the gap-contraction gate}
\label{sec:v10v30}

\subsection{Compulsory companion audit}

The five-version rule is applied before choosing the next step.  The newest
completed companion sources visible in \texttt{latex\_notes} were read at the
following exact checkpoints:
\begin{equation}
\begin{array}{c|r|r|l}
 \text{source}&\text{bytes}&\text{lines}&\text{SHA--256}\\ \hline
 \text{YM V41}&654609&18762&
 \texttt{2FE78F2F4DB4197B9D5617A03453B2755658FED9D0D890E643E0A0F42E02EDC5}\\
 \text{NS V26}&278283&5319&
 \texttt{82C887F8C2C69E4F28FAD7C3DC8DE5B9099CB5A4BF431EBA1FFF3E91935978AD}
\end{array}
\label{eq:t10v30-audit-record}
\end{equation}
The corresponding filenames are
\texttt{Renewal\_Yang\_Mills\_Mass\_Gap\_Research\_Notes\_V41.tex}
and
\texttt{Renewal\_NS\_Stability\_Research\_Notes\_V26.tex}.

YM V41 corrects an important quotient ledger.  The parity-tree regulator has
literal product Haar measure and no defining ghost variable.  It proves, in
finite dimensions, that a rooted tree slice has unit coordinate Jacobian and
that its positive Gaussian quotient determinant equals the usual
vector-minus-ghost expression only after the orbit Jacobian is displayed.
It then factors the reduced/coarse split by an exact Schur determinant and
records the coercivity and local-row hypotheses needed for a volume-free
background response.  NS V26 instead computes exact rank-one angular norms
for the first two Oseen-kernel derivatives; its improvement is confined to
one second-derivative channel.

\begin{proposition}[Audit decision]
\label{prop:t10v30-audit-decision}
The YM quotient identity is imported as a guard on the present determinant
programme: every vector/ghost magnitude used here must be derived from the
actual finite regulator quotient, including its orbit Jacobian.  A unit
rooted-tree Jacobian may then remove that debit.  The identity does not
trivialize a chiral determinant line or determine its phase.  The NS rank-one
calculation changes no estimate in the present topology block and is not
imported.
\end{proposition}

\begin{proof}
The YM theorem concerns a real positive Hessian on a gauge quotient.  Its
square-root determinant and Schur complement control magnitudes.  A chiral
fermion determinant is a complex section whose transition phases and zero
set are additional data, so the positive quotient identity cannot settle
them.  The NS tensors and their angular rank-one norm do not occur in the
finite chiral family below.
\end{proof}

The audit changes the order of attack.  Before trying to glue the local
current constructed in V29, one must identify the deleted zero-mode locus:
a local primitive on a regular patch cannot see loops linking a determinant
zero outside that patch.

\subsection{The finite-regulator discriminant}

Let $Q$ be a connected finite-dimensional smooth parameter manifold and let
\begin{equation}
 D:Q\longrightarrow {\rm Mat}_N(\mathbb C),\qquad
 f(q):=\det D(q),\qquad Z:=f^{-1}(0).
 \label{eq:t10v30-discriminant}
\end{equation}
All statements in this section concern a fixed finite regulator.  Infinite
volume and regulator removal are not hidden in the notation.

\begin{proposition}[Logarithmic derivative and winding]
\label{prop:t10v30-log-winding}
On $Q\setminus Z$ the one-form
\begin{equation}
 a:=\frac{1}{2\pi i}{\rm Tr}(D^{-1}dD)
   =\frac{1}{2\pi i}\frac{df}{f}
 \label{eq:t10v30-log-form}
\end{equation}
is closed.  For every piecewise $C^1$ closed loop $\gamma$ in
$Q\setminus Z$,
\begin{equation}
 \int_\gamma a={\rm wind}(f\circ\gamma,0)\in\mathbb Z.
 \label{eq:t10v30-winding}
\end{equation}
The integer is invariant under homotopies of $\gamma$ through
$Q\setminus Z$.
\end{proposition}

\begin{proof}
Jacobi's formula gives $df=f\,{\rm Tr}(D^{-1}dD)$.  Since $f$ is a scalar,
$d(df/f)=0$.  Parametrize the loop and choose a continuous argument of
$f\circ\gamma$ on successive subintervals.  Its total increment is an
integer multiple of $2\pi$, which proves \eqref{eq:t10v30-winding}.
Closedness and Stokes' theorem give homotopy invariance.
\end{proof}

This is an obstruction to a single-valued logarithm of the determinant.  It
is not, by itself, a nontrivial $U(1)$ holonomy: exponentiating an integral
$2\pi i k$ gives one.  Torsion and genuinely flat determinant-line
characters remain the separate V27--V29 obstruction.

\begin{theorem}[Boundary winding equals intersection with the zero divisor]
\label{thm:t10v30-intersection}
Let $\Sigma$ be a compact oriented smooth surface with
$\partial\Sigma=\gamma$, let $f$ be $C^1$ near $\Sigma$, and suppose that
$f$ is nonzero on $\gamma$ and transverse to $0$ in the interior.  Then
$f^{-1}(0)\cap\Sigma$ is finite and
\begin{equation}
 {\rm wind}(f\circ\gamma,0)
 =\sum_{p\in f^{-1}(0)\cap\Sigma}{\rm sign}\det_{\mathbb R}(df_p).
 \label{eq:t10v30-intersection}
\end{equation}
If $\Sigma$ is a complex curve and $f$ is holomorphic, the right side is the
sum of the positive zero multiplicities.
\end{theorem}

\begin{proof}
Transversality makes every inverse image point isolated; compactness makes
the set finite.  Remove mutually disjoint positively oriented discs
$B_p$ about the zeros.  Proposition~\ref{prop:t10v30-log-winding} and Stokes'
theorem on the punctured surface give
\[
 \int_\gamma\frac{df}{2\pi i f}
 =\sum_p\int_{\partial B_p}\frac{df}{2\pi i f}.
\]
For a sufficiently small disc, the last integral is the local topological
degree of $f$ at $p$, equal to the displayed real Jacobian sign at a regular
zero.  In the holomorphic case the local model is a nonzero unit times
$(z-p)^{m_p}$, so the local degree is $m_p>0$.
\end{proof}

\subsection{Why codimension two is the dangerous case}

\begin{lemma}[Regular determinant zeros]
\label{lem:t10v30-regular-zero}
At a complex matrix $A$ of corank one,
\begin{equation}
 d(\det)_A(B)={\rm Tr}({\rm adj}(A)B)
 \label{eq:t10v30-det-derivative}
\end{equation}
is a nonzero complex-linear functional.  Hence the corank-one part of
$\{\det=0\}$ is a smooth complex hypersurface, of real codimension two.  If
$D$ is transverse to this hypersurface, $Z$ is locally a real-codimension-two
submanifold of $Q$.
\end{lemma}

\begin{proof}
For corank one, ${\rm adj}(A)$ has rank one and is nonzero.  Choosing $B$
with a nonzero pairing in \eqref{eq:t10v30-det-derivative} proves
surjectivity onto $\mathbb C$.  The real implicit-function theorem gives the
codimension statement, and transversality pulls it back to $Q$.
\end{proof}

\begin{corollary}[Local linking generator]
\label{cor:t10v30-local-link}
Near a transverse corank-one zero there are real coordinates
$(x_1,x_2,y)$ and a nonvanishing function $u$ such that
$f=u(x_1+ix_2)$.  A small positively oriented normal circle has winding one.
Thus deleting the zero locus creates a local $\mathbb Z$ linking generator.
\end{corollary}

\begin{proof}
Use $({\rm Re},f,{\rm Im},f)$ as the two normal coordinates and absorb the
orientation and nonzero linear factor into $u$.  The factor $u$ has zero
winding on a sufficiently small disc, whereas $x_1+ix_2$ winds once.
\end{proof}

For comparison, a generic zero of a real determinant or real Pfaffian is
codimension one.  It separates sign chambers, and a path changes sign only
by crossing the zero set.  The complex chiral problem instead permits loops
that remain gapped while linking a real-codimension-two zero.

\subsection{Multiplets and the limit of local anomaly cancellation}

\begin{proposition}[Additivity of divisor charge]
\label{prop:t10v30-multiplet-additivity}
Let $f_s$ be nonzero along a closed loop and let
$F=\prod_s f_s^{\nu_s}$ with $\nu_s\in\mathbb Z$.  Then
\begin{equation}
 {\rm wind}(F\circ\gamma,0)
 =\sum_s\nu_s\,{\rm wind}(f_s\circ\gamma,0).
 \label{eq:t10v30-multiplet-winding}
\end{equation}
If the loop bounds a transverse surface, the same identity is the signed sum
of its intersections with the individual zero divisors.
\end{proposition}

\begin{proof}
On the loop, $dF/F=\sum_s\nu_s,df_s/f_s$.  Integrate and use
Proposition~\ref{prop:t10v30-log-winding}; then apply
Theorem~\ref{thm:t10v30-intersection}.
\end{proof}

The Standard Model anomaly sums recorded in V17 cancel the universal local
continuum curvature polynomial.  Proposition~\ref{prop:t10v30-multiplet-additivity}
shows what this does not prove: zero-divisor charges cancel only after their
actual finite-regulator intersections are paired with multiplicity.  The
closed one-form $dz/z$ on $\mathbb C^*$ already has zero curvature and unit
winding.  Therefore vanishing local curvature cannot be substituted for a
calculation on the discriminant complement.

\begin{proposition}[Change of quotient coordinates]
\label{prop:t10v30-coordinate-change}
Suppose on a common parameter patch
\begin{equation}
 \widetilde D(q)=L(q)D(q)R(q),
 \qquad L(q),R(q)\in {\rm GL}_N(\mathbb C).
 \label{eq:t10v30-coordinate-change}
\end{equation}
Then $D$ and $\widetilde D$ have the same zero set and
\begin{equation}
 {\rm wind}(\det\widetilde D\circ\gamma,0)
 ={
m wind}(\det D\circ\gamma,0)
 +{
m wind}(\det L\circ\gamma,0)
 +{
m wind}(\det R\circ\gamma,0).
 \label{eq:t10v30-coordinate-winding}
\end{equation}
In particular, unimodular or contractibly normalized quotient coordinates
preserve every divisor charge.
\end{proposition}

\begin{proof}
Take determinants in \eqref{eq:t10v30-coordinate-change}.  The invertibility
of $L$ and $R$ proves equality of zero sets, and
Proposition~\ref{prop:t10v30-multiplet-additivity} proves the winding formula.
\end{proof}

This is the precise point imported from YM V41.  Its rooted parity-tree map
has determinant one, so it introduces no additional divisor or winding.  A
formal ghost rewriting without the quotient identity would not justify that
conclusion.

\subsection{The gap-contraction theorem}

A spectral gap excludes $Z$ pointwise.  It does not determine the topology of
its complement.

\begin{theorem}[Gapped contraction removes linking]
\label{thm:t10v30-gap-contraction}
Let $S\subset Q$ and suppose there is a continuous contraction
$C:[0,1]\times S\to Q$ to $q_0$ such that
\begin{equation}
 \sigma_{\min}(D(C(t,q)))\geq m>0
 \quad\text{for all }(t,q)\in[0,1]\times S.
 \label{eq:t10v30-gap-contraction}
\end{equation}
Then every loop in $S$ has zero determinant winding.  Moreover $f|_S$ has a
continuous logarithm.  If $S$ is star-shaped and the straight radial
contraction obeys \eqref{eq:t10v30-gap-contraction}, this logarithm is fixed
uniquely after choosing its value at the centre.
\end{theorem}

\begin{proof}
The bound implies $C([0,1]\times S)\cap Z=\varnothing$.  Composing any loop
with $C$ gives a homotopy in $Q\setminus Z$ to the constant loop, so its
winding vanishes by Proposition~\ref{prop:t10v30-log-winding}.  The induced
map $f_*:\pi_1(S)\to\pi_1(\mathbb C^*)=\mathbb Z$ is therefore zero.  The
covering-space lifting criterion lifts $f:S\to\mathbb C^*$ through
$\exp:\mathbb C\to\mathbb C^*$, producing a continuous logarithm.  The
chosen central value fixes the lift on connected $S$.
\end{proof}

\begin{firewall}[A gap without a contraction is insufficient]
\label{fire:t10v30-gap-alone}
On the annulus $S=\{1/2<|z|<2\}$ take the one-by-one family $D(z)=z$.
Then $\sigma_{\min}(D)>1/2$, but the central circle has winding one and no
single-valued logarithm exists.  A uniform admissibility gap must therefore
be accompanied by a gapped contraction or by an explicit computation of the
remaining linking generators.
\end{firewall}

\begin{proposition}[Reflection pairs divisor charge]
\label{prop:t10v30-reflection-pairing}
Let $\Theta:Q\to Q$ be an involution satisfying
\begin{equation}
 f(\Theta q)=\overline{f(q)}.
 \label{eq:t10v30-reflection-real}
\end{equation}
For the reflected loop $(\Theta\gamma)(t)=\Theta(\gamma(t))$ with the same
parameter orientation,
\begin{equation}
 {\rm wind}(f\circ\Theta\gamma,0)
 =-{\rm wind}(f\circ\gamma,0).
 \label{eq:t10v30-reflection-winding}
\end{equation}
Consequently a reflection-stable family of loops has cancelling paired
charges, while a reflection-fixed loop can carry only zero integer winding.
\end{proposition}

\begin{proof}
Complex conjugation reverses the orientation of $\mathbb C^*$ and hence
negates winding.  If a loop is fixed as a parametrized loop by reflection,
its winding equals its negative.
\end{proof}

The last sentence does not eliminate a possible $\mathbb Z_2$ real-line
orientation obstruction.  Such a sign is torsion and is invisible to the
integer de Rham integral, exactly as anticipated by the flat-holonomy ledger
of V27--V29.

\subsection{Finite-regulator acquisition card}

\begin{theorem}[What is now reduced to finite checks]
\label{thm:t10v30-finite-card}
On a chosen finite Standard Model regulator sector, absence of a divisor
obstruction follows if all of the following are proved:
\begin{enumerate}
\item the determinant is that of the actual gauge quotient, with every orbit
      Jacobian included;
\item a uniform singular-value gap holds on the entire sector and on a stated
      contraction to a reference configuration;
\item the multiplet is reflection-real along that contraction; and
\item after the local curvature current of V29 is applied, the residual
      overlap cocycle and flat $H^1$ character, including torsion, vanish.
\end{enumerate}
Items 1--3 eliminate the zero-divisor and logarithm obstruction.  Item 4 is
independent and is required to trivialize the determinant line itself.
\end{theorem}

\begin{proof}
Item 1 makes Proposition~\ref{prop:t10v30-coordinate-change} applicable to
the defined regulator.  Item 2 invokes
Theorem~\ref{thm:t10v30-gap-contraction}; item 3 invokes
Proposition~\ref{prop:t10v30-reflection-pairing}.  The local primitive removes
the exact curvature component, but the V29 Čech criterion shows that a
global section still requires the stated overlap and flat-character checks.
\end{proof}

\begin{assessment}[V30 programme]
\label{ass:t10v30-programme}
The next upstream object is the gapped projector bundle.  V31 will construct
Kato transport on a star-shaped gapped sector, prove quantitative projection
and transport bounds, and identify the exact equivariance cocycle left by a
choice of transported frame.  This tests whether the contraction required by
Theorem~\ref{thm:t10v30-gap-contraction} can also preserve the locality and
gauge structure required by V29.
\end{assessment}

\begin{firewall}[Claim boundary after V30]
\label{fire:t10v30-claim}
V30 proves finite-dimensional topology and a sufficient gap-contraction
gate.  It does not prove that a Standard Model lattice sector obeys that gate,
that its chiral determinant line is globally trivial, or that the coupled
SM--Einstein continuum limit exists.
\end{firewall}

\subsection{Parent record}

This V30 note is appended to the validated V29 source, whose record is
\begin{equation}
 \text{bytes}=439999,qquad \text{lines}=11558,qquad
 \operatorname{SHA256}=
 \texttt{C575D32F110FC49A0EC5C8B8F579998C0EDC76519344144205167DF144D9016B}.
 \label{eq:t10v30-parent-record}
\end{equation}

% =============================================================================
% END APPENDED TENTH-SERIES V30 MATERIAL
% =============================================================================

% =============================================================================
% BEGIN APPENDED TENTH-SERIES V31 MATERIAL
% =============================================================================

\section{Tenth-series V31: gapped Kato transport, quasi-local frames,
and the equivariance cocycle}
\label{sec:v10v31}

\subsection{Purpose and finite setting}

V30 separated two assertions that are often conflated.  A pointwise spectral
gap removes determinant zeros, whereas a gapped contraction removes their
linking.  This note constructs the corresponding transported chiral frame
and measures its spatial tails.  The construction is exact at a fixed finite
regulator; its hypotheses must later be made uniform in volume and cutoff.

Let $\mathcal H=\bigoplus_{x\in\Lambda}\mathcal H_x$ be a finite-dimensional
lattice Hilbert space.  Let $S$ be star-shaped about $q_0$, write
$q_t=q_0+t(q-q_0)$, and let $H(q)=H(q)^*$ satisfy
\begin{equation}
 \operatorname{spec}H(q)subset[-M,-g]\cup[g,M],
 \qquad q\in S,qquad 0<g<M.
 \label{eq:t10v31-gap}
\end{equation}
Choose a fixed positively oriented contour $\Gamma$ enclosing $[-M,-g]$ and
no point of $[g,M]$, and put
\begin{equation}
 P(q)=\frac{1}{2\pi i}\int_\Gamma(z-H(q))^{-1}\,dz.
 \label{eq:t10v31-riesz}
\end{equation}
This is the occupied or negative-spectral projector.  The fixed contour may
be chosen at distance $\delta>0$ from the two spectral intervals.

\subsection{Projector estimates}

\begin{lemma}[Riesz derivative]
\label{lem:t10v31-riesz-derivative}
If $H$ is $C^1$, then $P$ is $C^1$ and, for a tangent vector $v$,
\begin{equation}
 dP[v]=\frac{1}{2\pi i}\int_\Gamma
 (z-H)^{-1}dH[v](z-H)^{-1}\,dz.
 \label{eq:t10v31-riesz-derivative}
\end{equation}
Consequently
\begin{equation}
 \|dP[v]\|
 \leq \frac{|\Gamma|}{2\pi\delta^2}\|dH[v]\|.
 \label{eq:t10v31-projector-bound}
\end{equation}
The rank of $P$ is constant on $S$.
\end{lemma}

\begin{proof}
Differentiate $(z-H)^{-1}$ by the inverse rule.  The resolvent norm on
$\Gamma$ is at most $\delta^{-1}$, which proves the bound.  The trace of a
continuous projection is integer-valued and therefore constant on connected
$S$.
\end{proof}

\begin{remark}[Higher derivatives]
\label{rem:t10v31-higher-riesz}
Repeated differentiation of \eqref{eq:t10v31-riesz} gives a sum over ordered
partitions of the derivative labels.  Every term is an alternating product
of resolvents and derivatives of $H$.  Thus analytic or Gevrey derivative
stocks for $H$, together with the fixed gap, give the corresponding stocks
for $P$ with one additional resolvent per differentiated block.  This is the
finite projector version of the V20 resolvent ledger.
\end{remark}

\subsection{Exact Kato trivialization}

Along the radial path define
\begin{equation}
 K_t=[\dot P_t,P_t],
 \qquad \dot U_t=K_tU_t,qquad U_0=I.
 \label{eq:t10v31-kato-ode}
\end{equation}

\begin{theorem}[Kato transport]
\label{thm:t10v31-kato}
Equation \eqref{eq:t10v31-kato-ode} has a unique unitary solution and
\begin{equation}
 P_tU_t=U_tP_0.
 \label{eq:t10v31-intertwining}
\end{equation}
If $V_0:\mathbb C^r\to\mathcal H$ is an isometry onto
$\operatorname{ran}P_0$, then
\begin{equation}
 V(q):=U_1(q)V_0
 \label{eq:t10v31-global-frame}
\end{equation}
is a global continuous orthonormal frame of $\operatorname{ran}P(q)$ on $S$.
If $P$ is $C^k$ or analytic, the frame has the same regularity.
\end{theorem}

\begin{proof}
Since $P_t=P_t^*$, the commutator $K_t$ is anti-Hermitian.  Therefore
$d(U_t^*U_t)/dt=0$, proving unitarity.  Differentiating $P_t^2=P_t$ gives
$P_t\dot P_tP_t=0$ and
$\dot P_t=P_t\dot P_t+\dot P_tP_t$.  Hence
\[
 [P_t,K_t]=[P_t,[\dot P_t,P_t]]=-\dot P_t.
\]
It follows that
\[
 \frac d{dt}(U_t^*P_tU_t)
 =U_t^*(\dot P_t+[P_t,K_t])U_t=0,
\]
which is \eqref{eq:t10v31-intertwining}.  Thus $V(q)$ is an isometry onto the
range of $P(q)$.  Standard parameter dependence for a finite-dimensional
linear ODE proves the last assertion.
\end{proof}

\begin{corollary}[Determinant bundle on the contracted sector]
\label{cor:t10v31-determinant-trivial}
The determinant line $\det(\operatorname{ran}P)\to S$ has the nowhere-zero
section
\begin{equation}
 s(q)=V_1(q)\wedge\cdots\wedge V_r(q).
 \label{eq:t10v31-det-section}
\end{equation}
Thus every topological obstruction on the unquotiented star-shaped gapped
sector vanishes.
\end{corollary}

\begin{proof}
The wedge of an orthonormal frame is nonzero at every point.  It is a global
section, which trivializes a complex line bundle.
\end{proof}

The adjective ``unquotiented'' is essential.  Gauge transformations identify
points of $S$ and impose an equivariance law on $s$; that law is treated
below.

\subsection{Weighted spatial locality}

For $\mu\geq0$ define the weighted row norm
\begin{equation}
 \|A\|_\mu:=\sup_{x\in\Lambda}
 \sum_{y\in\Lambda}e^{\mu d(x,y)}\|A_{xy}\|.
 \label{eq:t10v31-weighted-row}
\end{equation}
It obeys $\|AB\|_\mu\leq\|A\|_\mu\|B\|_\mu$ by the triangle inequality for
$d$.

\begin{theorem}[Quasi-local Kato transport]
\label{thm:t10v31-local-kato}
Assume, uniformly for $t\in[0,1]$ and $z\in\Gamma$,
\begin{equation}
 \|(z-H_t)^{-1}\|_\mu\leq C_R,qquad
 \|\dot H_t\|_\mu\leq C_H.
 \label{eq:t10v31-weighted-hypotheses}
\end{equation}
Set
\begin{equation}
 C_P:=\frac{|\Gamma|}{2\pi}C_R,qquad
 C_{\dot P}:=\frac{|\Gamma|}{2\pi}C_R^2C_H,qquad
 C_K:=2C_PC_{\dot P}.
 \label{eq:t10v31-local-constants}
\end{equation}
Then
\begin{equation}
 \|P_t\|_\mu\leq C_P,quad
 \|\dot P_t\|_\mu\leq C_{\dot P},quad
 \|K_t\|_\mu\leq C_K,quad
 \|U_t\|_\mu\leq e^{tC_K}.
 \label{eq:t10v31-local-bounds}
\end{equation}
In particular
\begin{equation}
 \|(U_t)_{xy}\|\leq e^{tC_K}e^{-\mu d(x,y)},
 \label{eq:t10v31-local-tail}
\end{equation}
so the transported frame has exponential spatial tails with constants
independent of volume whenever the hypotheses do.
\end{theorem}

\begin{proof}
Apply the weighted Banach-algebra norm under the contour integral in
\eqref{eq:t10v31-riesz} and
\eqref{eq:t10v31-riesz-derivative}.  The commutator estimate gives
$\|K_t\|_\mu\leq2\|P_t\|_\mu\|\dot P_t\|_\mu$.  Picard iteration of
\eqref{eq:t10v31-kato-ode} in the same Banach algebra gives
\[
 \|U_t\|_\mu
 \leq1+\sum_{n\geq1}\frac1{n!}
 \left(\int_0^t\|K_s\|_\mu ds\right)^n
 \leq e^{tC_K}.
\]
A single matrix block is bounded by its weighted row sum, which proves
\eqref{eq:t10v31-local-tail}.
\end{proof}

The weighted resolvent hypothesis is supplied by the V20
Combes--Thomas estimate when the family remains finite range and the gap is
uniform.  The factor $e^{C_K}$ displays the cost of the entire contraction;
a path whose weighted derivative length diverges with volume does not pass
this gate.

\begin{firewall}[A global frame need not be uniformly local]
\label{fire:t10v31-frame-locality}
Corollary~\ref{cor:t10v31-determinant-trivial} uses only contractibility and a
gap.  Theorem~\ref{thm:t10v31-local-kato} additionally needs a uniform
weighted resolvent and finite weighted path length.  Contracting all links at
once can satisfy the topological theorem while violating the analytic one.
\end{firewall}

\subsection{Curvature correction on the contracted sector}

Let
\begin{equation}
 A:=\operatorname{Tr}(V^*dV)
 \label{eq:t10v31-berry-connection}
\end{equation}
be the determinant-frame connection.  Its curvature is
$F=dA=\operatorname{Tr}(P\,dP\wedge dP)$.  Suppose the V29 local current $j$
has been constructed with $dj=-F$ on $S$.

\begin{proposition}[Radial counterphase]
\label{prop:t10v31-radial-counterphase}
The one-form $\alpha:=A+j$ is closed.  There is a unique scalar $\Phi$ with
$\Phi(q_0)=0$ and
\begin{equation}
 d\Phi=\alpha,qquad
 \Phi(q)=\int_0^1\alpha_{q_t}(q-q_0)\,dt.
 \label{eq:t10v31-counterphase}
\end{equation}
Consequently the corrected determinant frame has zero connection on the
unquotiented sector.  If $A$ and $j$ obey uniform polymer bounds, the radial
homotopy estimate of V29 gives the corresponding bound for $\Phi$.
\end{proposition}

\begin{proof}
The curvature identity and $dj=-F$ give $d\alpha=0$.  The degree-one radial
Poincare identity on the star-shaped set gives
$dH_1\alpha=\alpha-H_0d\alpha=\alpha$, and $H_1\alpha$ is exactly the
integral in \eqref{eq:t10v31-counterphase}.  The central normalization gives
uniqueness.
\end{proof}

\subsection{The gauge-equivariance cocycle}

Let a group $G$ act on $S$.  Suppose $R_g(q)$ is a unitary action on
$\mathcal H$ satisfying
\begin{equation}
 P(gq)=R_g(q)P(q)R_g(q)^{-1},qquad
 R_{gh}(q)=R_g(hq)R_h(q).
 \label{eq:t10v31-equivariant-projector}
\end{equation}
For the Kato frame define
\begin{equation}
 c(g,q):=V(gq)^*R_g(q)V(q)\in U(r),qquad
 \chi(g,q):=\det c(g,q)\in U(1).
 \label{eq:t10v31-cocycle}
\end{equation}

\begin{theorem}[Residual determinant cocycle]
\label{thm:t10v31-equivariance-cocycle}
The matrices and their determinants obey
\begin{equation}
 c(gh,q)=c(g,hq)c(h,q),qquad
 \chi(gh,q)=\chi(g,hq)\chi(h,q).
 \label{eq:t10v31-cocycle-law}
\end{equation}
Under a frame change $V(q)\mapsto V(q)b(q)$ with $b(q)\in U(r)$,
\begin{equation}
 \chi(g,q)\mapsto
 \det b(gq)^{-1}\,\chi(g,q)\,\det b(q).
 \label{eq:t10v31-coboundary}
\end{equation}
Hence a gauge-equivariant determinant section exists exactly when $\chi$ is
a coboundary in the stated coefficient system.
\end{theorem}

\begin{proof}
Insert $V(hq)V(hq)^*=P(hq)$ between $R_g(hq)$ and $R_h(q)$.  It acts as the
identity on $R_h(q)V(q)$, by
\eqref{eq:t10v31-equivariant-projector}.  This proves the first identity;
taking determinants proves the second.  Direct substitution of $Vb$ gives
\eqref{eq:t10v31-coboundary}.  If the cocycle is a coboundary, the
corresponding $b$ makes $\chi=1$ and the determinant wedge equivariant.
Conversely an equivariant wedge supplies such a frame phase, so the original
cocycle is a coboundary.
\end{proof}

\begin{firewall}[Contractibility before quotienting is not enough]
\label{fire:t10v31-quotient}
The sector $S$ may be star-shaped and its determinant bundle trivial while
$S/G$ carries a nontrivial determinant line.  The obstruction is exactly the
class of $\chi$ in Theorem~\ref{thm:t10v31-equivariance-cocycle}.  Kato
transport therefore converts the global anomaly question into a concrete
finite cocycle calculation; it does not answer that calculation.
\end{firewall}

\subsection{V31 acquisition card}

\begin{assessment}[What V31 changes]
\label{ass:t10v31-status}
The V30 gap-contraction gate now has an explicit frame.  Under weighted
resolvent control, the frame is exponentially quasi-local, and after the V29
curvature current its residual phase is radially integrable.  The sole
topological datum left on the quotient is the gauge cocycle $\chi$, together
with its possible reflection-real sign.  V32 will compute this cocycle on
loops of gauge transformations by spectral flow and separate perturbative
winding from the mod-two orientation obstruction.
\end{assessment}

\begin{firewall}[Claim boundary after V31]
\label{fire:t10v31-claim}
V31 proves a conditional finite-regulator trivialization and its locality
bound.  It does not establish a uniform Standard Model gap, a local gapped
contraction, triviality of $\chi$, or any continuum limit.
\end{firewall}

\subsection{Parent record}

This V31 note is appended to the validated V30 source, whose record is
\begin{equation}
 \text{bytes}=455462,qquad \text{lines}=11926,qquad
 \operatorname{SHA256}=
 \texttt{7EE9293930A4247E72DAA3AE9A93FC9017A7FDDC20D5C59A86A1304770B6B88A}.
 \label{eq:t10v31-parent-record}
\end{equation}

% =============================================================================
% END APPENDED TENTH-SERIES V31 MATERIAL
% =============================================================================

% =============================================================================
% BEGIN APPENDED TENTH-SERIES V32 MATERIAL
% =============================================================================

\section{Tenth-series V32: spectral-flow orientation, the weak
$\mathbb Z_2$ test, and the Standard Model count}
\label{sec:v10v32}

\subsection{Two components of the quotient cocycle}

The $U(1)$ cocycle $\chi$ isolated in V31 has a component visible to
infinitesimal gauge variations and a component supported on large gauge
transformations.  The first is the finite-regulator representative of the
perturbative anomaly.  The second can be torsion and is invisible to every
curvature integral.  This note treats the familiar four-dimensional
$SU(2)$ sign and performs the Standard Model representation count.

\begin{lemma}[Identity-component reduction]
\label{lem:t10v32-identity-component}
Let $G_0$ be the identity component of a Lie group acting on a connected
sector, and let $\chi(g,q)$ be a differentiable $U(1)$ cocycle normalized by
$\chi(1,q)=1$.  If
\begin{equation}
 \left.\frac d{dt}\right|_{t=0}
 \log\chi(\exp(tX),q)=0
 \label{eq:t10v32-infinitesimal-zero}
\end{equation}
for every Lie-algebra element $X$ and every $q$, then
$\chi(g,q)=1$ for every $g\in G_0$.
\end{lemma}

\begin{proof}
Join $1$ to $g$ by a piecewise smooth path $g_t$ in $G_0$.  Subdivide it so
that each increment is an exponential.  The cocycle law writes the
logarithmic derivative of $\chi(g_t,q)$ as the infinitesimal derivative
\eqref{eq:t10v32-infinitesimal-zero} evaluated at the transformed base point.
It vanishes on every segment.  The initial value is one, so the final value
is one.
\end{proof}

\begin{firewall}[Continuum anomaly cancellation is not yet the hypothesis]
\label{fire:t10v32-finite-infinitesimal}
V17 verifies the Standard Model representation sums multiplying the
universal continuum local anomaly polynomial.  Lemma~\ref{lem:t10v32-identity-component}
requires the exact finite-regulator identity
\eqref{eq:t10v32-infinitesimal-zero}, including the measure current and
quotient Jacobian.  Establishing that identity remains a regulator task.
\end{firewall}

\subsection{A finite spectral-flow orientation lemma}

Let $A_t$, $0\leq t\leq1$, be a $C^1$ path of real self-adjoint Fredholm
matrices with invertible endpoints.  After an arbitrarily small perturbation
fixing the endpoints, assume every zero crossing is simple and transverse.
Let $\operatorname{sf}_2(A_t)$ be the number of crossings modulo two.

\begin{lemma}[Orientation transport]
\label{lem:t10v32-orientation-transport}
The orientation line transported along $A_t$ changes by
\begin{equation}
 \operatorname{hol}_{\mathbb R}(A_t)
 =(-1)^{\operatorname{sf}_2(A_t)}.
 \label{eq:t10v32-orientation-holonomy}
\end{equation}
The parity is unchanged under endpoint-fixed homotopies for which the
endpoints remain invertible.
\end{lemma}

\begin{proof}
Near a simple crossing $t_*$, orthogonal spectral projection splits off one
real eigenline and $A_t$ is homotopic, without changing the other invertible
blocks, to $\operatorname{diag}(t-t_*,B)$.  Passing through $t_*$ reverses
the determinant-line orientation once.  Multiplication over crossings gives
\eqref{eq:t10v32-orientation-holonomy}.  During a generic two-parameter
endpoint-fixed homotopy, crossings are created or annihilated in pairs;
hence their parity is invariant.  Nongeneric families follow by a small
perturbation and the same homotopy argument.
\end{proof}

For a pseudoreal chiral representation, the relevant real line is the
Pfaffian orientation line.  A path whose endpoints are related by a gauge
transformation closes only after using that gauge transformation to identify
the endpoint fibres.  The resulting sign is the mod-two spectral-flow
cocycle.

\begin{proposition}[Quotient-loop sign]
\label{prop:t10v32-quotient-sign}
Let $q_t$ be a path from $q$ to $gq$ and let the endpoint operator be
identified with the initial operator by the gauge action.  Then the real
part of the determinant cocycle is
\begin{equation}
 \chi_{\mathbb R}(g,q)
 =(-1)^{\operatorname{sf}_2(A(q_t))}\,\varepsilon_g,
 \label{eq:t10v32-quotient-sign}
\end{equation}
where $\varepsilon_g$ is the orientation of the endpoint identification.
This sign is invariant under homotopies of the quotient loop that stay in
the invertible sector, and it is multiplicative under concatenation.
\end{proposition}

\begin{proof}
Lemma~\ref{lem:t10v32-orientation-transport} gives the transported
orientation.  Closing the line multiplies it by the endpoint-identification
orientation $\varepsilon_g$.  Homotopy invariance follows from the lemma,
and concatenating paths adds spectral flow modulo two and multiplies endpoint
identifications.
\end{proof}

In an infinite-dimensional continuum formulation the same parity is the
mod-two index of the corresponding five-dimensional mapping-torus operator.
The finite statement above is the regulator form needed for the present
programme.

\subsection{The weak $SU(2)$ generator input}

The following is a named external topological input, rather than a result
reproved by the renewal estimates.

\begin{theorem}[Witten generator theorem]
\label{thm:t10v32-witten-input}
For a four-dimensional left-handed Weyl fermion in the fundamental
representation of $SU(2)$, the nontrivial class of
$\pi_4(SU(2))\cong\mathbb Z_2$ has odd mod-two spectral flow.  Its Pfaffian
orientation therefore changes sign.  For $N$ independent fundamental
Weyl doublets the sign is $(-1)^N$.
\end{theorem}

The original source is E.~Witten, ``An $SU(2)$ anomaly,''
\emph{Physics Letters B} 117 (1982), 324--328,
doi:10.1016/0370-2693(82)90728-6.  The final multiplicity statement also
follows directly from the next lemma.

\begin{lemma}[Direct-sum parity]
\label{lem:t10v32-direct-sum}
For two real chiral families,
\begin{equation}
 \operatorname{sf}_2(A^{(1)}\oplus A^{(2)})
 =\operatorname{sf}_2(A^{(1)})+
  \operatorname{sf}_2(A^{(2)})\pmod2.
 \label{eq:t10v32-direct-sum}
\end{equation}
Hence their orientation signs multiply.
\end{lemma}

\begin{proof}
After a generic perturbation separating simultaneous crossings, the crossing
set of the direct sum is the disjoint union of the two crossing sets.  Count
modulo two and apply Lemma~\ref{lem:t10v32-orientation-transport}.
\end{proof}

\subsection{Standard Model doublet count}

Use only left-handed Weyl fields.  In one generation their weak
representations are
\begin{equation}
 Q_L:(\mathbf3,\mathbf2),\qquad
 L_L:(\mathbf1,\mathbf2),
 \label{eq:t10v32-left-doublets}
\end{equation}
while the charge-conjugates of the right-handed quarks and charged lepton
are weak singlets.  A right-handed neutrino, if added, is also a weak singlet.

\begin{theorem}[Cancellation of the familiar weak mod-two sign]
\label{thm:t10v32-sm-witten-count}
One Standard Model generation contains
\begin{equation}
 N_{\mathbf2}=3+1=4
 \label{eq:t10v32-one-generation-count}
\end{equation}
left-handed $SU(2)$ fundamental doublets, where the three copies of $Q_L$
are its colour multiplicity.  With three generations,
$N_{\mathbf2}=12$.  Therefore the Witten-generator sign is
\begin{equation}
 (-1)^{12}=+1.
 \label{eq:t10v32-sm-sign}
\end{equation}
It already cancels generation by generation.
\end{theorem}

\begin{proof}
Restricting $(\mathbf3,\mathbf2)$ to the weak factor gives three identical
fundamental doublets, and $(\mathbf1,\mathbf2)$ gives one.  All other Weyl
fermions are weak singlets and contribute zero to the direct-sum parity.
Lemma~\ref{lem:t10v32-direct-sum} and
Theorem~\ref{thm:t10v32-witten-input} give
$(-1)^4=+1$ per generation and \eqref{eq:t10v32-sm-sign} for three.
Scalars do not define a chiral Pfaffian orientation and contribute no sign.
\end{proof}

\begin{corollary}[Stability under singlet neutrinos]
\label{cor:t10v32-neutrino-singlet}
Adding any number of gauge-singlet right-handed neutrinos leaves
\eqref{eq:t10v32-sm-sign} unchanged.
\end{corollary}

\begin{proof}
Their restriction to $SU(2)$ is trivial, so their family is constant along a
pure weak gauge loop and has zero mod-two spectral flow.
\end{proof}

\subsection{What the count does and does not settle}

\begin{theorem}[Representation-level quotient ledger]
\label{thm:t10v32-representation-ledger}
Assume a finite chiral regulator has the following exact properties:
\begin{enumerate}
\item its infinitesimal cocycle equals the local anomaly polynomial with the
      V17 Standard Model coefficients;
\item its weak large-gauge sign agrees with the Witten-generator mod-two
      index; and
\item its quotient coordinates add no determinant winding, as required by
      V30 and the YM V41 audit.
\end{enumerate}
Then its cocycle is trivial on the identity component of the gauge group and
on the familiar $\pi_4(SU(2))$ generator.
\end{theorem}

\begin{proof}
The V17 sums make item 1 vanish; Lemma~\ref{lem:t10v32-identity-component}
then treats the identity component.  Theorem~\ref{thm:t10v32-sm-witten-count}
treats item 2.  Item 3 prevents a coordinate Jacobian from reintroducing a
phase.
\end{proof}

\begin{firewall}[This is not a classification of all global anomalies]
\label{fire:t10v32-global-classification}
Theorem~\ref{thm:t10v32-sm-witten-count} checks the familiar weak
$\pi_4(SU(2))$ sign on spin four-manifolds.  It does not classify torsion
anomalies of the full global gauge group, possible quotients such as a
central quotient of $SU(3)\times SU(2)\times U(1)$, mixed gauge--gravity
bordism classes, or the regulator cocycle on those classes.  Those questions
require a global-structure and spin-bordism ledger.
\end{firewall}

\subsection{V32 acquisition card}

\begin{assessment}[Next upstream step]
\label{ass:t10v32-status}
The best-known real sign obstruction cancels for the Standard Model matter
content, but the proof exposes a new exact input: the global form of the
gauge group and the tangential structure of spacetime.  V33 will construct
the associated mapping-torus/bordism ledger, prove which tests are invariant
under local counterterms, and state a finite list of global classes that a
regulator must evaluate.  This goes upstream of any attempt to declare the
V31 cocycle trivial.
\end{assessment}

\begin{firewall}[Claim boundary after V32]
\label{fire:t10v32-claim}
V32 proves the Standard Model multiplicity cancellation conditional only on
the named Witten generator theorem.  It does not construct a finite chiral
measure with that sign, classify the full Standard Model anomaly, prove a
uniform gap, or construct the SM--Einstein continuum.
\end{firewall}

\subsection{Parent record}

This V32 note is appended to the validated V31 source, whose record is
\begin{equation}
 \text{bytes}=467150,qquad \text{lines}=12254,qquad
 \operatorname{SHA256}=
 \texttt{7D56F8DF1F8AF16671C698B7328FBB9EA5271E60BD0CC58E3F9FF1ECEAE16A2B}.
 \label{eq:t10v32-parent-record}
\end{equation}

% =============================================================================
% END APPENDED TENTH-SERIES V32 MATERIAL
% =============================================================================

% =============================================================================
% BEGIN APPENDED TENTH-SERIES V33 MATERIAL
% =============================================================================

\section{Tenth-series V33: the global gauge group, mapping tori,
and the spin-bordism anomaly ledger}
\label{sec:v10v33}

\subsection{The four global forms}

The Lie algebra $\mathfrak{su}(3)\oplus\mathfrak{su}(2)\oplus\mathfrak u(1)$
does not determine the gauge group.  Put
\begin{equation}
 \widetilde G=SU(3)\times SU(2)\times U(1),
 \qquad
 \xi=(e^{2\pi i/3}I_3,-I_2,e^{i\pi/3}),
 \label{eq:t10v33-central-generator}
\end{equation}
and, for $n\in\{1,2,3,6\}$, let
\begin{equation}
 G_n:=\widetilde G/\Gamma_n,
 \qquad \Gamma_n:=\langle\xi^{6/n}\rangle\cong\mathbb Z/n.
 \label{eq:t10v33-global-groups}
\end{equation}

Use the integer hypercharge $q=6Y$.  For one generation, written entirely
with left-handed Weyl fields, the charge ledger is
\begin{equation}
\begin{array}{c|ccccc}
 &Q_L&u_R^c&d_R^c&L_L&e_R^c\\ \hline
 q&1&-4&2&-3&6\\
 SU(3)&\mathbf3&\overline{\mathbf3}&\overline{\mathbf3}&\mathbf1&\mathbf1\\
 SU(2)&\mathbf2&\mathbf1&\mathbf1&\mathbf2&\mathbf1
\end{array}.
\label{eq:t10v33-charge-table}
\end{equation}

\begin{lemma}[Descent of the Standard Model representations]
\label{lem:t10v33-representation-descent}
Every representation in \eqref{eq:t10v33-charge-table}, as well as the Higgs
$(\mathbf1,\mathbf2)_{q=3}$, is trivial on $\xi$.  Hence it is a bona fide
representation of every $G_n$ in \eqref{eq:t10v33-global-groups}.
\end{lemma}

\begin{proof}
If $t\in\mathbb Z/3$ is colour triality and $s\in\mathbb Z/2$ is weak
isospin parity, $\xi$ acts by
\begin{equation}
 \exp(2\pi it/3)(-1)^s\exp(i\pi q/3).
 \label{eq:t10v33-central-action}
\end{equation}
For $Q_L$ the product is
$e^{2\pi i/3}(-1)e^{i\pi/3}=1$.  For $u_R^c$ and $d_R^c$ it is respectively
$e^{-2\pi i/3}e^{-4\pi i/3}=1$ and
$e^{-2\pi i/3}e^{2\pi i/3}=1$.  For $L_L$ it is
$(-1)e^{-i\pi}=1$, and for $e_R^c$ it is $e^{2\pi i}=1$.
The Higgs gives $(-1)e^{i\pi}=1$.  Triviality on $\xi$ implies triviality on
every subgroup $\Gamma_n$.
\end{proof}

Thus the global form is genuine input even though the observed multiplets
descend to all four possibilities.

\subsection{Mapping tori and the anomaly character}

Let $M$ be a closed spin four-manifold, let $P\to M$ be a principal $G$
bundle with connection, and let $g$ be a gauge transformation.  Gluing the
two ends of $M\times[0,1]$ using $g$ gives a five-dimensional mapping torus
$Y_g$ with its induced spin structure and $G$-bundle.

\begin{proposition}[Bordism target for a flat anomaly]
\label{prop:t10v33-bordism-target}
After the perturbative anomaly polynomial vanishes, the exponentiated chiral
phase on closed five-dimensional tests defines a character
\begin{equation}
 \alpha_R:\Omega^{\rm Spin}_5(BG)\longrightarrow U(1).
 \label{eq:t10v33-anomaly-character}
\end{equation}
The determinant-line holonomy around the quotient loop generated by $g$ is
$\alpha_R([Y_g])$.  Direct sums of Weyl representations multiply these
characters.
\end{proposition}

\begin{proof}
This is the Dai--Freed mapping-torus theorem: when the local index density
cancels, the exponentiated reduced eta invariant is unchanged under spin
$G$-bordism and is multiplicative under disjoint union.  These are exactly
the relations defining a homomorphism
\eqref{eq:t10v33-anomaly-character}.  Gluing the fermion cylinder by $g$ is
the mapping-torus construction of determinant-line holonomy.  Eta invariants
add under direct sums, so their exponentials multiply.
\end{proof}

The proposition is a named external index-theoretic input.  It supplies the
correct finite list of topological tests; it does not construct the local
lattice measure.

\begin{corollary}[Zero bordism group]
\label{cor:t10v33-zero-bordism}
If $\Omega^{\rm Spin}_5(BG)=0$, every flat global anomaly character is
trivial.
\end{corollary}

\begin{proof}
The zero group has only the trivial homomorphism to $U(1)$.
\end{proof}

\subsection{The Standard Model bordism table}

The required low-dimensional topology has been computed by the
Atiyah--Hirzebruch spectral sequence.  We isolate it as a second external
input so that no spectral-sequence calculation is silently attributed to the
renewal argument.

\begin{theorem}[Spin-bordism table for the four Standard Model groups]
\label{thm:t10v33-bordism-table}
For the groups \eqref{eq:t10v33-global-groups},
\begin{equation}
\begin{array}{c|cccc}
 n&1&2&3&6\\ \hline
 \Omega^{\rm Spin}_5(BG_n)&\mathbb Z/2&0&\mathbb Z/2&0.
\end{array}
\label{eq:t10v33-bordism-table}
\end{equation}
For $n=1$ and $n=3$, the nonzero class is the familiar weak $SU(2)$ Witten
class; there is no additional independent global anomaly class.
\end{theorem}

A source with the full spectral-sequence proof is J.~Davighi, B.~Gripaios and
N.~Lohitsiri, ``Global anomalies in the Standard Model(s) and Beyond,''
\emph{JHEP} 07 (2020) 232, arXiv:1910.11277,
doi:10.1007/JHEP07(2020)232.  In particular,
$G_2\cong SU(3)\times U(2)$ and
$G_3\cong U(3)\times SU(2)$, which makes the difference between the zero and
$\mathbb Z/2$ entries visible.

\begin{theorem}[Global representation-level cancellation]
\label{thm:t10v33-sm-global-cancellation}
For the Standard Model fermion content on spin four-manifolds, the character
\eqref{eq:t10v33-anomaly-character} is trivial for each of the four groups
$G_n$, provided the perturbative anomaly polynomial has first been cancelled.
\end{theorem}

\begin{proof}
For $n=2,6$, apply Corollary~\ref{cor:t10v33-zero-bordism} to
Theorem~\ref{thm:t10v33-bordism-table}.  For $n=1,3$, the only possible
character is determined by the Witten $\mathbb Z/2$ generator.  V32 counted
four weak doublets per generation and proved that their product sign is
$+1$.  Hence the generator also maps to one.  A homomorphism from
$\mathbb Z/2$ that maps its generator to one is trivial.
\end{proof}

\begin{remark}[Interplay for the quotient groups]
\label{rem:t10v33-interplay}
For $G_2$ and $G_6$, an isolated hypercharge-zero weak doublet is not a
representation of the quotient group.  Weak isospin and hypercharge obey a
descent relation such as \eqref{eq:t10v33-central-action}.  This is why one
must cancel the local hypercharge anomalies before reading the zero bordism
entry; separating a formal weak sign from an inadmissible representation
would give a false test.
\end{remark}

\subsection{Pure and mixed gravitational placement}

\begin{lemma}[Ordinary spin gravity in degree five]
\label{lem:t10v33-spin-five}
The low spin-bordism group of a point satisfies
\begin{equation}
 \Omega^{\rm Spin}_5({\rm pt})=0.
 \label{eq:t10v33-spin-five}
\end{equation}
Thus there is no independent pure torsion character in degree five for an
ordinary spin four-dimensional fermion theory.  Mixed gauge--gravity global
classes, when present, are already included in
$\Omega^{\rm Spin}_5(BG)$.
\end{lemma}

\begin{proof}
Equation \eqref{eq:t10v33-spin-five} is the standard low-degree spin-bordism
calculation.  Forgetting the $G$-bundle maps a mixed test to the point group;
the full bundle-dependent information is retained only in $BG$, so it is
classified by the group in Theorem~\ref{thm:t10v33-bordism-table}.
\end{proof}

The local mixed $U(1)$--gravity anomaly is not torsion and is instead part of
the V17 polynomial ledger.  Its coefficient is the sum of integer
hypercharges.  In one generation,
\begin{equation}
 6(1)+3(-4)+3(2)+2(-3)+6=0,
 \label{eq:t10v33-mixed-gravity-sum}
\end{equation}
where the factors are representation dimensions.  This is the local input
needed before Proposition~\ref{prop:t10v33-bordism-target} applies.

\subsection{Consequences for the regulator programme}

\begin{theorem}[Topological target for a finite chiral measure]
\label{thm:t10v33-regulator-target}
Fix one $G_n$.  Suppose a finite regulator satisfies:
\begin{enumerate}
\item its multiplet representations descend as in
      Lemma~\ref{lem:t10v33-representation-descent};
\item its infinitesimal anomaly equals the continuum local polynomial and is
      cancelled by the Standard Model sums;
\item its quotient-loop cocycle converges, on each mapping-torus test, to the
      Dai--Freed character; and
\item the gapped local trivializations and overlap currents obey the V29--V31
      locality bounds uniformly.
\end{enumerate}
Then no representation-level global anomaly obstructs gluing its local
fermion measures.  The remaining problem is analytic: construct the finite
measure and prove the four uniform statements.
\end{theorem}

\begin{proof}
Items 1--3 identify the regulator cocycle with the character in
Theorem~\ref{thm:t10v33-sm-global-cancellation}, which is trivial.  Item 4
turns this topological triviality into a local gluing candidate rather than a
nonlocal choice of phase.  No topological class remains; all unproved inputs
are the stated uniform analytic estimates.
\end{proof}

\begin{firewall}[Dynamical gravity is still an additional construction]
\label{fire:t10v33-dynamical-gravity}
The bordism ledger permits curved spin backgrounds and includes mixed global
tests.  It does not define a sum over metrics or topologies, cure the
Euclidean conformal factor, prove Einstein constraint gluing, or supply
reflection positivity.  Those gravitational gates from V9--V25 remain
active.
\end{firewall}

\subsection{V33 acquisition card}

\begin{assessment}[Return from topology to analysis]
\label{ass:t10v33-status}
The representation-level local and global anomaly target is now closed for
all four standard global gauge groups, subject to the named index and bordism
theorems.  The live obstruction is construction of a uniformly local finite
regulator realizing that target.  V34 will build a conservative admissible
link chart with an explicit Hermitian-kernel gap, prove locality of its sign
projector, and test whether the chart contraction remains inside the compact
link manifold.
\end{assessment}

\begin{firewall}[Claim boundary after V33]
\label{fire:t10v33-claim}
V33 proves the consequence of the established spin-bordism classification
for Standard Model representations.  It does not rederive the AHSS table,
construct the regulator, prove its convergence to the Dai--Freed character,
or construct the SM--Einstein continuum.
\end{firewall}

\subsection{Parent record}

This V33 note is appended to the validated V32 source, whose record is
\begin{equation}
 \text{bytes}=478268,qquad \text{lines}=12523,qquad
 \operatorname{SHA256}=
 \texttt{304287AC684C65627830EF6D03C430813440E1AC791CAEFFC0BA0AB901A5DB9C}.
 \label{eq:t10v33-parent-record}
\end{equation}

% =============================================================================
% END APPENDED TENTH-SERIES V33 MATERIAL
% =============================================================================

% =============================================================================
% BEGIN APPENDED TENTH-SERIES V34 MATERIAL
% =============================================================================

\section{Tenth-series V34: an explicit Wilson-kernel gap chart and
exponential locality of the overlap projector}
\label{sec:v10v34}

\subsection{A conservative local regulator sector}

V33 closes the representation-level anomaly ledger but leaves open the
analytic realization.  This note constructs a small, explicit compact-link
chart on which the Hermitian Wilson kernel is uniformly gapped.  The chart is
intentionally conservative.  It proves that the V31 Kato hypotheses are
nonempty and quantitative; it does not cover every admissible topological
sector.

Work on a periodic $d$-dimensional unit lattice with $d=4$ at the end.  Let
the Euclidean gamma matrices be Hermitian, let the Wilson parameter be one,
and define
\begin{align}
 (D_W(U)\psi)(x)
 &:=(d-m_0)\psi(x)\notag\\
 &\quad-\frac12\sum_{\mu=1}^d
 \left[(1-\gamma_\mu)U_{x\mu}\psi(x+\widehat\mu)
 +(1+\gamma_\mu)U^*_{x-\widehat\mu,\mu}
       \psi(x-\widehat\mu)\right],
 \label{eq:t10v34-wilson-operator}\\
 H_W(U)&:=\gamma_{d+1}D_W(U).
 \label{eq:t10v34-hermitian-kernel}
\end{align}
Here $U_{x\mu}$ is the unitary matrix in one fixed fermion representation,
and $0<m_0<2$.  The notation incorporates the usual negative Wilson mass.

\subsection{The exact free gap}

\begin{lemma}[Free Wilson gap]
\label{lem:t10v34-free-gap}
At $U=I$,
\begin{equation}
 H_W(I)^2(p)=
 \sum_{\mu=1}^d\sin^2p_\mu+
 \left(\sum_{\mu=1}^d(1-\cos p_\mu)-m_0\right)^2.
 \label{eq:t10v34-free-square}
\end{equation}
Consequently
\begin{equation}
 \operatorname{dist}(0,\operatorname{spec}H_W(I))
 =g_0:=\min\{m_0,2-m_0\}>0.
 \label{eq:t10v34-free-gap}
\end{equation}
The value is independent of the periodic volume.
\end{lemma}

\begin{proof}
Fourier transformation gives the scalar Wilson term
$W(p)=\sum_\mu(1-\cos p_\mu)$ and the Clifford term
$i\sum_\mu\gamma_\mu\sin p_\mu$.  Gamma anticommutation proves
\eqref{eq:t10v34-free-square}.  Put $a_\mu=1-\cos p_\mu\in[0,2]$.  Its right
side is
\begin{equation}
 m_0^2+2(1-m_0)\sum_\mu a_\mu
 +2\sum_{\mu<\nu}a_\mu a_\nu.
 \label{eq:t10v34-multiaffine}
\end{equation}
This is affine in each $a_\mu$ separately, so its minimum on the hypercube is
attained at a vertex.  If exactly $k$ coordinates equal two, the value is
$(2k-m_0)^2$.  For $0<m_0<2$ the smallest values occur at $k=0,1$, giving
\eqref{eq:t10v34-free-gap}.  Restricting momenta to a finite periodic grid
cannot lower the continuous-torus minimum.
\end{proof}

\subsection{Compact-link perturbation and radial contraction}

Define the representation-link distance
\begin{equation}
 \varepsilon(U):=\sup_{x,\mu}\|U_{x\mu}-I\|.
 \label{eq:t10v34-link-distance}
\end{equation}

\begin{lemma}[Uniform hopping perturbation]
\label{lem:t10v34-hopping-bound}
For every compact-link field,
\begin{equation}
 \|H_W(U)-H_W(I)\|_{\ell^2\to\ell^2}
 \leq2d\,\varepsilon(U).
 \label{eq:t10v34-hopping-bound}
\end{equation}
\end{lemma}

\begin{proof}
The norm of $(1\pm\gamma_\mu)/2$ is one.  Each row of the difference has at
most $2d$ hopping blocks, each of norm at most $\varepsilon(U)$; the same is
true of each column.  The block Schur test gives the bound for $D_W$, and
multiplication by the unitary $\gamma_{d+1}$ preserves it.
\end{proof}

\begin{theorem}[Explicit principal-logarithm gap chart]
\label{thm:t10v34-gap-chart}
Suppose each link has a chosen principal logarithm
\begin{equation}
 U_{x\mu}=e^{A_{x\mu}},\qquad A_{x\mu}^*=-A_{x\mu},qquad
 \sup_{x,\mu}\|A_{x\mu}\|\leq\rho,qquad
 2d\rho<g_0.
 \label{eq:t10v34-log-chart}
\end{equation}
Then the radial compact-link path $U_t=e^{tA}$ satisfies
\begin{equation}
 \operatorname{dist}(0,\operatorname{spec}H_W(U_t))
 \geq g_*:=g_0-2d\rho>0,qquad0\leq t\leq1.
 \label{eq:t10v34-path-gap}
\end{equation}
Thus the chart is star-shaped in logarithmic link coordinates and gives the
gapped contraction required by V30--V31.
\end{theorem}

\begin{proof}
For an anti-Hermitian matrix with spectrum $i\theta_j$ in the principal
interval,
$\|e^{tA}-I\|=\max_j|e^{it\theta_j}-1|\leq t\max_j|\theta_j|leq\rho$.
Lemma~\ref{lem:t10v34-hopping-bound} and the spectral variation bound for
self-adjoint matrices give
\[
 \operatorname{dist}(0,\operatorname{spec}H_W(U_t))
 \geq g_0-\|H_W(U_t)-H_W(I)\|
 \geq g_0-2d\rho.
\]
The exponential path stays in the compact gauge group.
\end{proof}

For a product gauge group, \eqref{eq:t10v34-log-chart} is imposed on the
matrix in each actual fermion representation.  Representation-dependent
conversion from factorwise link radii to $\rho$ is a finite tensor-product
estimate and must be entered in the multiplet ledger.

\begin{corollary}[Small metric perturbations]
\label{cor:t10v34-metric-gap}
Let $H_W(U,h)$ be a Hermitian curved-background discretization satisfying
\begin{equation}
 \|H_W(U,h)-H_W(U,h_0)\|\leq C_h\eta
 \label{eq:t10v34-metric-lipschitz}
\end{equation}
when the tetrad, spin connection, and cell weights lie in a metric chart of
radius $\eta$.  If
\begin{equation}
 2d\rho+C_h\eta<g_0,
 \label{eq:t10v34-joint-gap-condition}
\end{equation}
then the joint gauge--metric radial path has gap at least
$g_0-2d\rho-C_h\eta$.
\end{corollary}

\begin{proof}
Apply the triangle inequality to the perturbation from $(I,h_0)$ and repeat
the self-adjoint spectral variation argument.
\end{proof}

This corollary records the exact missing input on the gravitational side:
$C_h$ must be derived from a specified curved lattice Dirac operator.  It is
not assigned a formal value here.

\subsection{An explicit exponential-locality estimate}

Let $H=H_W(U)$ on the gap chart, let $M_*\geq\|H\|$, and put
\begin{equation}
 q:=1-\frac{g_*^2}{M_*^2}\in[0,1),qquad
 X:=I-\frac{H^2}{M_*^2}.
 \label{eq:t10v34-binomial-data}
\end{equation}
The binomial series gives
\begin{equation}
 \operatorname{sign}(H)
 =\frac{H}{M_*}\sum_{n=0}^\infty c_nX^n,qquad
 c_n=\frac1{4^n}{2n\choose n}\leq1.
 \label{eq:t10v34-sign-series}
\end{equation}

\begin{theorem}[Volume-free overlap-projector tail]
\label{thm:t10v34-overlap-locality}
Assume $H$ has hopping range one.  Let
\begin{equation}
 P_-(U):=\frac12(I-\operatorname{sign}H_W(U)).
 \label{eq:t10v34-overlap-projector}
\end{equation}
For $x\ne y$ and $L=d(x,y)$,
\begin{equation}
 \|(P_-)_{xy}\|
 \leq\frac1{2(1-q)}q^{\lceil(L-1)/2\rceil}.
 \label{eq:t10v34-projector-tail}
\end{equation}
The constants depend only on the gap ratio $g_*/M_*$ and not on the lattice
volume.
\end{theorem}

\begin{proof}
The spectrum of $X$ lies in $[0,q]$, so the remainder after $N$ terms in
\eqref{eq:t10v34-sign-series} has operator norm at most
\begin{equation}
 \frac{\|H\|}{M_*}\sum_{n>N}c_nq^n
 \leq\frac{q^{N+1}}{1-q}.
 \label{eq:t10v34-sign-remainder}
\end{equation}
The polynomial $H\sum_{n=0}^Nc_nX^n$ has range at most $1+2N$, because
$H$ has range one and $X$ range two.  If $1+2N<L$, its $(x,y)$ block
vanishes.  Choose the largest such $N$ and use
$P_-=(I-\operatorname{sign}H)/2$.  The ceiling form in
\eqref{eq:t10v34-projector-tail} is the resulting harmless upper bound.
\end{proof}

\begin{corollary}[Ginsparg--Wilson algebra on the chart]
\label{cor:t10v34-gw-algebra}
With $\epsilon=\operatorname{sign}H_W$ and
$D_{\rm ov}=I+\gamma_{d+1}\epsilon$, one has
\begin{equation}
 \epsilon^2=I,qquad
 \gamma_{d+1}D_{\rm ov}+D_{\rm ov}\gamma_{d+1}
 =D_{\rm ov}\gamma_{d+1}D_{\rm ov}.
 \label{eq:t10v34-gw-relation}
\end{equation}
Moreover $D_{\rm ov}$ is exponentially local with the tail inherited from
\eqref{eq:t10v34-projector-tail}.
\end{corollary}

\begin{proof}
Functional calculus gives $\epsilon^2=I$.  Substitute
$D_{\rm ov}=I+\gamma_{d+1}\epsilon$ and expand both sides; the terms agree
using only $\gamma_{d+1}^2=I$ and $\epsilon^2=I$.  Locality follows by
multiplication with the on-site gamma matrix.
\end{proof}

\subsection{Gauge saturation and the quotient chart}

The logarithm chart about $U=I$ is not invariant under arbitrary lattice
gauge transformations.  Let a rooted tree slice provide a unique local
decomposition
\begin{equation}
 U=g\cdot e^A,qquad \|A\|_\infty<\rho,
 \label{eq:t10v34-slice-decomposition}
\end{equation}
with the root transformation retained.  Define the saturated chart by all
such $g\cdot e^A$.

\begin{proposition}[Quotient contraction]
\label{prop:t10v34-quotient-contraction}
On a patch where \eqref{eq:t10v34-slice-decomposition} is unique, the map
\begin{equation}
 C_t[g\cdot e^A]=[g\cdot e^{tA}]
 \label{eq:t10v34-quotient-contraction}
\end{equation}
is a well-defined contraction in the gauge quotient.  The kernel gap and the
projector-tail bound are unchanged along gauge orbits.
\end{proposition}

\begin{proof}
Uniqueness makes $A$ a function of the quotient point, so the formula is
well-defined.  At $t=0$ the endpoint $g\cdot I$ represents the flat quotient
point.  Gauge covariance gives
$H_W(g\cdot U)=G_gH_W(U)G_g^{-1}$; unitary conjugation preserves spectrum and
block decay in the covariantly transported frame.
\end{proof}

The rooted-tree unit Jacobian from the YM V41 audit ensures that this slice
introduces no separate Gaussian orbit determinant.  A proof that the chosen
chiral link chart and the YM parity tree are the same compatible slice is
still required.

\begin{firewall}[Kernel gap versus physical chiral zeros]
\label{fire:t10v34-kernel-vs-dirac}
The gap in \eqref{eq:t10v34-path-gap} makes
$\operatorname{sign}H_W$ and $P_-$ well-defined.  It does not imply that the
massless overlap operator $D_{\rm ov}$ is invertible.  Physical zero modes
and topological index sectors belong to $D_{\rm ov}$, not to the auxiliary
Wilson-kernel gap.  V30's determinant discriminant must therefore be checked
again after the retained physical zero modes are separated.
\end{firewall}

\subsection{V34 acquisition card}

\begin{assessment}[What is gained and what is too small]
\label{ass:t10v34-status}
The V31 hypotheses now hold on a nonempty finite gauge chart with explicit
volume-free constants, and they extend to a metric chart once
\eqref{eq:t10v34-metric-lipschitz} is proved.  The radius
$\rho<g_0/(2d)$ is a linkwise small-field condition and is far narrower than
a plaquette admissibility sector; it also chooses one local tree slice.
V35 will compare linkwise and plaquette control, prove the unavoidable flat
holonomy obstruction on a torus, and perform the next compulsory YM/NS
audit before choosing the larger admissible cover.
\end{assessment}

\begin{firewall}[Claim boundary after V34]
\label{fire:t10v34-claim}
V34 proves a conservative gap and locality theorem.  It does not cover
nontrivial topology, prove the curved-kernel Lipschitz bound, remove physical
fermion zeros, construct the chiral measure, or take a continuum limit.
\end{firewall}

\subsection{Parent record}

This V34 note is appended to the validated V33 source, whose record is
\begin{equation}
 \text{bytes}=489118,qquad \text{lines}=12789,qquad
 \operatorname{SHA256}=
 \texttt{223F93C1159C39C906932500C136BFACBAD0589370C1ECDB6827FBDC1746DDAE}.
 \label{eq:t10v34-parent-record}
\end{equation}

% =============================================================================
% END APPENDED TENTH-SERIES V34 MATERIAL
% =============================================================================

% =============================================================================
% BEGIN APPENDED TENTH-SERIES V35 MATERIAL
% =============================================================================

\section{Tenth-series V35: compulsory audit, plaquette admissibility,
and the topological-sector cover}
\label{sec:v10v35}

\subsection{Compulsory companion audit}

The newest completed companion files were inspected before selecting this
version's direction.  Their exact checkpoints are
\begin{equation}
\begin{array}{c|r|r|l}
 \text{source}&\text{bytes}&\text{lines}&\text{SHA--256}\\ \hline
 \text{YM V43}&699797&20016&
 \texttt{0430FCA45B706C2DB0B0A871024B19B9FE7DFFF8B59F8471548C06EB80D66CE1}\\
 \text{NS V26}&278283&5319&
 \texttt{82C887F8C2C69E4F28FAD7C3DC8DE5B9099CB5A4BF431EBA1FFF3E91935978AD}
\end{array}
\label{eq:t10v35-audit-record}
\end{equation}
The filenames are
\texttt{Renewal\_Yang\_Mills\_Mass\_Gap\_Research\_Notes\_V43.tex}
and
\texttt{Renewal\_NS\_Stability\_Research\_Notes\_V26.tex}.

YM V43 proves exact measured-slice invariance: if two complements of the
gauge directions are related by $I_1=I_0R+DC$, their Hessians and coordinate
Jacobians obey
$H_1=R^*H_0R$ and $J_1=J_0|\det R|$, so
$\frac12\log\det H_i-\log J_i$ is slice independent.  It also corrects an
earlier coefficient identification by separating the first Wilson--Schur
calibration from the invariant Gaussian reference shape.  NS V26 is
unchanged and remains unrelated to the chiral admissibility problem.

\begin{proposition}[Audit transfer: magnitude versus phase]
\label{prop:t10v35-audit-transfer}
For the complexified complement change $I_1=I_0R+DC$, the measured positive
Gaussian magnitude is invariant, but normalized determinant frames differ by
\begin{equation}
 \tau_{10}=\frac{\det R}{|\det R|}\in U(1).
 \label{eq:t10v35-slice-phase}
\end{equation}
Thus YM V43 closes the gauge-slice debit for vector/ghost magnitudes while
identifying, rather than removing, the chiral overlap cocycle.
\end{proposition}

\begin{proof}
The YM identities give
$(\det H_1)^{1/2}=|\det R|(\det H_0)^{1/2}$ and
$J_1=J_0|\det R|$, so their ratio is invariant.  On the quotient, an ordered
complex frame changes by $R$; dividing its top wedge by its norm removes
$|\det R|$ and leaves \eqref{eq:t10v35-slice-phase}.  On triple overlaps these
phases multiply by determinant multiplicativity, giving the V29 Čech
cocycle.
\end{proof}

This audit changes the V34 programme usefully.  One need not force a single
parity-tree slice on the whole admissible sector.  Local measured slices may
be used, provided their positive Jacobians and the phases
\eqref{eq:t10v35-slice-phase} are both retained.

\subsection{Gauge-invariant plaquette admissibility}

For a unitary fermion representation $R$, let $U_R(p)$ be the ordered product
around a plaquette and define
\begin{equation}
 \mathcal A_\epsilon^R
 :=\left\{U:\sup_p\|I-U_R(p)\|<\epsilon\right\}.
 \label{eq:t10v35-admissible-sector}
\end{equation}
Unlike the link chart of V34, this set is gauge invariant.

\begin{theorem}[Hernandez--Jansen--Luscher admissibility input]
\label{thm:t10v35-hjl-input}
For the standard four-dimensional Neuberger kernel with Wilson mass one, the
plaquette condition
\begin{equation}
 \sup_p\|I-U_R(p)\|<\epsilon<\frac1{30}
 \label{eq:t10v35-hjl-condition}
\end{equation}
implies a positive volume-independent lower bound on the squared Hermitian
Wilson kernel; in the standard normalization one may take
\begin{equation}
 H_W(U)^2\succeq(1-30\epsilon)I.
 \label{eq:t10v35-hjl-gap}
\end{equation}
The associated Neuberger operator and spectral projector are exponentially
local with constants depending on $\epsilon$ but not on volume.
\end{theorem}

This is an external lattice theorem.  A complete proof is P.~Hernandez,
K.~Jansen and M.~Luscher, ``Locality properties of Neuberger's lattice Dirac
operator,'' \emph{Nuclear Physics B} 552 (1999), 363--378,
arXiv:hep-lat/9808010, doi:10.1016/S0550-3213(99)00213-8.  V34 supplied a
self-contained weaker linkwise theorem and the polynomial-tail mechanism;
Theorem~\ref{thm:t10v35-hjl-input} supplies the gauge-invariant replacement.

\begin{corollary}[Uniform Kato constants on the admissible set]
\label{cor:t10v35-admissible-kato}
On every differentiable path contained in
$\mathcal A_\epsilon^R$, the Riesz projector and Kato transport of V31 are
well-defined.  If the weighted path length is bounded, their exponential
locality constants are volume independent.
\end{corollary}

\begin{proof}
The gap \eqref{eq:t10v35-hjl-gap} supplies a common Riesz contour, and the
finite hopping range supplies the weighted resolvent estimate.  Apply the
V31 weighted Banach-algebra theorem and integrate the Kato generator along
the path.
\end{proof}

For the full Standard Model, \eqref{eq:t10v35-hjl-condition} must hold in
every fermion representation, or be deduced from a factorwise admissibility
condition with explicit representation constants.  The smallest resulting
gap is the multiplet gap.

\subsection{Admissibility separates index sectors}

Let
\begin{equation}
 \nu(U):=-\frac12\operatorname{Tr}\operatorname{sign}H_W(U).
 \label{eq:t10v35-overlap-index}
\end{equation}
The conventional sign can be reversed with the convention for $H_W$; only
integrality and constancy are used below.

\begin{lemma}[Index constancy]
\label{lem:t10v35-index-constancy}
The integer $\nu(U)$ is constant on every connected component of
$\mathcal A_\epsilon^R$.
\end{lemma}

\begin{proof}
Theorem~\ref{thm:t10v35-hjl-input} prevents an eigenvalue of $H_W$ from
crossing zero.  Hence the ranks of its positive and negative spectral
projections are continuous integer-valued functions and are constant on a
connected component.  Their trace difference gives
\eqref{eq:t10v35-overlap-index}.
\end{proof}

\begin{corollary}[No global radial contraction across topology]
\label{cor:t10v35-no-global-contraction}
If the admissible set contains configurations with two distinct values of
$\nu$, it cannot admit a contraction to one reference configuration through
$\mathcal A_\epsilon^R$.
\end{corollary}

\begin{proof}
A contraction supplies a path from every configuration to the reference.
Lemma~\ref{lem:t10v35-index-constancy} would then give the same index to all
of them, a contradiction.
\end{proof}

Thus V30's gapped-contraction theorem applies component by component.  The
component label is physical retained data and cannot be discarded as a bad
polymer.

\subsection{Flat holonomy inside an index sector}

Even the zero-index sector need not be one logarithm chart.  For an abelian
factor on a periodic $d$-torus, constant links
\begin{equation}
 U_{x\mu}=e^{i\theta_\mu/L_\mu}
 \label{eq:t10v35-flat-toron}
\end{equation}
have every plaquette equal to one and Wilson lines $e^{i\theta_\mu}$.  Gauge
equivalence identifies $\theta_\mu$ modulo $2\pi$, so the flat quotient
contains the torus $(S^1)^d$.

\begin{proposition}[Plaquette control does not give one quotient chart]
\label{prop:t10v35-toron-obstruction}
For periodic boundary conditions, no single logarithm chart centred at the
identity parametrizes the whole flat quotient.  In particular, plaquette
admissibility alone does not provide the star-shaped quotient sector assumed
in V31.
\end{proposition}

\begin{proof}
The family \eqref{eq:t10v35-flat-toron} embeds $(S^1)^d$ into the flat
quotient.  A single-valued logarithm of all Wilson lines would lift the
identity map of $S^1$ to $\mathbb R$, which is impossible because the
fundamental loop has nonzero winding.  A star-shaped coordinate domain is
contractible, whereas $(S^1)^d$ is not.
\end{proof}

For nonabelian factors the flat locus is the space of commuting holonomies
modulo simultaneous conjugation; Weyl identifications add orbifold strata.
The same conclusion holds: flat holonomy must be retained as a slow boundary
coordinate and covered by patches.

\begin{firewall}[Small curvature is not a small potential uniformly]
\label{fire:t10v35-curvature-potential}
A small plaquette norm controls lattice curvature locally.  Reconstructing a
potential in a fixed gauge uses an inverse difference operator and can lose a
factor comparable to the torus diameter unless harmonic holonomies are
removed first.  Therefore the linkwise radius $\rho$ in V34 cannot be bounded
uniformly from \eqref{eq:t10v35-hjl-condition} without a harmonic router and
a local gauge theorem.
\end{firewall}

\subsection{A sectorwise good cover}

Fix an index $k$ and write
$\mathcal A_{\epsilon,k}^R=\{U\in\mathcal A_\epsilon^R:\nu(U)=k\}$.  Choose
holonomy patches $O_a$ on the retained flat moduli and local measured gauge
slices above them.  On each slice choose a reference $U_a$ and a gapped path
$C_a(t,U)$ that stays in the same admissible component.

\begin{theorem}[Conditional local-measure cover]
\label{thm:t10v35-sector-cover}
Assume the patches have uniformly bounded overlap order and that, on every
$O_a$,
\begin{enumerate}
\item the path $C_a$ has uniformly bounded weighted length;
\item Theorem~\ref{thm:t10v35-hjl-input} holds for all Standard Model
      representations;
\item the V29 curvature current has a uniform polymer norm; and
\item measured-slice transitions retain both the positive Jacobian and the
      phase \eqref{eq:t10v35-slice-phase}.
\end{enumerate}
Then each $O_a$ has a quasi-local Kato determinant frame, and all remaining
global information is the finite-overlap Čech cocycle and its flat holonomy
character.
\end{theorem}

\begin{proof}
Items 1--2 and the V31 theorem give a quasi-local frame on each patch.  Item 3
cancels its local curvature by the V29 primitive.  Item 4 makes the magnitude
independent of slice and identifies the exact normalized phase on overlaps.
The cocycle law follows on triple overlaps, leaving precisely the V29 gluing
and flat-character data.
\end{proof}

\begin{assessment}[Direction after the audit]
\label{ass:t10v35-status}
The global anomaly class is trivial by V33, so the overlap cocycle should be
a coboundary for the Standard Model multiplet.  The analytic task is to find
a coboundary with uniform locality, not merely an abstract global section.
V36 will first separate physical overlap zero modes in each index sector and
construct the reduced determinant/Pfaffian line.  This is necessary before
any local phase can be bounded near massless configurations.
\end{assessment}

\begin{firewall}[Claim boundary after V35]
\label{fire:t10v35-claim}
V35 imports a proved plaquette-gap theorem, proves sector and toron
obstructions to one global chart, and formulates a sectorwise local cover.  It
does not construct the cover with volume-uniform constants, trivialize its
Čech cocycle locally, treat massless physical zero modes, or take a coupled
continuum limit.
\end{firewall}

\subsection{Parent record}

This V35 note is appended to the validated V34 source, whose record is
\begin{equation}
 \text{bytes}=500464,qquad \text{lines}=13100,qquad
 \operatorname{SHA256}=
 \texttt{4FBD340C13DB6218AF968A1D6E9B4739CFD4BAD3B31F2F9512E9CFBBE4B8DD3B}.
 \label{eq:t10v35-parent-record}
\end{equation}

% =============================================================================
% END APPENDED TENTH-SERIES V35 MATERIAL
% =============================================================================

% =============================================================================
% BEGIN APPENDED TENTH-SERIES V36 MATERIAL
% =============================================================================

\section{Tenth-series V36: reduced chiral determinant lines, zero-mode
saturation, and the infrared handoff}
\label{sec:v10v36}

\subsection{Why the auxiliary gap is not the physical gap}

The admissibility theorem in V35 gaps the Hermitian Wilson kernel used to
define the overlap projector.  The massless physical Dirac operator can
still have exact zero modes and arbitrarily small nonzero modes.  Treating
its determinant as an ordinary nonvanishing function would therefore erase
topological sectors and invalidate the large-volume estimates.  The correct
object is the determinant line with its kernel and cokernel displayed.

Let $D(q):E_q\to F_q$ be a smooth finite-dimensional complex family.  Define
\begin{equation}
 \operatorname{Det}D
 :=\det(\ker D)\otimes\det(\operatorname{coker}D)^*,
 \qquad \det V:=\bigwedge^{\dim V}V.
 \label{eq:t10v36-determinant-line}
\end{equation}
This convention is dual to another common convention; all transition factors
below are consistent with \eqref{eq:t10v36-determinant-line}.

\subsection{Constant-rank reduction}

Fix a stratum $S_r$ on which $\operatorname{rank}D=r$ is constant.  Let
$P_K$ and $P_C$ be the orthogonal projectors onto
$K=\ker D$ and $C=\ker D^*=\operatorname{coker}D$, and put
\begin{equation}
 D_\perp:=(I-P_C)D|_{K^\perp}:K^\perp\longrightarrow C^\perp.
 \label{eq:t10v36-reduced-operator}
\end{equation}

\begin{lemma}[Smooth reduced inverse]
\label{lem:t10v36-smooth-reduction}
On a constant-rank stratum, $K$, $C$, $P_K$, $P_C$, and the Moore--Penrose
inverse
\begin{equation}
 D^+=D_\perp^{-1}(I-P_C)
 \label{eq:t10v36-pseudoinverse}
\end{equation}
are smooth.  If every nonzero singular value of $D$ is at least $\mu>0$,
then $\|D^+\|\leq\mu^{-1}$.
\end{lemma}

\begin{proof}
A constant-rank matrix has a nonvanishing $r\times r$ minor locally.  Gaussian
elimination using that minor gives smooth local bases for kernel, image, and
their orthogonal complements.  In those bases the nonzero $r\times r$ block
is invertible and its inverse is smooth.  This is
\eqref{eq:t10v36-pseudoinverse}.  The operator norm of the pseudoinverse is
the reciprocal of the smallest nonzero singular value.
\end{proof}

Choose orthonormal frames $k=(k_1,\ldots,k_a)$ of $K$ and
$c=(c_1,\ldots,c_b)$ of $C$.  The positive reduced magnitude is
\begin{equation}
 \Delta_\perp(D):=\det(D_\perp^*D_\perp)^{1/2}
 =\prod_{j=1}^r\sigma_j(D)>0.
 \label{eq:t10v36-reduced-magnitude}
\end{equation}

\begin{proposition}[Local determinant-line section]
\label{prop:t10v36-local-section}
On $S_r$, the data
\begin{equation}
 \mathfrak s_D
 :=(k_1\wedge\cdots\wedge k_a)
 \otimes(c_1\wedge\cdots\wedge c_b)^*
 \label{eq:t10v36-local-line-section}
\end{equation}
define a nonzero local section of $\operatorname{Det}D$.  Under unitary frame
changes $k\mapsto kA$ and $c\mapsto cB$,
\begin{equation}
 \mathfrak s_D\mapsto\det A\,\det B^{-1}\mathfrak s_D.
 \label{eq:t10v36-line-transition}
\end{equation}
The reduced magnitude \eqref{eq:t10v36-reduced-magnitude} is independent of
these frames, so it cannot cancel the phase in
\eqref{eq:t10v36-line-transition}.
\end{proposition}

\begin{proof}
The top exterior powers are one-dimensional.  Exterior multiplication gives
the factor $\det A$ on $\det K$ and the dual frame gives
$\det B^{-1}$ on $\det C^*$.  Singular values are invariant under unitary
basis changes.
\end{proof}

This is the physical-zero-mode analogue of the measured-slice distinction in
V35: magnitudes glue automatically after the Jacobian ledger, while phases
live in a line bundle.

\subsection{Berezin saturation of zero modes}

First take a square complex matrix with orthogonal decompositions
$E=K\oplus K^\perp$ and $F=C\oplus C^\perp$, and assume
$\dim K=\dim C=a$.  In adapted bases,
\begin{equation}
 D=\begin{pmatrix}0&0\\0&D_\perp\end{pmatrix}.
 \label{eq:t10v36-adapted-block}
\end{equation}

\begin{theorem}[Exact zero-mode saturation]
\label{thm:t10v36-berezin-saturation}
Let $(\psi_0,\psi_\perp)$ and
$(\overline\psi_0,\overline\psi_\perp)$ be the corresponding Grassmann
coordinates.  Then the vacuum integral vanishes when $a>0$, while
\begin{equation}
 \int d\overline\psi\,d\psi\;
 \left(\prod_{i=1}^a\overline\psi_{0,i}\psi_{0,i}\right)
 e^{-\overline\psi D\psi}
 =\det D_\perp
 \label{eq:t10v36-saturated-integral}
\end{equation}
up to the fixed orientation convention for the Berezin measure.  Insertions
of lower zero-mode degree give zero.
\end{theorem}

\begin{proof}
The exponent contains only
$-\overline\psi_\perp D_\perp\psi_\perp$.  A Berezin integral extracts the
coefficient of the product of every integration variable.  Without the
displayed insertion, or with fewer zero-mode variables, that coefficient is
zero.  With complete saturation, the zero-mode integral is one in the fixed
orientation and the remaining Gaussian Grassmann integral is
$\det D_\perp$.
\end{proof}

For a rectangular chiral operator or nonzero index, the insertion belongs to
$\det K^*\otimes\det C$ and pairs with
\eqref{eq:t10v36-determinant-line}.  The invariant result is again a scalar;
writing a bare determinant without that pairing is meaningless.

\begin{corollary}[Vacuum weight of a massless topological sector]
\label{cor:t10v36-vacuum-zero}
A massless fermion species with unsaturated exact zero modes gives zero to the
vacuum partition function in that background.  Correlation functions can be
nonzero only when interactions or external insertions saturate the complete
zero-mode wedge.
\end{corollary}

\begin{proof}
Apply Theorem~\ref{thm:t10v36-berezin-saturation} species by species and use
multiplicativity of independent Berezin integrals.
\end{proof}

This is the finite algebra behind the topological multi-fermion vertex.  It
is a selection rule, not permission to delete the gauge sector.

\subsection{Mass regularization of the magnitude}

\begin{lemma}[Reduced singular determinant as a mass limit]
\label{lem:t10v36-mass-limit}
If $a=\dim\ker D$ and the nonzero singular values are
$\sigma_1,\ldots,\sigma_r$, then
\begin{equation}
 \det(D^*D+m^2I)^{1/2}
 =m^a\prod_{j=1}^r(\sigma_j^2+m^2)^{1/2},
 \label{eq:t10v36-mass-factorization}
\end{equation}
and hence
\begin{equation}
 \lim_{m\downarrow0}m^{-a}
 \det(D^*D+m^2I)^{1/2}=\Delta_\perp(D).
 \label{eq:t10v36-reduced-mass-limit}
\end{equation}
\end{lemma}

\begin{proof}
Diagonalize the positive matrix $D^*D$.  It has $a$ zero eigenvalues and the
remaining eigenvalues are $\sigma_j^2$.  Multiply the shifted eigenvalues and
take the positive square root.
\end{proof}

\begin{firewall}[The positive mass limit loses the chiral phase]
\label{fire:t10v36-mass-phase}
Equation \eqref{eq:t10v36-reduced-mass-limit} reconstructs only
$|\det' D|$.  Since $D^*D$ is unchanged when $D$ is multiplied by a phase,
no argument based only on the positive mass-regularized determinant can
construct the section \eqref{eq:t10v36-local-line-section} or its anomaly
cocycle.
\end{firewall}

\subsection{Exact lifting by Yukawa or source couplings}

Let $D$ be square with the block form
\eqref{eq:t10v36-adapted-block} and let $B$ be a perturbation, written
\begin{equation}
 D+B=
 \begin{pmatrix}
 B_{00}&B_{0h}\\ B_{h0}&D_\perp+B_{hh}
 \end{pmatrix}.
 \label{eq:t10v36-lifted-block}
\end{equation}

\begin{theorem}[Zero-mode effective determinant]
\label{thm:t10v36-zero-schur}
If $D_\perp+B_{hh}$ is invertible, then
\begin{align}
 \det(D+B)
 &=\det(D_\perp+B_{hh})\det B_{\rm eff},
 \label{eq:t10v36-zero-schur-det}\\
 B_{\rm eff}
 &:=B_{00}-B_{0h}(D_\perp+B_{hh})^{-1}B_{h0}.
 \label{eq:t10v36-zero-effective}
\end{align}
Thus a Higgs--Yukawa field or an external source lifts the zero modes exactly
through the finite matrix $B_{\rm eff}$, including their mixing with the
nonzero sector.
\end{theorem}

\begin{proof}
Multiply \eqref{eq:t10v36-lifted-block} on the left and right by the usual
block unitriangular elimination matrices.  Their determinants are one and
the result is diagonal with blocks $B_{\rm eff}$ and
$D_\perp+B_{hh}$.  Taking determinants proves
\eqref{eq:t10v36-zero-schur-det}.
\end{proof}

\begin{corollary}[Quantitative lifting criterion]
\label{cor:t10v36-lifting-criterion}
If the smallest nonzero singular value of $D$ is at least $\mu$ and
$\|B_{hh}\|<\mu$, then
\begin{equation}
 \|(D_\perp+B_{hh})^{-1}\|
 \leq(\mu-\|B_{hh}\|)^{-1}.
 \label{eq:t10v36-lifted-inverse}
\end{equation}
If in addition $B_{00}$ is invertible and
\begin{equation}
 \|B_{00}^{-1}B_{0h}\|\,
 \|(D_\perp+B_{hh})^{-1}B_{h0}\|<1,
 \label{eq:t10v36-zero-small-gain}
\end{equation}
then $B_{\rm eff}$ and $D+B$ are invertible.
\end{corollary}

\begin{proof}
The first bound is the Neumann-series estimate.  Factor
$B_{\rm eff}=B_{00}[I-B_{00}^{-1}B_{0h}
(D_\perp+B_{hh})^{-1}B_{h0}]$ and apply the Neumann criterion again, then use
Theorem~\ref{thm:t10v36-zero-schur}.
\end{proof}

For the Standard Model, $B_{00}$ is the zero-mode overlap of the appropriate
Yukawa matrix with the Higgs field.  It can vanish at the symmetric point and
for exactly massless species.  No uniform inverse is assumed.

\subsection{The large-volume infrared obstruction}

\begin{proposition}[Why the physical pseudoinverse is not a local UV owner]
\label{prop:t10v36-ir-obstruction}
Suppose a sequence of volumes contains nonzero singular values
$\mu_L\downarrow0$.  Then
\begin{equation}
 \|D_L^+\|\geq\mu_L^{-1}\longrightarrow\infty.
 \label{eq:t10v36-pseudoinverse-divergence}
\end{equation}
Consequently no volume-uniform estimate that treats the complete physical
pseudoinverse as one gapped local covariance can hold.
\end{proposition}

\begin{proof}
The pseudoinverse maps a normalized left singular vector with singular value
$\mu_L$ to a vector of norm $\mu_L^{-1}$.  This proves
\eqref{eq:t10v36-pseudoinverse-divergence}.
\end{proof}

The repair is a Wilsonian handoff.  At scale $j$, choose a smooth spectral or
momentum partition
\begin{equation}
 I=\Pi_{<j}+\Pi_j+\Pi_{>j}
 \label{eq:t10v36-ir-partition}
\end{equation}
compatible with gauge covariance.  Exact zero modes remain in
$\Pi_{<j}$, the shell $\Pi_j$ is integrated only with its scale-dependent
gap, and no inverse is taken on the retained block.

\begin{theorem}[Schur handoff without an infrared inverse]
\label{thm:t10v36-ir-handoff}
Let $E=E_<\oplus E_>$ and write a perturbed operator as
\begin{equation}
 \mathcal D=
 \begin{pmatrix}D_{<<}&D_{<>}\\D_{><}&D_{>>}\end{pmatrix},
 \qquad D_{>>}\text{ invertible}.
 \label{eq:t10v36-ir-block}
\end{equation}
Then integrating the high Grassmann variables gives the exact retained
operator
\begin{equation}
 D_{\rm ret}=D_{<<}-D_{<>}D_{>>}^{-1}D_{><}
 \label{eq:t10v36-ir-schur}
\end{equation}
and the factor $\det D_{>>}$.  No inverse of $D_{<<}$ is used.
\end{theorem}

\begin{proof}
Translate the high Grassmann variables by the linear terms, or apply block
Gaussian elimination.  The Berezin Jacobian of the triangular translation
is one.  The high quadratic integral is $\det D_{>>}$ and the remaining
quadratic form is \eqref{eq:t10v36-ir-schur}.
\end{proof}

This is the fermionic analogue of retaining massless gauge and gravitational
modes before invoking the V20 gap theorem.

\begin{assessment}[V36 acquisition]
\label{ass:t10v36-status}
The determinant discriminant is now stratified correctly.  Exact topological
zero modes are determinant-line variables and require saturation; Yukawa and
source fields act through the finite effective matrix
\eqref{eq:t10v36-zero-effective}; closing nonzero modes require a shell
handoff rather than a global pseudoinverse.  V37 will build a smooth
finite-range shell partition for the overlap operator and prove a determinant
factorization with a local high-shell phase current and an explicitly
retained low block.
\end{assessment}

\begin{firewall}[Claim boundary after V36]
\label{fire:t10v36-claim}
V36 proves the finite algebra of zero modes and the infrared handoff.  It does
not construct a gauge-covariant smooth shell partition, bound its commutators,
prove uniform Yukawa lifting, sum topological sectors, or construct the
SM--Einstein continuum.
\end{firewall}

\subsection{Parent record}

This V36 note is appended to the validated V35 source, whose record is
\begin{equation}
 \text{bytes}=511768,qquad \text{lines}=13369,qquad
 \operatorname{SHA256}=
 \texttt{4A509CE0DD1609EF8BC5E4A6B6A6B065ADC3B2EEEA1482D0A99B906378993E85}.
 \label{eq:t10v36-parent-record}
\end{equation}

% =============================================================================
% END APPENDED TENTH-SERIES V36 MATERIAL
% =============================================================================

% =============================================================================
% BEGIN APPENDED TENTH-SERIES V37 MATERIAL
% =============================================================================

\section{Tenth-series V37: heat-time determinant shells and a bounded
ultraviolet chiral response}
\label{sec:v10v37}

\subsection{Replacing a sharp spectral cut}

A sharp spectral projector at a moving nonzero cutoff is discontinuous when
an eigenvalue crosses the cutoff.  A smooth function of the physical Dirac
operator is local, but it is not an idempotent and therefore does not give an
exact block determinant.  The repair is to split the logarithm itself by
heat time.  Short heat times are ultraviolet and local; the complete
long-time remainder is retained as infrared data.

Let $D:E\to F$ be a square finite-dimensional complex operator of constant
rank $r$ on the parameter patch, let
\begin{equation}
 A:=D^*D,qquad P_0:=\mathbf1_{\{0\}}(A),qquad
 \Delta_\perp(D)^2=\det{}'A,
 \label{eq:t10v37-positive-data}
\end{equation}
and use the reduced determinant-line frames of V36.

\subsection{Exact heat-time splitting of the magnitude}

\begin{lemma}[Scalar heat identity]
\label{lem:t10v37-scalar-heat}
For every $\lambda>0$ and $T>0$,
\begin{align}
 \log\lambda
 &=\int_0^T\frac{e^{-t}-e^{-t\lambda}}t\,dt
   +\int_T^\infty\frac{e^{-t}-e^{-t\lambda}}t\,dt.
 \label{eq:t10v37-scalar-heat}
\end{align}
Both integrals converge.
\end{lemma}

\begin{proof}
The difference in the numerator is $t(\lambda-1)+O(t^2)$ at zero and both
terms decay at infinity.  Differentiating the full right side with respect
to $\lambda$ gives
$\int_0^\infty e^{-t\lambda}dt=1/\lambda$.  Its value is zero at
$\lambda=1$, hence it equals $\log\lambda$.  Splitting at $T$ proves the
formula.
\end{proof}

\begin{theorem}[Exact ultraviolet/infrared magnitude factorization]
\label{thm:t10v37-heat-factorization}
Define
\begin{align}
 \Gamma_{\rm UV}(T;D)
 &:=\frac12\int_0^T\frac1t
 \left[r e^{-t}-\operatorname{Tr}(e^{-tA}-P_0)\right]dt,
 \label{eq:t10v37-uv-magnitude}\\
 \Gamma_{\rm IR}(T;D)
 &:=\frac12\int_T^\infty\frac1t
 \left[r e^{-t}-\operatorname{Tr}(e^{-tA}-P_0)\right]dt.
 \label{eq:t10v37-ir-magnitude}
\end{align}
Then
\begin{equation}
 \log\Delta_\perp(D)
 =\Gamma_{\rm UV}(T;D)+\Gamma_{\rm IR}(T;D).
 \label{eq:t10v37-exact-magnitude-split}
\end{equation}
No lower bound on the nonzero singular values is needed for the ultraviolet
term at fixed $T$.
\end{theorem}

\begin{proof}
Diagonalize $A$.  Since $e^{-tA}-P_0$ has eigenvalues $e^{-t\sigma_j^2}$ on
the $r$ nonzero singular directions and zero on the kernel, summing
Lemma~\ref{lem:t10v37-scalar-heat} with $\lambda=\sigma_j^2$ and multiplying
by $1/2$ proves the identity.  The short-time integrand is finite at zero by
the cancellation in the lemma, independently of how small $\sigma_j$ is.
\end{proof}

The long-time term contains all divergence as a nonzero singular value tends
to zero.  It is deliberately not estimated by a volume-uniform gap.

\subsection{A bounded ultraviolet pseudoinverse}

Define
\begin{equation}
 B_T:=\int_0^T e^{-tD^*D}D^*\,dt.
 \label{eq:t10v37-bounded-inverse}
\end{equation}

\begin{lemma}[Singular-value filter]
\label{lem:t10v37-filter}
On a right singular vector with singular value $\sigma>0$, $B_T$ has
singular value
\begin{equation}
 b_T(\sigma)=\frac{1-e^{-T\sigma^2}}\sigma,
 \label{eq:t10v37-filter-function}
\end{equation}
and it vanishes on the orthogonal cokernel.  In particular
\begin{equation}
 \|B_T\|\leq\sqrt T.
 \label{eq:t10v37-filter-bound}
\end{equation}
Moreover, on the nonzero subspaces,
\begin{equation}
 D^+=B_T+e^{-TD^*D}D^+.
 \label{eq:t10v37-inverse-split}
\end{equation}
\end{lemma}

\begin{proof}
Apply a singular-value decomposition $D u=\sigma v$ and
$D^*v=\sigma u$ in \eqref{eq:t10v37-bounded-inverse}.  Integration gives
\eqref{eq:t10v37-filter-function}.  With $x=T\sigma^2$,
\[
 b_T(\sigma)=\sqrt T\,\frac{1-e^{-x}}{\sqrt x}
 \leq\sqrt T\min\{\sqrt x,x^{-1/2}\}\leq\sqrt T.
\]
The scalar identity
$\sigma^{-1}=(1-e^{-T\sigma^2})/\sigma+e^{-T\sigma^2}/\sigma$
proves \eqref{eq:t10v37-inverse-split}.
\end{proof}

\begin{proposition}[Exact response split]
\label{prop:t10v37-response-split}
In fixed reduced determinant-line frames, the logarithmic derivative of the
reduced determinant obeys
\begin{align}
 d\log\det{}'D
 &=a_{\rm UV,T}+a_{\rm IR,T},
 \label{eq:t10v37-response-split}\\
 a_{\rm UV,T}&:=\operatorname{Tr}(B_T\,dD),
 \label{eq:t10v37-uv-response}\\
 a_{\rm IR,T}&:=\operatorname{Tr}(e^{-TD^*D}D^+\,dD),
 \label{eq:t10v37-ir-response}
\end{align}
with the kernel and cokernel frame connections added on changes of local
frame.  The ultraviolet response remains bounded as physical nonzero
singular values close.
\end{proposition}

\begin{proof}
On a constant-rank patch, the usual reduced Jacobi formula is
$d\log\det{}'D=\operatorname{Tr}(D^+dD)$ in fixed kernel and cokernel frames.
Insert \eqref{eq:t10v37-inverse-split}.  The bound
\eqref{eq:t10v37-filter-bound} controls the first term; every unbounded
$\sigma^{-1}$ factor in the second is multiplied by its long-time filter.
Frame changes add exactly the transition derivatives from V36.
\end{proof}

\begin{firewall}[The ultraviolet one-form need not be closed]
\label{fire:t10v37-uv-not-closed}
Although $a_{\rm UV,T}+a_{\rm IR,T}$ is the derivative of the reduced
logarithm on a fixed trivialization, its two summands need not be closed.
Their exterior derivatives cancel.  A shell anomaly current must therefore
be carried together with the compensating retained infrared curvature; one
may not integrate $a_{\rm UV,T}$ alone into a global determinant.
\end{firewall}

\subsection{Volume-free spatial locality at fixed heat time}

Use the weighted row norm of V31.  Assume
\begin{equation}
 \max\{\|D\|_\mu,\|D^*\|_\mu\}\leq C_D,\qquad \|D^*D\|_\mu\leq C_A
 \label{eq:t10v37-weighted-data}
\end{equation}
uniformly in volume on an admissible patch.

\begin{theorem}[Weighted locality of the ultraviolet inverse]
\label{thm:t10v37-weighted-filter}
Under \eqref{eq:t10v37-weighted-data},
\begin{equation}
 \|B_T\|_\mu
 \max\{\|D\|_\mu,\|D^*\|_\mu\}\leq C_D,\qquad \|D^*D\|_\mu\leq C_A
 =\begin{cases}
 \max\{\|D\|_\mu,\|D^*\|_\mu\}\leq C_D,\qquad \|D^*D\|_\mu\leq C_A
 \max\{\|D\|_\mu,\|D^*\|_\mu\}\leq C_D,\qquad \|D^*D\|_\mu\leq C_A
 \end{cases}
 \label{eq:t10v37-weighted-filter}
\end{equation}
Consequently
\begin{equation}
 \|(B_T)_{xy}\|
 \leq C_B(T)e^{-\mu d(x,y)}
 \label{eq:t10v37-filter-tail}
\end{equation}
with a volume-independent constant for every fixed $T$.
\end{theorem}

\begin{proof}
The weighted norm is a Banach algebra.  Its exponential series gives
$\|e^{-tA}\|_\mu\leq e^{t\|A\|_\mu}\leq e^{tC_A}$.  Apply this under \eqref{eq:t10v37-bounded-inverse} and multiply by the
assumed bound on $\|D^*\|_\mu$.  A matrix block is bounded by its
weighted row sum, proving \eqref{eq:t10v37-filter-tail}.
\end{proof}

For an overlap operator the hypotheses follow from the V35 admissibility
tail after reducing the exponential weight if necessary.  The estimate is
crude for large $T$ but exact in its logical role: every finite heat shell is
local without a physical infrared gap.

\begin{proposition}[Local variation of the ultraviolet magnitude]
\label{prop:t10v37-uv-variation}
On a constant-rank patch,
\begin{equation}
 d\Gamma_{\rm UV}(T;D)
 =\frac12\int_0^T
 \operatorname{Tr}(e^{-tD^*D}\,d(D^*D))\,dt.
 \label{eq:t10v37-uv-variation}
\end{equation}
If $dD$ is supported in a polymer $X$, the integrand has exponential tails
away from $X$ at every fixed $T$, with a bound obtained from
Theorem~\ref{thm:t10v37-weighted-filter}.
\end{proposition}

\begin{proof}
Duhamel's formula and cyclicity of the trace give
$d\operatorname{Tr}e^{-tA}=-t\operatorname{Tr}(e^{-tA}dA)$.  The rank and
$P_0$ trace are constant on the patch, so differentiating
\eqref{eq:t10v37-uv-magnitude} proves the formula.  Weighted multiplication
of the heat kernel with the locally supported $dA$ gives the tail.
\end{proof}

\subsection{A true multiscale heat partition}

Choose $0=T_0<T_1<\cdots<T_J<T_{J+1}=\infty$ and define
\begin{equation}
 B_j:=\int_{T_j}^{T_{j+1}}e^{-tD^*D}D^*dt
 \quad(0\leq j<J),
 \label{eq:t10v37-heat-shell}
\end{equation}
leaving the interval $[T_J,\infty)$ in the retained infrared block.

\begin{theorem}[Exact finite-shell response ledger]
\label{thm:t10v37-shell-ledger}
On the nonzero subspaces,
\begin{equation}
 D^+=\sum_{j=0}^{J-1}B_j+e^{-T_JD^*D}D^+.
 \label{eq:t10v37-shell-ledger}
\end{equation}
Each finite shell is exponentially local under
\eqref{eq:t10v37-weighted-data}; all failure of a volume-uniform physical gap
is confined to the final retained term.
\end{theorem}

\begin{proof}
Sum the integrals in \eqref{eq:t10v37-heat-shell} to obtain $B_{T_J}$ and use
\eqref{eq:t10v37-inverse-split}.  The locality proof on a finite interval is
identical to Theorem~\ref{thm:t10v37-weighted-filter}.
\end{proof}

The choice $T_j$ proportional to the square of the physical length scale
makes $e^{-tD^*D}$ the parabolic shell router.  Exact zeros are annihilated by
$D^*$ in every $B_j$ and remain solely in the V36 determinant-line wedge.

\subsection{Gauge and metric covariance}

\begin{lemma}[Covariance of the heat router]
\label{lem:t10v37-covariance}
If $D(q^g)=R_F(g,q)D(q)R_E(g,q)^{-1}$ with unitary endpoint actions, then
\begin{equation}
 B_T(q^g)=R_E(g,q)B_T(q)R_F(g,q)^{-1}.
 \label{eq:t10v37-filter-covariance}
\end{equation}
The same statement holds for a spin-frame or lattice-diffeomorphism change
implemented unitarily.
\end{lemma}

\begin{proof}
The positive operator transforms by
$D(q^g)^*D(q^g)=R_E(D^*D)R_E^{-1}$.  Functional calculus commutes with
unitary conjugation; insert this identity in
\eqref{eq:t10v37-bounded-inverse}.
\end{proof}

Thus the heat split does not choose momentum coordinates and extends to the
curved metric chart once the curved overlap operator and its weighted bounds
are supplied.

\begin{assessment}[V37 acquisition]
\label{ass:t10v37-status}
The physical infrared obstruction of V36 no longer contaminates the local
ultraviolet inverse.  The exact heat identity separates determinant
magnitude and logarithmic response into finite local shells plus one retained
long-time term, without a sharp spectral boundary.  V38 will convert the
finite-shell traces into connected polymer owners, prove the normalized-trace
volume bound, and quantify the reserve loss as $T_j$ grows.
\end{assessment}

\begin{firewall}[Claim boundary after V37]
\label{fire:t10v37-claim}
V37 proves finite heat-shell identities and locality bounds.  It does not
construct a convergent infinite shell expansion, trivialize the chiral phase,
bound the final infrared term, build the curved overlap operator, or take the
SM--Einstein continuum limit.
\end{firewall}

\subsection{Parent record}

This V37 note is appended to the validated V36 source, whose record is
\begin{equation}
 \text{bytes}=524566,qquad \text{lines}=13714,qquad
 \operatorname{SHA256}=
 \texttt{33E96C779DFA0C2F7D25CA7CCB0BC0C0D0D6EC9C597AE69BA0E89C6351FB902F}.
 \label{eq:t10v37-parent-record}
\end{equation}

% =============================================================================
% END APPENDED TENTH-SERIES V37 MATERIAL
% =============================================================================

% =============================================================================
% BEGIN APPENDED TENTH-SERIES V38 MATERIAL
% =============================================================================

\section{Tenth-series V38: heat-shell marked traces, spatial clustering,
and the polymer reserve}
\label{sec:v10v38}

\subsection{Weighted heat propagation}

V37 made every finite heat interval analytic without a physical infrared
gap.  To use it in the renewal expansion one still needs volume-free
connectedness: a variation supported in one polymer must have exponentially
small sensitivity to a distant variation.  This note proves that statement
from a weighted conjugation bound.

Let \(A=A^*\succeq0\) act on
\(\mathcal H=\bigoplus_{x\in\Lambda}\mathbb C^{n_f}\).  For a site \(z\),
let \(W_{\mu,z}\) be multiplication by
\(e^{\mu d(x,z)}\), with the weight truncated on a finite torus before the
bound is taken.  Assume
\begin{equation}
 \sup_z\|W_{\mu,z}AW_{\mu,z}^{-1}-A\|\leq v(\mu).
 \label{eq:t10v38-conjugation-cost}
\end{equation}

\begin{lemma}[Weighted heat-kernel bound]
\label{lem:t10v38-heat-kernel}
Under \eqref{eq:t10v38-conjugation-cost},
\begin{equation}
 \|(e^{-tA})_{xy}\|
 \leq e^{t v(\mu)}e^{-\mu d(x,y)},\qquad t\geq0.
 \label{eq:t10v38-heat-tail}
\end{equation}
The estimate is independent of volume when \(v(\mu)\) is.
\end{lemma}

\begin{proof}
Write \(WAW^{-1}=A+E\), with \(\|E\|\leq v(\mu)\).  Since
\(\|e^{-tA}\|\leq1\), the Duhamel series for
\(e^{-t(A+E)}\) gives
\[
 \|We^{-tA}W^{-1}\|=\|e^{-t(A+E)}\|\leq e^{t\|E\|}
 \leq e^{tv(\mu)}.
\]
Choose the weight centre \(z=y\) and take the \((x,y)\) block.  Removing the
truncation after the block estimate gives \eqref{eq:t10v38-heat-tail}.
\end{proof}

\begin{lemma}[Finite-range conjugation cost]
\label{lem:t10v38-finite-range-cost}
If \(A_{xy}=0\) for \(d(x,y)>R\) and
\begin{equation}
 J:=\max\left\{
 \sup_x\sum_y\|A_{xy}\|,\,
 \sup_y\sum_x\|A_{xy}\|
 \right\}<\infty,
 \label{eq:t10v38-row-column}
\end{equation}
then
\begin{equation}
 v(\mu)\leq J(e^{\mu R}-1).
 \label{eq:t10v38-explicit-cost}
\end{equation}
For an exponentially local \(A\), the same conclusion holds with
\(J(e^{\mu R}-1)\) replaced by its weighted hopping moment.
\end{lemma}

\begin{proof}
The conjugated block is multiplied by
\(e^{\mu(d(x,z)-d(y,z))}\).  The triangle inequality and the range condition
bound the absolute change of this factor by \(e^{\mu R}-1\).  The row and
column Schur test gives \eqref{eq:t10v38-explicit-cost}.  Summing the actual
factor against exponentially decaying blocks proves the last assertion.
\end{proof}

The growth \(e^{tv(\mu)}\) is the exact reserve cost in this elementary
bound.  It is harmless at fixed heat time and prevents one from sending
\(T\) to infinity inside a fixed locality norm.

\subsection{Duhamel derivatives with marked supports}

Let \(V_i=V_i^*\) be variations of \(A\), each supported on a finite set
\(X_i\) in the sense that \((V_i)_{xy}=0\) unless both endpoints lie in
\(X_i\).  Let \(\ell(X_1,\ldots,X_n)\) be the length of a minimum spanning
tree joining the sets, using the lattice distance between sets.

\begin{lemma}[Ordered Duhamel formula]
\label{lem:t10v38-duhamel}
For \(n\geq1\),
\begin{equation}
 D^n e^{-tA}[V_1,\ldots,V_n]
 =(-1)^n\sum_{\pi\in S_n}
 \int_{\substack{s_j\geq0\\s_0+\cdots+s_n=t}}
 e^{-s_0A}V_{\pi(1)}e^{-s_1A}\cdots
 V_{\pi(n)}e^{-s_nA}\,ds_1\cdots ds_n .
 \label{eq:t10v38-duhamel}
\end{equation}
The simplex has volume \(t^n/n!\).
\end{lemma}

\begin{proof}
The first derivative is Duhamel's formula.  Differentiate inductively; the
new insertion may occur in any of the existing heat intervals.  Relabelling
the resulting ordered simplices gives one term for every permutation.
\end{proof}

\begin{theorem}[Marked trace clustering]
\label{thm:t10v38-marked-clustering}
Assume \eqref{eq:t10v38-conjugation-cost} and
\(\|V_i\|_{\mu,1}\leq M_i\), where the norm controls both weighted row and
column sums.  Then there is a fixed local-fibre constant \(C_f\) such that
\begin{equation}
 \left|
 \operatorname{Tr}D^ne^{-tA}[V_1,\ldots,V_n]
 \right|
 \leq
 C_f^n t^n e^{tv(\mu)}
 e^{-\mu\ell(X_1,\ldots,X_n)}
 \prod_{i=1}^n |X_i|M_i .
 \label{eq:t10v38-marked-clustering}
\end{equation}
No factor of \(|\Lambda|\) occurs.
\end{theorem}

\begin{proof}
Expand one term of \eqref{eq:t10v38-duhamel} in position-space blocks.  Every
variation forces its two adjacent vertices into \(X_i\).  The heat kernels
therefore form a cyclic route visiting all marked sets.  The sum of the
route distances is at least the length of a spanning tree joining them.
Apply Lemma~\ref{lem:t10v38-heat-kernel} to every heat block.  The time
exponents multiply to \(e^{(s_0+\cdots+s_n)v(\mu)}=e^{tv(\mu)}\), and the
spatial exponents give at most \(e^{-\mu\ell}\).  Weighted row/column
summation absorbs all unmarked intermediate sites; the remaining choices in
each \(X_i\) give the displayed support factors and the fixed fibre factor.
The \(n!\) permutations cancel the simplex factor \(1/n!\), leaving \(t^n\).
\end{proof}

\begin{firewall}[Support is exponentially enlarged, not exactly preserved]
\label{fire:t10v38-exact-support}
Unlike the radial configuration homotopy of V29, a heat kernel propagates
between sites.  A marked heat owner has exponential spatial tails rather
than exact support \(X_i\).  Those tails must be allocated to connecting
polymers; calling the owner strictly supported on \(\bigcup_iX_i\) would be
false.
\end{firewall}

\subsection{Ultraviolet determinant derivatives}

Let \(\Gamma_{\rm UV}(T;D)\) be the V37 ultraviolet magnitude and
\(A=D^*D\).  For independent affine variations of \(A\), the \(n\)-th
derivative comes only from the heat exponential.

\begin{corollary}[Finite-shell analytic bound]
\label{cor:t10v38-shell-analytic}
For \(n\geq1\) and local affine variations \(V_i\),
\begin{equation}
 \left|
 D^n\Gamma_{\rm UV}(T)[V_1,\ldots,V_n]
 \right|
 \leq
 \frac{C_f^n}{2n}T^n e^{Tv(\mu)}
 e^{-\mu\ell(X_1,\ldots,X_n)}
 \prod_{i=1}^n|X_i|M_i .
 \label{eq:t10v38-shell-analytic}
\end{equation}
For a nonlinear family \(A(q)\), the same estimate is summed over set
partitions, with each \(V_i\) replaced by the corresponding derivative block
of \(A\).
\end{corollary}

\begin{proof}
Differentiate \eqref{eq:t10v37-uv-magnitude}.  The reference and kernel-rank
terms are constant, and the factor \(t^{-1}\) changes the bound in
Theorem~\ref{thm:t10v38-marked-clustering} to \(t^{n-1}\).  Bound
\(e^{tv(\mu)}\leq e^{Tv(\mu)}\) and integrate
\(\int_0^Tt^{n-1}dt=T^n/n\).  The higher chain rule is indexed by set
partitions, which gives the nonlinear statement.
\end{proof}

\begin{corollary}[Normalized trace bound]
\label{cor:t10v38-normalized-trace}
For a translation-covariant one-site variation \(V_x\) with uniform norm,
\begin{equation}
 \frac1{|\Lambda|}\sum_{x\in\Lambda}
 \left|D\Gamma_{\rm UV}(T)[V_x]\right|
 \leq \frac{C_f}{2}T e^{Tv(\mu)}M ,
 \label{eq:t10v38-normalized-trace}
\end{equation}
uniformly in volume.
\end{corollary}

\begin{proof}
Apply \eqref{eq:t10v38-shell-analytic} with \(n=1\) and \(|X_1|=1\), sum the
identical bound over \(x\), and divide by the volume.
\end{proof}

This is the heat-shell version of the volume-free normalized-row estimate
imported from YM.  It applies equally to gauge, Higgs, and local metric
variations once their derivative blocks satisfy the same row/column bounds.

\subsection{From marked tails to polymer owners}

For finite sets \(X_1,\ldots,X_n\), allocate the exponential
\(e^{-\mu\ell(X_1,\ldots,X_n)}\) to a chosen minimum tree and thicken each
tree edge into a connected lattice polymer \(Y\).  The number of possible
trees is controlled by the V21 Cayley/forest ledger.

\begin{proposition}[One reserve debit for support size]
\label{prop:t10v38-reserve-debit}
Let
\begin{equation}
 \|V\|_\kappa
 :=\sup_b\sum_{X\ni b}e^{\kappa|X|}\,|X|\,\|V_X\|.
 \label{eq:t10v38-owner-norm}
\end{equation}
For \(0\leq\kappa'<\kappa\), each additional factor \(|X|\) generated by a
marked trace obeys
\begin{equation}
 e^{\kappa'|X|}|X|
 \leq e^{\kappa|X|}
 \sup_{m\geq1}m e^{-(\kappa-\kappa')m}.
 \label{eq:t10v38-size-debit}
\end{equation}
Hence the marked support factors in
\eqref{eq:t10v38-shell-analytic} spend a fixed positive size reserve and no
volume factor.
\end{proposition}

\begin{proof}
Equation \eqref{eq:t10v38-size-debit} is immediate after multiplying and
dividing by \(e^{(\kappa-\kappa')|X|}\).  Sum over owners containing a fixed
block \(b\).  The supremum is finite for \(\kappa>\kappa'\).
\end{proof}

The connecting tails impose a second, geometric condition.  If the number of
connected polymers of size \(m\) through a fixed block is at most \(C_0^m\),
then
\begin{equation}
 \sum_{Y\ni b}e^{\kappa'|Y|}e^{-\mu\,{\rm length}(Y)}
 \label{eq:t10v38-tail-sum}
\end{equation}
is finite whenever the conversion from length to polymer size leaves an
exponential margin larger than \(\kappa'+\log C_0\).  This is the same
entropy-versus-decay comparison used in a standard polymer expansion.

\begin{theorem}[Conditional finite-heat polymer gate]
\label{thm:t10v38-polymer-gate}
Fix \(T<\infty\).  Suppose:
\begin{enumerate}
\item the derivative blocks of \(A\) have finite V19 Gevrey owner norm;
\item the spatial heat decay \(\mu\) beats the connected-polymer entropy as
      in \eqref{eq:t10v38-tail-sum};
\item a positive size reserve \(\kappa-\kappa'\) is available; and
\item the V21 tree activity with
      \(K_T=C_fT e^{Tv(\mu)}\) satisfies
      \begin{equation}
       q_T:=e zK_T<1 .
       \label{eq:t10v38-tree-smallness}
      \end{equation}
\end{enumerate}
Then the connected finite-heat determinant response has an absolutely
convergent, volume-uniform polymer expansion in the weakened norm
\(\kappa'\).
\end{theorem}

\begin{proof}
Corollary~\ref{cor:t10v38-shell-analytic} supplies an exponentially connected
marked activity.  Allocate it to spanning trees.  Proposition~
\ref{prop:t10v38-reserve-debit} absorbs support factors, and item 2 sums the
geometric thickenings.  The forest formula reduces connected combinatorics
to labelled trees; the V21 Cayley estimate sums their activities precisely
under \eqref{eq:t10v38-tree-smallness}.  The bounds are anchored at one
block, so no volume factor remains.
\end{proof}

\begin{firewall}[Large heat time consumes the smallness margin]
\label{fire:t10v38-large-time}
The gate constant \(K_T=C_fT e^{Tv(\mu)}\) grows with \(T\).  Therefore one
cannot put the entire long-time determinant into one polymer activity even
when every finite interval is local.  Heat time must track the RG scale, and
the V15 occupancy cycle must restore reserve between shells.
\end{firewall}

\subsection{Curvature and the retained block}

The same Duhamel formula applies to the exterior derivative of the
ultraviolet one-form \(a_{\rm UV,T}\).  Its marked kernels obey the same
spatial tree bound.  By V37,
\begin{equation}
 da_{\rm UV,T}=-da_{\rm IR,T}
 \label{eq:t10v38-curvature-transfer}
\end{equation}
in a fixed determinant-line chart.

\begin{proposition}[Local curvature transfer]
\label{prop:t10v38-curvature-transfer}
Under the hypotheses of Theorem~\ref{thm:t10v38-polymer-gate},
\(da_{\rm UV,T}\) is a volume-uniform local polymer two-form.  Subtracting
its V29 homotopy current transfers the equal and opposite curvature to the
retained infrared determinant line without changing the full determinant.
\end{proposition}

\begin{proof}
Differentiate the integral formula for \(B_T\); Duhamel's formula inserts one
additional local variation and Theorem~\ref{thm:t10v38-marked-clustering}
applies.  The polymer summation is the same as in
Theorem~\ref{thm:t10v38-polymer-gate}.  Equation
\eqref{eq:t10v38-curvature-transfer} shows that adding opposite local
connection currents to the two factors leaves their sum unchanged.
\end{proof}

\begin{assessment}[V38 acquisition]
\label{ass:t10v38-status}
The finite heat shells now meet the connected-owner requirements: their
mixed variations cluster by a spatial spanning-tree length, normalized
traces are volume free, and the exact reserve and tree-smallness debits are
displayed.  The remaining issue is scale choice.  V39 will choose
\(T_j\) relative to the block scale, prove a reblocking inequality that
restores the lost spatial exponent, and identify the marginal
four-dimensional power of each Standard Model field sector.
\end{assessment}

\begin{firewall}[Claim boundary after V38]
\label{fire:t10v38-claim}
V38 proves a conditional finite-heat polymer gate.  It does not prove the
gate uniformly for an infinite sequence of heat shells, supply the nonlinear
curved derivative stocks, control the retained infrared determinant, or
construct the SM--Einstein continuum.
\end{firewall}

\subsection{Parent record}

This V38 note is appended to the validated V37 source, whose record is
\begin{equation}
 \text{bytes}=536232,\qquad \text{lines}=14026,\qquad
 \operatorname{SHA256}=
 \texttt{3D0F23994DA2BA516911746FDBF590A072E46D46F79C56843CE613DA1E9442C9}.
 \label{eq:t10v38-parent-record}
\end{equation}

% =============================================================================
% END APPENDED TENTH-SERIES V38 MATERIAL
% =============================================================================

% =============================================================================
% BEGIN APPENDED TENTH-SERIES V39 MATERIAL
% =============================================================================

\section{Tenth-series V39: parabolic reblocking, a covariant heat bound,
and the four-dimensional scaling ledger}
\label{sec:v10v39}

\subsection{The scale problem left by V38}

The elementary conjugation estimate in V38 has
\(v(\mu)=O(\mu)\) for a generic finite-range matrix.  With heat time of order
the square of a block diameter, that estimate loses an exponential of the
diameter and cannot iterate.  A second-order operator has a stronger
parabolic bound, quadratic in \(\mu\).  This note proves it for the covariant
nearest-neighbour Laplacian and states the exact transfer condition required
of the overlap square.

\subsection{Covariant random-walk domination}

On a unit lattice let
\begin{equation}
 (\Delta_U\psi)(x)
 :=\sum_{\nu=1}^d
 \left(2\psi(x)-U_{x\nu}\psi(x+\widehat\nu)
 -U^*_{x-\widehat\nu,\nu}\psi(x-\widehat\nu)\right),
 \label{eq:t10v39-covariant-laplacian}
\end{equation}
where the link matrices are unitary in a fixed finite representation.

\begin{lemma}[Path domination]
\label{lem:t10v39-path-domination}
Let \(p_t(x,y)\) be the scalar heat kernel of the ordinary lattice
Laplacian.  Then
\begin{equation}
 \|(e^{-t\Delta_U})_{xy}\|\leq p_t(x,y).
 \label{eq:t10v39-path-domination}
\end{equation}
The statement is independent of the gauge field and the volume.
\end{lemma}

\begin{proof}
Write \(\Delta_U=2dI-\mathcal A_U\), where \(\mathcal A_U\) is the sum of the
\(2d\) covariant shifts.  The exponential series
\[
 e^{-t\Delta_U}=e^{-2dt}\sum_{n\geq0}\frac{t^n}{n!}\mathcal A_U^n
\]
is a sum over nearest-neighbour paths.  Every path amplitude is an ordered
product of unitary links and has norm one.  Taking norms term by term leaves
exactly the scalar path sum defining \(p_t\).
\end{proof}

\begin{theorem}[Quadratic weighted heat bound]
\label{thm:t10v39-quadratic-heat}
For every \(\mu\geq0\),
\begin{equation}
 \|(e^{-t\Delta_U})_{xy}\|
 \leq
 \exp\left\{-\mu |x-y|_1+2dt(\cosh\mu-1)\right\}.
 \label{eq:t10v39-quadratic-heat}
\end{equation}
In particular, for \(0\leq\mu\leq\mu_0\),
\begin{equation}
 2d(\cosh\mu-1)\leq d\cosh(\mu_0)\mu^2.
 \label{eq:t10v39-quadratic-cost}
\end{equation}
\end{theorem}

\begin{proof}
For the scalar continuous-time walk generated by the ordinary lattice
Laplacian,
\begin{equation}
 \sum_y e^{\eta\cdot(y-x)}p_t(x,y)
 =\exp\left\{2t\sum_{\nu=1}^d(\cosh\eta_\nu-1)\right\}.
 \label{eq:t10v39-walk-mgf}
\end{equation}
For fixed \(x,y\), choose
\(\eta_\nu=\mu\,\operatorname{sign}(y_\nu-x_\nu)\).  Positivity of the scalar
kernel and \eqref{eq:t10v39-walk-mgf} give
\[
 p_t(x,y)e^{\mu|x-y|_1}
 \leq e^{2dt(\cosh\mu-1)}.
\]
Apply Lemma~\ref{lem:t10v39-path-domination}.  Taylor's formula with positive
coefficients gives
\(\cosh\mu-1\leq\frac12\mu^2\cosh\mu_0\) on the stated interval.
\end{proof}

The result survives a nonnegative on-site potential.  A negative part
\(V_-\) contributes \(e^{t\|V_-\|}\), so scale iteration requires
\(\|V_-\|=O(\ell^{-2})\) or a separate relevant-parameter treatment.

\subsection{Uniformity under parabolic reblocking}

Let the \(j\)-th block diameter in fine lattice units be
\begin{equation}
 \ell_j=L^j,\qquad
 T_j=\tau\ell_j^2,\qquad
 \mu_j=\frac{\widehat\mu}{\ell_j},
 \qquad 1\leq\tau\leq L^2.
 \label{eq:t10v39-parabolic-scaling}
\end{equation}

\begin{corollary}[Restored spatial exponent]
\label{cor:t10v39-restored-exponent}
For the covariant Laplacian and \(t\leq T_j\),
\begin{equation}
 \|(e^{-t\Delta_U})_{xy}\|
 \leq C_{\rm par}
 e^{-\widehat\mu\,d_j(x,y)},
 \qquad
 d_j(x,y):=\frac{|x-y|_1}{\ell_j},
 \label{eq:t10v39-restored-exponent}
\end{equation}
where
\begin{equation}
 C_{\rm par}
 :=\exp\{d\tau\widehat\mu^2\cosh\widehat\mu\}
 \label{eq:t10v39-parabolic-constant}
\end{equation}
is independent of \(j\) and volume.
\end{corollary}

\begin{proof}
Insert \(\mu_j\) and \(t\leq\tau\ell_j^2\) into
\eqref{eq:t10v39-quadratic-heat}, and use
\eqref{eq:t10v39-quadratic-cost}.  The spatial exponent becomes
\(-\widehat\mu d_j(x,y)\), and the time cost is at most the exponent in
\eqref{eq:t10v39-parabolic-constant}.
\end{proof}

This is the analytic content of restoring locality by reblocking: the decay
rate is fixed in block distance even though it tends to zero in fine-site
distance.

\begin{theorem}[Scale-uniform marked heat gate]
\label{thm:t10v39-scale-uniform-gate}
Assume the positive operator \(A_j\) at scale \(j\) obeys the analogue of
\eqref{eq:t10v39-restored-exponent}, and its \(m\)-fold local derivative
blocks satisfy the parabolic scaling
\begin{equation}
 \|D^mA_j[V_1,\ldots,V_m]\|_{\mu_j,1}
 \leq C_m\ell_j^{-2}\prod_{a=1}^m\|V_a\|_j.
 \label{eq:t10v39-derivative-scaling}
\end{equation}
Then every factor \(T_j^n\) in the V38 \(n\)-marked bound is cancelled by
the \(\ell_j^{-2n}\) from \(n\) affine insertions.  The shell tree constant
may be chosen independent of \(j\):
\begin{equation}
 K_{\rm par}\leq C_f\tau C_1 C_{\rm par}.
 \label{eq:t10v39-uniform-tree-constant}
\end{equation}
Nonlinear derivative blocks are controlled by the corresponding
dimensionless sequence \(C_m\).
\end{theorem}

\begin{proof}
For affine insertions, substitute
\(T_j=\tau\ell_j^2\) and
\(M_{i,j}\leq C_1\ell_j^{-2}\|V_i\|_j\) into
\eqref{eq:t10v38-shell-analytic}.  The powers of \(\ell_j\) cancel exactly,
and Corollary~\ref{cor:t10v39-restored-exponent} replaces
\(e^{T_jv(\mu_j)}\) by \(C_{\rm par}\).  The same substitution in each block
of the set-partition chain rule proves the nonlinear statement.
\end{proof}

\begin{firewall}[The overlap transfer is a theorem still to prove]
\label{fire:t10v39-overlap-transfer}
Theorem~\ref{thm:t10v39-quadratic-heat} is proved for
\(\Delta_U\).  The physical operator in V37 is
\(A_{\rm ov}=D_{\rm ov}^*D_{\rm ov}\), which is exponentially local but not
nearest-neighbour.  Exponential locality alone gives only the V38 linear
conjugation cost.  A quadratic Davies bound and
\eqref{eq:t10v39-derivative-scaling} for \(A_{\rm ov}\) are separate
acquisitions.
\end{firewall}

\subsection{A sufficient paired-hopping extension}

The parabolic cancellation extends beyond nearest neighbours when the
operator has a positive paired-hopping form
\begin{equation}
 A=\sum_{\rho\in\mathcal R}c_\rho
 \left(2I-T_\rho-T_\rho^*\right)+V,
 \qquad c_\rho\geq0,\qquad V\succeq0,
 \label{eq:t10v39-paired-hopping}
\end{equation}
where \(T_\rho\) is a unitary covariant shift by \(\rho\).

\begin{proposition}[Quadratic bound for positive paired hopping]
\label{prop:t10v39-paired-hopping}
For \eqref{eq:t10v39-paired-hopping},
\begin{equation}
 \|(e^{-tA})_{xy}\|
 \leq
 \exp\left\{-\mu|x-y|_1+
 2t\sum_{\rho}c_\rho(\cosh(\mu|\rho|_1)-1)\right\}.
 \label{eq:t10v39-paired-bound}
\end{equation}
If
\(\sum_\rho c_\rho|\rho|_1^2e^{\mu_0|\rho|_1}<\infty\), the time exponent is
bounded by \(C\mu^2t\) for \(0\leq\mu\leq\mu_0\).
\end{proposition}

\begin{proof}
Expand the heat semigroup as a continuous-time walk whose jumps are
\(\pm\rho\) with rate \(c_\rho\).  Unit norm of every covariant path amplitude
gives scalar domination as in Lemma~\ref{lem:t10v39-path-domination}.  The
exponential moment of the scalar walk is the second term in
\eqref{eq:t10v39-paired-bound}.  Finally,
\(\cosh u-1\leq\frac12u^2e^{|u|}\) gives the quadratic estimate and the
stated moment condition.
\end{proof}

A proof that \(A_{\rm ov}\) admits \eqref{eq:t10v39-paired-hopping}, or a
direct substitute based on its functional calculus, would close the heat
locality side of the scale gate.

\subsection{Four-dimensional heat scaling}

The parabolic choice is also the canonical dimension ledger.  For a
Laplace-type continuum operator on a smooth four-manifold, the local heat
trace has the formal small-time structure
\begin{equation}
 \operatorname{Tr}e^{-tA}
 \sim (4\pi t)^{-2}
 \sum_{k\geq0}t^k\int_M a_{2k}(x)\,d{\rm vol}_g,
 \label{eq:t10v39-heat-asymptotic}
\end{equation}
where \(a_{2k}\) is a local invariant of engineering dimension \(2k\).

\begin{proposition}[Relevant, marginal, and irrelevant heat rows]
\label{prop:t10v39-heat-dimensions}
In four dimensions the determinant heat integral \(dt/t\) assigns:
\begin{enumerate}
\item \(a_0\) to the cosmological or vacuum row, with power divergence;
\item \(a_2\) to mass, Higgs quadratic, and Einstein--Hilbert curvature rows,
      with quadratic divergence;
\item \(a_4\) to gauge kinetic, scalar quartic/Yukawa-induced, and
      curvature-squared rows, with logarithmic scaling; and
\item \(a_{2k}\), \(k>2\), to irrelevant higher-derivative rows.
\end{enumerate}
Therefore a curved Standard Model determinant does not preserve a bare
Einstein-plus-renormalizable-matter action unless the required gravitational
counterterm rows, including curvature squared, are present.
\end{proposition}

\begin{proof}
Insert the \(k\)-th term of \eqref{eq:t10v39-heat-asymptotic} into
\(\int dt/t\).  Its power is \(t^{k-3}dt\): \(k=0,1\) are power divergent,
\(k=2\) is logarithmic, and \(k>2\) is suppressed at short time.  The
classification of local invariants follows from their dimensions:
\(1\) has dimension zero, \(R\) and mass squares dimension two, and
\(F^2\), \(R^2\), \(R_{\mu\nu}R^{\mu\nu}\), quartic scalar terms, and
dimension-four kinetic/Yukawa combinations lie in \(a_4\).
\end{proof}

This does not repeat the V17 gravity power-counting result.  It identifies
where those required rows occur in the actual heat-shell router and why the
vacuum, Einstein, and curvature-squared coefficients must be projected at
every reblocking step.

\begin{theorem}[Scale-step projection requirement]
\label{thm:t10v39-projection-requirement}
A closed finite-step SM--gravity RG map based on the V37 heat shells must,
at minimum, project each shell response onto
\begin{equation}
 \int\sqrt g,\quad
 \int\sqrt g\,R,\quad
 \int\sqrt g\,F_i^2,\quad
 \int\sqrt g\,|D\phi|^2,\quad
 \int\sqrt g\,V_{\leq4}(\phi),\quad
 \int\sqrt g\,\overline\psi\slashed D\psi,\quad
 \int\sqrt g\,R^2,\quad
 \int\sqrt g\,R_{\mu\nu}R^{\mu\nu},
 \label{eq:t10v39-required-rows}
\end{equation}
together with topological densities and the gauge-fixing/ghost quotient
ledger.  The complement may enter the irrelevant polymer norm only after
these rows are subtracted.
\end{theorem}

\begin{proof}
Proposition~\ref{prop:t10v39-heat-dimensions} lists all classes with
dimension at most four that are generated by a general Laplace-type matter
or gauge determinant, subject to the symmetries.  Leaving any generated row
in the irrelevant remainder gives it a nonnegative scaling exponent and
prevents contraction.  Projection and subtraction are therefore necessary
for the irrelevant norm to shrink.
\end{proof}

\begin{firewall}[Einstein gravity alone is not an invariant finite theory]
\label{fire:t10v39-einstein-only}
The presence of curvature-squared rows in \eqref{eq:t10v39-required-rows}
does not prove a nonperturbative ultraviolet completion.  It proves that an
Einstein-only finite-dimensional coupling surface is not closed under the
curved matter determinant.  One must use an effective tower with controlled
irrelevant tails or prove a fixed point with enough ultraviolet directions.
\end{firewall}

\subsection{V39 acquisition card}

\begin{assessment}[What is now scale uniform]
\label{ass:t10v39-status}
Parabolic reblocking restores a fixed spatial exponent and cancels the heat
time powers in marked responses, provided the squared Dirac operator obeys a
quadratic weighted heat bound and its field derivatives scale as a
second-order operator.  These facts are proved for covariant Laplacians and
positive paired hopping, but not yet for the overlap square.  V40 will
perform the compulsory companion audit and then attack that transfer through
the integral functional calculus of the sign operator.
\end{assessment}

\begin{firewall}[Claim boundary after V39]
\label{fire:t10v39-claim}
V39 proves the parabolic gate for specified positive hopping operators and
the necessary four-dimensional projection ledger.  It does not prove the
gate for the overlap square, close the gravitational coupling tower, control
the infrared determinant, or construct the continuum.
\end{firewall}

\subsection{Parent record}

This V39 note is appended to the validated V38 source, whose record is
\begin{equation}
 \text{bytes}=549512,\qquad \text{lines}=14371,\qquad
 \operatorname{SHA256}=
 \texttt{F3632804918D80B690C1999B4C6BE7D5CE495E8CB85D73DF0D19DD74A65BC317}.
 \label{eq:t10v39-parent-record}
\end{equation}

% =============================================================================
% END APPENDED TENTH-SERIES V39 MATERIAL
% =============================================================================

% =============================================================================
% BEGIN APPENDED TENTH-SERIES V40 MATERIAL
% =============================================================================

\section{Tenth-series V40: compulsory audit and the quadratic Davies
bound for the overlap square}
\label{sec:v10v40}

\subsection{Compulsory companion audit}

At the audit time the two newest YM filenames, V46 and V47, were
byte-identical, and both ended with the V46 claim boundary; no V47 appendix
was present.  The substantive newest checkpoint and the NS checkpoint are
\begin{equation}
\begin{array}{c|r|r|l}
 \text{source}&\text{bytes}&\text{lines}&\text{SHA--256}\\ \hline
 \text{YM V46/V47}&741227&21240&
 \texttt{9FE733C6324D19F6B6836E4F6AFAF41A151B64A4C1BC5A7968FC942B90C0F39E}\\
 \text{NS V26}&278283&5319&
 \texttt{82C887F8C2C69E4F28FAD7C3DC8DE5B9099CB5A4BF431EBA1FFF3E91935978AD}
\end{array}
\label{eq:t10v40-audit-record}
\end{equation}
The inspected filenames were
\texttt{Renewal\_Yang\_Mills\_Mass\_Gap\_Research\_Notes\_V47.tex} and
\texttt{Renewal\_NS\_Stability\_Research\_Notes\_V26.tex}, with YM V46
checked to establish the duplicate.

YM V46 proves that canonical rescaling converts a growing compliance contact
into a summable quartic precision germ.  It also proves an exact nested Schur
determinant cocycle: successive Gaussian shell log determinants add without
double counting.  Its contact telescope discards only terms lying exactly in
the kernel of the tested moment; a local two-derivative counterterm is
marginal and must be retained.  NS V26 is unchanged and supplies no estimate
for the present operator problem.

\begin{proposition}[Audit decision]
\label{prop:t10v40-audit-decision}
The YM nested-determinant and exact-contact rules are imported into the heat
router.  Adjacent heat intervals are combined algebraically before
estimation, and the dimension-at-most-four projection of V39 is removed
exactly rather than bounded as an irrelevant remainder.  No NS result is
imported.
\end{proposition}

\begin{proof}
Heat intervals partition one determinant logarithm just as nested Schur
blocks partition one finite determinant.  Additivity is exact, while
estimating two independently normalized determinants would duplicate their
shared part.  V39 proves that two-derivative and dimension-four local rows
have nonnegative scaling degree, so the YM contact firewall applies to them.
The NS Oseen-kernel rank-one norms do not enter this determinant identity.
\end{proof}

\subsection{The exact paired-unitary form}

Let
\begin{equation}
 \epsilon=\operatorname{sign}H_W,\qquad
 V=\gamma_5\epsilon,\qquad
 D_{\rm ov}=I+V .
 \label{eq:t10v40-overlap-unitary}
\end{equation}
Both \(\gamma_5\) and \(\epsilon\) are unitary and Hermitian, so \(V\) is
unitary, although it need not be Hermitian.

\begin{lemma}[Overlap square as a paired unitary]
\label{lem:t10v40-paired-unitary}
With \(T:=-V\),
\begin{equation}
 A_{\rm ov}:=D_{\rm ov}^*D_{\rm ov}
 =2I+V+V^*=2I-T-T^*,
 \qquad T^*T=TT^*=I.
 \label{eq:t10v40-paired-unitary}
\end{equation}
In particular \(A_{\rm ov}\succeq0\).
\end{lemma}

\begin{proof}
Expand \((I+V^*)(I+V)\) and use \(V^*V=I\).  For any vector \(u\),
\[
 \langle u,(2I-T-T^*)u\rangle
 =\|u-Tu\|^2\geq0.
\]
\end{proof}

This identity is stronger than generic exponential locality.  The first
weighted commutator is anti-Hermitian and therefore does not enter the real
growth rate of the heat semigroup.

\subsection{A quadratic Davies lemma}

Let \(\Phi\) be a bounded real diagonal multiplication operator and define
\begin{equation}
 A_\mu:=e^{\mu\Phi}Ae^{-\mu\Phi}.
 \label{eq:t10v40-weighted-operator}
\end{equation}
Assume \(A=A^*\succeq0\) and that, for \(0\leq s\mu\leq\mu_0\),
\begin{equation}
 \left\|
 e^{s\mu\Phi}[\Phi,[\Phi,A]]e^{-s\mu\Phi}
 \right\|\leq C_2 .
 \label{eq:t10v40-double-commutator}
\end{equation}

\begin{lemma}[Quadratic numerical-range bound]
\label{lem:t10v40-davies}
Under \eqref{eq:t10v40-double-commutator},
\begin{equation}
 \|e^{-tA_\mu}\|
 \leq\exp\left\{\frac12C_2\mu^2t\right\}.
 \label{eq:t10v40-davies}
\end{equation}
\end{lemma}

\begin{proof}
Taylor's formula with integral remainder gives
\begin{equation}
 A_\mu=A+\mu[\Phi,A]
 +\mu^2\int_0^1(1-s)
 e^{s\mu\Phi}[\Phi,[\Phi,A]]e^{-s\mu\Phi}\,ds.
 \label{eq:t10v40-conjugation-taylor}
\end{equation}
Because \(A\) and \(\Phi\) are self-adjoint, \([\Phi,A]\) is
anti-Hermitian and has zero real quadratic form.  Hence
\[
 \operatorname{Re}\langle u,A_\mu u\rangle
 \geq-\frac12C_2\mu^2\|u\|^2.
\]
The logarithmic-norm estimate for the finite-dimensional semigroup generated
by \(-A_\mu\) proves \eqref{eq:t10v40-davies}.
\end{proof}

\begin{lemma}[Exponential moment supplies the double commutator]
\label{lem:t10v40-second-moment}
Let \(A=2I-T-T^*\), and let
\(\Phi_x=d(x,z)\) with a bounded truncation.  If
\begin{equation}
 M_2(\mu_0):=
 \max\left\{
 \sup_x\sum_y d(x,y)^2e^{\mu_0d(x,y)}
  \bigl(\|T_{xy}\|+\|T^*_{xy}\|\bigr),
 \sup_y\sum_x d(x,y)^2e^{\mu_0d(x,y)}
  \bigl(\|T_{xy}\|+\|T^*_{xy}\|\bigr)
 \right\}<\infty,
 \label{eq:t10v40-second-moment}
\end{equation}
then \eqref{eq:t10v40-double-commutator} holds with
\(C_2=M_2(\mu_0)\), uniformly in \(z\), volume, and the truncation.
\end{lemma}

\begin{proof}
The \((x,y)\) block of \([\Phi,[\Phi,T]]\) is
\((\Phi_x-\Phi_y)^2T_{xy}\).  The one-Lipschitz property gives
\(|\Phi_x-\Phi_y|\leq d(x,y)\), and conjugation adds at most
\(e^{\mu_0d(x,y)}\).  Apply the row-and-column Schur test to the sum of the
\(T\) and \(T^*\) double commutators.  The diagonal \(2I\) commutes with
\(\Phi\).
\end{proof}

\begin{theorem}[Quadratic heat bound for an exponentially local unitary]
\label{thm:t10v40-unitary-heat}
Let \(T\) be unitary and satisfy
\eqref{eq:t10v40-second-moment}.  Then
\begin{equation}
 \left\|(e^{-t(2I-T-T^*)})_{xy}\right\|
 \leq
 \exp\left\{-\mu d(x,y)+\frac12M_2(\mu_0)\mu^2t\right\}
 \label{eq:t10v40-unitary-heat}
\end{equation}
for \(0\leq\mu\leq\mu_0\).
\end{theorem}

\begin{proof}
Apply Lemmas~\ref{lem:t10v40-davies} and
\ref{lem:t10v40-second-moment} with a distance weight centred at \(y\).
Taking the \((x,y)\) block of
\(e^{\mu\Phi}e^{-tA}e^{-\mu\Phi}\) supplies the factor
\(e^{\mu d(x,y)}\); move it to the other side.
\end{proof}

\subsection{Closing the overlap heat-transfer gate}

\begin{theorem}[Admissibility implies the quadratic overlap heat bound]
\label{thm:t10v40-overlap-heat}
On a V35 plaquette-admissible sector, suppose the sign-projector locality
bound holds uniformly with exponent \(\mu_0>0\).  Then the unitary
\(T=-\gamma_5\operatorname{sign}H_W\) has finite
\(M_2(\mu)\) for every \(0<\mu<\mu_0\), and
\begin{equation}
 \|(e^{-tD_{\rm ov}^*D_{\rm ov}})_{xy}\|
 \leq
 \exp\{-\mu d(x,y)+C_{\rm ov}(\mu)\mu^2t\}
 \label{eq:t10v40-overlap-heat}
\end{equation}
with constants independent of volume.
\end{theorem}

\begin{proof}
Multiplication by the on-site matrix \(\gamma_5\) does not change spatial
tails.  Exponential decay with exponent \(\mu_0\) makes the second moment
\eqref{eq:t10v40-second-moment} summable after replacing \(\mu_0\) there by
any smaller \(\mu\).  Lemma~\ref{lem:t10v40-paired-unitary} identifies the
overlap square with the operator in
Theorem~\ref{thm:t10v40-unitary-heat}.  Set
\(C_{\rm ov}=M_2/2\).
\end{proof}

\begin{corollary}[Scale-uniform overlap heat locality]
\label{cor:t10v40-scale-uniform-overlap}
With \(\ell_j=L^j\), \(t\leq\tau\ell_j^2\), and
\(\mu_j=\widehat\mu/\ell_j\),
\begin{equation}
 \|(e^{-tA_{\rm ov}})_{xy}\|
 \leq
 \exp\{-\widehat\mu d_j(x,y)+
 C_{\rm ov}\tau\widehat\mu^2\},
 \label{eq:t10v40-scale-uniform-overlap}
\end{equation}
uniformly in \(j\) and volume, provided the admissibility locality constants
are uniform.
\end{corollary}

\begin{proof}
Substitute the parabolic scaling into
\eqref{eq:t10v40-overlap-heat}; the factors \(\ell_j^2\) cancel.
\end{proof}

This closes the heat-propagation part of the V39 overlap transfer.  It does
not yet prove the scale-normalized derivative bound for variations of
\(A_{\rm ov}\).

\subsection{Derivative functional calculus}

For an invertible Hermitian \(H\),
\begin{equation}
 \operatorname{sign}H
 =\frac2\pi\int_0^\infty H(H^2+s^2)^{-1}\,ds.
 \label{eq:t10v40-sign-integral}
\end{equation}

\begin{proposition}[Gap-controlled sign derivative]
\label{prop:t10v40-sign-derivative}
If \(\operatorname{dist}(0,\operatorname{spec}H)\geq g>0\), then
\begin{equation}
 \|D(\operatorname{sign})_H[K]\|
 \leq \frac{2}{g}\|K\|.
 \label{eq:t10v40-sign-derivative}
\end{equation}
In weighted row/column norms, the same conclusion holds with a finite
constant depending on the weighted resolvent bound and \(g\).
\end{proposition}

\begin{proof}
Differentiate \eqref{eq:t10v40-sign-integral} using the inverse rule.  For
the operator-norm bound, decompose the space into the positive and negative
spectral subspaces and write \(H=H_+\oplus(-H_-)\), where
\(H_\pm\succeq gI\).  The diagonal blocks of the derivative vanish, and its
positive-to-negative block is
\[
 X=2\int_0^\infty e^{-tH_+}K_{+-}e^{-tH_-}\,dt.
\]
Therefore \(\|X\|\leq2\|K\|\int_0^\infty e^{-2gt}dt
=\|K\|/g\); the opposite block has the same bound.  This proves the stated
\(2\|K\|/g\) bound.  In the weighted algebra, differentiate the resolvent
integral and use the weighted resolvent estimate under the convergent
\(s\)-integral.
\end{proof}

\begin{corollary}[Overlap derivative bound]
\label{cor:t10v40-overlap-derivative}
For a local parameter variation,
\begin{equation}
 \|dA_{\rm ov}\|_\mu
 \leq 2\|dV\|_\mu
 \leq C(g,\mu)\|dH_W\|_\mu.
 \label{eq:t10v40-overlap-derivative}
\end{equation}
Higher derivatives obey the V20 Gevrey resolvent ledger.
\end{corollary}

\begin{proof}
Differentiate \(A_{\rm ov}=2+V+V^*\), use
\(V=\gamma_5\operatorname{sign}H_W\), and apply
Proposition~\ref{prop:t10v40-sign-derivative}.  Repeated resolvent
differentiation gives the last assertion.
\end{proof}

\begin{firewall}[Bounded derivatives are not yet scale-normalized]
\label{fire:t10v40-derivative-scaling}
Corollary~\ref{cor:t10v40-overlap-derivative} is volume uniform on an
admissible chart, but V39 requires
\(\|dA_j\|\lesssim\ell_j^{-2}\) in rescaled field coordinates.  That power
comes from how a coarse gauge, Higgs, or metric variation is embedded into
fine links.  Functional calculus alone does not supply it.
\end{firewall}

\subsection{Exact heat-shell cocycle and contact projection}

For \(0<T_0<T_1<T_2\), define the exact interval contribution
\begin{equation}
 \Gamma_{[T_a,T_b]}(D)
 :=\frac12\int_{T_a}^{T_b}\frac1t
 \left[re^{-t}-\operatorname{Tr}(e^{-tD^*D}-P_0)\right]dt.
 \label{eq:t10v40-heat-interval}
\end{equation}

\begin{lemma}[Heat determinant cocycle]
\label{lem:t10v40-heat-cocycle}
\begin{equation}
 \Gamma_{[T_0,T_2]}
 =\Gamma_{[T_0,T_1]}+\Gamma_{[T_1,T_2]}.
 \label{eq:t10v40-heat-cocycle}
\end{equation}
The identity persists after any number of background derivatives on a
constant-rank patch.
\end{lemma}

\begin{proof}
Split the same convergent integral at \(T_1\), then differentiate the
identity.
\end{proof}

This is the heat analogue of the audited YM nested Schur cocycle.  Relevant
and marginal local projections are applied to the combined interval before
its remainder is estimated.  A zero-displacement vacuum contact may vanish
under a derivative observable, but a two-derivative gauge or Einstein row
and a four-derivative curvature-squared row are retained with their exact
coefficients.

\begin{assessment}[V40 acquisition]
\label{ass:t10v40-status}
The overlap square has the exact paired-unitary structure needed for a
quadratic Davies estimate.  Plaquette admissibility now gives a
volume-independent parabolic heat bound, closing the principal locality
gap from V39.  The remaining scale gate is the embedding derivative power in
the last firewall.  V41 will define block-smooth gauge and metric variations,
prove their fine-link derivative powers, and identify which zero-momentum
holonomy directions must remain outside that estimate.
\end{assessment}

\begin{firewall}[Claim boundary after V40]
\label{fire:t10v40-claim}
V40 performs the companion audit and proves the quadratic overlap heat
bound.  It does not prove the block-embedding derivative scaling, construct
the local chiral phase, control the retained infrared determinant, close the
gravitational coupling tower, or construct the continuum.
\end{firewall}

\subsection{Parent record}

This V40 note is appended to the validated V39 source, whose record is
\begin{equation}
 \text{bytes}=562320,\qquad \text{lines}=14706,\qquad
 \operatorname{SHA256}=
 \texttt{A076E643AAC0983013332F8FF09370E785309DF8330CBE7E19401B19E52748EE}.
 \label{eq:t10v40-parent-record}
\end{equation}

% =============================================================================
% END APPENDED TENTH-SERIES V40 MATERIAL
% =============================================================================

% =============================================================================
% BEGIN APPENDED TENTH-SERIES V41 MATERIAL
% =============================================================================

\section{Tenth-series V41: block-smooth embeddings, the raw derivative
no-go, and Ward-relative scaling}
\label{sec:v10v41}

\subsection{Testing the remaining V40 gate}

V39 used the sufficient hypothesis that every derivative insertion of
\(A_{\rm ov}=D_{\rm ov}^*D_{\rm ov}\) scales as \(\ell^{-2}\).  That is true
for a scalar potential already normalized as a second-order perturbation,
but it is false for a gauge connection.  This note proves the failure and
replaces the hypothesis by a Ward-relative response bound.

Let \(\ell\) be a block diameter in fine-link units.  A block-smooth gauge
coordinate \(a_{B\mu}\) is embedded on a fine link \(e\subset B\) by
\begin{equation}
 U_e(a)=\exp\{\ell^{-1}a_{B\mu}\},
 \qquad a_{B\mu}^*=-a_{B\mu}.
 \label{eq:t10v41-block-embedding}
\end{equation}
Transition collars may interpolate neighbouring block values; their
derivatives obey the same power and are assigned to mixed-cell owners.

\begin{lemma}[Fine-link derivative power]
\label{lem:t10v41-link-derivative}
For a variation \(\dot a\),
\begin{equation}
 \|\dot U_e\|\leq\ell^{-1}\|\dot a_{B\mu}\|.
 \label{eq:t10v41-link-derivative}
\end{equation}
Consequently the Wilson-kernel derivative obeys
\begin{equation}
 \|\dot H_W\|\leq C_W\ell^{-1}\|\dot a\|_\infty,
 \label{eq:t10v41-wilson-derivative}
\end{equation}
with \(C_W\) independent of volume.
\end{lemma}

\begin{proof}
The derivative of the matrix exponential is
\[
 D(e^X)[Y]=\int_0^1e^{(1-s)X}Ye^{sX}\,ds.
\]
For anti-Hermitian \(X\), both exponential factors are unitary.  Set
\(X=\ell^{-1}a\) and \(Y=\ell^{-1}\dot a\) to obtain
\eqref{eq:t10v41-link-derivative}.  The Wilson operator contains finitely
many link multiplications per row, so the row-and-column Schur bound gives
\eqref{eq:t10v41-wilson-derivative}.
\end{proof}

\begin{corollary}[Raw overlap derivative]
\label{cor:t10v41-raw-overlap}
On a plaquette-admissible chart with Wilson-kernel gap \(g\),
\begin{equation}
 \|\dot A_{\rm ov}\|
 \leq C_{\rm ov}(g)\ell^{-1}\|\dot a\|_\infty.
 \label{eq:t10v41-raw-overlap}
\end{equation}
The same power holds in a weakened exponential row norm.
\end{corollary}

\begin{proof}
Apply the gap-controlled sign derivative and overlap derivative of V40 to
\eqref{eq:t10v41-wilson-derivative}.  Weighted functional calculus gives the
row-norm statement.
\end{proof}

\begin{theorem}[No-go for insertionwise parabolic scaling]
\label{thm:t10v41-insertion-no-go}
For a nontrivial gauge coupling, the bound
\begin{equation}
 \|\dot A_{\rm ov}\|\leq C\ell^{-2}\|\dot a\|
 \label{eq:t10v41-false-power}
\end{equation}
cannot hold uniformly for all block-smooth variations and arbitrarily large
\(\ell\).
\end{theorem}

\begin{proof}
At the flat background, choose a variation constant on a large block and a
fine spinor with momentum in a fixed ultraviolet region where the derivative
of the free overlap symbol with respect to a link phase is nonzero.  The
link phase changes by \(\ell^{-1}\dot a\), so differentiating the symbol gives
\(c\ell^{-1}\dot a+O(\ell^{-2})\) for some \(c\ne0\).  Its norm is bounded
below by \(c'\ell^{-1}\|\dot a\|\) for large \(\ell\), contradicting
\eqref{eq:t10v41-false-power}.
\end{proof}

Thus the powers of heat time cannot be cancelled before the gauge-invariant
trace and its Ward cancellations are formed.

\subsection{Exact gauge invariance of the heat shell}

Let
\begin{equation}
 \Gamma_{[T_0,T_1]}(U)
 :=\frac12\int_{T_0}^{T_1}\frac1t
 \left[re^{-t}-\operatorname{Tr}(e^{-tA_{\rm ov}(U)}-P_0(U))\right]dt
 \label{eq:t10v41-heat-shell}
\end{equation}
on a constant-rank patch.

\begin{lemma}[Finite-shell gauge invariance]
\label{lem:t10v41-shell-gauge-invariance}
For every lattice gauge transformation \(g\),
\begin{equation}
 \Gamma_{[T_0,T_1]}(U^g)=\Gamma_{[T_0,T_1]}(U).
 \label{eq:t10v41-shell-gauge-invariance}
\end{equation}
\end{lemma}

\begin{proof}
Gauge covariance gives
\(A_{\rm ov}(U^g)=G_gA_{\rm ov}(U)G_g^{-1}\) and the same conjugation law for
\(P_0\).  Functional calculus and cyclicity of the finite trace prove the
identity under the heat integral.
\end{proof}

Let \(d\) denote the linearized lattice gauge map from site parameters to
link parameters.  After the vacuum and one-point owners are projected, let
\(K\) be the Hessian of the shell at the flat background.

\begin{proposition}[Exact lattice Ward kernel]
\label{prop:t10v41-ward-kernel}
The Hessian satisfies
\begin{equation}
 Kd=0,\qquad d^*K=0.
 \label{eq:t10v41-ward-kernel}
\end{equation}
In translation-invariant coordinates this becomes
\begin{equation}
 \sum_\mu \widehat p_\mu^*K_{\mu\nu}(p)=0,\qquad
 \sum_\nu K_{\mu\nu}(p)\widehat p_\nu=0,
 \quad \widehat p_\mu=e^{ip_\mu}-1.
 \label{eq:t10v41-momentum-ward}
\end{equation}
\end{proposition}

\begin{proof}
Differentiate \eqref{eq:t10v41-shell-gauge-invariance} once in a gauge-orbit
direction and once in an arbitrary link direction.  At the projected
stationary background the term from the second derivative of the nonlinear
orbit is multiplied by the vanished first variation.  This gives \(Kd=0\).
Hermiticity of the Hessian gives the adjoint identity.  Fourier transformation
of the forward difference gives \(\widehat p_\mu\).
\end{proof}

\subsection{Contractible owners versus toron owners}

A local gauge-invariant owner supported in a contractible polymer is a
function of Wilson loops contained in that polymer.  On a constant commuting
flat connection every such loop has trivial holonomy.  Winding loops are
separate toron owners.

\begin{lemma}[Zero-momentum vanishing of a local gauge Hessian]
\label{lem:t10v41-zero-momentum}
Let \(K^{\rm loc}_{\mu\nu}(p)\) be the two-point kernel of a reflection-even,
translation-invariant, contractible gauge owner.  Then
\begin{equation}
 K^{\rm loc}(0)=0,\qquad
 K^{\rm loc}(-p)=K^{\rm loc}(p)^*.
 \label{eq:t10v41-zero-and-even}
\end{equation}
If its exponentially weighted second moment is finite, then
\begin{equation}
 \|K^{\rm loc}(p)\|\leq C_K|p|^2
 \label{eq:t10v41-quadratic-symbol}
\end{equation}
near \(p=0\).
\end{lemma}

\begin{proof}
A constant commuting link potential has identity holonomy on every
contractible loop, so the local owner is independent of it and its constant
mode Hessian vanishes.  Reflection gives the second identity.  Exponential
locality makes \(K^{\rm loc}(p)\) analytic.  Its Taylor expansion has no
constant term.  The linear term is odd and anti-Hermitian by the reflection
identity; the real even quadratic action has no such term.  Taylor's theorem
then gives \eqref{eq:t10v41-quadratic-symbol}, with \(C_K\) controlled by the
second spatial moment.
\end{proof}

\begin{remark}[Parity-odd topological row]
\label{rem:t10v41-parity-odd}
If reflection is not imposed, an allowed parity-odd local term must be
projected into the topological-density ledger rather than discarded.  The
quadratic estimate applies to the reflection-even kinetic response after
that exact projection.
\end{remark}

\begin{theorem}[Ward-relative two-point scaling]
\label{thm:t10v41-ward-relative}
Embed an external block momentum as \(p=\widehat p/\ell\), with
\(\widehat p\) in a fixed compact set.  Under the hypotheses of
Lemma~\ref{lem:t10v41-zero-momentum},
\begin{equation}
 \|K^{\rm loc}(\widehat p/\ell)\|
 \leq C_K\ell^{-2}|\widehat p|^2.
 \label{eq:t10v41-ward-relative}
\end{equation}
Thus two raw gauge insertions of size \(\ell^{-1}\) reorganize, after the
gauge-invariant trace is formed, into the second-order power required by
parabolic reblocking.
\end{theorem}

\begin{proof}
Substitute \(p=\widehat p/\ell\) into
\eqref{eq:t10v41-quadratic-symbol}.  The Ward identities
\eqref{eq:t10v41-momentum-ward} ensure that the resulting quadratic tensor
acts only on the transverse curvature channel; longitudinal terms vanish.
\end{proof}

\begin{firewall}[Torons are not controlled by the \(p^2\) factor]
\label{fire:t10v41-toron}
A harmonic link variation has \(p=0\) but can change a winding Wilson loop.
It is absent from contractible owners and belongs to the retained toron
block of V35.  Applying \eqref{eq:t10v41-ward-relative} to a winding owner
would erase genuine finite-volume holonomy dependence.
\end{firewall}

\subsection{Higher responses and the covariant local form}

The two-point theorem is exact and sufficient for the gauge kinetic row.
For higher responses the correct statement is not a separate
\(\ell^{-2}\) per insertion.  Gauge covariance groups the Taylor series into
local invariants built from
\begin{equation}
 F=dA+A\wedge A,\qquad D_A F,\qquad D_A^2F,\ldots .
 \label{eq:t10v41-curvature-generators}
\end{equation}

\begin{proposition}[Power after curvature compilation]
\label{prop:t10v41-curvature-power}
Assume a contractible shell owner admits an absolutely convergent local
expansion in the generators \eqref{eq:t10v41-curvature-generators}.  Under
the block embedding \eqref{eq:t10v41-block-embedding},
\begin{equation}
 F_{\rm fine}=O(\ell^{-2}),\qquad
 D_{\rm fine}^mF_{\rm fine}=O(\ell^{-2-m}).
 \label{eq:t10v41-curvature-power}
\end{equation}
Every monomial then has exactly its canonical derivative power, and the
dimension-at-most-four monomials are the projected gauge rows of V39.
\end{proposition}

\begin{proof}
A fine difference of a block-smooth potential contributes one factor
\(\ell^{-1}\), in addition to the \(\ell^{-1}\) size of the link potential.
The commutator of two potentials also has two such factors.  Each further
covariant derivative contributes one more fine difference or one commutator
with the \(O(\ell^{-1})\) connection.  Multiplication proves the monomial
statement.
\end{proof}

\begin{firewall}[The curvature compilation is still an acquisition]
\label{fire:t10v41-covariant-poincare}
Gauge invariance and exponential locality strongly constrain a contractible
owner, but Proposition~\ref{prop:t10v41-curvature-power} assumes the
convergent covariant expansion rather than proving it.  The needed result is
a quantitative lattice covariant Poincare lemma with polymer support and
analytic/Gevrey bounds.
\end{firewall}

\subsection{Metric and Higgs rows}

A block-smooth metric perturbation changes the principal symbol at order one,
so its raw operator derivative also fails the insertionwise V39 hypothesis.
Diffeomorphism invariance restores canonical powers only after the vacuum,
Einstein, and curvature-squared rows are projected.  Similarly, constant
Higgs variations generate the relevant quadratic and marginal potential
rows and cannot acquire derivative powers from a Ward identity.

\begin{proposition}[Revised scale gate]
\label{prop:t10v41-revised-gate}
The insertionwise condition \eqref{eq:t10v39-derivative-scaling} must be
replaced by the following response-level condition at every heat shell:
\begin{enumerate}
\item form the complete gauge-, spin-, and diffeomorphism-covariant trace;
\item separate toron, physical zero-mode, and other retained infrared blocks;
\item project exactly every local invariant of dimension at most four listed
      in V39; and
\item bound the connected remainder in block distance with its canonical
      negative scaling degree.
\end{enumerate}
These conditions are sufficient for the V38 forest gate when its tree
smallness holds.
\end{proposition}

\begin{proof}
Steps 1--2 make the Ward identities valid on the local complement.  Step 3
removes every noncontracting power identified by the heat scaling theorem.
Step 4 supplies the decaying activities and fixed block-distance exponent
used in V38.  The forest theorem then sums the connected owners under its
stated smallness condition.
\end{proof}

\begin{assessment}[V41 acquisition]
\label{ass:t10v41-status}
The apparent derivative bottleneck was caused by imposing a false
insertionwise power.  The raw gauge derivative is only \(O(\ell^{-1})\), but
the exact heat trace obeys Ward identities and its local even two-point
kernel is \(O(p^2)\), restoring the parabolic scale after compilation.
The next upstream task is the quantitative covariant Poincare lemma named in
the firewall.  V42 will prove its abelian lattice version with explicit
support enlargement and norm loss, then isolate the nonabelian commutator
iteration.
\end{assessment}

\begin{firewall}[Claim boundary after V41]
\label{fire:t10v41-claim}
V41 proves the raw derivative no-go and the local two-point Ward recovery.
It does not prove the all-order covariant local expansion, control winding
owners, establish the corresponding diffeomorphism theorem, close the
infrared determinant, or construct the continuum.
\end{firewall}

\subsection{Parent record}

This V41 note is appended to the validated V40 source, whose record is
\begin{equation}
 \text{bytes}=575594,\qquad \text{lines}=15069,\qquad
 \operatorname{SHA256}=
 \texttt{70BCB9DF7CAAEF85CEC0FCF6AD76E5B94D88691374DD64E2ECECCBC8B0EAA6C3}.
 \label{eq:t10v41-parent-record}
\end{equation}

% =============================================================================
% END APPENDED TENTH-SERIES V41 MATERIAL
% =============================================================================

% =============================================================================
% BEGIN APPENDED TENTH-SERIES V42 MATERIAL
% =============================================================================

\section{Tenth-series V42: a local lattice Poincare compiler and the
nonabelian curvature fixed point}
\label{sec:v10v42}

\subsection{The contractible cubical homotopy}

V41 reduced the all-order gauge scaling problem to a quantitative covariant
Poincare lemma.  Begin on a finite rectangular cubical complex \(B\) of
diameter \(r\), with real or Lie-algebra-valued cochains and coboundary \(d\).

\begin{lemma}[Cubical contraction]
\label{lem:t10v42-cubical-contraction}
There are linear maps
\begin{equation}
 h_k:C^k(B)\longrightarrow C^{k-1}(B),\qquad k\geq1,
 \label{eq:t10v42-homotopy-map}
\end{equation}
such that
\begin{equation}
 dh_k+h_{k+1}d=I\quad(k\geq1)
 \label{eq:t10v42-homotopy-identity}
\end{equation}
and
\begin{equation}
 \|h_k\omega\|_\infty\leq C_d r\,\|\omega\|_\infty.
 \label{eq:t10v42-homotopy-bound}
\end{equation}
The output is supported in the smallest coordinate box joining the support
of \(\omega\) to a fixed base face of \(B\).
\end{lemma}

\begin{proof}
Contract the box one coordinate at a time.  In one coordinate, the homotopy
is the partial sum
\[
 (h\omega)(x)=\sum_{s=x_0}^{x-1}\iota_{\partial_s}\omega(s),
\]
with the orientation sign dictated by the position of that coordinate.
The elementary telescoping identity is \(dh+hd=I-\pi\), where \(\pi\)
restricts to the base face.  Compose these contractions through all
coordinates.  On positive-degree cochains the final restriction to the
basepoint vanishes, giving \eqref{eq:t10v42-homotopy-identity}.  Each
coordinate partial sum contains at most \(r\) terms, and there are at most
\(d\) coordinate stages; their fixed combinatorial multiplicity gives
\(C_dr\).  Every term lies in the coordinate box described in the statement.
\end{proof}

For a one-cochain \(a\), \eqref{eq:t10v42-homotopy-identity} reads
\begin{equation}
 a=h_2(da)+d(h_1a).
 \label{eq:t10v42-abelian-decomposition}
\end{equation}
Thus \(h_1a=0\) is a concrete axial-type slice and \(h_2da\) is its
curvature representative.

\subsection{Exact abelian curvature compilation}

Let \(\mathcal F(a)\) be an analytic local owner depending on links in \(B\)
and invariant under
\begin{equation}
 a\longmapsto a+d\phi.
 \label{eq:t10v42-abelian-gauge}
\end{equation}

\begin{theorem}[Abelian local factorization]
\label{thm:t10v42-abelian-factorization}
On a contractible logarithm chart,
\begin{equation}
 \mathcal F(a)=\widetilde{\mathcal F}(f),
 \qquad f:=da,\qquad
 \widetilde{\mathcal F}(f):=\mathcal F(h_2f).
 \label{eq:t10v42-abelian-factorization}
\end{equation}
If
\begin{equation}
 \|D^n\mathcal F(a)\|
 \leq M n!R^{-n}\quad\text{for }\|a\|_\infty<R,
 \label{eq:t10v42-a-Cauchy}
\end{equation}
then
\begin{equation}
 \|D^n\widetilde{\mathcal F}(f)\|
 \leq M n!\left(\frac{C_dr}{R}\right)^n
 \label{eq:t10v42-f-Cauchy}
\end{equation}
for \(\|f\|_\infty<R/(C_dr)\).
\end{theorem}

\begin{proof}
Equation \eqref{eq:t10v42-abelian-decomposition} writes
\(a=h_2f+d\phi\) with \(\phi=h_1a\).  Gauge invariance proves
\eqref{eq:t10v42-abelian-factorization}.  Its \(n\)-th derivative is
\(D^n\mathcal F(h_2f)[h_2\dot f_1,\ldots,h_2\dot f_n]\).
Apply \eqref{eq:t10v42-a-Cauchy} and
\eqref{eq:t10v42-homotopy-bound}.
\end{proof}

For compact \(U(1)\) links, choose principal link and plaquette logarithms on
a patch with no branch crossing.  Then the plaquette logarithm is exactly
\(da\), and the theorem applies.  Nonzero flux and toron variables are not in
one such contractible chart and remain in the V35 retained ledger.

\begin{corollary}[Canonical block power]
\label{cor:t10v42-abelian-power}
If \(r\asymp\ell\) and a block-smooth curvature satisfies
\(\|f\|_\infty=O(\ell^{-2})\), then
\begin{equation}
 \|h_2f\|_\infty=O(\ell^{-1}).
 \label{eq:t10v42-abelian-power}
\end{equation}
The rescaled curvature variable \(\widehat f=\ell^2f\) has an analytic radius
independent of \(\ell\) whenever the link-coordinate radius \(R\) is fixed.
\end{corollary}

\begin{proof}
Use \eqref{eq:t10v42-homotopy-bound}.  In
\eqref{eq:t10v42-f-Cauchy}, each derivative with respect to
\(\widehat f\) contributes \(\ell^{-2}\), which cancels the
\(r\asymp\ell\) loss with one remaining canonical \(\ell^{-1}\) link power.
\end{proof}

\subsection{Polymer support and reserve}

If an owner has support \(X\), choose a contractible coordinate box
\(B_X\) containing \(X\), and denote its diameter by \(r_X\).  The homotopy
may enlarge \(X\) to \(B_X\), but no farther.

\begin{proposition}[Box-allocation bound]
\label{prop:t10v42-box-allocation}
Assume the boxes satisfy
\begin{equation}
 |B_X|\leq c_B|X|,\qquad r_X\leq c_r|X|.
 \label{eq:t10v42-box-geometry}
\end{equation}
For \(0\leq\kappa'<\kappa/c_B\),
\begin{equation}
 e^{\kappa'|B_X|}r_X
 \leq c_r\sup_{m\geq1}m
 e^{-(\kappa-\kappa'c_B)m}\,e^{\kappa|X|}.
 \label{eq:t10v42-box-reserve}
\end{equation}
Thus the support enlargement and the one homotopy length are absorbed by a
fixed positive polymer reserve.
\end{proposition}

\begin{proof}
Insert \eqref{eq:t10v42-box-geometry} into the left side, multiply and divide
by \(e^{\kappa|X|}\), and take the displayed supremum.  It is finite because
\(\kappa-\kappa'c_B>0\).
\end{proof}

Thin or multiply connected owners must first be filled by a box or a bounded
contractible hull.  If the hull-to-owner volume ratio is not uniform, this
construction is not a polymer-local proof.

\subsection{The nonabelian logarithmic curvature map}

On a small compact-group logarithm chart, the logarithm of plaquette
holonomy has the analytic form
\begin{equation}
 \mathcal K(a)=da+Q(a),
 \qquad Q(0)=0,\qquad DQ(0)=0.
 \label{eq:t10v42-nonabelian-curvature}
\end{equation}
The Baker--Campbell--Hausdorff series supplies \(Q\).  On a ball
\(\|a\|_\infty\leq R_0\), assume
\begin{align}
 \|Q(a)\|_\infty&\leq C_Q\|a\|_\infty^2,
 \label{eq:t10v42-Q-quadratic}\\
 \|Q(a)-Q(b)\|_\infty
 &\leq C_Q(\|a\|_\infty+\|b\|_\infty)\|a-b\|_\infty.
 \label{eq:t10v42-Q-Lipschitz}
\end{align}
These are finite-dimensional Cauchy bounds for the BCH map.

In the slice \(h_1a=0\), the homotopy identity and
\(f=\mathcal K(a)\) give the fixed-point equation
\begin{equation}
 a=h_2(f-Q(a)).
 \label{eq:t10v42-curvature-fixed-point}
\end{equation}
Put \(H=\|h_2\|_{\infty\leftarrow\infty}\leq C_dr\).

\begin{theorem}[Small-curvature nonabelian reconstruction]
\label{thm:t10v42-nonabelian-reconstruction}
Suppose
\begin{equation}
 4H^2C_Q\|f\|_\infty<1,\qquad
 2H\|f\|_\infty<R_0.
 \label{eq:t10v42-nonabelian-smallness}
\end{equation}
Then \eqref{eq:t10v42-curvature-fixed-point} has a unique solution in
\begin{equation}
 \|a\|_\infty\leq2H\|f\|_\infty.
 \label{eq:t10v42-a-bound}
\end{equation}
The solution is analytic in \(f\), uses only cells in the same contractible
box, and is equivariant under transformations preserving the chosen slice.
\end{theorem}

\begin{proof}
On the ball of radius \(R=2H\|f\|_\infty\), define
\(\mathcal T_f(a)=h_2(f-Q(a))\).  By
\eqref{eq:t10v42-Q-quadratic},
\[
 \|\mathcal T_f(a)\|
 \leq H\|f\|+HC_QR^2\leq R
\]
under \eqref{eq:t10v42-nonabelian-smallness}.  By
\eqref{eq:t10v42-Q-Lipschitz},
\[
 \|\mathcal T_f(a)-\mathcal T_f(b)\|
 \leq2HC_QR\|a-b\|<\|a-b\|.
\]
Banach's fixed-point theorem gives existence and uniqueness.  Analytic
dependence follows either by uniform convergence of the analytic iteration
or by the analytic implicit-function theorem, since
\(I+h_2DQ(a)\) is invertible by the same contraction bound.  The homotopy and
\(Q\) never use cells outside the box, proving support.  Naturality and
uniqueness give equivariance under slice-preserving transformations.
\end{proof}

\begin{corollary}[Nonabelian factorization on a local slice]
\label{cor:t10v42-nonabelian-factorization}
Let \(\mathcal F(a)\) be analytic and gauge invariant, and assume every small
orbit in the patch meets the slice \(h_1a=0\) uniquely.  Under
\eqref{eq:t10v42-nonabelian-smallness},
\begin{equation}
 \mathcal F(a)=\widehat{\mathcal F}(\mathcal K(a)),
 \qquad
 \widehat{\mathcal F}(f):=\mathcal F(a(f)).
 \label{eq:t10v42-nonabelian-factorization}
\end{equation}
Its analytic radius and polymer support obey the losses in
Theorem~\ref{thm:t10v42-nonabelian-reconstruction} and
Proposition~\ref{prop:t10v42-box-allocation}.
\end{corollary}

\begin{proof}
Gauge-transform \(a\) to its unique slice representative.  That
representative is the unique reconstructed \(a(f)\), so gauge invariance
gives the formula.  Analyticity and support follow from the theorem.
\end{proof}

\begin{firewall}[The smallness scales with the patch diameter]
\label{fire:t10v42-diameter-smallness}
Since \(H=O(r)\), condition
\eqref{eq:t10v42-nonabelian-smallness} requires
\(\|f\|=O(r^{-2})\).  This is exactly the block-smooth curvature power, but
it is not a uniform statement for arbitrary admissible plaquette fields on
arbitrarily large boxes.  Large configurations must be covered by bounded
patches and reblocked; a single global axial gauge would be nonlocal.
\end{firewall}

\subsection{Insertion into the heat-shell programme}

\begin{theorem}[Local all-order gauge compiler]
\label{thm:t10v42-gauge-compiler}
Let a contractible finite-heat owner satisfy the V38 analytic polymer gate
in link coordinates, be exactly gauge invariant, and lie in a curvature
patch obeying \eqref{eq:t10v42-nonabelian-smallness}.  Then it has an analytic
curvature-owner representation
\begin{equation}
 \mathcal F_X(U)=\widehat{\mathcal F}_{B_X}
 \bigl(\log U(p):p\subset B_X\bigr)
 \label{eq:t10v42-owner-curvature}
\end{equation}
with explicit homotopy-radius and reserve losses.  Under block-smooth
rescaling, every curvature and covariant-difference insertion has the
canonical power used in V41.
\end{theorem}

\begin{proof}
Use the logarithmic curvature map and
Corollary~\ref{cor:t10v42-nonabelian-factorization} on \(B_X\).
Proposition~\ref{prop:t10v42-box-allocation} controls the support and
polynomial homotopy factor.  The fixed-point equation expresses every link
coordinate as an analytic series in curvature with the bound
\eqref{eq:t10v42-a-bound}.  Differences of curvature arise from comparing
neighbouring cells, so the V41 canonical-power argument applies term by
term; absolute convergence permits regrouping.
\end{proof}

\begin{firewall}[Overlap and orbit compatibility]
\label{fire:t10v42-overlap-compatibility}
Different boxes use different axial slices.  V35 proves how their positive
measures compare and identifies the determinant phase
\(\det R/|\det R|\).  Theorem~\ref{thm:t10v42-gauge-compiler} does not prove
that these overlap phases form a local coboundary.  It supplies local
curvature owners to which the V29 Čech test can now be applied.
\end{firewall}

\begin{assessment}[V42 acquisition]
\label{ass:t10v42-status}
The abelian curvature compilation is exact, and the nonabelian compiler is
proved by a quantitative fixed point on each bounded contractible patch.
This closes the all-order local gauge-form assumption of V41 in the
small-curvature regime, with honest diameter and reserve costs.  V43 will
apply the same strategy to the metric sector: separate frame/diffeomorphism
directions, reconstruct a local connection from curvature in a normal
slice, and determine where the Einstein and curvature-squared rows consume
the available radius.
\end{assessment}

\begin{firewall}[Claim boundary after V42]
\label{fire:t10v42-claim}
V42 proves a local curvature compiler.  It does not build a global gauge
slice, control torons or nonzero flux in one chart, trivialize the chiral
overlap cocycle, prove the metric analogue, or construct the continuum.
\end{firewall}

\subsection{Parent record}

This V42 note is appended to the validated V41 source, whose record is
\begin{equation}
 \text{bytes}=588848,\qquad \text{lines}=15402,\qquad
 \operatorname{SHA256}=
 \texttt{C04DC27A93EF8E78FF9106285334F4F4C442A4D271896089AF1F3E596CD4401E}.
 \label{eq:t10v42-parent-record}
\end{equation}

% =============================================================================
% END APPENDED TENTH-SERIES V42 MATERIAL
% =============================================================================

% =============================================================================
% BEGIN APPENDED TENTH-SERIES V43 MATERIAL
% =============================================================================

\section{Tenth-series V43: a local DeTurck curvature chart for metric
owners}
\label{sec:v10v43}

\subsection{Relation to the earlier gravitational blocks}

V13--V14 already solve the curved symmetric-divergence repair modulo
approximate Killing data, test when that repair is variational, and quantify
the normal-coordinate curvature debit in stress moment cancellation.
V17--V19 already place the unbounded gravitational counterterm tower in a
graded analytic/Gevrey space.  The missing statement is different: a local
diffeomorphism-invariant heat owner must be written in a second-order metric
coordinate before its canonical block power can be read.  This note builds
that coordinate on a bounded metric patch.

Let \(B_r\subset\mathbb R^d\) be a fixed-shape smooth box of diameter \(r\),
let \(g_0\) be a smooth uniformly elliptic reference metric, and write
\(g=g_0+h\).  Boundary metric data, corner data, and the finite
approximate-Killing ledger are retained.  The eliminated interior
perturbation obeys homogeneous Dirichlet data.

\subsection{The gauge-fixed curvature map}

At a flat reference define the linear Bianchi operator
\begin{equation}
 \beta(h)_j:=\partial^i h_{ij}-\frac12\partial_j\operatorname{tr}h
 \label{eq:t10v43-bianchi-operator}
\end{equation}
and the symmetrized gradient
\begin{equation}
 (\delta^*\omega)_{ij}:=\frac12(\partial_i\omega_j+\partial_j\omega_i).
 \label{eq:t10v43-symmetric-gradient}
\end{equation}
Choose the sign of the DeTurck map so that
\begin{equation}
 \mathcal E(h):=\operatorname{Ric}(g_0+h)
 -\delta_{g_0+h}^*\beta_{g_0}(g_0+h)
 \label{eq:t10v43-deturck-map}
\end{equation}
has flat linearization
\begin{equation}
 D\mathcal E(0)h=-\frac12\Delta h=:Lh .
 \label{eq:t10v43-deturck-linearization}
\end{equation}
For a curved reference the principal part is the same and the lower terms
form the Lichnerowicz operator.

\begin{lemma}[Flat DeTurck cancellation]
\label{lem:t10v43-deturck-cancellation}
With the conventions above, the second-derivative gauge terms in the
linearized Ricci tensor cancel those in \(\delta^*\beta\), leaving
\eqref{eq:t10v43-deturck-linearization}.
\end{lemma}

\begin{proof}
At the Euclidean metric,
\[
 (D\operatorname{Ric}(0)h)_{ij}
 =\frac12\left(
 \partial_i\partial^kh_{kj}
 +\partial_j\partial^kh_{ki}
 -\Delta h_{ij}
 -\partial_i\partial_j\operatorname{tr}h\right).
\]
Using \eqref{eq:t10v43-bianchi-operator} in
\eqref{eq:t10v43-symmetric-gradient}, the first, second, and fourth terms are
exactly \((\delta^*\beta(h))_{ij}\).  Subtraction leaves
\(-\Delta h_{ij}/2\).
\end{proof}

\begin{lemma}[Scaled Dirichlet inverse]
\label{lem:t10v43-dirichlet-inverse}
On \(C^{2,\alpha}_0(B_r;S^2)\), the flat operator \(L=-\Delta/2\) is
invertible and
\begin{align}
 \|L^{-1}s\|_{C^{2,\alpha}_r}
 &\leq C_E r^2\|s\|_{C^{0,\alpha}_r},
 \label{eq:t10v43-dirichlet-Schauder}\\
 \|u\|_{C^{2,\alpha}_r}
 &:=\|u\|_\infty+r\|\partial u\|_\infty
   +r^2\|\partial^2u\|_{C^{0,\alpha}_r}.
 \label{eq:t10v43-scaled-holder}
\end{align}
The constant depends only on the rescaled box shape, \(d\), and \(\alpha\).
A sufficiently small \(C^{2,\alpha}\) perturbation of \(g_0\) preserves the
inverse with a uniformly enlarged constant.
\end{lemma}

\begin{proof}
Rescale \(x=ry\).  The equation \(Lu=s\) becomes
\(-\Delta_yu/2=r^2s(ry)\) on the unit reference box.  The Dirichlet Schauder
estimate gives \eqref{eq:t10v43-dirichlet-Schauder}.  The maximum principle
controls the zeroth-order term.  A curved Lichnerowicz operator differs from
the flat one by a bounded map
\(C^{2,\alpha}\to C^{0,\alpha}\); the Neumann-series criterion preserves
invertibility when the rescaled perturbation norm is smaller than
\(C_E^{-1}\).
\end{proof}

The Dirichlet inverse is an interior complement inverse.  Without boundary
conditions, Killing, affine, and propagating low modes must be retained
rather than divided out.

\subsection{Analytic nonlinear reconstruction}

Write
\begin{equation}
 \mathcal E(h)=Lh+Q(h),\qquad Q(0)=0,\qquad DQ(0)=0.
 \label{eq:t10v43-nonlinear-split}
\end{equation}
On a uniformly elliptic ball
\(\|h\|_{C^{2,\alpha}_r}\leq R_0\), analyticity of matrix inversion and of
the Christoffel/Ricci formula gives
\begin{align}
 \|Q(h)\|_{C^{0,\alpha}_r}
 &\leq C_Qr^{-2}\|h\|_{C^{2,\alpha}_r}^2,
 \label{eq:t10v43-Q-quadratic}\\
 \|Q(h)-Q(k)\|_{C^{0,\alpha}_r}
 &\leq C_Qr^{-2}
 (\|h\|_{C^{2,\alpha}_r}+\|k\|_{C^{2,\alpha}_r})
 \|h-k\|_{C^{2,\alpha}_r}.
 \label{eq:t10v43-Q-Lipschitz}
\end{align}

\begin{theorem}[Local DeTurck curvature chart]
\label{thm:t10v43-deturck-chart}
Let \(G=L^{-1}\), with the scaled bound in
\eqref{eq:t10v43-dirichlet-Schauder}.  If
\begin{equation}
 4C_E^2C_Q r^2\|s\|_{C^{0,\alpha}_r}<1,\qquad
 2C_Er^2\|s\|_{C^{0,\alpha}_r}<R_0,
 \label{eq:t10v43-source-smallness}
\end{equation}
then
\begin{equation}
 \mathcal E(h)=s,\qquad h|_{\partial B_r}=0
 \label{eq:t10v43-deturck-equation}
\end{equation}
has a unique solution in the ball
\begin{equation}
 \|h\|_{C^{2,\alpha}_r}
 \leq2C_Er^2\|s\|_{C^{0,\alpha}_r}.
 \label{eq:t10v43-metric-bound}
\end{equation}
The solution \(h=h(s)\) is analytic.
\end{theorem}

\begin{proof}
Equation \eqref{eq:t10v43-deturck-equation} is the fixed-point problem
\[
 h=\mathcal T_s(h):=G(s-Q(h)).
\]
Put \(R=2C_Er^2\|s\|\).  From
\eqref{eq:t10v43-dirichlet-Schauder} and
\eqref{eq:t10v43-Q-quadratic},
\[
 \|\mathcal T_s(h)\|
 \leq C_Er^2\|s\|+C_EC_Q\|h\|^2
 \leq R
\]
under \eqref{eq:t10v43-source-smallness}.  The Lipschitz constant on that
ball is at most \(2C_EC_QR<1\).  Banach's fixed-point theorem proves
existence and uniqueness.  The analytic implicit-function theorem applies
because \(I+G\,DQ(h)\) is invertible by the same contraction estimate.
\end{proof}

\begin{corollary}[Canonical metric power]
\label{cor:t10v43-metric-power}
For a block of diameter \(r\asymp\ell\), a source
\(s=O(\ell^{-2})\) gives \(h=O(1)\) in the scaled metric norm.  The rescaled
source
\begin{equation}
 \widehat s:=\ell^2s
 \label{eq:t10v43-rescaled-source}
\end{equation}
has an analytic radius independent of \(\ell\) whenever the rescaled
ellipticity and inverse constants are uniform.
\end{corollary}

\begin{proof}
Insert \(r\asymp\ell\) in
\eqref{eq:t10v43-source-smallness}--\eqref{eq:t10v43-metric-bound}.  All
powers occur through \(r^2s\), which is comparable to \(\widehat s\).
\end{proof}

\subsection{The local diffeomorphism slice}

The DeTurck map fixes the principal diffeomorphism degeneracy, but a physical
metric must first be moved to its Bianchi slice.  Let \(X|_{\partial B_r}=0\)
and define
\begin{equation}
 \mathcal B_g(X):=\beta_{g_0}((\exp X)^*g).
 \label{eq:t10v43-slice-map}
\end{equation}

\begin{lemma}[Linear slice operator]
\label{lem:t10v43-slice-linearization}
At \(g=g_0\) and \(X=0\), the derivative of
\(\mathcal B_{g_0}\) is a strongly elliptic vector Laplacian, with principal
part
\begin{equation}
 D_X\mathcal B_{g_0}(0)X=\Delta X
 \label{eq:t10v43-vector-laplacian}
\end{equation}
up to the fixed sign and curvature convention.  On the Dirichlet complement
it has an inverse of scaled norm \(O(r^2)\).
\end{lemma}

\begin{proof}
The infinitesimal metric change is
\(\mathcal L_Xg_0=2\delta^*X^\flat\).  Applying the Bianchi operator gives
the vector Laplacian plus the Ricci endomorphism; at a flat reference this
is \(\Delta X\).  Dirichlet elliptic theory and the rescaling argument of
Lemma~\ref{lem:t10v43-dirichlet-inverse} give the inverse bound.
\end{proof}

\begin{theorem}[Local Bianchi slice modulo retained data]
\label{thm:t10v43-bianchi-slice}
For a sufficiently small metric perturbation with fixed boundary data,
there is a unique small Dirichlet vector field \(X\) such that
\begin{equation}
 \beta_{g_0}((\exp X)^*g)=0.
 \label{eq:t10v43-bianchi-slice}
\end{equation}
If boundary diffeomorphisms or Killing fields are not fixed, their finite
coordinates must be retained and uniqueness holds only on the V13 Korn
complement.
\end{theorem}

\begin{proof}
Apply the analytic implicit-function theorem to
\eqref{eq:t10v43-slice-map} using
Lemma~\ref{lem:t10v43-slice-linearization}.  The Dirichlet inverse gives
existence and uniqueness on the complement.  If its kernel is restored, the
Fredholm alternative gives precisely the finite Killing and boundary
cokernel recorded in V13.
\end{proof}

This theorem is local in metric configuration space and within one bounded
spacetime box.  It does not assert a global slice on the full space of
metrics.

\subsection{Factorization of invariant metric owners}

Let \(\mathcal F(g,\Phi)\) be an analytic owner supported in \(B_r\),
invariant under diffeomorphisms that equal the identity on the boundary.
Here \(\Phi\) denotes matter fields transported with the diffeomorphism.

\begin{theorem}[Metric-source compiler]
\label{thm:t10v43-metric-compiler}
Under the smallness conditions of
Theorems~\ref{thm:t10v43-deturck-chart} and
\ref{thm:t10v43-bianchi-slice},
\begin{equation}
 \mathcal F(g,\Phi)
 =\widehat{\mathcal F}\bigl(s,\Phi_{\rm sl},
 g|_{\partial B_r},\mathfrak k\bigr),
 \qquad
 s=\mathcal E(h),
 \label{eq:t10v43-metric-factorization}
\end{equation}
where \(\Phi_{\rm sl}\) is the sliced matter field and
\(\mathfrak k\) denotes retained Killing/boundary coordinates.  The
dependence on \(s\) is analytic, and its \(n\)-th derivative loses at most
\((C_Er^2)^n\) relative to the metric-coordinate Cauchy bound.
\end{theorem}

\begin{proof}
Use Theorem~\ref{thm:t10v43-bianchi-slice} to choose the unique complement
representative.  Diffeomorphism invariance leaves the owner unchanged.
On that representative, Theorem~\ref{thm:t10v43-deturck-chart} gives the
analytic inverse \(h(s)\).  Substitute it into \(\mathcal F\).  Repeated
chain rule and the inverse bound
\(\|Dh(0)\|\leq C_Er^2\), together with the analytic fixed-point estimates,
give the stated radius loss.  Boundary and kernel data are arguments rather
than inverted variables.
\end{proof}

\begin{proposition}[Metric support and reserve debit]
\label{prop:t10v43-metric-reserve}
Suppose an owner support \(X\) is enlarged to a fixed-shape box \(B_X\) with
\begin{equation}
 |B_X|\leq c_B|X|,\qquad r_X\leq c_r|X|.
 \label{eq:t10v43-metric-box}
\end{equation}
Then the \(r_X^2\) inverse factor is absorbed by any reserve
\(\kappa-\kappa'c_B>0\):
\begin{equation}
 e^{\kappa'|B_X|}r_X^2
 \leq c_r^2
 \sup_{m\geq1}m^2e^{-(\kappa-\kappa'c_B)m}
 e^{\kappa|X|}.
 \label{eq:t10v43-metric-reserve}
\end{equation}
\end{proposition}

\begin{proof}
Use \eqref{eq:t10v43-metric-box}, insert \(e^{\kappa|X|}\), and take the
finite polynomial-exponential supremum.
\end{proof}

The Dirichlet Green operator fills the box and is not a finite-range map.
The theorem is polymer local only because the boxes have uniformly
controlled enlargement and are reblocked after each use.

\subsection{Projection of the gravitational heat rows}

On the Bianchi slice, \(s=\operatorname{Ric}(g)\).  The full Riemann tensor
is an analytic local differential expression in \(h(s)\) and the retained
boundary data.  The V39 heat rows therefore become analytic functions of
the rescaled source \(\widehat s\).

\begin{theorem}[Gravitational response-level scaling]
\label{thm:t10v43-gravity-scaling}
Assume a finite-heat owner satisfies the V38 marked trace bounds in metric
coordinates and the hypotheses of
Theorem~\ref{thm:t10v43-metric-compiler}.  After exact projection of
\begin{equation}
 \int\sqrt g,\qquad
 \int\sqrt g\,R,\qquad
 \int\sqrt g\,R^2,\qquad
 \int\sqrt g\,R_{\mu\nu}R^{\mu\nu},
 \label{eq:t10v43-gravity-projections}
\end{equation}
together with the topological curvature density, its remaining local metric
Taylor monomials have negative canonical scaling degree at the
four-dimensional Gaussian chart.  Their support and analytic losses are
those of Proposition~\ref{prop:t10v43-metric-reserve}.
\end{theorem}

\begin{proof}
The source \(s\) has two-derivative dimension and the reconstruction depends
analytically on \(r^2s\).  The constant, linear, and quadratic source
monomials give precisely the dimension-zero, dimension-two, and
dimension-four invariant rows in
\eqref{eq:t10v43-gravity-projections}, modulo Bianchi identities,
integrations by parts, and the topological density.  Every remaining
monomial has dimension greater than four and therefore contracts under
canonical reblocking.  The metric compiler and
Proposition~\ref{prop:t10v43-metric-reserve} provide the analytic and support
bounds.
\end{proof}

\begin{firewall}[Gaussian scaling does not close the Einstein ultraviolet]
\label{fire:t10v43-einstein-uv}
Theorem~\ref{thm:t10v43-gravity-scaling} is a local finite-shell statement.
V17 proves that higher gravity loops generate unbounded derivative order.
A full ultraviolet construction still needs the V18--V19 infinite graded
tail and a fixed-point or contraction theorem on it; projecting the four
rows in \eqref{eq:t10v43-gravity-projections} is not such a theorem.
\end{firewall}

\subsection{Reflection and conformal-factor checks}

If the box, boundary conditions, reference metric, and DeTurck operator are
reflection invariant, uniqueness gives
\begin{equation}
 h(\Theta s)=\Theta h(s).
 \label{eq:t10v43-reflection}
\end{equation}
Thus reflected metric owners compile equivariantly.  The determinant of the
slice operator and the Dirichlet elliptic determinant must still be placed
in the measured quotient ledger of V25 and V35.

\begin{firewall}[Coordinate reconstruction does not create coercivity]
\label{fire:t10v43-conformal-factor}
The local analytic inverse of the DeTurck map is a coordinate theorem.  It
does not change the sign of the Euclidean Einstein action.  The V22
conformal-factor obstruction and the V23--V25 constrained CMC/Schur route
remain the only admissible gravitational large-field branch currently
present in these notes.
\end{firewall}

\begin{assessment}[V43 acquisition]
\label{ass:t10v43-status}
The response-level metric power missing in V41 is now available on bounded
Dirichlet/Korn-complement patches: the metric is an analytic function of a
second-order DeTurck source with an explicit \(r^2\) loss, and low
gravitational heat rows can be projected before the remainder enters the
graded polymer tail.  The next bottleneck is compatibility with the
constrained CMC branch.  V44 will compare the DeTurck interior source,
the CMC scalar/vector constraint variables of V23--V25, and their boundary
Schur determinants, identifying the exact triangular transition and any
new Jacobian.
\end{assessment}

\begin{firewall}[Claim boundary after V43]
\label{fire:t10v43-claim}
V43 proves a local metric-source compiler on a fixed boundary complement.
It does not construct a global metric slice, solve the conformal-factor
problem, prove compatibility with the CMC constraint measure, contract the
infinite gravitational tower, or construct the SM--Einstein continuum.
\end{firewall}

\subsection{Parent record}

This V43 note is appended to the validated V42 source, whose record is
\begin{equation}
 \text{bytes}=601122,\qquad \text{lines}=15731,\qquad
 \operatorname{SHA256}=
 \texttt{7F89EF952886BACC7FDD69858C54D32E47BCEBDEFBAC4C13FFD0C8C8D3CE7D07}.
 \label{eq:t10v43-parent-record}
\end{equation}

% =============================================================================
% END APPENDED TENTH-SERIES V43 MATERIAL
% =============================================================================
\clearpage
\section{V44: The DeTurck--CMC fibre product and its nested interface determinant}
\label{sec:v10v44}

\subsection{Purpose and relation to the preceding construction}

The preceding note constructed a local DeTurck curvature chart for metric heat
owners.  Earlier notes constructed the constant-mean-curvature (CMC) scalar and
momentum constraint chart, including conformal-Killing quotients and scalar and
vector interface Schur complements.  These are not two names for the same
chart.  The first is a bulk elliptic chart for a spacetime metric (or for a
Riemannian filling used in the regulated Euclidean measure); the second is a
codimension-one chart for admissible initial or boundary data.  The purpose of
this note is to form their fibre product without integrating the conformal
metric degree twice.

All determinants in this note are initially finite-dimensional determinants at
a fixed regulator.  This removes domain and determinant-class ambiguities and
makes the block algebra exact.  Passage to Fredholm or zeta determinants is a
separate continuum debit.  Likewise, this note proves a local compatibility
theorem, not existence of a global Einstein filling.

\subsection{Trace splitting and the no-double-counting rule}

Let \(U\) be the regulated real vector space of bulk metric coefficients and
let \(I\) be the space of induced metric coefficients on the distinguished CMC
hypersurface \(\Sigma\).  Write
\[
 T_\Sigma:U\longrightarrow I
\]
for the regulated trace map.  We assume that a fixed extension operator
\(R_\Sigma:I\to U\) has been selected with
\(T_\Sigma R_\Sigma=\operatorname{id}_I\).  Then
\[
 U=U_0\oplus R_\Sigma I,
 \qquad U_0:=\ker T_\Sigma .
\]
Let \(Z\) denote the CMC constraint-owner space.  Its coordinates include the
momentum owner \(W\), the positive conformal factor \(\phi\), and any other
constrained matter owners.  Let \(B\) contain the retained trace, interface,
charge, and Killing data.  A prescribed-data map
\[
 \iota:Z\times B\longrightarrow I
\]
constructs the induced metric from the CMC variables.  The bulk metric is
parametrised by
\begin{equation}
 g(u,z,b)=g_0+u+R_\Sigma\iota(z,b),
 \qquad (u,z,b)\in U_0\times Z\times B .
 \label{eq:v44-trace-splitting}
\end{equation}

\begin{proposition}[No-double-counting trace chart]
\label{prop:v44-no-double-counting}
Suppose that \(\iota\) is injective in the conformal metric coordinates retained
in \(Z\), after quotienting the finite conformal-Killing stabiliser.  Then the
parametrisation \eqref{eq:v44-trace-splitting} is injective in \((u,z)\) for
fixed \(b\).  In particular, the conformal factor \(\phi\) is represented in
the trace term and is absent from the zero-trace bulk variable \(u\).
\end{proposition}

\begin{proof}
If \(g(u_1,z_1,b)=g(u_2,z_2,b)\), applying \(T_\Sigma\) gives
\(\iota(z_1,b)=\iota(z_2,b)\), because \(T_\Sigma u_j=0\) and
\(T_\Sigma R_\Sigma=\operatorname{id}_I\).  Injectivity of \(\iota\) on the
chosen quotient gives \(z_1=z_2\).  Subtracting the two bulk metrics then gives
\(u_1=u_2\).  Thus no coefficient belonging to the conformal trace is counted
in both variables.
\end{proof}

Boundary-fixed DeTurck diffeomorphisms act on \(U_0\).  Boundary diffeomorphisms
and Killing data that do not vanish at \(\Sigma\) belong to \(B\); this is the
same finite residual ledger retained in the earlier Ward--Korn and CMC notes.
Consequently, a boundary-changing diffeomorphism must not be silently absorbed
into the bulk DeTurck slice.

\subsection{The exact three-block elimination theorem}

Let \(Y,Q,J\) be regulated target spaces of the same respective dimensions as
\(U_0,Z,B\).  Consider analytic maps
\[
 E:U_0\times Z\times B\to Y,
 \qquad C:Z\times B\to Q,
 \qquad M:U_0\times Z\times B\to J.
\]
Here \(E\) is the fixed-trace DeTurck curvature map, \(C\) is the CMC
constraint map, and \(M\) is the remaining interface or matching equation.
The absence of \(u\) from \(C\) expresses the fact that the CMC data are solved
on \(\Sigma\) before the zero-trace bulk filling is varied.

Fix a base solution \((u_*,z_*,b_*)\) and set
\[
 A:=D_uE,
 \quad G:=D_zE,
 \quad H:=D_bE,
 \quad D:=D_zC,
 \quad K:=D_bC,
\]
\[
 P:=D_uM,
 \quad Q_1:=D_zM,
 \quad R:=D_bM,
\]
at that solution.  Define the twice-reduced interface operator
\begin{equation}
 S:=R-PA^{-1}H-\bigl(Q_1-PA^{-1}G\bigr)D^{-1}K.
 \label{eq:v44-nested-schur}
\end{equation}

\begin{theorem}[Nested DeTurck--CMC Schur theorem]
\label{thm:v44-nested-schur}
If \(A\), \(D\), and \(S\) are invertible, then the combined map
\[
 \mathcal F(u,z,b):=(E(u,z,b),C(z,b),M(u,z,b))
\]
is a local analytic diffeomorphism at \((u_*,z_*,b_*)\).  Its Jacobian has the
exact factorisation
\begin{equation}
 \det D\mathcal F=(\det A)(\det D)(\det S).
 \label{eq:v44-determinant-factorisation}
\end{equation}
The inverse of its derivative is obtained by first solving the bulk equation,
then the CMC equation, and finally the reduced interface equation.
\end{theorem}

\begin{proof}
In the order \((u,z,b)\), the derivative is
\[
 D\mathcal F=
 \begin{pmatrix}
 A&G&H\\
 0&D&K\\
 P&Q_1&R
 \end{pmatrix}.
\]
Left multiplication by the determinant-one block matrix
\[
 L_1=
 \begin{pmatrix}
 I&0&0\\
 0&I&0\\
 -PA^{-1}&0&I
 \end{pmatrix}
\]
removes the lower-left block and changes the last row to
\((0,Q_1-PA^{-1}G,R-PA^{-1}H)\).  A second determinant-one left multiplication,
with lower-middle block
\(-\bigl(Q_1-PA^{-1}G\bigr)D^{-1}\), removes the remaining lower-middle block.
The resulting matrix is block upper triangular with diagonal blocks
\(A,D,S\).  This proves \eqref{eq:v44-determinant-factorisation} and shows that
the derivative is invertible.  The finite-dimensional analytic inverse
function theorem gives the local analytic inverse.  The two row eliminations
also give the asserted order of solution.
\end{proof}

The formula has a useful invariant reading.  Holding the DeTurck target and the
CMC target fixed, the implicit derivatives are
\begin{equation}
 D_bz=-D^{-1}K,
 \qquad
 D_bu=-A^{-1}\bigl(H-GD^{-1}K\bigr).
 \label{eq:v44-implicit-responses}
\end{equation}
Substitution into \(D_bM\) gives exactly \(S\).  Thus \(S\) is the derivative
seen by the retained boundary owner after every stable interior owner has been
relaxed.  It is not the unreduced block \(R\).

\begin{corollary}[Quantitative inverse bound]
\label{cor:v44-inverse-bound}
Give all spaces fixed regulator norms and put
\[
 a=\lVert A^{-1}\rVert,
 \quad d=\lVert D^{-1}\rVert,
 \quad s=\lVert S^{-1}\rVert.
\]
Then \(D\mathcal F^{-1}\) has a bound depending polynomially on
\[
 a,d,s,\lVert G\rVert,\lVert H\rVert,
 \lVert K\rVert,\lVert P\rVert,\lVert Q_1\rVert.
\]
If this bound is \(N\) and
\(\lVert D\mathcal F(x)-D\mathcal F(x_*)\rVert\leq L\lVert x-x_*\rVert\)
on a ball of radius \(\rho\), then
\begin{equation}
 NL\rho<1
 \label{eq:v44-newton-radius}
\end{equation}
implies that every derivative in that ball remains invertible, with inverse
norm at most \(N/(1-NL\rho)\).
\end{corollary}

\begin{proof}
The block eliminations in the theorem express every block of the inverse as a
finite sum of products of the displayed operators and their three inverses;
this gives the polynomial bound.  Write
\(D\mathcal F(x)=D\mathcal F(x_*)[I+T(x)]\), where
\(\lVert T(x)\rVert\leq NL\rho<1\).  The Neumann series for
\((I+T(x))^{-1}\) proves the last assertion.
\end{proof}

In the DeTurck chart of V43 one has, at scale \(r\),
\(a\lesssim C_Er^2\).  In the sign chamber of the CMC scalar equation,
the scalar part of \(d\) is controlled by the coercivity constant
\(q_0^{-1}\); the vector part is controlled by the conformal Korn constant on
the conformal-Killing complement.  Formula \eqref{eq:v44-newton-radius} shows
where these constants multiply.  A small DeTurck radius cannot repair a
degenerating CMC coercivity constant or a vanishing interface singular value.

\subsection{CMC triangularity inside the middle block}

Write \(z=(W,\phi)\) and order the CMC constraints as momentum followed by
Hamiltonian.  At fixed freely specified data, the momentum equation does not
depend on \(\phi\), whereas the Hamiltonian equation depends on \(W\) through
the transverse tensor.  Therefore
\begin{equation}
 D=D_zC=
 \begin{pmatrix}
 \mathcal P&0\\
 \mathcal K&L_\phi
 \end{pmatrix},
 \label{eq:v44-cmc-triangular}
\end{equation}
where \(\mathcal P=-\operatorname{div}\mathcal L\) on the
conformal-Killing complement and \(L_\phi\) is the positive scalar
linearisation from the CMC sign chamber.

\begin{corollary}[Constraint determinant on the quotient]
\label{cor:v44-cmc-determinant}
After retaining conformal-Killing coordinates in \(B\) and restricting
\(\mathcal P\) to their chosen complement,
\begin{equation}
 \det D=(\det{}'\mathcal P)(\det L_\phi),
 \label{eq:v44-cmc-det-product}
\end{equation}
and the combined determinant becomes
\begin{equation}
 \det D\mathcal F
 = (\det A)(\det{}'\mathcal P)(\det L_\phi)(\det S).
 \label{eq:v44-full-det-product}
\end{equation}
\end{corollary}

\begin{proof}
Equation \eqref{eq:v44-cmc-triangular} is block lower triangular on the
conformal-Killing complement.  Taking its determinant and then using
Theorem~\ref{thm:v44-nested-schur} gives the two identities.
\end{proof}

This corollary combines, but does not repeat, the earlier scalar and vector
determinant calculations.  The new factor \(S\) includes the response of the
bulk DeTurck filling as well as the CMC response.  If the bulk block is removed,
\eqref{eq:v44-nested-schur} reduces to the interface Schur complement already
used in the CMC gluing note.

\subsection{Aligned graph coordinates and what their unit Jacobian means}

Fix nearby targets \((y,c)\).  By the inverse function theorem applied in
stages, let
\[
 z=Z(c,b)
 \]
be the solution of \(C(z,b)=c\), and, for arbitrary nearby \(z,b\), let
\[
 u=U(y,z,b)
\]
be the solution of \(E(u,z,b)=y\).  Introduce the graph coordinates
\begin{equation}
 \zeta=z-Z(c,b),
 \qquad e=u-U(y,z,b),
 \qquad b=b.
 \label{eq:v44-aligned-coordinates}
\end{equation}

\begin{proposition}[Unit graph-coordinate Jacobian]
\label{prop:v44-unit-graph-jacobian}
The change \((u,z,b)\mapsto(e,\zeta,b)\) in
\eqref{eq:v44-aligned-coordinates} has determinant one.  On the stable graph
\(e=\zeta=0\), the derivative of the remaining matching equation with respect
to \(b\) is \(S\).
\end{proposition}

\begin{proof}
The derivative of the coordinate change, in the order \((u,z,b)\), is block
upper triangular with identity diagonal blocks: \(D_ue=I\),
\(D_z\zeta=I\), and \(D_bb=I\).  Its determinant is therefore one.  On the
graph, differentiating \(C(Z(c,b),b)=c\) and
\(E(U(y,Z(c,b),b),Z(c,b),b)=y\) gives
\eqref{eq:v44-implicit-responses}.  The chain rule for \(M\) then gives
\eqref{eq:v44-nested-schur}.
\end{proof}

There is no contradiction between this unit Jacobian and
\eqref{eq:v44-determinant-factorisation}.  The unit Jacobian belongs to a
translation into graph deviations.  The factors \(\det A\), \(\det D\), and
\(\det S\) belong to the map from owners to equation residuals.  Confusing the
two would erase the coarea density by a coordinate convention.

\subsection{Three different determinant roles}

At a fixed regulator, ordinary change of variables gives the following exact
statements.

\begin{proposition}[Coordinate, constraint, and gauge factors]
\label{prop:v44-three-determinants}
Assume the hypotheses of Theorem~\ref{thm:v44-nested-schur} and choose
orientations, or take absolute values.
\begin{enumerate}
\item If the DeTurck source \(y=E(u,z,b)\) is used as a coordinate at fixed
      \((z,b)\), then
      \[
      du=|\det A|^{-1},dy.
      \]
      This is a coordinate density.
\item If the CMC equations are imposed by a delta constraint and \(z\) is
      integrated out at fixed \(b\), then
      \[
      dz\,\delta(C(z,b)-c)=|\det D|^{-1}
      \]
      evaluated at the local solution, with \(\det D\) replaced by the primed
      quotient determinant in the vector sector.
\item If the remaining matching equation is then imposed and \(b\) is
      integrated out, its density is \(|\det S|^{-1}\).
\end{enumerate}
Thus the local combined coarea density is
\begin{equation}
 |\det A|^{-1}|\det{}'\mathcal P|^{-1}
 |\det L_\phi|^{-1}|\det S|^{-1},
 \label{eq:v44-combined-density}
\end{equation}
only when all three residual maps are actually used in the indicated change of
variables or delta constraints.  A Faddeev--Popov determinant for quotienting
bulk diffeomorphisms is an additional gauge-orbit factor and is not any one of
the four factors in \eqref{eq:v44-combined-density}.
\end{proposition}

\begin{proof}
The first statement is the inverse-function change-of-variables formula.  The
second and third statements follow by applying the same formula to the unique
local zeros of the corresponding equal-dimensional maps.  The product is
\eqref{eq:v44-determinant-factorisation} together with
\eqref{eq:v44-cmc-det-product}.  The Faddeev--Popov map differentiates the
gauge condition along the diffeomorphism orbit, whereas \(A\) differentiates
the gauge-fixed curvature map along zero-trace metric coefficients.  Their
domains and maps differ, so their determinants cannot be identified.
\end{proof}

If \(E\) is used only to prove that every metric owner can be reconstructed
from a heat source, but the integration variable remains \(u\), then no
\(|\det A|^{-1}\) should be inserted.  Conversely, once the variable is changed
from \(u\) to \(y\), omitting that factor changes the regulated measure.  This
bookkeeping choice must be made once at the beginning of a continuum proof.

\subsection{Reflection compatibility}

Let \(\Theta\) exchange the two halves across the reflection hypersurface and
act linearly on all regulated owner and residual spaces.  Suppose
\(E,C,M\), the trace splitting, and the chosen complements are equivariant.

\begin{proposition}[Equivariance of the reduced response]
\label{prop:v44-reflection}
At a reflection-fixed base solution, the inverse blocks \(A^{-1}\), \(D^{-1}\)
and the reduced interface operator \(S\) intertwine the corresponding
reflection actions.  In particular, uniqueness forces the local solution maps
\(U\) and \(Z\) to be equivariant.
\end{proposition}

\begin{proof}
Differentiate the equivariance identities.  The resulting intertwining
relations for \(A,D\) pass to their inverses by uniqueness.  Substitution into
\eqref{eq:v44-nested-schur} gives the relation for \(S\).  Applying reflection
to a local solution of either implicit equation produces another solution with
the reflected data; local uniqueness identifies it with the reflected solution
map.
\end{proof}

Equivariance alone does not prove Osterwalder--Schrader positivity of the
determinant density.  Such a proof still requires the boundary localisation or
auxiliary-field representation developed earlier, now applied to the combined
operator \(S\).  The absolute determinant in
\eqref{eq:v44-combined-density} is positive as a finite density, but positivity
of a number is weaker than reflection positivity of the associated nonlocal
functional.

\subsection{Continuum debit and exact status}

The theorem isolates four analytic constants that a continuum construction
must keep uniform: the fixed-trace DeTurck inverse bound, the conformal Korn
bound, the CMC scalar coercivity bound, and the smallest singular value of the
twice-reduced interface operator.  It also isolates four determinant
renormalisations.  Uniformity of the first three does not imply uniformity of
the fourth.  In particular, a nearly singular boundary matching mode survives
after the two interior equations have been solved and cannot be hidden in an
interior elliptic estimate.

This note proves the exact local finite-regulator fibre-product algebra and the
no-double-counting parametrisation.  It does not prove that a prescribed CMC
slice has a global Einstein filling, that the Euclidean DeTurck operator stays
invertible along the full RG trajectory, or that the determinant factors have
regulator-independent renormalised limits.  These are genuine remaining
steps.

The next note will perform the scheduled comparison with the contemporaneous
Yang--Mills and Navier--Stokes programmes.  It will use that audit to choose a
concrete estimate for the new bottleneck: a lower bound for the reduced
interface operator \(S\) that survives local polymer gluing.

\clearpage
\section{V45: compulsory YM/NS audit and the interface-stability decision}
\label{sec:v10v45}

\subsection{Scope of the audit}

This is the scheduled five-version comparison.  The newest Yang--Mills
companion inspected was
\[
 \texttt{Renewal\_Yang\_Mills\_Mass\_Gap\_Research\_Notes\_V54.tex},
 \qquad
 \texttt{SHA--256 }388B846BD35A5B19F29870BEBBEF40A5DC0F2188892048B82E143E566FDE30C9,
\]
and the newest Navier--Stokes companion was
\[
 \texttt{Renewal\_NS\_Stability\_Research\_Notes\_V26.tex},
 \qquad
 \texttt{SHA--256 }82C887F8C2C69E4F28FAD7C3DC8DE5B9099CB5A4BF431EBA1FFF3E91935978AD.
\]
The comparison is confined to mechanisms that can change the direction of the
present SM--Einstein programme.  It does not import conclusions whose
hypotheses have not been verified here.

\subsection{The Yang--Mills transfer}

The Yang--Mills V53--V54 block contains two relevant mechanisms.  First, for a
positive block Hessian
\[
 {\mathsf H}_\theta=
 \begin{pmatrix}A_\theta&B_\theta\\B_\theta^*&C_\theta\end{pmatrix},
 \qquad
 {\mathsf S}_\theta=C_\theta-B_\theta^*A_\theta^{-1}B_\theta,
\]
it differentiates the harmonic router and Schur complement without expanding
the inverse entry by entry:
\[
 \dot L=-A^{-1}(\dot A L+\dot B),
\]
\[
 \dot{\mathsf S}
 =\dot C+\dot L_{\rm dir}^{\,*}B+B^*\dot L_{\rm dir}
   +L^*\dot A L,
\]
in the equivalent block notation of that note.  Second, V54 obtains the
higher derivatives from a finite word-derivative compiler and keeps crude
operator stocks separate from the Ward-summed physical mixing coefficient.

\begin{assessment}[What transfers from Yang--Mills]
\label{ass:v45-ym-transfer}
The inverse derivative identity
\[
 D(A^{-1})[\dot A]=-A^{-1}\dot A A^{-1}
\]
and the factorised differentiation of a Schur complement transfer verbatim to
the finite-regulator nested operator \(S\) of V44.  The positivity conclusion
does not transfer verbatim.  The Yang--Mills Schur kernel is the relaxation of
a positive block quadratic form.  The combined DeTurck--CMC system has not been
shown to be the Hessian of one positive functional, and the Einstein conformal
direction is precisely where such an assertion would require proof.
\end{assessment}

\begin{proof}
The inverse derivative identity follows by differentiating
\(A A^{-1}=I\).  It needs only invertibility and hence is independent of the
physical origin of \(A\).  By contrast, the variational identity
\[
 \langle b,{\mathsf S}b\rangle
 =\min_u\left\langle
 \binom ub,{\mathsf H}\binom ub\right\rangle
\]
uses self-adjointness and positivity of \({\mathsf H}\).  Neither follows from
the three-block inverse hypotheses of V44.  Therefore only the algebraic
derivative compiler is presently available.
\end{proof}

The methodological transfer is nevertheless strong: calculate a base
singular-value floor, differentiate the factorised nested response, and bound
the integrated perturbation.  As in YM V54, a crude absolute stock certifies
finiteness but must not be renamed a sharp physical coefficient.

\subsection{The Navier--Stokes transfer}

The Navier--Stokes V26 block computes exact pointwise norms of Oseen-kernel
derivatives on rank-one traceless inputs.  The first-derivative restricted norm
equals the full symmetric-traceless norm; the second-derivative norm is
strictly smaller only on an intermediate radial range, with a modest integrated
gain.  The relevant lesson is logical rather than numerical.

\begin{assessment}[What transfers from Navier--Stokes]
\label{ass:v45-ns-transfer}
An interface estimate may be sharpened on a restricted set of variations only
after the fibre-product equations prove that every admissible variation lies
in that set.  A rank-one, transverse-traceless, Ward-horizontal, or
constraint-tangent test is not a lower bound for the full operator unless that
set spans the actual domain on which invertibility is required.
\end{assessment}

\begin{proof}
For a linear operator \(T:B\to J\) and a proper subset
\(\mathcal A\) of the unit sphere,
\[
 \inf_{\|b\|=1}\|Tb\|
 \leq \inf_{b\in\mathcal A}\|Tb\|.
\]
Thus the restricted infimum can overestimate the smallest singular value.
It controls the inverse only if every required unit direction belongs to
\(\mathcal A\), or if a proved decomposition reduces all directions to
controlled pieces.  The exact rank-one calculation in the NS note respects
this distinction; the same distinction is mandatory here.
\end{proof}

For the SM--Einstein fibre product, the proven restrictions are presently the
zero-trace bulk splitting, the conformal-Killing quotient, and the linearised
constraint tangent equations.  They do not imply that every retained
interface metric variation is rank one.  No rank-one replacement will
therefore be made in the next estimate.

\subsection{Chosen direction: a pathwise singular-value certificate}

Let \(S(t):B\to J\), \(0\leq t\leq1\), be the V44 twice-reduced interface
operator along one regulated local gluing or RG path, after identifying the
finite-dimensional spaces by fixed isometries.  Put
\[
 \sigma_{\min}(S):=\inf_{\|b\|=1}\|Sb\|.
\]

\begin{lemma}[Pathwise floor]
\label{lem:v45-pathwise-floor}
If \(S\) is absolutely continuous in operator norm, then
\begin{equation}
 \sigma_{\min}(S(t))
 \geq \sigma_{\min}(S(0))
       -\int_0^t\|\dot S(\tau)\|\,d\tau .
 \label{eq:v45-pathwise-floor}
\end{equation}
Consequently, if the right-hand side is a positive number \(\gamma\), then
\(S(t)\) is injective and \(\|S(t)^{-1}\|\leq\gamma^{-1}\) whenever \(B\) and
\(J\) have the same finite dimension.
\end{lemma}

\begin{proof}
For every unit \(b\),
\[
 \|S(t)b\|
 \geq \|S(0)b\|-\|(S(t)-S(0))b\|
 \geq \sigma_{\min}(S(0))
       -\int_0^t\|\dot S(\tau)\|\,d\tau.
\]
Take the infimum over \(b\).  A positive lower bound gives injectivity.  In
equal finite dimensions this gives bijectivity and the stated inverse bound.
\end{proof}

\begin{decision}[V46 acquisition]
\label{dec:v45-v46}
The next note will derive an exact factorised formula for \(\dot S\) from the
seven V44 blocks and their first variations.  It will then give a computable
majorant in terms of the DeTurck inverse, the CMC inverse, and local block-jet
stocks.  Combining that majorant with
Lemma~\ref{lem:v45-pathwise-floor} will produce a finite-regulator interface
stability certificate.
\end{decision}

The certificate will have three layers:

\begin{enumerate}
\item a base floor \(\sigma_{\min}(S(0))\), obtained from an explicitly chosen
      local background rather than inferred from interior coercivity;
\item a factorised first-variation stock, keeping all inverse-response factors
      visible; and
\item a strict reserve inequality requiring the integrated stock to be smaller
      than the base floor.
\end{enumerate}

This direction uses both audits correctly.  It adopts the YM factorisation
instead of entrywise inverse algebra, and it permits restricted NS-style norms
only after a later geometric tangent-space theorem supplies the required
domain reduction.

\subsection{Status after the audit}

\begin{assessment}[Exact boundary of V45]
\label{ass:v45-status}
V45 proves the abstract pathwise singular-value estimate and fixes the next
calculation.  It does not yet evaluate a base Einstein interface floor or its
variation stock.  It does not promote DeTurck--CMC invertibility to a global
Einstein filling, renormalise any determinant, construct the full RG
trajectory, remove regulators, or establish the SM--Einstein continuum.
\end{assessment}

The post-audit direction is therefore to preserve full interface directions
and attack the nested operator directly.  If that estimate degenerates, the
upstream response will be to change the retained boundary owner or the
trace-extension operator and recompute the base Schur floor, rather than to
assume positivity absent from the combined system.

\clearpage
\section{V46: a factorised variation stock for the nested interface operator}
\label{sec:v10v46}

\subsection{Fixed quotient spaces along a regulated path}

Let \(t\in[0,T]\) parametrise one finite-regulator local gluing or RG
deformation.  The blocks of the V44 derivative are now functions of \(t\):
\[
 A(t),G(t),H(t),D(t),K(t),P(t),Q_1(t),R(t).
\]
They act between fixed finite-dimensional Hilbert spaces.  In the momentum
sector, fixed means that the conformal-Killing kernel has already been
split off and retained among the boundary owners.

\begin{firewall}[Kernel crossings are chart boundaries]
\label{fire:v46-kernel-crossing}
The following calculation assumes that the dimension of every retained gauge
or conformal-Killing space is constant on the path and that its complementary
projection is norm differentiable.  If an eigenvalue crosses zero, the primed
determinant and the inverse on the complement cannot be continued through that
point by the same chart.  Such a crossing requires a new finite-dimensional
owner and a new overlap map.
\end{firewall}

This condition can be checked at a regulator by a spectral gap around zero.
It is not supplied by the scalar CMC coercivity estimate.

Set
\[
 X=A^{-1},\qquad Y=D^{-1},\qquad
 T_1=Q_1-PXG.
\]
The twice-reduced interface operator of V44 can then be written
\begin{equation}
 S=R-PXH-T_1YK.
 \label{eq:v46-factorised-S}
\end{equation}
This form keeps the two relaxation stages visible.

\subsection{Exact first variation}

\begin{theorem}[Factorised derivative of the nested response]
\label{thm:v46-S-derivative}
Suppose all eight blocks are differentiable in operator norm and \(A,D\) are
invertible.  Then
\begin{align}
 \dot X&=-X\dot A X,
 &\dot Y&=-Y\dot D Y,
 \label{eq:v46-inverse-derivatives}\\
 \dot T_1
 &=\dot Q_1-\dot P XG+PX\dot A XG-PX\dot G,
 \label{eq:v46-T-derivative}\\
 \dot S
 &=\dot R-\dot P XH+PX\dot A XH-PX\dot H
   -\dot T_1YK+T_1Y\dot DYK-T_1Y\dot K.
 \label{eq:v46-S-derivative}
\end{align}
Every inverse in these expressions is one of the two inverses already required
to construct \(S\); differentiating introduces no new inverse.
\end{theorem}

\begin{proof}
Differentiate \(AX=I\) and \(DY=I\) to obtain
\eqref{eq:v46-inverse-derivatives}.  Differentiate
\(T_1=Q_1-PXG\), substitute \(\dot X=-X\dot A X\), and obtain
\eqref{eq:v46-T-derivative}.  Finally differentiate
\eqref{eq:v46-factorised-S} and substitute both inverse derivatives.  The
product rule gives exactly \eqref{eq:v46-S-derivative}.
\end{proof}

There is a second derivation that checks the signs.  Solve the differentiated
bulk and CMC response equations at fixed targets:
\[
 A\dot u+G\dot z+H\dot b=0,\qquad
 D\dot z+K\dot b=0.
\]
Thus
\[
 \dot z=-YK\dot b,\qquad
 \dot u=-X(H-GYK)\dot b.
\]
Substitution into the linearised matching equation gives
\[
 \dot M=S\dot b.
\]
Differentiating this identity with respect to the external path while keeping
\(\dot b\) fixed reproduces \eqref{eq:v46-S-derivative}.  This response check is
the analogue of the independent quadratic-row check in the audited YM note.

\subsection{A computable variation stock}

Use the abbreviations
\[
 x=\|X\|,\quad y=\|Y\|,\quad
 p=\|P\|,\quad g=\|G\|,\quad h=\|H\|,\quad
 k=\|K\|,\quad q=\|Q_1\|,
\]
and for differentiated blocks
\[
 a_1=\|\dot A\|,\ d_1=\|\dot D\|,\ p_1=\|\dot P\|,\ 
 g_1=\|\dot G\|,\ h_1=\|\dot H\|,\ k_1=\|\dot K\|,\ 
 q_1=\|\dot Q_1\|,\ r_1=\|\dot R\|.
\]
Define
\begin{align}
 \tau&:=q+pxg,
 \label{eq:v46-tau}\\
 \tau_1&:=q_1+p_1xg+px^2a_1g+pxg_1,
 \label{eq:v46-tau-one}\\
 \Omega&:=r_1+p_1xh+px^2a_1h+pxh_1
              +\tau_1yk+\tau y^2d_1k+\tau yk_1.
 \label{eq:v46-Omega}
\end{align}

\begin{proposition}[Variation-stock inequality]
\label{prop:v46-variation-stock}
At every point of the path,
\begin{equation}
 \|T_1\|\leq\tau,\qquad
 \|\dot T_1\|\leq\tau_1,\qquad
 \|\dot S\|\leq\Omega.
 \label{eq:v46-stock-bound}
\end{equation}
All constants in \(\Omega\) are norms of explicit regulated blocks, their
first jets, or the two existing inverse responses.
\end{proposition}

\begin{proof}
Apply submultiplicativity and the triangle inequality to
\(T_1=Q_1-PXG\), \eqref{eq:v46-T-derivative}, and
\eqref{eq:v46-S-derivative}, in that order.  Each displayed monomial is the
norm bound for exactly one product in the corresponding derivative formula.
\end{proof}

The placement of the inverse factors matters.  For example, a variation of the
bulk operator costs \(x^2a_1\), whereas a variation of the bulk-to-interface
coupling costs only \(xh_1\).  A construction that estimates every term by the
largest inverse power would lose the distinction between an operator jet and
a coupling jet and can consume the reserve unnecessarily.

\subsection{The strict interface certificate}

\begin{definition}[Nested interface stability certificate]
\label{def:v46-NISC}
A nested interface stability certificate (NISC) on \([0,T]\) consists of:
\begin{enumerate}
\item a certified base singular-value floor
      \(\sigma_{\min}(S(0))\geq\gamma_0>0\);
\item verified upper bounds for every stock in
      \eqref{eq:v46-Omega}, with the conformal-Killing complement fixed as
      above; and
\item a strict reserve
      \begin{equation}
       \int_0^T\Omega(t)\,dt\leq\gamma_0-\gamma_*,
       \qquad \gamma_*>0.
       \label{eq:v46-strict-reserve}
      \end{equation}
\end{enumerate}
\end{definition}

\begin{theorem}[Certified interface invertibility]
\label{thm:v46-certified-invertibility}
A NISC implies, for every \(t\in[0,T]\),
\begin{equation}
 \sigma_{\min}(S(t))\geq\gamma_*,
 \qquad
 \|S(t)^{-1}\|\leq\gamma_*^{-1}.
 \label{eq:v46-uniform-S-inverse}
\end{equation}
Together with uniform inverse bounds for \(A(t)\) and \(D(t)\), it gives a
uniform bound for the full three-block inverse \(D\mathcal F(t)^{-1}\).
\end{theorem}

\begin{proof}
Proposition~\ref{prop:v46-variation-stock} and the pathwise floor of V45 give
\[
 \sigma_{\min}(S(t))
 \geq\gamma_0-\int_0^t\Omega(\tau)\,d\tau
 \geq\gamma_*.
\]
The inverse bound follows in equal finite dimensions.  The two block Gaussian
eliminations in the V44 nested Schur theorem express the full inverse as finite
sums of products of \(A^{-1},D^{-1},S^{-1}\) and the coupling blocks.  Uniform
bounds for those factors give the last assertion.
\end{proof}

\begin{corollary}[Constant-stock form]
\label{cor:v46-constant-stock}
If \(\Omega(t)\leq\Omega_*\) on \([0,T]\), then
\[
 T\Omega_*<\gamma_0
\]
is sufficient for invertibility, and one may take
\(\gamma_*=\gamma_0-T\Omega_*\).
\end{corollary}

This is a local continuation criterion.  A long RG trajectory may be covered
by overlapping intervals, but at the start of each interval the new base
floor must be certified.  Merely subdividing time does not prevent the
accumulated physical floor from tending to zero.

\subsection{Consequences for the determinant density}

For an invertible real matrix path, the absolute determinant is differentiable
away from sign changes, and invertibility prevents such a sign change.

\begin{proposition}[Log-determinant variation]
\label{prop:v46-logdet}
Under a NISC,
\begin{equation}
 \frac{d}{dt}\log|\det S(t)|
 =\operatorname{tr}\bigl(S(t)^{-1}\dot S(t)\bigr),
 \label{eq:v46-logdet-identity}
\end{equation}
and, if \(n_B=\dim B\),
\begin{equation}
 \left|\log\frac{|\det S(t)|}{|\det S(0)|}\right|
 \leq n_B\gamma_*^{-1}\int_0^t\Omega(\tau)\,d\tau.
 \label{eq:v46-logdet-bound}
\end{equation}
\end{proposition}

\begin{proof}
Jacobi's determinant identity gives
\eqref{eq:v46-logdet-identity}.  For an \(n_B\)-dimensional space,
\[
 |\operatorname{tr}(S^{-1}\dot S)|
 \leq n_B\|S^{-1}\dot S\|
 \leq n_B\gamma_*^{-1}\Omega.
\]
Integrate this inequality.
\end{proof}

The factor \(n_B\) is harmless on one fixed interface cell but is not a
continuum estimate for an extensive boundary.  Polymer localisation must
replace it by a sum of connected local trace contributions before the
regulator is removed.

\subsection{Insertion of the elliptic stocks}

The DeTurck estimate of V43 gives schematically
\[
 x\leq C_Er^2
\]
on a ball of radius \(r\), at fixed trace.  The CMC block is triangular.  On
the conformal-Killing complement its inverse satisfies a bound of the form
\[
 y\leq C_{\rm tri}
 \left(
  \|\mathcal P^{-1}\|+\|L_\phi^{-1}\|
  +\|L_\phi^{-1}\mathcal K\mathcal P^{-1}\|
 \right),
\]
where
\[
 \|\mathcal P^{-1}\|\leq C_{\rm Korn},
 \qquad
 \|L_\phi^{-1}\|\leq q_0^{-1}
\]
in the previously established sign chamber.  Substitution into
\eqref{eq:v46-Omega} displays every possible loss in powers of
\(r^2\), \(C_{\rm Korn}\), and \(q_0^{-1}\).

\begin{assessment}[Upstream failure diagnosis]
\label{ass:v46-upstream}
If \eqref{eq:v46-strict-reserve} fails, the responsible monomial in
\(\Omega\) identifies the upstream repair:
\begin{enumerate}
\item terms containing \(x^2a_1\) require a smaller DeTurck chart, a better
      fixed-trace inverse, or a better bulk operator jet;
\item terms containing \(y^2d_1\) require a stronger CMC/Korn chamber or a
      different constrained owner;
\item large \(h_1,k_1,p_1,q_1,r_1\) require a better trace extension or
      interface owner; and
\item a small \(\gamma_0\) requires changing the base interface
      parametrisation itself.
\end{enumerate}
No reduction of the path step alone cures a base floor that is already zero.
\end{assessment}

\subsection{Locality debit and the next base calculation}

The NISC is an exact finite-regulator certificate.  To survive polymer gluing,
the block jets and inverse responses must also possess summable off-diagonal
decay.  Operator-norm finiteness on a large region is insufficient: it can
hide a factor proportional to boundary volume.  The required upgrade is a
weighted version of \eqref{eq:v46-Omega} in a convolution algebra whose norm
satisfies
\[
 \|UV\|_\kappa\leq\|U\|_\kappa\|V\|_\kappa.
\]
Then the same proof goes through with every norm replaced by
\(\|\cdot\|_\kappa\), while a small reserve in the exponential weight pays for
the enlarged support.

The missing input is now sharply located: a positive base floor
\(\gamma_0\) for a concrete combined interface chart.  The next note will go
upstream to the simplest nontrivial model, a flat product collar with a fixed
infrared regulator.  It will calculate the Dirichlet-to-Neumann symbols of the
scalar, vector, and gauge-fixed metric blocks, identify the retained zero
modes, and determine which boundary owner makes the base Schur operator
elliptic.

\begin{assessment}[Exact status after V46]
\label{ass:v46-status}
V46 proves the factorised derivative formula, the computable variation stock,
the strict finite-regulator invertibility certificate, and the determinant
variation bound.  It does not provide the numerical base floor, weighted
off-diagonal estimates, continuum determinant renormalisation, a global
Einstein filling, a full RG trajectory, or the SM--Einstein continuum.
\end{assessment}

\clearpage
\section{V47: the flat collar base floor and the massless interface zero mode}
\label{sec:v10v47}

\subsection{Why the base problem is moved upstream}

V46 reduced local continuation to a positive singular-value floor for one
concrete twice-reduced interface operator.  Interior ellipticity does not
supply that floor: the interface equation can still have a zero mode after
both interior problems have unique solutions.  We therefore choose the
simplest base geometry in which the interface owner and matching equation are
unambiguous.

Let
\[
 {\cal C}_+=[0,L_+]\times{\mathbb T}^d_\ell,\qquad
 {\cal C}_-=[-L_-,0]\times{\mathbb T}^d_\ell
\]
with common interface
\(\Gamma=\{0\}\times{\mathbb T}^d_\ell\).  At the outer faces impose
homogeneous Dirichlet data.  At \(\Gamma\), retain the Dirichlet trace \(b\);
the matching equation is the sum of outward conormal derivatives from the two
collars.  All formulae below also hold on a finite periodic lattice after
replacing \(|k|^2\) by the nonnegative lattice Laplacian symbol.

\subsection{The scalar Dirichlet-to-Neumann calculation}

For \(c>0\) and \(\mu\geq0\), consider
\[
 {\cal L}_{c,\mu}=c(-\partial_x^2-\Delta_y+\mu^2).
\]
For a tangential Fourier mode \(k\in(2\pi/\ell){\mathbb Z}^d\), put
\[
 \omega_k=(|k|^2+\mu^2)^{1/2}.
\]

\begin{proposition}[Two-sided scalar collar symbol]
\label{prop:v47-scalar-symbol}
The unique solution with interface value \(b_k\) and zero outer value is, on
the positive collar,
\[
 u_{+,k}(x)=
 b_k\frac{\sinh(\omega_k(L_+-x))}{\sinh(\omega_kL_+)},
\]
with the limiting affine formula \(b_k(1-x/L_+)\) when \(\omega_k=0\).
There is an analogous reflected formula on the negative collar.  The
two-sided Dirichlet-to-Neumann operator has Fourier multiplier
\begin{equation}
 s_{c,\mu}(k)
 =c\omega_k\{\coth(\omega_kL_+)+\coth(\omega_kL_-)\},
 \label{eq:v47-scalar-DtN}
\end{equation}
where the continuous value at \(\omega_k=0\) is
\begin{equation}
 s_{c,0}(0)=c(L_+^{-1}+L_-^{-1}).
 \label{eq:v47-zero-value}
\end{equation}
It is positive and satisfies
\begin{equation}
 s_{c,\mu}(k)\geq c(L_+^{-1}+L_-^{-1}).
 \label{eq:v47-scalar-floor}
\end{equation}
\end{proposition}

\begin{proof}
The Fourier-mode equation is
\(-u_k''+\omega_k^2u_k=0\).  The displayed hyperbolic-sine function has the
two boundary values and is unique by the positive Dirichlet energy.  At
\(x=0\), the positive-collar outward normal is \(-\partial_x\), so its outward
conormal derivative is
\(c\omega_k\coth(\omega_kL_+)b_k\).  The reflected negative collar contributes
the same expression with \(L_-\).  This proves
\eqref{eq:v47-scalar-DtN}.  Since
\[
 \lim_{z\downarrow0}z\coth(zL)=L^{-1}
\]
and \(z\coth(zL)\geq L^{-1}\) for \(z\geq0\), the remaining claims follow.
The last inequality is equivalent to \(w\cosh w\geq\sinh w\), whose derivative
is \(w\sinh w\geq0\) and whose value at \(w=0\) is zero.
\end{proof}

\begin{corollary}[Exact lowest scalar mode]
\label{cor:v47-lowest-scalar}
The function \(z\mapsto z\coth(zL)\) is increasing on
\([0,\infty)\).  Hence the smallest eigenvalue of the scalar collar operator is
\begin{equation}
 \gamma_{c,\mu}
 =c\mu\{\coth(\mu L_+)+\coth(\mu L_-)\},
 \label{eq:v47-exact-scalar-floor}
\end{equation}
with the continuous interpretation \eqref{eq:v47-zero-value} when \(\mu=0\).
\end{corollary}

\begin{proof}
For \(w=zL>0\),
\[
 \frac{d}{dz}\{z\coth(zL)\}
 =\coth w-w\,\operatorname{csch}^2w
 =\frac{\sinh(2w)-2w}{2\sinh^2w}>0.
\]
The tangential zero mode has the smallest \(\omega_k\), proving the claim.
\end{proof}

For the flat DeTurck metric block, the linearised operator is
\(-\frac12\Delta\) componentwise.  Thus \(c_D=\frac12,\mu_D=0\), and
\begin{equation}
 \gamma_D=\frac12(L_+^{-1}+L_-^{-1})
 \label{eq:v47-DeTurck-floor}
\end{equation}
on every gauge-fixed symmetric-tensor component.  For a constant-coefficient
CMC scalar linearisation \(-8\Delta+q_*\), one has
\[
 c_C=8,\qquad \mu_C=(q_*/8)^{1/2},
\]
and therefore
\begin{equation}
 \gamma_C
 =8(q_*/8)^{1/2}
 \left\{
 \coth\bigl(L_+(q_*/8)^{1/2}\bigr)
 +\coth\bigl(L_-(q_*/8)^{1/2}\bigr)
 \right\}.
 \label{eq:v47-CMC-floor}
\end{equation}
The continuous value at \(q_*=0\) is
\(8(L_+^{-1}+L_-^{-1})\).

\subsection{The momentum block from conformal Korn coercivity}

The flat conformal Killing operator in spatial dimension \(n\geq3\) is
\[
 ({\cal L}W)_{ij}
 =\partial_iW_j+\partial_jW_i-\frac{2}{n}(\operatorname{div}W)\delta_{ij},
\]
and the momentum operator is
\(\mathcal P=-\operatorname{div}{\cal L}\).  Its Dirichlet energy is
\[
 {\mathfrak a}(W,W)
 =\frac12\int_{\cal C}|{\cal L}W|^2,
\]
because integration by parts gives
\(\int\langle W,\mathcal PW\rangle=\frac12\int|{\cal L}W|^2\).

For each half collar, let \(\kappa_\pm>0\) be a conformal Korn constant on
fields vanishing at its outer face:
\begin{equation}
 \frac12\int_{{\cal C}_\pm}|{\cal L}W|^2
 \geq\kappa_\pm\int_{{\cal C}_\pm}|\nabla W|^2.
 \label{eq:v47-Korn}
\end{equation}
The outer Dirichlet condition removes the conformal-Killing kernel.  When a
different boundary condition is used, the kernel must instead be retained as
in the earlier CMC construction.

\begin{proposition}[Vector collar floor]
\label{prop:v47-vector-floor}
Let \(S_V\) be the sum of the two momentum Dirichlet-to-Neumann operators on
the common vector trace \(v\).  Then
\begin{equation}
 \langle v,S_Vv\rangle
 \geq
 \left(\frac{\kappa_+}{L_+}+\frac{\kappa_-}{L_-}\right)
 \|v\|_{L^2(\Gamma)}^2.
 \label{eq:v47-vector-DtN-floor}
\end{equation}
Consequently,
\[
 \gamma_V:=\frac{\kappa_+}{L_+}+\frac{\kappa_-}{L_-}
\]
is a certified base singular-value floor for the self-adjoint vector block.
\end{proposition}

\begin{proof}
The harmonic extension minimises \({\mathfrak a}\) among fields with the
prescribed interface and outer traces.  Green's formula identifies the
minimum two-sided energy with \(\langle v,S_Vv\rangle\).  For any competitor
on the positive collar, the fundamental theorem of calculus and
Cauchy--Schwarz give, pointwise in \(y\),
\[
 |v(y)|^2
 =\left|\int_0^{L_+}\partial_xW(x,y)\,dx\right|^2
 \leq L_+\int_0^{L_+}|\partial_xW(x,y)|^2\,dx.
\]
Integrate in \(y\), apply \eqref{eq:v47-Korn}, and repeat on the negative
collar.  Adding the two inequalities proves
\eqref{eq:v47-vector-DtN-floor}.
\end{proof}

This estimate is deliberately expressed through the Korn constants rather
than through a guessed componentwise symbol.  The normal and longitudinal
vector components mix in the half-space ODE; scalarising them would lose the
constraint geometry.

\subsection{The decoupled flat fibre-product base}

Use the \(L^2(\Gamma)\) Riesz maps to identify interface residuals with
interface owners.  At the flat, time-symmetric base with vanishing background
TT tensor, the first Hamiltonian variation in the momentum owner is zero,
because the quadratic term
\(|\sigma+{\cal L}W|^2\) has zero derivative at
\(\sigma=W=0\).  Choose the trace extension sectorwise by the harmonic
extensions above.  The resulting model interface operator is
\begin{equation}
 S_{\rm flat}=S_D\oplus S_C\oplus S_V.
 \label{eq:v47-flat-direct-sum}
\end{equation}

\begin{theorem}[Flat finite-collar base certificate]
\label{thm:v47-flat-certificate}
On the stated quotient and with outer Dirichlet conditions,
\(S_{\rm flat}\) is self-adjoint positive and
\begin{equation}
 \sigma_{\min}(S_{\rm flat})
 \geq
 \gamma_{\rm flat}:=\min\{\gamma_D,\gamma_C,\gamma_V\}>0.
 \label{eq:v47-flat-gamma}
\end{equation}
If the actual regulated base operator has the form
\[
 S(0)=S_{\rm flat}+{\cal E}_0
\]
under fixed isometric identifications and
\begin{equation}
 \|{\cal E}_0\|\leq\gamma_{\rm flat}-\gamma_0,
 \qquad \gamma_0>0,
 \label{eq:v47-base-defect}
\end{equation}
then \(\sigma_{\min}(S(0))\geq\gamma_0\), so it supplies item one of
the V46 NISC.
\end{theorem}

\begin{proof}
Each summand is a Dirichlet-to-Neumann operator of a positive Dirichlet form.
The scalar and DeTurck floors are
\eqref{eq:v47-exact-scalar-floor} and
\eqref{eq:v47-DeTurck-floor}; the vector floor is
\eqref{eq:v47-vector-DtN-floor}.  The smallest eigenvalue of a direct sum is
the minimum of the summand floors.  Finally, for every unit \(b\),
\[
 \|S(0)b\|
 \geq\|S_{\rm flat}b\|-\|{\cal E}_0b\|
 \geq\gamma_{\rm flat}-\|{\cal E}_0\|
 \geq\gamma_0.
\]
\end{proof}

The defect \({\cal E}_0\) contains the nonconstant background curvature,
nonzero TT data, matter coupling, differences between the chosen extension
and the sectorwise harmonic extension, and any residual mixing between the
DeTurck and CMC blocks.  It must be estimated; it is not set to zero by
notation.

\subsection{Why the Dirichlet trace is the correct retained owner}

\begin{proposition}[Owner choice and elliptic order]
\label{prop:v47-owner-choice}
For the scalar collar, the Dirichlet owner \(b\) with flux matching has
interface symbol \(s_{c,\mu}(k)\), of elliptic order one at high tangential
frequency.  If instead the conormal flux is retained and the trace is solved
for, the response has symbol \(s_{c,\mu}(k)^{-1}\), of order minus one.
Therefore the Dirichlet owner places the local matching equation in the
coercive orientation needed by the V46 singular-value certificate.
\end{proposition}

\begin{proof}
As \(|k|\to\infty\), \(\coth(\omega_kL_\pm)\to1\), so
\[
 s_{c,\mu}(k)=2c|k|+o(|k|).
\]
Its inverse is asymptotic to \((2c|k|)^{-1}\).  The former controls one
tangential derivative of the trace; the latter gains a derivative and has a
low-frequency amplification when the base floor is absent.
\end{proof}

This does not forbid Neumann coordinates on an overlap chart.  It fixes the
primary chart in which the interface equation has a positive-order principal
part.

\subsection{The infrared obstruction}

The finite-collar floor depends on the outer boundary.  Letting
\(L_+,L_-\to\infty\) in the massless scalar formula gives
\[
 s_{c,0}(k)\longrightarrow2c|k|.
\]
The tangential constant mode then has eigenvalue zero.

\begin{firewall}[No volume-independent massless floor from the collar]
\label{fire:v47-infrared}
For a massless metric or gauge block on an infinite collar, the
Dirichlet-to-Neumann operator has no positive floor on the complete trace
space.  One must either retain its constant mode as a macroscopic owner,
impose an independent infrared condition, or prove that the global matching
problem removes it.  The finite value in
\eqref{eq:v47-zero-value} cannot be used as a volume-independent continuum
gap.
\end{firewall}

For the gravitational metric trace, the constant mode includes changes of
macroscopic flat moduli and, depending on the component, residual gauge data.
The correct present choice is to project it into the finite macroscopic
ledger:
\[
 B=B_{\rm macro}\oplus B_{\rm loc},
\]
apply the NISC to \(B_{\rm loc}\), and carry \(B_{\rm macro}\) through the
global gluing equation.  This is the same logical treatment previously used
for harmonic and Killing modes.

\subsection{Locality implication and next step}

At a lattice regulator, the scalar multiplier is
\[
 f_{c,\mu,L_+,L_-}(\lambda)
 =c(\lambda+\mu^2)^{1/2}
 \left\{
 \coth\bigl(L_+(\lambda+\mu^2)^{1/2}\bigr)
 +\coth\bigl(L_-(\lambda+\mu^2)^{1/2}\bigr)
 \right\},
\]
evaluated on the local tangential lattice Laplacian.  Although the square root
is visible, the product \(z^{1/2}\coth(Lz^{1/2})\) is analytic at \(z=0\);
its first singularities occur when \(Lz^{1/2}\) reaches a nonzero imaginary
multiple of \(\pi\).  This suggests an exponentially local lattice kernel on
a strip fixed by the collar width, but the required weighted operator bound
has not yet been proved.

The next note will prove that statement by a resolvent or polynomial
functional-calculus argument and will quantify the loss of exponential
locality as \(L\) grows.  That is the missing bridge between the base floor
computed here and the weighted NISC required by polymer gluing.

\begin{assessment}[Exact status after V47]
\label{ass:v47-status}
V47 computes the exact scalar and DeTurck collar symbols, proves the vector
floor from conformal Korn coercivity, identifies a positive decoupled
finite-collar base, and proves a perturbative base certificate.  It also
proves that the massless infinite-collar zero mode must remain in the
macroscopic ledger.  It does not estimate the actual defect
\({\cal E}_0\), prove weighted off-diagonal locality, control the global
macroscopic modes, remove regulators, or establish the SM--Einstein
continuum.
\end{assessment}

\clearpage
\section{V48: exponential locality of the regulated collar response}
\label{sec:v10v48}

\subsection{Finite propagation and analytic functional calculus}

Let \(({\mathcal G},d)\) be the finite graph of tangential lattice sites at one
regulator.  For subsets \(X,Y\subset{\mathcal G}\), let
\(\mathbf 1_X\) and \(\mathbf 1_Y\) denote the coordinate projections.  A
self-adjoint matrix \(L\) has propagation at most \(r_0\) if
\[
 \mathbf 1_XL\mathbf 1_Y=0
 \quad\hbox{whenever}\quad d(X,Y)>r_0.
\]
Assume \(0\leq L\leq\Lambda I\).  Rescale it to
\[
 \widehat L=2L/\Lambda-I,
\]
whose spectrum lies in \([-1,1]\) and whose propagation is still at most
\(r_0\).

For \(\rho>1\), write \(E_\rho\) for the Bernstein ellipse with foci
\(\pm1\) and semiaxis sum \(\rho\).  Its leftmost point is
\(-(\rho+\rho^{-1})/2\).

\begin{theorem}[Analytic functional-calculus locality]
\label{thm:v48-chebyshev-locality}
Let \(f\) be analytic on and inside the image of \(E_\rho\) under
\(\lambda=\frac{\Lambda}{2}(z+1)\), and let
\[
 M_\rho=\sup_{z\in E_\rho}
 \left|f\left(\frac{\Lambda}{2}(z+1)\right)\right|.
\]
Then
\begin{equation}
 \|\mathbf 1_Xf(L)\mathbf 1_Y\|
 \leq
 \frac{2M_\rho\rho}{\rho-1}
 \exp\left\{-\frac{\log\rho}{r_0}d(X,Y)\right\}.
 \label{eq:v48-chebyshev-decay}
\end{equation}
The estimate is uniform in the number of lattice sites.
\end{theorem}

\begin{proof}
The Chebyshev expansion of the rescaled analytic function has coefficients
\(a_n\) satisfying
\[
 |a_n|\leq2M_\rho\rho^{-n}.
\]
Let \(p_{m-1}\) be its truncation through degree \(m-1\).  On \([-1,1]\),
\[
 \|f-p_{m-1}\|_\infty
 \leq2M_\rho\sum_{n=m}^\infty\rho^{-n}
 =\frac{2M_\rho\rho^{1-m}}{\rho-1}.
\]
A polynomial of degree \(m-1\) in \(\widehat L\) has propagation at most
\((m-1)r_0\).  Choose
\(m=\lceil d(X,Y)/r_0\rceil\).  Then
\(\mathbf 1_Xp_{m-1}(\widehat L)\mathbf 1_Y=0\), while the spectral theorem
bounds the remainder by its scalar supremum.  Since
\(m\geq d(X,Y)/r_0\), the displayed estimate follows.
\end{proof}

The theorem is a locality statement, not a smallness statement.
\(M_\rho\) grows when the chosen ellipse approaches a pole.  In applications
one chooses a strictly smaller ellipse and records both the prefactor and the
decay exponent.

\subsection{Location of the collar poles}

For one collar width \(L>0\), define
\[
 g_L(z)=z^{1/2}\coth(Lz^{1/2}).
\]
Although written with a square root, this is a single-valued meromorphic
function of \(z\).

\begin{lemma}[Meromorphic collar function]
\label{lem:v48-collar-meromorphic}
The apparent singularity of \(g_L\) at \(z=0\) is removable, with
\(g_L(0)=L^{-1}\).  Its poles are
\begin{equation}
 z=-\frac{\pi^2m^2}{L^2},\qquad m=1,2,\ldots.
 \label{eq:v48-collar-poles}
\end{equation}
\end{lemma}

\begin{proof}
The function \(w\coth(Lw)\) is even in \(w\), so its Laurent expansion contains
only powers of \(w^2=z\):
\[
 w\coth(Lw)=L^{-1}+\frac{L}{3}w^2+O(w^4).
\]
This removes the origin and proves single-valuedness in \(z\).  The nonzero
zeros of \(\sinh(Lw)\) occur at \(w=i\pi m/L\), which gives
\eqref{eq:v48-collar-poles}.
\end{proof}

For the two-sided scalar response of V47, set
\[
 f(\lambda)=c\{g_{L_+}(\lambda+\mu^2)
                  +g_{L_-}(\lambda+\mu^2)\}.
\]
Its closest pole to the nonnegative spectrum lies at
\[
 -\delta,\qquad
 \delta:=\mu^2+\frac{\pi^2}{L_{\max}^2},
 \qquad L_{\max}:=\max\{L_+,L_-\}.
\]
Under the rescaling of Theorem~\ref{thm:v48-chebyshev-locality}, this pole is
at
\[
 z_*=-1-\frac{2\delta}{\Lambda}.
\]
Define
\begin{equation}
 a_*:=1+\frac{2\delta}{\Lambda},
 \qquad
 \rho_*:=a_*+\sqrt{a_*^2-1}.
 \label{eq:v48-rho-star}
\end{equation}

\begin{corollary}[Scalar collar exponential decay]
\label{cor:v48-scalar-collar-decay}
For every \(1<\rho<\rho_*\), the regulated scalar
Dirichlet-to-Neumann matrix satisfies
\begin{equation}
 \|\mathbf 1_XS_{c,\mu}\mathbf 1_Y\|
 \leq C_\rho
 \exp\{-\alpha_\rho d(X,Y)\},
 \qquad
 \alpha_\rho=\frac{\log\rho}{r_0},
 \label{eq:v48-scalar-exp}
\end{equation}
where
\[
 C_\rho=\frac{2M_\rho\rho}{\rho-1}
\]
and \(M_\rho\) is the explicit ellipse supremum in
Theorem~\ref{thm:v48-chebyshev-locality}.
\end{corollary}

\begin{proof}
All poles of the shifted functions lie on the negative real axis at or to the
left of \(-\delta\).  The condition \(\rho<\rho_*\) is exactly
\[
 \frac{\rho+\rho^{-1}}{2}<1+\frac{2\delta}{\Lambda},
\]
so the rescaled Bernstein ellipse lies strictly to the right of the closest
pole.  Apply Theorem~\ref{thm:v48-chebyshev-locality}.
\end{proof}

For a massless collar with large \(L_{\max}\),
\[
 \log\rho_*
 =\operatorname{arcosh}\left(1+\frac{2\pi^2}
 {\Lambda L_{\max}^2}\right)
 =\frac{2\pi}{\sqrt{\Lambda}L_{\max}}+O(L_{\max}^{-3}).
\]
Thus the best available exponential rate deteriorates proportionally to
\(L_{\max}^{-1}\).  This is the locality counterpart of the infrared floor
loss found in V47.

\subsection{A general inverse-locality lemma}

The vector and coupled blocks need not be scalar functions of the tangential
Laplacian.  They can instead be treated as Schur complements of local positive
matrices.

\begin{proposition}[Exponential locality of a gapped finite-range inverse]
\label{prop:v48-inverse-locality}
Let \(A=A^*\) have propagation \(r_0\) and spectrum in
\([a,b]\), with \(0<a<b<\infty\).  Then, for every Bernstein ellipse that
avoids the rescaled pole of \(f(\lambda)=\lambda^{-1}\),
\begin{equation}
 \|\mathbf 1_XA^{-1}\mathbf 1_Y\|
 \leq C(a,b,\rho)
 \exp\left\{-\frac{\log\rho}{r_0}d(X,Y)\right\}.
 \label{eq:v48-inverse-decay}
\end{equation}
If \(B\) has propagation \(r_B\) and \(D\) has propagation \(r_D\), then
the Schur complement
\[
 S=D-B^*A^{-1}B
\]
obeys an exponential estimate with the same exponent and with its distance
shifted by at most \(2r_B\).
\end{proposition}

\begin{proof}
Rescale \([a,b]\) to \([-1,1]\).  The only pole of \(1/\lambda\) is at zero,
strictly to the left of the interval, so
Theorem~\ref{thm:v48-chebyshev-locality} applies to every ellipse not
containing that pole.  This proves \eqref{eq:v48-inverse-decay}.  For the Schur
term, if \(B_{zx}\) and \(B_{wy}\) are nonzero, then
\(d(z,x)\leq r_B\) and \(d(w,y)\leq r_B\).  Hence
\[
 d(z,w)\geq d(x,y)-2r_B.
\]
Insert the inverse decay between the two finite-range factors.  The local
\(D\) term vanishes when \(d(X,Y)>r_D\), proving the claim after enlarging the
prefactor.
\end{proof}

Applied to the collar interior matrix with outer Dirichlet data, this proves
exponential locality of the vector momentum Dirichlet-to-Neumann map.  The
lower spectral edge \(a\) is controlled by the conformal Korn and
one-dimensional Poincare constants.  It also treats mixed tensor blocks,
provided their fixed-gauge interior quadratic form is genuinely positive.
For the full Euclidean gravitational Hessian, positivity of such a form is
not assumed; the proposition is applied only to the already invertible
gauge-fixed elliptic blocks on contours separated from their spectrum.

\subsection{Weighted convolution algebra}

Let the graph have polynomial shell growth: there are constants
\(C_{\mathcal G},d_{\mathcal G}\), independent of volume, such that
\[
 \#\{y:d(x,y)=n\}\leq
 C_{\mathcal G}(1+n)^{d_{\mathcal G}-1}.
\]
For a block kernel \(K=(K_{xy})\), define
\begin{equation}
 \|K\|_\kappa
 :=\max\left\{
 \sup_x\sum_y e^{\kappa d(x,y)}\|K_{xy}\|,
 \sup_y\sum_x e^{\kappa d(x,y)}\|K_{xy}\|
 \right\}.
 \label{eq:v48-weighted-norm}
\end{equation}

\begin{proposition}[Weighted algebra and volume-free summation]
\label{prop:v48-weighted-algebra}
If
\[
 \|K_{xy}\|\leq C e^{-\alpha d(x,y)}
\]
with \(\alpha>\kappa\), then
\begin{equation}
 \|K\|_\kappa
 \leq CC_{\mathcal G}
 \sum_{n=0}^\infty(1+n)^{d_{\mathcal G}-1}
 e^{-(\alpha-\kappa)n}<\infty.
 \label{eq:v48-weighted-sum}
\end{equation}
Moreover,
\begin{equation}
 \|KL\|_\kappa\leq\|K\|_\kappa\|L\|_\kappa.
 \label{eq:v48-weighted-product}
\end{equation}
Both bounds are independent of the total number of sites.
\end{proposition}

\begin{proof}
Group the first row sum in \eqref{eq:v48-weighted-norm} by distance shells and
use the growth assumption; the column sum is identical.  For the product,
the triangle inequality and
\(d(x,z)\leq d(x,y)+d(y,z)\) give
\[
 e^{\kappa d(x,z)}\|(KL)_{xz}\|
 \leq\sum_y
 e^{\kappa d(x,y)}\|K_{xy}\|
 e^{\kappa d(y,z)}\|L_{yz}\|.
\]
Summation and the two Schur row/column bounds prove
\eqref{eq:v48-weighted-product}.
\end{proof}

This is precisely the algebra required in V46.  Every monomial in the
variation stock remains bounded when ordinary norms are replaced by
\(\|\cdot\|_\kappa\), provided all its factors have a common reserve
\(\alpha>\kappa\).

\subsection{Application to the nested DeTurck--CMC response}

At a fixed lattice regulator, assume:

\begin{enumerate}
\item the local blocks \(A,D,G,H,K,P,Q_1,R\) and their path derivatives have
      finite propagation or a uniform exponential kernel;
\item the fixed-trace DeTurck inverse and the scalar and vector CMC inverses
      have spectral or contour gaps sufficient for
      Proposition~\ref{prop:v48-inverse-locality}; and
\item harmonic, Killing, and massless constant modes have been extracted into
      \(B_{\rm macro}\), so their finite-rank projections are not included in
      the local kernel.
\end{enumerate}

\begin{theorem}[Weighted nested-interface certificate]
\label{thm:v48-weighted-NISC}
Under these assumptions, there exists \(\kappa>0\) for which
\(X=A^{-1}\), \(Y=D^{-1}\), \(T_1\), \(S\), and \(\dot S\) belong to the
weighted algebra \eqref{eq:v48-weighted-norm}.  The formulae and proofs of
V46 remain valid with every norm replaced by \(\|\cdot\|_\kappa\).  In
particular,
\begin{equation}
 \|\dot S\|_\kappa\leq\Omega_\kappa,
 \label{eq:v48-weighted-Omega}
\end{equation}
where \(\Omega_\kappa\) is the same polynomial as the V46 stock evaluated in
weighted norms, and its bound is independent of interface volume.
\end{theorem}

\begin{proof}
Choose \(\kappa\) smaller than the minimum decay exponent of all local inverse
and coupling kernels.  Propositions~\ref{prop:v48-inverse-locality} and
\ref{prop:v48-weighted-algebra} put the inverses and local blocks in the
weighted algebra.  The identities
\[
 T_1=Q_1-PXG,\qquad
 S=R-PXH-T_1YK
\]
and the exact derivative formula of V46 contain only sums and products in
that algebra.  Repeating the triangle and product estimates gives
\eqref{eq:v48-weighted-Omega}.  Volume independence follows from
\eqref{eq:v48-weighted-sum}.
\end{proof}

\begin{firewall}[Finite-rank global modes are not exponentially local]
\label{fire:v48-global-projection}
A projection onto a constant, harmonic, or Killing mode generally has a
kernel spread across the complete interface and does not have a
volume-uniform weighted norm.  Its small dimension does not make it local.
Such projections must remain in the explicit macroscopic block; adding them
back to \(S_{\rm loc}\) would invalidate the weighted certificate.
\end{firewall}

\subsection{What the locality theorem does and does not close}

The V47 base floor and Theorem~\ref{thm:v48-weighted-NISC} now give the two
abstract ingredients of a local polymer interface card: a strict local
singular-value reserve and a volume-free locality norm.  The remaining
model-specific input is to estimate the defect from the true
SM--Einstein blocks to the flat collar blocks in the same weighted norm.
The unweighted inequality \(\|{\cal E}_0\|<\gamma_{\rm flat}\) is not enough
for polymer summability.

The next note will construct a weighted defect ledger.  It will separate
background curvature, second fundamental form, TT/matter data, trace-extension
error, and gauge-slice variation, and will state a single sufficient
smallness chamber that feeds both the flat base floor and the V46 pathwise
reserve.

\begin{assessment}[Exact status after V48]
\label{ass:v48-status}
V48 proves regulator-uniform exponential off-diagonal decay for analytic
functions of finite-range lattice operators, locates the collar poles,
constructs the weighted convolution algebra, and lifts the V46 variation
stock into that algebra.  It does not estimate the actual curved
SM--Einstein defect, control macroscopic zero modes, renormalise the
extensive determinant, remove the lattice or volume regulators, or establish
the SM--Einstein continuum.
\end{assessment}

\clearpage
\section{V49: a weighted curved-defect chamber around the flat collar}
\label{sec:v10v49}

\subsection{Base and perturbed blocks}

Work in the weighted convolution algebra of V48 at one fixed regulator and
on the local interface space \(B_{\rm loc}\).  A subscript \(0\) denotes the
flat-collar base blocks, and no subscript denotes the curved SM--Einstein
blocks after fixed isometric identifications of the owner spaces.  Thus
\[
 A=A_0+\Delta A,\quad D=D_0+\Delta D,
\]
and similarly for \(G,H,K,P,Q_1,R\).  Put
\[
 X=A^{-1},\quad X_0=A_0^{-1},\qquad
 Y=D^{-1},\quad Y_0=D_0^{-1},
\]
\[
 T_1=Q_1-PXG,\qquad T_{1,0}=Q_{1,0}-P_0X_0G_0,
\]
and
\[
 S=R-PXH-T_1YK,\qquad
 S_0=R_0-P_0X_0H_0-T_{1,0}Y_0K_0.
\]
For this note, every norm is \(\|\cdot\|_\kappa\).

\subsection{Exact finite resolvent identities}

\begin{lemma}[Two inverse defects]
\label{lem:v49-inverse-defects}
If \(A,A_0,D,D_0\) are invertible, then
\begin{equation}
 X-X_0=-X(\Delta A)X_0,\qquad
 Y-Y_0=-Y(\Delta D)Y_0.
 \label{eq:v49-resolvent-identities}
\end{equation}
Consequently,
\begin{equation}
 \|X-X_0\|_\kappa\leq xx_0\delta_A,\qquad
 \|Y-Y_0\|_\kappa\leq yy_0\delta_D,
 \label{eq:v49-inverse-defect-bounds}
\end{equation}
where
\[
 x=\|X\|_\kappa,\quad x_0=\|X_0\|_\kappa,\quad
 y=\|Y\|_\kappa,\quad y_0=\|Y_0\|_\kappa,
\]
and \(\delta_Z=\|\Delta Z\|_\kappa\).
\end{lemma}

\begin{proof}
Multiply
\(A^{-1}-A_0^{-1}\) on the left by \(A\) and on the right by \(A_0\), or
expand
\[
 A^{-1}-A_0^{-1}=A^{-1}(A_0-A)A_0^{-1}.
\]
The identity for \(D\) is identical.  Apply the weighted algebra product
inequality from V48.
\end{proof}

The perturbed inverse norms themselves can be obtained from base data.  If
\begin{equation}
 x_0\delta_A<1,\qquad y_0\delta_D<1,
 \label{eq:v49-interior-neumann}
\end{equation}
then the Neumann series gives
\begin{equation}
 x\leq\frac{x_0}{1-x_0\delta_A},\qquad
 y\leq\frac{y_0}{1-y_0\delta_D}.
 \label{eq:v49-perturbed-inverse-bounds}
\end{equation}

\subsection{Exact non-infinitesimal interface defect}

Use the base stocks
\[
 p_0=\|P_0\|_\kappa,\qquad
 t_0=\|T_{1,0}\|_\kappa
\]
and the perturbed stocks
\[
 g=\|G\|_\kappa,\qquad h=\|H\|_\kappa,\qquad
 k=\|K\|_\kappa.
\]
Define
\begin{align}
 \varepsilon_T
 &:=
 \delta_{Q_1}+\delta_Pxg
 +p_0xx_0\delta_Ag+p_0x_0\delta_G,
 \label{eq:v49-epsilon-T}\\
 \varepsilon_S
 &:=
 \delta_R+\delta_Pxh
 +p_0xx_0\delta_Ah+p_0x_0\delta_H
 +\varepsilon_Tyk+t_0yy_0\delta_Dk+t_0y_0\delta_K.
 \label{eq:v49-epsilon-S}
\end{align}

\begin{theorem}[Weighted finite-defect estimate]
\label{thm:v49-finite-defect}
Under \eqref{eq:v49-interior-neumann},
\begin{equation}
 \|T_1-T_{1,0}\|_\kappa\leq\varepsilon_T,\qquad
 \|S-S_0\|_\kappa\leq\varepsilon_S.
 \label{eq:v49-finite-defect-bound}
\end{equation}
No omitted term is of higher order: the formulae
\eqref{eq:v49-epsilon-T}--\eqref{eq:v49-epsilon-S} bound exact finite
differences.
\end{theorem}

\begin{proof}
Insert and subtract intermediate products to obtain the exact identity
\[
 T_1-T_{1,0}
 =\Delta Q_1-(\Delta P)XG-P_0(X-X_0)G-P_0X_0\Delta G.
\]
For the interface block,
\[
\begin{split}
 S-S_0={}&\Delta R-(\Delta P)XH-P_0(X-X_0)H-P_0X_0\Delta H\\
 &-(T_1-T_{1,0})YK-T_{1,0}(Y-Y_0)K-T_{1,0}Y_0\Delta K.
\end{split}
\]
Apply the weighted product inequality, followed by
Lemma~\ref{lem:v49-inverse-defects}.  The resulting terms are exactly
\(\varepsilon_T\) and \(\varepsilon_S\).
\end{proof}

Other telescoping orders give different but equivalent stocks.  The order
above is chosen because it places the fully perturbed inverse only beside the
first changed factor and keeps the remaining response anchored at the flat
base.

\subsection{The geometric defect ledger}

Choose a collar radius \(r\) and use scale-invariant coefficient norms.  The
following five ledger entries are independent logical inputs:
\begin{align}
 \epsilon_{\rm geo}
 &:=\|g-g_{\rm flat}\|_{C^1_r}
    +r^2\|\operatorname{Rm}(g)\|_{C^0_r}
    +r\|\mathrm{II}_\Sigma\|_{C^0_r},
 \label{eq:v49-geo-ledger}\\
 \epsilon_{\rm CMC}
 &:=r\|\sigma\|_{C^0_r}
    +r\|{\cal L}W\|_{C^0_r}
    +\|\phi-\phi_*\|_{C^1_r}
    +r^2\|q_\phi-q_*\|_{C^0_r},
 \label{eq:v49-cmc-ledger}\\
 \epsilon_{\rm mat}
 &:=\text{the scaled norm of the matter stress, current, and their
 coefficient jets entering the constraints},
 \label{eq:v49-matter-ledger}\\
 \epsilon_{\rm ext}
 &:=\|R_\Sigma-R_{\Sigma,0}\|_\kappa
   +\|\nabla_nR_\Sigma-\nabla_nR_{\Sigma,0}\|_\kappa,
 \label{eq:v49-extension-ledger}\\
 \epsilon_{\rm gauge}
 &:=\|\Pi_{\rm slice}-\Pi_{\rm slice,0}\|_\kappa
   +\|\beta_g-\beta_{g_0}\|_{\mathrm{coeff},\kappa}.
 \label{eq:v49-gauge-ledger}
\end{align}
Here \(C^j_r\) denotes the norm in which each derivative is multiplied by
the corresponding power of \(r\).  The matter entry is a placeholder for the
already enumerated SM stress/current owners; it is not a claim that all
matter contributions have one tensor type.

\begin{proposition}[Finite-stencil coefficient reduction]
\label{prop:v49-coefficient-reduction}
Fix a compact coefficient chamber in which the lattice metric is uniformly
elliptic, \(\phi\) stays in a positive compact interval, and the local gauge
and constraint charts are defined.  For every block
\[
 Z\in\{A,D,G,H,K,P,Q_1,R\}
\]
there are regulator-stencil constants \(c_{Zj}\), independent of total
volume, such that
\begin{equation}
 \delta_Z
 \leq
 c_{Z{\rm g}}\epsilon_{\rm geo}
 +c_{Z{\rm c}}\epsilon_{\rm CMC}
 +c_{Z{\rm m}}\epsilon_{\rm mat}
 +c_{Z{\rm e}}\epsilon_{\rm ext}
 +c_{Z{\rm q}}\epsilon_{\rm gauge}.
 \label{eq:v49-block-ledger}
\end{equation}
A coefficient is zero when the corresponding physical entry does not occur
in that block.
\end{proposition}

\begin{proof}
At a fixed finite-range regulator, every matrix entry of a local block is a
smooth function of finitely many metric, connection, CMC, matter, extension,
and gauge coefficients in a bounded stencil.  The mean-value formula along
the line segment from the flat coefficient vector to the perturbed one bounds
the entry difference by the supremum of its coefficient derivative times the
five displayed differences.  Compactness of the chamber makes these
suprema finite.  Finite stencil overlap and the weighted row/column
definition turn the entry bounds into \(\|\cdot\|_\kappa\) bounds with
constants independent of total volume.
\end{proof}

The constants in \eqref{eq:v49-block-ledger} must be calculated from the
chosen discretisation.  The proposition proves that this calculation is
finite and local; it does not assign numerical values to the constants.

\subsection{One sufficient smallness chamber}

Let \(\delta_Z(\epsilon)\) be the right-hand side of
\eqref{eq:v49-block-ledger}, insert it into
\eqref{eq:v49-perturbed-inverse-bounds}, and then into
\eqref{eq:v49-epsilon-S}.  Denote the resulting explicit rational function by
\[
 {\mathfrak E}_S(
 \epsilon_{\rm geo},\epsilon_{\rm CMC},\epsilon_{\rm mat},
 \epsilon_{\rm ext},\epsilon_{\rm gauge}).
\]

\begin{theorem}[Curved base chamber]
\label{thm:v49-curved-chamber}
Let \(\gamma_{\rm flat}\) be the local flat-collar floor of V47.  Suppose
\begin{align}
 x_0\delta_A(\epsilon)&<1,\qquad
 y_0\delta_D(\epsilon)<1,
 \label{eq:v49-chamber-interior}\\
 {\mathfrak E}_S(\epsilon)
 &\leq\gamma_{\rm flat}-\gamma_0,
 \qquad \gamma_0>0.
 \label{eq:v49-chamber-interface}
\end{align}
Then the curved local interface operator is invertible and
\begin{equation}
 \sigma_{\min}(S)\geq\gamma_0,\qquad
 \|S^{-1}\|\leq\gamma_0^{-1}.
 \label{eq:v49-curved-floor}
\end{equation}
If, along a subsequent path, the weighted V46 variation stock satisfies
\begin{equation}
 \int_0^T\Omega_\kappa(t)\,dt
 \leq\gamma_0-\gamma_*,
 \qquad\gamma_*>0,
 \label{eq:v49-combined-reserve}
\end{equation}
then the complete local three-block chart remains uniformly invertible on
that path.
\end{theorem}

\begin{proof}
The first two inequalities give the perturbed interior inverse bounds.
Theorem~\ref{thm:v49-finite-defect} and
\eqref{eq:v49-chamber-interface} imply
\[
 \|S-S_0\|\leq\|S-S_0\|_\kappa
 \leq\gamma_{\rm flat}-\gamma_0.
\]
The singular-value perturbation inequality and the V47 floor give
\(\sigma_{\min}(S)\geq\gamma_0\).  The path assertion is the weighted NISC
of V46--V48.
\end{proof}

This theorem supplies one chamber, not the largest chamber.  Its purpose is
to expose which physical input consumes which part of the reserve.  In
particular, background smallness cannot compensate for a trace extension
whose weighted normal-derivative defect is large.

\subsection{A finite determinant comparison}

\begin{proposition}[Relative determinant series]
\label{prop:v49-relative-determinant}
Assume \(S_0\) and \(S\) are invertible and
\[
 \eta:=\|S_0^{-1}(S-S_0)\|<1.
\]
Then
\begin{equation}
 \log\frac{|\det S|}{|\det S_0|}
 =\sum_{m=1}^\infty\frac{(-1)^{m+1}}{m}
   \operatorname{tr}\{[S_0^{-1}(S-S_0)]^m\},
 \label{eq:v49-logdet-series}
\end{equation}
and
\begin{equation}
 \left|\log\frac{|\det S|}{|\det S_0|}\right|
 \leq n_B[-\log(1-\eta)].
 \label{eq:v49-logdet-finite-bound}
\end{equation}
\end{proposition}

\begin{proof}
Write
\(S=S_0(I+K)\), with \(K=S_0^{-1}(S-S_0)\).  The norm condition makes the
power series for \(\log(I+K)\) absolutely convergent.  Taking the trace gives
\eqref{eq:v49-logdet-series}.  Since
\(|\operatorname{tr}K^m|\leq n_B\|K\|^m\), summation gives
\eqref{eq:v49-logdet-finite-bound}.
\end{proof}

The factor \(n_B\) is extensive.  The weighted algebra makes each trace word
local enough for a connected-polymer reorganisation, but that reorganisation
and its counterterms are still required before a continuum determinant
density is claimed.

\subsection{Upstream response and the V50 audit}

The defect chamber now gives a deterministic response to failure.  If
\(\delta_A\) or \(\delta_D\) violates
\eqref{eq:v49-chamber-interior}, shrink the geometric/CMC coefficient chart
or change its owner coordinates.  If
\eqref{eq:v49-chamber-interface} fails while the interior conditions hold,
inspect the seven explicit terms in \(\varepsilon_S\); a dominant
\(\epsilon_{\rm ext}\) or \(\epsilon_{\rm gauge}\) calls for a different
extension or slice before any smaller RG step is attempted.  If
\(\gamma_{\rm flat}\) collapses, move the corresponding zero mode into the
macroscopic ledger.

V50 is the next compulsory companion audit.  It will inspect the then-current
YM and NS notes for a reusable implementation of either a finite local jet
certificate or a restricted physical direction norm.  The post-audit note
will then choose between calculating the actual stencil constants in
\eqref{eq:v49-block-ledger} and organising the relative determinant series
into connected boundary polymers.

\begin{assessment}[Exact status after V49]
\label{ass:v49-status}
V49 proves exact weighted finite-difference formulae for the two inverse
responses and the nested interface operator, constructs a five-entry physical
defect ledger, and gives one sufficient curved-background chamber.  It does
not compute the stencil constants, prove the chamber along a full RG
trajectory, control the macroscopic block, renormalise the determinant
series, remove regulators, or establish the SM--Einstein continuum.
\end{assessment}

\clearpage
\section{V50: compulsory companion audit and the measured-fibre correction}
\label{sec:v10v50}

\subsection{Files inspected}

The second scheduled audit in this cycle inspected
\[
 \texttt{Renewal\_Yang\_Mills\_Mass\_Gap\_Research\_Notes\_V57.tex},
 \qquad
 \texttt{SHA--256 }26321C107687AF284575E1BDD6B7DB35DD01CB45B3378DCB3E3001292A464522,
\]
and
\[
 \texttt{Renewal\_NS\_Stability\_Research\_Notes\_V26.tex},
 \qquad
 \texttt{SHA--256 }82C887F8C2C69E4F28FAD7C3DC8DE5B9099CB5A4BF431EBA1FFF3E91935978AD.
\]
The NS companion has not changed since the V45 audit.  The YM companion has
advanced from a finite mixed-jet compiler to a decomposition of local
irrelevant owners into an equation-of-motion direction and a lattice
anisotropy direction.

\subsection{New Yang--Mills lesson}

YM V57 proves, on a compact finite-dimensional integration fibre with smooth
measure \(d\mu\), the exact identity
\[
 \int\{DS[V]-\operatorname{div}_\mu V\}e^{-S}\,d\mu=0
\]
when \(V\) is tangent to that fibre.  It also identifies the obstruction:
a covariant equation-of-motion vector field that is vertically correct only
to linear order need not be tangent to the exact nonlinear block fibre.
Projecting it vertically can introduce a nonlocal inverse, while retaining
its horizontal part changes the coarse variable and carries another
Jacobian.

\begin{assessment}[Transfer to the SM--Einstein chart]
\label{ass:v50-ym-transfer}
A change of DeTurck gauge slice, CMC owner, or trace extension may be treated
as a redundant fibre coordinate change only after proving tangency to the
exact combined map
\[
 (E,C,M)=\hbox{fixed retained data}.
\]
Its action derivative alone is not a null owner.  The measure divergence of
the generating vector field is mandatory, and a nonzero horizontal component
must be recorded as a transformation of the retained interface owner.
\end{assessment}

\begin{proof}
The change-of-variables identity used in YM depends only on the divergence
theorem and therefore applies to any finite regulated fibre with a smooth
positive density.  The fibre in the present programme is cut out by the exact
DeTurck, CMC, and matching maps, so tangency means annihilation by their full
derivative, not by the flat principal symbols alone.  If the derivative has a
horizontal component, the flow leaves that fibre and compares different
retained data; the vertical divergence theorem no longer identifies the
result with zero at fixed retained owner.
\end{proof}

This directly precedes the determinant-polymer task.  A claimed cancellation
between a gauge or extension counterterm and a Jacobian would be invalid if
the coordinate generator were not exactly vertical.

\subsection{Navier--Stokes check}

The NS V26 conclusion is unchanged: restricting a tensor norm to rank-one
physical inputs gives no improvement at first derivative and only a modest
improvement at second derivative on a bounded radial range.  It continues to
support the V45 decision.

\begin{assessment}[No new restricted-domain replacement]
\label{ass:v50-ns-transfer}
The curved defect chamber of V49 must continue to use the full proved local
interface domain.  Neither the new YM redundancy direction nor the existing
NS rank-one calculation proves that SM--Einstein interface variations occupy
a smaller rank-one cone.
\end{assessment}

\subsection{Direction selected after the audit}

The V49 options were to compute stencil constants immediately or to organise
relative determinants into connected polymers.  The YM result shows that one
logical layer must come first: identify exactly which chart deformations are
vertical measured coordinate changes.  Otherwise either calculation can
assign a spurious local owner to a Jacobian that should cancel, or cancel an
owner whose generator changes the retained metric.

\begin{decision}[V51 measured-fibre task]
\label{dec:v50-v51}
V51 will:
\begin{enumerate}
\item define the regulated combined fibre
      \({\mathcal F}_{y,c,m}=\{(u,z,b):(E,C,M)=(y,c,m)\}\);
\item construct the vertical projection from the V44 block inverse;
\item prove the measured field-redefinition identity on that fibre;
\item decompose a general chart generator into vertical and horizontal
      parts; and
\item derive the exact retained-owner Jacobian for the horizontal part.
\end{enumerate}
Only after this decomposition will local determinant words be classified as
redundant or physical.
\end{decision}

The vertical projection will contain the inverse of the combined block
derivative.  V46--V49 now give the conditions under which that inverse is
bounded and exponentially local on the local sector.  Thus the audit does not
introduce an uncontrolled inverse; it explains why the preceding interface
work was necessary.

\subsection{Audit status}

\begin{assessment}[Exact boundary of V50]
\label{ass:v50-status}
V50 records the current companion files and proves that exact nonlinear
verticality plus measure divergence is required before a chart deformation
is declared redundant.  It does not yet construct the vertical projection,
evaluate its divergence, classify the actual SM--Einstein chart
deformations, compute the V49 stencil constants, renormalise determinant
polymers, control macroscopic modes, or establish the SM--Einstein continuum.
\end{assessment}

\clearpage
\section{V51: measured verticality in the combined local owner chart}
\label{sec:v10v51}

\subsection{The dimension correction}

V50 proposed writing a fibre by fixing all components of
\((E,C,M)\).  Under the V44 hypotheses that map is a local diffeomorphism
between equal-dimensional spaces, so such a level set is locally a point.
There is then no nonzero vertical field and no RG fibre integration.  The
correct construction must distinguish fine residual coordinates from
genuinely retained block coordinates.

Let \(X\) be the regulated bosonic owner manifold, of dimension \(N\), and let
\[
 {\cal B}:X\longrightarrow W
\]
be the exact retained block map, of constant rank \(m<N\).  Its fibre
\(X_w={\cal B}^{-1}(w)\) has dimension \(N-m\).  Choose a local fine-coordinate
map
\[
 {\cal R}:X\longrightarrow F,\qquad \dim F=N-m,
\]
made from the DeTurck source, CMC residual, stable interface residuals, and
any remaining local fine owners.  The square map
\begin{equation}
 \Psi=({\cal R},{\cal B}):X\longrightarrow F\times W
 \label{eq:v51-square-chart}
\end{equation}
is required to be a local diffeomorphism.

\begin{firewall}[Residual targets are not all retained owners]
\label{fire:v51-dimension}
The complete V44 map can be used as a square coordinate chart, but an RG
fibre is obtained only after declaring which coordinates are integrated and
which are retained.  Fixing every coordinate of a local diffeomorphism leaves
no fine integration.  Determinant factors must therefore be assigned after
the \(F/W\) split, not before it.
\end{firewall}

The V44 nested inverse supplies the inverse of \(D\Psi\) when its block rows
are ordered according to this split.  The V46--V49 certificates then control
that inverse and its locality.

\subsection{Horizontal lift and exact vertical projection}

Let
\[
 \iota_W:W\longrightarrow F\times W,\qquad
 \iota_W\xi=(0,\xi),
\]
and define, at each point of the chart,
\begin{equation}
 J:=D\Psi^{-1}\iota_W,\qquad
 \Pi_{\rm v}:=I-JD{\cal B},\qquad
 \Pi_{\rm h}:=JD{\cal B}.
 \label{eq:v51-projections}
\end{equation}

\begin{proposition}[Chart-induced vertical splitting]
\label{prop:v51-vertical-splitting}
The maps in \eqref{eq:v51-projections} obey
\[
 D{\cal B}\,J=I_W,\qquad
 \Pi_{\rm v}^2=\Pi_{\rm v},\qquad
 \Pi_{\rm h}^2=\Pi_{\rm h},\qquad
 \Pi_{\rm v}\Pi_{\rm h}=\Pi_{\rm h}\Pi_{\rm v}=0,
\]
and
\[
 \operatorname{ran}\Pi_{\rm v}=\ker D{\cal B}.
\]
Thus every vector field \(V\) has the exact decomposition
\begin{equation}
 V=V_{\rm v}+V_{\rm h},\qquad
 V_{\rm v}=\Pi_{\rm v}V,\qquad
 V_{\rm h}=J(D{\cal B}V).
 \label{eq:v51-vector-splitting}
\end{equation}
\end{proposition}

\begin{proof}
Because \(D\Psi J=\iota_W\), projection onto the second factor gives
\(D{\cal B}J=I_W\).  Hence
\[
 \Pi_{\rm h}^2=JD{\cal B}JD{\cal B}=JD{\cal B}=\Pi_{\rm h},
\]
and the corresponding identities for \(I-\Pi_{\rm h}\) follow.  Moreover
\(D{\cal B}\Pi_{\rm v}=0\), so
\(\operatorname{ran}\Pi_{\rm v}\subset\ker D{\cal B}\).  If
\(D{\cal B}v=0\), then \(\Pi_{\rm v}v=v\), proving the reverse inclusion.
The decomposition is the identity
\(I=\Pi_{\rm v}+\Pi_{\rm h}\).
\end{proof}

The projection is exact and nonlinear through its dependence on the base
point.  Replacing \(D\Psi\) by its flat principal part proves only approximate
verticality.

\subsection{The vertical measured identity}

Let \(d\mu_w\) be a smooth positive conditional measure on \(X_w\), and let
\(S:X\to{\mathbb R}\) be the regulated action.  Assume either that \(X_w\) is
compact without boundary, that the vector field below is tangent to its
boundary, or that the weighted boundary flux vanishes.

\begin{theorem}[Measured vertical redundancy]
\label{thm:v51-vertical-redundancy}
If \(V\) is a smooth vertical vector field,
\[
 D{\cal B}[V]=0,
\]
then
\begin{equation}
 \int_{X_w}
 \left\{DS[V]-\operatorname{div}_{\mu_w}V\right\}
 e^{-S}\,d\mu_w=0.
 \label{eq:v51-vertical-identity}
\end{equation}
Every derivative of this identity with respect to \(w\) also vanishes while
the same smooth fibre chart and boundary condition persist.
\end{theorem}

\begin{proof}
On \(X_w\),
\[
 \operatorname{div}_{\mu_w}(e^{-S}V)
 =e^{-S}\{\operatorname{div}_{\mu_w}V-DS[V]\}.
\]
The divergence theorem and the assumed boundary condition make its integral
zero, proving \eqref{eq:v51-vertical-identity}.  In a smooth local
trivialisation, the integral is the zero function of \(w\); differentiate
that function.
\end{proof}

\begin{corollary}[Projected chart generator]
\label{cor:v51-projected-generator}
For any smooth chart-deformation field \(V\), the owner
\begin{equation}
 {\cal O}_{V_{\rm v}}
 :=DS[\Pi_{\rm v}V]
   -\operatorname{div}_{\mu_w}(\Pi_{\rm v}V)
 \label{eq:v51-projected-owner}
\end{equation}
has zero conditional expectation at fixed \(w\).
\end{corollary}

This corollary is exact, but projection has a price.  The divergence contains
derivatives of \(D\Psi^{-1}\), and hence second derivatives of the DeTurck,
CMC, and matching maps.  The first-variation compiler of V46 is the beginning
of that calculation; a numerical cancellation cannot be claimed without the
projection-divergence term.

\subsection{General projectable generators and the retained Jacobian}

The horizontal component of a general vector field need not define a vector
field on \(W\).  Call \(V\) projectable if there is a vector field
\(\bar V\) on \(W\) such that
\begin{equation}
 D{\cal B}(x)V(x)=\bar V({\cal B}(x))
 \quad\hbox{for all }x.
 \label{eq:v51-projectability}
\end{equation}
In the local product coordinates \((r,w)=\Psi(x)\), write the regulated
measure as
\[
 d\mu(x)=j(r,w)\,dr\,dw
\]
and define the marginal density and effective action by
\begin{equation}
 Z(w)=\int_F e^{-S(r,w)}j(r,w)\,dr,\qquad
 S_{\rm eff}(w)=-\log Z(w).
 \label{eq:v51-effective-action}
\end{equation}
Boundary flux in the \(r\) variables is again assumed to vanish.

\begin{theorem}[Measured horizontal transport identity]
\label{thm:v51-horizontal-transport}
For every projectable \(V\),
\begin{equation}
 \int_F
 \{DS[V]-\operatorname{div}_{\mu}V\}
 e^{-S}j\,dr
 =-\operatorname{div}_w\{Z(w)\bar V(w)\}.
 \label{eq:v51-horizontal-density}
\end{equation}
Equivalently, after division by \(Z(w)\),
\begin{equation}
 {\mathbb E}_w
 [DS[V]-\operatorname{div}_{\mu}V]
 =
 DS_{\rm eff}(w)[\bar V(w)]
 -\operatorname{div}_w\bar V(w).
 \label{eq:v51-effective-owner}
\end{equation}
\end{theorem}

\begin{proof}
In product coordinates write
\(V=(V_F(r,w),\bar V(w))\).  Multiplication of the divergence identity by the
density gives
\[
 \{DS[V]-\operatorname{div}_{\mu}V\}e^{-S}j
 =
 -\operatorname{div}_{(r,w)}(e^{-S}jV).
\]
Integrate in \(r\).  The fine-coordinate divergence has zero boundary flux,
while projectability allows \(\bar V(w)\) to leave the integral:
\[
 -\int_F\operatorname{div}_{(r,w)}(e^{-S}jV)\,dr
 =-\operatorname{div}_w\left\{
 \bar V(w)\int_Fe^{-S}j\,dr\right\}.
\]
This is \eqref{eq:v51-horizontal-density}.  Since
\(D Z[\bar V]=-Z\,DS_{\rm eff}[\bar V]\), division by \(Z\) gives
\eqref{eq:v51-effective-owner}.
\end{proof}

The right side of \eqref{eq:v51-effective-owner} contains both the change in
the effective action and the Jacobian of the retained coordinate change.  A
chart deformation with nonzero projectable horizontal part is therefore not
a null owner at fixed retained coordinates.  It becomes a measured
field-redefinition owner for the effective measure on \(W\).

\begin{firewall}[Nonprojectable horizontal components]
\label{fire:v51-nonprojectable}
If \(D{\cal B}V\) varies along the fine fibre, it is not a deterministic
retained coordinate change.  Replacing it by an arbitrary value or by its
flat approximation invalidates \eqref{eq:v51-effective-owner}.  One may
project \(V\) vertically using \(\Pi_{\rm v}\); the remainder then requires an
explicit conditional, generally measure-dependent transport law.
\end{firewall}

\subsection{Divergence of the exact vertical projection}

In a local coordinate density \(d\mu=\varrho(x)\,dx\),
\[
 \operatorname{div}_{\mu}U
 =\operatorname{tr}(DU)+D\log\varrho[U].
\]
For \(U=\Pi_{\rm v}V\), the product rule gives
\begin{equation}
 \operatorname{div}_{\mu}(\Pi_{\rm v}V)
 =
 \operatorname{tr}\{\Pi_{\rm v}DV+(D\Pi_{\rm v})[V]\}
 +D\log\varrho[\Pi_{\rm v}V].
 \label{eq:v51-projected-divergence}
\end{equation}

\begin{proposition}[Derivative content of the projection]
\label{prop:v51-projection-derivative}
The derivative \(D\Pi_{\rm v}\) is a finite sum of products containing
\(D\Psi^{-1}\), \(D^2\Psi\), and \(D^2{\cal B}\).  More precisely,
for a direction \(h\),
\begin{equation}
 D J[h]
 =-D\Psi^{-1}\{D^2\Psi[h]\}J,
 \label{eq:v51-horizontal-lift-derivative}
\end{equation}
and
\begin{equation}
 D\Pi_{\rm v}[h]
 =-(DJ[h])D{\cal B}-J\,D^2{\cal B}[h].
 \label{eq:v51-projection-derivative}
\end{equation}
\end{proposition}

\begin{proof}
Differentiate \(D\Psi\,J=\iota_W\).  Since the inclusion is constant,
\[
 D^2\Psi[h]J+D\Psi\,DJ[h]=0.
\]
Multiplication by \(D\Psi^{-1}\) proves
\eqref{eq:v51-horizontal-lift-derivative}.  Differentiate
\(\Pi_{\rm v}=I-JD{\cal B}\) to obtain
\eqref{eq:v51-projection-derivative}.
\end{proof}

If \(D\Psi^{-1}\) belongs to the V48 weighted algebra and the second jets are
finite-range or exponentially local, then
\(D\Pi_{\rm v}\) is also exponentially local by the weighted product
inequality.  Thus exact vertical projection need not destroy polymer
locality inside the V49 chamber.  The global harmonic and Killing projections
remain excluded in the macroscopic block.

\subsection{Application to slice and extension changes}

Let \(V_{\rm slice}\) generate an infinitesimal change of DeTurck gauge slice,
and let \(V_{\rm ext}\) generate an infinitesimal change of the trace
extension \(R_\Sigma\).  Their correct decomposition is
\[
 V_{\rm slice}= \Pi_{\rm v}V_{\rm slice}
                 +J(D{\cal B}V_{\rm slice}),
\]
\[
 V_{\rm ext}= \Pi_{\rm v}V_{\rm ext}
                 +J(D{\cal B}V_{\rm ext}).
\]
The first terms give zero conditional owners only in the measured combinations
\eqref{eq:v51-projected-owner}.  The second terms either induce the effective
owner \eqref{eq:v51-effective-owner}, when projectable, or must remain as
explicit fine-dependent transport defects.

\begin{assessment}[Upstream test before determinant cancellation]
\label{ass:v51-upstream-test}
For each proposed slice or extension deformation, the following finite card
is necessary:
\begin{enumerate}
\item compute \(D{\cal B}V\) with the exact nonlinear block map;
\item test projectability along the fine fibre;
\item compute the weighted vertical projection through \(D\Psi^{-1}\);
\item include \((D\Pi_{\rm v})[V]\) and the measure-density term in the
      divergence; and
\item assign the projectable horizontal remainder to the retained effective
      action and its Jacobian.
\end{enumerate}
A flat linear Ward zero supplies only the first-order check of item one.
\end{assessment}

\subsection{Bosonic boundary and graded debits}

The proof above is an ordinary finite-dimensional bosonic divergence theorem.
For noncompact bosonic cutoffs, vanishing boundary flux must be verified
uniformly before the cutoff is removed.  For fermionic owners, the analogous
statement uses the Berezinian and graded superdivergence.  Those extensions
are required for the full Standard Model measure; they are not supplied by
renaming the trace in \eqref{eq:v51-projected-divergence} a supertrace.

The next note will build the graded version in a finite Grassmann algebra and
show how the bosonic determinant, Faddeev--Popov determinant, and fermionic
Berezinian enter one measured fibre identity without double counting.  This
is the necessary precursor to connected determinant polymers.

\begin{assessment}[Exact status after V51]
\label{ass:v51-status}
V51 corrects the fibre dimension, constructs the exact chart-induced
vertical projection, proves the vertical and projectable-horizontal measured
identities, and identifies the second-jet content of the projection
divergence.  It does not verify projectability for a specific nonlinear
DeTurck or trace-extension flow, compute the associated second jets, prove
uniform bosonic cutoff flux, construct the graded identity, renormalise
determinant polymers, control macroscopic modes, or establish the
SM--Einstein continuum.
\end{assessment}

\clearpage
\section{V52: graded measured fibres and the determinant-role ledger}
\label{sec:v10v52}

\subsection{What is new relative to the earlier fermion notes}

Earlier notes proved three finite-algebra facts that are retained here:
Grassmann Schur integration, Berezin saturation of exact zero modes, and
boundary-state gluing with its separate CAR positivity condition.  The
missing statement is different.  V51 derived the measured
field-redefinition identity only on a bosonic fine fibre.  The full
SM--Einstein owner chart also contains physical fermions and
Faddeev--Popov ghosts, so its coordinate Jacobian is a Berezinian and its
divergence is a superdivergence.  This note supplies that graded fibre
identity and separates three determinant roles that must not be counted
twice.

Everything is at a finite regulator and in a finite Grassmann algebra.
A fixed ordering of odd generators and a fixed orientation of every
Berezin integral are part of the data.

\subsection{Coordinate convention and superdivergence}

Let \(X\) be a local superdomain of dimension \(p|q\), with even coordinates
\(x^1,\ldots,x^p\) and odd coordinates
\(\theta^1,\ldots,\theta^q\).  Write
\[
 Dz=dx^1\cdots dx^p\,d\theta^q\cdots d\theta^1
\]
for the chosen coordinate Berezin volume.  Let
\[
 d\mu=\varrho(z)\,Dz
\]
be an even Berezin density whose body is nowhere zero.  Let \(V\) be an even
supervector field,
\[
 V=\sum_iV^i\partial_{x^i}+\sum_\alpha V^\alpha\partial_{\theta^\alpha},
\]
so \(V^i\) is even and \(V^\alpha\) is odd.  All odd derivatives below are
left derivatives.

\begin{definition}[Measured superdivergence]
\label{def:v52-superdivergence}
The superdivergence is defined invariantly by
\[
 {\cal L}_Vd\mu=(\operatorname{sdiv}_\mu V)d\mu.
\]
In the chosen coordinates,
\begin{equation}
 \operatorname{sdiv}_\mu V
 =
 \varrho^{-1}\left\{
 \sum_i\partial_{x^i}(\varrho V^i)
 -\sum_\alpha\partial_{\theta^\alpha}(\varrho V^\alpha)
 \right\}.
 \label{eq:v52-superdivergence}
\end{equation}
\end{definition}

\begin{lemma}[Infinitesimal Berezinian]
\label{lem:v52-infinitesimal-Berezinian}
For the infinitesimal even change \(z\mapsto z+\varepsilon V(z)\),
where \(\varepsilon^2=0\) is even,
\begin{equation}
 \operatorname{Ber}(I+\varepsilon DV)
 =1+\varepsilon\,\operatorname{Str}(DV),
 \label{eq:v52-infinitesimal-Ber}
\end{equation}
where
\[
 \operatorname{Str}(DV)
 =\sum_i\partial_{x^i}V^i
  -\sum_\alpha\partial_{\theta^\alpha}V^\alpha.
\]
Consequently,
\[
 \operatorname{sdiv}_\mu V
 =\operatorname{Str}(DV)+V(\log\varrho).
\]
\end{lemma}

\begin{proof}
For an invertible even supermatrix \(J\),
\(D\log\operatorname{Ber}(J)[H]=\operatorname{Str}(J^{-1}H)\).
Apply this at \(J=I\) with \(H=\varepsilon DV\).  Expanding the density
\(\varrho(z+\varepsilon V)\) supplies \(V(\log\varrho)\).  Formula
\eqref{eq:v52-superdivergence} has the same expansion; for the odd term,
the minus sign converts
\(-(\partial_{\theta^\alpha}\varrho)V^\alpha\) into
\(V^\alpha\partial_{\theta^\alpha}\varrho\), because both factors are odd.
\end{proof}

\subsection{The graded vertical identity}

Let
\[
 {\cal B}:X\longrightarrow W
\]
be an even retained supermap of constant rank, and assume the graded version
of the V51 square chart,
\[
 \Psi=({\cal R},{\cal B}):X\longrightarrow F\times W,
\]
has invertible Berezinian body.  The horizontal lift and vertical projection
are still
\begin{equation}
 J=D\Psi^{-1}\iota_W,\qquad
 \Pi_{\rm v}=I-JD{\cal B}.
 \label{eq:v52-graded-projection}
\end{equation}
These are even supermatrices, and the algebraic proof of the projection
identities in V51 applies without a sign change.

Let \(S\) be an even action.  Assume the even boundary flux of the fine fibre
vanishes.  Odd integration has no boundary term because it is coefficient
extraction.

\begin{theorem}[Graded measured vertical redundancy]
\label{thm:v52-graded-vertical}
If an even vector field \(V\) is vertical,
\[
 D{\cal B}[V]=0,
\]
then, on every fine superfibre \(X_w\),
\begin{equation}
 \int_{X_w}
 \{DS[V]-\operatorname{sdiv}_{\mu_w}V\}
 e^{-S}\,d\mu_w=0.
 \label{eq:v52-graded-vertical-identity}
\end{equation}
For an arbitrary even \(V\), the same identity holds with
\(V\) replaced by \(\Pi_{\rm v}V\).
\end{theorem}

\begin{proof}
The graded product rule gives
\[
 \operatorname{sdiv}_{\mu_w}(e^{-S}V)
 =e^{-S}\{\operatorname{sdiv}_{\mu_w}V-DS[V]\},
\]
because \(S\) and \(V\) are even.  Integrate.  Ordinary divergence terms
vanish by the stated fine-boundary condition.  Every odd derivative
integrates to zero under Berezin integration, since it cannot contain the
top monomial in that generator.  This proves
\eqref{eq:v52-graded-vertical-identity}.  The range of
\(\Pi_{\rm v}\) is \(\ker D{\cal B}\), as in V51.
\end{proof}

This theorem is the graded measured redundancy statement.  Neither
\(DS[V]\) nor \(-\operatorname{sdiv}V\) vanishes separately in general.

\subsection{Projectable horizontal transport}

Call \(V\) projectable when
\[
 D{\cal B}(z)V(z)=\bar V({\cal B}(z))
\]
for an even retained supervector field \(\bar V\).  In product coordinates
\((r,w)=\Psi(z)\), disintegrate
\[
 d\mu=j(r,w)\,Dr\,Dw
\]
and define the retained even density
\begin{equation}
 Z(w)=\int_F e^{-S(r,w)}j(r,w)\,Dr.
 \label{eq:v52-graded-Z}
\end{equation}

\begin{theorem}[Graded measured transport]
\label{thm:v52-graded-transport}
For every projectable even \(V\),
\begin{equation}
 \int_F
 \{DS[V]-\operatorname{sdiv}_\mu V\}e^{-S}j\,Dr
 =
 -\operatorname{sdiv}_{Dw}\{Z(w)\bar V(w)\}.
 \label{eq:v52-graded-density-transport}
\end{equation}
If the body of \(Z\) is nonzero, the Grassmann logarithm exists and
\(S_{\rm eff}=-\log Z\) satisfies
\begin{equation}
 {\mathbb E}_w[
 DS[V]-\operatorname{sdiv}_\mu V]
 =
 DS_{\rm eff}[\bar V]-\operatorname{sdiv}_{Dw}\bar V.
 \label{eq:v52-graded-effective-owner}
\end{equation}
\end{theorem}

\begin{proof}
In product coordinates, the left side is the negative total
superdivergence of \(e^{-S}jV\), integrated in the fine variables.  Its fine
even part has zero boundary flux and its fine odd part integrates to zero.
Projectability makes the retained component independent of \(r\), leaving
exactly the retained superdivergence on the right side of
\eqref{eq:v52-graded-density-transport}.  If the body of \(Z\) is nonzero,
division by \(Z\) and
\(D Z[\bar V]=-Z\,DS_{\rm eff}[\bar V]\) give
\eqref{eq:v52-graded-effective-owner}.
\end{proof}

\begin{firewall}[Zero-mode sectors may have no logarithm]
\label{fire:v52-zero-body}
In a sector with unsaturated exact fermion zero modes, the body of \(Z\) can
vanish.  Then \(-\log Z\) and the normalized expectation in
\eqref{eq:v52-graded-effective-owner} are undefined.  The density identity
\eqref{eq:v52-graded-density-transport} remains meaningful and must be used
together with the zero-mode saturation and determinant-line ledger proved
earlier.
\end{firewall}

\subsection{The exact Berezinian Schur formula}

Let an even supermatrix have block form
\[
 {\mathcal J}=
 \begin{pmatrix}
 B&\Gamma\\
 \Delta&F
 \end{pmatrix},
\]
where \(B,F\) have even entries and
\(\Gamma,\Delta\) have odd entries.

\begin{proposition}[Berezinian Schur factorisation]
\label{prop:v52-Ber-Schur}
If \(F\) is invertible, then
\begin{equation}
 \operatorname{Ber}{\mathcal J}
 =
 \frac{\det(B-\Gamma F^{-1}\Delta)}{\det F}.
 \label{eq:v52-Ber-Schur-F}
\end{equation}
If \(B\) is invertible, equivalently
\begin{equation}
 \operatorname{Ber}{\mathcal J}
 =
 \frac{\det B}{\det(F-\Delta B^{-1}\Gamma)}.
 \label{eq:v52-Ber-Schur-B}
\end{equation}
All determinants are finite Grassmann-algebra determinants with the fixed
orientation convention.
\end{proposition}

\begin{proof}
For invertible \(F\), factor
\[
 {\mathcal J}
 =
 \begin{pmatrix}I&\Gamma F^{-1}\\0&I\end{pmatrix}
 \begin{pmatrix}B-\Gamma F^{-1}\Delta&0\\
                 \Delta&F\end{pmatrix}.
\]
The first factor has Berezinian one; the second is block triangular and has
Berezinian
\(\det(B-\Gamma F^{-1}\Delta)/\det F\).
The second formula follows from the analogous factorisation using \(B^{-1}\).
\end{proof}

This is the graded analogue of the nested Schur determinant in V44.
The odd block appears in the denominator of the coordinate Berezinian.  Under
an inverse coordinate map, that denominator is inverted and can become the
usual fermionic determinant factor.

\subsection{Three determinant roles and the no-double-counting rule}

Consider first the block-diagonal linear coordinate change
\[
 y=A_Bu,\qquad \eta=A_F\theta,
\]
where \(u,y\) are even and \(\theta,\eta\) are odd.  The inverse map has
Berezinian
\begin{equation}
 Du\,D\theta
 =(\det A_B)^{-1}(\det A_F)\,Dy\,D\eta.
 \label{eq:v52-linear-coordinate-density}
\end{equation}

\begin{theorem}[Determinant-role ledger]
\label{thm:v52-determinant-ledger}
At a fixed regulator the following operations are distinct:
\begin{enumerate}
\item A bosonic residual coordinate change contributes the inverse ordinary
      Jacobian, such as \((\det A_B)^{-1}\).
\item An odd residual coordinate change contributes the inverse odd
      Berezinian denominator, such as \(\det A_F\) in
      \eqref{eq:v52-linear-coordinate-density}.
\item Integrating a physical quadratic fermion action
      \(\overline\psi D_{\rm phys}\psi\) contributes
      \(\det D_{\rm phys}\), with zero modes and phase treated by the earlier
      saturation and determinant-line construction.
\item Integrating the ghost action
      \(\overline c M_{\rm FP}c\) contributes
      \(\det M_{\rm FP}\), representing the gauge-orbit Jacobian.
\end{enumerate}
These factors multiply only when they arise from distinct variables or
distinct maps.  If a physical or ghost operator itself is used as the odd
coordinate map, its determinant in the Berezinian is the same
change-of-variables factor and must not be inserted a second time.
\end{theorem}

\begin{proof}
Items one and two are the Berezin change-of-variables formula
\eqref{eq:v52-linear-coordinate-density}.  Items three and four are the
finite Gaussian identity
\[
 \int D\overline\chi\,D\chi\,
 e^{-\overline\chi M\chi}=\det M
\]
for the corresponding independent odd generators.  If the same linear map
\(M\) is used to replace \(\chi\) by \(M\chi\), the Gaussian computation can
be derived by that coordinate change; multiplying by another copy would
perform the same Jacobian operation twice.  Physical and ghost fields are
independent families and their operators have different domains and
meanings, so their determinants are not identified merely because both have
Grassmann representations.
\end{proof}

\begin{firewall}[Faddeev--Popov is not the DeTurck coordinate Jacobian]
\label{fire:v52-FP-distinct}
The Faddeev--Popov operator differentiates the gauge condition along the
diffeomorphism or internal-gauge orbit.  The DeTurck block differentiates a
gauge-fixed curvature map along metric coefficients.  Their determinants
remain separate even in the graded ledger.  A cancellation requires an
explicit common coordinate transformation, not matching notation.
\end{firewall}

\subsection{Weighted locality with nilpotent coefficients}

Give a matrix with Grassmann coefficients the coefficient norm obtained by
summing the weighted kernel norms of all Grassmann monomials.  Because the
Grassmann algebra is finite, this is a finite-dimensional Banach algebra.

\begin{proposition}[Local inverse over a finite Grassmann algebra]
\label{prop:v52-Grassmann-local-inverse}
Let
\[
 {\mathcal A}=A_{\rm body}+N
\]
be an even matrix in the weighted algebra, where \(A_{\rm body}^{-1}\)
belongs to that algebra and \(N\) has positive Grassmann degree.  Then
\({\mathcal A}\) is invertible and
\begin{equation}
 {\mathcal A}^{-1}
 =
 \sum_{j=0}^{q}
 (-A_{\rm body}^{-1}N)^jA_{\rm body}^{-1},
 \label{eq:v52-finite-nilpotent-inverse}
\end{equation}
where \(q\) may be taken as the number of odd generators.  Hence the inverse
has the same positive exponential locality reserve, up to the finite product
loss in the weighted norm.
\end{proposition}

\begin{proof}
Every product of more than \(q\) positive-degree Grassmann factors vanishes.
Thus the Neumann series terminates and multiplication verifies that
\eqref{eq:v52-finite-nilpotent-inverse} is the inverse.  Closure and
submultiplicativity of the V48 weighted algebra control each finite term.
\end{proof}

Consequently, the graded horizontal lift, vertical projection, and
Berezinian Schur blocks inherit weighted locality whenever their bosonic
bodies lie in the V49 chamber and all nilpotent coefficient jets are local.
This is an algebraic finite-regulator result; it does not supply uniform
bounds as the number of fermion modes grows.

\subsection{Anomaly and continuum debits}

The finite identity
\eqref{eq:v52-graded-vertical-identity} is exact for the chosen regulated
measure.  In a chiral theory, the continuum limit of the regulated
supertrace can contain a local anomaly.  Gauge anomaly cancellation must be
proved by the Standard Model representation trace identities and by a
regulator that preserves the required gauge covariance.  The finite
divergence theorem alone does not prove that the limiting Jacobian is one.

Likewise, a nonzero chiral determinant is a section of a determinant line
before a trivialisation is chosen.  The earlier phase and zero-mode ledger
remains in force; the scalar notation \(\det D_{\rm phys}\) in
Theorem~\ref{thm:v52-determinant-ledger} abbreviates that already established
finite-regulator choice.

\subsection{Next acquisition}

The graded measured identity now makes it legitimate to classify a local
Jacobian word as vertical, retained, physical-fermionic, or ghost.  The next
note will use that classification to expand a relative superdeterminant into
connected boundary words.  It will prove an absolute convergence criterion
in the weighted algebra and separate the extensive local counterterm rows
from the retained macroscopic determinant line.

\begin{assessment}[Exact status after V52]
\label{ass:v52-status}
V52 proves the finite graded vertical and horizontal transport identities,
the Berezinian Schur formula, the determinant-role ledger, and weighted
locality of inverses with nilpotent coefficients.  It does not verify the
specific SM gauge-anomaly traces in this cycle, control a growing Grassmann
algebra uniformly, construct the connected superdeterminant polymers,
renormalise their local rows, remove regulators, control determinant-line
phases globally, or establish the SM--Einstein continuum.
\end{assessment}

\clearpage
\section{V53: connected relative-superdeterminant words at the interface}
\label{sec:v10v53}

\subsection{Relation to the earlier heat-shell polymer theorem}

The earlier finite-heat construction proves connected determinant polymers by
Duhamel expansions, spatial tree decay, and a forest summation.  That theorem
continues to govern overlap and heat-shell determinants.  The present problem
has a different small parameter.  V49 provides a directly small relative
interface defect, while V48 puts it in a weighted convolution algebra.  This
note proves the corresponding algebraic closed-word expansion.  It does not
replace the heat-shell theorem and does not count its determinant a second
time.

Let \({\cal G}\) be the finite interface block graph.  At every vertex the
owner fibre has fixed finite superdimension.  Grassmann coefficients are
given the coefficient norm of V52.

\subsection{Relative Berezinian logarithm}

Let \({\mathcal J}_0\) and \({\mathcal J}\) be even invertible local
supermatrices in the same determinant-line chart, and set
\begin{equation}
 I+K={\mathcal J}_0^{-1}{\mathcal J}.
 \label{eq:v53-relative-K}
\end{equation}
Assume
\begin{equation}
 \eta_\kappa:=\|K\|_\kappa<1
 \label{eq:v53-small-K}
\end{equation}
in the weighted row/column norm of V48.

\begin{proposition}[Relative supertrace series]
\label{prop:v53-supertrace-series}
Under \eqref{eq:v53-small-K},
\begin{equation}
 \log\frac{\operatorname{Ber}{\mathcal J}}
               {\operatorname{Ber}{\mathcal J}_0}
 =
 \operatorname{Str}\log(I+K)
 =
 \sum_{n=1}^{\infty}\frac{(-1)^{n+1}}{n}
       \operatorname{Str}(K^n).
 \label{eq:v53-relative-superlog}
\end{equation}
The series converges absolutely in every Grassmann coefficient.  It uses the
logarithm branch obtained by the straight path \(I+tK\), \(0\leq t\leq1\).
\end{proposition}

\begin{proof}
The body norm is smaller than one, so \(I+tK_{\rm body}\) is invertible on
the whole path and the ordinary logarithm power series converges.  Positive
Grassmann degree is nilpotent; analytic substitution into the finite
nilpotent extension produces only finitely many additional coefficient
terms at each body coefficient.  The standard identity
\(\log\operatorname{Ber}=\operatorname{Str}\log\) holds along the invertible
path and gives \eqref{eq:v53-relative-superlog}.
\end{proof}

\begin{firewall}[The small chart does not trivialise the determinant line]
\label{fire:v53-line-chart}
Equation \eqref{eq:v53-relative-superlog} is a relative formula inside one
chosen local determinant-line chart.  A transition between charts can carry
a nontrivial phase cocycle.  Exact zero modes and the retained infrared block
remain outside \(K\).
\end{firewall}

\subsection{Closed words and connected support}

Write \(K_{xy}\) for the block from the fibre at \(y\) to the fibre at \(x\).
Then
\begin{equation}
 \operatorname{Str}(K^n)
 =
 \sum_{x_0,\ldots,x_{n-1}}
 \operatorname{str}\{
 K_{x_0x_1}K_{x_1x_2}\cdots K_{x_{n-1}x_0}\}.
 \label{eq:v53-closed-words}
\end{equation}
A term is a closed word.  Its visited sites, together with a fixed chosen
shortest lattice path for every jump \(x_i\to x_{i+1}\), form a connected
lattice polymer \(X\).  Allocate each word to that canonical polymer and let
\(a(X)\) be the sum of its coefficients in
\eqref{eq:v53-relative-superlog}.  Then
\begin{equation}
 \log\frac{\operatorname{Ber}{\mathcal J}}
               {\operatorname{Ber}{\mathcal J}_0}
 =\sum_{X\subset{\cal G}\ {\rm connected}}a(X).
 \label{eq:v53-polymer-decomposition}
\end{equation}

Let \(\tau(X)\) be the length of a minimum spanning tree of \(X\) in the
lattice metric.  Let \(\nu\) be a constant such that the coefficient norm of
a local supertrace is at most \(\nu\) times the matrix coefficient norm.

\begin{theorem}[Anchored relative-superdeterminant bound]
\label{thm:v53-anchored-bound}
The activities in \eqref{eq:v53-polymer-decomposition} satisfy
\begin{equation}
 \sup_{x\in{\cal G}}
 \sum_{X\ni x}e^{\kappa\tau(X)}\|a(X)\|_{\rm coeff}
 \leq
 \nu\,\frac{\eta_\kappa}{1-\eta_\kappa}.
 \label{eq:v53-anchored-polymer}
\end{equation}
The bound is independent of the number of interface blocks.
\end{theorem}

\begin{proof}
For a closed word visiting \(X\), deleting one edge from its closed walk
leaves a connected spanning walk, so
\[
 \tau(X)\leq\sum_{i=0}^{n-1}d(x_i,x_{i+1}),
 \qquad x_n=x_0.
\]
The exponential tree weight is therefore bounded by the product of the
edge weights in the definition of \(\|K\|_\kappa\).

Fix \(x\).  A word whose polymer contains \(x\) has at least one occurrence
of \(x\).  Mark one such occurrence and cyclically rotate the word to start
there.  There are at most \(n\) choices, which cancels the factor \(1/n\) in
\eqref{eq:v53-relative-superlog}.  Successive weighted row sums are each
bounded by \(\eta_\kappa\), and the local supertrace costs at most \(\nu\).
Thus the total absolute contribution of length \(n\) is at most
\(\nu\eta_\kappa^n\).  Summing \(n\geq1\) proves
\eqref{eq:v53-anchored-polymer}.
\end{proof}

This is stronger than an extensive trace bound for local purposes: every
marked interface block has a volume-free activity sum.  The price is the
strict small-relative-defect condition \(\eta_\kappa<1\).

\begin{corollary}[Analytic marked responses]
\label{cor:v53-marked-response}
If \(K(q)\) is analytic in a local retained owner \(q\) and
\(\sup_{q}\|K(q)\|_\kappa\leq\eta_\kappa<1\) on a complex polydisc, then every
mixed \(q\)-derivative of the activities has an anchored bound obtained from
\eqref{eq:v53-anchored-polymer} by the corresponding Cauchy radii and
factorials.
\end{corollary}

\begin{proof}
The series is normally convergent in the smaller polydisc by
Theorem~\ref{thm:v53-anchored-bound}.  Differentiate by Cauchy's formula and
spend the chosen analytic-radius reserve.
\end{proof}

\subsection{Several determinant sectors}

The V52 ledger forbids combining physically different determinants by
notation.  Let the finite set of sectors be
\[
 {\mathscr S}
 =\{\text{chart Berezinian},\text{physical fermions},
   \text{ghosts},\text{bosonic Gaussian blocks}\}.
\]
For sector \(s\), let \(K_s\) be its relative local operator and let
\(\epsilon_s\) be its actual logarithmic multiplicity.  Examples are
\(\epsilon_s=1\) for a determinant,
\(\epsilon_s=-1\) for an inverse determinant, and
\(\epsilon_s=-\frac12\) for the usual positive bosonic Gaussian determinant
in the partition density.  Define
\begin{equation}
 {\cal L}_{\rm loc}
 :=\sum_{s\in{\mathscr S}}\epsilon_s
   \operatorname{Tr}_s\log(I+K_s),
 \label{eq:v53-sector-ledger}
\end{equation}
where \(\operatorname{Tr}_s\) is the supertrace only for a genuine
coordinate supermatrix and the ordinary trace for a Gaussian operator.

\begin{proposition}[Sectorwise polymer bound]
\label{prop:v53-sectorwise-bound}
If \(\eta_{s,\kappa}:=\|K_s\|_\kappa<1\) for every sector, then
\({\cal L}_{\rm loc}\) has connected activities \(a_{\rm tot}(X)\) obeying
\begin{equation}
 \sup_x\sum_{X\ni x}e^{\kappa\tau(X)}
 \|a_{\rm tot}(X)\|_{\rm coeff}
 \leq
 \sum_{s\in{\mathscr S}}
 |\epsilon_s|\nu_s
 \frac{\eta_{s,\kappa}}{1-\eta_{s,\kappa}}.
 \label{eq:v53-sectorwise-polymer}
\end{equation}
\end{proposition}

\begin{proof}
Apply Theorem~\ref{thm:v53-anchored-bound} to each sector and then use the
triangle inequality.
\end{proof}

The absolute bound deliberately takes no credit for boson--fermion or
ghost--gauge cancellation.  A cancellation may improve a coefficient only
after the two activities have been expressed in the same local owner basis
and an exact Ward or measured-fibre identity proves equality.

\subsection{Exact cancellations and commutator words}

\begin{lemma}[Graded cyclic cancellation]
\label{lem:v53-graded-commutator}
For homogeneous finite supermatrices,
\[
 \operatorname{Str}[A,B]_{\rm gr}=0,
 \qquad
 [A,B]_{\rm gr}=AB-(-1)^{|A||B|}BA.
\]
Consequently, a relative logarithmic variation that is an exact graded
commutator contributes zero to every complete trace coefficient.
\end{lemma}

\begin{proof}
The graded cyclicity identity
\(\operatorname{Str}(AB)=(-1)^{|A||B|}\operatorname{Str}(BA)\)
gives the first statement.  Apply it coefficientwise to the variation
series for the second.
\end{proof}

For an exact vertical coordinate flow, V52 gives the stronger measured
identity: the action variation cancels the superdivergence after fine-fibre
integration.  Lemma~\ref{lem:v53-graded-commutator} handles only the algebraic
commutator part of that cancellation.  The density derivative and the
derivative of the exact vertical projection must still be included.

\begin{firewall}[A zero total trace need not mean zero polymer by polymer]
\label{fire:v53-local-cancellation}
Cyclic or anomaly cancellation of the global trace does not automatically
cancel each chosen polymer activity.  A local cancellation requires a
support-preserving identity or a controlled coboundary that can be moved to
the polymer boundary.  Otherwise the cancellation may occur only after a
volume sum and is unusable for a local RG norm.
\end{firewall}

\subsection{Relevant projection and the retained block}

Let \({\cal P}_{\leq4}\) denote the already constructed covariant projection
onto the local SM--Einstein owner rows of engineering dimension at most four.
Apply it to each connected activity:
\[
 a(X)={\cal P}_{\leq4}a(X)
      +(I-{\cal P}_{\leq4})a(X).
\]
The first term feeds the running cosmological, Einstein, gauge, Higgs,
Yukawa, and other previously enumerated marginal/relevant rows.  The second
term is eligible for cofinal contraction only after its scale gain is proved.
The anchored bound controls the support sum but supplies no engineering
dimension by itself.

Exact zero modes, harmonic/Killing owners, and the low chiral determinant
line are not included in \(K_s\).  Their finite-dimensional factor is
\[
 {\cal D}_{\rm macro}(w),
\]
carried with the retained block.  Thus the regulated factorisation has the
form
\begin{equation}
 {\cal D}_{\rm total}
 =
 {\cal D}_{\rm macro}\,
 \exp\left\{\sum_Xa_{\rm tot}(X)\right\},
 \label{eq:v53-macro-local-factor}
\end{equation}
within one overlap chart.  No inverse of the macroscopic zero-mode block is
used.

\subsection{The remaining gate}

Theorem~\ref{thm:v53-anchored-bound} turns the V49 weighted defect into a
connected local determinant expansion whenever
\(\|{\mathcal J}_0^{-1}({\mathcal J}-{\mathcal J}_0)\|_\kappa<1\).
To use it along a full RG trajectory, the relative operator must be reset at
each scale so this smallness is restored.  Exact determinant cocycles then
combine adjacent steps without duplicating the shared base determinant, as
proved in the earlier shell notes.

The next note will formulate that scale-reset cocycle for the combined
sector ledger and identify the exact contact term produced when two
overlapping relative charts are composed.  It will keep all dimension at
most four contacts in the running owner bank and test whether vertical
measured contacts cancel before estimation.

\begin{assessment}[Exact status after V53]
\label{ass:v53-status}
V53 proves the relative Berezinian trace series, its exact connected
closed-word decomposition, a volume-free anchored polymer bound, and the
sectorwise determinant ledger.  It does not prove local anomaly cancellation,
evaluate the relevant owner projections, establish the scale-reset
smallness, control determinant-line transitions, remove regulators, or
establish the SM--Einstein continuum.
\end{assessment}

\clearpage
\section{V54: scale-reset cocycles and the exact reblocking contact}
\label{sec:v10v54}

\subsection{The missing cocycle}

The earlier notes contain a heat-interval cocycle and determinant-line gauge
cocycles.  V53 instead uses a small relative interface operator
\(I+K\).  A long trajectory need not remain in one ball
\(\|K\|_\kappa<1\), so the relative base must be reset.  The point of this
note is that resetting the base does not duplicate determinants, provided
the exact multiplicative cocycle is used before estimates are made.

Let
\[
 {\mathcal J}_0,{\mathcal J}_1,\ldots,{\mathcal J}_N
\]
be even invertible supermatrices in one transported local determinant-line
chart.  Define
\begin{equation}
 R_j={\mathcal J}_{j-1}^{-1}{\mathcal J}_j=I+K_j,
 \qquad 1\leq j\leq N.
 \label{eq:v54-step-relative}
\end{equation}

\begin{theorem}[Scale-reset Berezinian cocycle]
\label{thm:v54-reset-cocycle}
The exact identity
\begin{equation}
 \frac{\operatorname{Ber}{\mathcal J}_N}
      {\operatorname{Ber}{\mathcal J}_0}
 =
 \prod_{j=1}^N\operatorname{Ber}(I+K_j)
 \label{eq:v54-Ber-telescope}
\end{equation}
holds without commutativity assumptions.  If
\(\|K_j\|_\kappa<1\) for every \(j\), and every logarithm is chosen by the
straight path \(I+tK_j\), then the logarithm along the concatenated path is
\begin{equation}
 \log_{\rm path}
 \frac{\operatorname{Ber}{\mathcal J}_N}
      {\operatorname{Ber}{\mathcal J}_0}
 =
 \sum_{j=1}^N\operatorname{Str}\log(I+K_j).
 \label{eq:v54-log-telescope}
\end{equation}
\end{theorem}

\begin{proof}
Multiplicativity of the Berezinian gives
\[
 \prod_{j=1}^N
 \operatorname{Ber}({\mathcal J}_{j-1}^{-1}{\mathcal J}_j)
 =
 \operatorname{Ber}({\mathcal J}_0^{-1}{\mathcal J}_N),
\]
because all intermediate scalar Berezinians cancel.  This proves
\eqref{eq:v54-Ber-telescope}.  Each straight path stays invertible by the
small-norm condition.  Concatenating those paths defines a continuous
logarithm whose increment on step \(j\) is
\(\operatorname{Str}\log(I+K_j)\).  Adding the increments proves
\eqref{eq:v54-log-telescope}.
\end{proof}

Although
\[
 I+K_{02}=(I+K_{01})(I+K_{12})
\]
contains the product \(K_{01}K_{12}\), no extra determinant contact is
created.  Its apparent mixed trace words are exactly the re-expansion of the
sum of the two logarithms.  A genuine contact below comes from changing the
local owner projection during reblocking, not from failure of determinant
multiplicativity.

\subsection{The multi-step anchored bound}

Let \(a_j(X)\) be the connected activities supplied by V53 for step \(j\).
Let
\[
 {\mathscr R}_{j\to N}
\]
transport and reblock step-\(j\) activities to the final interface graph.
Assume its anchored norm bound is
\begin{equation}
 \|{\mathscr R}_{j\to N}a_j\|_{\rm anc,\kappa_N}
 \leq C_{jN}\|a_j\|_{\rm anc,\kappa_j},
 \label{eq:v54-reblock-bound}
\end{equation}
where
\[
 \|a\|_{\rm anc,\kappa}
 :=\sup_x\sum_{X\ni x}e^{\kappa\tau(X)}\|a(X)\|_{\rm coeff}.
\]

\begin{theorem}[Summable scale-reset gate]
\label{thm:v54-scale-sum}
Suppose
\[
 \eta_j:=\|K_j\|_{\kappa_j}<1
\]
and
\begin{equation}
 \sum_{j=1}^{\infty}
 C_{j\infty}\nu_j\frac{\eta_j}{1-\eta_j}<\infty.
 \label{eq:v54-scale-summability}
\end{equation}
Then the reblocked local logarithms
\[
 \sum_{j\geq1}{\mathscr R}_{j\to\infty}a_j
\]
converge absolutely in the final anchored polymer norm.  Every finite partial
sum equals the exact relative logarithm in
\eqref{eq:v54-log-telescope}, expressed on the final graph.
\end{theorem}

\begin{proof}
V53 gives
\[
 \|a_j\|_{\rm anc,\kappa_j}
 \leq\nu_j\frac{\eta_j}{1-\eta_j}.
\]
Apply \eqref{eq:v54-reblock-bound} and sum.  Condition
\eqref{eq:v54-scale-summability} is the Weierstrass criterion in the
anchored Banach space.  The equality of each finite partial sum with the
relative logarithm is Theorem~\ref{thm:v54-reset-cocycle}; reblocking changes
its representation, not its scalar value.
\end{proof}

\begin{firewall}[Resetting smallness is not summability]
\label{fire:v54-reset-not-sum}
The inequalities \(\eta_j<1\) make every individual chart expansion valid.
They do not imply \eqref{eq:v54-scale-summability}.  Marginal determinant
rows can contribute an order-one amount at every scale.  Those rows must be
extracted into running couplings before the remainder is summed.
\end{firewall}

\subsection{The exact projection--reblocking contact}

Let \({\cal P}_j\) be the covariant relevant/marginal owner projection at
scale \(j\), and let \({\mathscr R}_j\) reblock one scale.  For any local
logarithmic activity \(\ell_j\), define
\begin{equation}
 {\mathcal C}_j\ell_j
 :=
 {\cal P}_{j+1}{\mathscr R}_j\ell_j
 -{\mathscr R}_j{\cal P}_j\ell_j.
 \label{eq:v54-contact-definition}
\end{equation}

\begin{proposition}[Exact contact identity]
\label{prop:v54-contact}
The reblocked local logarithm splits exactly as
\begin{equation}
 {\cal P}_{j+1}{\mathscr R}_j\ell_j
 =
 {\mathscr R}_j{\cal P}_j\ell_j
 +{\mathcal C}_j\ell_j,
 \label{eq:v54-contact-identity}
\end{equation}
and
\begin{equation}
 (I-{\cal P}_{j+1}){\mathscr R}_j\ell_j
 =
 {\mathscr R}_j(I-{\cal P}_j)\ell_j
 -{\mathcal C}_j\ell_j.
 \label{eq:v54-remainder-contact}
\end{equation}
Thus the same contact is added to the running local owner and removed from
the irrelevant remainder.  Dropping it changes the exact determinant.
\end{proposition}

\begin{proof}
Equation \eqref{eq:v54-contact-identity} is the definition of
\({\mathcal C}_j\).  Subtract it from
\[
 {\mathscr R}_j\ell_j
 ={\mathscr R}_j{\cal P}_j\ell_j
  +{\mathscr R}_j(I-{\cal P}_j)\ell_j
\]
to obtain \eqref{eq:v54-remainder-contact}.
\end{proof}

The contact measures the failure of projection to commute with reblocking.
It includes, for example, a curved-background local term that is irrelevant
in fine coordinates but acquires a dimension at most four Taylor jet after
coarse rescaling.  Such a term is neither a new determinant nor an
estimation error; it is part of the exact running coupling.

\begin{corollary}[Sufficient remainder summability]
\label{cor:v54-remainder-sum}
Suppose the transported irrelevant part and contact obey
\[
 \|{\mathscr R}_j(I-{\cal P}_j)\ell_j\|_{\rm anc}
 \leq A_j,\qquad
 \|{\mathcal C}_j\ell_j\|_{\rm anc}\leq B_j,
\]
with
\[
 \sum_j(A_j+B_j)<\infty.
\]
Then the exact coarse irrelevant remainders in
\eqref{eq:v54-remainder-contact} are absolutely summable.  The projected
terms form the explicit running owner sequence and are not required to be
summable.
\end{corollary}

\begin{proof}
Apply the triangle inequality to
\eqref{eq:v54-remainder-contact} and sum.  The projected sequence is retained
by definition.
\end{proof}

\subsection{Sectorwise cocycles and matching bases}

Apply Theorem~\ref{thm:v54-reset-cocycle} separately to every sector of the
V52 determinant-role ledger.  For physical fermions, ghosts, bosonic
Gaussians, and the coordinate Berezinian, the intermediate base must be
exactly the terminal operator of the preceding step in the same transported
frame.  If a new frame \(C_j\) is inserted,
\[
 {\mathcal J}_j\mapsto{\mathcal J}_jC_j,
\]
then its Berezinian is an explicit transition factor.  It cancels between
adjacent steps only when the inverse transition is included on the other
side.

\begin{proposition}[Frame-transition ledger]
\label{prop:v54-frame-transition}
Let the right frame at the end of step \(j\) differ from the left frame of
step \(j+1\) by an even invertible \(C_j\).  Then exact telescoping requires
the pair of factors
\[
 \operatorname{Ber}(C_j)\operatorname{Ber}(C_j^{-1})=1.
\]
Omitting either factor leaves the transition cocycle in the total
determinant.  For a chiral low-mode block this factor belongs to the retained
determinant line and need not admit a single global logarithm.
\end{proposition}

\begin{proof}
Insert \(C_jC_j^{-1}\) between the two relative maps and use
multiplicativity of the Berezinian.  The determinant-line qualification is
the previously established obstruction to choosing one global frame.
\end{proof}

\subsection{Measured vertical contacts}

Suppose a complete step deformation is generated by an exactly vertical
even supervector field and obeys the fine-boundary conditions of V52.  Its
action and coordinate-density changes occur in the measured combination
\[
 DS[V]-\operatorname{sdiv}_\mu V.
\]
The conditional integral of that complete combination is zero.  Therefore
it must be removed before sectorwise absolute values are taken.

\begin{proposition}[Order of vertical cancellation and projection]
\label{prop:v54-vertical-order}
If a step owner is proved to be an exact measured vertical owner, its
effective retained density is zero before reblocking.  Hence its exact
projection, contact, and irrelevant remainder all vanish when computed from
the combined measured owner.  Projecting the action variation and
superdivergence separately need not give zero terms and can create a
spurious contact.
\end{proposition}

\begin{proof}
The graded vertical identity of V52 says that the combined conditional
density is identically zero as a function of the retained owner.  Linear
reblocking and projection send the zero function to zero.  The separate
summands are not zero by that theorem, so no corresponding conclusion holds
after separating them.
\end{proof}

This proposition is conditional on exact nonlinear verticality.  A flat Ward
zero or a principal-symbol commutator is insufficient.

\subsection{Cofinal scale ledger}

For each scale the combined determinant card now has the following exact
order:

\begin{enumerate}
\item pair action and superdivergence terms and remove only proved vertical
      measured owners;
\item factor high-shell and retained blocks without inverting the retained
      zero modes;
\item form each sector's small relative operator \(I+K_{s,j}\);
\item expand its connected activities by V53;
\item combine all sectors in the common local owner basis;
\item apply the scale projection and retain the contact
      \({\mathcal C}_j\); and
\item reblock only the irrelevant remainder and test the summability gate.
\end{enumerate}

The next note is the compulsory V55 YM/NS audit.  It will inspect whether the
current companion programmes supply a sharper verticality card, contact
cancellation, or restricted-direction norm.  The following post-audit note
will then attack the first model-specific item not closed by the abstract
ledger: the Standard Model gauge-anomaly trace and its locality at one
regulated cell, unless that calculation is already complete in the attached
notes.

\begin{assessment}[Exact status after V54]
\label{ass:v54-status}
V54 proves the scale-reset Berezinian cocycle, an absolute anchored
multi-scale gate, the exact projection--reblocking contact, and the
frame-transition ledger.  It does not establish the required scale gains,
compute the contacts, verify exact verticality for a specific chart flow,
reprove or complete all SM anomaly traces, control the retained determinant
line globally, remove regulators, or establish the SM--Einstein continuum.
\end{assessment}

\clearpage
\section{V55: compulsory audit and the divergence-compensation decision}
\label{sec:v10v55}

\subsection{Companion files inspected}

The mandatory audit inspected
\[
 \texttt{Renewal\_Yang\_Mills\_Mass\_Gap\_Research\_Notes\_V60.tex},
 \qquad
 \texttt{SHA--256 }CB265825C8CDF7D8359E23530A107B6A791CBD551429D75179806CCB072C13DB,
\]
and
\[
 \texttt{Renewal\_NS\_Stability\_Research\_Notes\_V26.tex},
 \qquad
 \texttt{SHA--256 }82C887F8C2C69E4F28FAD7C3DC8DE5B9099CB5A4BF431EBA1FFF3E91935978AD.
\]
The NS companion is unchanged.  The YM companion has advanced through a
conditional-current locality compiler and an exact compact-Haar
divergence calculation.

\subsection{The new YM mechanism}

YM V60 starts with a local equation-of-motion vector field having the desired
linearisation.  Its Haar divergence creates a nonzero quadratic marginal
owner.  A cubic gradient field has zero linearisation but a nonzero quadratic
divergence jet.  Choosing one explicit coefficient cancels the marginal
Jacobian while leaving the prescribed linearisation unchanged.  The note
also keeps two statements separate: the marginal is removable by a local
coordinate section, while its value in the unmodified raw pushed current
still requires a higher connected-current jet.

\begin{assessment}[Transfer to the combined owner chart]
\label{ass:v55-ym-transfer}
The transferable principle is a divergence-jet compensation in exact fine
fibre coordinates.  If a chart deformation is already vertical for the
nonlinear retained map, a higher-order vertical correction can cancel an
unwanted marginal divergence jet without changing its linearisation.  This
does not calculate the raw pushforward coefficient and does not make a
nonvertical field redundant.
\end{assessment}

\begin{proof}
The V51--V52 measured identity applies to any exact vertical vector field.
Adding a second vertical field preserves verticality.  If the correction
vanishes to third order at the reference point, its first derivative is zero.
Its ordinary divergence begins at quadratic order, while its contraction
with the score of a stationary smooth density begins at fourth order.
Therefore the quadratic divergence jet can be adjusted independently of the
linearisation.  V56 will prove the surjectivity and give the explicit
correction.
\end{proof}

This mechanism enters before the V54 relevant projection.  A proved measured
vertical owner is cancelled as a combined action--divergence pair before
sectorwise bounds; a merely removable coordinate coefficient is recorded as
a choice of retained section.

\subsection{The unchanged NS result}

NS V26 still shows that a rank-one restriction produces no first-derivative
gain and only a modest second-derivative gain on an intermediate radial
range.  It supplies no new domain reduction for the present fibre chart.

\begin{assessment}[No NS change of direction]
\label{ass:v55-ns-transfer}
The V49 weighted defect and V53 determinant activities continue to use the
full proved local owner domain.  No rank-one replacement or improved
interface floor follows from the current NS note.
\end{assessment}

\subsection{Nonduplication check against the SM anomaly ledger}

The accumulated SM notes already contain:

\begin{enumerate}
\item the perturbative anomaly-polynomial ledger, including the Standard
      Model hypercharge sums;
\item the even weak-doublet count and the resulting cancellation of the
      \(SU(2)\) mod-two sign;
\item the four possible global gauge groups
      \((SU(3)\times SU(2)\times U(1))/\Gamma_n\);
\item the spin-bordism anomaly-character test and its representation-level
      cancellation; and
\item the mixed \(U(1)\)--gravity hypercharge sum.
\end{enumerate}

Those results are representation-level and topological.  Their remaining
debit is the analytic construction of a local regulated measure and the
uniform locality of its phase/current.  Repeating the charge arithmetic in
V56 would not advance that debit.

\subsection{Post-audit direction}

Use the exact product chart of V51,
\[
 (r,w)=\Psi(x),
\]
where \(r\in{\mathbb R}^n\) are bosonic fine-fibre coordinates and \(w\) is
retained.  Let the reference point be \(r=0\), let the conditional density be
\(\rho(r,w)\,dr\), and assume its score vanishes at the reference point:
\[
 D_r\log\rho(0,w)=0.
\]
For a homogeneous quadratic local owner \(q_2(r,w)\), the candidate cubic
correction is
\[
 W_3(r,w)=\frac{r\,q_2(r,w)}{n+2},
\]
up to the sign required by the measured owner.

\begin{decision}[V56 acquisition]
\label{dec:v55-v56}
V56 will prove:
\begin{enumerate}
\item the homogeneous identity
      \(\operatorname{div}(r q_2)=(n+2)q_2\);
\item the quadratic-jet surjectivity of cubic vertical fields for a general
      stationary smooth conditional density;
\item preservation of the prescribed vector-field linearisation and absence
      of a quadratic action contribution at a stationary action;
\item exact nonlinear verticality after pullback through the V51 fibre chart;
      and
\item a weighted locality bound for the pulled-back correction.
\end{enumerate}
It will distinguish cancellation of a chosen coordinate-section marginal
from calculation of the raw pushed-current coefficient.
\end{decision}

This goes upstream of the connected determinant projection.  It supplies a
finite card for removing chart-dependent quadratic Jacobian rows before
testing the cofinal summability of the physical remainder.

\subsection{Audit status}

\begin{assessment}[Exact boundary of V55]
\label{ass:v55-status}
V55 records the current companion state, identifies the YM
divergence-compensation mechanism, and avoids repeating the completed
representation-level anomaly ledger.  It does not yet prove the local
compensation theorem, calculate a raw SM pushforward current, establish
four-point cluster bounds, control the determinant line globally, remove
regulators, or establish the SM--Einstein continuum.
\end{assessment}

\clearpage
\section{V56: local vertical compensation of marginal divergence jets}
\label{sec:v10v56}

\subsection{Fine-fibre coordinates and the problem}

Use the exact V51 product chart
\[
 (r,w)=\Psi(x),
\]
where \(r=(r_1,\ldots,r_N)\) are bosonic fine-fibre coordinates and \(w\)
is retained.  The actual conditional density is
\[
 d\mu_w=\rho(r,w)\,dr,
 \qquad \rho>0.
\]
Its measured divergence is
\begin{equation}
 \operatorname{div}_{\mu_w}W
 =\operatorname{div}_rW+D_r\log\rho[W].
 \label{eq:v56-measured-divergence}
\end{equation}
A raw vertical chart generator can have the desired linearisation but an
unwanted quadratic or higher local divergence jet.  The task is to alter
that coordinate section without altering the generator's first jet.

The construction is local.  If a homogeneous owner \(q(r)\) depends only on
a coordinate set \(I\subset\{1,\ldots,N\}\), write
\[
 d_I=|I|,\qquad
 E_I=\sum_{i\in I}r_i\partial_{r_i}
\]
for the local Euler field.

\subsection{The homogeneous divergence right inverse}

\begin{lemma}[Local Euler divergence identity]
\label{lem:v56-Euler-divergence}
Let \(q_m(r_I,w)\) be homogeneous of degree \(m\geq0\) in the coordinates
\(r_I\).  Define the vector field, supported in the same coordinate set,
by
\begin{equation}
 ({\mathcal H}_{I,m}q_m)_i
 =
 \begin{cases}
 \displaystyle\frac{r_iq_m}{d_I+m},&i\in I,\\
 0,&i\notin I.
 \end{cases}
 \label{eq:v56-Euler-homotopy}
\end{equation}
Then
\begin{equation}
 \operatorname{div}_r({\mathcal H}_{I,m}q_m)=q_m.
 \label{eq:v56-right-inverse}
\end{equation}
\end{lemma}

\begin{proof}
Only coordinates in \(I\) contribute.  Euler's homogeneous identity gives
\[
 \sum_{i\in I}\partial_{r_i}(r_iq_m)
 =d_Iq_m+E_Iq_m=(d_I+m)q_m.
\]
Division by \(d_I+m\) proves \eqref{eq:v56-right-inverse}.
\end{proof}

For \(m=2\), the correction is cubic and therefore has zero value, first
derivative, and second derivative at \(r=0\).  Explicitly,
\begin{equation}
 W_3={\mathcal H}_{I,2}q_2
 =\frac{r_Iq_2}{d_I+2}.
 \label{eq:v56-cubic-correction}
\end{equation}
This is the Euclidean fine-fibre analogue of the cubic Haar compensation
found in the audited YM note.

\subsection{Measured divergence is triangular in polynomial degree}

Let
\[
 s(r,w)=D_r\log\rho(r,w)
\]
be the score.  If \(W_{m+1}\) is homogeneous of degree \(m+1\), then
\[
 \operatorname{div}_rW_{m+1}
\]
has degree \(m\), while \(s[W_{m+1}]\) has degree at least \(m+1\).
If \(s(0,w)=0\), it has degree at least \(m+2\).

\begin{theorem}[Finite measured-divergence jet surjectivity]
\label{thm:v56-jet-surjectivity}
Fix \(M\geq2\).  Let
\[
 q_{\leq M}=\sum_{m=2}^Mq_m
\]
be a local polynomial with each \(q_m\) homogeneous of degree \(m\) and
supported in a fixed finite coordinate set \(I\).  Then there is a polynomial
vector field
\[
 W=\sum_{m=2}^MW_{m+1},
\]
supported in \(I\), with \(W_{m+1}\) homogeneous of degree \(m+1\), such that
\begin{equation}
 {\mathcal T}_{[2,M]}
 \operatorname{div}_{\mu_w}W=q_{\leq M}.
 \label{eq:v56-prescribed-divergence-jet}
\end{equation}
Here \({\mathcal T}_{[2,M]}\) denotes the Taylor projection onto degrees
two through \(M\) in \(r\).  In particular, \(W(0)=0\) and \(DW(0)=0\).
If \(s(0,w)=0\), the quadratic correction
\(W_3={\mathcal H}_{I,2}q_2\) creates no cubic divergence term.
\end{theorem}

\begin{proof}
Proceed by degree.  At degree two, set
\(W_3={\mathcal H}_{I,2}q_2\).  By
Lemma~\ref{lem:v56-Euler-divergence}, the ordinary divergence contributes
exactly \(q_2\).  The score term has degree at least three, so it does not
alter the degree-two equation.

Suppose \(W_3,\ldots,W_m\) have been chosen so that the measured divergence
agrees with \(q_{\leq M}\) through degree \(m-1\).  Let \(e_m\) be the
degree-\(m\) residual still required.  Set
\[
 W_{m+1}={\mathcal H}_{I,m}e_m.
\]
Its ordinary divergence is \(e_m\), while its score contribution begins in
degree \(m+1\).  Thus it corrects degree \(m\) without changing any lower
degree.  Induction through \(M\) proves
\eqref{eq:v56-prescribed-divergence-jet}.  Every summand has degree at least
three, proving the first-jet statement.  If \(s(0,w)=0\), the score times
\(W_3\) begins in degree four, proving the last assertion.
\end{proof}

The solution is not unique.  One may add a measured-divergence-free local
field or any field whose divergence begins above degree \(M\).  This freedom
is a coordinate-section freedom, not a new physical coupling.

\subsection{Cancellation of a raw marginal Jacobian jet}

Let \(V\) be an exact vertical generator in the fine chart and write
\[
 j_{\leq M}
 :={\mathcal T}_{[2,M]}
 \operatorname{div}_{\mu_w}V.
\]
Apply Theorem~\ref{thm:v56-jet-surjectivity} with
\(q_{\leq M}=-j_{\leq M}\), and let \(W\) be the resulting correction.

\begin{corollary}[Divergence-compensated vertical generator]
\label{cor:v56-compensated-generator}
The field
\[
 V^{\rm comp}=V+W
\]
is vertical, has the same value and linearisation as \(V\) at \(r=0\), and
satisfies
\begin{equation}
 {\mathcal T}_{[2,M]}
 \operatorname{div}_{\mu_w}V^{\rm comp}=0.
 \label{eq:v56-cancelled-divergence}
\end{equation}
\end{corollary}

\begin{proof}
Both fields have only fine-coordinate components, hence are vertical in the
product chart.  The correction has zero value and first derivative by the
theorem.  Linearity of divergence and
\eqref{eq:v56-prescribed-divergence-jet} give
\eqref{eq:v56-cancelled-divergence}.
\end{proof}

Suppose also that the action is stationary in the fine directions at the
reference point:
\[
 D_rS(0,w)=0.
\]
Since \(W=O(|r|^3)\),
\begin{equation}
 DS[W]=O(|r|^4).
 \label{eq:v56-action-order}
\end{equation}
Therefore the cubic correction that removes a quadratic Jacobian marginal
does not change the quadratic or cubic action jet.  If the action is not
stationary, it still changes no quadratic jet, but a cubic action term may
appear and must be retained.

\begin{proposition}[Measured-owner effect at quadratic order]
\label{prop:v56-measured-owner-quadratic}
At a stationary action and density reference point, let \(q_2\) be the
quadratic divergence jet of \(V\), and choose
\[
 W_3=-{\mathcal H}_{I,2}q_2.
\]
Then
\[
 {\mathcal T}_2
 \{DS[V+W_3]-\operatorname{div}_{\mu_w}(V+W_3)\}
 =
 {\mathcal T}_2DS[V].
\]
Thus the quadratic coordinate-Jacobian owner is removed while the action
quadratic jet is unchanged.
\end{proposition}

\begin{proof}
Theorem~\ref{thm:v56-jet-surjectivity} gives
\({\mathcal T}_2\operatorname{div}_{\mu_w}W_3=-q_2\).
Equation \eqref{eq:v56-action-order} gives
\({\mathcal T}_2DS[W_3]=0\).  Add these identities to the measured owner.
\end{proof}

This proposition proves removability in a chosen local coordinate section.
It does not assert that the raw generator had \(q_2=0\).

\subsection{Pullback and exact nonlinear verticality}

Return to the original owner coordinates \(x\).  Let
\[
 \widehat W(x)
 :=D\Psi(x)^{-1}(W(r,w),0),
 \qquad (r,w)=\Psi(x).
 \label{eq:v56-pulled-field}
\]

\begin{theorem}[Exact vertical pullback]
\label{thm:v56-exact-pullback}
The pulled field satisfies
\begin{equation}
 D{\cal B}(x)[\widehat W(x)]=0
 \label{eq:v56-exact-verticality}
\end{equation}
at every point of the chart.  It has zero first jet at the reference
section \(r=0\).  The measured-divergence cancellation obtained in product
coordinates is coordinate invariant and hence holds for the pulled
conditional density on the original fibre.
\end{theorem}

\begin{proof}
By definition,
\[
 D\Psi[\widehat W]=(W,0).
\]
Projection onto the retained component gives
\eqref{eq:v56-exact-verticality}.  Smoothness of \(D\Psi^{-1}\) and
\(W=O(|r|^3)\) show that the pullback is also third order in the fine
displacement, hence has zero first jet.  Superdivergence was defined by the
Lie derivative of the density in V52; its bosonic restriction is therefore
coordinate invariant.  The product-coordinate identity pulls back without
an additional Jacobian term because that Jacobian is already part of the
transformed density.
\end{proof}

This is the nonlinear verticality missing from a principal-symbol
equation-of-motion argument.  The inverse chart used in
\eqref{eq:v56-pulled-field} is the same inverse controlled by the V46--V49
nested-interface chamber.

\subsection{Polymer locality of the compensation}

Decompose the unwanted jet into connected local owners,
\[
 q_{\leq M}=\sum_Xq_X,
\]
where \(q_X\) depends only on fine coordinates attached to \(X\).  Apply the
Euler homotopy on that coordinate set, producing \(W_X\).  On a complex
fine-field polydisc \(|r_i|\leq R\), use the component-sum norm.  The
denominator \(d_I+m\) cancels the sum over the \(d_I\) vector components and
gives the elementary bound
\begin{equation}
 \|{\mathcal H}_{I,m}q_m\|
 \leq R\|q_m\|.
 \label{eq:v56-local-homotopy-bound}
\end{equation}

\begin{proposition}[Weighted compensation bound]
\label{prop:v56-weighted-compensation}
Suppose the score jets, \(D\Psi^{-1}\), and their required derivatives lie
in the V48 weighted algebra at exponent \(\kappa\), and the prescribed
owner jet has finite anchored norm \(\|q_{\leq M}\|_{\rm anc,\kappa}\).
Then the finite recursion of
Theorem~\ref{thm:v56-jet-surjectivity} and its pullback have a finite
anchored norm at every \(\kappa'<\kappa\).  More precisely, for fixed \(M\),
\begin{equation}
 \|\widehat W\|_{\rm anc,\kappa'}
 \leq
 C_{M,\rho,\Psi,\kappa-\kappa'}\,R\,
 \|q_{\leq M}\|_{\rm anc,\kappa},
 \label{eq:v56-weighted-pullback}
\end{equation}
where \(C\) is polynomial in the finite score-jet and inverse-chart stocks
and contains only finitely many reserve factors
\[
 \sup_{n\geq1}n^a e^{-(\kappa-\kappa')n}.
\]
\end{proposition}

\begin{proof}
The first Euler step obeys
\eqref{eq:v56-local-homotopy-bound} support by support.  At each later
degree, the new residual is a finite sum of products of a score jet with an
earlier correction.  Weighted-algebra submultiplicativity bounds those
products.  The recursion stops at degree \(M\).  Pullback multiplies by
\(D\Psi^{-1}\) and composes with the local chart, again a finite collection
of weighted products.  Differentiating or marking a support creates only a
fixed power of its size; weakening \(\kappa\) to \(\kappa'\) absorbs each
such power by the displayed finite supremum.
\end{proof}

No total-volume factor appears.  However, uniformity as \(M\to\infty\) is
not claimed; the theorem is a finite relevant-jet compensation.

\subsection{Raw current versus chosen section}

Let \(V_{\rm raw}\) be the vector field obtained by pushing a microscopic
deformation through the exact block map.  Computing its quadratic divergence
jet requires the corresponding conditional-current derivatives.  The
existence of \(W_3\) proves that any resulting local quadratic Jacobian owner
can be removed by changing the fine coordinate section.  It does not compute
that raw coefficient and does not justify setting it to zero in the
unmodified pushforward.

\begin{firewall}[Compensation is not a physical cancellation]
\label{fire:v56-section-vs-physical}
The compensated field gives the same conditional integral only in the full
measured combination
\[
 DS[V^{\rm comp}]-\operatorname{div}_{\mu_w}V^{\rm comp}.
\]
Removing the divergence jet from one displayed row transfers higher-order
terms into \(DS[W]\) and into the higher divergence remainder.  Those terms
must enter the V53--V54 local activity and projection ledgers.
\end{firewall}

\subsection{Next acquisition}

V56 closes the finite marginal compensation step.  The remaining analytic
choice is now between computing the raw conditional-current jet and proving
that the compensated higher-order remainder contracts.  The latter is more
directly tied to the continuum gate.  The next note will derive the analytic
Taylor remainder of the measured compensation, place it in the V53
connected activity norm, and state a cofinal scale inequality under which
the transferred degree-four-and-higher owner is summable.

\begin{assessment}[Exact status after V56]
\label{ass:v56-status}
V56 proves an explicit local right inverse for divergence jets, a finite
triangular measured-divergence compensation, preservation of the prescribed
linearisation, exact nonlinear verticality after pullback, and a weighted
locality bound.  It does not calculate the raw pushed-current coefficient,
establish an all-order compensation, prove cofinal contraction of the
higher remainder, control the determinant line globally, remove regulators,
or establish the SM--Einstein continuum.
\end{assessment}

\clearpage
\section{V57: the compensation remainder in the graded gravitational tail}
\label{sec:v10v57}

\subsection{The remainder left by quadratic compensation}

Work first in one exact fine-fibre product chart of V51.  Let the action and
conditional density be analytic in the fine coordinate \(r\), and assume the
reference section is stationary:
\begin{equation}
 D_rS(0,w)=0,\qquad D_r\log\rho(0,w)=0.
 \label{eq:v57-stationary-section}
\end{equation}
Let \(q_2(r,w)\) be the unwanted homogeneous quadratic divergence jet and
choose the cubic correction from V56 with the cancelling sign,
\begin{equation}
 W_3=-{\mathcal H}_{I,2}q_2.
 \label{eq:v57-cubic-W}
\end{equation}
Then
\[
 \operatorname{div}_rW_3=-q_2
\]
and the measured owner added by this change of section is
\begin{align}
 \Delta{\cal O}
 &:=DS[W_3]-\operatorname{div}_{\mu_w}W_3\\
 &=q_2+{\cal R}_4,
 \qquad
 {\cal R}_4:=DS[W_3]-D_r\log\rho[W_3].
 \label{eq:v57-transfer-remainder}
\end{align}
The \(q_2\) term cancels the raw Jacobian marginal.  The transferred owner is
\({\cal R}_4\).

\begin{proposition}[Fourth-order onset]
\label{prop:v57-fourth-order}
Under \eqref{eq:v57-stationary-section},
\begin{equation}
 {\mathcal T}_{\leq3}{\cal R}_4=0.
 \label{eq:v57-fourth-order}
\end{equation}
If either stationarity condition is omitted, the corresponding summand can
begin at degree three but still has no quadratic term.
\end{proposition}

\begin{proof}
The cubic field satisfies \(W_3=O(|r|^3)\).  Each differential in
\eqref{eq:v57-stationary-section} is \(O(|r|)\); contracting it with \(W_3\)
gives \(O(|r|^4)\).  Without stationarity the differential has a constant
term, producing order three but not order two.
\end{proof}

This is a field-amplitude statement.  It is not yet an RG statement.

\subsection{A support--grade--analytic norm}

Let \(E_n\) be the covariant local invariant space of gravitational excess
grade \(n\) from V18--V19.  Let \(X\) be a connected interface polymer and
let \(\alpha\) be a fine-field multi-index.  Expand a local analytic owner as
\[
 F=\sum_{X,n,\alpha}F_{X,n,\alpha}r^\alpha.
\]
For support exponent \(\kappa>0\), grade radius \(R>0\), and fine analytic
radius \(\varrho>0\), define
\begin{equation}
 \|F\|_{\kappa,R,\varrho}
 :=
 \sup_b\sum_{X\ni b}e^{\kappa\tau(X)}
 \sum_{n\geq0}R^n
 \sum_\alpha\varrho^{|\alpha|}
 \|F_{X,n,\alpha}\|_{E_n}.
 \label{eq:v57-triple-norm}
\end{equation}
The basis multiplicity at each grade is included in the \(E_n\) norm, as
required by the V18 firewall.

\begin{lemma}[Analytic tail and derivative]
\label{lem:v57-analytic-bounds}
For \(0<\varrho'<\varrho\),
\begin{align}
 \|D_rF\|_{\kappa,R,\varrho'}
 &\leq
 \frac{1}{\varrho-\varrho'}
 \|F\|_{\kappa,R,\varrho},
 \label{eq:v57-Cauchy-derivative}\\
 {\mathcal T}_{<m}F=0
 \quad&\Longrightarrow\quad
 \|F\|_{\kappa,R,\theta\varrho}
 \leq\theta^m\|F\|_{\kappa,R,\varrho},
 \qquad0<\theta<1.
 \label{eq:v57-Taylor-tail}
\end{align}
\end{lemma}

\begin{proof}
For a coefficient of total degree \(k\),
\[
 k(\varrho')^{k-1}
 \leq\frac{\varrho^k}{\varrho-\varrho'},
\]
because, with \(t=\varrho'/\varrho\),
\(kt^{k-1}\leq1+t+t^2+\cdots=(1-t)^{-1}\).
Sum the coefficient majorant to get
\eqref{eq:v57-Cauchy-derivative}.  If every degree is at least \(m\), then
\((\theta\varrho)^k\leq\theta^m\varrho^k\), proving
\eqref{eq:v57-Taylor-tail}.
\end{proof}

Assume the local composition compiler has the same grade bound as V18 and a
finite connected-support convolution constant \(C_{\rm loc}\).  Then, after
the already declared finite grade-radius loss \(R\mapsto C_1R\),
\begin{equation}
 \|FG\|_{\kappa,R,\varrho}
 \leq C_{\rm loc}
 \|F\|_{\kappa,C_1R,\varrho}
 \|G\|_{\kappa,C_1R,\varrho}.
 \label{eq:v57-triple-product}
\end{equation}
Indeed, intersecting local supports have
\(\tau(X\cup Y)\leq\tau(X)+\tau(Y)\), field degrees add, and the V18 Cauchy
product controls the operator grade.

\subsection{An explicit analytic remainder stock}

Choose
\[
 0<\varrho_w<\varrho_s
\]
and put
\[
 S_*=\|S-S(0,w)\|_{\kappa,C_1R,\varrho_s},
 \qquad
 \sigma_*=\|D_r\log\rho\|_{\kappa,C_1R,\varrho_w},
\]
\[
 Q_*=\|q_2\|_{\kappa,C_1R,\varrho_w}.
\]
The local Euler homotopy bound of V56 gives
\begin{equation}
 \|W_3\|_{\kappa,C_1R,\varrho_w}
 \leq\varrho_w Q_*.
 \label{eq:v57-W-stock}
\end{equation}
Define
\begin{equation}
 B_{\rm comp}
 :=
 C_{\rm loc}\varrho_wQ_*
 \left\{
 \frac{S_*}{\varrho_s-\varrho_w}+\sigma_*
 \right\}.
 \label{eq:v57-Bcomp}
\end{equation}

\begin{theorem}[Analytic compensation-remainder bound]
\label{thm:v57-compensation-bound}
Under the preceding hypotheses,
\begin{equation}
 \|{\cal R}_4\|_{\kappa,R,\varrho_w}
 \leq B_{\rm comp}.
 \label{eq:v57-remainder-bound}
\end{equation}
Moreover, for \(0<\theta<1\),
\begin{equation}
 \|{\cal R}_4\|_{\kappa,R,\theta\varrho_w}
 \leq\theta^4B_{\rm comp}.
 \label{eq:v57-fourth-small-field}
\end{equation}
\end{theorem}

\begin{proof}
Stationarity allows \(DS[W_3]\) to be computed using
\(S-S(0,w)\).  Apply the derivative estimate
\eqref{eq:v57-Cauchy-derivative}, the product estimate
\eqref{eq:v57-triple-product}, and
\eqref{eq:v57-W-stock}.  The score term is bounded by the second product in
\eqref{eq:v57-Bcomp}.  This proves
\eqref{eq:v57-remainder-bound}.  Proposition~\ref{prop:v57-fourth-order} and
the Taylor-tail bound with \(m=4\) prove
\eqref{eq:v57-fourth-small-field}.
\end{proof}

The constant is deliberately expressed in stocks already required by the
analytic owner chart.  It is finite at a fixed regulator.  Uniformity in
scale remains a hypothesis to be verified.

\subsection{Why Taylor order cannot replace gravitational grade}

The metric coefficient \(g_{\mu\nu}\), or a dimensionless perturbation of it,
has canonical dimension zero.  Therefore multiplying a two-derivative
covariant density by additional powers of a metric coordinate does not
increase its derivative grade.  Covariantly, those powers resum into the
same owner, such as
\[
 \sqrt g,\qquad \sqrt g\,R,\qquad
 \sqrt g\,R^2,\qquad \sqrt g\,R_{\mu\nu}R^{\mu\nu}.
\]

\begin{firewall}[Fourth field order is not automatically irrelevant]
\label{fire:v57-field-vs-grade}
The factor \(\theta^4\) in
\eqref{eq:v57-fourth-small-field} is a small-field analytic gain.  It is not
a canonical RG factor for pure metric powers.  Cofinal contraction may be
claimed only after the complete covariant projection onto every
noncontracting gravitational derivative grade has been removed.
\end{firewall}

This is why the projection must act on the covariant owner before a finite
Taylor truncation in the metric chart.

\subsection{Complete low-grade projection}

Let
\[
 {\cal P}_{<N}:\mathfrak E_R\longrightarrow
 \bigoplus_{n<N}E_n
\]
be the V18 low-grade covariant projection, extended coefficientwise to the
support and analytic norm.  Decompose the full analytic remainder, not merely
its first few metric Taylor coefficients:
\begin{equation}
 c_{\rm run}:={\cal P}_{<N}{\cal R}_4,\qquad
 e_{\rm tail}:=(I-{\cal P}_{<N}){\cal R}_4.
 \label{eq:v57-grade-split}
\end{equation}

\begin{proposition}[No hidden low-grade Taylor tail]
\label{prop:v57-complete-projection}
The split \eqref{eq:v57-grade-split} is exact and obeys
\begin{equation}
 \|c_{\rm run}\|_{\kappa,R,\varrho_w}
 +\|e_{\rm tail}\|_{\kappa,R,\varrho_w}
 \leq C_PB_{\rm comp}
 \label{eq:v57-projection-bound}
\end{equation}
for the declared bounded projection constant \(C_P\).  Every analytic metric
power belonging to a grade \(n<N\) remains in \(c_{\rm run}\); none is placed
in \(e_{\rm tail}\) because its field degree is large.
\end{proposition}

\begin{proof}
Boundedness of the projection and its complement gives
\eqref{eq:v57-projection-bound}.  Since the projection acts on the covariant
grade index \(n\) before any Taylor truncation, membership in either summand
is decided by grade alone.
\end{proof}

The running term \(c_{\rm run}\) must be inserted into the V54 exact contact
and shooting ledgers.  It can be coordinate-section dependent, but it is not
discarded until an exact measured vertical cancellation is proved.

\subsection{Insertion into the existing stable-tail theorem}

Let \({\cal T}\) be the reblocked high-grade map from V18 on the radius-\(r\)
tail ball, with Lipschitz constant
\begin{equation}
 \lambda_{\rm tail}
 =L^{-N}+2C_Qr+\epsilon_N<1.
 \label{eq:v57-tail-lambda}
\end{equation}
Adding the compensation remainder changes its source by \(e_{\rm tail}\).

\begin{theorem}[Compensation-tail reserve]
\label{thm:v57-tail-reserve}
Suppose the original high-grade map sends the closed radius-\(r_0\) ball into
itself with an unused radial reserve \(r_{\rm res}>0\), and suppose
\begin{equation}
 \|e_{\rm tail}\|_{\kappa,R,\varrho_w}
 \leq(1-\lambda_{\rm tail})r_{\rm res}.
 \label{eq:v57-tail-reserve}
\end{equation}
Then the compensated high-grade map is a contraction on the enlarged
radius-\((r_0+r_{\rm res})\) ball.  Its fixed point differs from the
uncompensated fixed point by at most
\begin{equation}
 \frac{\|e_{\rm tail}\|_{\kappa,R,\varrho_w}}
      {1-\lambda_{\rm tail}}.
 \label{eq:v57-fixed-point-shift}
\end{equation}
\end{theorem}

\begin{proof}
The added source is independent of the tail variable, so it does not change
the Lipschitz constant.  Its norm is no larger than the available
self-map margin by \eqref{eq:v57-tail-reserve}.  Banach's theorem gives the
new fixed point.  If \(v\) and \(\widetilde v\) are the two fixed points,
\[
 \|\widetilde v-v\|
 \leq\lambda_{\rm tail}\|\widetilde v-v\|
      +\|e_{\rm tail}\|,
\]
which rearranges to \eqref{eq:v57-fixed-point-shift}.
\end{proof}

Combining \eqref{eq:v57-projection-bound} with
\eqref{eq:v57-tail-reserve} gives the concrete sufficient inequality
\begin{equation}
 C_PB_{\rm comp}
 \leq(1-\lambda_{\rm tail})r_{\rm res}.
 \label{eq:v57-concrete-gate}
\end{equation}
This is the cofinal compensation gate.  It invokes the V18--V19 stable-tail
hypotheses rather than reproving them.

\subsection{Scale and covariance debits}

The preceding theorem is conditional on three model-specific facts:

\begin{enumerate}
\item the fine-fibre chart, score, and action stocks in
      \(B_{\rm comp}\) are uniform under canonical rescaling;
\item the covariant projection \({\cal P}_{<N}\) is compatible with BRST,
      diffeomorphism, and interface gluing identities; and
\item the high-grade contraction constant and reserve are uniform along the
      shot trajectory.
\end{enumerate}

If the first fails, shrinking the fine analytic radius can reduce the
small-field factor but cannot repair a divergent covariant grade coefficient.
If the second fails, go upstream to the covariant owner basis.  If the third
fails, the compensation construction is not the cause; the underlying
gravitational fixed-point tail gate remains open.

\subsection{Next acquisition}

The local compensation is now compatible with the established infinite
gravitational tower.  The next missing quantity is a scale-uniform bound for
the conditional density score and its retained derivatives in
\(B_{\rm comp}\).  The next note will derive a score identity from a local
fibre partition function, express its derivatives as connected cumulants,
and state the precise cluster condition that returns a volume-free
\(\sigma_*\) in the V53 anchored norm.

\begin{assessment}[Exact status after V57]
\label{ass:v57-status}
V57 proves the fourth-order analytic compensation remainder, a combined
support--grade--analytic norm, the complete low-grade projection rule, and a
quantitative insertion into the existing stable gravitational tail theorem.
It does not prove the uniform score bound, verify the covariant projection
for every SM--Einstein owner, establish the ultraviolet tail gate, calculate
the raw current, control the determinant line globally, remove regulators,
or establish the SM--Einstein continuum.
\end{assessment}

\clearpage
\section{V58: chart scores versus interacting pushforward cumulants}
\label{sec:v10v58}

\subsection{Two scores that must be distinguished}

The compensation bound in V57 contains
\[
 \sigma_*=\|D_r\log\rho\|.
\]
Here \(\rho(r,w)\,dr\) is the conditional reference measure in the exact
fine-fibre chart.  This is a geometric chart or coarea density.  It is not
the normalized interacting Gibbs density.  By contrast, derivatives of the
retained effective action come from
\[
 Z(w)=\int e^{-S(r,w)}\rho(r,w)\,dr
\]
and are connected cumulants.  Conflating the two would charge an interacting
cluster bound for a coordinate Jacobian, or omit the cumulants needed for a
raw pushforward current.

This note derives both identities.

\subsection{Exact chart-score formula}

Let \(x=X(r,w)=\Psi^{-1}(r,w)\) be the inverse of the V51 square chart.  Let
the original bosonic reference density be
\[
 d\mu_0(x)=m(x)\,dx,\qquad m(x)>0.
\]
In product coordinates its density is
\begin{equation}
 j(r,w)=m(X(r,w))\,|\det DX(r,w)|.
 \label{eq:v58-chart-density}
\end{equation}
Normalizing \(j(r,w)\,dr\) by a factor depending only on \(w\) does not
change its fine score \(D_r\log j\).

Put
\[
 H=DX=(D\Psi(X))^{-1},
\]
and let \(H_a=H\,\iota_{r_a}\) be the inverse-chart column associated with
the fine coordinate \(r_a\).

\begin{theorem}[Fine chart-score identity]
\label{thm:v58-fine-score}
For every fine coordinate \(r_a\),
\begin{equation}
 \partial_{r_a}\log j
 =
 D_x\log m(X)[H_a]
 -
 \operatorname{tr}\left\{
 H\,D^2\Psi(X)[H_a]
 \right\}.
 \label{eq:v58-fine-score}
\end{equation}
In the graded chart, determinant and trace are replaced by Berezinian and
supertrace, with the same formula.
\end{theorem}

\begin{proof}
Differentiate the first factor in
\eqref{eq:v58-chart-density} by the chain rule.  For the second factor,
Jacobi's identity gives
\[
 \partial_{r_a}\log|\det H|
 =\operatorname{tr}(H^{-1}\partial_{r_a}H).
\]
Differentiate \(H=(D\Psi(X))^{-1}\):
\[
 \partial_{r_a}H
 =-H\,D^2\Psi(X)[H_a]\,H.
\]
Since \(H^{-1}=D\Psi(X)\), cyclicity of the trace gives the second term in
\eqref{eq:v58-fine-score}.  The Berezinian identity
\(D\log\operatorname{Ber}=\operatorname{Str}(J^{-1}DJ)\) proves the graded
version.
\end{proof}

The formula uses the same inverse chart and second chart jet that appeared
in the exact vertical projection divergence of V51.  It requires no
interacting partition function.

\subsection{Weighted locality of the chart score}

A derivative in one fine coordinate marks one local block.  The trace in
\eqref{eq:v58-fine-score} is therefore a marked trace rather than an
unmarked extensive trace.

\begin{lemma}[Marked weighted trace]
\label{lem:v58-marked-trace}
Let \(A\) and \(B[h]\) be block kernels on the interface graph, with
\(h\) supported at one marked block.  Suppose
\[
 \|A\|_\kappa\leq a,\qquad
 \|B[h]\|_\kappa\leq b\|h\|.
\]
If the local matrix fibre dimension is at most \(\nu\), then the coefficient
of the marked trace satisfies an anchored bound
\begin{equation}
 \|\operatorname{tr}(AB[h])\|_{\rm anc,\kappa}
 \leq \nu ab\|h\|.
 \label{eq:v58-marked-trace-bound}
\end{equation}
The bound is independent of total volume.
\end{lemma}

\begin{proof}
Expand the trace as
\[
 \operatorname{tr}(AB[h])
 =\sum_{x,y}\operatorname{tr}_{x}(A_{xy}B[h]_{yx}).
\]
The mark in \(h\) anchors the support.  Weight the two kernels by
\(e^{\kappa d(x,y)}\) and use the row/column Schur sums in the V48 algebra.
The finite local trace costs at most \(\nu\).  No free spatial sum remains.
\end{proof}

Let
\[
 h_*=\|H\|_\kappa,\qquad
 b_*=\|D^2\Psi\|_\kappa,\qquad
 m_*=\|D\log m\|_\kappa.
\]

\begin{corollary}[Volume-free fine-score stock]
\label{cor:v58-fine-score-stock}
The fine chart score obeys
\begin{equation}
 \|D_r\log j\|_{\rm anc,\kappa}
 \leq
 C_{\rm col}\,m_*h_*
 +C_{\rm tr}\,\nu h_*^2b_*,
 \label{eq:v58-fine-score-stock}
\end{equation}
where \(C_{\rm col},C_{\rm tr}\) are the fixed local column and support
conventions.  This supplies the \(\sigma_*\) input in V57 whenever the three
stocks on the right are uniform.
\end{corollary}

\begin{proof}
Bound the first term of \eqref{eq:v58-fine-score} by one weighted product.
For the second, \(H_a\) costs one inverse-chart column, and
Lemma~\ref{lem:v58-marked-trace} costs the other \(H\) and the second jet.
\end{proof}

Higher retained or fine derivatives of the chart score contain only a finite
sum of products of \(H\) and jets \(D^k\Psi,D^k\log m\).  They follow by
repeatedly using
\[
 DH[h]=-H\,D^2\Psi[Hh]\,H.
\]
Thus a \(k\)-th score jet uses chart jets through order \(k+2\); it is not
obtained from the first two jets alone.

\subsection{The interacting retained score}

Define
\[
 L(r,w):=-S(r,w)+\log j(r,w),
 \qquad
 Z(w)=\int e^{L(r,w)}\,dr.
\]
For a retained direction \(h\), put
\[
 L_h=D_wL[h].
\]
Expectations and cumulants below are with respect to the normalized
conditional density \(Z(w)^{-1}e^Ldr\), whenever the body of \(Z\) is
nonzero.

\begin{proposition}[First and second retained score identities]
\label{prop:v58-first-second-score}
For retained directions \(h_1,h_2\),
\begin{align}
 D\log Z[h_1]
 &=\langle L_{h_1}\rangle,
 \label{eq:v58-first-score}\\
 D^2\log Z[h_1,h_2]
 &=\langle L_{h_1h_2}\rangle
   +\langle L_{h_1};L_{h_2}\rangle_c.
 \label{eq:v58-second-score}
\end{align}
\end{proposition}

\begin{proof}
Differentiate under the finite regulated integral:
\[
 DZ[h]=\int L_he^Ldr.
\]
Division by \(Z\) gives \eqref{eq:v58-first-score}.  Differentiate the
normalized expectation.  The derivative of the observable gives
\(\langle L_{h_1h_2}\rangle\), while differentiating the density subtracts
the product of means and gives the connected covariance in
\eqref{eq:v58-second-score}.
\end{proof}

For a nonempty set \(B\subset\{1,\ldots,n\}\), write
\[
 L_B=D_w^{|B|}L[h_i:i\in B].
\]
Let \(\kappa_m(A_1,\ldots,A_m)\) denote the joint connected cumulant.

\begin{theorem}[All-order retained score formula]
\label{thm:v58-all-score}
At every finite regulator,
\begin{equation}
 D_w^n\log Z[h_1,\ldots,h_n]
 =
 \sum_{\pi\in{\mathfrak P}_n}
 \kappa_{|\pi|}(L_B:B\in\pi),
 \label{eq:v58-partition-cumulant}
\end{equation}
where \({\mathfrak P}_n\) is the set of set partitions of
\(\{1,\ldots,n\}\).
\end{theorem}

\begin{proof}
Introduce commuting parameters \(t_1,\ldots,t_n\) and write
\[
 L(t)=L(r,w+t_1h_1+\cdots+t_nh_n).
\]
The logarithm of
\(\int e^{L(t)}dr\) is the cumulant-generating functional for the collection
of Taylor coefficients of \(L(t)-L(0)\).  Extracting the coefficient of
\(t_1\cdots t_n\) groups the indices first into the derivative blocks \(B\)
acting on one copy of \(L\), and then takes the connected cumulant of those
copies.  These groupings are exactly the set partitions \(\pi\), with
coefficient one because all directions are distinctly labelled.  This gives
\eqref{eq:v58-partition-cumulant}.  The cases \(n=1,2\) are
\eqref{eq:v58-first-score}--\eqref{eq:v58-second-score}.
\end{proof}

This theorem is the correct source of the raw pushed-current derivatives.
A third derivative of a conditional current can require a fourth cumulant
because differentiating the current adds one marked observable.

\subsection{A finite tree-cluster gate}

Assume each marked observable decomposes into connected local pieces,
\[
 L_B=\sum_XL_{B,X}.
\]
For \(m\leq n\), suppose the interacting truncated cumulants obey the tree
bound
\begin{equation}
 \left\|
 \kappa_m(L_{B_1,X_1},\ldots,L_{B_m,X_m})
 \right\|
 \leq
 C_c^m
 \left(\prod_{i=1}^m a_{B_i}(X_i)\right)
 \sum_{T\in{\mathcal T}_m}
 \prod_{\{i,j\}\in T}e^{-\mu d(X_i,X_j)}.
 \label{eq:v58-tree-cumulant}
\end{equation}
Let
\[
 v_\mu
 :=\sup_X\sum_Y e^{-\mu d(X,Y)}a(Y),
 \qquad
 a(Y):=\max_{1\leq |B|\leq n}a_B(Y),
\]
and let
\[
 z:=\sup_x\sum_{X\ni x}e^{\kappa\tau(X)}a(X).
\]

\begin{theorem}[Finite-order volume-free score gate]
\label{thm:v58-finite-cluster-gate}
If \(z<\infty\), \(v_\mu<\infty\), and
\eqref{eq:v58-tree-cumulant} holds through cumulant order \(n\), then every
retained derivative in
\eqref{eq:v58-partition-cumulant} through order \(n\) has a finite anchored
polymer norm independent of volume.  One explicit bound is
\begin{equation}
 \|D_w^n\log Z\|_{\rm anc,\kappa'}
 \leq
 \sum_{\pi\in{\mathfrak P}_n}
 C_c^{|\pi|}
 |\mathcal T_{|\pi|}|\,
 z\,v_\mu^{|\pi|-1}\,
 C_{\rm geom}(\kappa-\kappa',n),
 \label{eq:v58-score-cluster-bound}
\end{equation}
for every \(\kappa'<\kappa\), where
\(|\mathcal T_m|=m^{m-2}\) for \(m\geq2\), with the usual value one for
\(m=1\).  The finite factor \(C_{\rm geom}\) spends the support-size reserve
needed to thicken the connecting tree.
\end{theorem}

\begin{proof}
Fix one polymer piece containing the anchored block.  For each labelled tree
in \eqref{eq:v58-tree-cumulant}, sum the remaining supports by removing
leaves successively.  Every removed leaf costs at most \(v_\mu\), leaving
the anchored activity \(z\).  Cayley's formula counts the labelled trees.
Thickening their connector paths creates only a fixed support-size power for
fixed \(n\), absorbed by weakening \(\kappa\) to \(\kappa'\).  Sum the
finitely many partitions in \eqref{eq:v58-partition-cumulant}.
\end{proof}

No small-activity inequality is needed for one fixed derivative order once
the tree cumulant bound is assumed.  An all-order analytic or Gevrey score
requires additional control of the Bell and Cayley growth; that is the
stronger forest gate developed earlier.

\subsection{What is now proved and what remains an input}

Corollary~\ref{cor:v58-fine-score-stock} reduces the V57 compensation score
to local chart jets.  Theorem~\ref{thm:v58-finite-cluster-gate} reduces the
raw retained current jets to a declared interacting tree-cumulant bound.
These are separate acquisition routes.

For the current programme, the chart-score route can proceed inside the
V49 inverse chamber.  The raw pushforward route still needs the
four-marked-observable cluster estimate at the order identified by the YM
audit.  Neither the algebraic cumulant identity nor the finite tree
summation proves that interacting estimate.

\begin{firewall}[A normalized score identity is not a cluster theorem]
\label{fire:v58-score-not-cluster}
Equations \eqref{eq:v58-first-score}--\eqref{eq:v58-partition-cumulant} are
exact at finite volume even in a phase with long-range correlations.  They
do not imply the tree decay
\eqref{eq:v58-tree-cumulant}.  Uniform clustering must come from the
small-activity/large-field polymer construction or another independent
mixing theorem.
\end{firewall}

\subsection{Next acquisition}

The most direct remaining step is to populate the geometric score stock in
\eqref{eq:v58-fine-score-stock} for the actual nested DeTurck--CMC chart.
The next note will differentiate the V44 block inverse to bound \(H\) and
\(D^2\Psi\) in the weighted algebra, producing an explicit score polynomial
in the block inverse, first block jets, and second block jets.  This avoids
waiting for the harder interacting four-cumulant gate when only the
coordinate compensation is required.

\begin{assessment}[Exact status after V58]
\label{ass:v58-status}
V58 proves the exact chart-score formula, a volume-free marked-trace bound,
the all-order interacting retained-score identity, and a finite-order
tree-cluster implication.  It does not prove the interacting cumulant decay,
populate the actual DeTurck--CMC chart-jet constants, calculate the raw
pushforward coefficient, establish the ultraviolet gravitational tail gate,
control the determinant line globally, remove regulators, or establish the
SM--Einstein continuum.
\end{assessment}

\clearpage
\section{V59: explicit nested-block stocks for the chart score}
\label{sec:v10v59}

\subsection{The square derivative and notation}

Order the fine and retained coordinates so that the derivative of the exact
square chart has the V44 form
\begin{equation}
 {\mathcal J}:=D\Psi=
 \begin{pmatrix}
 A&G&H\\
 0&D&K\\
 P&Q&R
 \end{pmatrix}.
 \label{eq:v59-chart-J}
\end{equation}
The letters \(H\) and \({\mathcal J}^{-1}\) must not be confused; in this
note write
\[
 X=A^{-1},\qquad Y=D^{-1},\qquad
 T=Q-PXG,
\]
\[
 S=R-PXH-TYK,\qquad Z=S^{-1}.
\]
All norms below are the weighted algebra norm at an exponent already
weakened enough to pay the finite support marks.

\subsection{Exact inverse response}

For target increments \((y,c,m)\), solve
\[
 {\mathcal J}
 \begin{pmatrix}u\\z\\b\end{pmatrix}
 =
 \begin{pmatrix}y\\c\\m\end{pmatrix}.
\]

\begin{proposition}[Nested inverse formula]
\label{prop:v59-nested-inverse}
The solution is
\begin{align}
 b&=Z\{m-PXy-TYc\},
 \label{eq:v59-b-solution}\\
 z&=Y(c-Kb),
 \label{eq:v59-z-solution}\\
 u&=X(y-Gz-Hb).
 \label{eq:v59-u-solution}
\end{align}
Hence no inverse other than \(X,Y,Z\) occurs in the full chart inverse.
\end{proposition}

\begin{proof}
The first two block rows give
\[
 u=X(y-Gz-Hb),\qquad z=Y(c-Kb).
\]
Insert both expressions into the third row.  The coefficient of \(c\) is
\((Q-PXG)Y=TY\), and the coefficient of \(b\) is
\[
 R-PXH-(Q-PXG)YK=S.
\]
Thus
\[
 m=PXy+TYc+Sb,
\]
which gives \eqref{eq:v59-b-solution}; back substitution gives the other two
formulae.
\end{proof}

Put
\[
 x=\|X\|,\quad y_0=\|Y\|,\quad z_0=\|Z\|,
 \quad p=\|P\|,\quad g=\|G\|,\quad h=\|H\|,
 \quad k=\|K\|,\quad t=\|T\|.
\]
Define
\begin{align}
 \beta_b&:=z_0(1+px+ty_0),
 \label{eq:v59-beta-b}\\
 \beta_z&:=y_0(1+k\beta_b),
 \label{eq:v59-beta-z}\\
 \beta_u&:=x(1+g\beta_z+h\beta_b),
 \label{eq:v59-beta-u}\\
 \beta&:=\beta_u+\beta_z+\beta_b.
 \label{eq:v59-beta}
\end{align}

\begin{corollary}[Full inverse-chart stock]
\label{cor:v59-full-inverse-stock}
With sum norms on the three domain and target factors,
\begin{equation}
 \|{\mathcal J}^{-1}\|\leq\beta.
 \label{eq:v59-full-inverse}
\end{equation}
If the V46 interface certificate gives
\(\|S^{-1}\|\leq\gamma_*^{-1}\), then one may set
\(z_0=\gamma_*^{-1}\) in
\eqref{eq:v59-beta-b}--\eqref{eq:v59-beta}.
\end{corollary}

\begin{proof}
Apply the triangle and product inequalities successively to
\eqref{eq:v59-b-solution}--\eqref{eq:v59-u-solution}.  The response of each
component to a unit sum-norm target is bounded by the corresponding
\(\beta\) row.  Sum the three component bounds.
\end{proof}

The polynomial \(\beta\) makes the loss transparent.  A small interface
floor enters every response through \(\beta_b\), then propagates to the CMC
and bulk rows.  It cannot be hidden inside the DeTurck inverse constant.

\subsection{Second and third chart-jet stocks}

Define the multilinear weighted stocks
\begin{equation}
 j_2:=\|D{\mathcal J}\|=\|D^2\Psi\|,
 \qquad
 j_3:=\|D^2{\mathcal J}\|=\|D^3\Psi\|.
 \label{eq:v59-jets}
\end{equation}
They are finite sums of derivatives of the eight blocks in
\eqref{eq:v59-chart-J}.  More explicitly, if
\[
 \delta_1A,\delta_1G,\ldots,\delta_1R
\]
are uniform first block-jet stocks in a unit owner direction, and similarly
\(\delta_2A,\ldots,\delta_2R\) for second block jets, then for the chosen
sum norms
\begin{equation}
 j_2\leq C_J\sum_{B\in{\mathscr B}}\delta_1B,
 \qquad
 j_3\leq C_J\sum_{B\in{\mathscr B}}\delta_2B,
 \quad
 {\mathscr B}=\{A,G,H,D,K,P,Q,R\},
 \label{eq:v59-block-jet-sums}
\end{equation}
where \(C_J\) is the fixed block-norm convention.

\begin{proposition}[Inverse-chart derivative]
\label{prop:v59-inverse-derivative}
For a target direction \(\xi\),
\begin{equation}
 D({\mathcal J}^{-1})[\xi]
 =
 -{\mathcal J}^{-1}
 \{D{\mathcal J}[{\mathcal J}^{-1}\xi]\}
 {\mathcal J}^{-1},
 \label{eq:v59-inverse-derivative}
\end{equation}
and hence
\begin{equation}
 \|D({\mathcal J}^{-1})[\xi]\|
 \leq\beta^3j_2\|\xi\|.
 \label{eq:v59-inverse-derivative-bound}
\end{equation}
\end{proposition}

\begin{proof}
The target direction changes the base point by
\({\mathcal J}^{-1}\xi\).  Differentiate
\({\mathcal J}{\mathcal J}^{-1}=I\) in that base-point direction and solve
for the inverse derivative, obtaining
\eqref{eq:v59-inverse-derivative}.  Three inverse factors are counted:
one in the base-point displacement and two in the inverse derivative.
Submultiplicativity gives
\eqref{eq:v59-inverse-derivative-bound}.
\end{proof}

This is the exact response hidden in a notation that differentiates the
inverse with respect to target coordinates as though source and target
directions were the same.

\subsection{The score stock}

Let the original reference density be \(m(x)\,dx\), and put
\[
 m_1=\|D\log m\|,\qquad m_2=\|D^2\log m\|.
\]
Let \(\nu\) be the finite local trace dimension.  The V58 formula is
\[
 s_a
 =D\log m[{\mathcal J}^{-1}\iota_{r_a}]
 -\operatorname{tr}\{
 {\mathcal J}^{-1}D{\mathcal J}
   [{\mathcal J}^{-1}\iota_{r_a}]
 \}.
\]

\begin{theorem}[Explicit fine-score polynomial]
\label{thm:v59-score-polynomial}
The anchored fine-score norm obeys
\begin{equation}
 \|D_r\log j\|_{\rm anc}
 \leq
 \Sigma_0,
 \qquad
 \Sigma_0:=C_{\rm mark}
 \{m_1\beta+\nu\beta^2j_2\}.
 \label{eq:v59-Sigma-zero}
\end{equation}
The constant \(C_{\rm mark}\) depends only on the fixed local column and
support conventions and is independent of total volume.
\end{theorem}

\begin{proof}
A fine coordinate column of \({\mathcal J}^{-1}\) has norm at most
\(\beta\).  The reference-density term costs \(m_1\beta\).  The marked trace
contains two inverse columns/factors and one first derivative of
\({\mathcal J}\), costing \(\nu\beta^2j_2\) by the V58 marked-trace lemma.
The mark anchors every spatial sum.
\end{proof}

Combining \eqref{eq:v59-Sigma-zero} with V57 gives the explicit compensation
stock
\begin{equation}
 B_{\rm comp}
 \leq
 C_{\rm loc}\varrho_wQ_*
 \left\{
 \frac{S_*}{\varrho_s-\varrho_w}
 +\Sigma_0
 \right\}.
 \label{eq:v59-explicit-Bcomp}
\end{equation}
Thus the score denominator in the cofinal gate is now expressed through
\(X,Y,S^{-1}\), coupling blocks, and first block jets.

\subsection{One retained derivative of the score}

A raw conditional-current calculation needs retained derivatives.  Let
\(\xi\) be a unit retained target direction.  Differentiating the score uses
\(m_2,j_3\) and the inverse derivative of
Proposition~\ref{prop:v59-inverse-derivative}.

\begin{theorem}[First retained score-jet stock]
\label{thm:v59-score-one}
For a fixed unit fine mark and a unit retained direction,
\begin{equation}
 \|D_wD_r\log j\|_{\rm anc}
 \leq\Sigma_1,
 \label{eq:v59-Sigma-one-bound}
\end{equation}
where one valid polynomial stock is
\begin{equation}
 \Sigma_1
 :=
 C_{\rm mark}\left\{
 m_2\beta^2+m_1\beta^3j_2
 +\nu\left(2\beta^4j_2^2+\beta^3j_3\right)
 \right\}.
 \label{eq:v59-Sigma-one}
\end{equation}
\end{theorem}

\begin{proof}
The retained direction moves the base point by a vector of norm at most
\(\beta\).  Differentiating the reference-density term produces
\(m_2\beta^2\) and a differentiated inverse column, bounded by
\(m_1\beta^3j_2\).

For the trace term, differentiating either inverse occurrence gives two
terms.  Each costs at most
\(\nu\beta^4j_2^2\): the differentiated inverse costs
\(\beta^3j_2\), the remaining chart jet costs \(j_2\), and the remaining
column/factor accounts for the displayed power after the base-point
response is included as in
\eqref{eq:v59-inverse-derivative-bound}.  Differentiating
\(D{\mathcal J}=D^2\Psi\) gives \(D^2{\mathcal J}=D^3\Psi\) and costs
\(\nu\beta^3j_3\).  The marked weighted trace bound removes total-volume
dependence.  Summing the four families gives
\eqref{eq:v59-Sigma-one}.
\end{proof}

The polynomial is not asserted sharp.  Its value is that every inverse
power and every required chart jet is visible.

\subsection{Higher score jets and the derivative count}

Repeated differentiation gives a finite sum over rooted trees.  Every leaf
is a target direction, every internal chart vertex is a jet
\(D^k\Psi\), and every inverse edge is \({\mathcal J}^{-1}\).  A \(k\)-th
retained derivative of the fine score requires chart jets through order
\(k+2\) and reference-density jets through order \(k+1\).

\begin{proposition}[No hidden lower jet closure]
\label{prop:v59-jet-count}
For \(k\geq0\), \(D_w^kD_r\log j\) is determined by
\[
 {\mathcal J}^{-1},\quad
 D^2\Psi,\ldots,D^{k+2}\Psi,\quad
 D\log m,\ldots,D^{k+1}\log m.
\]
In general it is not determined by a smaller jet list.
\end{proposition}

\begin{proof}
The assertion follows inductively from
\[
 D({\mathcal J}^{-1})
 =-{\mathcal J}^{-1}(D{\mathcal J}){\mathcal J}^{-1}
\]
with the base-point response inserted.  Differentiating the highest existing
chart jet raises its order by one, so \(D^{k+2}\Psi\) appears.  A scalar
example in which only that highest derivative is changed while all lower
jets agree changes the resulting score derivative, proving the final
statement.
\end{proof}

This derivative count agrees with the YM audit's warning that a marginal
Jacobian Hessian can spend one more current jet than a naive quadratic
calculation suggests.

\subsection{Local acquisition card}

At a fixed regulator, the remaining chart-score acquisition is finite:

\begin{enumerate}
\item certify \(x=\|A^{-1}\|\), \(y_0=\|D^{-1}\|\), and the V46 interface
      floor \(\gamma_*\);
\item populate the coupling stocks \(p,g,h,k,t\);
\item compute the first and second local block jets in
      \eqref{eq:v59-block-jet-sums};
\item insert them into \(\beta,\Sigma_0,\Sigma_1\); and
\item test \eqref{eq:v59-explicit-Bcomp} against the V57 tail reserve.
\end{enumerate}

All of these are weighted, volume-free quantities on the local sector.
Harmonic, Killing, constant metric, and chiral zero modes remain in the
macroscopic block and do not enter \(\beta\).

\begin{firewall}[A polynomial stock is not a verified numerical chamber]
\label{fire:v59-stock-not-number}
Equations \eqref{eq:v59-Sigma-zero} and
\eqref{eq:v59-Sigma-one} prove how certified block constants combine.  They
do not supply numerical values for those constants or prove that the V57
reserve inequality holds.  Those facts require the actual regulated
stencil and background chamber.
\end{firewall}

\subsection{Next acquisition}

V60 is the next compulsory YM/NS audit.  It will check whether either
companion has produced a reusable four-marked cluster gate or a sharper
local current-jet compiler.  The post-audit note will then choose between
populating one explicit flat-collar block-jet row and formulating the
four-cumulant small-activity inequality needed for the raw pushforward.

\begin{assessment}[Exact status after V59]
\label{ass:v59-status}
V59 proves the exact nested inverse response, an explicit full
inverse-chart norm, and volume-free zeroth and first retained chart-score
stocks.  It does not populate the block-jet constants numerically, prove the
four-cumulant interacting cluster gate, verify the V57 reserve along a full
trajectory, control macroscopic modes or the determinant line, remove
regulators, or establish the SM--Einstein continuum.
\end{assessment}

\clearpage
\section{V60: compulsory audit and the aligned-selector decision}
\label{sec:v10v60}

\subsection{Companion state}

The mandatory audit inspected
\[
 \texttt{Renewal\_Yang\_Mills\_Mass\_Gap\_Research\_Notes\_V64.tex},
 \qquad
 \texttt{SHA--256 }714F28870215CA2E78A8113070166BF448C9DEB712B64CF96E477964127905B4,
\]
and the unchanged
\[
 \texttt{Renewal\_NS\_Stability\_Research\_Notes\_V26.tex},
 \qquad
 \texttt{SHA--256 }82C887F8C2C69E4F28FAD7C3DC8DE5B9099CB5A4BF431EBA1FFF3E91935978AD.
\]

YM V64 replaces a fixed weak-field radius by a running window, proves
simultaneous decay of an explicit good-field nonlinear remainder and an
explicit bad-set primitive, and uses a moving Haar coordinate so retained
background derivatives do not hit the selector.  Its remaining debit is a
source-marked Gaussian forest extraction.  The NS note supplies no new
operator or clustering input.

\subsection{What transfers and what does not}

\begin{assessment}[Transferable coordinate mechanism]
\label{ass:v60-coordinate-transfer}
The moving-coordinate selector convention transfers to the exact
DeTurck--CMC graph chart.  If fine deviations from the moving stable solution
are used as integration variables, a selector depending only on those
deviations has zero retained derivative.  All retained scores then occur in
the transported action, chart density, block map, and current.  The
horizontal current and its Jacobian remain mandatory.
\end{assessment}

\begin{proof}
This is an ordinary parameter-dependent change of variables.  Holding the
new fine coordinate fixed annihilates every selector that is a fixed
function of that coordinate.  The retained dependence moves into the inverse
chart and its determinant.  V51's projectable horizontal transport identity
shows that this determinant and current are the terms required to make the
two coordinate descriptions equal.
\end{proof}

The explicit YM window
\[
 r_\beta=\beta^{-2/5}
\]
and its \(O(\beta^{-1/5})\) remainder use compact Wilson coercivity and
asymptotic weak coupling.  No corresponding gravitational bare parameter or
non-Gaussian fixed-point estimate has been proved in this programme.

\begin{firewall}[No import of Yang--Mills weak-coupling powers]
\label{fire:v60-no-beta-import}
The powers and numerical constants in YM V64 cannot be used as an Einstein
small-field chamber.  Gravity still requires the V18 ultraviolet tail gate
or another proved fixed-point mechanism.  Only the exact moving-coordinate
identity is imported.
\end{firewall}

\subsection{Relation to the V58--V59 scores}

V58 distinguished the fine chart score from the interacting retained score.
The audit sharpens that distinction.  A moving selector removes an
artificial retained derivative of the cutoff, but it does not make
\[
 D_r\log j
\]
or
\[
 D_w\log Z
\]
vanish.  The former remains the geometric score bounded in V59; the latter
contains the action and chart observables in the connected cumulant formula.

\begin{assessment}[NS check]
\label{ass:v60-ns-check}
The unchanged rank-one Oseen calculation does not justify restricting
DeTurck, CMC, or interface deviations.  The selector must be defined on the
full proved local fine-coordinate domain.
\end{assessment}

\subsection{Post-audit acquisition}

Let the stable graph be
\[
 x=X(r,w),
 \qquad X(0,w)=x_*(w),
\]
where \(r\) contains the aligned deviations of V44 and \(w\) the retained
owners.  Choose a smooth fixed profile \(\chi\) and a scale radius \(R_j\),
and define
\[
 \chi_j(r)=\chi(\|r\|/R_j).
\]

\begin{decision}[V61 target]
\label{dec:v60-v61}
V61 will prove:
\begin{enumerate}
\item exact good/bad decomposition in aligned graph coordinates;
\item vanishing of all retained derivatives of \(\chi_j\) at fixed \(r\);
\item equality with the fixed-original-coordinate formula, including the
      horizontal transport current and chart Jacobian;
\item all-order retained derivative formulae with no artificial
      \(R_j^{-1}\) selector factors; and
\item the precise bad-set suppression input still needed for the
      SM--Einstein chart.
\end{enumerate}
\end{decision}

This precedes the four-cumulant acquisition.  Removing spurious selector
marks reduces the observable bank that the eventual cluster theorem must
control.

\subsection{Audit boundary}

\begin{assessment}[Exact status after V60]
\label{ass:v60-status}
V60 records the current companion state and imports only the justified
moving-coordinate mechanism.  It does not prove an Einstein running
small-field window, a bad-set suppression estimate, the four-marked cluster
gate, the ultraviolet gravitational tail gate, determinant-line control,
regulator removal, or the SM--Einstein continuum.
\end{assessment}

\clearpage
\section{V61: aligned selectors, horizontal transport, and the bad-set input}
\label{sec:v10v61}

\subsection{Aligned good and bad coordinates}

Let
\[
 x=X(r,w)=\Psi^{-1}(r,w)
\]
be the exact product chart, with
\[
 X(0,w)=x_*(w)
\]
the moving stable DeTurck--CMC solution.  The fine coordinate \(r\) consists
of the aligned deviations from that graph, while \(w\) is retained.  Let
\(\chi:[0,\infty)\to[0,1]\) be smooth, equal to one on \([0,1]\), and equal
to zero on \([2,\infty)\).  At scale \(j\), choose a positive number \(R_j\)
that depends on scale but not on the retained field and put
\begin{equation}
 \chi_j(r)=\chi(\|r\|/R_j).
 \label{eq:v61-selector}
\end{equation}

Let
\[
 j(r,w)=m(X(r,w))|\det D_{(r,w)}X|
\]
be the chart density and set
\[
 \Lambda(r,w)=-S(X(r,w),w)+\log j(r,w).
\]
Define
\begin{align}
 Z_{{\rm g},j}(w)
 &=\int\chi_j(r)e^{\Lambda(r,w)}\,dr,
 \label{eq:v61-good}\\
 Z_{{\rm b},j}(w)
 &=\int(1-\chi_j(r))e^{\Lambda(r,w)}\,dr.
 \label{eq:v61-bad}
\end{align}

\begin{proposition}[Exact good--bad decomposition]
\label{prop:v61-good-bad}
At every finite regulator,
\begin{equation}
 Z(w)=Z_{{\rm g},j}(w)+Z_{{\rm b},j}(w).
 \label{eq:v61-good-bad}
\end{equation}
No overlap correction is present.
\end{proposition}

\begin{proof}
The identity is the integral of
\(\chi_j+(1-\chi_j)=1\).  Smooth overlap is already encoded in the two
complementary weights.
\end{proof}

\subsection{Retained derivatives in moving coordinates}

For a retained direction \(h\), let
\[
 \Lambda_h=D_w\Lambda[h]
\]
with \(r\) held fixed.

\begin{theorem}[Stationary selector convention]
\label{thm:v61-stationary-selector}
For every retained direction and every \(n\geq1\),
\begin{equation}
 D_w^n\chi_j(r)=0
 \quad\hbox{at fixed }r.
 \label{eq:v61-no-selector-derivative}
\end{equation}
Consequently,
\begin{equation}
 D_wZ_{{\rm g},j}[h]
 =\int\chi_j e^\Lambda\Lambda_h\,dr,
 \label{eq:v61-first-good-derivative}
\end{equation}
and, for labelled directions \(h_1,\ldots,h_n\),
\begin{equation}
 D_w^nZ_{{\rm g},j}[h_1,\ldots,h_n]
 =
 \int\chi_je^\Lambda
 \sum_{\pi\in{\mathfrak P}_n}
 \prod_{B\in\pi}\Lambda_B\,dr,
 \label{eq:v61-Bell-good}
\end{equation}
where
\(\Lambda_B=D_w^{|B|}\Lambda[h_i:i\in B]\).
No factor \(R_j^{-1}\) is generated by a retained derivative.
\end{theorem}

\begin{proof}
The selector in \eqref{eq:v61-selector} has no \(w\)-dependence when \(r\)
is the independent variable, proving
\eqref{eq:v61-no-selector-derivative}.  Differentiate
\eqref{eq:v61-good} under the finite regulated integral.  The first
derivative gives \eqref{eq:v61-first-good-derivative}; repeated product
differentiation of \(e^\Lambda\) is indexed by set partitions and gives
\eqref{eq:v61-Bell-good}.
\end{proof}

If the body of \(Z_{{\rm g},j}\) is nonzero, its logarithmic derivatives are
the cumulant formula of V58 with the good conditional measure.  The selector
changes that measure, but it is not a marked retained observable.

\subsection{Equality with the fixed-coordinate formula}

The absence of a selector derivative does not omit a current.  It moves that
current into the inverse chart and its Jacobian.

Let
\[
 {\cal R}(x,w)=r
\]
be the fine component of \(\Psi(x,w)\).  In original coordinates,
\begin{equation}
 Z_{{\rm g},j}(w)
 =
 \int
 \chi_j({\cal R}(x,w))e^{-S(x,w)}m(x)\,dx.
 \label{eq:v61-fixed-x-good}
\end{equation}
Define the horizontal lift at fixed \(r\) by
\begin{equation}
 V_h(x,w)
 :=
 D_wX({\cal R}(x,w),w)[h].
 \label{eq:v61-horizontal-current}
\end{equation}

\begin{lemma}[Selector transport]
\label{lem:v61-selector-transport}
At fixed \(x\),
\begin{equation}
 D_w\{\chi_j({\cal R}(x,w))\}[h]
 =
 -D_x\{\chi_j({\cal R}(x,w))\}[V_h].
 \label{eq:v61-selector-transport}
\end{equation}
\end{lemma}

\begin{proof}
Differentiate
\[
 X({\cal R}(x,w),w)=x
\]
at fixed \(x\).  The result says that the change of the fine coordinate
cancels the horizontal motion of \(X\).  Apply the chain rule to
\(\chi_j\).
\end{proof}

Let
\(\operatorname{div}_{m}V=m^{-1}\operatorname{div}(mV)\).

\begin{lemma}[Chart Jacobian transport]
\label{lem:v61-Jacobian-transport}
At fixed \(r\),
\begin{equation}
 D_w\log j(r,w)[h]
 =
 \operatorname{div}_{m}V_h(X(r,w),w).
 \label{eq:v61-Jacobian-current}
\end{equation}
\end{lemma}

\begin{proof}
The parameter-dependent diffeomorphism
\(r\mapsto X(r,w)\) transports the density \(m(x)\,dx\).  Its logarithmic
Lie derivative along the velocity \(V_h\) is
\(\operatorname{div}_{m}V_h\).  Equivalently, differentiate
\(m(X)\det DX\), use the chain rule for \(m\), and Jacobi's determinant
identity for \(DX\).
\end{proof}

\begin{theorem}[Exact cutoff-current cancellation]
\label{thm:v61-cutoff-current}
The derivative of \eqref{eq:v61-fixed-x-good} at fixed \(x\) equals
\eqref{eq:v61-first-good-derivative}.  More explicitly,
\begin{align}
 &\int \chi_j e^{-S}
 \{-D_wS|_x-D_xS[V_h]+\operatorname{div}_mV_h\}\,d\mu_0
 \label{eq:v61-moving-derivative}\\
 &\qquad =
 \int e^{-S}
 \{D_w\chi_j({\cal R})-\chi_jD_wS|_x\}\,d\mu_0.
 \label{eq:v61-fixed-derivative}
\end{align}
\end{theorem}

\begin{proof}
The total action derivative at fixed \(r\) is
\[
 D_wS|_r=D_wS|_x+D_xS[V_h].
\]
Combine this with
Lemma~\ref{lem:v61-Jacobian-transport} to obtain the integrand on the first
line.  The measured divergence identity gives
\[
 \operatorname{div}_{m}(e^{-S}V_h)
 =e^{-S}\{\operatorname{div}_mV_h-D_xS[V_h]\}.
\]
Multiply by \(\chi_j\) and integrate by parts.  The compact support of the
smooth selector removes the boundary flux and yields
\[
 \int\chi_je^{-S}
 \{\operatorname{div}_mV_h-D_xS[V_h]\}\,d\mu_0
 =
 -\int D_x\chi_j[V_h]e^{-S}\,d\mu_0.
\]
Use Lemma~\ref{lem:v61-selector-transport}.  The remaining explicit action
derivative is identical on both sides.
\end{proof}

Thus the apparent \(R_j^{-1}\) cutoff derivative in fixed coordinates is
exactly paired with the horizontal current and chart Jacobian.  Holding both
\(x\) and \(r\) fixed would count neither coordinate system correctly.

\begin{firewall}[The horizontal current is not vertical redundancy]
\label{fire:v61-horizontal-not-vertical}
The field \(V_h\) satisfies
\[
 D{\cal B}[V_h]=h.
\]
It changes the retained owner and is projectable rather than vertical.
Its action and Jacobian terms produce the retained transport identity of
V51; they cannot be deleted by the vertical measured cancellation.
\end{firewall}

\subsection{Which derivatives can still cost the radius}

Theorem~\ref{thm:v61-stationary-selector} concerns retained derivatives.
A fine-coordinate derivative of \(\chi_j(r)\) does cost \(R_j^{-1}\), and a
derivative with respect to the scale parameter can differentiate \(R_j\).
Such marks must be included when an integration by parts, a fine Ward
identity, or a continuous scale flow actually produces them.  They are
absent from the retained current jets in
\eqref{eq:v61-Bell-good}.

If a norm \(\|\cdot\|_{\rm an}\) takes absolute values of fine derivatives
up to order \(k\), the selector stock is
\[
 \|D_r^k\chi_j\|\leq C_{\chi,k}R_j^{-k}.
\]
This is a genuine fine-boundary cost and is not removed by moving
coordinates.

\subsection{A precise bad-set domination input}

Because the full chiral/Euclidean weight need not be positive, bad-set
suppression must be stated relative to a positive reference measure.  Let
\(d\nu_j\) be such a measure on one local fine block, normalized to one, and
let \(\|\cdot\|_{\rm coeff}\) denote the absolute coefficient norm of the
regulated boson--fermion weight.

\begin{definition}[Local bad-set domination card]
\label{def:v61-bad-card}
The scale-\(j\) chart passes the bad-set card if there are
\(A_j<\infty\) and \(q_j\in[0,1]\), independent of total volume, such that
\begin{align}
 \|e^{\Lambda(r,w)}dr\|_{\rm coeff}
 &\leq A_j\,d\nu_j(r)
 \quad\hbox{on the chart},
 \label{eq:v61-positive-domination}\\
 \int(1-\chi_j(r))\,d\nu_j(r)
 &\leq q_j.
 \label{eq:v61-reference-tail}
\end{align}
For \(n\) retained marks, the same card includes a dominating observable
\(M_{n,j}\) with
\[
 \left\|
 \sum_{\pi\in{\mathfrak P}_n}\prod_{B\in\pi}\Lambda_B
 \right\|_{\rm coeff}
 \leq M_{n,j}(r)
\]
and a tail moment
\[
 \int(1-\chi_j)M_{n,j}\,d\nu_j\leq q_{n,j}.
\]
\end{definition}

\begin{proposition}[Bad primitive and marked bounds]
\label{prop:v61-bad-bound}
Under Definition~\ref{def:v61-bad-card},
\begin{equation}
 \|Z_{{\rm b},j}(w)\|_{\rm coeff}\leq A_jq_j,
 \label{eq:v61-bad-primitive}
\end{equation}
and the \(n\)-th retained derivative satisfies the corresponding bound
\(A_jq_{n,j}\).
\end{proposition}

\begin{proof}
Take the coefficient norm inside the finite integral, use
\eqref{eq:v61-positive-domination}, and apply
\eqref{eq:v61-reference-tail}.  For derivatives use
\eqref{eq:v61-Bell-good} with \(1-\chi_j\) and the marked tail moment.
\end{proof}

For a \(d_0\)-dimensional local Gaussian reference with precision
\(A\geq\lambda I\), Chernoff's inequality gives one explicit primitive:
\begin{equation}
 \nu_j\{\|r\|\geq R\}
 \leq
 2^{d_0/2}\exp\{-\lambda R^2/4\}.
 \label{eq:v61-Gaussian-tail}
\end{equation}

\begin{proof}
For the normalized Gaussian,
\[
 {\mathbb E}\exp\{\lambda\|r\|^2/4\}
 \leq2^{d_0/2}.
\]
Apply Markov's inequality.
\end{proof}

The dimension \(d_0\) is the fixed number of fine coordinates in one local
block.  Bad blocks must then be joined into contours and summed against
their entropy, as in the established polymer ledger.

\begin{firewall}[Euclidean gravity has no supplied Gaussian domination]
\label{fire:v61-conformal-mode}
The Einstein--Hilbert Euclidean action is not bounded below in the conformal
direction.  Equation \eqref{eq:v61-Gaussian-tail} applies to a gravitational
block only after a positive contour, a conformal-factor treatment, or a
different constructive reference measure has been specified and proved
compatible with the Ward, reflection, and determinant-line ledgers.  The
DeTurck elliptic chart alone does not provide that positivity.
\end{firewall}

This is now the leading upstream obstruction to importing the complete YM
good/bad argument into the gravitational sector.

\subsection{Running-window gate}

A scale-dependent radius \(R_j\) is useful only if three independent
conditions hold:

\begin{enumerate}
\item \(2R_j\) remains inside the analytic and inverse-chart radius;
\item the good-field nonlinear remainder and chart-score stocks tend to the
      required small values on \(\|r\|\leq2R_j\); and
\item the bad primitive \(A_jq_j\), including marked moments and contour
      entropy, tends to zero or fits the stable-tail source reserve.
\end{enumerate}

The moving-coordinate theorem removes artificial retained
\(R_j^{-1}\) factors but proves none of these three scale inequalities.

\subsection{Next acquisition}

The aligned selector is now exact.  The next note will go upstream to the
conformal-mode obstruction in
Firewall~\ref{fire:v61-conformal-mode}.  It will compare the three legitimate
finite-regulator choices---real Euclidean contour with stabilizing
higher-curvature terms, a rotated conformal contour, and a Lorentzian
oscillatory construction---and state the precise positivity, locality,
reflection, and gluing conditions each would need.  No contour will be
declared valid without those conditions.

\begin{assessment}[Exact status after V61]
\label{ass:v61-status}
V61 proves the aligned good--bad decomposition, all-order absence of
retained selector derivatives, the exact cutoff--horizontal-current
identity, and a precise bad-set domination card.  It does not supply a
positive gravitational reference measure, prove bad-contour entropy,
establish the four-cumulant cluster gate, verify an Einstein running window,
control the determinant line globally, remove regulators, or establish the
SM--Einstein continuum.
\end{assessment}

\clearpage
\section{V62: the bulk conformal coverage obstruction and a reduction card}
\label{sec:v10v62}

\subsection{What V22 and V23 do not yet combine}

V22 proved that the real Euclidean Einstein--Hilbert action is unbounded
below along conformal directions.  V23--V26 constructed a rigorous CMC
constraint chart for the conformal factor of a distinguished spatial slice.
V44 then prevented double counting by putting that slice conformal factor
in the retained trace and the bulk DeTurck variable in the kernel of the
slice trace.

That bookkeeping is correct, but it does not remove all bulk conformal
directions.  This note proves the coverage gap and states the additional
reduction theorem that would be needed before the V61 positive bad-set card
can apply to gravity.

\subsection{The slice constraint leaves bulk conformal directions}

Let \(M\) be a four-dimensional regulated Euclidean block with distinguished
hypersurface \(\Sigma\).  Let
\[
 T_\Sigma:U\longrightarrow I
\]
be the metric trace map used in V44 and \(U_0=\ker T_\Sigma\).  Let \(g_0\)
be a smooth background.

\begin{proposition}[Bulk conformal modes remain in the zero-trace space]
\label{prop:v62-bulk-conformal-remains}
For every nonzero smooth scalar \(\sigma\) compactly supported in
\(M\setminus\Sigma\),
\begin{equation}
 h_\sigma=2\sigma g_0
 \label{eq:v62-conformal-variation}
\end{equation}
belongs to \(U_0\).  Hence fixing the CMC conformal trace on \(\Sigma\) does
not remove the bulk conformal subspace.
\end{proposition}

\begin{proof}
The support condition gives
\(h_\sigma|_\Sigma=0\).  Therefore
\(T_\Sigma h_\sigma=0\), so \(h_\sigma\in\ker T_\Sigma=U_0\).
There are infinitely many linearly independent compactly supported choices
of \(\sigma\), even at successively finer regulators.
\end{proof}

The DeTurck gauge condition removes diffeomorphism directions.  A generic
compactly supported conformal rescaling is not a diffeomorphism, so gauge
fixing does not change this conclusion.

\begin{corollary}[Persistence of the V22 instability]
\label{cor:v62-persistent-instability}
The real bulk integration over \(U_0\) still contains the conformal
instability proved in V22.  Consequently, the CMC trace solve and the local
DeTurck inverse do not by themselves supply the positive domination
hypothesis in the V61 bad-set card.
\end{corollary}

\begin{proof}
Apply the V22 conformal-action calculation to a scalar supported away from
\(\Sigma\).  Its entire conformal path preserves the fixed slice trace, so it
lies inside the V44 bulk integration domain.  The negative kinetic term can
be made arbitrarily large in magnitude there.
\end{proof}

This is a coverage obstruction, not a defect in the CMC theorem.

\subsection{The three alternative routes}

The current notes now distinguish three logically possible routes.

\begin{enumerate}
\item \emph{Bulk conformal reduction.}  Solve a global scalar equation for
      the bulk conformal factor with the CMC trace as boundary data, then
      integrate only the reduced variables with the exact coarea density.
\item \emph{Complex contour.}  Specify a middle-dimensional contour in the
      complexified metric space and prove convergence, homological
      invariance, reflection reality, and reconstruction of real
      observables.
\item \emph{Modified physical action or Lorentzian construction.}  Retain a
      stabilizing interaction as part of the physical continuum theory, or
      construct an oscillatory Lorentzian measure and prove a Wick-rotation
      theorem into the desired Schwinger functions.
\end{enumerate}

The earlier notes already show that rotating one negative Gaussian
coordinate does not produce a probability characteristic function for the
real metric, and that a simple removable fourth-order stabilizer has a
reflection-positive spectral obstruction.  Those shortcuts are not reopened
here.

\begin{firewall}[No transfer between incompatible measure categories]
\label{fire:v62-measure-categories}
An absolutely convergent complex contour integral or a Lorentzian
oscillatory distribution is not the positive reference measure required by
the V61 probability tail estimate.  It can replace that framework only
after a separate reconstruction theorem supplies the positivity and
continuation properties used later.
\end{firewall}

\subsection{A candidate bulk trace equation}

Write
\[
 g=e^{2\sigma}\widehat g
\]
with fixed boundary value for \(\sigma\) determined by the CMC trace, and
choose \(\widehat g\) in a declared unimodular or conformally transverse
slice.  In four dimensions,
\begin{equation}
 R_g=e^{-2\sigma}
 \left(
 R_{\widehat g}-6\widehat\Delta\sigma
 -6|\widehat\nabla\sigma|^2
 \right).
 \label{eq:v62-conformal-curvature}
\end{equation}
The vacuum trace Einstein equation \(R_g=4\Lambda\) becomes
\begin{equation}
 {\cal C}_{\rm bulk}(\sigma;\widehat g)
 :=
 R_{\widehat g}-6\widehat\Delta\sigma
 -6|\widehat\nabla\sigma|^2
 -4\Lambda e^{2\sigma}=0.
 \label{eq:v62-bulk-trace}
\end{equation}
Matter adds its appropriately rescaled trace source to this equation.

At a solution, the scalar linearisation is
\begin{equation}
 L_\sigma\eta
 =
 -6\widehat\Delta\eta
 -12\langle\widehat\nabla\sigma,\widehat\nabla\eta\rangle
 -8\Lambda e^{2\sigma}\eta
 +D_\sigma{\cal T}_{\rm mat}[\eta],
 \label{eq:v62-bulk-linearization}
\end{equation}
with homogeneous Dirichlet data for variations when the boundary conformal
trace is held fixed.

\begin{proposition}[Principal ellipticity]
\label{prop:v62-principal-ellipticity}
The principal symbol of \(L_\sigma\) is
\[
 6|\xi|_{\widehat g}^2.
\]
Thus the candidate trace equation is a second-order elliptic scalar equation
on a Riemannian block.  Ellipticity alone does not imply invertibility or
global uniqueness.
\end{proposition}

\begin{proof}
Only the term \(-6\widehat\Delta\eta\) contains two derivatives of the
variation.  Its principal symbol is \(6|\xi|^2\).  The drift and potential
terms can create a kernel or destroy a maximum-principle sign, so the final
qualification follows.
\end{proof}

\subsection{Local bulk conformal graph}

Let \(q\) collect \(\widehat g\), matter, CMC boundary data, and retained
owners.  Work between fixed regulated Banach spaces or at a finite
regulator.

\begin{theorem}[Local bulk conformal reduction]
\label{thm:v62-local-bulk-graph}
Suppose
\[
 {\cal C}_{\rm bulk}(\sigma_0;q_0)=0
\]
and the Dirichlet linearisation
\[
 L_{\sigma_0}=D_\sigma{\cal C}_{\rm bulk}(\sigma_0;q_0)
\]
is invertible.  Then there are neighbourhoods in which a unique analytic
map
\[
 \sigma=\Sigma_{\rm bulk}(q)
\]
solves the bulk trace equation with the prescribed boundary value.  Its
first response is
\begin{equation}
 D\Sigma_{\rm bulk}[h]
 =
 -L_\sigma^{-1}D_q{\cal C}_{\rm bulk}[h].
 \label{eq:v62-bulk-response}
\end{equation}
Higher responses are finite rooted-tree expressions containing one
\(L_\sigma^{-1}\) on every response edge and higher local jets of
\({\cal C}_{\rm bulk}\).
\end{theorem}

\begin{proof}
Apply the analytic implicit-function theorem to
\({\cal C}_{\rm bulk}\) in the \(\sigma\) variable.  Differentiate the
identity
\[
 {\cal C}_{\rm bulk}(\Sigma_{\rm bulk}(q);q)=0
\]
to obtain \eqref{eq:v62-bulk-response}.  Repeated differentiation and
solution by \(L_\sigma^{-1}\) give the rooted-tree structure.
\end{proof}

This theorem is local.  V61's bad-set completion requires a global graph or
a finite compatible cover whose bad complement has a proved cost.

\subsection{Coarea density and no double counting}

At a finite regulator, suppose the scalar owner \(\sigma\) and scalar
constraint have the same dimension.  Inserting the exact delta constraint
gives
\begin{equation}
 d\sigma\,
 \delta({\cal C}_{\rm bulk}(\sigma;q))
 =
 |\det L_\sigma|^{-1}
 \label{eq:v62-bulk-coarea}
\end{equation}
at the unique local solution.  The reduced metric is
\begin{equation}
 g(q)=e^{2\Sigma_{\rm bulk}(q)}\widehat g(q).
 \label{eq:v62-reduced-metric}
\end{equation}

\begin{proposition}[Bulk conformal no-double-counting rule]
\label{prop:v62-bulk-no-double}
Assume the conformal slice for \(\widehat g\) is transverse and the boundary
value of \(\Sigma_{\rm bulk}\) is exactly the CMC conformal trace.  Then local
variations split uniquely into:
\begin{enumerate}
\item a bulk conformal response determined by
      \(\Sigma_{\rm bulk}\);
\item a conformally transverse bulk variation of \(\widehat g\); and
\item retained boundary, Killing, and macroscopic data.
\end{enumerate}
The determinant in \eqref{eq:v62-bulk-coarea} is a new bulk scalar
constraint factor.  It is not the CMC scalar determinant, the DeTurck
coordinate determinant, or the Faddeev--Popov determinant.
\end{proposition}

\begin{proof}
Transversality gives a direct local splitting between the scalar conformal
direction and the chosen \(\widehat g\) slice.  The Dirichlet boundary value
identifies, rather than duplicates, the slice conformal datum.  The implicit
graph determines the scalar interior.  Each determinant differentiates a
different map on a different domain, proving the final statement.
\end{proof}

The full nested Jacobian therefore acquires another scalar block.  Its Schur
factorisation and locality must be proved; it is not obtained by appending
\(|\det L_\sigma|^{-1}\) informally.

\subsection{Reflection and positivity conditions}

Let \(\Theta\) exchange the two halves of a reflected block and preserve the
boundary data.

\begin{proposition}[Reflection of a unique bulk graph]
\label{prop:v62-reflection-graph}
If \({\cal C}_{\rm bulk}\) is reflection equivariant and the solution in the
chosen chamber is unique, then
\begin{equation}
 \Sigma_{\rm bulk}(\Theta q)=\Theta\Sigma_{\rm bulk}(q).
 \label{eq:v62-reflection-graph}
\end{equation}
\end{proposition}

\begin{proof}
Apply \(\Theta\) to the constraint equation for
\(\Sigma_{\rm bulk}(q)\).  Equivariance produces a solution for the reflected
data.  Uniqueness identifies it with
\(\Sigma_{\rm bulk}(\Theta q)\).
\end{proof}

This equivariance is necessary but not sufficient for
Osterwalder--Schrader positivity.  The reduced action and the coarea
determinant must factor through retained boundary data so that the two
halves form reflected conditional copies.  A determinant depending
nonlocally on both halves can spoil that factorisation.

\begin{definition}[Bulk conformal reduction card]
\label{def:v62-bulk-card}
A regulated block passes the card if:
\begin{enumerate}
\item the bulk trace equation has a unique global solution, or a finite
      reflection-compatible cover with controlled overlaps;
\item \(L_\sigma^{-1}\) has a uniform weighted locality bound on the chosen
      complement and all kernel modes are retained;
\item the bulk coarea determinant has a local connected expansion and its
      phase/orientation is controlled;
\item the reduced remaining variables have a positive reference measure and
      satisfy the V61 good/bad domination card;
\item reflection factorisation holds conditional on all cut data; and
\item imposing the bulk scalar equation is derived from the intended
      gravitational theory rather than adopted as an unproved change of
      theory.
\end{enumerate}
\end{definition}

\subsection{The physical-equivalence debit}

The CMC Hamiltonian constraint is part of the canonical Einstein constraint
system.  By contrast, the four-dimensional equation
\({\cal C}_{\rm bulk}=0\) is the trace of the bulk Einstein equation and is
an equation of motion for off-shell Euclidean metrics.  Inserting its delta
function generally changes the Euclidean path integral.

\begin{firewall}[Stationary reduction needs a derivation]
\label{fire:v62-physical-equivalence}
The local graph theorem and coarea formula do not prove that the reduced
measure is physically equivalent to quantum Einstein gravity.  Such
equivalence requires a derivation from the Lorentzian constraint system, an
exact renewal identity, a BRST/Batalin--Vilkovisky gauge construction, or
another stated principle.  Treating an equation of motion as a gauge
condition without that derivation would remove physical off-shell
fluctuations.
\end{firewall}

This is the decisive distinction between a mathematically normalizable
reduced model and the requested SM--Einstein continuum.

\subsection{Next acquisition}

V62 has moved the conformal bottleneck upstream.  The next note will extend
the V44 nested Schur theorem by one bulk scalar block.  It will give the exact
four-stage determinant and response operator for bulk conformal reduction,
DeTurck reconstruction on the conformally transverse slice, CMC constraints,
and interface matching.  This algebra is needed whichever physical
derivation eventually justifies the scalar reduction.

\begin{assessment}[Exact status after V62]
\label{ass:v62-status}
V62 proves that CMC trace reduction leaves infinitely many bulk conformal
directions, formulates the bulk scalar equation and its local analytic
graph, and states the exact coarea, reflection, and positivity requirements.
It does not prove a global bulk conformal solve, justify equation-of-motion
reduction as the intended quantum theory, construct a valid complex or
Lorentzian alternative, establish bad-set domination, remove regulators, or
establish the SM--Einstein continuum.
\end{assessment}

\clearpage
\section{V63: the four-stage bulk-conformal--DeTurck--CMC Schur map}
\label{sec:v10v63}

\subsection{The hierarchical derivative}

Let the unknowns be ordered as
\[
 (\sigma,u,z,b),
\]
where \(\sigma\) is the bulk conformal scalar with homogeneous variation at
the boundary, \(u\) is the conformally transverse DeTurck bulk variable,
\(z\) contains the CMC scalar/vector owners on the distinguished slice, and
\(b\) contains retained interface and macroscopic data.  Order the equations
as bulk scalar constraint, DeTurck equation, CMC constraints, and matching
equation.

At a base solution, write the derivative as
\begin{equation}
 {\mathbb J}=
 \begin{pmatrix}
 L&B_u&B_z&B_b\\
 E_\sigma&A&G&H\\
 0&0&D&K\\
 M_\sigma&P&Q&R
 \end{pmatrix}.
 \label{eq:v63-four-block}
\end{equation}
Here
\[
 L=D_\sigma{\cal C}_{\rm bulk}.
\]
The zero blocks express the chosen hierarchy: after the CMC boundary
conformal owner has been placed in \(z\), the CMC equations do not depend on
the homogeneous-boundary bulk scalar variation or on the conformally
transverse bulk filling.  If a regulator violates this hierarchy, the
general four-block Schur complement must be used and the simpler theorem
below does not apply.

Assume \(L\) is invertible and define the scalar-eliminated blocks
\begin{align}
 \widehat A&=A-E_\sigma L^{-1}B_u,
 &
 \widehat G&=G-E_\sigma L^{-1}B_z,
 &
 \widehat H&=H-E_\sigma L^{-1}B_b,
 \label{eq:v63-hat-E}\\
 \widehat P&=P-M_\sigma L^{-1}B_u,
 &
 \widehat Q&=Q-M_\sigma L^{-1}B_z,
 &
 \widehat R&=R-M_\sigma L^{-1}B_b.
 \label{eq:v63-hat-M}
\end{align}
When \(\widehat A\) and \(D\) are invertible, put
\begin{equation}
 \widehat T=\widehat Q-\widehat P\widehat A^{-1}\widehat G
 \label{eq:v63-hat-T}
\end{equation}
and
\begin{equation}
 \widehat S
 =
 \widehat R-\widehat P\widehat A^{-1}\widehat H
 -\widehat T D^{-1}K.
 \label{eq:v63-hat-S}
\end{equation}

\subsection{Exact determinant and inverse}

\begin{theorem}[Four-stage nested Schur factorisation]
\label{thm:v63-four-stage}
If
\[
 L,\qquad\widehat A,\qquad D,\qquad\widehat S
\]
are invertible, then \({\mathbb J}\) is invertible and
\begin{equation}
 \det{\mathbb J}
 =
 (\det L)(\det\widehat A)(\det D)(\det\widehat S).
 \label{eq:v63-four-determinants}
\end{equation}
On the conformal-Killing quotient,
\[
 \det D=(\det{}'\mathcal P)(\det L_\phi)
\]
as in V44.
\end{theorem}

\begin{proof}
Left multiply \eqref{eq:v63-four-block} by
\[
 {\mathbb L}_0=
 \begin{pmatrix}
 I&0&0&0\\
 -E_\sigma L^{-1}&I&0&0\\
 0&0&I&0\\
 -M_\sigma L^{-1}&0&0&I
 \end{pmatrix}.
\]
This determinant-one block operation removes the first-column entries below
\(L\).  Its lower-right three-by-three block is
\[
 \begin{pmatrix}
 \widehat A&\widehat G&\widehat H\\
 0&D&K\\
 \widehat P&\widehat Q&\widehat R
 \end{pmatrix}.
\]
The V44 nested Schur theorem factors its determinant as
\[
 (\det\widehat A)(\det D)(\det\widehat S).
\]
Multiplication by the first diagonal pivot \(\det L\) gives
\eqref{eq:v63-four-determinants}.  Invertibility follows from invertibility
of the four pivots.
\end{proof}

Let the equation increments be
\[
 (y_\sigma,y_E,y_C,y_M).
\]
Define
\begin{equation}
 \widetilde y_E=y_E-E_\sigma L^{-1}y_\sigma,
 \qquad
 \widetilde y_M=y_M-M_\sigma L^{-1}y_\sigma.
 \label{eq:v63-shifted-targets}
\end{equation}

\begin{proposition}[Four-stage response formula]
\label{prop:v63-four-response}
Under the theorem's hypotheses, the unique response is
\begin{align}
 b&=\widehat S^{-1}
 \left\{
 \widetilde y_M
 -\widehat P\widehat A^{-1}\widetilde y_E
 -\widehat T D^{-1}y_C
 \right\},
 \label{eq:v63-b-response}\\
 z&=D^{-1}(y_C-Kb),
 \label{eq:v63-z-response}\\
 u&=\widehat A^{-1}
 (\widetilde y_E-\widehat Gz-\widehat Hb),
 \label{eq:v63-u-response}\\
 \sigma&=L^{-1}
 (y_\sigma-B_uu-B_zz-B_bb).
 \label{eq:v63-sigma-response}
\end{align}
\end{proposition}

\begin{proof}
The scalar row first gives
\eqref{eq:v63-sigma-response}.  Substitution into the remaining rows gives
the scalar-eliminated three-block system with targets
\((\widetilde y_E,y_C,\widetilde y_M)\).
Apply the V59 inverse response formula to that system.
\end{proof}

This order matters.  Solving the DeTurck block with \(A^{-1}\) before
eliminating \(\sigma\) misses the feedback
\(E_\sigma L^{-1}B_u\) and generally gives the wrong interface operator.

\subsection{The complete local density}

If all four equation maps are used as equal-dimensional residual coordinates
or delta constraints, the local absolute density contains
\begin{equation}
 |\det L|^{-1}
 |\det\widehat A|^{-1}
 |\det{}'\mathcal P|^{-1}
 |\det L_\phi|^{-1}
 |\det\widehat S|^{-1}.
 \label{eq:v63-complete-density}
\end{equation}

\begin{firewall}[Five factors, five maps]
\label{fire:v63-five-factors}
The factors in \eqref{eq:v63-complete-density} differentiate, respectively,
the bulk scalar equation, scalar-eliminated DeTurck equation, momentum
constraint, CMC Hamiltonian constraint, and fully reduced interface
matching.  None is a Faddeev--Popov determinant or a physical chiral
determinant.  A factor is included only when its corresponding residual is
actually used in the measure construction.
\end{firewall}

The physical-equivalence firewall of V62 still applies: the algebra is exact
for the reduced model but does not justify inserting the bulk trace delta
constraint into quantum gravity.

\subsection{A perturbative pivot chamber}

The scalar elimination changes the DeTurck pivot.  Put
\[
 \ell=\|L^{-1}\|_\kappa,\quad
 e=\|E_\sigma\|_\kappa,\quad
 b_u=\|B_u\|_\kappa,\quad
 x=\|A^{-1}\|_\kappa.
\]

\begin{proposition}[Persistence of the DeTurck pivot]
\label{prop:v63-DeTurck-pivot}
If
\begin{equation}
 x e\ell b_u<1,
 \label{eq:v63-DeTurck-pivot-small}
\end{equation}
then \(\widehat A\) is invertible and
\begin{equation}
 \|\widehat A^{-1}\|_\kappa
 \leq
 \frac{x}{1-xe\ell b_u}.
 \label{eq:v63-hat-A-inverse}
\end{equation}
\end{proposition}

\begin{proof}
Factor
\[
 \widehat A
 =A\{I-A^{-1}E_\sigma L^{-1}B_u\}.
\]
The weighted norm of the second term is smaller than one by
\eqref{eq:v63-DeTurck-pivot-small}.  The Neumann series gives the inverse and
the displayed bound.
\end{proof}

The interface pivot \(\widehat S\) must then pass the V46 pathwise
singular-value certificate with all hatted blocks.  An unhatted floor for
\(S\) is insufficient unless the scalar-elimination defect is separately
bounded.

\begin{proposition}[Hatted-block defect stocks]
\label{prop:v63-hat-defects}
The scalar corrections obey
\begin{align}
 \|\widehat A-A\|_\kappa&\leq e\ell b_u,
 &
 \|\widehat G-G\|_\kappa&\leq e\ell b_z,
 &
 \|\widehat H-H\|_\kappa&\leq e\ell b_b,
 \label{eq:v63-E-defects}\\
 \|\widehat P-P\|_\kappa&\leq m\ell b_u,
 &
 \|\widehat Q-Q\|_\kappa&\leq m\ell b_z,
 &
 \|\widehat R-R\|_\kappa&\leq m\ell b_b,
 \label{eq:v63-M-defects}
\end{align}
where
\[
 m=\|M_\sigma\|_\kappa,\quad
 b_z=\|B_z\|_\kappa,\quad b_b=\|B_b\|_\kappa.
\]
Substitution of these six bounds into the V49 finite-defect formula gives an
explicit sufficient bound for
\(\|\widehat S-S\|_\kappa\).
\end{proposition}

\begin{proof}
Apply weighted submultiplicativity to
\eqref{eq:v63-hat-E}--\eqref{eq:v63-hat-M}.  The V49 theorem is an exact
finite-difference bound for a three-block Schur operator, so it applies with
these block differences and the perturbed inverse
\eqref{eq:v63-hat-A-inverse}.
\end{proof}

This provides a direct upstream diagnosis.  Failure can arise from the bulk
scalar inverse \(\ell\), from its coupling into the DeTurck row \(e\), from
its coupling into the matching row \(m\), or from the scalar row couplings
\(b_u,b_z,b_b\).

\subsection{Weighted locality and the new zero-mode ledger}

\begin{theorem}[Locality of the four-stage response]
\label{thm:v63-four-locality}
Assume \(L^{-1}\), \(A^{-1}\), \(D^{-1}\), and
\(\widehat S^{-1}\) belong to the V48 weighted algebra, and every coupling
block in \eqref{eq:v63-four-block} is finite range or exponentially local
with a common positive reserve.  Then all hatted blocks and every response
operator in
\eqref{eq:v63-b-response}--\eqref{eq:v63-sigma-response} belong to a
slightly weakened weighted algebra, with volume-independent norms.
\end{theorem}

\begin{proof}
Equations \eqref{eq:v63-hat-E}--\eqref{eq:v63-hat-M} contain only sums and
products in the weighted algebra.  The four response equations contain the
same operations and the four declared inverses.  Closure and
submultiplicativity prove the claim; a finite weakening absorbs support
marks from derivatives.
\end{proof}

If \(L\) has a kernel, its modes cannot be hidden in a pseudoinverse.  They
must be added to the retained macroscopic block, and the scalar determinant
becomes primed in a fixed constant-rank chart.  A kernel crossing is a chart
boundary, exactly as for the conformal-Killing and chiral zero modes.

\subsection{Reflection compatibility}

Suppose every block in \eqref{eq:v63-four-block} intertwines the reflected
half-space actions and the four pivots are unique inverses on their chosen
complements.

\begin{proposition}[Reflection of the nested reduction]
\label{prop:v63-reflection}
The hatted blocks, \(\widehat T\), \(\widehat S\), and all four response maps
intertwine reflection.  Paired interior determinant factors are complex
conjugates in reflected determinant-line frames.
\end{proposition}

\begin{proof}
Intertwining passes to unique inverses.  Every hatted and Schur block is a
sum or product of intertwining maps, so it also intertwines.  The response
formulae then give reflected solutions by uniqueness.  Adjoint reflection
complex conjugates the finite determinant in paired frames.
\end{proof}

The interface and bulk scalar determinant functionals still require the
boundary-state or local-polymer positivity test; conjugate pairing of numbers
alone is not OS positivity.

\subsection{Next acquisition}

The new first pivot is \(L\), the bulk scalar linearisation.  The next note
will calculate its flat-collar Dirichlet symbol, including cosmological and
matter-trace potentials, and determine an explicit sign chamber for a
positive scalar inverse.  It will also show where that chamber fails and
which zero or negative modes must be retained rather than inverted.

\begin{assessment}[Exact status after V63]
\label{ass:v63-status}
V63 proves the exact four-stage determinant and inverse response, a
perturbative hatted-pivot chamber, weighted locality, and reflection
equivariance.  It does not prove invertibility or global uniqueness of the
bulk scalar equation, justify its physical insertion, establish reduced
measure positivity, control all zero modes, remove regulators, or establish
the SM--Einstein continuum.
\end{assessment}

\clearpage
\section{V64: flat-collar spectrum of the bulk conformal pivot}
\label{sec:v10v64}

\subsection{Constant-coefficient base}

Linearise the V62 bulk scalar equation at a flat conformal background on
\[
 {\cal C}_+=[0,L_+]\times{\mathbb T}^3_\ell,
 \qquad
 {\cal C}_-=[-L_-,0]\times{\mathbb T}^3_\ell.
\]
Let the constant zeroth-order coefficient be
\begin{equation}
 \nu:=-8\Lambda+v_{\rm mat},
 \label{eq:v64-nu}
\end{equation}
where \(v_{\rm mat}\) is the derivative of the rescaled matter-trace source
with respect to the bulk conformal scalar at the base point.  No sign is
assigned to \(v_{\rm mat}\).  The pivot is
\begin{equation}
 L_0=-6(\partial_x^2+\Delta_y)+\nu.
 \label{eq:v64-L-zero}
\end{equation}
Interior variations have homogeneous Dirichlet data at the interface and at
the outer face when \(L_0^{-1}\) is formed.  The interface response is
obtained by prescribing the interface trace and retaining the outer
Dirichlet condition.

\subsection{Interior spectrum}

For tangential momentum
\(k\in(2\pi/\ell){\mathbb Z}^3\) and normal Dirichlet index \(n\geq1\), the
positive-collar eigenvalues are
\begin{equation}
 \lambda_{n,k}^{(+)}
 =
 6\left\{\frac{\pi^2n^2}{L_+^2}+|k|^2\right\}+\nu,
 \label{eq:v64-interior-spectrum}
\end{equation}
and similarly with \(L_-\).

\begin{proposition}[Interior pivot chamber]
\label{prop:v64-interior-chamber}
Put \(L_{\max}=\max\{L_+,L_-\}\).  Both collar Dirichlet pivots are strictly
positive if and only if
\begin{equation}
 \nu>-\frac{6\pi^2}{L_{\max}^2}.
 \label{eq:v64-interior-positive}
\end{equation}
In that chamber,
\begin{equation}
 \|L_0^{-1}\|_{L^2\to L^2}
 \leq
 \left(\nu+\frac{6\pi^2}{L_{\max}^2}\right)^{-1}.
 \label{eq:v64-interior-inverse}
\end{equation}
At equality a zero mode occurs on every longest collar.  Below the threshold
the quadratic form has a negative mode, although the matrix can remain
invertible away from the discrete crossing values.
\end{proposition}

\begin{proof}
Separation of variables gives
\eqref{eq:v64-interior-spectrum}.  Its minimum over \(n,k\) and both collars
is attained at \(n=1,k=0\) on a longest collar.  Positivity and the inverse
bound follow from the spectral theorem.  Equality gives the corresponding
kernel.  Decreasing \(\nu\) through it makes that eigenvalue negative;
further kernels occur only at the displayed discrete spectral values.
\end{proof}

Invertibility alone is sufficient for a local implicit graph.  Positivity is
stronger and is useful for maximum-principle and orientation control.

\subsection{Exact two-sided Dirichlet-to-Neumann symbol}

For fixed \(k\), put
\begin{equation}
 \lambda_k=|k|^2+\nu/6.
 \label{eq:v64-lambda-k}
\end{equation}
Define
\begin{equation}
 d_L(\lambda)=
 \begin{cases}
 6\sqrt{\lambda}\coth(L\sqrt{\lambda}),&\lambda>0,\\
 6/L,&\lambda=0,\\
 6\sqrt{-\lambda}\cot(L\sqrt{-\lambda}),&\lambda<0,
 \end{cases}
 \label{eq:v64-dL}
\end{equation}
whenever the last denominator is nonzero.

\begin{theorem}[Bulk scalar interface symbol]
\label{thm:v64-interface-symbol}
The sum of the two outward conormal derivatives has multiplier
\begin{equation}
 s_\nu(k)=d_{L_+}(\lambda_k)+d_{L_-}(\lambda_k).
 \label{eq:v64-interface-symbol}
\end{equation}
On the component of the resolvent set satisfying
\[
 \lambda_k>-\pi^2/L_{\max}^2,
\]
the function \(s_\nu(k)\) is increasing as a function of \(|k|^2\).
Consequently its smallest tangential mode is \(k=0\).
\end{theorem}

\begin{proof}
The mode equation is
\[
 -u''+\lambda_ku=0.
\]
For positive \(\lambda_k\), the solution with interface value \(b_k\) and
zero outer value is the hyperbolic-sine solution used in V47; its outward
conormal derivative is the first row of \eqref{eq:v64-dL}.  At zero it is
affine.  For negative \(\lambda_k=-a^2\), the solution is
\[
 u(x)=b_k\frac{\sin(a(L-x))}{\sin(aL)},
\]
and the conormal derivative is \(6a\cot(aL)b_k\).  Adding the two collars
gives \eqref{eq:v64-interface-symbol}.

For \(\lambda>0\), monotonicity was proved in the V47 scalar collar
calculation.  For \(\lambda=-a^2\) with \(0<aL<\pi\),
\[
 \frac{d}{da}\{a\cot(aL)\}
 =\cot(aL)-aL\operatorname{csc}^2(aL)<0,
\]
because \(\sin x\cos x-x<0\) on \((0,\pi)\).  Since
\(d\lambda/da=-2a<0\), the derivative with respect to \(\lambda\) is
positive.  Sum the two increasing functions.
\end{proof}

\subsection{Three distinct thresholds}

\begin{proposition}[Strong, interface, and interior chambers]
\label{prop:v64-three-chambers}
The following are sufficient nested chambers:
\begin{enumerate}
\item If \(\nu\geq0\), the interior operator is positive, the maximum
      principle has a nonnegative potential, and
      \begin{equation}
       s_\nu(k)\geq6(L_+^{-1}+L_-^{-1}).
       \label{eq:v64-strong-floor}
      \end{equation}
\item If
      \begin{equation}
       \nu>-\frac{3\pi^2}{2L_{\max}^2},
       \label{eq:v64-interface-positive}
      \end{equation}
      each one-sided zero-momentum Dirichlet-to-Neumann multiplier is
      positive, hence \(s_\nu(k)>0\) for every \(k\).
\item If only \eqref{eq:v64-interior-positive} holds, the interior pivot is
      positive, but the interface multiplier can be negative.
\end{enumerate}
\end{proposition}

\begin{proof}
Item one follows from
\(z\coth(Lz)\geq L^{-1}\).  For item two, if \(\nu<0\), set
\(a=\sqrt{-\nu/6}\).  The condition gives
\(aL_{\max}<\pi/2\), so
\(a\cot(aL_\pm)>0\).  Monotonicity in Theorem
\ref{thm:v64-interface-symbol} handles all higher \(k\).
For item three, choose
\[
 \frac{\pi}{2}<aL_{\max}<\pi.
\]
Then the longest-collar interior first eigenvalue remains positive, while
its zero-mode conormal multiplier \(6a\cot(aL_{\max})\) is negative.  The
other collar contribution need not compensate it.
\end{proof}

Thus a positive bulk scalar inverse does not automatically give a positive
interface Schur row.  The V63 pivot and interface certificates are genuinely
separate.

When \(\nu<0\) and the interface chamber holds, the exact zero-mode floor is
\begin{equation}
 \gamma_{\sigma,\partial}
 =
 6a\{\cot(aL_+)+\cot(aL_-)\},
 \qquad a=\sqrt{-\nu/6}.
 \label{eq:v64-negative-potential-floor}
\end{equation}
For \(\nu\geq0\), it is
\begin{equation}
 \gamma_{\sigma,\partial}
 =
 6\mu\{\coth(\mu L_+)+\coth(\mu L_-)\},
 \qquad \mu=\sqrt{\nu/6},
 \label{eq:v64-positive-potential-floor}
\end{equation}
with the continuous value in
\eqref{eq:v64-strong-floor} at \(\nu=0\).

\subsection{Weighted locality and collapse at a crossing}

At a lattice regulator, the tangential response is the meromorphic function
\(d_L(\lambda_{\rm lat}+\nu/6)\).  Its first pole occurs when
\[
 \lambda_{\rm lat}
 =
 -\left(\frac{\nu}{6}+\frac{\pi^2}{L^2}\right).
\]
Define the pole distance
\begin{equation}
 \delta_\sigma
 :=
 \frac{\nu}{6}+\frac{\pi^2}{L_{\max}^2}.
 \label{eq:v64-pole-distance}
\end{equation}
It is positive exactly in the interior positivity chamber.

\begin{proposition}[Locality reserve of the scalar response]
\label{prop:v64-locality-reserve}
If \(\delta_\sigma>0\), the V48 Bernstein-ellipse argument gives exponential
off-diagonal decay of the regulated scalar inverse and
Dirichlet-to-Neumann response for every ellipse lying before the pole.
As \(\delta_\sigma\downarrow0\), both the inverse bound and the best
available exponential locality exponent degenerate.
\end{proposition}

\begin{proof}
The interior operator is finite range with spectral distance at least
\(6\delta_\sigma\) from zero.  Apply the V48 inverse-locality theorem.
The response function is meromorphic with nearest pole at the stated
negative tangential value, so the V48 analytic functional calculus applies.
The admissible Bernstein ellipse shrinks to the spectral interval as the
pole approaches zero, forcing its logarithmic decay exponent to vanish.
\end{proof}

A near-kernel therefore consumes two reserves at once: response size and
polymer locality.

\subsection{Curved and variable-potential perturbations}

Let the true scalar pivot be
\[
 L=L_0+{\cal V},
\]
where \({\cal V}\) includes curvature, drift, variable matter response, and
coefficient discretization defects.  In the weighted algebra, put
\[
 \varepsilon_\sigma=\|L_0^{-1}{\cal V}\|_\kappa.
\]

\begin{corollary}[Perturbed scalar pivot]
\label{cor:v64-perturbed-pivot}
If \(\varepsilon_\sigma<1\), then
\begin{equation}
 \|L^{-1}\|_\kappa
 \leq
 \frac{\|L_0^{-1}\|_\kappa}{1-\varepsilon_\sigma}.
 \label{eq:v64-perturbed-inverse}
\end{equation}
The determinant orientation and kernel dimension remain constant along the
straight perturbation path.
\end{corollary}

\begin{proof}
Factor
\[
 L=L_0(I+L_0^{-1}{\cal V})
\]
and use the Neumann series.  Every matrix on the straight path is
invertible by the same bound, so no eigenvalue crosses zero and the
determinant orientation is constant.
\end{proof}

The hatted DeTurck chamber in V63 then requires
\[
 \|A^{-1}\|_\kappa
 \|E_\sigma\|_\kappa
 \|L^{-1}\|_\kappa
 \|B_u\|_\kappa<1.
\]
This displays the direct cost of approaching the scalar threshold.

\subsection{Mode ledger}

At a crossing
\[
 \nu=-6\left\{\frac{\pi^2n^2}{L_\pm^2}+|k|^2\right\},
\]
the corresponding scalar mode must be retained in a constant-rank
macroscopic block or the chart must end.  Continuing with a pseudoinverse
while keeping the determinant unprimed would violate both the V63 response
formula and determinant ledger.  Negative but nonzero modes do not prevent
a local implicit graph, but they destroy the positive-pivot interpretation
and change the Morse index relative to the strong chamber.

\begin{firewall}[A positive scalar pivot does not justify the reduction]
\label{fire:v64-positive-not-physical}
Even the strong chamber \(\nu\geq0\) proves only analytic and positivity
properties of the candidate bulk trace equation.  It does not prove that
imposing this equation is physically equivalent to the intended quantum
Einstein theory.  The V62 physical-equivalence debit remains.
\end{firewall}

\subsection{Next acquisition}

V64 supplies the flat scalar floor and locality reserve required by V63.
V65 is the next compulsory companion audit.  It will inspect the newest YM
and NS notes for a finite marked-forest estimate, a moving-window
large-field mechanism, or a treatment of a sign-indefinite mode.  The
post-audit note will then choose between a curved scalar maximum-principle
chamber and a physical derivation of the bulk conformal reduction.

\begin{assessment}[Exact status after V64]
\label{ass:v64-status}
V64 computes the exact flat interior spectrum and interface symbol, separates
the three positivity thresholds, and proves perturbative weighted locality.
It does not prove a curved global scalar solve, a positive gravitational
measure, physical equivalence of the reduction, bad-set domination,
regulator removal, or the SM--Einstein continuum.
\end{assessment}

\section{Companion audit: analytic selectors and the forest layer}
\label{sec:v10v65}

This is the compulsory fifth-version audit following V60.  It was performed
against the exact files
\begin{align*}
 &\texttt{Renewal\_Yang\_Mills\_Mass\_Gap\_Research\_Notes\_V65.tex},\\
 &\qquad
 \operatorname{SHA256}
 =\texttt{7A5AC49398BF5FE00E454FA469459BC6216E6E1109C79FC376A15576C72C858C},\\
 &\texttt{Renewal\_NS\_Stability\_Research\_Notes\_V26.tex},\\
 &\qquad
 \operatorname{SHA256}
 =\texttt{82C887F8C2C69E4F28FAD7C3DC8DE5B9099CB5A4BF431EBA1FFF3E91935978AD}.
\end{align*}
The Yang--Mills companion is newer than the file inspected at V60.  The
Navier--Stokes companion is unchanged.  This section records only the
transferable mathematical mechanism; coupling powers, Haar positivity, and
Wilson-action estimates remain properties of the Yang--Mills programme.

\subsection{The obstruction found by the companion}

V61 used a compact smooth selector in a finite retained-derivative card.
That use is legitimate when one fixes in advance the largest derivative
order.  It is not legitimate inside the all-order analytic norm used by the
V18--V19 forest layer and the V53 determinant activities.

\begin{proposition}[No compact analytic selector]
\label{prop:v65-no-compact-analytic-selector}
Let \(U\) be a connected real analytic manifold and let
\(\chi:U\to\mathbb R\) be real analytic.  If \(\chi\) vanishes on a
nonempty open subset of \(U\), then \(\chi\) vanishes identically.
Consequently there is no nonzero real analytic function on
\(\mathbb R^n\) with compact support.
\end{proposition}

\begin{proof}
In a real analytic coordinate ball, all Taylor coefficients at every point
of the open zero set vanish.  The identity theorem propagates the zero germ
along every chain of overlapping coordinate balls.  Connectedness then
gives \(\chi=0\) on \(U\).  A compactly supported function vanishes on the
nonempty open complement of its support, so the final assertion follows.
\end{proof}

There is a second, independent failure.  A forest formula can create
arbitrarily many unmarked differentiation vertices as the polymer size
grows.  Hence a bound involving only
\(\max_{k\leq k_0}\|\partial^k\chi\|\) for fixed \(k_0\) cannot close an
all-polymer analytic summation.  The compact selector is therefore confined
to finite-order large-field bookkeeping.  It is removed from every
all-order analytic activity.

\subsection{Transferable replacement}

The companion replaces the compact selector by an exact soft split.  In
gravity the split must be made in the aligned, reduced coordinates of the
V54--V61 transport chart.  Let \(r_y\) denote the local fluctuation in one
block, let \(Q_y(r_y)\geq0\) be a real analytic reference distance with
\(Q_y(0)=0\), and let \(R_y>0\).  Define
\begin{equation}
 g_y(r_y)=\exp\{-Q_y(r_y)/R_y^2\},
 \qquad
 b_y(r_y)=1-g_y(r_y).
 \label{eq:v65-soft-split}
\end{equation}
Then
\begin{equation}
 g_y+b_y=1,
 \qquad
 0\leq g_y,b_y\leq1,
 \qquad
 b_y\leq Q_y/R_y^2.
 \label{eq:v65-soft-elementary}
\end{equation}
The last bound follows from \(1-e^{-t}\leq t\) for \(t\geq0\).

The quantity \(Q_y\) is an auxiliary positive reference distance.  It is
not the Euclidean Einstein action and does not repair the conformal sign
problem identified at V61--V64.  Its use is conditional on the construction
of a positive reference measure and on comparison of the interacting
density with that measure.

For a finite block set \(\Lambda\), the split is the exact identity
\begin{equation}
 1
 =
 \prod_{y\in\Lambda}(g_y+b_y)
 =
 \sum_{Y\subset\Lambda}
 \prod_{y\in Y}b_y
 \prod_{y\in\Lambda\setminus Y}g_y .
 \label{eq:v65-subset-resolution}
\end{equation}
Thus soft-bad marks are never discarded.  A far large-field contour and the
soft transition are different owners in the activity ledger.

\subsection{What may and may not be imported}

Suppose temporarily that the positive reference law factors,
\(\nu_\Lambda=\bigotimes_{y\in\Lambda}\nu_y\), and
\begin{equation}
 \int Q_y\,d\nu_y\leq m_y.
 \label{eq:v65-reference-moment}
\end{equation}
For every \(Y\subset\Lambda\),
\begin{equation}
 \int\prod_{y\in Y}b_y\,d\nu_\Lambda
 \leq
 \prod_{y\in Y}\frac{m_y}{R_y^2}.
 \label{eq:v65-bad-product}
\end{equation}
This follows directly from \eqref{eq:v65-soft-elementary} and factorization.
For a correlated reference law, \eqref{eq:v65-bad-product} is not available
without a mixing, conditional-moment, or tree-decoupling estimate.  That
estimate is a named debit.

Let \({\cal A}\) be any unital submultiplicative analytic Banach algebra and
put \(q_y=Q_y/R_y^2\).  The exponential series gives
\begin{equation}
 \|g_y\|_{\cal A}\leq e^{\|q_y\|_{\cal A}},
 \qquad
 \|b_y\|_{\cal A}
 \leq e^{\|q_y\|_{\cal A}}-1.
 \label{eq:v65-analytic-selector-bound}
\end{equation}
Unlike a compact smooth cutoff, this estimate controls all derivatives
encoded by the analytic norm.  It also displays the price: the soft-good
factor inserts the local analytic vertex \(q_y\), and its norm must be paid
in the polymer activity.

If \(Q_y\) and \(R_y\) are fixed functions of the aligned coordinate
\(r_y\), their retained horizontal derivative is zero at fixed \(r_y\).
If either depends on the retained variable, its derivative belongs to the
horizontal-current/Jacobian identity of V61 and cannot be omitted.

\begin{ledger}[Corrected selector ownership]
\label{led:v65-selector-ownership}
The roles are now separated as follows.
\begin{enumerate}
\item A compact \(C^\infty\) selector may be used in a finite-order
      deterministic chart or in an exterior contour estimate.
\item The all-order forest and determinant norms use
      \eqref{eq:v65-soft-split}; their new local activity is bounded by
      \(e^{\|Q_y/R_y^2\|_{\cal A}}-1\).
\item A soft-bad mark is paid by a reference moment such as
      \(m_y/R_y^2\), together with whatever domination or mixing comparison
      the interacting measure requires.
\item A far-bad contour is paid by a separate large-field tail estimate.
      It is not identified with the soft-transition activity.
\end{enumerate}
\end{ledger}

The numerical Yang--Mills primitive of order \(\beta^{-1/5}\) does not
transfer: it uses the compact gauge group, Haar probability, Wilson
positivity, and its own coupling scaling.  Likewise the Navier--Stokes note
adds no new positive gravitational reference measure.  The only imported
object is the analytic soft-split architecture and its ownership ledger.

\subsection{Direction selected by the audit}

V66 will prove the finite-volume soft-split theorem needed by the SM forest
layer.  It will include the exact subset resolution, an all-order Banach
algebra estimate, a conditional bad-mark product bound, and the resulting
entropy gate.  The theorem will state explicitly the simultaneous running
window required of \(R_y\).  It will not assume that the missing positive
Einstein reference law already exists.

\begin{assessment}[Exact status after V65]
\label{ass:v65-status}
The audit detects and corrects an analytic-regularity mismatch in the V61
selector.  It supplies a valid replacement mechanism but no
gravity-specific moment estimate, mixing theorem, positive conformal
measure, regulator removal, or SM--Einstein continuum.
\end{assessment}
\section{Finite-volume analytic soft-split theorem}
\label{sec:v10v66}

V65 identified a mismatch between compact smooth selectors and the
all-order analytic forest norm.  This section proves the replacement card
in a form that allows correlations.  Its hypotheses are deliberately
measure-theoretic: the absent positive gravitational law is not smuggled
into the proof.

\subsection{Set-up}

Let \(\Lambda\) be a finite set of regulator blocks and let
\[
 (\Omega_\Lambda,{\cal F}_\Lambda,\mu_\Lambda)
\]
be a probability space carrying aligned local coordinates
\(r=(r_y)_{y\in\Lambda}\).  For each \(y\), let \(Q_y(r_y)\geq0\) on the real
slice, let \(R_y>0\), and write
\begin{equation}
 q_y=\frac{Q_y}{R_y^2},
 \qquad
 g_y=e^{-q_y},
 \qquad
 b_y=1-e^{-q_y}.
 \label{eq:v66-definitions}
\end{equation}
Let
\[
 {\cal F}_{-y}:=\sigma(r_z:z\neq y).
\]
The correlated conditional-moment hypothesis is
\begin{equation}
 \operatorname*{ess\,sup}
 {\mathbb E}_{\mu_\Lambda}
 \bigl[q_y\mid{\cal F}_{-y}\bigr]
 \leq p_y.
 \label{eq:v66-conditional-moment}
\end{equation}
This condition is uniform in the exterior field.  For a product reference
law it reduces to
\(\int Q_y\,d\mu_y/R_y^2\leq p_y\), but no factorization is assumed below.

\begin{theorem}[Exact soft split with correlated bad-mark bound]
\label{thm:v66-soft-split}
Under \eqref{eq:v66-conditional-moment}, the following statements hold.
\begin{enumerate}
\item For every nonnegative measurable \(H\),
      \begin{equation}
       {\mathbb E}_{\mu_\Lambda}H
       =
       \sum_{Y\subset\Lambda}
       {\mathbb E}_{\mu_\Lambda}
       \left[
        H
        \prod_{y\in Y}b_y
        \prod_{z\in\Lambda\setminus Y}g_z
       \right].
       \label{eq:v66-exact-resolution}
      \end{equation}
\item For every \(Y\subset\Lambda\),
      \begin{equation}
       {\mathbb E}_{\mu_\Lambda}
       \prod_{y\in Y}b_y
       \leq
       \prod_{y\in Y}p_y.
       \label{eq:v66-correlated-product}
      \end{equation}
\item If \(0\leq H\leq K\), then the \(Y\)-marked summand in
      \eqref{eq:v66-exact-resolution} is at most
      \(K\prod_{y\in Y}p_y\).
\end{enumerate}
\end{theorem}

\begin{proof}
The pointwise relations \(g_y+b_y=1\) and \(0\leq g_y,b_y\leq1\) give
\[
 H=H\prod_{y\in\Lambda}(g_y+b_y).
\]
Expansion of this finite product and integration prove
\eqref{eq:v66-exact-resolution}.

For the second statement, use \(1-e^{-t}\leq t\) for \(t\geq0\), so
\(b_y\leq q_y\).  Enumerate \(Y=\{y_1,\ldots,y_n\}\).  The product of the
factors with indices different from \(y_1\) is
\({\cal F}_{-y_1}\)-measurable.  Therefore
\begin{align*}
 {\mathbb E}\prod_{i=1}^n b_{y_i}
 &=
 {\mathbb E}\left[
  \prod_{i=2}^n b_{y_i}\,
  {\mathbb E}(b_{y_1}\mid{\cal F}_{-y_1})
 \right]\\
 &\leq
 p_{y_1}{\mathbb E}\prod_{i=2}^n b_{y_i}.
\end{align*}
Iteration proves \eqref{eq:v66-correlated-product}.  The last statement
uses \(g_z\leq1\), \(H\leq K\), and the same bound.
\end{proof}

The theorem shows exactly what correlation estimate is required.  An
unconditional one-block moment is insufficient: the exterior field could
concentrate the conditional law of one block on a large-\(Q_y\) region.

\subsection{All-order analytic estimate}

For each block let \({\cal A}_y\) be a unital Banach algebra of functions on
a complex neighborhood of the aligned chart, normalized by
\(\|1\|_{{\cal A}_y}=1\), with
\[
 \|uv\|_{{\cal A}_y}
 \leq
 \|u\|_{{\cal A}_y}\|v\|_{{\cal A}_y}.
\]
The norm may be the strong support--grade--analytic norm of V57 restricted
to one block.

\begin{theorem}[All-order soft-selector activity]
\label{thm:v66-analytic-activity}
If \(q_y\in{\cal A}_y\) and
\(\|q_y\|_{{\cal A}_y}\leq\epsilon_y\), then
\begin{align}
 \|g_y\|_{{\cal A}_y}
 &\leq e^{\epsilon_y},
 \label{eq:v66-g-bound}\\
 \|g_y-1\|_{{\cal A}_y}
 =
 \|b_y\|_{{\cal A}_y}
 &\leq e^{\epsilon_y}-1.
 \label{eq:v66-b-bound}
\end{align}
Consequently the exact replacement \(g_y=1+(g_y-1)\) creates a local
analytic activity with size
\begin{equation}
 \eta_{{\rm soft},y}:=e^{\epsilon_y}-1.
 \label{eq:v66-soft-activity}
\end{equation}
For distinct blocks in a tensor-product submultiplicative norm,
\begin{equation}
 \left\|
  \prod_{y\in Y}(g_y-1)
 \right\|
 \leq
 \prod_{y\in Y}\eta_{{\rm soft},y}.
 \label{eq:v66-soft-product-analytic}
\end{equation}
\end{theorem}

\begin{proof}
The exponential series converges absolutely in a unital Banach algebra.
Submultiplicativity gives
\[
 \|e^{-q_y}\|
 \leq
 \sum_{n\geq0}\frac{\|q_y\|^n}{n!}
 \leq e^{\epsilon_y}.
\]
Removing the \(n=0\) term gives
\[
 \|e^{-q_y}-1\|
 \leq
 \sum_{n\geq1}\frac{\|q_y\|^n}{n!}
 \leq e^{\epsilon_y}-1.
\]
Since \(b_y=-(g_y-1)\), this proves
\eqref{eq:v66-g-bound}--\eqref{eq:v66-b-bound}.  Repeated
submultiplicativity proves \eqref{eq:v66-soft-product-analytic}.
\end{proof}

All Taylor derivatives encoded by \({\cal A}_y\) are controlled at once.
There is no fixed derivative ceiling and no derivative of a nonanalytic
boundary.

The same factor can instead be absorbed into the local exponent:
\[
 g_y=e^{-q_y}.
\]
In that convention \(q_y\), rather than \(g_y-1\), is the added local
vertex.  These are two exact expansions of the same object and must not be
charged simultaneously.

\subsection{Contour entropy gate}

Let the regulator blocks carry a bounded-degree adjacency graph.  Assume
that the number \(N_n(y_0)\) of connected \(n\)-block subsets containing a
fixed block \(y_0\) satisfies
\begin{equation}
 N_n(y_0)\leq C_{\rm ent}^{\,n-1}.
 \label{eq:v66-connected-count}
\end{equation}
Let \(p=\sup_y p_y\).

\begin{corollary}[Annealed soft-bad entropy bound]
\label{cor:v66-entropy-gate}
For every \(\kappa\geq0\), if
\begin{equation}
 C_{\rm ent}e^\kappa p<1,
 \label{eq:v66-entropy-condition}
\end{equation}
then
\begin{equation}
 \sum_{\substack{Y\ni y_0\\Y\ {\rm connected}}}
 e^{\kappa|Y|}
 {\mathbb E}_{\mu_\Lambda}
 \prod_{y\in Y}b_y
 \leq
 \frac{e^\kappa p}
 {1-C_{\rm ent}e^\kappa p}.
 \label{eq:v66-entropy-sum}
\end{equation}
If an additional local contour factor is bounded by \(K^{|Y|}\), replace
\(p\) by \(Kp\) in
\eqref{eq:v66-entropy-condition}--\eqref{eq:v66-entropy-sum}.
\end{corollary}

\begin{proof}
Group the connected sets by \(n=|Y|\), use
\eqref{eq:v66-correlated-product} and
\eqref{eq:v66-connected-count}, and sum the geometric series:
\[
 \sum_{n\geq1}
 C_{\rm ent}^{\,n-1}e^{\kappa n}p^n
 =
 \frac{e^\kappa p}
 {1-C_{\rm ent}e^\kappa p}.
\]
The extra factor \(K^n\) changes \(p\) to \(Kp\).
\end{proof}

This is an annealed mark estimate.  A complete polymer theorem still needs
the V53 determinant activities, the V57 grade tail, the V59 chart score,
and the interaction comparison to fit in one norm budget.

\subsection{Far-field owner}

Fix \(A_y>0\) and set
\[
 F_y(A_y)=\{Q_y\geq A_yR_y^2\}.
\]
On this event \(b_y\geq1-e^{-A_y}\), so
\begin{equation}
 \mu_\Lambda(F_y(A_y))
 \leq
 \frac{{\mathbb E}b_y}{1-e^{-A_y}}
 \leq
 \frac{p_y}{1-e^{-A_y}}.
 \label{eq:v66-far-from-soft}
\end{equation}
A stronger far-tail estimate requires stronger input.  If
\begin{equation}
 {\mathbb E}_{\mu_\Lambda}e^{\theta_yQ_y}\leq M_y
 \qquad(\theta_y>0),
 \label{eq:v66-exponential-moment}
\end{equation}
then Markov's inequality gives
\begin{equation}
 \mu_\Lambda(F_y(A_y))
 \leq
 M_y e^{-\theta_yA_yR_y^2}.
 \label{eq:v66-far-exponential}
\end{equation}
The event indicator in this last decomposition belongs to the
measure-theoretic contour layer, outside the analytic forest
differentiation.  Equation \eqref{eq:v66-definitions} remains the analytic
identity.

\subsection{Aligned retained transport}

Let \(w\) be a retained coordinate and take the horizontal derivative at
fixed aligned fluctuation \(r\).  If \(Q_y=Q_y(r_y)\) and \(R_y\) are
independent of \(w\), then
\begin{equation}
 \nabla_w^{\rm hor}q_y
 =
 \nabla_w^{\rm hor}g_y
 =
 \nabla_w^{\rm hor}b_y
 =0.
 \label{eq:v66-horizontal-zero}
\end{equation}
The selector then creates no retained cutoff derivative.  The moving chart
still contributes its V61 horizontal current and measured Jacobian.  If
\(Q_y\) or \(R_y\) depends on \(w\), the chain rule gives
\begin{equation}
 \nabla_w^{\rm hor}g_y
 =
 -g_y\nabla_w^{\rm hor}q_y,
 \qquad
 \nabla_w^{\rm hor}b_y
 =
 g_y\nabla_w^{\rm hor}q_y,
 \label{eq:v66-horizontal-dependent}
\end{equation}
and these terms must be entered in that current.

\subsection{Running window}

At regulator scale \(j\), a sufficient selector window is
\begin{align}
 p_j
 &:=
 \sup_y\operatorname*{ess\,sup}
 {\mathbb E}(Q_{j,y}/R_{j,y}^2\mid{\cal F}_{-y}),
 \label{eq:v66-window-p}\\
 \epsilon_j
 &:=
 \sup_y\|Q_{j,y}/R_{j,y}^2\|_{{\cal A}_{j,y}}.
 \label{eq:v66-window-epsilon}
\end{align}
It must obey
\begin{equation}
 C_{{\rm ent},j}e^{\kappa_j}K_jp_j<1,
 \qquad
 e^{\epsilon_j}-1<\eta_{{\rm soft},j}^{\rm budget},
 \label{eq:v66-running-window}
\end{equation}
and the complex domain defining \({\cal A}_{j,y}\) must remain inside the
aligned DeTurck--CMC--scalar chart.  These are simultaneous requirements,
not consequences of choosing \(R_{j,y}\) large: enlarging a window beyond
the valid chart is forbidden, while shrinking it worsens both ratios.
A continuum proof must produce scale-dependent fluctuation estimates that
leave a nonempty interval of admissible \(R_{j,y}\).

\begin{firewall}[Conditional theorem]
\label{fire:v66-positive-law}
Theorem \ref{thm:v66-soft-split} is rigorous for any positive probability
law satisfying \eqref{eq:v66-conditional-moment}.  V66 does not construct
such a law for Euclidean Einstein fluctuations.  In particular, the
positive auxiliary distance \(Q_y\) cannot be substituted for the
sign-indefinite physical action.
\end{firewall}

\subsection{Next acquisition}

V67 will place
\[
 \eta_{\rm soft},\quad
 p,\quad
 \eta_{\rm det},\quad
 \eta_{\rm tail},\quad
 \Sigma_0,\quad
 \Sigma_1
\]
in a single finite-volume budget card.  The purpose is to expose which
smallness conditions already follow from V44--V66 and which still require
a positive gravitational reference measure or a physical derivation of
the bulk conformal reduction.

\begin{assessment}[Exact status after V66]
\label{ass:v66-status}
V66 proves the exact correlated soft-bad product bound, the all-order
analytic selector estimate, and the associated contour entropy gate.  The
result repairs the V61 selector at the forest level.  The conditional
moment and exponential-tail hypotheses, the positive gravitational law,
physical equivalence of the scalar reduction, uniform regulator bounds,
and the SM--Einstein continuum remain open.
\end{assessment}

\section{A noncircular master budget for the finite-volume step}
\label{sec:v10v67}

The preceding notes produced several quantitative gates, but they play
different roles.  Adding every displayed constant into one number would
double count the chart Jacobian and would mix a measure activity with an RG
source.  This section gives a single card whose rows are coupled but whose
owners are disjoint.

\subsection{Anchored activity stocks}

For a connected-polymer activity \(a=(a(X))\), define
\begin{equation}
 \|a\|_{\kappa,\circ}
 :=
 \sup_x\sum_{X\ni x}e^{\kappa\tau(X)}
 \|a(X)\|_{\rm coeff}.
 \label{eq:v67-anchored-norm}
\end{equation}
Let \(z_{\rm core}\) be the anchored strong analytic norm of the already
renormalized local SM--Einstein interaction, excluding all of the following:
relative determinants, the V66 selector, the V54 scale-reset contact, and
far-field contours.  Let \(z_{\rm contact}\) denote the noncancelled,
nonprojected part of the exact V54 contact in the same norm.

For every determinant sector \(s\), retain the V53 quantities
\[
 \eta_{s,\kappa}=\|K_s\|_\kappa,
 \qquad
 d_{s,\kappa}
 =
 |\epsilon_s|\nu_s
 \frac{\eta_{s,\kappa}}{1-\eta_{s,\kappa}},
\]
and put
\begin{equation}
 d_\kappa=\sum_s d_{s,\kappa}.
 \label{eq:v67-determinant-stock}
\end{equation}
This includes a chart Berezinian only when that Berezinian occurs in the V53
sector ledger.  The same object is not reinserted through its score.

Choose the V66 convention
\[
 g_y=1+(g_y-1)
\]
rather than absorption of \(q_y\) into the exponent, and set
\begin{equation}
 s_\kappa=e^\epsilon-1,
 \qquad
 \epsilon=\sup_y\|Q_y/R_y^2\|_{{\cal A}_y}.
 \label{eq:v67-selector-stock}
\end{equation}
There is exactly one selector stock.  Using \(q_y\) as an exponent vertex
instead would replace \(s_\kappa\) by its norm \(\epsilon\), not add a second
term.

Define the forest input
\begin{equation}
 z_{\rm for}
 :=
 z_{\rm core}+z_{\rm contact}+d_\kappa+s_\kappa.
 \label{eq:v67-forest-input}
\end{equation}

\begin{proposition}[Disjoint-owner activity bound]
\label{prop:v67-disjoint-owner}
Suppose the four families in \eqref{eq:v67-forest-input} are expressed in
the same anchored support--grade--analytic norm and each polymer coefficient
is assigned to exactly one family.  Then their sum has norm at most
\(z_{\rm for}\).  If the fixed labelled-tree estimate has edge stock \(K_T\),
the connected forest series converges absolutely whenever
\begin{equation}
 q_{\rm for}:=eK_Tz_{\rm for}<1.
 \label{eq:v67-forest-gate}
\end{equation}
Its weak-norm connected sum is bounded by
\begin{equation}
 z_{\rm for}
 +
 \frac{e^2z_{\rm for}^2K_T}{1-q_{\rm for}}.
 \label{eq:v67-forest-sum}
\end{equation}
\end{proposition}

\begin{proof}
The first assertion is the triangle inequality in
\eqref{eq:v67-anchored-norm}.  The V53 sectorwise theorem bounds the
determinant family by \(d_\kappa\), and V66 bounds the selected local family
by \(s_\kappa\).  The two remaining bounds hold by definition of their
stocks.  The hypotheses of the V21 absolute connected-tree theorem then
hold with \(z=z_{\rm for}\); that theorem gives
\eqref{eq:v67-forest-gate}--\eqref{eq:v67-forest-sum}.
\end{proof}

This is a conditional composition theorem.  It does not assert that
\(z_{\rm core}\), \(z_{\rm contact}\), or \(K_T\) already have
scale-uniform numerical values.

\subsection{Bad-mark row}

Let \(p\) be the uniform V66 conditional-moment stock and let \(K_{\rm bad}\)
bound every additional local factor carried by a soft-bad contour.  The
independent bad-mark gate is
\begin{equation}
 q_{\rm bad}
 :=
 C_{\rm ent}e^\kappa K_{\rm bad}p<1.
 \label{eq:v67-bad-gate}
\end{equation}
It gives
\begin{equation}
 \sup_{y_0}
 \sum_{\substack{Y\ni y_0\\Y\ {\rm connected}}}
 e^{\kappa|Y|}K_{\rm bad}^{|Y|}
 {\mathbb E}\prod_{y\in Y}b_y
 \leq
 \frac{e^\kappa K_{\rm bad}p}{1-q_{\rm bad}}.
 \label{eq:v67-bad-sum}
\end{equation}
This row is not included in \(z_{\rm for}\): it controls a marked
measure-theoretic subseries.  If a construction instead converts each bad
mark into an ordinary analytic polymer vertex, it may move the corresponding
bound into \(z_{\rm for}\), but it must then delete the row
\eqref{eq:v67-bad-sum} for that same contribution.

Far-field events use the V66 exponential-moment row
\[
 M_y e^{-\theta_yA_yR_y^2}.
\]
They remain separate from both \(s_\kappa\) and \(p\).

\subsection{Score and compensation rows}

The V59 score stocks are derivative bounds for the chart density:
\[
 \|D_r\log j\|_{\rm anc}\leq\Sigma_0,
 \qquad
 \|D_wD_r\log j\|_{\rm anc}\leq\Sigma_1.
\]
They are not new logarithmic determinant activities.  The scalar activity
\(\log j\), when present, is already in the appropriate V53 determinant
sector.

Introduce
\begin{align}
 Q_0&=\|q_2\|_{\kappa,C_1R,\varrho_w},
 &
 Q_1&=\|D_wq_2\|_{\kappa,C_1R,\varrho_w},
 \label{eq:v67-Q-stocks}\\
 A_0&=
 \frac{\|S-S(0,w)\|_{\kappa,C_1R,\varrho_s}}
 {\varrho_s-\varrho_w},
 &
 A_1&=
 \frac{\|D_w(S-S(0,w))\|_{\kappa,C_1R,\varrho_s}}
 {\varrho_s-\varrho_w}.
 \label{eq:v67-action-stocks}
\end{align}
The V57 compensation bound with the V59 score insertion is
\begin{equation}
 B_0
 :=
 C_{\rm loc}\varrho_wQ_0(A_0+\Sigma_0).
 \label{eq:v67-B-zero}
\end{equation}

\begin{proposition}[One retained derivative of the compensation source]
\label{prop:v67-B-one}
Assume the local Euler inverse and product map commute with the retained
derivative in the fixed aligned chart and obey the same constant
\(C_{\rm loc}\).  Then
\begin{equation}
 \|D_w{\cal R}_4\|_{\kappa,R,\varrho_w}
 \leq B_1,
 \label{eq:v67-B-one-bound}
\end{equation}
where
\begin{equation}
 B_1
 :=
 C_{\rm loc}\varrho_w
 \left\{
 Q_1(A_0+\Sigma_0)
 +
 Q_0(A_1+\Sigma_1)
 \right\}.
 \label{eq:v67-B-one}
\end{equation}
\end{proposition}

\begin{proof}
The V56 Euler homotopy gives
\[
 \|W_3\|\leq\varrho_wQ_0,
 \qquad
 \|D_wW_3\|\leq\varrho_wQ_1.
\]
The compensation remainder is the sum of the action derivative paired with
\(W_3\) and the score paired with \(W_3\).  Differentiate each pairing.
Terms in which the derivative hits \(W_3\) are bounded by
\[
 C_{\rm loc}\varrho_wQ_1(A_0+\Sigma_0).
\]
Terms in which it hits the action derivative or score are bounded by
\[
 C_{\rm loc}\varrho_wQ_0(A_1+\Sigma_1).
\]
The analytic Cauchy loss in the action terms is exactly the denominator
already included in \(A_0,A_1\).  Adding the two classes proves
\eqref{eq:v67-B-one-bound}.
\end{proof}

If the horizontal derivative also changes the analytic radii, projection, or
Euler inverse, their connection terms must be added to \(B_1\).  The
proposition applies to the fixed aligned chart specified in its hypothesis.

The high-grade row is
\begin{equation}
 \lambda_{\rm tail}
 =
 L^{-N}+2C_Qr+\epsilon_N<1,
 \qquad
 C_PB_0
 \leq
 (1-\lambda_{\rm tail})r_{\rm res}.
 \label{eq:v67-tail-gate}
\end{equation}
By V57, this gives a compensated fixed-point displacement at most
\[
 \frac{C_PB_0}{1-\lambda_{\rm tail}}.
\]
The derivative stock \(B_1\) is retained for differentiability of the shot
and for the next scale-reset estimate; it is not added to the scalar
self-map radius in \eqref{eq:v67-tail-gate}.

\subsection{Pivot and locality prerequisite}

All stocks above are downstream of the V63 four-stage inverse.  Before
testing them, one must certify
\begin{align}
 \|L^{-1}\|_\kappa&<\infty,
 &
 \|\widehat A^{-1}\|_\kappa&<\infty,
 &
 \|D^{-1}\|_\kappa&<\infty,
 &
 \|\widehat S^{-1}\|_\kappa&<\infty,
 \label{eq:v67-four-pivots}\\
 \|A^{-1}\|_\kappa
 \|E_\sigma\|_\kappa
 \|L^{-1}\|_\kappa
 \|B_u\|_\kappa&<1.
 \label{eq:v67-hatted-A-gate}
\end{align}
The flat scalar stock may be supplied by V64 only in its stated chamber and
after the curved perturbation satisfies
\[
 \|L_0^{-1}{\cal V}\|_\kappa<1.
\]
These inverse bounds enter \(K_T\), the relative defects
\(\eta_{s,\kappa}\), and the score polynomials \(\Sigma_0,\Sigma_1\).
Testing a downstream gate with an uncertified inverse would be circular.

\subsection{The master card}

\begin{theorem}[Finite-volume budget implication]
\label{thm:v67-master-card}
At one regulator scale, assume:
\begin{enumerate}
\item the four pivot and locality bounds
      \eqref{eq:v67-four-pivots}--\eqref{eq:v67-hatted-A-gate};
\item \(\eta_{s,\kappa}<1\) in every determinant sector, with one
      determinant-line chart and its declared transition data;
\item the disjoint-owner forest estimate and
      \(q_{\rm for}<1\);
\item a positive normalized law satisfying the V66 uniform conditional
      moments and \(q_{\rm bad}<1\);
\item the far-contour moment estimate required by the chosen large-field
      layer; and
\item the complete covariant projection and tail gate
      \eqref{eq:v67-tail-gate}.
\end{enumerate}
Then the finite-volume connected good/soft-bad expansion is absolutely
summable in its weakened anchored norm, every declared relative determinant
is locally summable, and the compensated high-grade RG map has a unique
fixed point in the enlarged tail ball.  All three conclusions have no
factor proportional to the total number of blocks.
\end{theorem}

\begin{proof}
Item two and the V53 anchored theorem give \(d_\kappa\).  V66 gives
\(s_\kappa\) and the soft-bad product estimate.  Item three and Proposition
\ref{prop:v67-disjoint-owner} sum the analytic connected forest.  Item four
and \eqref{eq:v67-bad-sum} sum the soft-bad contours.  Item five handles the
separate far-field contours.  Finally item six and the V57 contraction
theorem produce the unique high-grade fixed point.  Every bound is anchored
at one block or is a contraction bound, so none introduces a volume factor.
The owner conventions above ensure that these applications use disjoint
contributions.
\end{proof}

\begin{ledger}[No-double-counting rules]
\label{led:v67-no-double-counting}
\begin{enumerate}
\item Charge a relative chart determinant through \(d_\kappa\); use
      \(\Sigma_0,\Sigma_1\) only when its derivatives enter compensation.
\item Charge the soft selector through either \(g-1\) or \(q\), once.
\item Charge a soft-bad family through either the marked expectation row or
      an ordinary analytic activity row, once.
\item Insert \({\cal P}_{<N}{\cal R}_4\) into running couplings and
      \((I-{\cal P}_{<N}){\cal R}_4\) into the tail map.  Do not put the
      unsplit remainder in both.
\item Keep heat-shell, overlap, interface-relative, Faddeev--Popov,
      coarea, and physical-fermion determinants in their distinct V52--V53
      sectors.
\end{enumerate}
\end{ledger}

\subsection{Acquired and open rows}

At fixed regulator, the following implications are now proved: the V53
formula for \(d_\kappa\), the V66 formula for \(s_\kappa\) and the bad-mark
sum, the V59 polynomials \(\Sigma_0,\Sigma_1\), the compensation stocks
\(B_0,B_1\), and the V57 tail implication.  The following numerical inputs
are not yet acquired uniformly in scale:
\[
 z_{\rm core},\quad z_{\rm contact},\quad K_T,\quad
 p,\quad K_{\rm bad},\quad M_y,\quad
 \eta_{s,\kappa},\quad
 1-\lambda_{\rm tail}.
\]
Most seriously, no positive Einstein reference law has been constructed to
supply \(p\) and the far-tail moment.  The candidate bulk scalar reduction
has also not passed the V62 physical-equivalence row.

\subsection{Next acquisition}

The master card shows that the first upstream bottleneck is measure
positivity, not another forest estimate.  V68 will test whether a
reflection-compatible positive conditional law can be obtained from the
canonical constrained Einstein system without imposing the unproved
four-dimensional bulk trace equation.  It will begin with the linearized
ADM conformal decomposition and identify exactly which scalar modes are
constraints, which are gauge, and which remain physical or unstable.

\begin{assessment}[Exact status after V67]
\label{ass:v67-status}
V67 proves a volume-free finite-regulator implication with disjoint owners
and adds the first retained derivative bound for the compensation source.
It also localizes the remaining bottleneck at the positive gravitational
conditional law and the physical status of the bulk scalar reduction.
Uniform hypotheses of the master card, Osterwalder--Schrader positivity,
projective consistency, regulator removal, and the full SM--Einstein
continuum remain open.
\end{assessment}
\section{Linearized canonical reduction and the positive TT reference}
\label{sec:v10v68}

V23--V26 solved and glued a nonlinear CMC conformal constraint in a
restricted sign chamber.  V62--V64 then considered a different,
four-dimensional bulk trace equation and recorded its unresolved physical
status.  The present note does not repeat either construction.  It asks a
more upstream question: does the canonical Einstein constraint system
itself remove the conformal instability at the linearized level, without
inserting the bulk trace equation?

The answer is yes for every nonzero mode about flat, time-symmetric vacuum
data.  The resulting transverse-traceless Gaussian is positive and
reflection positive.  A second result identifies the remaining obstruction:
the natural Gaussian innovation coordinates need not be spatially local, so
positivity alone does not fill the V66 polymer row.

\subsection{Fourier--York decomposition}

Let the spatial slice be the flat torus \({\mathbb T}^3\), and linearize the
ADM phase space at
\[
 g_{ij}=\delta_{ij},\qquad \pi^{ij}=0.
\]
Write \(h_{ij}\) for the metric perturbation and \(\pi^{ij}\) for its
conjugate momentum.  On every nonzero Fourier mode, use the orthogonal
decompositions
\begin{align}
 h_{ij}
 ={}&
 h_{ij}^{\rm TT}
 +2\partial_{(i}v^{T}_{j)}
 +\left(\partial_i\partial_j-\frac13\delta_{ij}\Delta\right)s
 +\frac13\delta_{ij}\tau,
 \label{eq:v68-h-York}\\
 \pi^{ij}
 ={}&
 \pi_{\rm TT}^{ij}
 +2\partial^{(i}p_T^{j)}
 +\left(\partial^i\partial^j-\frac13\delta^{ij}\Delta\right)\rho
 +\frac13\delta^{ij}\varpi.
 \label{eq:v68-pi-York}
\end{align}
Here
\[
 \partial^ih_{ij}^{\rm TT}=0,\quad
 \delta^{ij}h_{ij}^{\rm TT}=0,\quad
 \partial_iv_T^i=\partial_ip_T^i=0,
\]
and similarly for \(\pi_{\rm TT}\).  The variables \(\tau,\varpi\) are the
two traces.  Inversion of this decomposition uses \(\Delta^{-1}\), so the
zero mode is excluded at this stage.

The vacuum linearized constraints are
\begin{align}
 {\cal H}^{(1)}
 &=
 \partial_i\partial_jh_{ij}-\Delta\tau,
 \label{eq:v68-H-linear}\\
 {\cal H}^{(1)}_i
 &=
 -2\partial_j\pi^j{}_i.
 \label{eq:v68-M-linear}
\end{align}

\begin{lemma}[Constraint equations in York coordinates]
\label{lem:v68-York-constraints}
On a nonzero Fourier mode,
\begin{align}
 {\cal H}^{(1)}
 &=
 \frac23\Delta(\Delta s-\tau),
 \label{eq:v68-H-York}\\
 {\cal H}^{(1)}_i
 &=
 -2\left\{
 \Delta p_i^T+
 \partial_i\left(\frac23\Delta\rho+\frac13\varpi\right)
 \right\}.
 \label{eq:v68-M-York}
\end{align}
Consequently the constraint surface satisfies
\begin{equation}
 \tau=\Delta s,\qquad
 p_i^T=0,\qquad
 \varpi=-2\Delta\rho.
 \label{eq:v68-constraint-solutions}
\end{equation}
\end{lemma}

\begin{proof}
The TT and transverse-vector terms in
\(\partial_i\partial_jh_{ij}\) vanish.  Applying two divergences to the
traceless scalar term gives
\[
 \left(1-\frac13\right)\Delta^2s=\frac23\Delta^2s,
\]
while the trace term gives \(\Delta\tau/3\).  Subtracting
\(\Delta\tau\) proves \eqref{eq:v68-H-York}.  One divergence of
\eqref{eq:v68-pi-York} gives
\[
 \Delta p_i^T+\frac23\partial_i\Delta\rho
 +\frac13\partial_i\varpi.
\]
The transverse and longitudinal parts vanish separately on a nonzero mode,
which proves the remaining assertions.
\end{proof}

\subsection{Gauge flows and complete scalar removal}

A spatial infinitesimal diffeomorphism has
\[
 \xi_i=\xi_i^T+\partial_i\alpha,\qquad \partial_i\xi_T^i=0,
\]
and acts at the base point by
\begin{equation}
 \delta_\xi h_{ij}=\partial_i\xi_j+\partial_j\xi_i.
 \label{eq:v68-spatial-gauge}
\end{equation}
In the variables \eqref{eq:v68-h-York},
\begin{equation}
 \delta v_i^T=\xi_i^T,\qquad
 \delta s=2\alpha,\qquad
 \delta\tau=2\Delta\alpha.
 \label{eq:v68-spatial-York}
\end{equation}
The Hamiltonian constraint smeared with a lapse perturbation \(n\) generates
\begin{equation}
 \delta_n\pi^{ij}
 =
 -\partial^i\partial^jn+\delta^{ij}\Delta n,
 \label{eq:v68-normal-gauge}
\end{equation}
up to an overall convention for the sign of \(n\).  With the convention in
\eqref{eq:v68-normal-gauge},
\begin{equation}
 \delta_n\rho=-n,\qquad
 \delta_n\varpi=2\Delta n.
 \label{eq:v68-normal-York}
\end{equation}

\begin{theorem}[Nonzero-mode canonical reduction]
\label{thm:v68-canonical-reduction}
For every nonzero Fourier mode, impose the constraints
\({\cal H}^{(1)}={\cal H}^{(1)}_i=0\) and quotient their gauge flows.  Every
orbit has a unique representative satisfying
\begin{equation}
 v_i^T=0,\qquad s=0,\qquad\rho=0.
 \label{eq:v68-gauge-section}
\end{equation}
On this representative,
\begin{equation}
 \tau=0,\qquad p_i^T=0,\qquad\varpi=0.
 \label{eq:v68-removed-components}
\end{equation}
Thus the reduced nonzero-mode phase space consists exactly of
\((h^{\rm TT},\pi_{\rm TT})\), with two configuration polarizations and
their two conjugate momenta.
\end{theorem}

\begin{proof}
Choose \(\xi_i^T=-v_i^T\) and \(\alpha=-s/2\) in
\eqref{eq:v68-spatial-York}; this sets \(v_i^T=s=0\).  The Hamiltonian
constraint in \eqref{eq:v68-constraint-solutions} then gives \(\tau=0\).
Choose \(n=\rho\) in \eqref{eq:v68-normal-York}; this sets \(\rho=0\), and
the scalar momentum constraint gives \(\varpi=0\).  Its transverse part
already gives \(p_i^T=0\).

For uniqueness, a gauge transformation preserving
\eqref{eq:v68-gauge-section} has
\(\xi_i^T=0\), \(\alpha=0\), and \(n=0\) on a nonzero mode.  The TT
components are orthogonal to all three gauge directions and do not occur in
the linearized constraints.  A symmetric tensor in three dimensions has
six components; the decomposition allocates two TT, two transverse-vector,
and two scalar components, proving the polarization count.
\end{proof}

The theorem identifies the scalar roles precisely.  The invariant
configuration combination \(\tau-\Delta s\) is killed by the Hamiltonian
constraint; \(s\) is removed by the scalar spatial diffeomorphism.  The
invariant momentum combination \(\varpi+2\Delta\rho\) is killed by the
momentum constraint; \(\rho\) is removed by the normal gauge flow.  No
nonzero scalar phase-space variable remains to be integrated on a real
Euclidean contour.

The entire zero Fourier sector is retained.  It contains the volume and
torus-shape modes as well as their momenta; constant diffeomorphisms are
stabilizers rather than invertible gauge directions.  Hiding this sector in
\(\Delta^{-1}\) would violate the V24 and V64 zero-mode ledgers.

\subsection{Positive reduced quadratic Hamiltonian}

\begin{proposition}[TT coercivity]
\label{prop:v68-TT-coercivity}
The quadratic ADM Hamiltonian restricted to the reduced nonzero-mode phase
space has the form
\begin{equation}
 H_{\rm red}^{(2)}
 =
 \sum_{k\neq0}
 \left\{
 c_\pi|\widehat\pi_{\rm TT}(k)|^2
 +
 c_h|k|^2|\widehat h_{\rm TT}(k)|^2
 \right\},
 \qquad c_\pi,c_h>0.
 \label{eq:v68-positive-Hamiltonian}
\end{equation}
After a canonical rescaling of each polarization it is a sum of harmonic
oscillator Hamiltonians.
\end{proposition}

\begin{proof}
In the ADM Hamiltonian, the momentum quadratic form is proportional to
\[
 \pi^{ij}\pi_{ij}-\frac12\varpi^2.
\]
On the reduced section \(\varpi=0\), it is
\(|\pi_{\rm TT}|^2\) with a positive coefficient.  The second variation of
the spatial curvature term, after integration by parts, is a positive
multiple of
\[
 \int_{{\mathbb T}^3}|\partial_kh_{ij}^{\rm TT}|^2.
\]
All divergence and trace terms vanish on TT tensors.  Parseval's identity
gives \eqref{eq:v68-positive-Hamiltonian}.  Positive linear rescalings of
\(h_{\rm TT}\) and \(\pi_{\rm TT}\) put both coefficients into the standard
oscillator normalization.
\end{proof}

This positivity is a statement about the canonically reduced Lorentzian
phase space and its oscillator Wick rotation.  It does not make the
unreduced real Euclidean Einstein action coercive.

Integrating the reduced Gaussian momenta gives, on Euclidean time,
\begin{equation}
 S_{\rm TT}^{(2)}
 =
 \frac12
 \sum_{\lambda=1}^2\sum_{k\neq0}
 \int_{\mathbb R}
 \left\{
 |\partial_tq_{\lambda,k}|^2
 +\omega_k^2|q_{\lambda,k}|^2
 \right\}dt,
 \qquad
 \omega_k=c|k|>0.
 \label{eq:v68-TT-Euclidean}
\end{equation}

\begin{theorem}[Reflection positivity of the reduced Gaussian]
\label{thm:v68-TT-OS}
The centered Gaussian with action \eqref{eq:v68-TT-Euclidean} is
reflection positive under \(t\mapsto-t\).  More explicitly, its covariance
in one polarization and spatial mode is
\begin{equation}
 C_k(t-s)=\frac{e^{-\omega_k|t-s|}}{2\omega_k},
 \label{eq:v68-oscillator-covariance}
\end{equation}
and, for every test function \(f_k\) supported in \(t\geq0\),
\begin{equation}
 \int_0^\infty\!\!\int_0^\infty
 \overline{f_k(t)}C_k(-t-s)f_k(s)\,dt\,ds
 =
 \frac1{2\omega_k}
 \left|
  \int_0^\infty e^{-\omega_kt}f_k(t)\,dt
 \right|^2
 \geq0.
 \label{eq:v68-OS-square}
\end{equation}
\end{theorem}

\begin{proof}
The inverse of
\(-\partial_t^2+\omega_k^2\) on the line is
\eqref{eq:v68-oscillator-covariance}.  For \(t,s\geq0\),
\[
 C_k(-t-s)=\frac{e^{-\omega_k(t+s)}}{2\omega_k}.
\]
Substitution factors the double integral into the square in
\eqref{eq:v68-OS-square}.  Summing nonnegative squares over modes and
polarizations proves reflection positivity.
\end{proof}

The strict inequality \(\omega_k>0\) explains again why the spatial zero
mode must be retained outside this Gaussian.

\subsection{Gaussian conditional moments and innovations}

The preceding theorem gives a positive finite-regulator reference after
spatial and temporal mode cutoffs.  The V66 card also requires a
block-local conditional-moment estimate.  The following finite-dimensional
calculation shows what a Gaussian can supply.

\begin{proposition}[One-block Gaussian innovation]
\label{prop:v68-Gaussian-innovation}
Let \(x=(x_y,x_{-y})\) have a centered Gaussian density proportional to
\[
 \exp\{-\tfrac12x^TAx\},
 \qquad A=A^T>0.
\]
Write the corresponding blocks of \(A\) as \(A_{yy}\) and \(A_{y,-y}\).
Conditioned on \(x_{-y}\),
\begin{equation}
 r_y
 :=
 x_y+A_{yy}^{-1}A_{y,-y}x_{-y}
 \label{eq:v68-local-innovation}
\end{equation}
is centered Gaussian with covariance \(A_{yy}^{-1}\).  If
\(d_y=\dim x_y\) and
\begin{equation}
 Q_y=r_y^TA_{yy}r_y,
 \label{eq:v68-innovation-Q}
\end{equation}
then
\begin{align}
 {\mathbb E}(Q_y\mid x_{-y})&=d_y,
 \label{eq:v68-innovation-moment}\\
 {\mathbb E}(e^{\theta Q_y}\mid x_{-y})
 &=(1-2\theta)^{-d_y/2},
 \qquad0\leq\theta<\frac12.
 \label{eq:v68-innovation-exponential}
\end{align}
Both bounds are uniform in the exterior value.
\end{proposition}

\begin{proof}
Complete the square:
\begin{align*}
 x^TAx
 ={}&
 r_y^TA_{yy}r_y\\
 &+
 x_{-y}^T
 \left(
 A_{-y,-y}-A_{-y,y}A_{yy}^{-1}A_{y,-y}
 \right)x_{-y}.
\end{align*}
The conditional density of \(r_y\) is therefore the Gaussian with precision
\(A_{yy}\).  The change of variables
\(z=A_{yy}^{1/2}r_y\) makes \(Q_y=|z|^2\) for a standard
\(d_y\)-dimensional Gaussian.  Its first moment is \(d_y\), and its
factorized Gaussian integral gives
\((1-2\theta)^{-d_y/2}\).
\end{proof}

For one conditioned block, the V66 stocks would be
\begin{equation}
 p_y=\frac{d_y}{R_y^2},
 \qquad
 M_y(\theta)=(1-2\theta)^{-d_y/2}.
 \label{eq:v68-formal-V66-stocks}
\end{equation}
This is a genuine positive-law estimate and does not use the Euclidean
Einstein conformal action.

\subsection{Why the innovation does not yet close the polymer card}

Proposition \ref{prop:v68-Gaussian-innovation} conditions on the original
exterior variables.  If one defines every \(r_y\) simultaneously by
\eqref{eq:v68-local-innovation}, the factors attached to neighboring blocks
overlap.  In particular, \(b_z(r_z)\) can depend on \(x_y\), so it is not
measurable with respect to the exterior sigma algebra used when integrating
\(x_y\).  The iterative proof of
\eqref{eq:v66-correlated-product} then does not apply verbatim.

A triangular Gaussian factorization repairs independence.  Choose an
ordering of all coordinates and factor
\begin{equation}
 A=L^TDL
 \label{eq:v68-LDL}
\end{equation}
with \(D>0\) block diagonal and \(L\) unit lower triangular.  Then
\begin{equation}
 r=Lx
 \label{eq:v68-triangular-innovation}
\end{equation}
has independent Gaussian blocks with precisions given by \(D\), so V66
applies exactly to \(Q_y=r_y^TD_{yy}r_y\).

\begin{firewall}[Gaussian positivity versus spatial locality]
\label{fire:v68-innovation-locality}
Sparse positive \(A\) need not have a uniformly sparse factor \(L\);
elimination can create fill-in, and its range can grow with the volume.
For the TT Gaussian there is an additional source of nonlocality: the
Fourier TT projector contains inverse powers of the spatial Laplacian.
Therefore \eqref{eq:v68-triangular-innovation} proves a product positive
reference in innovation coordinates, but does not prove the uniformly
local owner structure required by the V53, V66, and V67 anchored norms.
\end{firewall}

This is the precise new bottleneck.  A Fourier product measure is excellent
for positivity and poor for spatial polymers.  A raw local metric Gaussian
is spatially local but includes gauge and constrained scalar directions.
The next construction must preserve locality while enforcing the canonical
constraints.

\subsection{Matter and background limitations}

Theorem \ref{thm:v68-canonical-reduction} uses flat time-symmetric vacuum
data.  If the background matter fields or extrinsic curvature have a
nonzero first variation, the scalar and vector constraints acquire source
and mixing blocks.  The nonzero gravitational scalar is then determined by
the coupled constraint rather than set to zero.  Its elimination requires
the V44--V49 Schur chamber, with the matter variables retained in their
actual sector.  The vacuum theorem does not justify deleting the Higgs
radial mode, charge densities, or scalar matter excitations.

Likewise, nonlinear constraint reduction can have multiple charts, Gribov
stabilizers, Yamabe boundaries, and zero-mode crossings.  The V23--V26
global sign chamber remains the only nonlinear result presently proved in
these notes.

\begin{ledger}[Effect on the V67 master card]
\label{led:v68-master-effect}
\begin{enumerate}
\item The linearized nonzero-mode gravitational law is positive and
      reflection positive after canonical reduction.
\item Gaussian innovations give explicit conditional and exponential
      moments, \eqref{eq:v68-formal-V66-stocks}.
\item No scale-uniform spatially local simultaneous innovation chart has
      been proved; hence \(p\) is not yet entered as an acquired local
      polymer stock.
\item The entire zero-mode sector remains in the macroscopic ledger.
\item The four-dimensional bulk trace equation is unnecessary for the
      linearized vacuum reduction, but V68 does not prove a nonlinear
      replacement for it.
\end{enumerate}
\end{ledger}

\subsection{Next acquisition}

V69 will seek a local representation of the constrained TT Gaussian.
The natural candidates are a reflection-compatible constrained metric
density with lapse and shift retained as Lagrange multipliers, or local
tensor potentials whose double-curl image is transverse and whose kernel is
kept as a gauge complex.  The test will be strict: the representation must
have a finite-range or exponentially local precision, a positive physical
quotient, explicit zero-mode owners, and a conditional-moment estimate that
survives spatial reblocking.  If neither candidate passes, the programme
must use a multiscale rather than one-step Gaussian innovation.

\begin{assessment}[Exact status after V68]
\label{ass:v68-status}
V68 proves complete removal of all nonzero scalar and vector phase-space
modes in linearized vacuum ADM gravity, positivity of the reduced TT
Hamiltonian, reflection positivity of its Euclidean Gaussian, and exact
Gaussian innovation moments.  It also proves that these facts do not by
themselves give a spatially local V66 reference.  Nonlinear constrained
positivity, local innovations, coupled matter backgrounds, zero-mode
dynamics, uniform RG bounds, and the full SM--Einstein continuum remain
open.
\end{assessment}
\section{Local constraint penalties and colored Gaussian innovations}
\label{sec:v10v69}

V68 produced a positive transverse-traceless Gaussian but found that a
global triangular innovation can destroy spatial locality.  This section
replaces the global factorization by a bounded coloring of the finite-range
precision graph.  The price is explicit: a one-block bad primitive \(p\)
becomes \(p^{1/q}\), where \(q\) is the number of colors.  Since \(q\) is
bounded by the local degree plus one, the estimate is uniform in volume.

A second theorem constructs a local positive penalty for the linearized
trace and divergence constraints and proves convergence to the TT Gaussian.
This closes the finite-regulator positive-reference row for the linearized
vacuum sector.  The penalty limit in the strong analytic norm remains a
separate scale problem.

\subsection{Finite-range Gaussian graph}

Let \(\Lambda\) be a finite block set and let
\[
 x=(x_y)_{y\in\Lambda},\qquad x_y\in{\mathbb R}^{d_y},
\]
have centered Gaussian law
\begin{equation}
 d\mu_A(x)
 =
 Z_A^{-1}\exp\{-\tfrac12 x^TAx\}\,dx,
 \qquad A=A^T>0.
 \label{eq:v69-Gaussian-law}
\end{equation}
Define the precision graph \(G_A\) by joining \(y\neq z\) exactly when
\(A_{yz}\neq0\).  Assume its maximum degree is at most \(\Delta_A\).  Choose
a proper coloring
\begin{equation}
 \Lambda=C_1\sqcup\cdots\sqcup C_q,
 \qquad q\leq\Delta_A+1.
 \label{eq:v69-coloring}
\end{equation}
No two vertices of one \(C_c\) are joined.

For \(y\in C_c\), define the color-conditional innovation
\begin{equation}
 r_y
 :=
 x_y+A_{yy}^{-1}
 \sum_{z\notin C_c}A_{yz}x_z.
 \label{eq:v69-color-innovation}
\end{equation}
Because the coloring is proper, the sum is unchanged if
\(z\notin C_c\) is replaced by \(z\neq y\).  Set
\begin{equation}
 Q_y=r_y^TA_{yy}r_y,\qquad
 g_y=e^{-Q_y/R_y^2},\qquad
 b_y=1-g_y.
 \label{eq:v69-color-selector}
\end{equation}

\begin{lemma}[Conditional factorization within one color]
\label{lem:v69-color-factorization}
Conditioned on \(x_{C_c^c}\), the variables
\((r_y)_{y\in C_c}\) are independent centered Gaussians with covariance
\(A_{yy}^{-1}\).  Consequently
\begin{align}
 {\mathbb E}(Q_y\mid x_{C_c^c})&=d_y,
 \label{eq:v69-color-first-moment}\\
 {\mathbb E}(e^{\theta Q_y}\mid x_{C_c^c})
 &=(1-2\theta)^{-d_y/2},
 \qquad0\leq\theta<\frac12.
 \label{eq:v69-color-exp-moment}
\end{align}
\end{lemma}

\begin{proof}
Since \(A_{yz}=0\) for distinct \(y,z\in C_c\), the terms in the exponent
that depend on \(x_{C_c}\) are
\[
 \sum_{y\in C_c}
 \left\{
 \frac12x_y^TA_{yy}x_y
 +
 x_y^T\sum_{z\notin C_c}A_{yz}x_z
 \right\}.
\]
Complete the square separately for each \(y\).  The result is
\[
 \frac12\sum_{y\in C_c}r_y^TA_{yy}r_y
\]
plus a term depending only on \(x_{C_c^c}\).  Thus the conditional density
factorizes.  The moment identities are the standard Gaussian calculation
already used in Proposition \ref{prop:v68-Gaussian-innovation}.
\end{proof}

Put
\begin{equation}
 \bar p_y:=\min\{1,d_y/R_y^2\},
 \qquad
 \bar p:=\sup_y\bar p_y.
 \label{eq:v69-pbar}
\end{equation}

\begin{theorem}[Colored soft-bad product bound]
\label{thm:v69-colored-product}
For every \(Y\subset\Lambda\),
\begin{equation}
 {\mathbb E}_{\mu_A}\prod_{y\in Y}b_y
 \leq
 \prod_{y\in Y}\bar p_y^{\,1/q}
 \leq
 \bar p^{\,|Y|/q}.
 \label{eq:v69-colored-product}
\end{equation}
The exact subset identity
\begin{equation}
 1=
 \sum_{Y\subset\Lambda}
 \prod_{y\in Y}b_y
 \prod_{z\in\Lambda\setminus Y}g_z
 \label{eq:v69-exact-subset}
\end{equation}
holds pointwise.
\end{theorem}

\begin{proof}
For a fixed color, let
\[
 B_c=\prod_{y\in Y\cap C_c}b_y.
\]
Condition on \(x_{C_c^c}\).  Lemma
\ref{lem:v69-color-factorization}, the inequality
\(1-e^{-t}\leq t\), and \(0\leq b_y\leq1\) give
\[
 {\mathbb E}(B_c\mid x_{C_c^c})
 \leq
 \prod_{y\in Y\cap C_c}
 \min\{1,d_y/R_y^2\}.
\]
The right side is deterministic, so the same bound holds for
\({\mathbb E}B_c\).

Now
\[
 \prod_{y\in Y}b_y=\prod_{c=1}^qB_c.
\]
Generalized H\"older inequality with all exponents equal to \(q\) gives
\[
 {\mathbb E}\prod_{c=1}^qB_c
 \leq
 \prod_{c=1}^q({\mathbb E}B_c^q)^{1/q}.
\]
Since \(0\leq B_c\leq1\), one has \(B_c^q\leq B_c\).  Inserting the
one-color bounds proves \eqref{eq:v69-colored-product}.  The pointwise
relations \(g_y+b_y=1\) prove \eqref{eq:v69-exact-subset}.
\end{proof}

\begin{corollary}[Volume-free colored contour gate]
\label{cor:v69-colored-entropy}
Assume the connected-set count
\(N_n(y_0)\leq C_{\rm ent}^{\,n-1}\), and let every marked block carry an
additional factor at most \(K_{\rm bad}\).  If
\begin{equation}
 C_{\rm ent}e^\kappa K_{\rm bad}\bar p^{1/q}<1,
 \label{eq:v69-colored-gate}
\end{equation}
then
\begin{equation}
 \sup_{y_0}
 \sum_{\substack{Y\ni y_0\\Y\ {\rm connected}}}
 e^{\kappa|Y|}K_{\rm bad}^{|Y|}
 {\mathbb E}\prod_{y\in Y}b_y
 \leq
 \frac{e^\kappa K_{\rm bad}\bar p^{1/q}}
 {1-C_{\rm ent}e^\kappa K_{\rm bad}\bar p^{1/q}}.
 \label{eq:v69-colored-sum}
\end{equation}
\end{corollary}

\begin{proof}
Use Theorem \ref{thm:v69-colored-product}, group by
\(|Y|=n\), and sum the resulting geometric series.
\end{proof}

\subsection{Local analytic support}

If \(A\) has range \(r_A\), then \(r_y\) and \(Q_y\) in
\eqref{eq:v69-color-innovation}--\eqref{eq:v69-color-selector} depend only
on the closed \(r_A\)-neighborhood of \(y\).  Thus the selector enlarges a
one-block support by a fixed amount and does not introduce a volume-length
edge.

For a concrete analytic bound, suppose
\begin{equation}
 a_0:=\sup_y\|A_{yy}\|<\infty,
 \qquad
 m_A:=
 1+\sup_y\sum_{z\neq y}
 \|A_{yy}^{-1}A_{yz}\|<\infty.
 \label{eq:v69-row-stocks}
\end{equation}
On a coefficient-majorant polydisc with
\(\|x_z\|\leq\varrho\), there is a fixed fibre constant \(C_d\) such that
\begin{equation}
 \|Q_y/R_y^2\|_{\cal A}
 \leq
 \epsilon_{{\rm col},y}
 :=
 \frac{C_da_0m_A^2\varrho^2}{R_y^2}.
 \label{eq:v69-local-analytic-q}
\end{equation}
Indeed, \(\|r_y\|\leq m_A\varrho\), and the quadratic contraction with
\(A_{yy}\) costs at most \(C_da_0\).  The V66 Banach-algebra theorem then
gives
\begin{equation}
 \|g_y-1\|_{\cal A}
 \leq e^{\epsilon_{{\rm col},y}}-1.
 \label{eq:v69-local-selector-activity}
\end{equation}
Both the support radius and the norm bound are independent of
\(|\Lambda|\) when \(r_A,\Delta_A,a_0,m_A,d_y\) are uniform.

\subsection{A local positive TT penalty}

Let \({\cal H}_0\) be the finite-dimensional space of zero-mean symmetric
tensor fields on a periodic spatial lattice.  Let \(D\) be the chosen
finite-range lattice divergence and \(T\) the pointwise trace.  Define
\begin{equation}
 B:=D^*D+T^*T,
 \qquad
 {\cal H}_{\rm TT}:=\ker B=\ker D\cap\ker T.
 \label{eq:v69-B-operator}
\end{equation}
Assume that the lattice symbols of \(D\) have no unintended zero away from
zero momentum; any exceptional modes are otherwise added to the retained
ledger.  Let
\[
 K_0=-\Delta_{\rm lat}\otimes I
\]
on \({\cal H}_0\), and for \(\varepsilon>0\) put
\begin{equation}
 K_\varepsilon
 :=
 K_0+\varepsilon^{-1}B.
 \label{eq:v69-penalized-precision}
\end{equation}
Both operators are finite range and translation invariant, hence commute.
On the zero-mean space, \(K_0>0\), so
\(K_\varepsilon>0\).

\begin{theorem}[Positive local constraint penalty]
\label{thm:v69-penalty-limit}
Let \(P_{\rm TT}\) be the orthogonal projection onto
\({\cal H}_{\rm TT}\).  At fixed finite lattice,
\begin{equation}
 K_\varepsilon^{-1}
 \longrightarrow
 P_{\rm TT}K_0^{-1}P_{\rm TT}
 \quad\hbox{in operator norm as }\varepsilon\downarrow0.
 \label{eq:v69-covariance-limit}
\end{equation}
The normalized centered Gaussians with covariance
\(K_\varepsilon^{-1}\) converge weakly to the TT Gaussian with covariance
\(P_{\rm TT}K_0^{-1}P_{\rm TT}\), supported on
\({\cal H}_{\rm TT}\).
\end{theorem}

\begin{proof}
Since \(K_0\) and \(B\) are commuting self-adjoint matrices, they have a
common orthonormal eigenbasis.  On \(\ker B\), the inverse eigenvalue is
exactly that of \(K_0^{-1}\).  On \((\ker B)^\perp\), let
\(\beta_{\min}>0\) be the least positive eigenvalue of \(B\).  Then
\[
 \|(I-P_{\rm TT})K_\varepsilon^{-1}
   (I-P_{\rm TT})\|
 \leq
 \frac{\varepsilon}{\beta_{\min}}.
\]
This proves \eqref{eq:v69-covariance-limit}.  In finite dimension,
convergence of the covariance matrices gives pointwise convergence of
Gaussian characteristic functions.  The limiting covariance vanishes on
\((\ker B)^\perp\), so the limiting law is the TT Gaussian times a point
mass at zero in the orthogonal directions.
\end{proof}

The corresponding Euclidean-time action is
\begin{equation}
 S_\varepsilon(h)
 =
 \frac12\int_{\mathbb R}
 \left\{
 \|\partial_th\|^2+
 \langle h,K_\varepsilon h\rangle
 \right\}dt.
 \label{eq:v69-penalty-action}
\end{equation}

\begin{theorem}[Reflection positivity before and after the penalty limit]
\label{thm:v69-penalty-OS}
For every \(\varepsilon>0\), the Gaussian defined by
\eqref{eq:v69-penalty-action} is reflection positive in Euclidean time.
For a positive-time test function \(f\), its reflection form is
\begin{equation}
 \frac12
 \left\|
 \int_0^\infty
 K_\varepsilon^{-1/4}
 e^{-tK_\varepsilon^{1/2}}f(t)\,dt
 \right\|^2
 \geq0.
 \label{eq:v69-penalty-OS-square}
\end{equation}
On test functions with values in \({\cal H}_{\rm TT}\), this form is
independent of \(\varepsilon\) and equals the V68 TT reflection form.
\end{theorem}

\begin{proof}
The covariance is
\[
 C_\varepsilon(t-s)
 =
 \frac12K_\varepsilon^{-1/2}
 e^{-|t-s|K_\varepsilon^{1/2}}.
\]
For \(t,s\geq0\), insert \(C_\varepsilon(-t-s)\) into the reflected
quadratic form and factor the two commuting quarter powers.  This gives
\eqref{eq:v69-penalty-OS-square}.  On \({\cal H}_{\rm TT}\), \(B=0\) and
\(K_\varepsilon=K_0\).
\end{proof}

After discretizing Euclidean time by a reflection-positive nearest-neighbor
quadratic form, the full spacetime precision remains finite range and has a
uniformly bounded graph degree.  Therefore Theorem
\ref{thm:v69-colored-product} applies at every
\(\varepsilon>0\).  With a fixed local fibre dimension \(d\),
\begin{equation}
 \bar p_{\rm eff}
 =
 \min\{1,d/R^2\}^{1/q}
 \label{eq:v69-effective-p}
\end{equation}
is a volume-independent soft-bad primitive.

\subsection{The singular analytic-radius debit}

Although the probabilistic primitive
\eqref{eq:v69-effective-p} is independent of the eigenvalues of the
precision, the analytic stock in
\eqref{eq:v69-local-analytic-q} is not.  For
\eqref{eq:v69-penalized-precision},
\[
 a_0(\varepsilon)=O(\varepsilon^{-1})
\]
in any block touched by \(D^*D+T^*T\).  At fixed original-field analytic
radius \(\varrho\), the quantity
\(\epsilon_{{\rm col},y}\) can therefore diverge as
\(\varepsilon\downarrow0\).

A sufficient local window is
\begin{equation}
 \frac{d}{R^2}\ll1,
 \qquad
 \frac{a_0(\varepsilon)m_A(\varepsilon)^2\varrho_\varepsilon^2}
      {R^2}
 \leq c_{\rm an},
 \label{eq:v69-penalty-window}
\end{equation}
where \(c_{\rm an}\) is the allocated forest budget.  If
\(m_A(\varepsilon)=O(1)\), this requires
\(\varrho_\varepsilon=O(\sqrt{\varepsilon}\,R)\) in the stiff directions.
Shrinking the analytic radius of the physical TT directions would destroy
the intended continuum chart.  The needed norm must therefore distinguish
constraint-violating directions from physical directions without inserting
the nonlocal TT projector into every local owner.

\begin{firewall}[Reference regulator versus physical measure]
\label{fire:v69-penalty-physical}
The penalty family proves a local positive reference and its normalized
linear TT limit.  At finite \(\varepsilon\) it is not the Einstein
Faddeev--Popov density.  Its unnormalized determinant has a singular
constraint-sector factor as \(\varepsilon\downarrow0\), and interacting
observables can couple TT and stiff directions.  Physical equivalence
requires a uniform decoupling theorem plus the correct constraint, orbit,
and determinant owners; it is not implied by weak Gaussian convergence.
\end{firewall}

\begin{ledger}[Acquisition after V69]
\label{led:v69-acquisition}
At finite regulator and \(\varepsilon>0\):
\begin{enumerate}
\item the reference precision is positive, finite range, and reflection
      positive;
\item colored innovations are local and give the exact bound
      \eqref{eq:v69-colored-product};
\item the contour primitive \eqref{eq:v69-effective-p} is uniform in total
      volume;
\item the normalized free measure converges to the V68 TT Gaussian; and
\item the remaining singularity is isolated in the strong analytic radius
      and in the physical determinant/decoupling comparison.
\end{enumerate}
\end{ledger}

\subsection{Next acquisition}

V70 is the next compulsory companion audit.  It will inspect the newest
Yang--Mills and Navier--Stokes notes for a colored or multiscale innovation,
a stiff-direction analytic norm, or an exact decoupling estimate.  The
post-audit note will then construct an anisotropic local analytic norm for
the penalty complex or replace the penalty by a multiscale conditional
factorization.

\begin{assessment}[Exact status after V69]
\label{ass:v69-status}
V69 proves a local colored-innovation theorem, a volume-free Gaussian
bad-contour gate, a positive finite-range constraint penalty, its
reflection positivity, and convergence to the linear TT Gaussian.  It does
not prove a uniform strong analytic penalty limit, determinant
renormalization, interacting stiff-mode decoupling, nonlinear constraint
positivity, regulator removal, or the full SM--Einstein continuum.
\end{assessment}
\section{Companion audit: degree-weighted forests and exact rebasing}
\label{sec:v10v70}

This is the compulsory audit after V65.  The newest active companion files
were inspected read-only:
\begin{align*}
 &\texttt{Renewal\_Yang\_Mills\_Mass\_Gap\_Research\_Notes\_V1.tex},\\
 &\qquad
 \operatorname{SHA256}
 =\texttt{E54AC40849413570B3EBD015FB8F18DADBB2A0C0392A22CA330552DEEB1A441A},\\
 &\texttt{Renewal\_NS\_Stability\_Research\_Notes\_V26.tex},\\
 &\qquad
 \operatorname{SHA256}
 =\texttt{82C887F8C2C69E4F28FAD7C3DC8DE5B9099CB5A4BF431EBA1FFF3E91935978AD}.
\end{align*}
The Yang--Mills programme has begun a compact new cycle at V1; that file is
newer than the V65 file inspected at the preceding audit.  The
Navier--Stokes file is unchanged.

Two Yang--Mills mechanisms bear directly on the present bottleneck.  First,
an all-good analytic soft factor can sometimes be absorbed exactly into a
stronger positive reference law, leaving only marked difference
insertions.  Second, an ordinary Taylor--Wiener norm pays a factorial for
all forest derivatives incident at one vertex.  The factorials have an
exact labelled-tree sum and change the sufficient forest constant.  Neither
statement depends on Yang--Mills coupling powers.

\subsection{The degree factorial}

Let a labelled tree \(T\) on \([n]\) have vertex degrees \(d_T(i)\).  If
\(F_i\) is analytic on a strong radius \(R\), and \(r<R\) with
\(\delta=R-r\), the one-time Cauchy estimate is
\begin{equation}
 \|\partial^{\nu_i}F_i\|_r
 \leq
 \frac{|\nu_i|!}{\delta^{|\nu_i|}}\|F_i\|_R.
 \label{eq:v70-one-time-Cauchy}
\end{equation}
For a forest integrand, \(|\nu_i|=d_T(i)\).  Thus a fixed-tree bound derived
only from this analytic estimate contains
\begin{equation}
 \prod_{i=1}^nd_T(i)!.
 \label{eq:v70-degree-factor}
\end{equation}
Spending the radius once at the whole vertex avoids an iterative chain of
radii, but it does not remove \eqref{eq:v70-degree-factor}.

\begin{lemma}[Exact degree-weighted labelled-tree sum]
\label{lem:v70-degree-tree-sum}
For \(n\geq2\),
\begin{equation}
 \sum_{T\ {\rm tree\ on}\ [n]}
 \prod_{i=1}^nd_T(i)!
 =
 (n-2)!\binom{3n-3}{n-2}.
 \label{eq:v70-degree-tree-sum}
\end{equation}
\end{lemma}

\begin{proof}
Under the Pruefer correspondence, label \(i\) occurs
\(m_i=d_T(i)-1\) times in a word of length \(n-2\).  For a fixed occurrence
vector \(m_i\geq0\), \(\sum_i m_i=n-2\), the number of words is
\((n-2)!/\prod_i m_i!\).  Its contribution after multiplication by the
degree factorials is
\[
 (n-2)!\prod_i(m_i+1).
\]
The generating function
\[
 \sum_{m\geq0}(m+1)x^m=(1-x)^{-2}
\]
shows that the remaining sum is the coefficient of \(x^{n-2}\) in
\((1-x)^{-2n}\), namely \(\binom{3n-3}{n-2}\).
\end{proof}

\begin{theorem}[Corrected ordinary-analytic forest gate]
\label{thm:v70-corrected-forest}
Assume a fixed labelled tree satisfies
\begin{equation}
 \|{\cal A}_T(F_1,\ldots,F_n)\|_{\rm weak}
 \leq
 z^nK^{n-1}\prod_{i=1}^nd_T(i)!.
 \label{eq:v70-fixed-tree}
\end{equation}
If
\begin{equation}
 8zK<1,
 \label{eq:v70-eight-gate}
\end{equation}
then the connected exponential is absolutely convergent and
\begin{equation}
 \left\|
 \sum_{n\geq1}\frac1{n!}\kappa_C(F,\ldots,F)
 \right\|_{\rm weak}
 \leq
 z+\frac{4z^2K}{1-8zK}.
 \label{eq:v70-eight-sum}
\end{equation}
\end{theorem}

\begin{proof}
For \(n\geq2\), sum \eqref{eq:v70-fixed-tree} over labelled trees, divide by
\(n!\), and apply Lemma \ref{lem:v70-degree-tree-sum}.  The \(n\)-th term is
at most
\[
 z^nK^{n-1}
 \frac1{n(n-1)}
 \binom{3n-3}{n-2}.
\]
Since
\(\binom{3n-3}{n-2}\leq2^{3n-3}\) and \(n(n-1)\geq2\), it is at most
\[
 \frac1{16K}(8zK)^n.
\]
Summing from \(n=2\) gives the second term in
\eqref{eq:v70-eight-sum}; the one-vertex term is at most \(z\).
\end{proof}

\begin{correction}[Scope of the V21 and V67 gates]
\label{cor:v70-v67-gate}
The abstract V21 theorem remains valid under its stated hypothesis
\[
 \|{\cal A}_T\|\leq z^nK_T^{n-1}.
\]
It does not prove that hypothesis for an ordinary analytic Cauchy norm.
Whenever the SM--Einstein fixed-tree estimate is obtained from
\eqref{eq:v70-one-time-Cauchy}, the V67 condition
\(eK_Tz_{\rm for}<1\) must be replaced by
\begin{equation}
 8K_Tz_{\rm for}<1,
 \label{eq:v70-corrected-V67-gate}
\end{equation}
and its connected bound by \eqref{eq:v70-eight-sum}.  The smaller constant
may be used only after a separate norm theorem proves a degree-factor-free
fixed-tree estimate.
\end{correction}

This correction tightens a sufficient constant; it does not invalidate the
finite-volume constructions or create a factorial divergence.

\subsection{Exact rebase and its transfer boundary}

The Yang--Mills companion uses the elementary identity
\begin{equation}
 1=e^{-q}+\left(1-e^{-q}\right)
   =e^{-q}\{1+(e^q-1)\}.
 \label{eq:v70-rebase-identity}
\end{equation}
Multiplying all all-good factors into a positive reference law removes
their derivatives from the forest bank.  The remaining marked insertion is
\(e^q-1\).

For a product Haar law with disjoint cells, both the normalization ratio and
the marked expectation factorize exactly.  Those properties use compact
Haar probability and do not transfer to a general Gaussian graph.  In V69,
neighboring Gaussian innovations overlap, so rebasing all colors at once
can create both fill-in and correlated marked insertions.

One color is different.  Conditioned on its complement, the V69 innovations
are independent centered Gaussians.  Multiplication by
\[
 \exp\left\{-s\sum_{y\in C_c}Q_y\right\}
\]
only strengthens each conditional precision by the scalar factor
\(1+2s\).  Its conditional normalization is independent of the exterior.
Therefore the exact Yang--Mills rebase architecture transfers to one
Gaussian color without importing Haar structure.  V71 will prove this
statement and combine it with Theorem
\ref{thm:v70-corrected-forest}.

\subsection{Navier--Stokes audit result}

The Navier--Stokes V26 computes rank-one norms for derivatives of the Oseen
kernel.  Its improvement is tied to a restricted rank-one input channel and
a certified radial calculation.  The V69 Gaussian constraint penalty must
control the full local symmetric-tensor fibre before the constraints are
taken to their limit.  No theorem in NS V26 proves that all admissible
penalty variations lie in its rank-one channel.  Consequently no constant
or inequality is imported from that note.

\begin{ledger}[V70 transfer decision]
\label{led:v70-transfer}
\begin{enumerate}
\item Import the degree-weighted tree sum and use the corrected constant
      \(8\) for ordinary analytic Cauchy norms.
\item Import the all-good rebase architecture only on a Gaussian independent
      color, where its normalization can be proved exactly.
\item Do not import the Yang--Mills powers of its coupling, Haar integral
      bounds, or Wilson residual estimate.
\item Do not import the Navier--Stokes rank-one constants.
\end{enumerate}
\end{ledger}

\subsection{Post-audit direction}

V71 will derive the exact one-color Gaussian normalization, the rebased
conditional law, and the marked primitive
\[
 (1+2s)^{d_y/2}-1.
\]
It will place that primitive in the corrected degree-weighted forest bank.
The unresolved step will then be a bounded-range schedule that rebases all
colors while keeping the analytic radii and covariance row sums uniform.

\begin{assessment}[Exact status after V70]
\label{ass:v70-status}
V70 corrects the instantiated ordinary-analytic forest constant and
identifies an exact positive Gaussian rebase available on one independent
color.  It does not yet construct the all-color schedule, a uniform
penalty-limit analytic norm, interacting constraint decoupling, regulator
removal, or the full SM--Einstein continuum.
\end{assessment}
\section{Exact Gaussian rebasing and the corrected forest bank}
\label{sec:v10v71}

The V70 audit identified an exact way to remove all-good selector factors
from the forest differentiation.  This section first proves the promised
one-color conditional rebase.  It then exploits additional structure of the
V69 penalty: its precision is a sum of local positive squares.  A uniform
rebase of those energy cells scales the entire Gaussian precision, preserves
finite range, and avoids color-induced fill-in.

\subsection{One independent color}

Use the Gaussian precision graph and coloring of V69.  Fix one independent
color \(C\).  For \(y\in C\), let \(r_y,Q_y\) be the conditional innovations
and quadratic forms in
\eqref{eq:v69-color-innovation}--\eqref{eq:v69-color-selector}.  Choose
numbers \(s_y>0\) and define
\begin{equation}
 G_C
 :=
 \exp\left\{-\sum_{y\in C}s_yQ_y\right\},
 \qquad
 h_y:=e^{s_yQ_y}-1.
 \label{eq:v71-color-rebase-factors}
\end{equation}
Put
\begin{equation}
 c_C:=\prod_{y\in C}(1+2s_y)^{-d_y/2}.
 \label{eq:v71-color-normalization}
\end{equation}

\begin{theorem}[Exact one-color Gaussian rebase]
\label{thm:v71-one-color-rebase}
The probability law
\begin{equation}
 d\mu_C^{\rm reb}
 :=
 c_C^{-1}G_C\,d\mu_A
 \label{eq:v71-rebased-law}
\end{equation}
has the same marginal law on \(x_{C^c}\) as \(\mu_A\).  Conditioned on
\(x_{C^c}\), its innovations \(r_y\), \(y\in C\), are independent centered
Gaussians with covariance
\begin{equation}
 \bigl\{(1+2s_y)A_{yy}\bigr\}^{-1}.
 \label{eq:v71-rebased-covariance}
\end{equation}
Moreover the exact density identity is
\begin{equation}
 d\mu_A
 =
 c_C
 \prod_{y\in C}(1+h_y)\,
 d\mu_C^{\rm reb}.
 \label{eq:v71-color-density-identity}
\end{equation}
For \(Y\subset C\),
\begin{equation}
 {\mathbb E}_{\mu_C^{\rm reb}}
 \prod_{y\in Y}h_y
 =
 \prod_{y\in Y}
 \left\{(1+2s_y)^{d_y/2}-1\right\}.
 \label{eq:v71-color-mark-product}
\end{equation}
\end{theorem}

\begin{proof}
Conditioned on \(x_{C^c}\), Lemma
\ref{lem:v69-color-factorization} writes the \(C\)-dependent exponent as
\[
 \frac12\sum_{y\in C}r_y^TA_{yy}r_y.
\]
Multiplication by \(G_C\) changes this to
\[
 \frac12\sum_{y\in C}(1+2s_y)r_y^TA_{yy}r_y.
\]
The conditional Gaussian integral is multiplied by
\(\prod_y(1+2s_y)^{-d_y/2}=c_C\), independently of
\(x_{C^c}\).  Hence \eqref{eq:v71-rebased-law} is normalized, its exterior
marginal is unchanged, and its conditional covariance is
\eqref{eq:v71-rebased-covariance}.

Since \(G_C^{-1}=\prod_y e^{s_yQ_y}=\prod_y(1+h_y)\), rearranging
\eqref{eq:v71-rebased-law} proves
\eqref{eq:v71-color-density-identity}.  Under the rebased conditional law,
\((1+2s_y)Q_y\) is chi-squared with \(d_y\) degrees of freedom.  Therefore
\[
 {\mathbb E}_{\rm reb}e^{s_yQ_y}
 =
 \left(
  1-\frac{2s_y}{1+2s_y}
 \right)^{-d_y/2}
 =
 (1+2s_y)^{d_y/2}.
\]
Conditional independence and the fact that this moment is independent of
the exterior prove \eqref{eq:v71-color-mark-product}.
\end{proof}

Thus the exact one-color marked primitive is
\begin{equation}
 q_{{\rm reb},y}
 :=
 (1+2s_y)^{d_y/2}-1
 =
 d_ys_y+O(s_y^2).
 \label{eq:v71-color-primitive}
\end{equation}
The all-good factor \(G_C\) belongs to the new reference law and has no
forest derivative.  The local normalization cost
\[
 -\log c_C
 =
 \sum_{y\in C}\frac{d_y}{2}\log(1+2s_y)
\]
is an explicit vacuum owner.

\subsection{Positive-square energy cells}

The penalty precision has a more useful representation.  Let
\({\cal I}\) be a finite set of local energy cells and suppose
\begin{equation}
 A=\sum_{a\in{\cal I}}E_a,
 \qquad
 E_a=B_a^*B_a\succeq0,
 \label{eq:v71-square-decomposition}
\end{equation}
where each \(B_a\) has uniformly bounded support and the number of cells
meeting one block is at most \(C_{\rm cell}\).  On the zero-mode complement,
assume \(A>0\).  Define
\begin{equation}
 Q_a(x)=\langle x,E_ax\rangle=\|B_ax\|^2.
 \label{eq:v71-cell-energy}
\end{equation}
Then
\[
 \sum_{a\in{\cal I}}Q_a=x^TAx.
\]
The lattice gradient, time-difference, divergence-penalty, and trace-penalty
terms in V69 all have this form.

Choose one common \(s>0\), set
\begin{equation}
 G_{\rm all}=e^{-s\sum_aQ_a},
 \qquad
 h_a=e^{sQ_a}-1,
 \qquad
 c_s=(1+2s)^{-N/2},
 \label{eq:v71-global-factors}
\end{equation}
where \(N\) is the dimension of the integrated zero-mode complement.

\begin{theorem}[Range-preserving all-cell rebase]
\label{thm:v71-global-rebase}
Let \(\mu_A\) be the normalized Gaussian with precision \(A\), and let
\(\mu_s\) be the normalized Gaussian with precision
\begin{equation}
 A_s=(1+2s)A.
 \label{eq:v71-scaled-precision}
\end{equation}
Then
\begin{equation}
 d\mu_A
 =
 c_s\prod_{a\in{\cal I}}(1+h_a)\,d\mu_s.
 \label{eq:v71-global-density}
\end{equation}
The precision range and graph degree are unchanged, and
\begin{equation}
 A_s^{-1}=(1+2s)^{-1}A^{-1}.
 \label{eq:v71-scaled-covariance}
\end{equation}
If \(A\) is reflection positive through the V69 oscillator construction,
then so is \(A_s\).
\end{theorem}

\begin{proof}
The Gaussian change of variables \(u=A^{1/2}x\) gives
\[
 {\mathbb E}_{\mu_A}e^{-sx^TAx}
 =
 {\mathbb E}e^{-s|u|^2}
 =
 (1+2s)^{-N/2}=c_s.
\]
Consequently
\[
 d\mu_s=c_s^{-1}G_{\rm all}\,d\mu_A.
\]
Since
\[
 G_{\rm all}^{-1}
 =
 \prod_ae^{sQ_a}
 =
 \prod_a(1+h_a),
\]
rearrangement proves \eqref{eq:v71-global-density}.  Scalar multiplication
of \(A\) proves the range, degree, and covariance claims.  In the
reflection-square formula, scalar multiplication merely rescales the
positive spatial operator and its positive spectral measure, so positivity
is preserved.
\end{proof}

Unlike successive color rebasing, Theorem
\ref{thm:v71-global-rebase} creates no neighbor-of-neighbor fill-in.  The
factor \(c_s\) is extensive, but
\[
 -\log c_s=\frac N2\log(1+2s)
\]
is a known local vacuum-energy owner.  It cancels exactly against the
reference normalization when \eqref{eq:v71-global-density} is used; it must
not be counted as a physical determinant anomaly.

\subsection{One-cell marked moment}

The energy cells need not be probabilistically independent.  They still
have a uniform one-cell moment under the stronger reference.

\begin{proposition}[Uniform marked cell primitive]
\label{prop:v71-cell-primitive}
Let \(r_a=\operatorname{rank}E_a\).  Under \(\mu_s\),
\begin{equation}
 {\mathbb E}_{\mu_s}h_a
 \leq
 (1+2s)^{r_a/2}-1.
 \label{eq:v71-cell-primitive}
\end{equation}
\end{proposition}

\begin{proof}
From \eqref{eq:v71-square-decomposition},
\(0\preceq E_a\preceq A\).  Hence
\[
 0\preceq
 M_a:=A^{-1/2}E_aA^{-1/2}
 \preceq I.
\]
For a centered Gaussian with covariance
\((1+2s)^{-1}A^{-1}\), the quadratic Gaussian determinant formula gives
\[
 {\mathbb E}_{\mu_s}e^{sQ_a}
 =
 \det\left(
 I-\frac{2s}{1+2s}M_a
 \right)^{-1/2}.
\]
There are at most \(r_a\) nonzero eigenvalues of \(M_a\), each in
\([0,1]\).  Every determinant factor is therefore at most
\[
 \left(1-\frac{2s}{1+2s}\right)^{-1/2}
 =(1+2s)^{1/2}.
\]
Subtract one to obtain \eqref{eq:v71-cell-primitive}.
\end{proof}

This estimate is not a product bound for overlapping cells.  In the route
below, the \(h_a\) are analytic forest vertices, so no product-moment claim
is needed.

\subsection{Analytic marked bank}

Let the strong and weak field radii be \(R_{\rm an}>r_{\rm an}>0\), with
\[
 \delta_{\rm an}=R_{\rm an}-r_{\rm an}.
\]
Assume
\begin{equation}
 \sup_a\|sQ_a\|_{R_{\rm an}}\leq\epsilon_{\rm cell}.
 \label{eq:v71-cell-analytic-stock}
\end{equation}
The Banach-algebra exponential estimate gives
\begin{equation}
 \|h_a\|_{R_{\rm an}}
 \leq
 \eta_{\rm cell}:=e^{\epsilon_{\rm cell}}-1.
 \label{eq:v71-cell-analytic-activity}
\end{equation}
Because at most \(C_{\rm cell}\) cell owners meet one anchor, their total
anchored local stock is bounded by
\begin{equation}
 z_{\rm cell}:=C_{\rm cell}\eta_{\rm cell}.
 \label{eq:v71-cell-bank}
\end{equation}

Let
\begin{equation}
 z_{\rm reb}
 :=
 z_{\rm core}+z_{\rm contact}+d_\kappa+z_{\rm cell}.
 \label{eq:v71-rebased-bank}
\end{equation}
The V67 selector stock \(s_\kappa\) is absent: its all-good part is already
in \(\mu_s\), and its marked part is exactly \(z_{\rm cell}\).

If \(G_\mu(A^{-1})\) is the weighted covariance row sum, then
\eqref{eq:v71-scaled-covariance} gives
\begin{equation}
 G_\mu(A_s^{-1})
 =
 \frac{G_\mu(A^{-1})}{1+2s}.
 \label{eq:v71-row-improvement}
\end{equation}
With endpoint multiplicity \(h_*\), define
\begin{equation}
 K_{\mu,s}
 :=
 \frac{h_*G_\mu(A^{-1})}
 {(1+2s)\delta_{\rm an}^2}.
 \label{eq:v71-rebased-edge}
\end{equation}

\begin{theorem}[Corrected rebased forest gate]
\label{thm:v71-rebased-forest}
Suppose every fixed labelled tree is bounded by
\[
 z_{\rm reb}^nK_{\mu,s}^{n-1}
 \prod_{i=1}^nd_T(i)!
\]
in the common support--grade--Taylor--Wiener norm.  If
\begin{equation}
 8z_{\rm reb}K_{\mu,s}<1,
 \label{eq:v71-rebased-gate}
\end{equation}
then the complete connected rebased expansion converges absolutely and is
bounded by
\begin{equation}
 z_{\rm reb}
 +
 \frac{4z_{\rm reb}^2K_{\mu,s}}
 {1-8z_{\rm reb}K_{\mu,s}}.
 \label{eq:v71-rebased-sum}
\end{equation}
No all-good selector derivative or extra V66 soft-selector activity occurs.
\end{theorem}

\begin{proof}
The exact density identity
\eqref{eq:v71-global-density} expresses every marked cell as the analytic
vertex \(h_a\) and puts the common all-good factor in \(\mu_s\).
Equations \eqref{eq:v71-cell-bank} and
\eqref{eq:v71-rebased-bank} bound the disjoint owner bank.
Equation \eqref{eq:v71-row-improvement} bounds every covariance edge.
The degree-weighted hypotheses of Theorem
\ref{thm:v70-corrected-forest} now hold with
\(z=z_{\rm reb}\) and \(K=K_{\mu,s}\), which proves the result.
\end{proof}

\begin{ledger}[Choice of selector route]
\label{led:v71-route-choice}
\begin{enumerate}
\item The one-color route uses
      \eqref{eq:v71-color-density-identity} and may charge the exact product
      primitive \eqref{eq:v71-color-mark-product}.
\item The all-cell route uses
      \eqref{eq:v71-global-density} and charges \(z_{\rm cell}\) in the
      analytic forest.
\item A construction chooses one route for a given selector factor.  It
      does not add the one-color primitive, the V69 colored primitive, and
      \(z_{\rm cell}\) for the same marks.
\item Far-field contours remain a separate large-field owner in either
      route.
\end{enumerate}
\end{ledger}

\subsection{Application to the constraint penalty}

For the V69 precision, take cells for temporal differences, spatial
differences, lattice divergence, and trace:
\[
 x^TA_\varepsilon x
 =
 \sum_aQ_{a,{\rm kin}}
 +
 \varepsilon^{-1}\sum_aQ_{a,{\rm div}}
 +
 \varepsilon^{-1}\sum_aQ_{a,{\rm tr}}.
\]
Theorem \ref{thm:v71-global-rebase} preserves this precision range and the
TT penalty limit.  It also improves the covariance edge stock by
\((1+2s)^{-1}\).

The stiff analytic debit remains visible:
\begin{equation}
 \|sQ_{a,{\rm div}}/\varepsilon\|_{R_{\rm an}}
 +
 \|sQ_{a,{\rm tr}}/\varepsilon\|_{R_{\rm an}}
 \leq\epsilon_{\rm cell}
 \label{eq:v71-stiff-cell-debit}
\end{equation}
requires an \(O(\sqrt{\varepsilon})\) radius in the
constraint-violating field combinations at fixed \(s\).  Increasing \(s\)
improves the covariance row but worsens the marked analytic norm.  The
choice of \(s\) therefore has a genuine two-sided window.

\begin{firewall}[Rebasing is an identity, not decoupling]
\label{fire:v71-no-decoupling}
Equation \eqref{eq:v71-global-density} is an exact change of reference at
fixed \(\varepsilon\).  It does not prove that interacting vertices are
uniform when constraint-violating radii shrink, nor that the singular
Gaussian normalization is the physical Einstein coarea determinant.
Those statements require an anisotropic interaction theorem and the
constraint determinant ledger.
\end{firewall}

\subsection{Next acquisition}

The global energy-cell rebase removes the fill-in problem anticipated at
V70.  The remaining bottleneck is now the anisotropic analytic norm.  V72
will split each local field jet into a soft physical owner and the images of
the local constraint operators \(D\) and \(T\), assign the stiff variables
their natural \(\sqrt{\varepsilon}\) radii, and determine which nonlinear
SM--Einstein vertices remain uniformly analytic under that scaling.  A
vertex with a negative power of \(\sqrt{\varepsilon}\) after constraint
projection will be identified as a nondecoupling obstruction rather than
estimated away.

\begin{assessment}[Exact status after V71]
\label{ass:v71-status}
V71 proves the exact one-color and all-energy-cell Gaussian rebases,
preservation of range and reflection positivity, a uniform one-cell marked
moment, and the corrected degree-weighted forest gate.  It eliminates the
all-good selector from the analytic bank at fixed penalty.  A uniform
anisotropic interaction norm, singular determinant matching, nonlinear
constraint decoupling, regulator removal, and the full SM--Einstein
continuum remain open.
\end{assessment}
\section{Exact constraint-residual scaling and the anisotropic norm}
\label{sec:v10v72}

V71 left a stiff analytic factor \(1/\varepsilon\) in the original metric
coordinates.  That factor is artificial if the penalty is intended to
approximate an exact constraint delta function.  The correct stiff
coordinate is the nonlinear constraint residual itself.  In that coordinate
the normalized delta penalty is a standard Gaussian, and every analytic
dependence transverse to the constraint surface gains powers of
\(\sqrt{\varepsilon}\).

This section proves the finite-dimensional theorem and its local analytic
version.  Its application to gravity is conditional on selecting an
independent constant-rank constraint and gauge residual, with all Bianchi,
Killing, boundary, and zero-mode dependencies removed into their existing
ledgers.

\subsection{Constant-rank residual chart}

Let \(X\) be an \(n\)-dimensional regulated field space and let
\[
 {\cal C}:X\longrightarrow{\mathbb R}^m,\qquad m<n,
\]
be a real analytic residual map.  Fix a chart in which
\(D{\cal C}\) has rank \(m\).  Suppose there is a real analytic
diffeomorphism
\begin{equation}
 \Psi:(u,v)\longmapsto x,
 \qquad
 u\in U\subset{\mathbb R}^{n-m},\quad
 v\in V\subset{\mathbb R}^m,
 \label{eq:v72-residual-chart}
\end{equation}
such that
\begin{equation}
 {\cal C}(\Psi(u,v))=v.
 \label{eq:v72-exact-residual}
\end{equation}
Write
\begin{equation}
 dx=J(u,v)\,du\,dv,
 \qquad J(u,v)>0
 \label{eq:v72-chart-Jacobian}
\end{equation}
on the real chart.  All orbit, coarea, and physical determinant factors
that are not part of \(J\) remain separate owners.

For \(\varepsilon>0\), define the normalized Gaussian delta approximant
\begin{equation}
 \delta_\varepsilon(v)
 :=
 (2\pi\varepsilon)^{-m/2}
 e^{-|v|^2/(2\varepsilon)}.
 \label{eq:v72-delta-epsilon}
\end{equation}

\begin{theorem}[Exact residual scaling]
\label{thm:v72-exact-scaling}
For every integrable \(F\) supported in the residual chart,
\begin{equation}
 \int_XF(x)\delta_\varepsilon({\cal C}(x))\,dx
 =
 \int_U\int_{\mathbb R^m}
 F(\Psi(u,\sqrt{\varepsilon}y))
 J(u,\sqrt{\varepsilon}y)\,d\gamma_m(y)\,du,
 \label{eq:v72-scaled-integral}
\end{equation}
where
\[
 d\gamma_m(y)=(2\pi)^{-m/2}e^{-|y|^2/2}\,dy,
\]
with the integrand extended by zero outside
\(\sqrt{\varepsilon}y\in V\).  If the right side is dominated uniformly and
the boundary of \(V\) stays a positive distance from \(v=0\), then
\begin{equation}
 \lim_{\varepsilon\downarrow0}
 \int_XF(x)\delta_\varepsilon({\cal C}(x))\,dx
 =
 \int_UF(\Psi(u,0))J(u,0)\,du.
 \label{eq:v72-coarea-limit}
\end{equation}
\end{theorem}

\begin{proof}
Use \eqref{eq:v72-exact-residual} and
\eqref{eq:v72-chart-Jacobian}, then substitute
\(v=\sqrt{\varepsilon}y\).  The factor
\(\varepsilon^{m/2}\) from \(dv\) cancels the corresponding factor in
\eqref{eq:v72-delta-epsilon}, proving
\eqref{eq:v72-scaled-integral}.  For every fixed \(u,y\), the integrand
converges to \(F(\Psi(u,0))J(u,0)\).  The declared domination and the
Gaussian tail allow dominated convergence.  Since
\(\gamma_m\) has total mass one, the limit is
\eqref{eq:v72-coarea-limit}.
\end{proof}

There is no divergent normalization left after the residual is scaled.
The factor \(J(u,0)\) is exactly the chart/coarea density on the constraint
surface.  Whether it is the correct physical Einstein density depends on
which constraints and gauge equations were used to define \({\cal C}\).

\subsection{A quantitative convergence estimate}

Let
\begin{equation}
 H(u,v):=F(\Psi(u,v))J(u,v).
 \label{eq:v72-H-density}
\end{equation}

\begin{proposition}[Even Gaussian improvement]
\label{prop:v72-even-error}
Assume \(H\) is twice differentiable in \(v\), the \(u\)-integral of
\[
 M_2(u):=
 \sup_{v\in V}\|D_v^2H(u,v)\|
\]
is finite, and the omitted Gaussian tail outside \(V\) is bounded by
\(T_\varepsilon\).  Then
\begin{equation}
 \left|
 \int F\delta_\varepsilon({\cal C})\,dx
 -
 \int_UH(u,0)\,du
 \right|
 \leq
 \frac{\varepsilon m}{2}
 \int_UM_2(u)\,du+T_\varepsilon.
 \label{eq:v72-even-error}
\end{equation}
If \(V\) contains \(\{|v|\leq r_v\}\) and \(H\) is bounded there, the local
tail has the form
\begin{equation}
 T_\varepsilon
 \leq C_H e^{-c r_v^2/\varepsilon}
 \label{eq:v72-Gaussian-tail}
\end{equation}
for fixed positive constants \(C_H,c\).
\end{proposition}

\begin{proof}
Taylor's formula gives
\[
 H(u,\sqrt{\varepsilon}y)
 =
 H(u,0)
 +
 \sqrt{\varepsilon}D_vH(u,0)y
 +
 \frac{\varepsilon}{2}
 D_v^2H(u,\theta\sqrt{\varepsilon}y)[y,y]
\]
for a pointwise \(0<\theta<1\) while the segment remains in \(V\).
The centered Gaussian integral of the linear term is zero, and
\(\int|y|^2d\gamma_m(y)=m\).  This proves the first term in
\eqref{eq:v72-even-error}.  The standard radial Gaussian tail bound outside
\(|y|\leq r_v/\sqrt{\varepsilon}\) gives
\eqref{eq:v72-Gaussian-tail}, after absorbing its polynomial prefactor into
a smaller positive exponent constant.
\end{proof}

The compact boundary device used to isolate the tail is a
measure-theoretic large-field owner.  It is not differentiated inside the
all-order analytic forest.

\subsection{Anisotropic analytic scaling}

Let \({\cal A}(R_u,R_v)\) be a coefficient-majorant analytic Banach algebra
on a complex product polydisc with radii \(R_u,R_v\):
\begin{equation}
 H(u,v)=\sum_{\alpha,\beta}H_{\alpha\beta}u^\alpha v^\beta,
 \qquad
 \|H\|_{R_u,R_v}
 =
 \sum_{\alpha,\beta}
 \|H_{\alpha\beta}\|R_u^{|\alpha|}R_v^{|\beta|}.
 \label{eq:v72-anisotropic-norm}
\end{equation}
Support, covariant grade, and owner weights may be inserted coefficientwise
as in V57.

Define
\begin{equation}
 H_\varepsilon(u,y):=H(u,\sqrt{\varepsilon}y).
 \label{eq:v72-rescaled-H}
\end{equation}

\begin{theorem}[Uniform stiff-direction analytic bound]
\label{thm:v72-anisotropic-bound}
If \(\sqrt{\varepsilon}R_y\leq R_v\), then
\begin{equation}
 \|H_\varepsilon\|_{R_u,R_y}
 =
 \|H\|_{R_u,\sqrt{\varepsilon}R_y}
 \leq
 \|H\|_{R_u,R_v}.
 \label{eq:v72-uniform-bound}
\end{equation}
Moreover,
\begin{equation}
 \|H_\varepsilon-H(\,\cdot\,,0)\|_{R_u,R_y}
 \leq
 \frac{\sqrt{\varepsilon}R_y}{R_v}
 \|H-H(\,\cdot\,,0)\|_{R_u,R_v},
 \label{eq:v72-transverse-gain}
\end{equation}
and every \(k\)-th derivative in \(y\) obeys the exact identity
\begin{equation}
 D_y^kH_\varepsilon
 =
 \varepsilon^{k/2}
 (D_v^kH)(u,\sqrt{\varepsilon}y).
 \label{eq:v72-derivative-gain}
\end{equation}
\end{theorem}

\begin{proof}
Substitution in \eqref{eq:v72-anisotropic-norm} multiplies the coefficient
of transverse degree \(|\beta|\) by
\(\varepsilon^{|\beta|/2}\).  This proves the equality and inequality in
\eqref{eq:v72-uniform-bound}.  Every term in
\(H-H(\,\cdot\,,0)\) has \(|\beta|\geq1\), so
\[
 (\sqrt{\varepsilon}R_y)^{|\beta|}
 \leq
 \frac{\sqrt{\varepsilon}R_y}{R_v}R_v^{|\beta|}.
\]
Summation proves \eqref{eq:v72-transverse-gain}.  The chain rule proves
\eqref{eq:v72-derivative-gain}.
\end{proof}

Thus the stiff variable has a fixed \(y\)-radius while its physical
residual \(v\) has radius \(\sqrt{\varepsilon}R_y\).  The soft variable
\(u\) keeps its full radius \(R_u\).  No nonlocal TT projector occurs in the
definition of the norm.

\subsection{Uniform selector rebase in residual coordinates}

Partition the independent residual coordinates into fixed local cells
\(y_a\).  The delta Gaussian is
\[
 d\gamma(y)=\prod_a d\gamma_{m_a}(y_a).
\]
For \(s>0\), define
\begin{equation}
 g_a=e^{-s|y_a|^2},
 \qquad
 h_a=e^{s|y_a|^2}-1.
 \label{eq:v72-residual-selector}
\end{equation}
The all-good rebase changes the standard Gaussian precision from \(I\) to
\((1+2s)I\).  The exact V71 identities give
\begin{equation}
 d\gamma(y)
 =
 (1+2s)^{-m/2}
 \prod_a(1+h_a)\,d\gamma_s(y),
 \label{eq:v72-residual-rebase}
\end{equation}
where \(d\gamma_s\) is the normalized Gaussian with covariance
\((1+2s)^{-1}I\).

On the \(y\)-polydisc of radius \(R_y\),
\begin{equation}
 \|h_a\|_{R_y}
 \leq
 e^{sC_{m_a}R_y^2}-1,
 \label{eq:v72-uniform-mark}
\end{equation}
where \(C_{m_a}\) depends only on the fixed cell dimension.  This bound is
independent of \(\varepsilon\).  It replaces the divergent original-field
bound \eqref{eq:v71-stiff-cell-debit}.

\begin{corollary}[Uniform corrected forest row]
\label{cor:v72-uniform-forest-row}
Assume the residual chart \(\Psi\), its Jacobian, the pulled-back physical
interaction, and their anchored support--grade norms are bounded uniformly
in \(\varepsilon\) on
\[
 {\cal D}(R_u,\sqrt{\varepsilon}R_y).
\]
Let \(z_{\rm phys}\) denote their disjoint owner stock after the complete
low-grade projection, and let
\begin{equation}
 z_{\rm mark}
 :=
 C_{\rm cell}\{e^{sC_mR_y^2}-1\}.
 \label{eq:v72-mark-stock}
\end{equation}
If the rebased Gaussian edge stock is \(K_s\) and
\begin{equation}
 8K_s(z_{\rm phys}+z_{\rm mark})<1,
 \label{eq:v72-uniform-forest-gate}
\end{equation}
then the degree-weighted connected forest is uniformly summable in
\(\varepsilon\), with the bound
\begin{equation}
 z_{\rm phys}+z_{\rm mark}
 +
 \frac{4K_s(z_{\rm phys}+z_{\rm mark})^2}
 {1-8K_s(z_{\rm phys}+z_{\rm mark})}.
 \label{eq:v72-uniform-forest-sum}
\end{equation}
\end{corollary}

\begin{proof}
Theorem \ref{thm:v72-anisotropic-bound} makes the pulled-back interaction
uniform, and \eqref{eq:v72-uniform-mark} bounds the exact rebased mark bank.
Apply Theorem \ref{thm:v70-corrected-forest}.
\end{proof}

\subsection{Why the exact residual is necessary}

Suppose instead that the nonlinear constraint is written in an unaligned
linear coordinate \(w\):
\begin{equation}
 {\cal C}(u,w)=Lw+N(u,w).
 \label{eq:v72-unaligned-constraint}
\end{equation}
The penalty contains
\begin{equation}
 \frac1{2\varepsilon}\|Lw+N(u,w)\|^2.
 \label{eq:v72-unaligned-penalty}
\end{equation}
Even if \(Lw=\sqrt{\varepsilon}y\), its cross term is
\[
 \varepsilon^{-1/2}\langle y,N(u,w)\rangle,
\]
and the last term is
\(\|N(u,w)\|^2/(2\varepsilon)\).  They are not uniform at fixed soft
radius unless nongeneric powers of \(\varepsilon\) are inserted into
\(N\).

\begin{proposition}[Graph centering removes all negative penalty powers]
\label{prop:v72-graph-centering}
If the exact residual chart
\eqref{eq:v72-exact-residual} exists, then
\[
 \frac1{2\varepsilon}
 \|{\cal C}(\Psi(u,v))\|^2
 =
 \frac12|y|^2
 \qquad(v=\sqrt{\varepsilon}y).
\]
No negative power of \(\sqrt{\varepsilon}\) occurs in the remaining analytic
density \(H_\varepsilon\).  Conversely, using
\eqref{eq:v72-unaligned-constraint} leaves the displayed negative powers
unless the nonlinear remainder vanishes along the chosen soft chart.
\end{proposition}

\begin{proof}
The first assertion is immediate from
\({\cal C}(\Psi(u,v))=v\).  The remaining density is
\(H(u,\sqrt{\varepsilon}y)\), whose coefficients have only nonnegative
powers by Theorem \ref{thm:v72-anisotropic-bound}.  Expansion of
\eqref{eq:v72-unaligned-penalty} proves the converse statement.
\end{proof}

This is the penalty analogue of the aligned graph coordinates in
V25, V54, and V61.  Recentring only at the linear stationary point is
insufficient; the full nonlinear residual must be a coordinate.

\subsection{Gravity-specific constant-rank card}

For the SM--Einstein application, the residual vector must contain an
independent set chosen from:
\[
 \hbox{Hamiltonian constraint},\quad
 \hbox{momentum constraint},\quad
 \hbox{gauge conditions},\quad
 \hbox{interface matching}.
\]
The number of residual components must equal the number of eliminated
directions.  Differential identities are removed before forming the
Jacobian.  Conformal-Killing, harmonic, chiral, boundary-charge, and
constant metric modes remain retained coordinates.

\begin{definition}[Uniform residual-chart card]
\label{def:v72-residual-card}
A regulator scale passes the card if:
\begin{enumerate}
\item the selected constraint and gauge residual has constant rank on a
      reflection-compatible chart;
\item its complement inverse is the V63 nested inverse, with all four pivot
      and locality bounds uniform;
\item \(\Psi\), \(J\), and their retained derivatives have uniformly
      bounded anchored support--grade--analytic norms on
      \({\cal D}(R_u,R_v)\);
\item the residual chart intertwines reflection and its Jacobian is assigned
      to the correct coarea, orbit, and interface sectors;
\item the pulled-back physical interaction passes the low-grade projection
      and \eqref{eq:v72-uniform-forest-gate}; and
\item the exterior of the chart has the separate large-field estimate
      required by V22 and V66.
\end{enumerate}
\end{definition}

At a fixed regulator in the V23 sign chamber and the V49/V63 Schur chamber,
items one and two are local implicit-function statements already available
sectorwise.  Their simultaneous scale-uniform version, item three, is not
yet proved.

\begin{firewall}[Delta constraints require a physical derivation]
\label{fire:v72-physical-delta}
Theorems \ref{thm:v72-exact-scaling} and
\ref{thm:v72-anisotropic-bound} apply to any declared residual map.  They do
not license inserting a new equation of motion into the path integral.  The
canonical Hamiltonian and momentum constraints have a physical constraint
origin.  The V62 four-dimensional bulk trace equation does not yet have
that status.  It may not be included in \({\cal C}\) until the physical
equivalence debit is resolved.
\end{firewall}

\subsection{Next acquisition}

V72 removes the singular analytic-radius debit conditionally on the exact
nonlinear residual chart.  The next upstream issue is uniform locality of
that chart as the constraint penalty and lattice cutoff are removed.  V73
will differentiate the identity
\({\cal C}(\Psi(u,v))=v\), express all residual-chart jets through the V63
four-stage inverse, and derive a majorant recursion for their anchored
analytic norms.  The recursion will reveal whether a fixed complex radius
survives or collapses at the scalar, DeTurck, CMC, or interface pivot.

\begin{assessment}[Exact status after V72]
\label{ass:v72-status}
V72 proves exact normalized constraint-delta scaling, an
\(O(\varepsilon)\) integrated convergence estimate, uniform anisotropic
analytic bounds, a penalty-independent selector activity, and the corrected
forest implication in exact residual coordinates.  The uniform
gravity-specific residual chart, its physical determinant matching,
nonlinear large-field coercivity, regulator removal, and the full
SM--Einstein continuum remain open.
\end{assessment}
\section{Majorant construction of the nonlinear residual chart}
\label{sec:v10v73}

V72 made the constraint-penalty limit uniform once an exact nonlinear
residual chart with uniform anchored analytic bounds is available.  This
section constructs such a chart from quantitative data.  The theorem is an
analytic implicit-function result in the weighted owner algebra, with a
scalar majorant that controls every derivative at once.

\subsection{Eliminated and retained variables}

At one regulator scale, let \(u\) collect the retained physical, boundary,
zero-mode, and running-coupling coordinates, and let \(w\) collect the
directions eliminated by the independent constraint, gauge, and matching
residual.  After translation to a base solution, write
\begin{equation}
 {\cal C}(u,w)
 =
 Lw+Au+N(u,w),
 \label{eq:v73-residual-expansion}
\end{equation}
where
\begin{equation}
 N(0,0)=0,\qquad DN(0,0)=0.
 \label{eq:v73-N-quadratic}
\end{equation}
The exact residual coordinate is \(v={\cal C}(u,w)\).  All spaces are
finite-dimensional at this stage, but their norms are the volume-free
weighted support--grade--analytic norms.

Assume
\begin{equation}
 \|L^{-1}\|_\kappa\leq\beta,
 \qquad
 \|A\|_\kappa\leq a.
 \label{eq:v73-linear-stocks}
\end{equation}
The constant \(\beta\) is supplied by the complete V63 four-stage inverse,
not by one pivot alone.

Expand the nonlinear part as
\begin{equation}
 N(u,w)=
 \sum_{\substack{p,q\geq0\\p+q\geq2}}N_{pq}[u^{\otimes p},w^{\otimes q}],
 \label{eq:v73-N-series}
\end{equation}
and choose nonnegative coefficients \(n_{pq}\) that include all fixed
support convolution and owner multiplicities and satisfy
\begin{equation}
 \|N_{pq}[u^{\otimes p},w^{\otimes q}]\|_\kappa
 \leq
 n_{pq}\|u\|_\kappa^p\|w\|_\kappa^q.
 \label{eq:v73-N-majorization}
\end{equation}
Define the scalar majorant
\begin{equation}
 {\mathfrak n}(r,t)
 :=
 \sum_{p+q\geq2}n_{pq}r^pt^q.
 \label{eq:v73-scalar-majorant}
\end{equation}

\subsection{The radius window}

Choose positive \(r_u,r_v,r_w\) inside the convergence domain of
\({\mathfrak n}\), and put
\begin{equation}
 q_{\rm imp}
 :=
 \beta\,\partial_t{\mathfrak n}(r_u,r_w).
 \label{eq:v73-implicit-q}
\end{equation}

\begin{theorem}[Weighted analytic residual chart]
\label{thm:v73-residual-chart}
Suppose
\begin{align}
 \beta\{r_v+ar_u+{\mathfrak n}(r_u,r_w)\}
 &\leq r_w,
 \label{eq:v73-self-map}\\
 q_{\rm imp}&<1.
 \label{eq:v73-contraction}
\end{align}
Then for every complex
\(\|u\|_\kappa\leq r_u\), \(\|v\|_\kappa\leq r_v\), there is a unique
\(\|w\|_\kappa\leq r_w\) satisfying
\begin{equation}
 {\cal C}(u,w)=v.
 \label{eq:v73-residual-equation}
\end{equation}
The solution \(w=\sigma(u,v)\) is analytic, and
\begin{equation}
 \Psi(u,v):=(u,\sigma(u,v))
 \label{eq:v73-Psi}
\end{equation}
obeys the exact identity
\({\cal C}(\Psi(u,v))=v\).
\end{theorem}

\begin{proof}
For fixed \(u,v\), define
\[
 {\cal T}_{u,v}(w)
 =
 L^{-1}\{v-Au-N(u,w)\}.
\]
If \(\|w\|_\kappa\leq r_w\), the coefficient majorant gives
\[
 \|{\cal T}_{u,v}(w)\|_\kappa
 \leq
 \beta\{r_v+ar_u+{\mathfrak n}(r_u,r_w)\}
 \leq r_w.
\]
For \(w_1,w_2\) in this ball, integrate \(D_wN\) along their line segment.
The nonnegative scalar majorant gives
\[
 \|N(u,w_1)-N(u,w_2)\|_\kappa
 \leq
 \partial_t{\mathfrak n}(r_u,r_w)
 \|w_1-w_2\|_\kappa.
\]
Thus \({\cal T}_{u,v}\) is a contraction with constant
\(q_{\rm imp}<1\).  Banach's theorem gives existence and uniqueness.

The Picard iterates are analytic in \(u,v\) and converge uniformly on every
strictly smaller closed polydisc because the contraction constant is
uniform.  Their limit is analytic.  Substitution gives the exact residual
identity.
\end{proof}

The inequalities are a genuine two-radius window.  The first controls the
image of the ball, while the second controls the moving inverse.  Neither
one follows merely from invertibility at the base point.

\subsection{All-order coefficient majorant}

Let \(\widehat w(r,s)\) be the smallest nonnegative solution obtained by
monotone iteration from zero in
\begin{equation}
 \widehat w
 =
 \beta\{s+ar+{\mathfrak n}(r,\widehat w)\}.
 \label{eq:v73-scalar-fixed-point}
\end{equation}
Under \eqref{eq:v73-self-map}--\eqref{eq:v73-contraction}, it is analytic
for \(0\leq r<r_u\), \(0\leq s<r_v\).

\begin{proposition}[Coefficientwise implicit majorant]
\label{prop:v73-coefficient-majorant}
The Taylor series of \(\sigma(u,v)\) at the origin is coefficientwise
majorized in the declared norm by the nonnegative Taylor series of
\(\widehat w(r,s)\).  In particular,
\begin{equation}
 \|\sigma\|_{r_u,r_v}\leq\widehat w(r_u,r_v)\leq r_w
 \label{eq:v73-sigma-majorant}
\end{equation}
in the coefficient-majorant analytic norm.
\end{proposition}

\begin{proof}
Start both Picard iterations at zero.  Suppose the coefficients of the
vector iterate are bounded by the corresponding coefficients of the scalar
iterate.  The bound for \(L^{-1}\), the coefficient bounds
\eqref{eq:v73-N-majorization}, and nonnegativity of every \(n_{pq}\) preserve
that ordering after one substitution.  Induction proves it for every
iterate.  Uniform convergence permits passage to the limit.  The scalar
iteration stays in \([0,r_w]\) and is a contraction, proving the final
bound.
\end{proof}

This proposition controls arbitrary star degrees before the forest formula
is applied.  Forest differentiation must still use the V70
degree-weighted constant.

\subsection{First residual-chart jets}

Differentiate
\({\cal C}(u,\sigma(u,v))=v\).  Let
\[
 M(u,v):=L+D_wN(u,\sigma(u,v)).
\]
The Neumann condition follows from
\eqref{eq:v73-contraction}:
\begin{equation}
 \|L^{-1}D_wN\|_\kappa\leq q_{\rm imp}<1.
 \label{eq:v73-Neumann-condition}
\end{equation}

\begin{proposition}[Explicit first-jet stocks]
\label{prop:v73-first-jets}
Throughout the chart,
\begin{align}
 \|M^{-1}\|_\kappa
 &\leq
 \frac{\beta}{1-q_{\rm imp}},
 \label{eq:v73-M-inverse}\\
 \|D_v\sigma\|_\kappa
 &\leq
 \frac{\beta}{1-q_{\rm imp}},
 \label{eq:v73-Dv-sigma}\\
 \|D_u\sigma\|_\kappa
 &\leq
 \frac{\beta\{a+\partial_r{\mathfrak n}(r_u,r_w)\}}
 {1-q_{\rm imp}}.
 \label{eq:v73-Du-sigma}
\end{align}
\end{proposition}

\begin{proof}
Factor
\[
 M=L\{I+L^{-1}D_wN\}
\]
and use the Neumann series to prove
\eqref{eq:v73-M-inverse}.  Implicit differentiation gives
\[
 D_v\sigma=M^{-1},
 \qquad
 D_u\sigma=-M^{-1}(A+D_uN).
\]
The scalar majorant bounds
\(\|D_uN\|_\kappa\) by
\(\partial_r{\mathfrak n}(r_u,r_w)\).  Insert
\eqref{eq:v73-M-inverse}.
\end{proof}

In the V72 penalty scaling \(v=\sqrt{\varepsilon}y\),
\begin{equation}
 D_y\sigma(u,\sqrt{\varepsilon}y)
 =
 \sqrt{\varepsilon}\,D_v\sigma(u,\sqrt{\varepsilon}y),
 \label{eq:v73-scaled-chart-jet}
\end{equation}
so every transverse chart response gains the expected
\(\sqrt{\varepsilon}\).

\subsection{Jacobian and determinant owner}

In the product coordinates \(x=(u,w)\), the derivative of
\(\Psi(u,v)=(u,\sigma(u,v))\) is block triangular.  Therefore
\begin{equation}
 J(u,v)
 =
 |\det D_v\sigma(u,v)|
 =
 |\det M(u,v)|^{-1}
 \label{eq:v73-Jacobian}
\end{equation}
inside a fixed orientation chart.

Put
\begin{equation}
 K_J:=L^{-1}D_wN(u,\sigma(u,v)).
 \label{eq:v73-KJ}
\end{equation}
Then \(\|K_J\|_\kappa\leq q_{\rm imp}<1\) and
\begin{equation}
 \log J
 =
 -\log|\det L|
 -\operatorname{Tr}\log(I+K_J).
 \label{eq:v73-log-J}
\end{equation}

\begin{theorem}[Anchored residual-Jacobian bound]
\label{thm:v73-Jacobian-bound}
At fixed regulator and in the chosen determinant-line orientation, the
nonconstant connected activity of \(\log J\) obeys
\begin{equation}
 \|\log J+\log|\det L|\|_{\kappa,\circ}
 \leq
 \nu_J\frac{q_{\rm imp}}{1-q_{\rm imp}},
 \label{eq:v73-Jacobian-activity}
\end{equation}
where \(\nu_J\) is the fixed local trace multiplicity.  Its retained and
residual derivatives obey the V53 Cauchy bounds on every smaller polydisc.
\end{theorem}

\begin{proof}
Equation \eqref{eq:v73-Neumann-condition} puts \(K_J\) in the V53
relative-determinant chamber.  Apply the anchored closed-word theorem to
\(-\operatorname{Tr}\log(I+K_J)\).  The constant
\(-\log|\det L|\) has no nonconstant polymer activity.  Normal convergence
and Cauchy's formula give the derivative statement.
\end{proof}

The constant determinant belongs to the base Gaussian/coarea normalization.
The relative determinant in \eqref{eq:v73-Jacobian-activity} belongs to the
chart/coarea sector.  It is not also charged through the V59 score except
when its derivative enters the compensation source.

\subsection{Pullback of the physical interaction}

Let \(H(x)\) be an analytic local density with scalar coefficient majorant
\({\mathfrak h}(t)\).  Then
\[
 H\circ\Psi(u,v)=H(u,\sigma(u,v))
\]
has coefficient majorant
\begin{equation}
 {\mathfrak h}\bigl(r_u+\widehat w(r_u,r_v)\bigr)
 \label{eq:v73-composition-majorant}
\end{equation}
after inserting the fixed coordinate-splitting constant into
\({\mathfrak h}\).

\begin{corollary}[Uniform V72 interaction stock]
\label{cor:v73-V72-stock}
Suppose \(\beta,a,{\mathfrak n},{\mathfrak h}\) and the support convolution
constants are uniform in regulator scale, and suppose
\eqref{eq:v73-self-map}--\eqref{eq:v73-contraction} hold with common strict
margins.  Then the pulled-back density
\[
 H(\Psi(u,\sqrt{\varepsilon}y))J(u,\sqrt{\varepsilon}y)
\]
has a uniform V72 anisotropic analytic norm for
\(\sqrt{\varepsilon}R_y\leq r_v\).  Its nonconstant residual dependence
vanishes at least as \(O(\sqrt{\varepsilon})\) in the coefficient-majorant
norm.
\end{corollary}

\begin{proof}
Proposition \ref{prop:v73-coefficient-majorant} and
\eqref{eq:v73-composition-majorant} bound the physical pullback.
Theorem \ref{thm:v73-Jacobian-bound} bounds the Jacobian activity.
Products are controlled by the V57 support--grade algebra.  Apply Theorem
\ref{thm:v72-anisotropic-bound}; its transverse-gain estimate supplies the
last assertion.
\end{proof}

\subsection{Insertion of the four-stage inverse}

Let \({\mathbb J}\) be the V63 derivative ordered by bulk scalar, DeTurck,
CMC, and interface variables.  On the declared complements, its exact
response formula gives a computable weighted bound
\begin{equation}
 \beta_{63}
 :=
 \|{\mathbb J}^{-1}\|_\kappa
 \label{eq:v73-beta63}
\end{equation}
as a polynomial in
\[
 \|L_\sigma^{-1}\|_\kappa,\quad
 \|\widehat A^{-1}\|_\kappa,\quad
 \|D^{-1}\|_\kappa,\quad
 \|\widehat S^{-1}\|_\kappa
\]
and the finitely many coupling-block norms.  The theorem above uses
\(\beta=\beta_{63}\).

\begin{ledger}[Radius-collapse diagnosis]
\label{led:v73-collapse}
The residual chart can lose its uniform radius only if at least one of the
following declared stocks degenerates:
\begin{enumerate}
\item the bulk scalar, hatted DeTurck, CMC, or interface inverse entering
      \(\beta_{63}\);
\item the linear retained coupling \(a\);
\item the nonlinear majorant \({\mathfrak n}\) or its \(w\)-derivative;
\item the support/grade convolution constants hidden in \(n_{pq}\); or
\item the determinant-line orientation or constant-rank complement.
\end{enumerate}
A failure of \eqref{eq:v73-contraction} is sent to the first upstream row
that degenerates.  It is not repaired by shrinking the RG step after the
complex chart radius has already collapsed.
\end{ledger}

The flat scalar threshold and locality loss of V64 are therefore seen
directly: as its pole distance tends to zero,
\(\beta_{63}\) grows, both inequalities
\eqref{eq:v73-self-map}--\eqref{eq:v73-contraction} fail, and the nonlinear
residual radius closes.

\subsection{Reflection and zero modes}

If \({\cal C}\) intertwines reflection, \(L^{-1}\) intertwines it on the
chosen complement, and the three radius balls are reflection invariant,
then uniqueness in Theorem \ref{thm:v73-residual-chart} gives
\begin{equation}
 \sigma(\Theta u,\Theta v)=\Theta\sigma(u,v).
 \label{eq:v73-reflection-chart}
\end{equation}
The Jacobian determinants on paired halves are conjugate in their declared
frames.

Every kernel mode of \(L\) must be removed before
\eqref{eq:v73-linear-stocks} is written.  Adding a pseudoinverse while
leaving its coordinate in \(w\) makes the contraction map nonunique and
invalidates the determinant formula.

\begin{firewall}[Uniform coefficients are not yet acquired]
\label{fire:v73-uniform-debit}
V73 proves that uniform V63 inverse and nonlinear coefficient stocks imply a
uniform exact residual chart.  It does not prove those stocks for the full
nonlinear SM--Einstein action.  In particular, the bulk scalar equation may
enter \(L\) only if its V62 physical-equivalence debit is resolved; the
canonical-constraint route must otherwise use only the Hamiltonian,
momentum, and legitimate gauge residuals.
\end{firewall}

\subsection{Next acquisition}

The majorant theorem reduces the chart problem to a finite list of
scale-normalized coefficients.  V74 will apply canonical dimensions to
\({\mathfrak n}\), separate relevant and marginal residual nonlinearities
from contracting ones, and formulate a scale-reset inequality that restores
the radius spent by \(\beta_{63}\).  This is the point at which a marginal
Einstein or matter-constraint coefficient must be entered into the running
owner bank rather than assumed small.

\begin{assessment}[Exact status after V73]
\label{ass:v73-status}
V73 proves a quantitative weighted analytic residual chart, an all-order
coefficient majorant, explicit first-jet and Jacobian activity bounds, and
a precise radius-collapse diagnosis tied to the V63 pivots.  Uniform
scale-normalized nonlinear stocks, a global canonical constraint cover,
physical determinant matching, large-field coercivity, regulator removal,
and the full SM--Einstein continuum remain open.
\end{assessment}
\section{Scale-normalized residual majorants and radius restoration}
\label{sec:v10v74}

V73 reduces the nonlinear residual-chart problem to the dimensionless stocks
\(\beta,a,{\mathfrak n}\).  This section proves a scale-reset theorem for
those stocks.  It combines the V18 canonical radius restoration with the
V54 projection--reblocking contact, while preserving a crucial
gravity-specific fact: extra powers of the dimensionless metric do not
create canonical contraction.

\subsection{Canonical residual grading}

At scale \(j\), let
\[
 {\cal C}_j(u,w)=L_jw+A_ju+N_j(u,w)
\]
be the normalized residual map.  The field and residual scaling maps are
denoted by
\[
 S_{u,j},\qquad S_{w,j},\qquad S_{v,j}.
\]
After one reblocking step, the canonically transported residual is
\begin{equation}
 {\cal C}_{j+1}^{\rm can}(u',w')
 :=
 S_{v,j}{\cal C}_j(S_{u,j}^{-1}u',S_{w,j}^{-1}w').
 \label{eq:v74-canonical-residual}
\end{equation}
Thus
\begin{align}
 L_{j+1}^{\rm can}
 &=
 S_{v,j}L_jS_{w,j}^{-1},
 \label{eq:v74-canonical-L}\\
 A_{j+1}^{\rm can}
 &=
 S_{v,j}A_jS_{u,j}^{-1}.
 \label{eq:v74-canonical-A}
\end{align}
The dimensionless inverse stock is
\begin{equation}
 \beta_j
 :=
 \|S_{w,j}L_j^{-1}S_{v,j}^{-1}\|_\kappa
 \label{eq:v74-dimensionless-beta}
\end{equation}
in the scale-\(j\) normalized norm.  This is the quantity that must remain
bounded; an unscaled inverse norm can carry a harmless canonical power.

Choose a covariant local residual-owner basis modulo Bianchi identities,
integrations by parts, redundant field redefinitions, and the retained
kernel complex.  Split the nonlinear residual exactly as
\begin{equation}
 N_j=N_j^{\rm run}+E_j,
 \qquad
 N_j^{\rm run}:={\cal P}_jN_j,
 \qquad
 E_j:=(I-{\cal P}_j)N_j.
 \label{eq:v74-residual-split}
\end{equation}
Here \({\cal P}_j\) contains every residual owner with nonnegative canonical
scaling exponent.  The complementary owners have a strictly positive excess
dimension at least \(\sigma>0\), so their canonical transport is bounded by
\begin{equation}
 \lambda_{\rm can}:={\mathscr L}^{-\sigma}<1
 \label{eq:v74-canonical-lambda}
\end{equation}
for block factor \({\mathscr L}>1\).

\begin{firewall}[Metric powers remain in the running owner]
\label{fire:v74-metric-powers}
The metric has canonical dimension zero.  Therefore the full analytic metric
dependence multiplying a fixed covariant tensor, such as the Einstein
tensor, a stress tensor, or a curvature-squared variation, has the same
canonical exponent as that tensor.  It must be resummed inside
\(N_j^{\rm run}\).  Sorting high metric Taylor powers into \(E_j\) would
produce a false contraction and repeat the error excluded at V57.
\end{firewall}

The running residual bank includes the cosmological and
Einstein--Hilbert rows, all renormalizable Standard Model stress and current
rows, curvature-squared variations generated by the heat-shell ledger,
legitimate gauge and constraint contacts, and their complete covariant
metric dependence.  Topological bulk variations that vanish identically
remain boundary or determinant-line owners rather than zero estimates.

\subsection{Exact scale contact}

Let \({\mathscr R}_j\) be the nonlinear reblocking map on residual
activities.  Define the exact contact
\begin{equation}
 {\cal K}_jN_j
 :=
 {\cal P}_{j+1}{\mathscr R}_jN_j
 -
 {\mathscr R}_j{\cal P}_jN_j.
 \label{eq:v74-residual-contact}
\end{equation}
The same algebra as V54 gives
\begin{align}
 {\cal P}_{j+1}{\mathscr R}_jN_j
 &=
 {\mathscr R}_jN_j^{\rm run}+{\cal K}_jN_j,
 \label{eq:v74-running-contact}\\
 (I-{\cal P}_{j+1}){\mathscr R}_jN_j
 &=
 {\mathscr R}_jE_j-{\cal K}_jN_j.
 \label{eq:v74-tail-contact}
\end{align}
Thus the contact is added to the running residual and removed from the tail.
It cannot be discarded or charged to both.

Let \({\mathfrak r}_j(r,t)\),
\({\mathfrak e}_j(r,t)\), and
\({\mathfrak c}_j(r,t)\) be nonnegative coefficient majorants for
\(N_j^{\rm run}\), \(E_j\), and the irrelevant part of the contact and other
new tail sources, respectively.  Assume the exact tail update implies the
coefficientwise bound
\begin{equation}
 {\mathfrak e}_{j+1}
 \preceq
 \lambda_{\rm can}{\mathfrak e}_j+{\mathfrak c}_j.
 \label{eq:v74-majorant-recurrence}
\end{equation}
The notation \(\preceq\) means coefficientwise domination of power series
with nonnegative coefficients.

\subsection{Invariant tail envelope}

\begin{lemma}[Scale-uniform tail envelope]
\label{lem:v74-tail-envelope}
Suppose there is a nonnegative analytic majorant
\({\mathfrak e}_*\) such that
\begin{equation}
 {\mathfrak e}_0\preceq{\mathfrak e}_*,
 \qquad
 {\mathfrak c}_j
 \preceq
 (1-\lambda_{\rm can}){\mathfrak e}_*
 \quad\hbox{for every }j.
 \label{eq:v74-envelope-hypothesis}
\end{equation}
Then
\begin{equation}
 {\mathfrak e}_j\preceq{\mathfrak e}_*
 \quad\hbox{for every }j.
 \label{eq:v74-envelope-conclusion}
\end{equation}
More generally,
\begin{equation}
 {\mathfrak e}_j
 \preceq
 \lambda_{\rm can}^j{\mathfrak e}_0
 +
 \sum_{k=0}^{j-1}
 \lambda_{\rm can}^{j-1-k}{\mathfrak c}_k.
 \label{eq:v74-Duhamel-majorant}
\end{equation}
\end{lemma}

\begin{proof}
Iterating \eqref{eq:v74-majorant-recurrence} gives
\eqref{eq:v74-Duhamel-majorant}.  Under
\eqref{eq:v74-envelope-hypothesis}, the right side is bounded by
\[
 \lambda_{\rm can}^j{\mathfrak e}_*
 +(1-\lambda_{\rm can})
 \sum_{k=0}^{j-1}\lambda_{\rm can}^{j-1-k}{\mathfrak e}_*
 ={\mathfrak e}_*.
\]
\end{proof}

The source \({\mathfrak c}_j\) includes the irrelevant portion of the V54
contact, shell integration errors, and compensation tails.  Every
noncontracting component of those objects is removed into the running bank
before \eqref{eq:v74-majorant-recurrence} is asserted.

\subsection{Uniform residual-chart theorem}

Assume dimensionless bounds
\begin{equation}
 \beta_j\leq\beta_*,
 \qquad
 a_j:=\|A_j\|_\kappa\leq a_*,
 \qquad
 {\mathfrak r}_j\preceq{\mathfrak r}_*
 \label{eq:v74-uniform-linear-running}
\end{equation}
and the tail envelope of Lemma \ref{lem:v74-tail-envelope}.  Put
\begin{equation}
 {\mathfrak n}_*:={\mathfrak r}_*+{\mathfrak e}_*.
 \label{eq:v74-total-envelope}
\end{equation}

\begin{theorem}[Scale-reset residual-radius cycle]
\label{thm:v74-radius-cycle}
Suppose fixed positive radii \(r_u,r_v,r_w\) obey
\begin{align}
 \beta_*
 \{r_v+a_*r_u+{\mathfrak n}_*(r_u,r_w)\}
 &\leq r_w,
 \label{eq:v74-uniform-self-map}\\
 \beta_*\partial_t{\mathfrak n}_*(r_u,r_w)
 &\leq1-\delta_{\rm imp}
 \label{eq:v74-uniform-contraction}
\end{align}
for some \(\delta_{\rm imp}>0\).  Then every scale \(j\) has an exact
analytic residual chart
\[
 w=\sigma_j(u,v)
\]
on the common polydisc
\(\|u\|\leq r_u,\|v\|\leq r_v\), with
\[
 \|\sigma_j\|\leq r_w.
\]
Its first jets obey the uniform bounds
\begin{align}
 \|D_v\sigma_j\|_\kappa
 &\leq
 \frac{\beta_*}{\delta_{\rm imp}},
 \label{eq:v74-uniform-Dv}\\
 \|D_u\sigma_j\|_\kappa
 &\leq
 \frac{\beta_*
 \{a_*+\partial_r{\mathfrak n}_*(r_u,r_w)\}}
 {\delta_{\rm imp}}.
 \label{eq:v74-uniform-Du}
\end{align}
The relative chart-Jacobian activity is uniformly bounded by
\begin{equation}
 \nu_J
 \frac{1-\delta_{\rm imp}}{\delta_{\rm imp}}.
 \label{eq:v74-uniform-Jacobian}
\end{equation}
\end{theorem}

\begin{proof}
Equations \eqref{eq:v74-uniform-linear-running} and
\eqref{eq:v74-envelope-conclusion} imply that the scale-\(j\) V73
majorant is bounded coefficientwise by \({\mathfrak n}_*\).
Hence \eqref{eq:v74-uniform-self-map} and
\eqref{eq:v74-uniform-contraction} imply the two V73 radius conditions at
every scale.  Apply Theorem \ref{thm:v73-residual-chart}.
The V73 first-jet estimates have denominator
\(1-q_{{\rm imp},j}\geq\delta_{\rm imp}\), proving
\eqref{eq:v74-uniform-Dv}--\eqref{eq:v74-uniform-Du}.
Finally, the V73 determinant bound has
\(q_{{\rm imp},j}\leq1-\delta_{\rm imp}\), which gives
\eqref{eq:v74-uniform-Jacobian}.
\end{proof}

This theorem restores the residual-chart radius after each exact
projection and canonical reblocking.  It does not require the running
coefficients to contract.  It requires their trajectory to remain in the
compact chamber represented by
\(\beta_*,a_*,{\mathfrak r}_*\).

\subsection{Strict reserve and perturbations}

Define the self-map and derivative reserves
\begin{align}
 \Delta_{\rm map}
 &:=
 r_w-\beta_*
 \{r_v+a_*r_u+{\mathfrak n}_*(r_u,r_w)\},
 \label{eq:v74-map-reserve}\\
 \Delta_{\rm der}
 &:=
 1-\beta_*\partial_t{\mathfrak n}_*(r_u,r_w).
 \label{eq:v74-derivative-reserve}
\end{align}
A useful chamber has both reserves positive.  If a new shell changes the
dimensionless stocks by
\(\delta\beta,\delta a,\delta{\mathfrak n}\), it remains in the chamber
whenever
\begin{align}
 \delta\beta
 \{r_v+a_*r_u+{\mathfrak n}_*(r_u,r_w)\}
 +
 \beta_*\{\delta a\,r_u+
 \delta{\mathfrak n}(r_u,r_w)\}
 &<
 \Delta_{\rm map},
 \label{eq:v74-map-perturbation}\\
 \delta\beta\,\partial_t{\mathfrak n}_*(r_u,r_w)
 +
 \beta_*\partial_t\delta{\mathfrak n}(r_u,r_w)
 &<
 \Delta_{\rm der}.
 \label{eq:v74-derivative-perturbation}
\end{align}
These inequalities allocate rather than hide the scale-reset cost.

\subsection{Running rows versus contracting rows}

\begin{proposition}[Necessary extraction of noncontracting residuals]
\label{prop:v74-necessary-extraction}
Suppose a residual coefficient has canonical scaling factor
\({\mathscr L}^{\omega}\) with \(\omega\geq0\).  If it is left in \(E_j\),
no inequality of the form
\[
 \|E_{j+1}\|\leq\lambda\|E_j\|+\hbox{summable source},
 \qquad\lambda<1,
\]
can follow from canonical scaling alone.  It must be inserted into
\(N_j^{\rm run}\), cancelled by an exact identity, or controlled by a
separate noncanonical mechanism.
\end{proposition}

\begin{proof}
Under repeated canonical transport, the coefficient is multiplied by
\({\mathscr L}^{k\omega}\), which is constant when \(\omega=0\) and grows
when \(\omega>0\).  A strict contraction factor cannot dominate that
sequence without an additional estimate.  Projection into the running bank
removes it from the assertion.
\end{proof}

For four-dimensional SM--Einstein theory the noncontracting action rows
already identified at V39 include
\[
 \int\sqrt g,\quad
 \int\sqrt g\,R,\quad
 \int\sqrt g\,R^2,\quad
 \int\sqrt g\,R_{\mu\nu}R^{\mu\nu},
\]
the renormalizable gauge, Higgs, Yukawa, and fermion kinetic rows, plus
topological, gauge, boundary, and constraint owners.  Their variations with
respect to eliminated fields populate \(N_j^{\rm run}\).  The list is a
necessary bank, not a proof that its beta functions possess the required
trajectory.

\subsection{Combination with determinant scale reset}

Let the determinant tail at scale \(j\) have relative norm
\(\eta_j<1\) after all running local rows are projected.  V54 requires
\[
 \sum_jC_{j\infty}\nu_j\frac{\eta_j}{1-\eta_j}<\infty.
\]
The residual chart adds the uniform Jacobian activity
\eqref{eq:v74-uniform-Jacobian}, but its noncontracting projection is again
part of the running bank.  Only its irrelevant remainder enters the
summability condition.

\begin{ledger}[One scale-reset order]
\label{led:v74-scale-order}
At each scale the operations occur in this order:
\begin{enumerate}
\item solve and orient the V63 pivot complement;
\item construct the exact residual chart by V73;
\item compute determinant, score, and interaction activities in that chart;
\item project every noncontracting covariant row;
\item apply the exact V54 contact identity;
\item canonically reblock the positive-excess residual and determinant
      tails; and
\item test \eqref{eq:v74-envelope-hypothesis} and the strict reserves
      \eqref{eq:v74-map-reserve}--\eqref{eq:v74-derivative-reserve}.
\end{enumerate}
Changing this order can either double count the Jacobian or apply canonical
contraction to a marginal row.
\end{ledger}

\subsection{The remaining trajectory problem}

Theorem \ref{thm:v74-radius-cycle} is uniform if the running sequence stays
in a compact chamber with
\[
 \beta_j\leq\beta_*,
 \quad
 a_j\leq a_*,
 \quad
 {\mathfrak r}_j\preceq{\mathfrak r}_*.
\]
Constructing such a sequence is the finite-dimensional shooting problem for
all relevant and marginal SM--Einstein owners.  In particular, a
non-Gaussian ultraviolet fixed point or another ultraviolet mechanism is
needed if the dimensionless Newton and curvature couplings do not remain in
this chamber.  Canonical power counting does not prove that mechanism.

\begin{firewall}[Scale restoration is conditional on a running chamber]
\label{fire:v74-running-chamber}
V74 proves that a compact running chamber plus a contracting residual tail
gives a common nonlinear chart.  It does not prove that the SM--Einstein
beta functions enter or remain in that chamber.  Nor does it resolve the
V62 physical status of a four-dimensional bulk trace residual.
\end{firewall}

\subsection{Next acquisition}

V75 is the next compulsory companion audit.  It will inspect the newest
Yang--Mills and Navier--Stokes notes for a concrete marginal shooting
argument, an invariant envelope recurrence, or a sharp way to preserve a
strict radius reserve.  The post-audit note will then formulate the
finite-dimensional running-coupling shooting map associated with
\eqref{eq:v74-uniform-linear-running}, including its stable and unstable
directions.

\begin{assessment}[Exact status after V74]
\label{ass:v74-status}
V74 proves a scale-normalized residual grading, an exact contact recurrence,
a uniform tail envelope, and a common-radius theorem with explicit reserves.
It also proves that noncontracting metric and SM--Einstein residual rows must
be retained.  The required running-coupling trajectory, uniform V63 pivots,
nonlinear large-field measure, regulator removal, and the full
SM--Einstein continuum remain open.
\end{assessment}
\section{Companion audit: global rate matching before shooting}
\label{sec:v10v75}

This is the compulsory audit following V70.  The newest active companion
files were inspected read-only:
\begin{align*}
 &\texttt{Renewal\_Yang\_Mills\_Mass\_Gap\_Research\_Notes\_V4.tex},\\
 &\qquad
 \operatorname{SHA256}
 =\texttt{B581452CFFBA4B2D8305C7B6E81E3AF5BFC59B3C0AF94738DF894139E94809A2},\\
 &\texttt{Renewal\_NS\_Stability\_Research\_Notes\_V26.tex},\\
 &\qquad
 \operatorname{SHA256}
 =\texttt{82C887F8C2C69E4F28FAD7C3DC8DE5B9099CB5A4BF431EBA1FFF3E91935978AD}.
\end{align*}
The Yang--Mills companion advanced from compact-cycle V1 to V4.  The
Navier--Stokes companion is unchanged.

\subsection{Transferable obstruction}

The Yang--Mills note compares a physical rate function with the squared
distance defining its heat-kernel reference.  It proves that agreement of
their Hessians at the vacuum gives no global sign for the difference of the
two rates.  If the reference rate exceeds the physical rate somewhere by a
fixed positive amount, the normalized density ratio grows exponentially in
the inverse small parameter, so it cannot be a small analytic activity even
in zeroth order.

The same logical obstruction applies to V69--V74.  Those notes construct a
positive local Gaussian reference, a constraint residual chart, and
uniform local analytic estimates.  They do not prove that the globally
reduced SM--Einstein density is dominated by that reference.  A local
residual-chart radius cannot control configurations outside its domain.

\begin{proposition}[Local quadratic matching is insufficient]
\label{prop:v75-local-not-global}
Let \(V,I\) be continuous nonnegative functions on a configuration space,
with a common minimum \(x_0\), and suppose
\[
 V(x)=I(x)+O(\|x-x_0\|^3)
\]
near \(x_0\).  This local relation does not imply \(V\geq I\) globally.
If there is a point \(x_1\) with
\begin{equation}
 I(x_1)-V(x_1)=\Delta>0,
 \label{eq:v75-positive-rate-defect}
\end{equation}
then the unnormalized ratio
\[
 e^{-V(x_1)/g^2}/e^{-I(x_1)/g^2}
\]
equals \(e^{\Delta/g^2}\) and is exponentially large as \(g\downarrow0\).
\end{proposition}

\begin{proof}
The first assertion follows because one may add to either function a
continuous term supported outside a neighborhood of \(x_0\), without
changing its local expansion.  The second assertion is the displayed
identity.
\end{proof}

Normalization constants can shift the rate by a configuration-independent
number, but cannot remove variation of a positive rate defect across
configurations.  V76 will include them explicitly.

\subsection{What transfers from the covariance path}

The Yang--Mills companion broadens its compact heat-kernel covariance until
the reference rate is globally dominated, then transports the order-one
quadratic mismatch back to the physical covariance by an exact semigroup
identity.  Three parts of this architecture transfer:

\begin{enumerate}
\item a marginal quadratic covariance mismatch is transported exactly
      rather than inserted as a small polymer vertex;
\item the covariance interpolation has its own uniform edge stock; and
\item the Gaussian normalization determinant is a separate vacuum or
      retained-background owner.
\end{enumerate}

On a Euclidean vector space the transferred identity is the standard
Gaussian formula
\begin{equation}
 \frac{d}{dt}{\mathbb E}_{C_t}F
 =
 \frac12{\mathbb E}_{C_t}
 \left[
  \sum_{a,b}\dot C_t^{ab}\partial_a\partial_bF
 \right].
 \label{eq:v75-Gaussian-path-preview}
\end{equation}
It needs no compact group.  V76 will prove it with the determinant and
weighted locality ledger.

The compact Yang--Mills proof that a sufficiently broad reference is
globally rate dominated does not transfer.  It uses compactness, a finite
diameter, and a physical action with a unique zero.  Reduced gravitational
and Higgs variables are noncompact, and the Euclidean conformal obstruction
can destroy coercivity.  A gravitational theorem must assume or prove a
proper many-block lower bound after legitimate canonical constraint
reduction.

\subsection{Navier--Stokes audit result}

The NS V26 rank-one Oseen-kernel estimates remain unrelated to the present
rate-comparison problem.  No theorem there supplies a constrained
gravitational rate, Gaussian covariance path, or marginal shooting
trajectory.  No NS constant is imported.

\begin{ledger}[V75 transfer decision]
\label{led:v75-transfer}
\begin{enumerate}
\item Import the global rate-defect test.
\item Import exact covariance transport and its separate normalization
      determinant.
\item Do not import compact-diameter, Haar, Wilson-action, or Yang--Mills
      coupling estimates.
\item Retain V74's scale-reset theorem, but delay the running-coupling
      shooting construction until the reference comparison has a valid
      global branch.
\end{enumerate}
\end{ledger}

\subsection{Post-audit direction}

V76 will formulate a noncompact constrained rate-comparison theorem in the
V72 residual chart.  It will:
\begin{enumerate}
\item compute the normalized small-parameter density-ratio rate, including
      partition functions;
\item prove an exponential obstruction when the rate defect is positive;
\item give a sufficient coercive domination condition for a broad positive
      Gaussian reference;
\item derive the exact Gaussian covariance-path identity and determinant
      owner; and
\item state the reflection and locality requirements along the path.
\end{enumerate}
Only after this branch is available can the marginal shooting map be tested
against a globally meaningful reference law.

\begin{assessment}[Exact status after V75]
\label{ass:v75-status}
V75 finds a missing global reference-comparison row upstream of the V74
shooting problem and imports the exact covariance-transport architecture.
It does not prove gravitational rate domination, a global constrained
measure, the marginal trajectory, regulator removal, or the full
SM--Einstein continuum.
\end{assessment}
\section{Noncompact constrained rate matching and covariance transport}
\label{sec:v10v76}

V75 showed that the local residual chart and Hessian-matched Gaussian do not
control the global reduced density.  This section supplies the missing
finite-dimensional comparison theorem.  It applies after a legitimate
canonical constraint reduction and before the local forest expansion.

\subsection{Reduced target and Gaussian reference}

Let \(u\in{\mathbb R}^d\) denote the retained bosonic coordinates after the
V72 residual \(v\) has been set to zero.  Define the reduced rate
\begin{equation}
 V(u):=S(\Psi(u,0))-S(\Psi(0,0)),
 \qquad V(0)=0,
 \label{eq:v76-reduced-rate}
\end{equation}
and let \(J_g(u)>0\) contain the positive coarea, orbit, and magnitude
density owners.  The normalized target is
\begin{equation}
 d\mu_g(u)
 =
 Z_V(g)^{-1}J_g(u)e^{-V(u)/g^2}\,du.
 \label{eq:v76-target-law}
\end{equation}
A sign or chiral phase is not part of \(J_g\); it remains in its determinant
line.

Let \(C>0\) and
\begin{equation}
 I_C(u):=\frac12\langle u,C^{-1}u\rangle.
 \label{eq:v76-Gaussian-rate}
\end{equation}
For a broadening factor \(\varsigma\geq1\), put
\begin{equation}
 I_{\varsigma C}(u)=\varsigma^{-1}I_C(u)
 \label{eq:v76-broad-rate}
\end{equation}
and let \(\nu_{g,\varsigma C}\) be the centered normalized Gaussian with
covariance \(g^2\varsigma C\).

Assume the partition functions have subexponential rates:
\begin{equation}
 \lim_{g\downarrow0}g^2\log Z_V(g)=0,
 \qquad
 \lim_{g\downarrow0}g^2\log Z_{\varsigma C}(g)=0.
 \label{eq:v76-partition-rate}
\end{equation}
For a finite-dimensional nondegenerate minimum these statements follow, for
example, from two-sided polynomial bounds
\(c g^M\leq Z\leq Cg^{-M}\).  Assume also
\begin{equation}
 \lim_{g\downarrow0}g^2\log J_g(u)=0
 \label{eq:v76-density-subexponential}
\end{equation}
at every fixed \(u\), locally uniformly on the chart.

\begin{theorem}[Normalized density-ratio rate]
\label{thm:v76-ratio-rate}
The Radon--Nikodym ratio
\begin{equation}
 {\cal W}_{g,\varsigma}(u)
 :=
 \frac{d\mu_g}{d\nu_{g,\varsigma C}}(u)
 \label{eq:v76-density-ratio}
\end{equation}
satisfies
\begin{equation}
 \lim_{g\downarrow0}
 g^2\log{\cal W}_{g,\varsigma}(u)
 =
 I_{\varsigma C}(u)-V(u)
 \label{eq:v76-rate-limit}
\end{equation}
locally uniformly wherever
\eqref{eq:v76-density-subexponential} is locally uniform.
\end{theorem}

\begin{proof}
The normalized Gaussian density is
\[
 Z_{\varsigma C}(g)^{-1}e^{-I_{\varsigma C}(u)/g^2}.
\]
Therefore
\[
 \log{\cal W}_{g,\varsigma}
 =
 \log\frac{Z_{\varsigma C}(g)}{Z_V(g)}
 +\log J_g(u)
 +\frac{I_{\varsigma C}(u)-V(u)}{g^2}.
\]
Multiply by \(g^2\) and use
\eqref{eq:v76-partition-rate}--\eqref{eq:v76-density-subexponential}.
\end{proof}

\begin{corollary}[Exponential global obstruction]
\label{cor:v76-rate-obstruction}
If there is \(u_*\) such that
\begin{equation}
 I_{\varsigma C}(u_*)-V(u_*)=\Delta>0
 \label{eq:v76-positive-defect}
\end{equation}
and \(J_g\) has zero rate there, then
\begin{equation}
 \|{\cal W}_{g,\varsigma}-1\|_\infty
 \geq
 \exp\{(\Delta+o(1))/g^2\}-1.
 \label{eq:v76-exponential-obstruction}
\end{equation}
Thus the global density ratio cannot be a small activity even in zeroth
analytic order.
\end{corollary}

\begin{proof}
Evaluate \eqref{eq:v76-rate-limit} at \(u_*\) and use the definition of the
supremum norm.
\end{proof}

This obstruction survives exact local Hessian matching at \(u=0\).

\subsection{When a broad Gaussian is globally dominated}

Define
\begin{equation}
 \sigma_*:=
 \sup_{u\neq0}\frac{I_C(u)}{V(u)}
 \label{eq:v76-sigma-star}
\end{equation}
with value \(+\infty\) if \(V(u)=0<I_C(u)\) anywhere.

\begin{lemma}[A noncompact finiteness criterion]
\label{lem:v76-sigma-finite}
Suppose:
\begin{enumerate}
\item \(V\) is continuous, \(V(u)>0\) for \(u\neq0\);
\item near zero,
      \begin{equation}
       V(u)\geq c_0 I_C(u)
       \label{eq:v76-local-quadratic-floor}
      \end{equation}
      for some \(c_0>0\); and
\item outside a compact set,
      \begin{equation}
       V(u)\geq c_\infty I_C(u)
       \label{eq:v76-infinity-quadratic-floor}
      \end{equation}
      for some \(c_\infty>0\).
\end{enumerate}
Then \(\sigma_*<\infty\).
\end{lemma}

\begin{proof}
The quotient is bounded by \(c_0^{-1}\) near zero and by
\(c_\infty^{-1}\) outside the chosen compact set.  On the remaining compact
annulus, \(V\) has a positive minimum and \(I_C\) has a finite maximum, so
the quotient is bounded there as well.
\end{proof}

\begin{theorem}[Globally rate-dominated broad reference]
\label{thm:v76-broad-domination}
If \(\sigma_*<\infty\), then every
\(\varsigma>\sigma_*\) satisfies
\begin{equation}
 V(u)-I_{\varsigma C}(u)
 \geq
 \left(1-\frac{\sigma_*}{\varsigma}\right)V(u)
 \geq0.
 \label{eq:v76-global-rate-floor}
\end{equation}
The rate in \eqref{eq:v76-rate-limit} is therefore nonpositive everywhere.
\end{theorem}

\begin{proof}
Definition \eqref{eq:v76-sigma-star} gives
\(I_C(u)\leq\sigma_*V(u)\).  Divide by \(\varsigma\) and rearrange.
\end{proof}

The theorem prevents an exponentially growing likelihood ratio.  It does
not make the ratio close to one: the broad Gaussian has an order-one
quadratic mismatch near the minimum.

For a family of finite block sets \(\Lambda\), the useful hypothesis is the
volume-uniform many-block inequality
\begin{equation}
 I_{C,\Lambda}(u)
 \leq
 \sigma_*V_\Lambda(u)
 \quad\hbox{with one }\sigma_*<\infty.
 \label{eq:v76-uniform-many-block}
\end{equation}
A constant growing with \(|\Lambda|\) does not yield a polymer estimate.
Negative interface energies and all boundary owners must be included in
\(V_\Lambda\) before \eqref{eq:v76-uniform-many-block} is tested.

\subsection{A sufficient partition-rate lemma}

\begin{proposition}[Coercive zero partition rate]
\label{prop:v76-partition-rate}
Suppose, uniformly for \(0<g\leq g_0\),
\begin{align}
 V(u)&\geq a\|u\|^p-b,
 \label{eq:v76-coercive-rate}\\
 0<J_g(u)&\leq
 C\exp\{c\|u\|^q\},
 \qquad q<p,
 \label{eq:v76-density-growth}
\end{align}
and near a minimum there is a ball of radius \(g^\alpha\) on which
\[
 V(u)\leq C_0g^2,\qquad
 J_g(u)\geq c_0g^M
\]
for fixed \(\alpha,M\).  If the additive constant is chosen so that
\(\inf V=0\), then
\[
 \lim_{g\downarrow0}g^2\log Z_V(g)=0.
\]
\end{proposition}

\begin{proof}
The lower bound follows by integrating over the stated ball:
\[
 Z_V(g)\geq
 c'g^{M+\alpha d}e^{-C_0},
\]
whose logarithm has zero \(g^2\)-rate.  For the upper bound, split at a
fixed radius large enough that
\(c\|u\|^q\leq a\|u\|^p/(2g_0^2)\).  The compact part is at most
\(C'g^{-M'}\) after enlarging a harmless polynomial constant.  On the tail,
\[
 J_g(u)e^{-V(u)/g^2}
 \leq
 C e^{b/g^2}e^{-a\|u\|^p/(2g^2)}.
\]
After the normalization \(\inf V=0\), refine the compact/tail split using
the actual sublevel sets of \(V\); Laplace's elementary upper bound gives
\(\limsup g^2\log Z_V\leq0\).  The lower bound gives the reverse inequality.
\end{proof}

The role of \eqref{eq:v76-coercive-rate} is exactly the V22 many-block
large-field debit.  The proposition does not establish it for gravity.

\subsection{Exact covariance transport}

Let \(C_t\), \(0\leq t\leq1\), be a \(C^1\) path of positive covariance
matrices.  Denote expectation for the centered Gaussian with covariance
\(g^2C_t\) by \({\mathbb E}_{g,C_t}\).

\begin{theorem}[Gaussian covariance-path identity]
\label{thm:v76-covariance-path}
For every smooth \(F\) whose first two derivatives have sufficient Gaussian
growth,
\begin{equation}
 \frac{d}{dt}{\mathbb E}_{g,C_t}F
 =
 \frac{g^2}{2}
 {\mathbb E}_{g,C_t}
 \left[
  \operatorname{Tr}\{\dot C_tD^2F\}
 \right].
 \label{eq:v76-covariance-derivative}
\end{equation}
Consequently
\begin{equation}
 {\mathbb E}_{g,C_1}F-{\mathbb E}_{g,C_0}F
 =
 \frac{g^2}{2}\int_0^1
 {\mathbb E}_{g,C_t}
 \left[
  \operatorname{Tr}\{\dot C_tD^2F\}
 \right]dt.
 \label{eq:v76-covariance-Duhamel}
\end{equation}
\end{theorem}

\begin{proof}
It suffices first to take
\(F(u)=e^{i\langle k,u\rangle}\).  Its expectation is
\[
 \exp\{-\tfrac12g^2\langle k,C_tk\rangle\}.
\]
Differentiation gives
\[
 -\frac{g^2}{2}\langle k,\dot C_tk\rangle
 e^{-\frac12g^2\langle k,C_tk\rangle},
\]
which equals the expectation of the right side of
\eqref{eq:v76-covariance-derivative}.  Fourier approximation and Gaussian
domination extend the identity to the stated class.  Integration in \(t\)
proves \eqref{eq:v76-covariance-Duhamel}.
\end{proof}

For the scalar path
\begin{equation}
 C_t=\varsigma_tC,
 \qquad
 \varsigma_t=\varsigma-(\varsigma-1)t,
 \label{eq:v76-scalar-path}
\end{equation}
the broad reference is at \(t=0\) and the Hessian-matched reference at
\(t=1\).  Equation \eqref{eq:v76-covariance-Duhamel} transports their full
quadratic mismatch.  It is not a small interaction vertex.

The Gaussian normalization is
\[
 Z_C(g)=(2\pi g^2)^{d/2}(\det C)^{1/2}.
\]
Hence the path carries the exact determinant owner
\begin{equation}
 \log\frac{Z_{\varsigma C}(g)}{Z_C(g)}
 =
 \frac d2\log\varsigma.
 \label{eq:v76-path-determinant}
\end{equation}
If \(C=C(w)\) depends on retained data, derivatives of
\eqref{eq:v76-path-determinant} belong to the bosonic Gaussian determinant
sector of V52--V54.  The constant is not added to the small connected
activity.

\subsection{Weighted locality and reflection along the path}

Let
\[
 G_\kappa(C_t)
 :=
 \sup_x\sum_y e^{\kappa d(x,y)}\|(C_t)_{xy}\|.
\]
For \eqref{eq:v76-scalar-path},
\begin{equation}
 G_\kappa(C_t)
 =
 \varsigma_tG_\kappa(C)
 \leq
 \varsigma G_\kappa(C).
 \label{eq:v76-path-row}
\end{equation}
If all incident analytic derivatives spend one radius gap
\(\delta_{\rm an}\), the covariance-path insertion has edge stock bounded by
\begin{equation}
 K_{\rm path}
 :=
 \frac{g^2(\varsigma-1)h_*G_\kappa(C)}
 {2\delta_{\rm an}^2}.
 \label{eq:v76-path-edge-stock}
\end{equation}
The degree-weighted forest constant of V70 must be used if several such
insertions join one analytic vertex.

\begin{proposition}[Reflection-positive covariance path]
\label{prop:v76-path-reflection}
Suppose every \(C_t\) intertwines the reflection and its reflected covariance
form is nonnegative on positive-half test functions.  Then every Gaussian
on the path is reflection positive.  In particular, positive scalar
multiples of the V68--V69 TT covariance satisfy this hypothesis.
\end{proposition}

\begin{proof}
Gaussian reflection positivity is equivalent to nonnegativity of the
reflected covariance quadratic form, followed by the standard Gaussian
exponential completion.  This is assumed for every \(t\).  A positive scalar
multiple preserves the sign of the quadratic form.
\end{proof}

A general interpolation between two positive matrices need not preserve the
specific time-reflection structure merely because both endpoints are
positive as ordinary matrices.  The reflected covariance condition is
tested along the whole path.

\subsection{Reference-comparison card}

\begin{definition}[Constrained global-reference card]
\label{def:v76-reference-card}
A finite regulated block family passes the card if:
\begin{enumerate}
\item the physical residual is a legitimate canonical constraint/gauge map
      and the V72 coarea limit gives a positive reduced bosonic density;
\item \(V_\Lambda\) and \(J_{g,\Lambda}\) satisfy the uniform coercive and
      subexponential partition-rate hypotheses;
\item one Hessian covariance \(C_\Lambda\) obeys the uniform many-block rate
      comparison \eqref{eq:v76-uniform-many-block};
\item a fixed \(\varsigma>\sigma_*\) leaves enough weighted covariance and
      analytic-radius reserve for the V70 forest gate;
\item the scalar covariance path has uniform locality and reflection
      positivity, and its determinant is entered once; and
\item the Hessian-matched physical ratio is used only on the local analytic
      chart, while its complement is paid by the many-block coercive
      large-field owner.
\end{enumerate}
\end{definition}

The broad reference and the V71 residual-selector rebase have different
roles.  Broadening \(C\) controls the global retained-field rate and worsens
the covariance row by \(\varsigma\).  Strengthening the residual Gaussian
controls analytic soft marks and narrows its covariance by \(1+2s\).  They
belong to different coordinate sectors and are not algebraically cancelled.

\begin{firewall}[The global gravitational floor remains unproved]
\label{fire:v76-gravity-floor}
No previous note proves
\eqref{eq:v76-local-quadratic-floor} and
\eqref{eq:v76-infinity-quadratic-floor} for the complete nonlinear
constraint-reduced SM--Einstein action.  V23 supplies a scalar constraint
sign chamber, and V68 supplies the linear TT quadratic floor.  Neither is a
many-block nonlinear coercivity theorem for all retained gravitational and
matter variables.
\end{firewall}

\subsection{Next acquisition}

The global-reference card now separates the rate branch from the local
forest and residual-radius branches.  V77 will return to the finite
running-coupling shooting problem of V74.  It will formulate a graph
transform around a candidate scale orbit, with the global-reference,
pivot, forest, determinant, and tail reserves included in the admissible
parameter chamber.  A hyperbolicity or cone condition will be stated as a
hypothesis and traced to the actual linearized beta map; it will not be
inferred from power counting alone.

\begin{assessment}[Exact status after V76]
\label{ass:v76-status}
V76 proves the normalized constrained rate formula, an exponential mismatch
obstruction, a sufficient noncompact broad-reference criterion, exact
Gaussian covariance transport, and its determinant, locality, and
reflection ledgers.  The uniform gravitational rate floor, the marginal
running orbit, nonlinear OS positivity, regulator removal, and the full
SM--Einstein continuum remain open.
\end{assessment}
\section{Running-orbit graph transform with admissible reserves}
\label{sec:v10v77}

V16 proved a constant-coefficient relevant--irrelevant shooting theorem and
a scalar marginal estimate.  V74 and V76 add two requirements not present
there: the linear splitting can vary along a running orbit, and the orbit
must remain inside the simultaneous pivot, global-rate, forest,
determinant, tail, and residual-chart chambers.  This section proves the
corresponding graph-transform card.  It also extends the scalar marginal
estimate to a coupled cone.

\subsection{Candidate orbit and linear splitting}

Let \(g_j^*\) be a candidate trajectory in the projected running-owner space:
\begin{equation}
 g_{j+1}^*={\cal B}_j(g_j^*).
 \label{eq:v77-candidate-orbit}
\end{equation}
The exact Ward, conservation, and anomaly-restoring coordinates are treated
as prescribed affine data and are not shooting variables.

At \(g_j^*\), split the tangent space as
\begin{equation}
 {\cal X}_j
 =
 X_j^s\oplus X_j^c\oplus X_j^u,
 \label{eq:v77-splitting}
\end{equation}
where \(X^s\) includes the polymer and positive-excess gravitational tail,
\(X^u\) contains forward-expanding relevant directions, and \(X^c\)
contains marginal directions selected by the actual beta map.  Canonical
dimension alone does not determine this splitting at a non-Gaussian orbit.

Let
\[
 A_j=D{\cal B}_j(g_j^*)
\]
preserve \eqref{eq:v77-splitting}.  Write the stable forward evolution as
\[
 \Phi^s(n,k)=A_{n-1}^s\cdots A_k^s,\qquad n\geq k,
\]
and the unstable backward evolution as
\[
 \Phi^u(n,k)=
 (A_n^u)^{-1}\cdots(A_{k-1}^u)^{-1},
 \qquad k>n.
\]
Assume the exponential dichotomy bounds
\begin{align}
 \|\Phi^s(n,k)\|
 &\leq M_sq_s^{\,n-k},
 \qquad0<q_s<1,
 \label{eq:v77-stable-dichotomy}\\
 \|\Phi^u(n,k)\|
 &\leq M_uq_u^{\,k-n},
 \qquad0<q_u<1.
 \label{eq:v77-unstable-dichotomy}
\end{align}

The center boundary condition depends on the actual marginal normal form.
Represent its linear inhomogeneous solution operator by
\begin{equation}
 {\cal G}^c:\mathcal Y_c\longrightarrow\mathcal Z_c,
 \qquad
 \|{\cal G}^c\|\leq G_c,
 \label{eq:v77-center-Green}
\end{equation}
on a declared polynomially weighted sequence space.  Existence of
\({\cal G}^c\) is a physical beta-map hypothesis, not a consequence of
\eqref{eq:v77-stable-dichotomy}--\eqref{eq:v77-unstable-dichotomy}.

\subsection{The full Green operator}

Let \({\cal Z}\) be the product sequence space for the three sectors, with
the center weight built into its norm.  For a forcing sequence
\(f=(f^s,f^c,f^u)\), define
\begin{align}
 ({\cal G}^sf)_n
 &:=
 \sum_{k=0}^{n-1}\Phi^s(n,k+1)f_k^s,
 \label{eq:v77-stable-Green}\\
 ({\cal G}^uf)_n
 &:=
 -\sum_{k=n}^{\infty}\Phi^u(n,k+1)f_k^u,
 \label{eq:v77-unstable-Green}\\
 {\cal G}f
 &:=
 {\cal G}^sf+{\cal G}^cf+{\cal G}^uf.
 \label{eq:v77-full-Green}
\end{align}
The first sum enforces the prescribed stable initial datum; the second
removes the forward-growing unstable homogeneous component.

\begin{lemma}[Green norm]
\label{lem:v77-Green-norm}
On the corresponding maximum product norm,
\begin{equation}
 \|{\cal G}\|
 \leq
 G_*:=
 \max\left\{
 \frac{M_s}{1-q_s},
 G_c,
 \frac{M_uq_u}{1-q_u}
 \right\}.
 \label{eq:v77-G-star}
\end{equation}
\end{lemma}

\begin{proof}
Sum the stable geometric series in
\eqref{eq:v77-stable-Green}.  In the unstable sum,
\(k+1-n\geq1\), so its geometric sum is
\(M_uq_u/(1-q_u)\).  The center estimate is
\eqref{eq:v77-center-Green}.  The maximum product norm gives the stated
maximum.
\end{proof}

Let \(h\in{\cal Z}\) be the linear homogeneous sequence encoding the
prescribed stable initial datum, the chosen marginal renormalization datum,
and zero unstable growing mode.  Write the exact perturbation equation
around \(g^*\) as
\begin{equation}
 z=h+{\cal G}{\cal F}(z),
 \label{eq:v77-Lyapunov-Perron}
\end{equation}
where \({\cal F}\) contains nonlinear terms, scale contacts, and the
difference between the exact map and the declared linearization.

\begin{theorem}[Nonautonomous three-sector shooting]
\label{thm:v77-three-sector-shooting}
Suppose on the radius-\(r\) sequence ball,
\begin{align}
 \|{\cal F}(z)-{\cal F}(z')\|_{\cal Y}
 &\leq L_F\|z-z'\|_{\cal Z},
 \label{eq:v77-F-Lipschitz}\\
 \|{\cal F}(0)\|_{\cal Y}&\leq b_F.
 \label{eq:v77-F-forcing}
\end{align}
Define
\begin{equation}
 \kappa_{\rm shot}:=G_*L_F.
 \label{eq:v77-shot-kappa}
\end{equation}
If
\begin{equation}
 \kappa_{\rm shot}<1,
 \qquad
 \|h\|_{\cal Z}+G_*b_F+\kappa_{\rm shot}r\leq r,
 \label{eq:v77-shot-gate}
\end{equation}
then \eqref{eq:v77-Lyapunov-Perron} has a unique solution in the radius-\(r\)
ball.  This solution is the unique trajectory satisfying the prescribed
stable and center data and carrying no forward-growing unstable homogeneous
component.
\end{theorem}

\begin{proof}
The map
\[
 {\mathfrak T}(z)=h+{\cal G}{\cal F}(z)
\]
has Lipschitz constant at most \(G_*L_F=\kappa_{\rm shot}\) by Lemma
\ref{lem:v77-Green-norm}.  At zero its norm is at most
\(\|h\|+G_*b_F\).  The second inequality in
\eqref{eq:v77-shot-gate} makes the ball invariant, and the first makes the
map a contraction.  Banach's theorem gives the unique fixed point.

Taking the discrete difference of the three Green formulae verifies the
forward perturbation equation.  Conversely, iteration forward in the
stable sector and backward in the unstable sector produces
\eqref{eq:v77-stable-Green}--\eqref{eq:v77-unstable-Green}.  The chosen
center boundary condition is equivalent to
\eqref{eq:v77-center-Green}.  Hence every admissible trajectory solves the
same fixed-point equation and is unique.
\end{proof}

The bare relevant counterterms are the \(X_0^u\) component of the fixed
point.  They are selected by the infinite future condition; they are not
forward-contracted.

\subsection{A coupled marginal cone}

The center Green operator must be derived from the marginal beta map.  The
following vector estimate supplies a checkable sufficient normal form.

Let \(K\subset{\mathbb R}^m\) be a closed convex pointed cone and let
\(\ell\) be a linear functional strictly positive on \(K\setminus\{0\}\).
Use
\[
 |m|_\ell:=\ell(m)
\]
as the radial coordinate on \(K\).  Consider
\begin{equation}
 m_{j+1}
 =
 m_j-B(m_j,m_j)+r_j,
 \label{eq:v77-marginal-vector}
\end{equation}
where \(B\) is bilinear.

\begin{theorem}[Coupled marginally irrelevant cone]
\label{thm:v77-marginal-cone}
Suppose \(K\) is invariant under
\eqref{eq:v77-marginal-vector} for
\(|m|_\ell\leq\delta\), and there are
\(0<b_-\leq b_+\) and \(c\geq0\) such that
\begin{align}
 b_-|m|_\ell^2
 \leq
 \ell(B(m,m))
 &\leq
 b_+|m|_\ell^2,
 \label{eq:v77-cone-quadratic}\\
 |\ell(r_j)|
 &\leq
 c|m_j|_\ell^3.
 \label{eq:v77-cone-remainder}
\end{align}
If
\begin{equation}
 c\delta\leq\frac{b_-}{2},
 \qquad
 c\delta\leq\frac{b_+}{2},
 \qquad
 \frac{3b_+}{2}\delta\leq\frac12,
 \label{eq:v77-cone-smallness}
\end{equation}
then every \(m_0\in K\) with
\(0<|m_0|_\ell\leq\delta\) stays in \(K\), decreases radially, and obeys
\begin{equation}
 \frac1{|m_0|_\ell^{-1}+3b_+j}
 \leq
 |m_j|_\ell
 \leq
 \frac1{|m_0|_\ell^{-1}+(b_-/2)j}.
 \label{eq:v77-cone-decay}
\end{equation}
\end{theorem}

\begin{proof}
Put \(u_j=|m_j|_\ell\).  Applying \(\ell\) to
\eqref{eq:v77-marginal-vector} and using
\eqref{eq:v77-cone-quadratic}--\eqref{eq:v77-cone-smallness} gives
\begin{equation}
 \frac{b_-}{2}u_j^2
 \leq
 u_j-u_{j+1}
 \leq
 \frac{3b_+}{2}u_j^2.
 \label{eq:v77-cone-decrement}
\end{equation}
The upper bound is at most \(u_j/2\), while the lower bound is positive.
Cone invariance closes the induction and gives
\(0<u_{j+1}\leq u_j\leq\delta\).

Now
\[
 u_{j+1}^{-1}-u_j^{-1}
 =
 \frac{u_j-u_{j+1}}{u_ju_{j+1}}.
\]
Since \(u_{j+1}\leq u_j\), the lower bound is \(b_-/2\).  Since
\(u_{j+1}\geq u_j/2\), the upper bound is \(3b_+\).  Sum and invert to obtain
\eqref{eq:v77-cone-decay}.
\end{proof}

This theorem permits coupled marginal rows only when the actual quadratic
beta tensor preserves a cone and has a strictly dissipative radial
component there.  Gauge asymptotic freedom in one coordinate does not prove
the condition for the Higgs, Yukawa, Newton, and curvature-squared system.

To use Theorem \ref{thm:v77-three-sector-shooting}, one must additionally
linearize \eqref{eq:v77-marginal-vector} along the chosen cone orbit and
prove the weighted Green estimate \eqref{eq:v77-center-Green}.  Radial
decay alone does not control transverse rotations or neutral directions.

\subsection{Admissible reserve barrier}

Let \(\lambda\) denote the full running owner vector.  For each scale define
validated certificate functions:
\begin{align}
 {\cal R}_{\rm piv}(\lambda)
 &=
 \min\{\gamma_\sigma,\gamma_{\widehat A},
              \gamma_D,\gamma_{\widehat S}\},
 \label{eq:v77-pivot-reserve}\\
 {\cal R}_{\rm imp}(\lambda)
 &=
 \min\{\Delta_{\rm map},\Delta_{\rm der}\},
 \label{eq:v77-implicit-reserve}\\
 {\cal R}_{\rm for}(\lambda)
 &=
 1-8K_Tz_{\rm reb},
 \label{eq:v77-forest-reserve}\\
 {\cal R}_{\rm det}(\lambda)
 &=
 \min_s\{1-\eta_s\},
 \label{eq:v77-determinant-reserve}\\
 {\cal R}_{\rm tail}(\lambda)
 &=
 1-\lambda_{\rm tail},
 \label{eq:v77-tail-reserve}\\
 {\cal R}_{\rm rate}(\lambda)
 &=
 1-\sigma_*/\varsigma.
 \label{eq:v77-rate-reserve}
\end{align}
A certificate may contain additional stocks, but it must imply the
corresponding theorem's hypotheses.  In particular,
\({\cal R}_{\rm rate}>0\) uses a proved many-block value of \(\sigma_*\),
not a local Hessian.

Let the candidate orbit satisfy
\begin{equation}
 {\cal R}_a(g_j^*)\geq\rho_a>0
 \quad\hbox{for all }j
 \label{eq:v77-candidate-reserves}
\end{equation}
and every listed row \(a\).  Assume that on a radius-\(r\) tube,
\begin{equation}
 |{\cal R}_a(\lambda)-{\cal R}_a(g_j^*)|
 \leq
 L_a\|\lambda-g_j^*\|.
 \label{eq:v77-reserve-Lipschitz}
\end{equation}

\begin{theorem}[Reserve-preserving shot]
\label{thm:v77-reserve-preserving}
If the fixed point of Theorem
\ref{thm:v77-three-sector-shooting} satisfies
\begin{equation}
 \|z\|_{\cal Z}
 \leq
 r_{\rm adm}
 :=
 \min_a\frac{\rho_a}{2L_a},
 \label{eq:v77-admissible-radius}
\end{equation}
with the convention that a row with \(L_a=0\) imposes no restriction, then
along the shot trajectory
\begin{equation}
 {\cal R}_a(g_j^*+z_j)\geq\frac{\rho_a}{2}>0
 \label{eq:v77-shot-reserves}
\end{equation}
for every scale and every certificate row.
\end{theorem}

\begin{proof}
Equations \eqref{eq:v77-candidate-reserves},
\eqref{eq:v77-reserve-Lipschitz}, and
\eqref{eq:v77-admissible-radius} give
\[
 {\cal R}_a(g_j^*+z_j)
 \geq
 \rho_a-L_a\|z_j\|
 \geq
 \rho_a/2.
\]
\end{proof}

Thus a shooting theorem is useful to the continuum programme only when its
sequence ball fits inside \(r_{\rm adm}\).  The abstract contraction radius
cannot be chosen without reference to the analytic and measure chambers.

Combining \eqref{eq:v77-shot-gate} and
\eqref{eq:v77-admissible-radius} gives the explicit sufficient condition
\begin{equation}
 \|h\|_{\cal Z}+G_*b_F
 \leq
 (1-\kappa_{\rm shot})r_{\rm adm}.
 \label{eq:v77-combined-shot-gate}
\end{equation}

\subsection{Finite-horizon stability}

Let \({\cal G}_N\) truncate the unstable future sum at \(N\), with bounded
terminal unstable datum.  The missing homogeneous term is bounded by
\[
 M_uq_u^{N-n}\|y_N^{\rm term}\|.
\]
If the constants in
\eqref{eq:v77-combined-shot-gate} are uniform in \(N\), every fixed scale
window has the same terminal-removal estimate as V16.  Extraction of a
measure subsequence still requires the V10 tightness and equicontinuity
cards; the graph transform does not supply them.

\subsection{What must be calculated}

The theorem turns the retained-flow bottleneck into four concrete inputs:
\begin{enumerate}
\item a candidate projected SM--Einstein orbit \(g_j^*\);
\item the spectral projectors and dichotomy constants of
      \(D{\cal B}_j(g_j^*)\);
\item the coupled marginal quadratic tensor \(B\), its invariant cone, and
      the center Green bound; and
\item uniform lower bounds for every reserve in
      \eqref{eq:v77-pivot-reserve}--\eqref{eq:v77-rate-reserve}.
\end{enumerate}
None follows from the existence of formal beta functions.  The beta map
must be computed in the same regulator, owner basis, and determinant scheme
used by the forest and constraint charts.

\begin{firewall}[No candidate orbit has been constructed]
\label{fire:v77-no-orbit}
V77 proves a graph transform around a candidate orbit with verified
dichotomy, center, and reserve data.  It does not construct that orbit or
calculate the coupled SM--Einstein beta tensor.  In particular, it does not
assert asymptotic safety, asymptotic freedom of the whole marginal sector,
or a positive global gravitational rate floor.
\end{firewall}

\subsection{Next acquisition}

V78 will construct the finite-dimensional derivative card for the projected
beta map without pretending to know its spectrum.  It will show how local
shell activities, determinant contacts, and coarea-score terms contribute
to \(D{\cal B}_j\), and will derive interval bounds that certify an
exponential dichotomy when a numerical or analytic approximate beta matrix
is supplied.  If the marginal cluster touches the unit circle, it will be
isolated for the cone/normal-form calculation rather than forced into a
hyperbolic block.

\begin{assessment}[Exact status after V77]
\label{ass:v77-status}
V77 proves a nonautonomous stable--center--unstable graph transform, a
coupled marginal cone decay theorem, and a reserve barrier preserving all
current analytic and measure cards.  The candidate SM--Einstein orbit,
coupled beta matrix, center Green bound, global gravitational rate floor,
measure compactness, regulator removal, and the full continuum remain open.
\end{assessment}
\section{Auditable beta-Jacobian and dichotomy certification}
\label{sec:v10v78}

V77 requires the derivative of the projected RG map along a candidate orbit.
A formal list of beta functions is insufficient: determinant, residual
Jacobian, projection-contact, and covariance-transport responses must be
differentiated in the same owner frame.  This section gives the exact
derivative ledger and two certification theorems.  One treats a
scale-dependent approximate linear sequence; the other treats a candidate
fixed point by resolvent contours.

\subsection{Disjoint projected beta owners}

Let \(c\) be the finite vector of noncontracting physical couplings after
Ward-fixed and redundant coordinates have been removed.  In a transported
owner frame, write one exact projected step as
\begin{equation}
 {\cal B}_j(c)
 =
 S_jc+
 \sum_{o\in{\mathscr O}}b_{j,o}(c),
 \label{eq:v78-beta-decomposition}
\end{equation}
where \(S_j\) is canonical transport and
\begin{equation}
 {\mathscr O}
 =
 \{
 {\rm core,\ determinant,\ residual\ Jacobian,\ covariance\ path,\
 contact,\ compensation}
 \}.
 \label{eq:v78-owner-set}
\end{equation}
Each contribution is already projected into its actual running owner.  A
physical fermion determinant, a ghost determinant, a bosonic Gaussian
determinant, and a coordinate Jacobian remain separate subowners inside the
determinant and residual-Jacobian rows.

The exact derivative is
\begin{equation}
 D{\cal B}_j(c)
 =
 S_j+\sum_{o\in{\mathscr O}}Db_{j,o}(c).
 \label{eq:v78-exact-Jacobian}
\end{equation}
If an owner is represented before projection as
\begin{equation}
 b_{j,o}(c)=P_{j,o}(c)\ell_{j,o}(c),
 \label{eq:v78-moving-projection}
\end{equation}
then
\begin{equation}
 Db_{j,o}[\dot c]
 =
 (DP_{j,o}[\dot c])\ell_{j,o}
 +
 P_{j,o}D\ell_{j,o}[\dot c].
 \label{eq:v78-projection-derivative}
\end{equation}
The first term is part of the V54 contact/frame connection.  Omitting it
changes the beta matrix.

\begin{ledger}[Derivative ownership]
\label{led:v78-derivative-ownership}
\begin{enumerate}
\item Derivatives of a relative determinant use its V53 marked response.
\item Derivatives of the residual-chart determinant use V73, and its fine
      score is not added a second time.
\item The order-one covariance mismatch is differentiated through the exact
      V76 path and its Gaussian determinant, not through a small quadratic
      vertex.
\item The V54 projection contact is differentiated once through
      \eqref{eq:v78-projection-derivative}.
\item The V57 compensation source contributes its projected running part;
      its high-grade part stays in the stable tail.
\end{enumerate}
\end{ledger}

\subsection{Cauchy interval for the beta matrix}

Fix a candidate point \(c_j^*\) and complex coupling polydiscs
\[
 {\cal D}_r(c_j^*)\Subset{\cal D}_R(c_j^*),
 \qquad
 \delta_c:=R-r>0.
\]
For each owner choose an explicitly evaluated affine approximation
\begin{equation}
 b_{j,o}^{\rm app}(c)
 =
 b_{j,o}^{(0)}
 +
 B_{j,o}^{\rm app}(c-c_j^*)
 \label{eq:v78-owner-affine}
\end{equation}
and analytic remainder
\[
 \rho_{j,o}=b_{j,o}-b_{j,o}^{\rm app}.
\]
Suppose
\begin{equation}
 \sup_{{\cal D}_R(c_j^*)}
 \|\rho_{j,o}\|
 \leq z_{j,o}^{\rm rem}.
 \label{eq:v78-owner-remainder}
\end{equation}
Set
\begin{equation}
 A_j^{\rm app}
 :=
 S_j+\sum_oB_{j,o}^{\rm app}.
 \label{eq:v78-approximate-matrix}
\end{equation}

\begin{theorem}[Beta-Jacobian interval]
\label{thm:v78-Jacobian-interval}
On \({\cal D}_r(c_j^*)\),
\begin{equation}
 \|D{\cal B}_j(c)-A_j^{\rm app}\|
 \leq
 \epsilon_j,
 \qquad
 \epsilon_j:=
 \frac1{\delta_c}
 \sum_{o\in{\mathscr O}}z_{j,o}^{\rm rem}.
 \label{eq:v78-Jacobian-interval}
\end{equation}
For anisotropic coupling radii, replace
\(\delta_c^{-1}\) by the corresponding column radius losses.
\end{theorem}

\begin{proof}
Cauchy's integral formula on the coupling polydisc gives
\[
 \|D\rho_{j,o}\|_{{\cal D}_r}
 \leq
 \delta_c^{-1}z_{j,o}^{\rm rem}.
\]
Differentiate \eqref{eq:v78-beta-decomposition}, subtract
\eqref{eq:v78-approximate-matrix}, and sum the owner bounds.
\end{proof}

The theorem is useful only if the affine coefficients in
\eqref{eq:v78-owner-affine} and the remainder norms are computed in the same
scheme.  Moving a marginal quadratic term from the affine matrix into a
remainder can make \(\epsilon_j\) artificially small at one radius and
destroy it after reblocking.

The V53, V70, V72, and V73 analytic estimates provide valid sources for the
numbers \(z_{j,o}^{\rm rem}\), after the noncontracting local projections
have been removed.  They do not evaluate those numbers automatically.

\subsection{Roughness of a scale-dependent Green operator}

Let \({\cal Z}\) and \({\cal Y}\) be the sequence spaces and boundary
conditions used in V77.  Define
\begin{equation}
 ({\cal L}_{\rm app}z)_j
 :=
 z_{j+1}-A_j^{\rm app}z_j.
 \label{eq:v78-L-app}
\end{equation}
Assume it has a right inverse, with the prescribed stable, center, and
unstable boundary conditions,
\begin{equation}
 {\cal G}_{\rm app}:{\cal Y}\to{\cal Z},
 \qquad
 \|{\cal G}_{\rm app}\|\leq G_{\rm app}.
 \label{eq:v78-G-app}
\end{equation}
Let
\begin{equation}
 E_j:=D{\cal B}_j(c_j^*)-A_j^{\rm app},
 \qquad
 \|E\|_{{\cal Z}\to{\cal Y}}
 :=
 \sup_j\|E_j\|
 \leq\epsilon.
 \label{eq:v78-E-sequence}
\end{equation}

\begin{theorem}[Certified nonautonomous Green operator]
\label{thm:v78-Green-roughness}
If
\begin{equation}
 G_{\rm app}\epsilon<1,
 \label{eq:v78-Green-gate}
\end{equation}
then the exact linearized difference operator
\begin{equation}
 ({\cal L}z)_j
 =
 z_{j+1}-D{\cal B}_j(c_j^*)z_j
 \label{eq:v78-exact-L}
\end{equation}
has the same boundary-value Green operator
\begin{equation}
 {\cal G}
 =
 (I-{\cal G}_{\rm app}E)^{-1}{\cal G}_{\rm app},
 \label{eq:v78-exact-Green}
\end{equation}
and
\begin{equation}
 \|{\cal G}\|
 \leq
 \frac{G_{\rm app}}
 {1-G_{\rm app}\epsilon}.
 \label{eq:v78-Green-bound}
\end{equation}
If the sequence spaces carry exponential stable/unstable and polynomial
center weights, the same weights are preserved.
\end{theorem}

\begin{proof}
The exact equation
\({\cal L}z=f\) is equivalent to
\[
 {\cal L}_{\rm app}z=f+Ez.
\]
Applying the approximate Green operator gives
\[
 z={\cal G}_{\rm app}f+{\cal G}_{\rm app}Ez.
\]
Condition \eqref{eq:v78-Green-gate} makes
\(I-{\cal G}_{\rm app}E\) invertible by its Neumann series.  This proves
\eqref{eq:v78-exact-Green}.  Summing that series gives
\eqref{eq:v78-Green-bound}.  Since all operators act on the already weighted
spaces, no weight is lost.
\end{proof}

The theorem certifies the V77 center and dichotomy Green stock with
\[
 G_*=
 \frac{G_{\rm app}}
 {1-G_{\rm app}\epsilon}.
\]
It is stronger than checking the eigenvalues of each \(D{\cal B}_j\)
separately.  Rapidly rotating nonnormal eigenspaces can invalidate a
nonautonomous dichotomy even when every one-step spectrum avoids the unit
circle; the sequence Green test detects that failure.

\subsection{Fixed-point spectral certificate}

At an autonomous candidate fixed point, a resolvent enclosure can first
construct the approximate splitting.  Let \(A^{\rm app}\) be a finite
matrix and let \(\Gamma\) be a positively oriented contour disjoint from its
spectrum.  Put
\begin{equation}
 M_\Gamma
 :=
 \sup_{z\in\Gamma}
 \|(zI-A^{\rm app})^{-1}\|.
 \label{eq:v78-resolvent-stock}
\end{equation}

\begin{theorem}[Riesz-projector interval]
\label{thm:v78-Riesz-interval}
If
\[
 A=A^{\rm app}+E,
 \qquad
 \|E\|\leq\epsilon,
 \qquad
 M_\Gamma\epsilon<1,
\]
then \(\Gamma\) is disjoint from the spectrum of \(A\), and the Riesz
projectors
\[
 P=\frac1{2\pi i}\int_\Gamma(zI-A)^{-1}dz,
 \qquad
 P^{\rm app}=
 \frac1{2\pi i}\int_\Gamma(zI-A^{\rm app})^{-1}dz
\]
have the same rank.  Moreover
\begin{equation}
 \|P-P^{\rm app}\|
 \leq
 \frac{|\Gamma|}{2\pi}
 \frac{M_\Gamma^2\epsilon}
 {1-M_\Gamma\epsilon}.
 \label{eq:v78-projector-bound}
\end{equation}
\end{theorem}

\begin{proof}
Factor
\[
 zI-A=
 \{I-E(zI-A^{\rm app})^{-1}\}(zI-A^{\rm app}).
\]
The first factor is invertible by a Neumann series uniformly on
\(\Gamma\).  The resolvent identity gives
\[
 (zI-A)^{-1}-(zI-A^{\rm app})^{-1}
 =
 (zI-A)^{-1}E(zI-A^{\rm app})^{-1},
\]
whose norm is at most
\(M_\Gamma^2\epsilon/(1-M_\Gamma\epsilon)\).  Integrating gives
\eqref{eq:v78-projector-bound}.  The projectors vary continuously along
\(A^{\rm app}+tE\); their integer rank is therefore constant.
\end{proof}

Use disjoint contours for stable, center, and unstable clusters.  After
their ranks are certified, their actual evolution still must pass Theorem
\ref{thm:v78-Green-roughness} along a nonconstant orbit.

\subsection{Marginal cluster rule}

An eigenvalue enclosure meeting the unit circle is not assigned to the
stable or unstable block.  It enters the center cluster.  To close that
cluster one needs:
\begin{enumerate}
\item the quadratic and cubic beta tensors in its Riesz range;
\item the invariant-cone test of V77 or another proved normal form;
\item the polynomially weighted approximate Green operator; and
\item the interval inequality \eqref{eq:v78-Green-gate}.
\end{enumerate}
Rounding an interval away from one or classifying it by canonical dimension
would invalidate the shooting proof.

\subsection{Reserve derivatives from the same card}

For each certificate reserve \({\cal R}_a(c)\) in V77, choose an analytic
extension to \({\cal D}_R(c_j^*)\).  Cauchy's formula gives
\begin{equation}
 L_a
 \leq
 \frac{
 \sup_{{\cal D}_R}
 |{\cal R}_a(c)-{\cal R}_a(c_j^*)|
 }{\delta_c}.
 \label{eq:v78-reserve-Lipschitz}
\end{equation}
Thus the admissible radius
\[
 r_{\rm adm}=\min_a\rho_a/(2L_a)
\]
can be certified from the same coupling polydisc used for
\eqref{eq:v78-Jacobian-interval}.  The rate reserve requires an analytic or
interval certificate for the many-block constant \(\sigma_*(c)\); a
pointwise optimization at the base coupling is insufficient.

\subsection{Auditable derivative card}

\begin{definition}[Projected beta-derivative card]
\label{def:v78-beta-card}
A regulator step passes the card if:
\begin{enumerate}
\item every term in \eqref{eq:v78-beta-decomposition} has a disjoint owner
      and an affine approximation in one transported frame;
\item every analytic remainder has the strong-polydisc bound
      \eqref{eq:v78-owner-remainder};
\item moving projections include \eqref{eq:v78-projection-derivative};
\item a fixed-point Riesz enclosure or a scale-dependent approximate Green
      operator identifies the stable, center, and unstable boundary data;
\item the interval error passes \eqref{eq:v78-Green-gate}; and
\item the reserve derivatives pass
      \eqref{eq:v78-reserve-Lipschitz} and the V77 combined shooting gate.
\end{enumerate}
\end{definition}

\begin{firewall}[Interval machinery does not supply beta coefficients]
\label{fire:v78-no-beta-values}
V78 proves how verified shell and determinant bounds certify a beta
linearization and dichotomy.  It does not calculate the Standard
Model--Einstein affine coefficients \(B_{j,o}^{\rm app}\), the marginal
quadratic tensor, or the candidate orbit.  A zero-width numerical matrix
without error bounds is not evidence for the card.
\end{firewall}

\subsection{Next acquisition}

V79 will move from coupling trajectories to measures.  It will formulate a
projective-consistency and tightness theorem for the finite-regulator
constraint-reduced measures along a V77 shot orbit.  The theorem will keep
the chiral determinant-line phase, boundary states, and retained zero modes
in the projective system and will identify which uniform moment and
equicontinuity estimates are still absent.  This will prevent a successful
coupling-space shooting theorem from being mistaken for construction of a
continuum field measure.

\begin{assessment}[Exact status after V78]
\label{ass:v78-status}
V78 proves an exact beta-Jacobian owner ledger, analytic interval bounds,
roughness of the scale-dependent Green operator, fixed-point Riesz
certification, and reserve-derivative bounds.  The actual SM--Einstein beta
coefficients, candidate orbit, marginal center calculation, globally
coercive reduced measure, tightness, regulator removal, and the full
continuum remain open.
\end{assessment}
\section{Projective consistency, tightness, and the measure limit}
\label{sec:v10v79}

V77--V78 concern trajectories in coupling space.  Even a complete shot
trajectory would not by itself construct a continuum field theory.  This
section supplies the missing measure-level theorem.  It treats exact
projective consistency and the summably approximate consistency generated
by adjacent-density comparisons.  It then separates inverse-limit
existence from tightness in the desired distribution space.

\subsection{Finite regulator tower}

Let
\[
 E_0\longleftarrow E_1\longleftarrow E_2\longleftarrow\cdots
\]
be standard Borel regulator spaces with measurable coarse maps
\begin{equation}
 p_{n+1,n}:E_{n+1}\to E_n,
 \qquad
 p_{m,n}:=p_{n+1,n}\circ\cdots\circ p_{m,m-1}.
 \label{eq:v79-coarse-maps}
\end{equation}
Each \(E_n\) includes all retained boundary, zero-mode, charge, branch, and
determinant-line frame data needed to define the next coarse map.  Removing
such a coordinate before it has a compatible pushforward changes the
projective system.

Let \(\mu_n\) be the normalized positive bosonic body at regulator \(n\).
Fermionic and chiral data will be added coefficientwise below.

\begin{theorem}[Exact inverse-limit measure]
\label{thm:v79-exact-projective}
Suppose
\begin{equation}
 (p_{n+1,n})_*\mu_{n+1}=\mu_n
 \quad\hbox{for every }n.
 \label{eq:v79-exact-consistency}
\end{equation}
Then there is a unique probability measure \(\mu_\infty^{\rm inv}\) on the
cylinder sigma algebra of
\begin{equation}
 E_\infty^{\rm inv}
 :=
 \{(x_n):p_{n+1,n}x_{n+1}=x_n\hbox{ for every }n\}
 \label{eq:v79-inverse-limit}
\end{equation}
whose \(n\)-th marginal is \(\mu_n\).
\end{theorem}

\begin{proof}
For each \(N\), push \(\mu_N\) to
\(E_0\times\cdots\times E_N\) by
\[
 x_N\longmapsto
 (p_{N,0}x_N,\ldots,p_{N,N-1}x_N,x_N).
\]
Equation \eqref{eq:v79-exact-consistency} makes these finite-dimensional
laws consistent.  The Kolmogorov extension theorem on the countable product
of standard Borel spaces gives a probability measure with those marginals.
Every relation
\(p_{n+1,n}x_{n+1}=x_n\) holds with probability one.  Their countable
intersection is \eqref{eq:v79-inverse-limit}, so the measure is supported
there.  Cylinder marginals determine uniqueness.
\end{proof}

This theorem constructs an abstract inverse-limit random object.  It does
not show that the object is a distribution, has the required regularity, or
obeys continuum symmetries.

\subsection{Summably approximate consistency}

Exact block integration can be replaced by a controlled adjacent
comparison.  Assume
\begin{equation}
 \left\|
 (p_{n+1,n})_*\mu_{n+1}-\mu_n
 \right\|_{\rm TV}
 \leq\epsilon_n.
 \label{eq:v79-adjacent-error}
\end{equation}

\begin{theorem}[Summable projective repair]
\label{thm:v79-summable-repair}
If
\begin{equation}
 \sum_{n=0}^{\infty}\epsilon_n<\infty,
 \label{eq:v79-error-sum}
\end{equation}
then, for every fixed \(k\), the measures
\begin{equation}
 \nu_{N,k}:=(p_{N,k})_*\mu_N,\qquad N\geq k,
 \label{eq:v79-coarse-tail-law}
\end{equation}
converge in total variation to a probability measure \(\nu_k\).  The limit
family is exactly consistent:
\begin{equation}
 (p_{k+1,k})_*\nu_{k+1}=\nu_k,
 \label{eq:v79-repaired-consistency}
\end{equation}
and
\begin{equation}
 \|\nu_k-\mu_k\|_{\rm TV}
 \leq
 \sum_{n=k}^{\infty}\epsilon_n.
 \label{eq:v79-repair-distance}
\end{equation}
Consequently \((\nu_k)\) has the inverse-limit measure of Theorem
\ref{thm:v79-exact-projective}.
\end{theorem}

\begin{proof}
Total variation contracts under measurable pushforward.  Hence
\begin{align*}
 \|\nu_{N+1,k}-\nu_{N,k}\|_{\rm TV}
 &=
 \left\|
 (p_{N,k})_*
 \{(p_{N+1,N})_*\mu_{N+1}-\mu_N\}
 \right\|_{\rm TV}\\
 &\leq\epsilon_N.
\end{align*}
The tail of \(\sum\epsilon_N\) makes
\((\nu_{N,k})_{N\geq k}\) Cauchy in the Banach space of finite signed
measures with total-variation norm.  Its limit is nonnegative and has mass
one.  The same estimate summed from \(N=k\) proves
\eqref{eq:v79-repair-distance}.  Pushforward is continuous in total
variation, so
\[
 (p_{k+1,k})_*\nu_{k+1}
 =
 \lim_{N\to\infty}
 (p_{N,k})_*\mu_N
 =
 \nu_k.
\]
Apply Theorem \ref{thm:v79-exact-projective}.
\end{proof}

The adjacent error \(\epsilon_n\) must be computed from the complete
density difference: action, determinant, coarea, ambient transport,
constraint Jacobian, covariance path, local counterterm, and branch
owners.  The V54 cocycle prevents determinant duplication; the V73 score
prevents omission of the moving residual chart.

\subsection{Embedding in a continuum distribution space}

Let \({\cal S}\) be a Polish distribution space, for example a chosen
negative Sobolev or Besov space on every compact spacetime window with the
usual projective metric.  Suppose continuous reconstruction maps
\begin{equation}
 \iota_n:E_n\to{\cal S}
 \label{eq:v79-reconstruction}
\end{equation}
and continuous cylinder projections
\(\pi_k:{\cal S}\to E_k\) satisfy
\begin{equation}
 \pi_k\iota_n=p_{n,k}
 \qquad(n\geq k)
 \label{eq:v79-reconstruction-consistency}
\end{equation}
on all retained physical coordinates.  Put
\[
 \widetilde\mu_n=(\iota_n)_*\mu_n.
\]

\begin{theorem}[Tight projective continuum limit]
\label{thm:v79-tight-projective-limit}
Assume:
\begin{enumerate}
\item the reconstructed laws \((\widetilde\mu_n)\) are uniformly tight on
      \({\cal S}\);
\item the adjacent errors obey \eqref{eq:v79-error-sum}; and
\item the countable family \((\pi_k)\) separates points and generates the
      Borel sigma algebra of \({\cal S}\).
\end{enumerate}
Then \(\widetilde\mu_n\) converges weakly to a unique Radon probability
measure \(\widetilde\mu_\infty\) on \({\cal S}\), and
\begin{equation}
 (\pi_k)_*\widetilde\mu_\infty=\nu_k
 \label{eq:v79-limit-marginals}
\end{equation}
for every \(k\).
\end{theorem}

\begin{proof}
Uniform tightness and Prokhorov's theorem give a weakly convergent
subsequence.  Fix \(k\).  For \(n\geq k\),
\[
 (\pi_k)_*\widetilde\mu_n
 =
 (p_{n,k})_*\mu_n
 =
 \nu_{n,k},
\]
which converges in total variation, hence weakly, to \(\nu_k\) by Theorem
\ref{thm:v79-summable-repair}.  The continuous mapping theorem shows that
every weak subsequential limit has \(k\)-th marginal \(\nu_k\).
Since the projections generate the Borel sigma algebra, these marginals
determine at most one probability measure on \({\cal S}\).  Thus every
subsequence has the same possible limit.  Tightness then implies convergence
of the full sequence, and \eqref{eq:v79-limit-marginals} holds.
\end{proof}

\begin{corollary}[Compact-embedding tightness criterion]
\label{cor:v79-compact-tightness}
Let \({\cal S}_0\) be a separable Banach space compactly embedded in
\({\cal S}\).  If, for some \(p>0\),
\begin{equation}
 \sup_n
 {\mathbb E}_{\widetilde\mu_n}
 \|X\|_{{\cal S}_0}^p
 \leq M,
 \label{eq:v79-S0-moment}
\end{equation}
then \((\widetilde\mu_n)\) is uniformly tight on \({\cal S}\).
\end{corollary}

\begin{proof}
Markov's inequality gives probability at least
\(1-M/R^p\) to the radius-\(R\) ball of \({\cal S}_0\).  Its closure is
compact in \({\cal S}\).  Choose \(R\) uniformly for any prescribed tail
probability.
\end{proof}

On a noncompact spacetime, apply the corollary on an increasing sequence of
compact windows and use a diagonal choice of radii whose tail probabilities
are summable.

The V76 normalized coercivity/free-energy row can imply
\eqref{eq:v79-S0-moment} only after a deterministic reconstruction estimate
bounds the \({\cal S}_0\) norm by the controlled energy.  Neither coercivity
nor that reconstruction has yet been proved for the full nonlinear
constraint-reduced gravitational field.

\subsection{Constraint support in the limit}

Let \({\cal Q}:{\cal S}\to[0,\infty]\) be a lower semicontinuous continuum
constraint defect, including only legitimate canonical constraints and
gauge-quotient residuals.

\begin{proposition}[Vanishing residual support]
\label{prop:v79-constraint-support}
If
\begin{equation}
 \int_{\cal S}{\cal Q}(X)\,d\widetilde\mu_n(X)
 \longrightarrow0,
 \label{eq:v79-constraint-moment}
\end{equation}
then the weak limit in Theorem
\ref{thm:v79-tight-projective-limit} satisfies
\begin{equation}
 {\cal Q}(X)=0
 \quad\hbox{for }\widetilde\mu_\infty\hbox{-almost every }X.
 \label{eq:v79-constraint-support}
\end{equation}
\end{proposition}

\begin{proof}
The Portmanteau theorem for nonnegative lower semicontinuous functions gives
\[
 \int{\cal Q}\,d\widetilde\mu_\infty
 \leq
 \liminf_n\int{\cal Q}\,d\widetilde\mu_n=0.
\]
A nonnegative function with zero integral vanishes almost surely.
\end{proof}

For the V72 normalized residual Gaussian,
\[
 {\mathbb E}|v|^2=m\varepsilon.
\]
Thus \eqref{eq:v79-constraint-moment} follows for a fixed finite residual
block as \(\varepsilon\downarrow0\).  The continuum statement additionally
requires a mesh-uniform reconstruction of all residual components and
control of their number and compatibility identities.

\subsection{Ward identities and reflection positivity}

Let \({\cal A}_{\rm cyl}\) be a bounded continuous physical cylinder
observable algebra.  Suppose a regulated Ward operator \(W_n\) and a
continuum operator \(W\) satisfy, for every \(F\in{\cal A}_{\rm cyl}\),
\begin{align}
 \left|
 \int W_nF\,d\widetilde\mu_n
 \right|&\leq\delta_n(F),
 \qquad\delta_n(F)\to0,
 \label{eq:v79-Ward-error}\\
 W_nF&\longrightarrow WF
 \quad\hbox{uniformly on tight compact sets}.
 \label{eq:v79-Ward-convergence}
\end{align}

\begin{proposition}[Passage of Ward and OS identities]
\label{prop:v79-Ward-OS-limit}
Under Theorem \ref{thm:v79-tight-projective-limit} and
\eqref{eq:v79-Ward-error}--\eqref{eq:v79-Ward-convergence},
\begin{equation}
 \int WF\,d\widetilde\mu_\infty=0.
 \label{eq:v79-limit-Ward}
\end{equation}
If every finite regulator is reflection positive,
\begin{equation}
 \int\overline{\Theta F}\,F\,d\widetilde\mu_n\geq0
 \label{eq:v79-finite-OS}
\end{equation}
for positive-half bounded continuous cylinder \(F\), and reflection and the
cylinder embedding converge, then
\begin{equation}
 \int\overline{\Theta F}\,F\,d\widetilde\mu_\infty\geq0.
 \label{eq:v79-limit-OS}
\end{equation}
\end{proposition}

\begin{proof}
Choose a compact set carrying uniformly all but an arbitrarily small
probability.  Uniform convergence there, boundedness on its complement, and
weak convergence pass the Ward integral to the limit; then use
\eqref{eq:v79-Ward-error}.  For reflection positivity, the reflected
quadratic cylinder observable is bounded and continuous under the stated
compatibility.  Weak convergence passes its nonnegative expectation to the
limit.
\end{proof}

For polynomial or other unbounded observables, replace boundedness by a
uniform \(L^{1+\delta}\) bound so that the family is uniformly integrable.
Reflection covariance alone is insufficient; the finite inequality
\eqref{eq:v79-finite-OS} must come from the boundary-square construction.

\subsection{Fermionic coefficient projective system}

Grassmann variables do not define a positive probability law.  Let
\[
 {\mathfrak F}_n^{(r)}
\]
be the degree-\(r\) fermionic coefficient functional over the bosonic body,
written in a transported determinant-line frame.  For every fixed cylinder
test tuple \({\bf f}\), assume
\begin{align}
 \left|
 {\mathfrak F}_{n+1}^{(r)}({\bf f})
 -
 {\mathfrak F}_n^{(r)}({\bf f})
 \right|
 &\leq\epsilon_{n,r}({\bf f}),
 \label{eq:v79-fermion-adjacent}\\
 \sum_n\epsilon_{n,r}({\bf f})&<\infty,
 \label{eq:v79-fermion-sum}\\
 |{\mathfrak F}_n^{(r)}({\bf f})|
 &\leq
 C_r\prod_{a=1}^rp_r(f_a)
 \label{eq:v79-fermion-bound}
\end{align}
for continuous test seminorms \(p_r\).

\begin{theorem}[Coefficientwise fermionic limit]
\label{thm:v79-fermion-limit}
Under
\eqref{eq:v79-fermion-adjacent}--\eqref{eq:v79-fermion-bound}, every fixed
degree coefficient converges to a continuous multilinear functional
\({\mathfrak F}_\infty^{(r)}\).  The collection is compatible with the
bosonic limit and its cylinder projections.  If the finite determinant-line
frames differ by a cocycle, this conclusion is global only when that
cocycle has a compatible trivialization; otherwise the limit is a section,
not a scalar functional.
\end{theorem}

\begin{proof}
The summable adjacent bound makes every coefficient Cauchy.  The uniform
multilinear estimate passes to the limit and proves continuity.  Cylinder
compatibility passes through the same telescoping argument as Theorem
\ref{thm:v79-summable-repair}.  A nontrivial frame cocycle changes scalar
representatives on overlaps, so only a compatible trivialization makes the
coefficient globally scalar.
\end{proof}

The V17 local Standard Model anomaly sums remove the local curvature
obstruction, but do not by themselves supply the global trivialization
required in the last sentence.

\subsection{Retained boundary and zero-mode tightness}

Gradient energies do not control constant metric, toron, harmonic, Killing,
or global charge coordinates.  For each retained finite-dimensional
coordinate \(z_n\), tightness requires a separate condition such as
\begin{equation}
 \sup_n{\mathbb E}|z_n|^p<\infty
 \label{eq:v79-zero-mode-moment}
\end{equation}
or compactness of its physical quotient.  These variables must be included
in \(\iota_n\), \(\pi_k\), and the boundary reflection algebra.  Treating
them as absent from the reconstruction can produce apparent tightness while
mass escapes in a silent kernel direction.

\subsection{Regulator independence}

\begin{corollary}[Common-refinement uniqueness]
\label{cor:v79-common-refinement}
Suppose two regulator towers admit a cofinal common-refinement tower, and
the total-variation comparison errors from each tower to the refinement are
summable.  If all three reconstructed families are tight in the same
\({\cal S}\) and use the same physical cylinder projections, their continuum
limits coincide.
\end{corollary}

\begin{proof}
Theorem \ref{thm:v79-summable-repair} gives identical limiting cylinder
marginals along each comparison because the comparison tails tend to zero.
Theorem \ref{thm:v79-tight-projective-limit} gives uniqueness in
\({\cal S}\).
\end{proof}

\subsection{Measure-construction card}

\begin{definition}[V79 projective measure card]
\label{def:v79-measure-card}
A shot regulator family passes the card if:
\begin{enumerate}
\item the positive bosonic bodies obey exact or summably approximate
      projective consistency;
\item reconstruction maps satisfy
      \eqref{eq:v79-reconstruction-consistency};
\item a compact-embedding, characteristic-functional, or equivalent
      estimate proves uniform tightness;
\item canonical constraint residuals vanish in the sense of
      \eqref{eq:v79-constraint-moment};
\item Ward and reflection-positive cylinder forms pass Proposition
      \ref{prop:v79-Ward-OS-limit};
\item every retained zero and boundary mode has its own tightness estimate;
\item fermionic coefficients pass Theorem \ref{thm:v79-fermion-limit} in a
      globally compatible determinant-line frame; and
\item cofinal regulator families satisfy the common-refinement comparison.
\end{enumerate}
\end{definition}

\begin{firewall}[No measure tightness has been acquired]
\label{fire:v79-no-tightness}
V79 proves the implication from summable adjacent errors and uniform
tightness to a continuum measure.  The programme still lacks the nonlinear
gravitational coercivity and reconstruction estimates needed for
\eqref{eq:v79-S0-moment}, the actual summable adjacent error along an
SM--Einstein shot orbit, and the global chiral determinant-line
trivialization.
\end{firewall}

\subsection{Next acquisition}

V80 is the next compulsory companion audit.  It will inspect the newest
Yang--Mills and Navier--Stokes notes for a normalized comparison estimate,
compactness mechanism, or common-refinement construction.  The post-audit
note will attack the first missing quantitative V79 row, expected to be a
uniform negative-Sobolev moment derived from the constrained reference and
the globally dominated density ratio.

\begin{assessment}[Exact status after V79]
\label{ass:v79-status}
V79 proves exact and summably repaired projective consistency, a tight
continuum-measure theorem, passage of constraint support, Ward identities,
reflection positivity, fermionic coefficient limits, and common-refinement
uniqueness.  The hypotheses requiring global gravitational coercivity,
uniform reconstruction moments, actual adjacent comparison sums, chiral
line gluing, and symmetry restoration remain unproved.  The full
SM--Einstein continuum is not constructed.
\end{assessment}
% =============================================================================
% BEGIN APPENDED TENTH-SERIES V80 MATERIAL
% =============================================================================

\section{Tenth-series V80: companion audit and local-Gaussian route correction}
\label{sec:v10v80}

\subsection{Audit identity and scope}

This is the compulsory five-version audit following V75.  The following
newest active companion notes were read without modification:
\begin{align*}
 &\texttt{Renewal\_Yang\_Mills\_Mass\_Gap\_Research\_Notes\_V6.tex},
 \\
 &\qquad
 \operatorname{SHA256}
 =
 \texttt{14B8E1D7A4E622034681495B9A94F68581D8EB43E1F1904755CE7DD12795198F},
 \\
 &\texttt{Renewal\_NS\_Stability\_Research\_Notes\_V26.tex},
 \\
 &\qquad
 \operatorname{SHA256}
 =
 \texttt{82C887F8C2C69E4F28FAD7C3DC8DE5B9099CB5A4BF431EBA1FFF3E91935978AD}.
\end{align*}
The Yang--Mills note is newer than the file inspected at the V75 gate.
The Navier--Stokes note is unchanged.  The purpose of the audit is to
transfer only arguments whose hypotheses can be stated in the
SM--Einstein variables.  Numerical powers arising from compact
Yang--Mills weak-coupling coordinates and estimates specific to the Oseen
kernel are not imported.

\subsection{What the Yang--Mills update changes}

The useful correction in the Yang--Mills companion is structural.  A
Morse normal form produced in one conditional fibre and for one fixed
boundary filling is a local chart theorem.  It does not provide a
volume-uniform, finite-range square root of the global Hessian.  A global
square root can have long range even when the Hessian is local.

This correction applies verbatim to the gravitational constraint chart of
V72--V73.  The weighted inverse used there solves the residual equation in
a local analytic chart.  It must not be reinterpreted as a local square
root of the complete gauge-fixed bosonic Hessian.  Thus the admissible
coordinate operations are:
\begin{enumerate}
\item local analytic residual coordinates with a quantitatively bounded
      inverse;
\item Schur complements and Gaussian conditionals whose covariance bounds
      are proved in the retained and eliminated sectors; and
\item exact changes of variables inside a specified chart.
\end{enumerate}
No conclusion about locality of an arbitrary spectral square root follows.

\begin{firewall}[Local normal form is not a global locality theorem]
\label{fire:v80-local-not-global}
The V72--V73 residual chart controls the nonlinear constraint map on its
proved implicit radius.  It does not diagonalize the full interacting
SM--Einstein action, does not prove a volume-uniform mass gap, and does not
make a global Hessian square root finite range.
\end{firewall}

The second transferable device is the exact add--subtract decomposition
around a Gaussian two-jet.  Let \(b\) denote local reduced coordinates,
let \(D\) be the validity domain of the nonlinear chart, and write a local
density schematically as
\[
 Q(b)=Q_2(b)\exp R_3(b)
 \qquad (b\in D),
\]
where \(Q_2\) contains the constant, linear, and quadratic Taylor terms.
For any exterior contribution \(E_{\rm phys}\), the elementary identity
\begin{equation}
 \int_D Q\,d\gamma+E_{\rm phys}
 =
 \int Q_2\,d\gamma
 +\int_D Q_2\bigl(e^{R_3}-1\bigr)\,d\gamma
 -\int_{D^c}Q_2\,d\gamma
 +E_{\rm phys}
 \label{eq:v80-add-subtract}
\end{equation}
separates four distinct owners:
\begin{enumerate}
\item an all-space Gaussian two-jet, to be evaluated exactly;
\item an interior analytic Taylor defect;
\item a Gaussian chart-exit defect; and
\item the physical exterior contribution.
\end{enumerate}
The indicator of \(D\) occurs only in the last three signed or positive
defects.  It is therefore never differentiated by the analytic forest
calculus.  Equation \eqref{eq:v80-add-subtract} is purely algebraic and
does not require a compact field space.

The third transferable ingredient is a joint exit estimate.  If a
centered Gaussian vector indexed by disjoint blocks \(y\in Y\), each of
fixed dimension \(d_*\), has covariance operator bounded by \(C_0 I\),
then exponential Markov inequality gives a factor per occupied block:
\begin{equation}
 {\mathbb P}\{\|X_y\|>R\ \hbox{for every }y\in Y\}
 \le
 \left(
  2^{d_*/2}\exp\left\{-\frac{R^2}{4C_0}\right\}
 \right)^{|Y|}.
 \label{eq:v80-joint-exit-preview}
\end{equation}
V81 will prove the estimate and its shifted, two-jet version.  This is
more useful than a union bound because it can pay the entropy of a
connected bad-block contour.

\subsection{What does not transfer}

The Yang--Mills powers in its coupling parameter use compact plaquette
coercivity, a Wilson exterior estimate, and a model-specific small-field
radius.  The corresponding gravitational exterior is noncompact and
indefinite before the conformal and gauge contours are fixed.  Consequently
the Yang--Mills exterior term cannot be renamed as an SM--Einstein bound.
It identifies the required slot but does not fill it.

Likewise, a covariance bound for a massive eliminated shell does not imply
a volume-uniform operator bound for every retained massless mode.  In
\eqref{eq:v80-joint-exit-preview}, \(C_0\) must come from a shell
covariance, a conditional covariance, or a retained-sector norm with the
zero and boundary modes removed and separately controlled.

The Navier--Stokes V26 note adds no new transferable input at this gate.
Its rank-one Oseen closure concerns a deterministic parabolic kernel and
does not supply the missing Euclidean gravitational coercivity or chiral
determinant-line data.

\subsection{Adjustment of the immediate route}

V79 proposed deriving a negative-Sobolev moment from a globally dominated
density ratio.  Such a density ratio has not been proved, so using it now
would merely move the principal noncompact exterior problem into a
hypothesis.  The audit gives a less circular order:
\begin{enumerate}
\item prove the exact finite-dimensional Gaussian two-jet formula;
\item prove interior Taylor and joint chart-exit estimates with explicit
      constants;
\item convert the resulting unnormalised signed defects into an adjacent
      total-variation comparison, conditional only on a partition lower
      bound and the still-visible physical exterior term;
\item derive reconstruction moments only after the exterior term and
      retained infrared modes have their own coercive estimates.
\end{enumerate}
The decomposition will expose exactly which part of the V79 error
\(\epsilon_n\) is already summable and which part remains an honest
gravitational bottleneck.

\begin{definition}[V80 audited transfer card]
\label{def:v80-audit-card}
At a regulator scale \(n\), the local-Gaussian route passes the audited
card when:
\begin{enumerate}
\item the nonlinear residual chart has a proved radius \(t_0\);
\item the real part of the Gaussian two-jet precision is bounded below
      uniformly on the eliminated shell;
\item its conditional covariance on each disjoint block is bounded by one
      constant \(C_0\);
\item the cubic Taylor majorant is uniform in volume;
\item the chart radius pays connected-set entropy through a joint exit
      estimate; and
\item the physical exterior and every retained kernel direction have
      estimates stated independently of the local chart.
\end{enumerate}
\end{definition}

\begin{assessment}[Status after the V80 audit]
\label{ass:v80-status}
The audit corrects the scope of the local normal form and imports the exact
add--subtract architecture and the correlated Gaussian exit mechanism.
It does not import Yang--Mills exterior coercivity.  V79's projective-limit
theorem remains valid, but none of its missing hypotheses is declared
proved.  V81 will construct the local two-jet and chart-exit bounds and
will keep the physical exterior as an explicit remainder.  The full
SM--Einstein continuum remains open.
\end{assessment}

\paragraph{Parent-version record.}
V80 was formed from
\texttt{Renewal\_SM\_Einstein\_Continuum\_Research\_Notes\_V79.tex}.
The V79 parent had \(28\,157\) counted source lines, \(1\,030\,961\) bytes,
and SHA--256
\[
 \texttt{3DCC0FD3C0ADEAD59D0BBD6E6ABB247AFDB5295217522F5ED7B31D58672F34F6}.
\]

% =============================================================================
% END APPENDED TENTH-SERIES V80 MATERIAL
% =============================================================================
% =============================================================================
% BEGIN APPENDED TENTH-SERIES V81 MATERIAL
% =============================================================================

\section{Tenth-series V81: exact Gaussian two-jet gluing and chart-exit products}
\label{sec:v10v81}

V80 isolated a local mechanism that can contribute a genuine summable
adjacent-scale error to V79 without presupposing global gravitational
coercivity.  This note proves that mechanism in finite dimension.  The
result applies to one conditioned eliminated shell or one finite family of
local residual blocks.  Retained massless, boundary, and zero modes remain
outside its scope.

\subsection{Exact absorption of the two-jet}

Let \(E\) be a real Euclidean space of dimension \(d\), and let
\(\gamma\) be its standard centred Gaussian probability measure.  For a
real symmetric operator \(B:E\to E\), a vector \(a\in E\), a constant
\(c\in{\mathbb R}\), and \(g>0\), put
\begin{equation}
 A_g:=I-g^2B,
 \qquad
 Q_{2,g}(b):=
 \exp\left\{c+g\langle a,b\rangle
       +\frac{g^2}{2}\langle b,Bb\rangle\right\}.
 \label{eq:v81-twojet-def}
\end{equation}

\begin{lemma}[Exact two-jet Gaussian]
\label{lem:v81-exact-twojet}
If \(A_g\) is positive definite, then
\begin{equation}
 Z_{2,g}:=\int_E Q_{2,g}\,d\gamma
 =
 e^c\det(A_g)^{-1/2}
 \exp\left\{\frac{g^2}{2}
       \langle a,A_g^{-1}a\rangle\right\}.
 \label{eq:v81-twojet-partition}
\end{equation}
Moreover,
\[
 Z_{2,g}^{-1}Q_{2,g}\,d\gamma
 =
 {\cal N}(m_g,A_g^{-1}),
 \qquad
 m_g=gA_g^{-1}a.
\]
\end{lemma}

\begin{proof}
The exponent relative to Lebesgue measure is
\[
 -\frac12\langle b,A_gb\rangle+g\langle a,b\rangle+c
 =
 -\frac12\langle b-m_g,A_g(b-m_g)\rangle
 +c+\frac{g^2}{2}\langle a,A_g^{-1}a\rangle.
\]
Translation and the standard Gaussian determinant formula give
\eqref{eq:v81-twojet-partition}.  The completed square also identifies the
normalised mean and covariance.
\end{proof}

Thus the constant, current, and Hessian counterterms are absorbed exactly.
There is no perturbative expansion of the determinant in this step.

\subsection{Analytic remainder and exact gluing}

Let
\[
 \ell(x)=c+\langle a,x\rangle+\frac12\langle x,Bx\rangle+R_3(x)
\]
be real analytic for \(\|x\|<t_0\).  Write its homogeneous expansion as
\[
 \ell(x)=\sum_{k\geq0}L_k[x^{\otimes k}],
 \qquad
 {\cal M}_{t_0}:=\sum_{k\geq0}\|L_k\|t_0^k<\infty.
\]
For \(0<q<1\), analyticity gives
\begin{equation}
 \sup_{\|x\|\leq qt_0}|R_3(x)|
 \leq {\cal M}_{t_0}q^3.
 \label{eq:v81-cubic-majorant}
\end{equation}
Indeed, every term of degree \(k\geq3\) gains at least the factor \(q^3\).

Choose \(R>0\) so that \(q:=gR/t_0<1\), and put
\(D_R=\{b:\|b\|\leq R\}\).  Suppose a physical finite-volume integral has
the form
\begin{equation}
 Z_{\rm phys}
 =
 \int_{D_R}Q_{2,g}(b)e^{R_3(gb)}\,d\gamma(b)
 +E_{\rm ext},
 \label{eq:v81-physical-split}
\end{equation}
where \(E_{\rm ext}\) is the integral over the complement of the proved
coordinate chart, expressed in the original variables.

\begin{proposition}[Exact local Gaussian gluing]
\label{prop:v81-exact-gluing}
Under the preceding hypotheses,
\begin{align}
 Z_{\rm phys}
 &=Z_{2,g}+D_{\rm Tay}-D_{\rm exit}+E_{\rm ext},
 \label{eq:v81-exact-gluing}\\
 D_{\rm Tay}
 &:=
 \int_{D_R}Q_{2,g}\bigl(e^{R_3(gb)}-1\bigr)\,d\gamma,
 \qquad
 D_{\rm exit}:=
 \int_{D_R^c}Q_{2,g}\,d\gamma.
 \label{eq:v81-defects}
\end{align}
Furthermore,
\begin{equation}
 |D_{\rm Tay}|
 \leq
 Z_{2,g}\left(
  \exp\{{\cal M}_{t_0}q^3\}-1
 \right).
 \label{eq:v81-Taylor-bound}
\end{equation}
\end{proposition}

\begin{proof}
Add and subtract
\(\int_{D_R}Q_{2,g}\,d\gamma\) in
\eqref{eq:v81-physical-split}, and then use
\[
 \int_{D_R}Q_{2,g}\,d\gamma
 =
 Z_{2,g}-\int_{D_R^c}Q_{2,g}\,d\gamma.
\]
This proves the identity.  On \(D_R\),
\eqref{eq:v81-cubic-majorant} and
\(|e^u-1|\leq e^{|u|}-1\) give a uniform bound for the bracket in
\(D_{\rm Tay}\).  Positivity of \(Q_{2,g}\) then gives
\eqref{eq:v81-Taylor-bound}.
\end{proof}

A useful scale choice is \(R_g=g^{-\alpha}\), with \(0<\alpha<1\).
Then
\begin{equation}
 q_g=t_0^{-1}g^{1-\alpha},
 \qquad
 |D_{\rm Tay}|/Z_{2,g}
 \leq
 \exp\{{\cal M}_{t_0}t_0^{-3}g^{3(1-\alpha)}\}-1.
 \label{eq:v81-scale-Taylor}
\end{equation}
The chart expands in the standard Gaussian coordinate while its image
\(gD_{R_g}\) shrinks inside the analytic physical chart.  This is the
source of simultaneous Taylor smallness and Gaussian exit suppression.

\subsection{One-block and joint exit estimates}

\begin{lemma}[Shifted Gaussian exit]
\label{lem:v81-shifted-exit}
Let \(X\) be a centred Gaussian vector in \({\mathbb R}^d\) with
covariance \(\Sigma\leq C_0 I\).  If \(m\in{\mathbb R}^d\) and
\(R\geq2\|m\|\), then
\begin{equation}
 {\mathbb P}\{\|X+m\|>R\}
 \leq
 2^{d/2}\exp\left\{-\frac{R^2}{16C_0}\right\}.
 \label{eq:v81-shifted-exit}
\end{equation}
\end{lemma}

\begin{proof}
The event on the left implies \(\|X\|>R/2\).  For
\(s=(4C_0)^{-1}\), diagonalisation of \(\Sigma\) gives
\[
 {\mathbb E}e^{s\|X\|^2}
 =
 \det(I-2s\Sigma)^{-1/2}
 \leq2^{d/2}.
\]
Exponential Markov inequality at threshold \(R^2/4\) proves the claim.
\end{proof}

Applying the lemma to Lemma \ref{lem:v81-exact-twojet}, if
\(A_g\geq a_0I\) and \(R\geq2\|m_g\|\), then
\begin{equation}
 D_{\rm exit}
 \leq
 Z_{2,g}\,2^{d/2}
 \exp\left\{-\frac{a_0R^2}{16}\right\}.
 \label{eq:v81-twojet-exit}
\end{equation}

The product estimate needed for contours does not require independence.

\begin{theorem}[Correlated joint block exit]
\label{thm:v81-joint-exit}
Let \(Y\) be finite, and let
\(X=(X_y)_{y\in Y}\) be a centred Gaussian vector with
\(X_y\in{\mathbb R}^{d_*}\).  If its full covariance on
\({\mathbb R}^{d_*|Y|}\) satisfies
\(\Sigma_Y\leq C_0I\), then
\begin{equation}
 {\mathbb P}\{\|X_y\|>R\ \text{for all }y\in Y\}
 \leq
 \left[
  2^{d_*/2}\exp\left\{-\frac{R^2}{4C_0}\right\}
 \right]^{|Y|}.
 \label{eq:v81-joint-centred}
\end{equation}
If deterministic means obey \(\|m_y\|\leq R/2\), then
\begin{equation}
 {\mathbb P}\{\|X_y+m_y\|>R\ \text{for all }y\in Y\}
 \leq
 \left[
  2^{d_*/2}\exp\left\{-\frac{R^2}{16C_0}\right\}
 \right]^{|Y|}.
 \label{eq:v81-joint-shifted}
\end{equation}
\end{theorem}

\begin{proof}
On the event in \eqref{eq:v81-joint-centred},
\(\|X\|^2=\sum_y\|X_y\|^2>|Y|R^2\).  With
\(s=(4C_0)^{-1}\),
\[
 {\mathbb E}e^{s\|X\|^2}
 \leq2^{d_*|Y|/2}.
\]
Markov inequality proves \eqref{eq:v81-joint-centred}.  In the shifted
case, every displayed event implies \(\|X_y\|>R/2\).  Apply the same
argument with total threshold \(|Y|R^2/4\).
\end{proof}

If the number of connected \(k\)-block sets containing a fixed root is at
most \(C_{\rm ent}^k\), the centred exit contours are summable whenever
\begin{equation}
 C_{\rm ent}\,2^{d_*/2}
 \exp\left\{-\frac{R^2}{4C_0}\right\}<1.
 \label{eq:v81-exit-entropy}
\end{equation}
For shifted two-jet contours, \(4C_0\) is replaced by \(16C_0\).
This gives a volume-free activity criterion after root normalisation.

\subsection{From unnormalised defects to probability measures}

We use the total-variation norm
\[
 \|\sigma\|_{\rm TV}:=
 \sup_{\|f\|_\infty\leq1}\left|\int f\,d\sigma\right|
\]
for a finite signed measure \(\sigma\).

\begin{lemma}[Normalisation cost]
\label{lem:v81-normalisation}
Let \(\alpha,\beta\) be positive finite measures with masses
\(A=\alpha(1)\), \(B=\beta(1)\), and
\(\min(A,B)\geq z_*>0\).  If
\(\|\alpha-\beta\|_{\rm TV}\leq\delta\), then
\begin{equation}
 \left\|\frac{\alpha}{A}-\frac{\beta}{B}\right\|_{\rm TV}
 \leq\frac{2\delta}{z_*}.
 \label{eq:v81-normalisation-cost}
\end{equation}
\end{lemma}

\begin{proof}
Since \(|A-B|\leq\delta\),
\[
 \left\|\frac{\alpha}{A}-\frac{\beta}{B}\right\|_{\rm TV}
 \leq
 \frac{\|\alpha-\beta\|_{\rm TV}}{A}
 +B\left|\frac1A-\frac1B\right|
 =
 \frac{\delta+|A-B|}{A}
 \leq\frac{2\delta}{z_*}.
\]
\end{proof}

\begin{proposition}[Adjacent comparison owner bound]
\label{prop:v81-adjacent-owner}
Suppose an unnormalised adjacent comparison at scale \(n\) has the exact
signed decomposition
\[
 (p_{n+1,n})_*\nu_{n+1}-\nu_n
 =
 \Delta_{{\rm Tay},n}
 +\Delta_{{\rm exit},n}
 +\Delta_{{\rm ext},n}
 +\Delta_{{\rm alg},n},
\]
where the first two terms arise from
\eqref{eq:v81-exact-gluing}, the exterior term is measured in the original
variables, and \(\Delta_{{\rm alg},n}\) contains the already-owned
finite-dimensional reblocking and counterterm discrepancies.  If both
masses are at least \(z_*\), then the normalised measures satisfy
\begin{equation}
 \left\|(p_{n+1,n})_*\mu_{n+1}-\mu_n\right\|_{\rm TV}
 \leq\frac{2}{z_*}
 \sum_{\bullet\in\{{\rm Tay},{\rm exit},{\rm ext},{\rm alg}\}}
 \|\Delta_{\bullet,n}\|_{\rm TV}.
 \label{eq:v81-adjacent-TV}
\end{equation}
\end{proposition}

\begin{proof}
The triangle inequality bounds the unnormalised difference by the sum in
\eqref{eq:v81-adjacent-TV}.  Lemma
\ref{lem:v81-normalisation} then applies.
\end{proof}

Equations \eqref{eq:v81-scale-Taylor},
\eqref{eq:v81-twojet-exit}, and
\eqref{eq:v81-exit-entropy} can now fill the Taylor and chart-exit rows of
the V79 summable-repair hypothesis once their constants are uniform along
the shot orbit and the forest converts local activities to the stated
global signed-measure norms.  The other two rows remain separate, so no
exterior estimate is hidden in the local normal form.

\subsection{Application boundary for the SM--Einstein shell}

For the V72 constraint coordinate, \(g\) denotes the scale-dependent
small amplitude after \(v=\sqrt{\varepsilon}\,y\), and \(E\) is a
conditioned finite-dimensional eliminated shell.  Lemma
\ref{lem:v81-exact-twojet} applies if the gauge-fixed, contour-chosen real
part of the shell precision obeys \(A_g\geq a_0I\).  Theorem
\ref{thm:v81-joint-exit} applies if every selected family of disjoint
shell blocks has conditional covariance at most \(C_0I\).

\begin{firewall}[Infrared and exterior hypotheses remain visible]
\label{fire:v81-infrared-exterior}
Neither \(a_0\) nor \(C_0\) is asserted for the unprojected massless
graviton, photon, boundary, or global gauge sector.  Those modes must stay
in the retained field or receive an independent quotient and moment
estimate.  The term \(\Delta_{{\rm ext},n}\) in
\eqref{eq:v81-adjacent-TV} is not bounded by Gaussian chart exit; it is the
original noncompact physical exterior and still requires gravitational
coercivity on the chosen contour.
\end{firewall}

\begin{assessment}[Exact status after V81]
\label{ass:v81-status}
V81 proves exact two-jet absorption, exact local add--subtract gluing,
analytic Taylor control, correlated joint Gaussian exit bounds, contour
entropy payment, and the normalisation lemma converting unnormalised
defects into adjacent total variation.  These are finite-dimensional
shell theorems with explicit hypotheses.  Uniform shell covariance,
physical exterior coercivity, reconstruction moments, determinant-line
gluing, symmetry restoration, and the full SM--Einstein continuum remain
unproved.
\end{assessment}

\subsection{Next acquisition}

V82 will insert the V72 nonlinear residual chart into the V81 split.  It
will derive a scale-explicit conditional error card and identify a
summability window for the chart radius \(R_n=g_n^{-\alpha}\).  It will
then test whether the same window yields a uniform negative-Sobolev moment
for the eliminated shells without assuming control of the retained
infrared variables.

\paragraph{Parent-version record.}
V81 was formed from
\texttt{Renewal\_SM\_Einstein\_Continuum\_Research\_Notes\_V80.tex}.
The V80 parent had \(28\,348\) counted source lines, \(1\,039\,324\) bytes,
and SHA--256
\[
 \texttt{856330D04B93730ADEA3A0F919068F7672DB72FD75FD03C942241C2D238A32D8}.
\]

% =============================================================================
% END APPENDED TENTH-SERIES V81 MATERIAL
% =============================================================================
% =============================================================================
% BEGIN APPENDED TENTH-SERIES V82 MATERIAL
% =============================================================================

\section{Tenth-series V82: residual-shell scale window and the global-TV obstruction}
\label{sec:v10v82}

This note inserts the exact V72 residual scaling into V81.  It produces a
summable local chart activity for every geometric penalty schedule.  It
also finds an upstream obstruction: total variation of the entire growing
lattice is generally the wrong topology for a local continuum
construction.  A small error per cell can make the global total variation
tend to its maximal value.  The V79 repair theorem is correct as a
sufficient statement, but its global hypothesis must be replaced by a
cylinder-local comparison theorem before it can serve as the main
construction route.

\subsection{Conditional two-jet in exact residual coordinates}

Fix the retained variable \(u\).  In the exact residual coordinate of V72,
write
\begin{equation}
 H_\varepsilon(u,y)
 =
 H(u,\sqrt{\varepsilon}y),
 \qquad
 \ell_u(v):=\log H(u,v),
 \label{eq:v82-residual-log}
\end{equation}
on a region where \(H>0\).  Suppose, uniformly in the admissible \(u\),
\begin{equation}
 \ell_u(v)
 =
 c_u+\langle a_u,v\rangle
 +\frac12\langle v,B_uv\rangle+R_{3,u}(v),
 \label{eq:v82-residual-twojet}
\end{equation}
with
\begin{equation}
 \|a_u\|\leq M_1,\qquad
 \|B_u\|\leq M_2,\qquad
 \sum_{k\geq0}\|L_{k,u}\|t_0^k\leq{\cal M}.
 \label{eq:v82-uniform-analytic-data}
\end{equation}
Here \(L_{k,u}\) are the homogeneous Taylor tensors of \(\ell_u\).
The Gaussian penalty is exactly \(d\gamma_m(y)\), by Theorem
\ref{thm:v72-exact-scaling}.  Therefore the V81 coupling is
\begin{equation}
 g=\sqrt{\varepsilon},
 \qquad
 A_{\varepsilon,u}=I-\varepsilon B_u,
 \qquad
 m_{\varepsilon,u}
 =
 \sqrt{\varepsilon}\,A_{\varepsilon,u}^{-1}a_u.
 \label{eq:v82-residual-Gaussian}
\end{equation}

\begin{theorem}[Residual chart scale window]
\label{thm:v82-scale-window}
Assume \eqref{eq:v82-uniform-analytic-data}.  Fix
\(0<\alpha<1\) and \(0<a_0<1\), and suppose
\(\varepsilon M_2\leq1-a_0\).  Put
\begin{equation}
 R_\varepsilon:=\varepsilon^{-\alpha/2},
 \qquad
 q_\varepsilon:=
 \frac{\sqrt{\varepsilon}R_\varepsilon}{t_0}
 =
 t_0^{-1}\varepsilon^{(1-\alpha)/2}.
 \label{eq:v82-radius-choice}
\end{equation}
For all sufficiently small \(\varepsilon\), uniformly in \(u\):
\begin{enumerate}
\item \(A_{\varepsilon,u}\geq a_0I\);
\item the ball \(\{\|y\|\leq R_\varepsilon\}\) lies in the residual chart;
\item \(\|m_{\varepsilon,u}\|\leq R_\varepsilon/2\);
\item the relative Taylor defect obeys
\begin{equation}
 \frac{|D_{{\rm Tay},\varepsilon}(u)|}
      {Z_{2,\varepsilon}(u)}
 \leq
 \exp\left\{
  {\cal M}t_0^{-3}
  \varepsilon^{3(1-\alpha)/2}
 \right\}-1;
 \label{eq:v82-Taylor-activity}
\end{equation}
and
\item the relative one-block chart-exit defect in residual dimension \(m\)
obeys
\begin{equation}
 \frac{D_{{\rm exit},\varepsilon}(u)}
      {Z_{2,\varepsilon}(u)}
 \leq
 2^{m/2}
 \exp\left\{-\frac{a_0}{16}
                  \varepsilon^{-\alpha}\right\}.
 \label{eq:v82-exit-activity}
\end{equation}
\end{enumerate}
\end{theorem}

\begin{proof}
The operator estimate follows from
\[
 \langle z,A_{\varepsilon,u}z\rangle
 \geq(1-\varepsilon\|B_u\|)\|z\|^2
 \geq a_0\|z\|^2.
\]
Since \(q_\varepsilon\to0\), the image radius
\(\sqrt{\varepsilon}R_\varepsilon=q_\varepsilon t_0\) is eventually
strictly smaller than \(t_0\).  Also,
\[
 \|m_{\varepsilon,u}\|
 \leq\frac{M_1}{a_0}\sqrt{\varepsilon},
 \qquad
 \frac{\|m_{\varepsilon,u}\|}{R_\varepsilon}
 \leq
 \frac{M_1}{a_0}\varepsilon^{(1+\alpha)/2}\longrightarrow0.
\]
Equations \eqref{eq:v82-Taylor-activity} and
\eqref{eq:v82-exit-activity} are respectively
\eqref{eq:v81-Taylor-bound} and \eqref{eq:v81-twojet-exit} with
\(g=\sqrt{\varepsilon}\) and \(R=R_\varepsilon\).
\end{proof}

\begin{corollary}[Geometric-scale summability]
\label{cor:v82-geometric-sum}
If \(\varepsilon_n\leq C_\varepsilon\rho^n\) for some \(0<\rho<1\),
then, for every \(0<\alpha<1\),
\begin{align}
 \sum_n\sup_u
 \frac{|D_{{\rm Tay},n}(u)|}{Z_{2,n}(u)}
 &<\infty,
 \label{eq:v82-Taylor-sum}\\
 \sum_n\sup_u
 \frac{D_{{\rm exit},n}(u)}{Z_{2,n}(u)}
 &<\infty.
 \label{eq:v82-exit-sum}
\end{align}
For fixed block dimension \(m_*\) and connected-set entropy
\(C_{\rm ent}\), the shifted contour condition
\begin{equation}
 C_{\rm ent}2^{m_*/2}
 \exp\left\{-\frac{a_0}{16}\varepsilon_n^{-\alpha}\right\}<1
 \label{eq:v82-eventual-contour-gate}
\end{equation}
holds for all sufficiently large \(n\).
\end{corollary}

\begin{proof}
For small \(x\), \(e^x-1\leq2x\).  The first series is therefore bounded
in its tail by a constant times
\(\sum_n\rho^{3(1-\alpha)n/2}\).  The second is superexponentially small
in \(n\).  The left side of
\eqref{eq:v82-eventual-contour-gate} tends to zero.
\end{proof}

This closes the local penalty-radius question: there is no conflict
between analytic Taylor smallness and Gaussian exit entropy.  The open
question is whether the constants in
\eqref{eq:v82-uniform-analytic-data}, the conditional shell precision,
and the physical exterior are uniform for the actual coupled
SM--Einstein shot orbit.

\subsection{What shell Gaussian control gives in negative Sobolev norm}

Let \({\cal H}_{-s}\) be the regulated negative-Sobolev reconstruction
space.  If \(T_n:E_n\to{\cal H}_{-s}\) embeds the \(n\)-th Gaussian shell,
the shifted two-jet law has covariance at most \(a_0^{-1}I\).

\begin{proposition}[One-shell reconstruction moment]
\label{prop:v82-one-shell-moment}
For \(Y_n\sim{\cal N}(m_n,\Sigma_n)\) with
\(\Sigma_n\leq a_0^{-1}I\),
\begin{equation}
 {\mathbb E}\|T_nY_n\|_{-s}^2
 \leq
 a_0^{-1}\|T_n\|_{\rm HS}^2+\|T_nm_n\|_{-s}^2.
 \label{eq:v82-one-shell-moment}
\end{equation}
\end{proposition}

\begin{proof}
Write \(Y_n=m_n+X_n\), where \(X_n\) is centred.  The cross term has zero
expectation, and
\[
 {\mathbb E}\|T_nX_n\|_{-s}^2
 =
 \operatorname{Tr}(T_n\Sigma_nT_n^*)
 \leq a_0^{-1}\operatorname{Tr}(T_nT_n^*).
\]
\end{proof}

The usual free order-counting consequence can be stated without asserting
that the interacting gravitational covariance has that order.

\begin{corollary}[Order-minus-two shell criterion]
\label{cor:v82-order-two}
Suppose a \(D\)-dimensional dyadic shell of frequency
\(\Lambda_n\) contains at most \(C_N\Lambda_n^D\) modes, its centred
covariance eigenvalues are at most \(C_C\Lambda_n^{-2}\), and the
\({\cal H}_{-s}\) weight is at most \(C_s\Lambda_n^{-2s}\).
Then
\begin{equation}
 {\mathbb E}\|X_n\|_{-s}^2
 \leq
 C_NC_CC_s\Lambda_n^{D-2-2s}.
 \label{eq:v82-symbol-moment}
\end{equation}
Consequently the ultraviolet shell moments are summable when
\begin{equation}
 s>\frac{D-2}{2}.
 \label{eq:v82-free-regularity-threshold}
\end{equation}
In four dimensions this condition is \(s>1\).
\end{corollary}

\begin{proof}
Multiply the number of shell modes by the covariance and Sobolev weights.
For dyadic \(\Lambda_n\), the resulting geometric series converges exactly
when \(D-2-2s<0\).
\end{proof}

This proves a reference-shell tightness mechanism.  It does not yet prove
a moment for the physical measure: a total-variation error controls
bounded observables, whereas \(\|\phi\|_{-s}^2\) is unbounded.  A weighted
comparison norm or a uniform higher-moment truncation argument is needed.

\subsection{Why global total variation is too strong}

The local activities in Theorem \ref{thm:v82-scale-window} are
volume-independent.  They do not imply that the total variation of two
entire \(N\)-cell measures is small uniformly in \(N\).

\begin{proposition}[Product-measure obstruction]
\label{prop:v82-product-obstruction}
Let \(0<p\neq q<1\), and let
\(\mu_N=\operatorname{Bernoulli}(p)^{\otimes N}\) and
\(\nu_N=\operatorname{Bernoulli}(q)^{\otimes N}\).  With the norm
convention of V81,
\begin{equation}
 \lim_{N\to\infty}\|\mu_N-\nu_N\|_{\rm TV}=2,
 \label{eq:v82-product-TV}
\end{equation}
although the one-site total variation can be made arbitrarily small by
choosing \(q\) close to \(p\).
\end{proposition}

\begin{proof}
Choose \(\delta<|p-q|/3\), and let \(A_N\) be the event that the empirical
mean lies within \(\delta\) of \(p\).  The weak law of large numbers gives
\(\mu_N(A_N)\to1\) and \(\nu_N(A_N)\to0\).  For
\(f_N=2{\bf1}_{A_N}-1\), one has \(\|f_N\|_\infty=1\) and
\[
 \left|\int f_N\,d(\mu_N-\nu_N)\right|
 =
 2|\mu_N(A_N)-\nu_N(A_N)|\longrightarrow2.
\]
The total variation norm of a difference of probabilities is at most
\(2\), proving the limit.
\end{proof}

Therefore a per-block V81 chart error can accumulate extensively even
when its rooted polymer expansion is perfectly summable.  Requiring the
global V79 error
\(\|(p_{n+1,n})_*\mu_{n+1}-\mu_n\|_{\rm TV}\) to vanish uniformly in the
growing volume would exclude ordinary locally perturbed product systems.

The correct projective statement must test observables depending on a
fixed bounded set of coarse blocks.  If \(A\) is such a set, define the
cylinder seminorm
\begin{equation}
 \|\sigma\|_A
 :=
 \sup\left\{
  \left|\int f\,d\sigma\right|:
  \|f\|_\infty\leq1,\ 
  f\text{ depends only on fields in }A
 \right\}.
 \label{eq:v82-cylinder-seminorm}
\end{equation}
A local comparison should have a bound such as
\begin{equation}
 \|\,(p_{n+1,n})_*\mu_{n+1}-\mu_n\,\|_A
 \leq C|A|\eta_n,
 \qquad
 \sum_n\eta_n<\infty,
 \label{eq:v82-local-comparison-target}
\end{equation}
or a weighted quasi-local variant with exponential decay away from \(A\).
This is compatible with an extensive number of harmless remote defects.

\begin{firewall}[Two distinct gaps]
\label{fire:v82-two-gaps}
The residual scale window closes neither of the following gaps:
\begin{enumerate}
\item the physical exterior \(\Delta_{{\rm ext},n}\) still needs
      noncompact gravitational coercivity; and
\item the reference-shell moment
      \eqref{eq:v82-symbol-moment} still needs transfer to the interacting
      measure in a weighted norm.
\end{enumerate}
Moreover, V79's global total-variation repair is now retained only as a
strong sufficient theorem.  It is not assumed as the natural
infinite-volume comparison topology.
\end{firewall}

\begin{assessment}[Exact status after V82]
\label{ass:v82-status}
V82 proves that the V72 residual chart and V81 Gaussian split admit the
summable window \(R_n=\varepsilon_n^{-\alpha/2}\) for every
\(0<\alpha<1\) under uniform analytic and conditional-covariance
constants.  It proves the reference order-minus-two regularity threshold
\(s>(D-2)/2\), and proves by a product example that global total variation
is unsuitable for accumulating local errors.  Uniform gravitational
exterior control, local-cylinder projective repair, weighted moment
transfer, infrared closure, determinant-line gluing, and the full
SM--Einstein continuum remain open.
\end{assessment}

\subsection{Next acquisition}

V83 will replace the global repair hypothesis by a compatible family of
cylinder seminorm estimates.  It will prove existence and uniqueness of
limiting cylinder marginals under summable local errors, including a
quasi-local leakage kernel, and state the additional tightness condition
needed to extend that cylinder functional to a Radon measure on the
distribution space.

\paragraph{Parent-version record.}
V82 was formed from
\texttt{Renewal\_SM\_Einstein\_Continuum\_Research\_Notes\_V81.tex}.
The V81 parent had \(28\,720\) counted source lines, \(1\,051\,201\) bytes,
and SHA--256
\[
 \texttt{3093988BF52EC9871E8A0B5FED6A20AC3D5226439EC214AF0744C7A8294BE9CC}.
\]

% =============================================================================
% END APPENDED TENTH-SERIES V82 MATERIAL
% =============================================================================
% =============================================================================
% BEGIN APPENDED TENTH-SERIES V83 MATERIAL
% =============================================================================

\section{Tenth-series V83: cylinder-local projective repair}
\label{sec:v10v83}

V82 showed that global total variation is too strong for an extensive
system.  This note proves the local replacement.  The construction has
two logically separate parts.  Summable cylinder errors determine a
unique compatible family of local marginals.  Tightness in the chosen
distribution topology is then required to realize that family as a Radon
continuum field measure.

\subsection{Local marginals and seminorms}

Let \((X_n,{\cal B}_n)\) be standard Borel regulator spaces with measurable
coarse maps
\[
 p_{m,n}:X_m\longrightarrow X_n,
 \qquad
 p_{\ell,n}=p_{m,n}\circ p_{\ell,m}
 \quad(\ell\geq m\geq n).
\]
Let \(\mu_n\) be probability measures on \(X_n\).  At scale \(k\), let
\(A\) be a finite set of coarse cells, let \(X_{k,A}\) be its local state
space, and let
\[
 r_{k,A}:X_k\longrightarrow X_{k,A}
\]
be the restriction map.  Define the local seminorm of a finite signed
measure \(\sigma\) on \(X_k\) by
\begin{equation}
 \|\sigma\|_{k,A}
 :=
 \left\|(r_{k,A})_*\sigma\right\|_{\rm TV}.
 \label{eq:v83-local-seminorm}
\end{equation}
Equivalently, this is the supremum against bounded measurable functions
of the \(A\)-variables with sup norm at most one.

For \(n\geq k\), put
\begin{equation}
 \nu_{n;k,A}
 :=
 (r_{k,A}\circ p_{n,k})_*\mu_n.
 \label{eq:v83-descended-marginal}
\end{equation}

\begin{theorem}[Cylinder-local repair]
\label{thm:v83-cylinder-repair}
Suppose that for every pair \((k,A)\) there is a summable sequence
\(e_n(k,A)\), \(n\geq k\), such that
\begin{equation}
 \left\|
 (p_{n+1,k})_*\mu_{n+1}
 -
 (p_{n,k})_*\mu_n
 \right\|_{k,A}
 \leq e_n(k,A).
 \label{eq:v83-adjacent-local-error}
\end{equation}
Then:
\begin{enumerate}
\item \(\nu_{n;k,A}\) converges in total variation to a probability
      measure \(\bar\mu_{k,A}\) on \(X_{k,A}\);
\item for \(N>n\geq k\),
\begin{equation}
 \|\nu_{N;k,A}-\nu_{n;k,A}\|_{\rm TV}
 \leq\sum_{j=n}^{N-1}e_j(k,A);
 \label{eq:v83-local-Cauchy}
\end{equation}
and
\item all limit marginals are compatible under restriction and
      coarse-graining maps.
\end{enumerate}
Hence they define a unique probability measure on the cylinder
sigma-algebra of the countable projective system.
\end{theorem}

\begin{proof}
Equation \eqref{eq:v83-adjacent-local-error} is exactly
\(\|\nu_{n+1;k,A}-\nu_{n;k,A}\|_{\rm TV}\leq e_n(k,A)\).
The triangle inequality gives \eqref{eq:v83-local-Cauchy}.  Finite signed
measures are complete in total variation, so a limit
\(\bar\mu_{k,A}\) exists.  Positivity and total mass one pass to the
total-variation limit.

Let \(s:X_{k,A}\to X_{\ell,B}\) be any local map induced by restriction
or coarse graining.  At every finite \(n\) for which both sides are
defined,
\[
 s_*\nu_{n;k,A}=\nu_{n;\ell,B}.
\]
Pushforward contracts total variation.  Taking the limit gives
\(s_*\bar\mu_{k,A}=\bar\mu_{\ell,B}\).  The countable standard-Borel
projective-limit theorem now gives a cylinder probability measure.
Uniqueness holds because its values on every finite cylinder are fixed by
the displayed marginals.
\end{proof}

No global total-variation estimate appears.  The error may grow like
\(|A|\); for each fixed cylinder it remains summable.

\subsection{Rooted polymers imply a cylinder estimate}

The local Gaussian chart errors of V82 enter through polymers.  The
following elementary contact estimate records the needed cancellation of
remote polymers after normalization.

\begin{lemma}[Rooted contact bound]
\label{lem:v83-rooted-contact}
Let \(f\) be supported in a finite cell set \(A\), and suppose an adjacent
expectation difference has an absolutely convergent linked representation
\[
 \Delta_n(f)=\sum_P\Phi_{n,f}(P),
\]
over connected polymers \(P\), with
\[
 \Phi_{n,f}(P)=0\quad\text{if }P\cap A=\varnothing,
 \qquad
 |\Phi_{n,f}(P)|\leq2\|f\|_\infty w_n(P),
\]
where \(w_n(P)\geq0\).  If
\begin{equation}
 \sup_x\sum_{P\ni x}w_n(P)\leq\eta_n,
 \label{eq:v83-rooted-activity}
\end{equation}
then
\begin{equation}
 |\Delta_n(f)|
 \leq2\|f\|_\infty |A|\eta_n.
 \label{eq:v83-contact-estimate}
\end{equation}
\end{lemma}

\begin{proof}
Only polymers meeting \(A\) contribute.  Overcount each such polymer by
one selected contact cell:
\[
 \sum_{P:P\cap A\neq\varnothing}w_n(P)
 \leq
 \sum_{x\in A}\sum_{P\ni x}w_n(P)
 \leq |A|\eta_n.
\]
Multiply by \(2\|f\|_\infty\).
\end{proof}

\begin{corollary}[Local chart errors feed projective repair]
\label{cor:v83-chart-to-repair}
Suppose the normalized adjacent comparison has the linked representation
of Lemma \ref{lem:v83-rooted-contact}, and suppose its rooted activity is
bounded by
\begin{align}
 \eta_n
 \leq C_{\rm F}\bigg[
 &\exp\left\{
  C_{\rm Tay}\varepsilon_n^{3(1-\alpha)/2}
 \right\}-1
 +C_{\rm exit}
  \exp\{-c_{\rm exit}\varepsilon_n^{-\alpha}\}
 \nonumber\\
 &\quad+\eta_{{\rm ext},n}
 +\eta_{{\rm alg},n}
 \bigg].
 \label{eq:v83-owner-activity}
\end{align}
If \(\varepsilon_n\) is geometric and
\[
 \sum_n(\eta_{{\rm ext},n}+\eta_{{\rm alg},n})<\infty,
\]
then Theorem \ref{thm:v83-cylinder-repair} applies with
\(e_n(k,A)=2|A|\eta_n\) for exactly local coarse maps.
\end{corollary}

\begin{proof}
Corollary \ref{cor:v82-geometric-sum} makes the first two rows of
\eqref{eq:v83-owner-activity} summable.  The hypotheses make the other
rows summable.  Lemma \ref{lem:v83-rooted-contact} gives
\eqref{eq:v83-adjacent-local-error}.
\end{proof}

The constant \(C_{\rm F}\) is the actual connected-forest summation
constant.  The corollary does not assert that the SM--Einstein exterior
activity \(\eta_{{\rm ext},n}\) has been bounded.

\subsection{Quasi-local coarse maps}

Exact locality may be replaced by summable leakage.  Suppose the
scale-\(n\) pullback of every \(A\)-cylinder function \(f\) admits, for
each integer \(r\geq0\), an approximation \(f_{n,r}\) supported in the
\(r\)-neighbourhood \(A^{+r}\) and satisfying
\begin{equation}
 \|f\circ p_{n,k}-f_{n,r}\|_\infty
 \leq C_{\rm ql}|A|e^{-\kappa r}\|f\|_\infty.
 \label{eq:v83-quasilocal-approx}
\end{equation}
Assume also that the adjacent signed difference
\[
 \delta_n:=(p_{n+1,n})_*\mu_{n+1}-\mu_n
\]
obeys
\begin{equation}
 \|\delta_n\|_{n,B}\leq\eta_n|B|
 \label{eq:v83-local-delta}
\end{equation}
for every finite \(B\), and that lattice neighbourhoods satisfy
\begin{equation}
 |A^{+r}|\leq C_D|A|(1+r)^D.
 \label{eq:v83-neighbourhood-growth}
\end{equation}

\begin{proposition}[Quasi-local leakage budget]
\label{prop:v83-quasilocal-budget}
For any \(r_n\geq0\),
\begin{equation}
 \left|
 \int f\circ p_{n,k}\,d\delta_n
 \right|
 \leq
 \|f\|_\infty\left[
 C_D|A|\eta_n(1+r_n)^D
 +2C_{\rm ql}|A|e^{-\kappa r_n}
 \right].
 \label{eq:v83-quasilocal-error}
\end{equation}
If \(0<\eta_n\leq e^{-1}\) and
\begin{equation}
 \sum_n
 \eta_n\left(1+\log\frac1{\eta_n}\right)^D<\infty,
 \label{eq:v83-log-summability}
\end{equation}
then the right side is summable for
\(r_n=\kappa^{-1}\log(1/\eta_n)\).
\end{proposition}

\begin{proof}
Test \(\delta_n\) first against \(f_{n,r_n}\) and use
\eqref{eq:v83-local-delta} and
\eqref{eq:v83-neighbourhood-growth}.  Since \(\delta_n\) is the
difference of two probability measures, its total-variation norm is at
most \(2\); hence the approximation error contributes at most the second
term of \eqref{eq:v83-quasilocal-error}.  With the stated choice,
\(e^{-\kappa r_n}=\eta_n\), and
\eqref{eq:v83-log-summability} proves summability.
\end{proof}

For a geometric scale error, the logarithmic factor in
\eqref{eq:v83-log-summability} is harmless.  Thus exponential locality of
the coarse map is sufficient for cylinder repair.  A merely algebraic
locality tail requires a separate exponent calculation and is not covered
by this proposition.

\subsection{From cylinders to a Radon distribution measure}

Let \({\cal S}\) be a Polish distribution space, let
\(\iota_n:X_n\to{\cal S}\) be reconstruction maps, and let
\(\Pi_{k,A}:{\cal S}\to X_{k,A}\) be continuous cylinder projections.
Assume
\begin{equation}
 \Pi_{k,A}\circ\iota_n
 =
 r_{k,A}\circ p_{n,k}
 \qquad(n\geq k).
 \label{eq:v83-reconstruction-cylinder}
\end{equation}

\begin{theorem}[Tight realization of the repaired cylinders]
\label{thm:v83-tight-realization}
Suppose Theorem \ref{thm:v83-cylinder-repair} holds,
\(\{(\iota_n)_*\mu_n\}\) is tight in \({\cal S}\), and the countable family
\(\{\Pi_{k,A}\}\) separates probability measures on \({\cal S}\).
Then \((\iota_n)_*\mu_n\) converges weakly to a unique Radon probability
measure \(\mu\) on \({\cal S}\), and
\begin{equation}
 (\Pi_{k,A})_*\mu=\bar\mu_{k,A}.
 \label{eq:v83-realized-marginal}
\end{equation}
\end{theorem}

\begin{proof}
Tightness and Prokhorov's theorem give a weakly convergent subsequence with
Radon limit \(\mu\).  By
\eqref{eq:v83-reconstruction-cylinder}, its \((k,A)\)-marginal is
\(\nu_{n;k,A}\).  For bounded continuous local test functions, weak
convergence and Theorem \ref{thm:v83-cylinder-repair} identify the limit
as \(\bar\mu_{k,A}\).  Thus every subsequential limit has the same
separating cylinder marginals.  It follows that all subsequential limits
coincide.  Tightness then implies convergence of the full sequence, and
\eqref{eq:v83-realized-marginal} follows.
\end{proof}

\begin{corollary}[Local passage of identities]
\label{cor:v83-local-identities}
Any Ward identity, reflection-positive quadratic form, or constraint
identity expressed through finitely many bounded continuous cylinder
observables passes to the repaired limit.  An unbounded identity passes
only after a uniform-integrability estimate for its insertions.
\end{corollary}

\begin{proof}
The bounded case follows from weak convergence in Theorem
\ref{thm:v83-tight-realization}.  Uniform integrability permits the
standard truncation limit in the unbounded case.
\end{proof}

\begin{firewall}[Cylinder repair does not prove tightness]
\label{fire:v83-no-tightness}
Summable local comparison controls every fixed observation window but does
not prevent escape in a global zero mode, a boundary charge, or an
ultraviolet distribution norm.  Theorem
\ref{thm:v83-tight-realization} therefore retains tightness as a separate
hypothesis.  The reference-shell estimate of V82 is one input toward that
hypothesis, not its proof for the interacting measure.
\end{firewall}

\begin{assessment}[Exact status after V83]
\label{ass:v83-status}
V83 replaces the unsuitable global-TV route by a rigorous cylinder-local
repair theorem.  It proves the rooted-polymer contact bound, handles
exponentially quasi-local coarse maps, and identifies the exact tightness
condition that realizes the repaired cylinders as a unique Radon measure.
The actual exterior activity, weighted moment transfer, retained
infrared-mode tightness, determinant-line gluing, symmetry restoration,
and the full SM--Einstein continuum remain open.
\end{assessment}

\subsection{Next acquisition}

V84 will construct a weighted local comparison norm that transfers the
order-minus-two reference moment of V82 to the interacting cylinder
family.  It will distinguish ultraviolet shell moments, retained
massless modes, and finite-dimensional zero or boundary modes, so that
failure in one sector cannot be concealed by tightness in another.

\paragraph{Parent-version record.}
V83 was formed from
\texttt{Renewal\_SM\_Einstein\_Continuum\_Research\_Notes\_V82.tex}.
The V82 parent had \(29\,063\) counted source lines, \(1\,062\,976\) bytes,
and SHA--256
\[
 \texttt{3B8F0F1CE1B875B38722161F196194EA29245C05098D1BB8A11B4B7C6E6D5E40}.
\]

% =============================================================================
% END APPENDED TENTH-SERIES V83 MATERIAL
% =============================================================================
% =============================================================================
% BEGIN APPENDED TENTH-SERIES V84 MATERIAL
% =============================================================================

\section{Tenth-series V84: weighted comparison and Sobolev tightness transfer}
\label{sec:v10v84}

V82 obtained the free order-minus-two shell moment, while V83 repaired
bounded cylinder observables.  A moment is unbounded and therefore does
not pass under ordinary total variation without uniform integrability.
This note supplies the missing weighted norm and proves a sectorwise
tightness theorem.

\subsection{Weighted variation}

For a measurable weight \(W\geq1\) and a finite signed measure \(\sigma\),
define
\begin{equation}
 \|\sigma\|_W
 :=
 \sup_{|f|\leq W}\left|\int f\,d\sigma\right|
 =
 \int W\,d|\sigma|.
 \label{eq:v84-weighted-TV}
\end{equation}
The equality is the polar decomposition of a signed measure.  In
particular, \(\|\sigma\|_{\rm TV}\leq\|\sigma\|_W\).

\begin{lemma}[Moment transfer]
\label{lem:v84-moment-transfer}
If \(\mu,\nu\) are probabilities and
\(\|\mu-\nu\|_W\leq\delta\), then
\begin{equation}
 \int W\,d\mu\leq\int W\,d\nu+\delta.
 \label{eq:v84-moment-transfer}
\end{equation}
More generally, every measurable \(F\) with \(|F|\leq W\) satisfies
\[
 |\mu(F)-\nu(F)|\leq\delta.
\]
\end{lemma}

\begin{proof}
Both statements follow immediately from the definition with \(f=W\) and
\(f=F\), using monotone truncations if the functions are unbounded.
\end{proof}

Normalisation has a visible weighted cost.

\begin{lemma}[Weighted normalisation]
\label{lem:v84-weighted-normalisation}
Let \(\alpha,\beta\) be positive finite measures of masses \(A,B\geq z_*>0\).
Suppose
\[
 \|\alpha-\beta\|_W\leq\delta,
 \qquad
 \frac{\beta(W)}{B}\leq M_W.
\]
Then
\begin{equation}
 \left\|\frac{\alpha}{A}-\frac{\beta}{B}\right\|_W
 \leq
 \frac{(1+M_W)\delta}{z_*}.
 \label{eq:v84-weighted-normalisation}
\end{equation}
\end{lemma}

\begin{proof}
Because \(W\geq1\), \(|A-B|\leq\delta\).  Therefore
\begin{align*}
 \left\|\frac{\alpha}{A}-\frac{\beta}{B}\right\|_W
 &\leq
 \frac{\|\alpha-\beta\|_W}{A}
 +\left|\frac1A-\frac1B\right|\beta(W)\\
 &\leq
 \frac{\delta}{z_*}
 +\frac{|A-B|}{A}\frac{\beta(W)}B
 \leq\frac{(1+M_W)\delta}{z_*}.
\end{align*}
\end{proof}

\subsection{Weighted local Gaussian defects}

Let \(\nu_{2,g}\) be the normalized shifted Gaussian of Lemma
\ref{lem:v81-exact-twojet}.  On one fixed-dimensional shell block choose
\begin{equation}
 W_T(b):=1+\|Tb\|_{\cal H}^2,
 \label{eq:v84-block-weight}
\end{equation}
where \(T:E\to{\cal H}\) is the reconstruction into a Hilbert norm.

\begin{proposition}[Weighted Taylor and exit bounds]
\label{prop:v84-weighted-local-defects}
Assume the hypotheses of Proposition \ref{prop:v81-exact-gluing}.  If
\[
 \delta_{\rm Tay}:=
 \exp\{{\cal M}_{t_0}q^3\}-1,
\]
then
\begin{equation}
 \frac1{Z_{2,g}}
 \int_{D_R}
 W_TQ_{2,g}|e^{R_3(gb)}-1|\,d\gamma
 \leq
 \delta_{\rm Tay}\,\nu_{2,g}(W_T).
 \label{eq:v84-weighted-Taylor}
\end{equation}
Suppose in addition that the shifted Gaussian has covariance
\(\Sigma\leq C_0I\), its mean obeys \(\|m\|\leq R/2\), and
\begin{equation}
 \nu_{2,g}(W_T^2)\leq K_W.
 \label{eq:v84-fourth-stock}
\end{equation}
Then
\begin{equation}
 \frac1{Z_{2,g}}
 \int_{D_R^c}W_TQ_{2,g}\,d\gamma
 \leq
 K_W^{1/2}2^{d/4}
 \exp\left\{-\frac{R^2}{32C_0}\right\}.
 \label{eq:v84-weighted-exit}
\end{equation}
\end{proposition}

\begin{proof}
The uniform Taylor bound used in
\eqref{eq:v81-Taylor-bound} gives
\eqref{eq:v84-weighted-Taylor}.  For the exit term, Cauchy--Schwarz gives
\[
 \nu_{2,g}(W_T{\bf1}_{D_R^c})
 \leq
 \nu_{2,g}(W_T^2)^{1/2}
 \nu_{2,g}(D_R^c)^{1/2}.
\]
Apply Lemma \ref{lem:v81-shifted-exit} and take the square root.
\end{proof}

For fixed \(d\), bounded \(C_0\), bounded mean, and bounded
\(\|T\|_{\rm op}\), the fourth stock
\eqref{eq:v84-fourth-stock} is finite uniformly.  This follows from the
fourth Gaussian moment.  With the residual choice
\(R_\varepsilon=\varepsilon^{-\alpha/2}\), the weighted Taylor activity is
\(O(\varepsilon^{3(1-\alpha)/2})\), and the weighted exit remains
\(\exp\{-c\varepsilon^{-\alpha}\}\).  Thus both are summable for a
geometric penalty schedule.  The weighted physical exterior remains an
independent owner.

\subsection{Weighted cylinder repair}

For a local state space \(X_{k,A}\), let \(W_{k,A}\geq1\) be a measurable
weight.  Define
\[
 \|\sigma\|_{k,A;W}
 :=
 \int_{X_{k,A}}W_{k,A}\,
 d\left|(r_{k,A})_*\sigma\right|.
\]

\begin{theorem}[Weighted local repair]
\label{thm:v84-weighted-repair}
In the setting of Theorem \ref{thm:v83-cylinder-repair}, assume
\begin{equation}
 \left\|
 (p_{n+1,k})_*\mu_{n+1}
 -
 (p_{n,k})_*\mu_n
 \right\|_{k,A;W}
 \leq e_n^W(k,A),
 \qquad
 \sum_{n\geq k}e_n^W(k,A)<\infty.
 \label{eq:v84-weighted-adjacent}
\end{equation}
Then the descended local marginals converge in weighted variation:
\begin{equation}
 \|\nu_{N;k,A}-\nu_{n;k,A}\|_{W_{k,A}}
 \leq\sum_{j=n}^{N-1}e_j^W(k,A).
 \label{eq:v84-weighted-Cauchy}
\end{equation}
Their limit \(\bar\mu_{k,A}\) obeys
\begin{equation}
 \bar\mu_{k,A}(W_{k,A})
 \leq
 \nu_{n;k,A}(W_{k,A})
 +\sum_{j=n}^{\infty}e_j^W(k,A).
 \label{eq:v84-limit-moment}
\end{equation}
\end{theorem}

\begin{proof}
The proof of \eqref{eq:v84-weighted-Cauchy} is the telescoping argument in
V83 in the Banach space of signed measures with finite \(W_{k,A}\)-norm.
Apply Lemma \ref{lem:v84-moment-transfer} between
\(\nu_{n;k,A}\) and the limit to obtain
\eqref{eq:v84-limit-moment}.
\end{proof}

This theorem also passes unbounded local Ward or stress-energy insertions
when their absolute value is dominated by the chosen weight.  The weight
must be assigned before the forest estimate; an ordinary activity bound
cannot be promoted to a weighted one after the fact.

\subsection{Summing ultraviolet shell moments}

Let \(P_j\) be ultraviolet shell projections and choose
\(s_0>(D-2)/2\).  Suppose the reconstructed cutoff field
\(\phi^{(N)}\) obeys a uniform Littlewood--Paley comparison
\begin{equation}
 \|P_{\rm UV}\phi^{(N)}\|_{H^{-s_0}}^2
 \leq
 C_{\rm LP}\sum_{j=j_0}^N
 \|P_j\phi^{(N)}\|_{H^{-s_0}}^2.
 \label{eq:v84-LP-comparison}
\end{equation}

\begin{theorem}[Ultraviolet moment transfer]
\label{thm:v84-UV-transfer}
For each \(j\), let \(\nu_j\) be a reference shell law and let
\(\mu_j^{(N)}\) be the \(j\)-shell marginal of the interacting cutoff law,
\(N\geq j\).  Put
\[
 W_j=1+\|P_j\phi\|_{H^{-s_0}}^2.
\]
Assume
\begin{align}
 \sup_{N\geq j}\|\mu_j^{(N)}-\nu_j\|_{W_j}
 &\leq d_j,
 \label{eq:v84-shell-weighted-error}\\
 \sum_{j\geq j_0}
 \left[
  \nu_j\!\left(\|P_j\phi\|_{H^{-s_0}}^2\right)+d_j
 \right]
 &<\infty.
 \label{eq:v84-shell-sum}
\end{align}
Then
\begin{equation}
 \sup_N
 {\mathbb E}_{\mu^{(N)}}
 \|P_{\rm UV}\phi^{(N)}\|_{H^{-s_0}}^2
 <\infty.
 \label{eq:v84-UV-uniform-moment}
\end{equation}
\end{theorem}

\begin{proof}
Lemma \ref{lem:v84-moment-transfer} gives
\[
 {\mathbb E}_{\mu^{(N)}}
 \|P_j\phi^{(N)}\|_{H^{-s_0}}^2
 \leq
 \nu_j\!\left(\|P_j\phi\|_{H^{-s_0}}^2\right)+d_j.
\]
Sum over \(j\leq N\), apply
\eqref{eq:v84-LP-comparison}, and use
\eqref{eq:v84-shell-sum}.
\end{proof}

Under the order-minus-two hypothesis of Corollary
\ref{cor:v82-order-two}, the reference part of
\eqref{eq:v84-shell-sum} is finite for \(s_0>(D-2)/2\).  The remaining
task is to prove that the weighted polymer defects \(d_j\), including the
physical exterior, are summable.

\subsection{Sectorwise tightness}

Write the reconstructed field schematically as
\begin{equation}
 \phi^{(N)}
 =
 \phi_{\rm UV}^{(N)}
 +\phi_{\rm ret}^{(N)}
 +z^{(N)},
 \label{eq:v84-sector-split}
\end{equation}
where \(\phi_{\rm UV}^{(N)}\) contains the eliminated high-frequency
shells, \(\phi_{\rm ret}^{(N)}\) contains retained massless and infrared
fields, and \(z^{(N)}\) contains finite-dimensional zero, boundary, and
charge modes.

\begin{theorem}[Direct-sum tightness card]
\label{thm:v84-sector-tightness}
Let \(s>s_0\).  Suppose:
\begin{enumerate}
\item the ultraviolet family has the moment bound
      \eqref{eq:v84-UV-uniform-moment};
\item \(\{\phi_{\rm ret}^{(N)}\}\) is tight in a Polish space continuously
      embedded in \(H^{-s}\);
\item \(z^{(N)}\) is tight in its finite-dimensional physical quotient;
      and
\item the addition map in \eqref{eq:v84-sector-split} is continuous into
      \(H^{-s}\).
\end{enumerate}
Then \(\{\phi^{(N)}\}\) is tight in \(H^{-s}\).
\end{theorem}

\begin{proof}
The embedding \(H^{-s_0}\hookrightarrow H^{-s}\) is compact on a compact
base manifold.  Markov's inequality and
\eqref{eq:v84-UV-uniform-moment} put the ultraviolet field, uniformly in
\(N\), in a compact subset of \(H^{-s}\) up to arbitrarily small
probability.  By the other two tightness hypotheses, choose compact sets
for the retained and finite-dimensional sectors with equally small
exceptional probabilities.  The continuous image of their product under
addition is compact.  The union bound proves tightness of the sum.
\end{proof}

No independence between the three sectors is required.  On a noncompact
base manifold the compact Sobolev embedding must be replaced by local
compactness plus an explicit control of spatial escape.

\subsection{A weighted owner card}

\begin{definition}[V84 moment-transfer card]
\label{def:v84-moment-card}
A shell family passes the weighted card if:
\begin{enumerate}
\item its two-jet Gaussian has uniform second and fourth weighted moments;
\item the analytic Taylor and chart-exit defects satisfy Proposition
      \ref{prop:v84-weighted-local-defects};
\item the rooted forest is summable with one \(W_j\)-marked root;
\item the weighted physical exterior and algebraic reblocking errors form
      a summable sequence \(d_j\);
\item the reference covariance passes the order-minus-two count; and
\item retained infrared, zero, boundary, and charge modes pass their
      separate tightness rows.
\end{enumerate}
\end{definition}

\begin{firewall}[The weighted exterior is still missing]
\label{fire:v84-weighted-exterior}
The Gaussian Taylor and chart-exit rows now pass with a quadratic
Sobolev weight.  No estimate in V84 bounds
\[
 \int_{\text{physical chart exterior}}
 \bigl(1+\|\phi\|_{H^{-s_0}}^2\bigr)e^{-S_{\rm phys}}\,D\phi.
\]
That estimate requires a lower bound for the real part of the
gauge-fixed gravitational action on the chosen contour and compatible
control of the matter interaction.  It cannot be inferred from the local
Morse chart.
\end{firewall}

\begin{assessment}[Exact status after V84]
\label{ass:v84-status}
V84 proves weighted normalisation, weighted local Gaussian defect bounds,
weighted cylinder repair, ultraviolet moment transfer, and a direct-sum
tightness theorem.  Conditional on a summable weighted exterior and the
retained-sector rows, the order-minus-two reference estimate yields
Radon tightness in \(H^{-s}\) for \(s>s_0>(D-2)/2\).  Those conditions,
chiral determinant-line gluing, symmetry restoration, and the full
SM--Einstein continuum remain unproved.
\end{assessment}

\subsection{Next acquisition}

V85 is the next compulsory companion audit.  It will inspect the newest
Yang--Mills and Navier--Stokes notes for a weighted large-field estimate,
a retained-sector compactness mechanism, or a correction to the
cylinder-local route.  V86 will then attack the first transferable
bottleneck identified by that audit.

\paragraph{Parent-version record.}
V84 was formed from
\texttt{Renewal\_SM\_Einstein\_Continuum\_Research\_Notes\_V83.tex}.
The V83 parent had \(29\,410\) counted source lines, \(1\,074\,874\) bytes,
and SHA--256
\[
 \texttt{E513221835CDA55068B6028742E5F89434418686729F472B12E1F9DDD374C4E7}.
\]

% =============================================================================
% END APPENDED TENTH-SERIES V84 MATERIAL
% =============================================================================
% =============================================================================
% BEGIN APPENDED TENTH-SERIES V85 MATERIAL
% =============================================================================

\section{Tenth-series V85: companion audit and conditional-kernel route}
\label{sec:v10v85}

\subsection{Compulsory companion audit}

The newest active companions at this gate were inspected read-only:
\begin{align*}
 &\texttt{Renewal\_Yang\_Mills\_Mass\_Gap\_Research\_Notes\_V7.tex},
 \\
 &\qquad
 \operatorname{SHA256}
 =
 \texttt{09B911B8C7EBCD784059FEDB8743150C38459CEBB654DC3A42339294106206E1},
 \\
 &\texttt{Renewal\_NS\_Stability\_Research\_Notes\_V26.tex},
 \\
 &\qquad
 \operatorname{SHA256}
 =
 \texttt{82C887F8C2C69E4F28FAD7C3DC8DE5B9099CB5A4BF431EBA1FFF3E91935978AD}.
\end{align*}
The Yang--Mills file has advanced from V6 to V7 since the V80 audit.
The Navier--Stokes file is unchanged.

\subsection{Transferable Yang--Mills architecture}

The V7 companion normalizes each exact local conditional measure before
forming products.  In abstract notation, for a block \(y\),
\[
 P_y=G_y+\Delta_y,
 \qquad
 \Delta_y{\bf1}=0,
 \qquad
 \|\Delta_y\|_{\rm var}\leq q_y,
\]
where \(P_y\) and \(G_y\) are probability kernels on the same local state
space.  When blocks of one colour have disjoint interaction
neighbourhoods, conditioning on the exterior gives exact factorization:
\begin{equation}
 P_C=\bigotimes_{y\in C}P_y
 =
 \sum_{S\subset C}
 \left(\bigotimes_{y\in S}\Delta_y\right)
 \otimes
 \left(\bigotimes_{y\in C\setminus S}G_y\right).
 \label{eq:v85-colour-product}
\end{equation}
Each block occurs exactly once.  Scalar two-jet determinants were already
used to normalize \(G_y\); they do not reappear in
\eqref{eq:v85-colour-product}.

This architecture fits V81--V84.  The unnormalised Taylor, chart-exit,
physical-exterior, and algebraic defects first pass through the
normalisation lemmas.  Only the resulting zero-mass kernel
\(\Delta_y\) is inserted in a colour product.  Remote factors are
probability kernels of operator norm one, so a local observable sees only
defects in its backward dependence neighbourhood.  V86 will prove this
volume-free local statement and the corresponding invariant-kernel
resolvent estimate.

The V7 companion also makes a useful ownership distinction:
\begin{enumerate}
\item a regular local collar has one moving stationary centre and a
      Gaussian defect kernel;
\item failure of collar regularity has an explicit local witness and is
      assigned to a saturated bad component; and
\item a competing centre outside the regular chart must pay a separate
      action barrier.
\end{enumerate}
For the SM--Einstein programme this is a template, not an acquired
estimate.  The residual chart gives the local moving centre on its
implicit radius, but noncompact gravitational fibres do not provide the
compact-complement minimum used in the Yang--Mills barrier argument.

\subsection{Nontransferable hypotheses}

Three Yang--Mills inputs remain model-specific:
\begin{enumerate}
\item compactness of the reduced \(SU(3)\) fibre;
\item the positive Wilson plaquette identity that turns collar curvature
      into a Peierls action; and
\item the unproved all-colour schedule condition that prevents later
      updates from creating a second owner.
\end{enumerate}
None follows from V81's Gaussian chart-exit estimate.  In gravity, the
analogue of the first two rows requires a real-part lower bound for the
gauge-fixed action on the chosen conformal contour, including the matter
potential and retained boundary data.

The Navier--Stokes V26 rank-one Oseen calculation supplies no conditional
probability kernel, weighted large-field estimate, or compactness result
for this programme.  There is no transfer from it at V85.

\subsection{Route selected after the audit}

The most useful next upstream theorem is abstract.  It should not assume a
Yang--Mills collar or a gravitational coercivity result.  It will state
precisely what those model-specific inputs must buy:
\begin{enumerate}
\item a normalized local kernel difference \(P_y-G_y\) of small variation;
\item a volume-free bound on a one-colour product tested by a local
      observable;
\item a finite propagation rule through a complete colour sweep;
\item a strict contraction of the reference sweep on the physical
      quotient, with retained slow modes treated as parameters; and
\item a resolvent comparison of the exact and reference invariant
      conditional laws.
\end{enumerate}
Once proved abstractly, the theorem converts the V84 weighted kernel
activities into the cylinder error required by V83.  It will also expose
whether the next genuinely gravitational task is the exterior kernel
bound, the reference-sweep contraction, or the identification of the
coarse pushforward with the exact invariant conditional law.

\begin{definition}[V85 conditional-kernel card]
\label{def:v85-kernel-card}
For one regulator step and fixed retained data, the card consists of:
\begin{enumerate}
\item probability kernels \(P_y,G_y\) on the same local fibre;
\item zero-mass defects \(\Delta_y=P_y-G_y\) with ordinary and weighted
      variation bounds;
\item a fixed finite colouring of the dependence graph;
\item exact factorization within each frozen colour;
\item an ownership schedule for irregular or saturated components;
\item a reference-sweep contraction on the reduced physical sector; and
\item an exact identity relating the invariant exact sweep to the desired
      fine-to-coarse conditional pushforward.
\end{enumerate}
\end{definition}

\begin{firewall}[No compact collar import]
\label{fire:v85-no-compact-import}
The Yang--Mills fixed barrier outside a local stationary chart uses
compactness of a reduced group fibre and positivity of the Wilson action.
V85 does not import that barrier into the metric, conformal, Higgs, or
gauge-potential directions.  The gravitational physical exterior and its
quadratic Sobolev weight remain unbounded owners.
\end{firewall}

\begin{assessment}[Status after the V85 audit]
\label{ass:v85-status}
The audit imports the normalized conditional-kernel ownership pattern and
the exact one-colour product identity.  It identifies the all-colour
schedule and reference contraction as separate requirements and rejects
the compact Yang--Mills barrier as evidence for gravitational coercivity.
V86 will prove the abstract local-product, sweep-perturbation, and
invariant-resolvent statements.  The full SM--Einstein continuum remains
open.
\end{assessment}

\paragraph{Parent-version record.}
V85 was formed from
\texttt{Renewal\_SM\_Einstein\_Continuum\_Research\_Notes\_V84.tex}.
The V84 parent had \(29\,783\) counted source lines, \(1\,086\,735\) bytes,
and SHA--256
\[
 \texttt{BCEC22EACDD42601175CCE30E3D78DFDFEEA2D8DD7FBB4D86FA76D1E84377F8E}.
\]

% =============================================================================
% END APPENDED TENTH-SERIES V85 MATERIAL
% =============================================================================
% =============================================================================
% BEGIN APPENDED TENTH-SERIES V86 MATERIAL
% =============================================================================

\section{Tenth-series V86: local colour sweeps and invariant-kernel resolvents}
\label{sec:v10v86}

V85 reduced the next step to an operator theorem.  This note proves the
volume-free action of a normalized colour product on local observables,
propagates it through a finite sweep, and compares invariant laws through
the resolvent of a contractive reference sweep.  The theorem is abstract:
gravity still has to supply the small conditional kernels and the
contraction.

\subsection{One colour and a local observable}

Let \(\Lambda\) be a finite block set and
\(X=\prod_{y\in\Lambda}X_y\).  For a colour \(C\subset\Lambda\), assume
the update neighbourhoods are disjoint and the simultaneous exact and
reference kernels factorize, conditional on the frozen exterior:
\[
 P_C=\bigotimes_{y\in C}P_y,
 \qquad
 G_C=\bigotimes_{y\in C}G_y.
\]
Put
\[
 \Delta_y=P_y-G_y,
 \qquad
 \Delta_y{\bf1}=0,
 \qquad
 \|\Delta_y\|_{\infty\to\infty}\leq q_y.
\]
The norm bound is uniform in the frozen exterior.  Assume also that
updates at distinct sites of \(C\) neither alter nor use each other's
coordinates.

\begin{theorem}[Local colour-product bound]
\label{thm:v86-local-colour}
If \(f\) depends only on the coordinate set \(A\), then
\begin{equation}
 \|(P_C-G_C)f\|_\infty
 \leq
 \|f\|_\infty
 \left[
  \prod_{y\in A\cap C}(1+q_y)-1
 \right].
 \label{eq:v86-local-colour}
\end{equation}
In particular, if \(q_y\leq q\), then
\begin{equation}
 \|(P_C-G_C)f\|_\infty
 \leq
 \|f\|_\infty\{e^{q|A\cap C|}-1\}.
 \label{eq:v86-local-colour-exp}
\end{equation}
\end{theorem}

\begin{proof}
Expand the exact identity
\[
 P_C-G_C
 =
 \sum_{\varnothing\neq S\subset C}
 \left(\bigotimes_{y\in S}\Delta_y\right)
 \otimes
 \left(\bigotimes_{y\in C\setminus S}G_y\right).
\]
If \(S\) contains \(y\notin A\), the factor \(\Delta_y\) acts on a
function independent of \(x_y\) and gives zero.  The disjointness
hypothesis ensures that the other colour factors do not create
\(x_y\)-dependence.  Hence only nonempty \(S\subset A\cap C\) remain.
Every reference factor has sup-norm operator bound one, so the surviving
term has norm at most
\(\|f\|_\infty\prod_{y\in S}q_y\).  Summation gives
\eqref{eq:v86-local-colour}; \(1+q\leq e^q\) gives
\eqref{eq:v86-local-colour-exp}.
\end{proof}

The bound depends on the observation set, not on the total volume.  The
zero-mass identity \(\Delta_y{\bf1}=0\), obtained by normalizing before
multiplication, is essential.

\subsection{A complete finite colour sweep}

Let
\[
 P=P_1P_2\cdots P_\chi,
 \qquad
 G=G_1G_2\cdots G_\chi
\]
be exact and reference colour sweeps.  Products act on observables from
right to left.  Suppose each layer \(P_j\) and \(G_j\) sends a function
supported in \(B\) to one supported in a fixed backward neighbourhood
\({\cal N}_j(B)\).  For an \(A\)-supported \(f\), define recursively
\[
 B_\chi=A,
 \qquad
 B_j={\cal N}_{j+1}(B_{j+1})
 \quad(1\leq j<\chi).
\]
Thus \(B_j\) contains the support seen by the \(j\)-th layer after the
reference layers to its right have acted.

\begin{proposition}[Local sweep telescoping]
\label{prop:v86-sweep-telescope}
If the hypotheses of Theorem \ref{thm:v86-local-colour} hold for each
layer, then
\begin{equation}
 \|(P-G)f\|_\infty
 \leq
 \|f\|_\infty
 \sum_{j=1}^{\chi}
 \left[
  \prod_{y\in B_j\cap C_j}(1+q_{j,y})-1
 \right].
 \label{eq:v86-sweep-bound}
\end{equation}
If \(|B_j|\leq L_*|A|\) and \(q_{j,y}\leq q\), then
\begin{equation}
 \|(P-G)f\|_\infty
 \leq
 \chi\|f\|_\infty
 \{e^{L_*|A|q}-1\}.
 \label{eq:v86-uniform-sweep}
\end{equation}
For \(L_*|A|q\leq1\), the right side is at most
\(2\chi L_*|A|q\|f\|_\infty\).
\end{proposition}

\begin{proof}
Use the exact telescoping identity
\[
 P-G
 =
 \sum_{j=1}^{\chi}
 P_1\cdots P_{j-1}(P_j-G_j)G_{j+1}\cdots G_\chi.
\]
The left exact layers are Markov contractions in sup norm.  The right
reference layers produce support contained in \(B_j\).  Apply Theorem
\ref{thm:v86-local-colour} to the middle factor and sum.  The last two
bounds follow from \(1+q\leq e^q\) and \(e^x-1\leq2x\) for
\(0\leq x\leq1\).
\end{proof}

A fixed finite colouring and fixed-range neighbourhoods make \(\chi\) and
\(L_*\) volume-independent.  Therefore a summable local kernel activity
\(q_n\) produces the V83 cylinder error \(O(|A|q_n)\).

\subsection{The invariant-law resolvent}

Let \({\cal M}_0\) be a Banach space of finite signed measures of total
mass zero, with norm \(\|\cdot\|_*\).  Markov kernels act on the right.
The norm may be a weighted influence, Wasserstein, or local oscillation
dual norm; ordinary global total variation is not required.

\begin{theorem}[Invariant-kernel comparison]
\label{thm:v86-invariant-comparison}
Let \(G\) and \(P\) be Markov kernels with invariant probabilities
\(\pi_G\) and \(\pi_P\).  Assume
\begin{equation}
 \|\sigma G\|_*\leq\theta\|\sigma\|_*
 \quad
 (\sigma\in{\cal M}_0),
 \qquad
 0\leq\theta<1,
 \label{eq:v86-reference-contraction}
\end{equation}
and
\begin{equation}
 \|\pi_P(P-G)\|_*\leq\delta.
 \label{eq:v86-invariant-defect}
\end{equation}
Then
\begin{equation}
 \|\pi_P-\pi_G\|_*
 \leq\frac{\delta}{1-\theta}.
 \label{eq:v86-invariant-resolvent}
\end{equation}
More precisely,
\begin{equation}
 \pi_P-\pi_G
 =
 \sum_{r=0}^{\infty}\pi_P(P-G)G^r
 \label{eq:v86-resolvent-series}
\end{equation}
with convergence in \({\cal M}_0\).
\end{theorem}

\begin{proof}
Invariance gives the exact identity
\begin{align*}
 \pi_P-\pi_G
 &=\pi_PP-\pi_GG\\
 &=\pi_P(P-G)+(\pi_P-\pi_G)G.
\end{align*}
Iterating \(N\) times yields
\[
 \pi_P-\pi_G
 =
 \sum_{r=0}^{N-1}\pi_P(P-G)G^r
 +(\pi_P-\pi_G)G^N.
\]
The first term in each summand has total mass zero.  By
\eqref{eq:v86-reference-contraction}, its norm is at most
\(\delta\theta^r\), while the remainder tends to zero.  Summing the
geometric series proves both claims.
\end{proof}

A statewise kernel bound implies
\eqref{eq:v86-invariant-defect} whenever the norm is convex:
\begin{equation}
 \sup_x\|P(x,\cdot)-G(x,\cdot)\|_*\leq\delta
 \quad\Longrightarrow\quad
 \|\pi_P(P-G)\|_*\leq\delta.
 \label{eq:v86-statewise-defect}
\end{equation}

\begin{corollary}[Stability and uniqueness of the exact sweep]
\label{cor:v86-robust-contraction}
Suppose, in addition,
\begin{equation}
 \|\sigma(P-G)\|_*\leq\delta_*\|\sigma\|_*
 \qquad(\sigma\in{\cal M}_0)
 \label{eq:v86-operator-defect}
\end{equation}
and \(\theta+\delta_*<1\).  Then
\[
 \|\sigma P\|_*
 \leq(\theta+\delta_*)\|\sigma\|_*.
\]
Consequently \(P\) has at most one invariant probability whose difference
from another admissible probability belongs to \({\cal M}_0\).
\end{corollary}

\begin{proof}
Use the triangle inequality with
\(P=G+(P-G)\).  If \(\pi\) and \(\pi'\) are invariant, their difference
satisfies
\[
 \|\pi-\pi'\|_*
 =
 \|(\pi-\pi')P^N\|_*
 \leq(\theta+\delta_*)^N\|\pi-\pi'\|_*,
\]
which forces equality to zero.
\end{proof}

\subsection{Retained variables as parameters}

Let \(r\) denote retained massless, boundary, or zero-mode data.  It must
not be placed in a norm in which the physical sweep is falsely claimed
to contract.

\begin{corollary}[Lipschitz moving invariant family]
\label{cor:v86-moving-invariant}
Let \(G_r\) have invariant probability \(\pi_r\) for each retained
parameter \(r\).  Suppose the contraction
\eqref{eq:v86-reference-contraction} holds uniformly with the same
\(\theta<1\), and
\begin{equation}
 \sup_x
 \|G_r(x,\cdot)-G_{r'}(x,\cdot)\|_*
 \leq L_r d(r,r').
 \label{eq:v86-parameter-kernel}
\end{equation}
Then
\begin{equation}
 \|\pi_r-\pi_{r'}\|_*
 \leq\frac{L_r}{1-\theta}d(r,r').
 \label{eq:v86-parameter-invariant}
\end{equation}
\end{corollary}

\begin{proof}
Apply Theorem \ref{thm:v86-invariant-comparison} with
\(P=G_r\) and \(G=G_{r'}\), and use
\eqref{eq:v86-statewise-defect}.
\end{proof}

This is the appropriate regularity statement for a moving reference:
fast physical fibres contract conditionally, while retained data
parameterize their invariant laws.

\subsection{Connection to a coarse pushforward}

\begin{proposition}[Conditional sweep-to-coarse comparison]
\label{prop:v86-sweep-to-coarse}
Fix retained data \(r\).  Suppose:
\begin{enumerate}
\item the exact conditional fine law \(\rho_r\) is invariant under the
      complete exact sweep \(P_r\);
\item the reference conditional law \(\widehat\rho_r\) is invariant under
      \(G_r\);
\item \(G_r\) satisfies
      \eqref{eq:v86-reference-contraction};
\item the complete sweep defect satisfies
      \eqref{eq:v86-invariant-defect} with \(\delta_r\); and
\item the coarse map \(Q_r\) contracts the chosen signed-measure norm by
      a factor at most \(L_Q\).
\end{enumerate}
Then
\begin{equation}
 \|(Q_r)_*\rho_r-(Q_r)_*\widehat\rho_r\|_*
 \leq
 \frac{L_Q\delta_r}{1-\theta}.
 \label{eq:v86-coarse-comparison}
\end{equation}
\end{proposition}

\begin{proof}
Theorem \ref{thm:v86-invariant-comparison} compares the two conditional
invariant laws.  Apply the coarse pushforward bound.
\end{proof}

If \(\delta_r\) is the local V84 weighted kernel activity and is summable
over regulator scales, \eqref{eq:v86-coarse-comparison} supplies the
cylinder comparison input for V83.  This conclusion requires the exact
invariance identities in items one and two; similarity of one sweep is
not enough.

\subsection{Weighted drift requirement}

The V84 quadratic Sobolev weight is compatible with the preceding
argument if every colour kernel has a uniform Lyapunov bound
\begin{equation}
 K W\leq\lambda W+b,
 \qquad
 \lambda<1,
 \label{eq:v86-Lyapunov-drift}
\end{equation}
and the signed defects are small in the induced weighted operator norm.
Then invariant laws have \(W\)-moment at most \(b/(1-\lambda)\), and the
same telescoping and resolvent proofs apply in the corresponding
zero-mass weighted space.  Indeed, invariance and
\eqref{eq:v86-Lyapunov-drift} give
\[
 \pi(W)\leq\lambda\pi(W)+b.
\]
V86 records this as the precise weighted sweep target; it does not assert
the drift for the gravitational physical exterior.

\begin{firewall}[Schedule and contraction are model inputs]
\label{fire:v86-schedule-contraction}
Proposition \ref{prop:v86-sweep-telescope} assumes that each layer has a
uniform conditional kernel bound for every exterior state to which the
schedule applies.  It does not prove that a moving regular/bad
classification is preserved through later colours.  Theorem
\ref{thm:v86-invariant-comparison} assumes a strict contraction on the
reduced physical sector.  Neither property follows from the one-block
Gaussian chart.  Retained massless and boundary variables remain
parameters and require their own tightness estimates.
\end{firewall}

\begin{assessment}[Exact status after V86]
\label{ass:v86-status}
V86 proves a volume-free local colour-product estimate, a finite-sweep
telescoping bound, the invariant-kernel resolvent, robust contraction,
Lipschitz dependence on retained data, and the conditional
sweep-to-coarse comparison.  These theorems convert normalized local
defect kernels into V83 cylinder errors once a uniform all-colour schedule
and a reference contraction are supplied.  The gravitational exterior
kernel, that schedule, the physical contraction, retained-sector
tightness, determinant-line gluing, symmetry restoration, and the full
SM--Einstein continuum remain open.
\end{assessment}

\subsection{Next acquisition}

V87 will construct a checkable Dobrushin influence matrix for the
reference Gaussian conditional sweep.  It will derive the contraction
constant from block precision matrices by a Schur-complement estimate and
will identify exactly which massless and gauge directions must be removed
from the contracted sector.

\paragraph{Parent-version record.}
V86 was formed from
\texttt{Renewal\_SM\_Einstein\_Continuum\_Research\_Notes\_V85.tex}.
The V85 parent had \(29\,949\) counted source lines, \(1\,093\,769\) bytes,
and SHA--256
\[
 \texttt{74C547EBE85552E06A15CE5851755947D6675F2DB898A5D0B817FBE016E72383}.
\]

% =============================================================================
% END APPENDED TENTH-SERIES V86 MATERIAL
% =============================================================================
% =============================================================================
% BEGIN APPENDED TENTH-SERIES V87 MATERIAL
% =============================================================================

\section{Tenth-series V87: Gaussian influence matrix and the massless obstruction}
\label{sec:v10v87}

V86 requires a strict contraction of the reference conditional sweep.
For a Gaussian reference this can be read directly from the block
precision.  The result below gives an explicit weighted Dobrushin
constant and shows why an unprojected massless lattice field cannot pass
the test uniformly.

\subsection{Block Gaussian conditionals}

Let
\[
 E=\prod_{i\in\Lambda}E_i
\]
be a finite product of Euclidean block spaces.  Let \(A=A^*>0\) have block
entries \(A_{ij}:E_j\to E_i\), and let
\[
 d\gamma_A(x)=Z_A^{-1}
 \exp\left\{-\frac12\langle x,Ax\rangle\right\}\,dx.
\]
The conditional law at block \(i\), with all other blocks fixed, is
\begin{equation}
 x_i\mid x_{\ne i}
 \sim
 {\cal N}\left(
  -A_{ii}^{-1}\sum_{j\ne i}A_{ij}x_j,\,
  A_{ii}^{-1}
 \right).
 \label{eq:v87-Gaussian-conditional}
\end{equation}
Choose norms \(\|\cdot\|_i\) and weights \(w_i>0\), and define
\begin{equation}
 c_{ij}:=
 \|A_{ii}^{-1}A_{ij}\|_{i\leftarrow j},
 \qquad
 \kappa_w:=
 \sup_i\frac1{w_i}\sum_{j\ne i}c_{ij}w_j.
 \label{eq:v87-influence}
\end{equation}
On \(E\), use the weighted maximum metric
\begin{equation}
 d_w(x,x'):=\max_i\frac{\|x_i-x_i'\|_i}{w_i}.
 \label{eq:v87-block-metric}
\end{equation}

Colour the precision graph so that \(A_{ij}=0\) for distinct blocks of
one colour.  Let \(G\) be one systematic heat-bath sweep through all
colours, updating every block exactly once.

\begin{theorem}[Precision-matrix Dobrushin contraction]
\label{thm:v87-Gaussian-Dobrushin}
If \(\kappa_w<1\), then for all probability measures \(\mu,\nu\) with a
finite first \(d_w\)-moment,
\begin{equation}
 W_{1,d_w}(\mu G,\nu G)
 \leq
 \kappa_w W_{1,d_w}(\mu,\nu).
 \label{eq:v87-Wasserstein-contraction}
\end{equation}
In particular, \(\gamma_A\) is the unique invariant probability with a
finite first \(d_w\)-moment, and the reference contraction constant in
V86 may be taken to be
\(\theta=\kappa_w\).
\end{theorem}

\begin{proof}
Start two sweeps at \(x,x'\) and use the same Gaussian noise in each
corresponding conditional update.  If block \(i\) is updated while the
current configurations are \(z,z'\), equation
\eqref{eq:v87-Gaussian-conditional} gives
\[
 \|z_i^{\rm new}-{z_i'}^{\rm new}\|_i
 \leq
 \sum_{j\ne i}c_{ij}\|z_j-z_j'\|_j
 \leq
 \kappa_w w_i\,d_w(z,z').
\]
The other block discrepancies are unchanged.  During a sweep the current
maximum discrepancy never exceeds its initial value, because an updated
coordinate becomes at most \(\kappa_w\) times that value.  After the
complete sweep every coordinate has been updated, hence
\[
 d_w(X_1,X_1')\leq\kappa_wd_w(x,x')
\]
under this coupling.  Integrate over an arbitrary initial coupling and
take the infimum to obtain
\eqref{eq:v87-Wasserstein-contraction}.  Every individual heat-bath
kernel preserves \(\gamma_A\), so the sweep does also.  Contraction gives
uniqueness.
\end{proof}

The theorem is uniform in volume whenever the weighted row bound
\(\kappa_w<1\) is uniform.  It uses the full conditional mean; no
independence between different colours is asserted beyond the zero
precision blocks required for simultaneous same-colour updates.

\subsection{A checkable diagonal-dominance certificate}

\begin{corollary}[Block diagonal-dominance card]
\label{cor:v87-diagonal-card}
Suppose
\[
 A_{ii}\geq\lambda_iI
\]
and
\begin{equation}
 \sum_{j\ne i}\|A_{ij}\|_{i\leftarrow j}w_j
 \leq\kappa\lambda_iw_i
 \qquad\text{for every }i
 \label{eq:v87-diagonal-dominance}
\end{equation}
with \(0\leq\kappa<1\).  Then
\(\kappa_w\leq\kappa\), and Theorem
\ref{thm:v87-Gaussian-Dobrushin} applies.
\end{corollary}

\begin{proof}
The spectral lower bound gives
\(\|A_{ii}^{-1}\|\leq\lambda_i^{-1}\).  Insert this in
\eqref{eq:v87-influence} and use
\eqref{eq:v87-diagonal-dominance}.
\end{proof}

The weights may absorb nonuniform block sizes or a controlled distance
from an observation set.  Finding them is a positive-vector problem for
the nonnegative influence matrix.

\subsection{Schur complements and loss of locality}

The physical fast precision may be a Schur complement.  Write
\[
 {\cal A}
 =
 \begin{pmatrix}
  D+R&B\\
  B^*&H
 \end{pmatrix},
 \qquad
 H>0,
\]
where \(D\) is block diagonal on the fast variables.  After eliminating
the auxiliary block, the fast precision is
\begin{equation}
 A_{\rm eff}=D+R-S,
 \qquad
 S:=BH^{-1}B^*.
 \label{eq:v87-effective-precision}
\end{equation}

\begin{proposition}[Schur influence certificate]
\label{prop:v87-Schur-influence}
Suppose \(D_{ii}\geq\lambda_iI\) and
\(\lambda_i>\|S_{ii}\|\).  If, for some \(0\leq\kappa<1\),
\begin{equation}
 \sum_{j\ne i}
 \bigl(\|R_{ij}\|+\|S_{ij}\|\bigr)w_j
 \leq
 \kappa\bigl(\lambda_i-\|S_{ii}\|\bigr)w_i,
 \label{eq:v87-Schur-row}
\end{equation}
then the Gaussian heat-bath sweep for \(A_{\rm eff}\) contracts in
\(W_{1,d_w}\) with factor at most \(\kappa\).
\end{proposition}

\begin{proof}
The diagonal block satisfies
\[
 (A_{\rm eff})_{ii}
 =
 D_{ii}+R_{ii}-S_{ii}.
\]
Absorb any declared self-adjoint \(R_{ii}\) into \(D_{ii}\); the assumed
lower bound is then
\((A_{\rm eff})_{ii}\geq
(\lambda_i-\|S_{ii}\|)I\).
For \(i\ne j\),
\[
 \|(A_{\rm eff})_{ij}\|
 \leq\|R_{ij}\|+\|S_{ij}\|.
\]
Equation \eqref{eq:v87-Schur-row} is therefore the diagonal-dominance
condition of Corollary \ref{cor:v87-diagonal-card}.
\end{proof}

This proposition exposes a locality debit.  The crude estimate
\[
 \|S_{ij}\|
 \leq
 \|B_i\|\,\|H^{-1}\|\,\|B_j\|
\]
does not give a summable row in a large volume.  A usable certificate
needs an exponentially or sufficiently fast algebraically decaying
kernel for \(H^{-1}\), or it must keep the corresponding slow auxiliary
directions in the retained sector.  This is the same nested-inverse
locality issue isolated in V63 and V74.

\subsection{Retained parameters}

Partition the Gaussian variables into fast fields \(x\) and retained
fields \(r\):
\[
 \frac12
 \begin{pmatrix}x\\r\end{pmatrix}^{\!*}
 \begin{pmatrix}A_{FF}&A_{FR}\\A_{RF}&A_{RR}\end{pmatrix}
 \begin{pmatrix}x\\r\end{pmatrix}.
\]
At fixed \(r\), the fast invariant conditional law is
\begin{equation}
 \pi_r
 =
 {\cal N}\left(-A_{FF}^{-1}A_{FR}r,\,
                A_{FF}^{-1}\right).
 \label{eq:v87-fast-invariant}
\end{equation}

\begin{proposition}[Exact retained-parameter Lipschitz bound]
\label{prop:v87-retained-Lipschitz}
For a norm whose translation coupling has cost given by \(d_w\),
\begin{equation}
 W_{1,d_w}(\pi_r,\pi_{r'})
 \leq
 \|A_{FF}^{-1}A_{FR}(r-r')\|_w.
 \label{eq:v87-retained-Lipschitz}
\end{equation}
If
\(\|A_{FF}^{-1}A_{FR}\|_{w\leftarrow R}\leq L_R\), then the right side is
at most \(L_R\|r-r'\|_R\).
\end{proposition}

\begin{proof}
Couple the two Gaussian laws with the same centred Gaussian sample.  Their
difference is the deterministic difference of the means in
\eqref{eq:v87-fast-invariant}.
\end{proof}

The retained variables are not contracted.  Proposition
\ref{prop:v87-retained-Lipschitz} supplies the parameter regularity needed
by Corollary \ref{cor:v86-moving-invariant}.

\subsection{The massless obstruction}

\begin{proposition}[Nearest-neighbour massless row]
\label{prop:v87-massless-row}
For the scalar nearest-neighbour precision
\[
 A_h=-\Delta_h+m^2
\]
on a \(D\)-dimensional periodic lattice of mesh \(h\), with one scalar per
site and unweighted blocks,
\begin{equation}
 \kappa_h
 =
 \frac{2D h^{-2}}{2D h^{-2}+m^2}
 =
 \frac{2D}{2D+m^2h^2}.
 \label{eq:v87-massive-kappa}
\end{equation}
Thus \(\kappa_h=1\) when \(m=0\), and for fixed physical \(m>0\),
\(\kappa_h\uparrow1\) as \(h\downarrow0\).
\end{proposition}

\begin{proof}
The diagonal precision is \(2Dh^{-2}+m^2\), while each of the \(2D\)
nearest-neighbour off-diagonal entries has magnitude \(h^{-2}\).  Insert
these entries in \eqref{eq:v87-influence}.
\end{proof}

Therefore single-site heat-bath Dobrushin contraction is not uniform in
the continuum limit for an unprojected massless field.  A gauge zero mode
makes the full precision singular before the row test is even defined.
The contraction required by V86 must instead use one of the following
proved mechanisms:
\begin{enumerate}
\item retain the infrared and gauge kernel and contract only a
      high-frequency shell;
\item use blocks whose physical diameter scales with the correlation
      length and prove a block influence gap;
\item use a multilevel reference sweep with a separate coarse correction;
      or
\item use a spectral or hypocoercive norm adapted to the projected shell.
\end{enumerate}

\begin{firewall}[Positive precision is not enough]
\label{fire:v87-positive-not-contractive}
A positive finite-volume Gaussian precision has an invariant probability,
but positivity alone does not give a volume- and cutoff-uniform
Dobrushin constant.  Proposition \ref{prop:v87-massless-row} rules out
that inference even for the free scalar lattice field.  The TT, photon,
gauge, and boundary kernels must be projected or retained before a strict
uniform contraction can be claimed.
\end{firewall}

\begin{definition}[V87 reference-contraction card]
\label{def:v87-contraction-card}
A Gaussian reference sweep passes the card if:
\begin{enumerate}
\item its fast conditional precision is positive after all declared
      gauge and contour choices;
\item the fast block influence matrix admits weights with
      \(\kappa_w<1\) uniformly in volume and cutoff;
\item every Schur inverse entering the row has a proved locality bound;
\item the same-colour precision blocks vanish or are handled by a
      nonfactorized layer estimate;
\item the retained-to-fast mean map is uniformly bounded; and
\item massless, gauge, zero, and boundary modes excluded from the
      contracted sector appear in the V84 tightness ledger.
\end{enumerate}
\end{definition}

\begin{assessment}[Exact status after V87]
\label{ass:v87-status}
V87 derives the Dobrushin contraction directly from the block precision,
gives a Schur-complement certificate, proves retained-parameter
Lipschitz dependence, and exhibits the exact massless single-site
obstruction.  It does not establish a uniform influence gap for the
SM--Einstein fast shell.  The physical exterior, projected-shell
contraction, weighted drift, retained-sector tightness,
determinant-line gluing, symmetry restoration, and the full
SM--Einstein continuum remain open.
\end{assessment}

\subsection{Next acquisition}

V88 will replace the failed single-site row by a projected-shell
contraction theorem.  It will use a spectral window and a block
preconditioner to determine a sufficient condition for a cutoff-uniform
fast-sector gap, while retaining all modes below the shell threshold.

\paragraph{Parent-version record.}
V87 was formed from
\texttt{Renewal\_SM\_Einstein\_Continuum\_Research\_Notes\_V86.tex}.
The V86 parent had \(30\,336\) counted source lines, \(1\,106\,067\) bytes,
and SHA--256
\[
 \texttt{D2BFA14B2F834ADEC427A74663DA23E5AC72753A93DD552F013FF1CE3EFC3BDC}.
\]

% =============================================================================
% END APPENDED TENTH-SERIES V87 MATERIAL
% =============================================================================
% =============================================================================
% BEGIN APPENDED TENTH-SERIES V88 MATERIAL
% =============================================================================

\section{Tenth-series V88: projected-shell Ornstein--Uhlenbeck contraction}
\label{sec:v10v88}

V87 rules out a cutoff-uniform single-site Dobrushin constant for an
unprojected massless field.  The correct fast variable is a spectral
shell whose precision is of the same order as the square of its shell
frequency.  This note constructs an exact Gaussian transition on that
shell and proves a uniform contraction.  Locality of the shell
coordinates remains a separate requirement.

\subsection{Fast-shell spectral window}

Let \(F_n\) be a finite-dimensional real fast-shell space, with
characteristic frequency \(\Lambda_n\).  Retained infrared, gauge-kernel,
boundary, and zero modes lie in a complementary space \(R_n\).  At fixed
retained datum \(r\), suppose the Gaussian fast conditional law is
\begin{equation}
 \gamma_{n,r}
 =
 {\cal N}(m_n(r),A_n(r)^{-1})
 \label{eq:v88-fast-Gaussian}
\end{equation}
and that its precision obeys the uniform shell window
\begin{equation}
 c_-\Lambda_n^2I
 \leq A_n(r)\leq
 c_+\Lambda_n^2I
 \qquad
 (0<c_-\leq c_+<\infty).
 \label{eq:v88-shell-window}
\end{equation}
All inequalities are on \(F_n\) and uniform in the admissible \(r\).

For a fixed transition time \(\tau>0\), define
\begin{align}
 E_{n,r}
 &:=
 \exp\{-\tau\Lambda_n^{-2}A_n(r)\},
 \label{eq:v88-OU-mean}\\
 Q_{n,r}
 &:=
 A_n(r)^{-1}\bigl(I-E_{n,r}^2\bigr).
 \label{eq:v88-OU-noise}
\end{align}
Because \(E_{n,r}\) is a function of \(A_n(r)\), the factors commute and
\(Q_{n,r}\geq0\).  Define the transition
\begin{equation}
 X'
 =
 m_n(r)+E_{n,r}\{X-m_n(r)\}+\xi,
 \qquad
 \xi\sim{\cal N}(0,Q_{n,r}).
 \label{eq:v88-OU-kernel}
\end{equation}
This is the time-\(\tau\) kernel of the preconditioned Ornstein--Uhlenbeck
equation with scalar shell preconditioner \(\Lambda_n^{-2}I\).

\begin{theorem}[Uniform projected-shell contraction]
\label{thm:v88-OU-contraction}
Under \eqref{eq:v88-shell-window}, the kernel \(G_{n,r}\) in
\eqref{eq:v88-OU-kernel} preserves \(\gamma_{n,r}\) and satisfies
\begin{equation}
 W_2(\mu G_{n,r},\nu G_{n,r})
 \leq
 \theta W_2(\mu,\nu),
 \qquad
 \theta:=e^{-\tau c_-}<1,
 \label{eq:v88-uniform-W2}
\end{equation}
for all probability measures on \(F_n\) with finite second moment.  The
constant is independent of \(n\), the volume, and \(r\).
\end{theorem}

\begin{proof}
If \(X-m_n(r)\) has covariance \(A_n(r)^{-1}\), then the new centred
covariance is
\[
 E_{n,r}A_n(r)^{-1}E_{n,r}+Q_{n,r}
 =
 A_n(r)^{-1}.
\]
The mean is unchanged, proving invariance.  Couple two transitions with
the same noise.  Their difference is
\[
 X'-Y'=E_{n,r}(X-Y).
\]
The spectral window gives
\[
 \|E_{n,r}\|_{\rm op}
 \leq e^{-\tau c_-}=\theta.
\]
Take the quadratic transportation cost, integrate over an arbitrary
coupling of \(\mu,\nu\), and minimize.
\end{proof}

Unlike the single-site constant in
\eqref{eq:v87-massive-kappa}, \(\theta\) does not approach one as the
cutoff is removed.  The factor \(\Lambda_n^{-2}\) exactly compensates the
order-two shell precision.

\subsection{Quadratic Lyapunov drift}

Let \(d_n=\dim F_n\), and put
\[
 W_{n,r}(x)
 =
 1+a_n\|x-m_n(r)\|^2,
 \qquad a_n>0.
\]

\begin{proposition}[Exact shell drift]
\label{prop:v88-shell-drift}
The OU kernel satisfies
\begin{equation}
 G_{n,r}W_{n,r}(x)
 \leq
 \theta^2W_{n,r}(x)+b_n,
 \label{eq:v88-shell-drift}
\end{equation}
where
\begin{equation}
 b_n
 :=
 1-\theta^2+\frac{a_nd_n}{c_-\Lambda_n^2}.
 \label{eq:v88-drift-constant}
\end{equation}
\end{proposition}

\begin{proof}
Using \(\|E_{n,r}\|\leq\theta\) and
\(0\leq Q_{n,r}\leq A_n(r)^{-1}\leq
(c_-\Lambda_n^2)^{-1}I\),
\begin{align*}
 {\mathbb E}\|X'-m_n(r)\|^2
 &=
 \|E_{n,r}\{x-m_n(r)\}\|^2+\operatorname{Tr}Q_{n,r}\\
 &\leq
 \theta^2\|x-m_n(r)\|^2+
 \frac{d_n}{c_-\Lambda_n^2}.
\end{align*}
Multiply by \(a_n\), add one, and rearrange.
\end{proof}

For a negative-Sobolev shell weight take
\(a_n\simeq\Lambda_n^{-2s_0}\).  If
\(d_n=O(\Lambda_n^D)\), the nonconstant term in \(b_n\) has order
\(\Lambda_n^{D-2-2s_0}\), exactly the V82 summability exponent.

\subsection{A block-preconditioned form}

The scalar shell preconditioner can be replaced by a symmetric positive
operator \(M_n\).  Set
\[
 B_n=M_n^{1/2}A_nM_n^{1/2}.
\]
In the coordinate \(z=M_n^{-1/2}(x-m_n)\), the OU drift is
\(-B_nz\).

\begin{corollary}[Preconditioned spectral certificate]
\label{cor:v88-preconditioned}
If
\begin{equation}
 c_MI\leq B_n\leq C_MI
 \label{eq:v88-preconditioned-window}
\end{equation}
uniformly, then the exact preconditioned OU kernel contracts in the
\(M_n^{-1}\)-metric with factor \(e^{-\tau c_M}\).
\end{corollary}

\begin{proof}
In the \(z\)-coordinate the synchronous-coupling difference is
\(e^{-\tau B_n}(z-z')\).  Its operator norm is at most
\(e^{-\tau c_M}\).  Transform back.
\end{proof}

For example, if
\[
 A_n=A_{0,n}^{1/2}(I+K_n)A_{0,n}^{1/2},
 \qquad
 \|K_n\|\leq\delta<1,
\]
then \(M_n=A_{0,n}^{-1}\) gives
\(c_M=1-\delta\) and \(C_M=1+\delta\).  This spectral statement does not
by itself make \(M_n\) local.

\subsection{Locality of the transition}

Assume temporarily that the normalized generator
\begin{equation}
 B_n:=\Lambda_n^{-2}A_n
 \label{eq:v88-normalized-generator}
\end{equation}
has block range at most \(R_0\) in a graph metric and absolute block-row
norm at most \(J\):
\begin{equation}
 \sup_i\sum_j\|(B_n)_{ij}\|\leq J.
 \label{eq:v88-generator-row}
\end{equation}
Both constants are independent of \(n\).

\begin{proposition}[Matrix-exponential locality]
\label{prop:v88-exponential-locality}
Let \(d(i,j)\) be graph distance and
\(r=\lceil d(i,j)/R_0\rceil\).  Then
\begin{equation}
 \|(E_{n,r_0})_{ij}\|
 \leq
 \sum_{k\geq r}\frac{(\tau J)^k}{k!},
 \label{eq:v88-exponential-tail}
\end{equation}
where \(r_0\) denotes fixed retained data, not the integer \(r\).
Moreover, for \(r\geq1\),
\begin{equation}
 \sum_{k\geq r}\frac{(\tau J)^k}{k!}
 \leq
 e^{\tau J}
 \left(\frac{e\tau J}{r}\right)^r.
 \label{eq:v88-Poisson-tail}
\end{equation}
\end{proposition}

\begin{proof}
The block \((B_n^k)_{ij}\) vanishes when \(kR_0<d(i,j)\).  Expanding
\(e^{-\tau B_n}\) and taking the absolute block-row majorant gives
\eqref{eq:v88-exponential-tail}.  For a Poisson random variable
\(N\) of mean \(\tau J\), the left side is
\(e^{\tau J}{\mathbb P}\{N\geq r\}\).  The exponential Markov bound
\({\mathbb P}\{N\geq r\}\leq(e\tau J/r)^r\) proves
\eqref{eq:v88-Poisson-tail}.
\end{proof}

Thus the dependence of the updated mean on distant initial blocks has a
superexponential fixed-time tail.  For unbounded fields, converting this
kernel tail to V83's sup-norm quasi-locality requires a weighted
Lipschitz norm and the V84 moment stock.  It is not a uniform
approximation on the whole noncompact field space.

\subsection{Sharp shells can destroy locality}

Let \(P_n\) be a projection onto a fast spectral band.  Even if the full
precision \(A\) is finite range, the compressed operator
\(P_nAP_n\), expressed back in site coordinates, can have a long-range
kernel because \(P_n\) itself is nonlocal.  Therefore
\eqref{eq:v88-generator-row} does not follow from locality of \(A\).

A usable fast-shell realization must provide one of:
\begin{enumerate}
\item local multiresolution detail coordinates in which the conditional
      precision has uniform range or exponential decay;
\item a smooth quasi-local band filter with a quantified tail;
\item a block Schur complement satisfying Proposition
      \ref{prop:v87-Schur-influence}; or
\item a multilevel transition whose coarse correction is retained
      explicitly.
\end{enumerate}

\begin{proposition}[Fourier model check]
\label{prop:v88-Fourier-check}
For a translation-invariant order-two scalar precision with symbol
\(a(k)\) satisfying
\[
 c_-|k|^2\leq a(k)\leq c_+|k|^2
\]
on the band \(c_0\Lambda_n\leq|k|\leq C_0\Lambda_n\), the shell window
\eqref{eq:v88-shell-window} holds after changing \(c_\pm\) by the fixed
band constants.  Hence the OU contraction is uniform on that band.
\end{proposition}

\begin{proof}
On the stated band, \(|k|^2\) is bounded above and below by fixed
multiples of \(\Lambda_n^2\).  The Fourier multiplier inequalities give
\eqref{eq:v88-shell-window}, and Theorem
\ref{thm:v88-OU-contraction} applies.
\end{proof}

This check explains the mechanism but does not establish the required
variable-coefficient, gauge-reduced, boundary-compatible shell estimate
for the full SM--Einstein precision.

\subsection{Bridge to the physical invariant law}

The reference OU kernel has the required contraction and invariant
two-jet Gaussian.  To invoke V86, one must construct a physical Markov
kernel \(P_{n,r}\) with:
\begin{enumerate}
\item the exact conditional interacting law as invariant probability;
\item the same retained parameters \(r\);
\item a small local or weighted difference from \(G_{n,r}\); and
\item a quasi-local dependence kernel compatible with V83.
\end{enumerate}
A Metropolis correction of the OU proposal is one possible exact bridge.
Its rejection probability must be estimated by the V81 Taylor,
chart-exit, and physical-exterior owners rather than by a global
supremum of the extensive interaction.

\begin{firewall}[A contractive reference is not yet a close reference]
\label{fire:v88-contractive-not-close}
Theorem \ref{thm:v88-OU-contraction} proves a uniform fast-shell
contraction, but V86 also requires a small exact-to-reference kernel
difference.  Replacing the local Gaussian heat-bath sweep by the OU
kernel does not make that difference small automatically.  The sharp
spectral projector can also violate locality.  Both the exact physical
bridge and a local shell realization remain to be constructed.
\end{firewall}

\begin{assessment}[Exact status after V88]
\label{ass:v88-status}
V88 constructs an exact invariant fast-shell OU kernel with
cutoff-independent Wasserstein contraction, proves its quadratic
Lyapunov drift and matrix-exponential locality under a finite-range
normalized generator, and isolates the nonlocal sharp-projector
obstruction.  The physical-kernel comparison, local multiresolution
shell, gravitational exterior, retained-sector tightness,
determinant-line gluing, symmetry restoration, and the full
SM--Einstein continuum remain open.
\end{assessment}

\subsection{Next acquisition}

V89 will analyze the exact Metropolis correction of the Gaussian shell
proposal.  It will prove a one-block rejection and kernel-difference
bound from local oscillation of the nonlinear remainder, and it will show
why an extensive global supremum cannot be used.  The physical exterior
will remain a separate rejection owner.

\paragraph{Parent-version record.}
V88 was formed from
\texttt{Renewal\_SM\_Einstein\_Continuum\_Research\_Notes\_V87.tex}.
The V87 parent had \(30\,681\) counted source lines, \(1\,117\,435\) bytes,
and SHA--256
\[
 \texttt{2EA56F3C6A620C7C51157C6BFDF175B2AB5942A8ADC9295B792DB9CBB2FE62CE}.
\]

% =============================================================================
% END APPENDED TENTH-SERIES V88 MATERIAL
% =============================================================================
% =============================================================================
% BEGIN APPENDED TENTH-SERIES V89 MATERIAL
% =============================================================================

\section{Tenth-series V89: exact Metropolis bridge and local rejection owners}
\label{sec:v10v89}

V88 constructs a contractive Gaussian shell proposal.  This note applies
the Metropolis correction to obtain an exact kernel for a perturbed
physical density.  The rejection probability is controlled on a local
chart by the oscillation of the nonlinear remainder and a proposal-exit
term.  A separate example proves that a global refresh against an
extensive interaction has asymptotically total rejection.

\subsection{Exact reversible correction}

Let \(\nu\) be a probability measure on a standard Borel space \(X\), and
let \(G\) be reversible with respect to \(\nu\):
\[
 \nu(dx)G(x,dy)=\nu(dy)G(y,dx).
\]
Let \(V:X\to{\mathbb R}\) be measurable with
\(0<Z_V:=\nu(e^{-V})<\infty\), and define
\begin{equation}
 \mu(dx)=Z_V^{-1}e^{-V(x)}\nu(dx).
 \label{eq:v89-target}
\end{equation}
Set
\begin{align}
 a(x,y)&:=1\wedge e^{V(x)-V(y)},
 \label{eq:v89-acceptance}\\
 r(x)&:=\int\{1-a(x,y)\}\,G(x,dy),
 \label{eq:v89-rejection}\\
 P(x,dy)&:=a(x,y)G(x,dy)+r(x)\delta_x(dy).
 \label{eq:v89-MH-kernel}
\end{align}

\begin{theorem}[Exact Metropolis bridge]
\label{thm:v89-exact-MH}
The kernel \(P\) is reversible with respect to \(\mu\), hence
\(\mu P=\mu\).  Moreover,
\begin{equation}
 P(x,\cdot)-G(x,\cdot)
 =
 -\{1-a(x,\cdot)\}G(x,\cdot)+r(x)\delta_x,
 \label{eq:v89-kernel-difference}
\end{equation}
and, with the V81 total-variation convention,
\begin{equation}
 \|P(x,\cdot)-G(x,\cdot)\|_{\rm TV}
 \leq2r(x).
 \label{eq:v89-TV-rejection}
\end{equation}
\end{theorem}

\begin{proof}
Multiplying the off-diagonal transition by the target density gives
\begin{align*}
 \mu(dx)G(x,dy)a(x,y)
 =
 Z_V^{-1}
 \min\{e^{-V(x)},e^{-V(y)}\}
 \nu(dx)G(x,dy),
\end{align*}
which is symmetric in \(x,y\).  The rejection term is diagonal and hence
symmetric.  This proves detailed balance and invariance.
Equation \eqref{eq:v89-kernel-difference} follows from the definition.
The negative off-diagonal measure has mass \(r(x)\), and the positive
diagonal measure has the same mass, so the triangle inequality gives
\eqref{eq:v89-TV-rejection}.
\end{proof}

\subsection{Buffered local chart}

Let \(D_{\rm in}\subset D_{\rm out}\subset X\).  Think of
\(D_{\rm in}\) as the running residual ball and \(D_{\rm out}\) as a
larger ball still contained in the proved analytic chart.

\begin{lemma}[Local rejection bound]
\label{lem:v89-local-rejection}
Suppose
\begin{equation}
 \operatorname{osc}_{D_{\rm out}}V
 :=
 \sup_{x,y\in D_{\rm out}}|V(x)-V(y)|
 \leq\omega
 \label{eq:v89-local-oscillation}
\end{equation}
and
\begin{equation}
 \sup_{x\in D_{\rm in}}G(x,D_{\rm out}^c)\leq\tau.
 \label{eq:v89-proposal-exit}
\end{equation}
Then, for every \(x\in D_{\rm in}\),
\begin{equation}
 r(x)\leq\omega+\tau,
 \qquad
 \|P(x,\cdot)-G(x,\cdot)\|_{\rm TV}
 \leq2(\omega+\tau).
 \label{eq:v89-local-kernel-bound}
\end{equation}
\end{lemma}

\begin{proof}
For \(t\geq0\), \(1-e^{-t}\leq t\).  Hence
\[
 1-a(x,y)
 \leq[V(y)-V(x)]_+
 \leq\omega
\]
when \(x,y\in D_{\rm out}\).  Split the rejection integral into
\(D_{\rm out}\) and its complement.  The first part is at most
\(\omega\), and the second is at most \(\tau\).  Apply
\eqref{eq:v89-TV-rejection}.
\end{proof}

For the V81 residual chart, take \(V=-R_3\).  On a ball of physical
radius \(qt_0\),
\[
 \omega\leq2{\cal M}_{t_0}q^3.
\]
With \(q=q_\varepsilon\) from V82, this is
\(O(\varepsilon^{3(1-\alpha)/2})\).  If the OU proposal has a buffered
Gaussian exit
\(\tau\leq C\exp\{-c\varepsilon^{-\alpha}\}\), the local kernel defect is
summable for a geometric penalty schedule.

The buffer matters.  A starting point arbitrarily close to the boundary
of \(D_{\rm out}\) need not have a small proposal-exit probability even
when the stationary Gaussian tail is tiny.

\subsection{Weighted rejection}

Let \(W\geq1\).  Define the weighted exit owner
\begin{equation}
 \tau_W(x)
 :=
 \int_{D_{\rm out}^c}\{W(y)+W(x)\}\,G(x,dy).
 \label{eq:v89-weighted-exit}
\end{equation}

\begin{proposition}[Weighted Metropolis defect]
\label{prop:v89-weighted-defect}
Under \eqref{eq:v89-local-oscillation}, for \(x\in D_{\rm in}\),
\begin{equation}
 \|P(x,\cdot)-G(x,\cdot)\|_W
 \leq
 \omega\{GW(x)+W(x)\}+\tau_W(x).
 \label{eq:v89-weighted-kernel}
\end{equation}
If
\[
 GW\leq\lambda W+b
\]
on \(D_{\rm in}\), then
\begin{equation}
 \|P(x,\cdot)-G(x,\cdot)\|_W
 \leq
 \omega\{(1+\lambda)W(x)+b\}+\tau_W(x).
 \label{eq:v89-weighted-drift-bound}
\end{equation}
\end{proposition}

\begin{proof}
From \eqref{eq:v89-kernel-difference},
\[
 \|P(x,\cdot)-G(x,\cdot)\|_W
 \leq
 \int\{1-a(x,y)\}W(y)\,G(x,dy)+r(x)W(x).
\]
On \(D_{\rm out}\), bound \(1-a\) by \(\omega\).  On its complement,
use \(1-a\leq1\).  The two exterior terms are exactly bounded by
\(\tau_W(x)\), while the interior terms are at most
\(\omega\{GW(x)+W(x)\}\).  The drift inequality gives the last claim.
\end{proof}

For the shell weight of V88, Proposition
\ref{prop:v88-shell-drift} supplies \(\lambda=\theta^2\).
A Gaussian fourth-moment estimate bounds \(\tau_W\) by the same
stretched-exponential exit with a polynomial prefactor, as in V84.

\subsection{A local invariant comparison}

Let the reference OU kernel \(G_{n,r}\) satisfy the V88 contraction in a
signed-measure norm \(\|\cdot\|_*\).  Suppose the Metropolis kernel
\(P_{n,r}\) is built from a local interaction remainder and that the
ordinary or weighted version of
\eqref{eq:v89-local-kernel-bound} gives
\begin{equation}
 \sup_{x\in D_{\rm in}}
 \|P_{n,r}(x,\cdot)-G_{n,r}(x,\cdot)\|_*
 \leq q_n.
 \label{eq:v89-good-kernel-q}
\end{equation}
If current states outside \(D_{\rm in}\) are assigned to a bad-state
measure of norm \(b_n\), then the invariant defect in
\eqref{eq:v86-invariant-defect} is at most \(q_n+2b_n\): split the
invariant law over \(D_{\rm in}\) and its complement and use the maximal
difference \(2\) on the latter in ordinary variation.  The weighted
version uses the weighted bad-state mass instead.

Consequently, Theorem \ref{thm:v86-invariant-comparison} gives
\begin{equation}
 \|\pi_{P,n,r}-\pi_{G,n,r}\|_*
 \leq
 \frac{q_n+2b_n}{1-\theta}.
 \label{eq:v89-local-invariant-comparison}
\end{equation}
This is a theorem once \(b_n\) is known.  In the SM--Einstein
application, \(b_n\) contains the physical chart exterior and has not
been bounded.

\subsection{Why a global refresh fails}

The small local oscillation cannot be replaced by the supremum of an
extensive interaction.

\begin{proposition}[Extensive independence-proposal rejection]
\label{prop:v89-extensive-rejection}
Let
\[
 \nu_N=\operatorname{Bernoulli}(1/2)^{\otimes N},
 \qquad
 V_N(x)=h\sum_{i=1}^Nx_i,
 \qquad h>0,
\]
and let \(\mu_N\propto e^{-V_N}\nu_N\).  Use the independence proposal
\(G_N(x,dy)=\nu_N(dy)\) in the Metropolis kernel.  If
\(X\sim\mu_N\), then
\begin{equation}
 {\mathbb E}a(X,Y)\longrightarrow0,
 \qquad
 {\mathbb E}r(X)\longrightarrow1
 \label{eq:v89-global-rejection}
\end{equation}
for an independent \(Y\sim\nu_N\).
\end{proposition}

\begin{proof}
Under \(\mu_N\), the coordinates are Bernoulli with
\[
 p_h=\frac1{1+e^h}<\frac12.
\]
Choose \(\delta>0\) so that
\(p_h+\delta<1/2-\delta\).  With probability tending to one,
\[
 \frac1N\sum_iX_i\leq p_h+\delta,
 \qquad
 \frac1N\sum_iY_i\geq\frac12-\delta
\]
by the law of large numbers.  On this event,
\[
 a(X,Y)
 \leq
 \exp\{-hN(1/2-p_h-2\delta)\},
\]
which tends to zero.  Since \(0\leq a\leq1\), dominated convergence on
the high-probability event and its complement gives
\({\mathbb E}a\to0\).  The rejection identity gives the second limit.
\end{proof}

Thus even a harmless product perturbation makes a global independence
proposal asymptotically singular.  The correct bridge uses local or
fixed-shell proposals, normalized conditional kernels, and the
cylinder-local topology of V83.

\subsection{Gravity-specific obligations}

For the proposed SM--Einstein shell bridge one must still prove:
\begin{enumerate}
\item the target conditional density is positive on the chosen real or
      deformed contour, or an alternative nonprobabilistic coefficient
      argument is provided;
\item the OU proposal and target share a measurable local shell state
      space;
\item the analytic chart has a uniform buffer;
\item the weighted probability of a current bad state is summable;
\item the exterior action makes \(Z_V\) finite; and
\item the Metropolis update respects reflection and all retained gauge,
      chiral, boundary, and determinant-line data.
\end{enumerate}

\begin{firewall}[Metropolis does not create coercivity]
\label{fire:v89-no-created-coercivity}
The identity \(\mu P=\mu\) assumes that
\(Z_V=\nu(e^{-V})\) is finite.  Metropolis correction does not prove this
integrability and does not bound the invariant mass outside the analytic
chart.  The missing weighted gravitational exterior reappears as the
bad-state owner \(b_n\) in
\eqref{eq:v89-local-invariant-comparison}.
\end{firewall}

\begin{assessment}[Exact status after V89]
\label{ass:v89-status}
V89 constructs an exact reversible physical kernel from the V88 Gaussian
proposal, proves ordinary and weighted local rejection bounds, derives
the invariant comparison with an explicit bad-state owner, and proves
the extensive global-refresh obstruction.  The local Taylor and proposal
exit pieces are summable under the V82 scale window.  The physical
bad-state mass, local shell realization, positivity on the physical
contour, retained-sector tightness, determinant-line gluing, symmetry
restoration, and the full SM--Einstein continuum remain open.
\end{assessment}

\subsection{Next acquisition}

V90 is the next compulsory companion audit.  It will inspect the newest
Yang--Mills and Navier--Stokes notes for an all-colour schedule,
bad-state invariant-mass estimate, or local influence comparison.  The
post-audit note will continue with the first transferable result and will
not treat the Metropolis bridge as closing the physical exterior.

\paragraph{Parent-version record.}
V89 was formed from
\texttt{Renewal\_SM\_Einstein\_Continuum\_Research\_Notes\_V88.tex}.
The V88 parent had \(31\,022\) counted source lines, \(1\,128\,696\) bytes,
and SHA--256
\[
 \texttt{03649B82E04D2A88275655D02A5063E12CA8A738EEFF2771A6A5414A4CD393D5}.
\]

% =============================================================================
% END APPENDED TENTH-SERIES V89 MATERIAL
% =============================================================================
% =============================================================================
% BEGIN APPENDED TENTH-SERIES V90 MATERIAL
% =============================================================================

\section{Tenth-series V90: companion audit and measured-coordinate correction}
\label{sec:v10v90}

\subsection{Compulsory companion audit}

The newest active companion notes at this gate were read without
modification:
\begin{align*}
 &\texttt{Renewal\_Yang\_Mills\_Mass\_Gap\_Research\_Notes\_V10.tex},
 \\
 &\qquad
 \operatorname{SHA256}
 =
 \texttt{EB228387210698AF0D6AA82B4194A843FFE0B920033EDF201C9232AEC36A517D},
 \\
 &\texttt{Renewal\_NS\_Stability\_Research\_Notes\_V26.tex},
 \\
 &\qquad
 \operatorname{SHA256}
 =
 \texttt{82C887F8C2C69E4F28FAD7C3DC8DE5B9099CB5A4BF431EBA1FFF3E91935978AD}.
\end{align*}
The Yang--Mills companion advanced from V7 to V10 after the V85 audit.
The Navier--Stokes companion remains unchanged.

\subsection{Exact triangular integration, not repeated sampling}

The Yang--Mills V8 construction resolves its colour schedule by
successive exact integration.  If the disjoint fibre sets are ordered by
colours \(C_1,\ldots,C_\chi\), it defines
\[
 W_a(x_{R_a};s)
 =
 \int W_{a-1}(x_{C_a},x_{R_a};s)\,dH_{C_a},
 \qquad
 R_a={\cal B}\setminus\bigcup_{j\leq a}C_j.
\]
At physical source, the ratios \(W_{a-1}/W_a\) are normalized
conditional kernels and telescope to the exact joint law.  A fibre is
integrated once and is never updated or reclassified.

This corrects the interpretation of V86--V89.  Their contractive Markov
kernel and invariant-resolvent route is mathematically valid, but it is
not required to identify a finite-step exact coarse pushforward.  For the
RG integration itself, triangular disintegration is the more direct
route.  A Markov bridge remains useful for comparison, sampling, or
mixing estimates only if its invariant-law hypotheses can be verified.

The V8 historical saturation rule is also transferable as abstract owner
bookkeeping: pass owners remain unchanged, hit owners are consumed and
replaced by connected outputs, and immutable birth labels prevent later
colours from charging a block twice.  Its Yang--Mills numerical smallness,
compact collar stability, and Wilson bad-block witness do not transfer to
gravity.  The analogous gravitational schedule still needs the weighted
physical exterior activity of V84.

\subsection{Measured coordinate changes}

The V9 companion proves a second structural point.  If a vector field
\({\cal W}\) generates a coordinate flow for a measure
\(e^{-S}d\omega\), the infinitesimal logarithmic action is
\begin{equation}
 {\cal O}_{\cal W}
 =
 DS[{\cal W}]-\operatorname{div}_{\omega}{\cal W}.
 \label{eq:v90-measured-owner}
\end{equation}
The action-gradient and density-Jacobian terms are inseparable.  A
postcomposition of an exact block map by that flow changes the
coordinate representative of the coarse measure while transporting
observables.  It does not delete a physical term at fixed coordinates.

This mechanism is directly relevant to the V78 beta-Jacobian owner
ledger.  A raw coefficient that lies in the image of an admissible local
measured field redefinition should be sectioned out before a physical
stability or symmetry-restoration matrix is estimated.  The Yang--Mills
companion constructs one scalar compensation direction.  The
SM--Einstein problem requires a multi-direction theorem with metric,
gauge, Higgs, fermionic, coarea, and boundary densities and with an
explicit invertibility condition for the compensation matrix.  V91 will
prove that abstract theorem.

\subsection{Finite-jet contraction instead of massless influence}

The V10 companion confirms the obstruction found in V87: a raw
single-site total-variation influence gap is not uniform in a massless
sector.  Its successful contraction acts instead on one finite
hypercubic anisotropy coordinate after the measured redundant direction
has been removed.  The derivative consists of:
\begin{enumerate}
\item the engineering factor \(b^{-2}\); and
\item a connected remainder with one distinguished anisotropy mark and
      at least one nonreference primitive.
\end{enumerate}
Normal convergence gives a marked bound of the form
\[
 \frac{Cw}{(1-Aw)^2},
\]
so the exact derivative is strictly less than one when the activity
\(w\) is sufficiently small.  This does not contract the massless field.

The transferable SM--Einstein target is therefore a finite local
symmetry-breaking jet after quotienting admissible measured-coordinate
directions.  Candidate rows include regulator-induced Lorentz or
diffeomorphism anisotropies and irrelevant gauge-breaking completions.
Physical relevant and marginal couplings, anomaly classes, global
topological sectors, and infrared fields must remain outside that
contraction.

\subsection{No Navier--Stokes transfer}

The Navier--Stokes V26 pointwise Oseen-kernel derivatives do not provide
a measured field-coordinate quotient, an exact quantum block
disintegration, a weighted gravitational exterior estimate, or a
finite-jet SM symmetry-restoration matrix.  No result is transferred
from that file at V90.

\subsection{Adjusted dependency order}

The corrected order is:
\begin{enumerate}
\item use exact triangular integration for the finite regulator
      pushforward, with normalized conditional factors only as owner
      devices;
\item retain the V86--V89 kernel route as an optional comparison theorem,
      not as the definition of the pushforward;
\item construct a multi-component measured-coordinate section and
      transport observables exactly;
\item project the finite local response onto physical and redundant
      quotient rows;
\item use one-mark connected estimates to contract irrelevant
      symmetry-breaking jets;
\item separately prove the weighted physical exterior, retained-sector
      tightness, and local projective comparison.
\end{enumerate}

\begin{definition}[V90 exact-pushforward card]
\label{def:v90-pushforward-card}
One SM--Einstein regulator step passes the corrected card if:
\begin{enumerate}
\item the chosen fine fibres form a genuine disjoint product after all
      constraints, dependencies, and retained modes are declared;
\item successive colour integrals are finite and satisfy Fubini or
      Tonelli at physical source;
\item normalized conditional factors carry no duplicated determinant;
\item pass, hit, jet, and bad-component owners follow an immutable
      triangular history;
\item every coordinate postcomposition acts on the full measure density
      and on observables;
\item the measured redundant response matrix has a uniform inverse; and
\item finite-jet contractions are asserted only after the redundant
      section is fixed.
\end{enumerate}
\end{definition}

\begin{firewall}[Audit transfer limits]
\label{fire:v90-transfer-limits}
The Yang--Mills V8 schedule relies on finite-volume integrability, a
compact local fibre, and a Wilson-action witness for bad collars.  V90
imports its exact triangular algebra and unique-payment logic, not those
coercive inputs.  The V10 contraction is a finite irrelevant-jet result;
it gives no SM--Einstein field-space mass gap, no gravitational exterior
bound, and no retained infrared tightness.
\end{firewall}

\begin{assessment}[Status after the V90 audit]
\label{ass:v90-status}
The audit replaces an unnecessary invariant-sampler dependency by exact
triangular pushforward, imports measured-coordinate ownership, and
redirects contraction toward finite symmetry-breaking jets after
quotienting redundant rows.  V91 will prove the multi-component
measured-coordinate section.  The weighted gravitational exterior,
physical quotient identification, anomaly and determinant-line data,
retained-sector tightness, symmetry restoration, and the full
SM--Einstein continuum remain open.
\end{assessment}

\paragraph{Parent-version record.}
V90 was formed from
\texttt{Renewal\_SM\_Einstein\_Continuum\_Research\_Notes\_V89.tex}.
The V89 parent had \(31\,346\) counted source lines, \(1\,139\,359\) bytes,
and SHA--256
\[
 \texttt{1A2120223763833B76EDF92D5A03778040DCD1608EE37A5D651A773BB0A94A0A}.
\]

% =============================================================================
% END APPENDED TENTH-SERIES V90 MATERIAL
% =============================================================================
% =============================================================================
% BEGIN APPENDED TENTH-SERIES V91 MATERIAL
% =============================================================================

\section{Tenth-series V91: multi-direction measured-coordinate section}
\label{sec:v10v91}

V90 identified measured field redefinitions as the correct quotient to
take before contracting finite symmetry-breaking jets.  This note proves
the required multi-direction theorem.  It applies to a finite regulator
with a declared positive bosonic density.  A chiral fermionic application
requires the separate determinant-line and anomaly hypotheses stated
below.

\subsection{Measured flows}

Let \(M\) be an oriented finite-dimensional manifold, possibly with a
declared boundary-condition submanifold, and let \(d\omega\) be a smooth
positive density.  Let \(S:M\to{\mathbb R}\) be smooth and suppose
\(e^{-S}d\omega\) is finite.  For a vector field \({\cal W}\) whose local
flow \(\varphi_t\) exists on the integration domain, define the transformed
measured action by
\begin{equation}
 e^{-S^{[t]}}d\omega
 :=
 (\varphi_{-t})_*(e^{-S}d\omega).
 \label{eq:v91-measured-action}
\end{equation}
Equivalently,
\begin{equation}
 S^{[t]}(y)
 =
 S(\varphi_t(y))
 -
 \log\operatorname{Jac}_{\omega}\varphi_t(y).
 \label{eq:v91-action-Jacobian}
\end{equation}

\begin{lemma}[Measured tangent]
\label{lem:v91-measured-tangent}
At \(t=0\),
\begin{equation}
 \left.\frac d{dt}\right|_{t=0}S^{[t]}
 =
 DS[{\cal W}]
 -
 \operatorname{div}_{\omega}{\cal W}
 =:{\cal O}_{\cal W}.
 \label{eq:v91-measured-owner}
\end{equation}
For every integrable observable \(F\),
\begin{equation}
 \int_MF\,e^{-S^{[t]}}d\omega
 =
 \int_MF\circ\varphi_{-t}\,e^{-S}d\omega.
 \label{eq:v91-observable-transport}
\end{equation}
\end{lemma}

\begin{proof}
Differentiate \eqref{eq:v91-action-Jacobian} and use
\[
 \left.\frac d{dt}\right|_{t=0}
 \log\operatorname{Jac}_{\omega}\varphi_t
 =
 \operatorname{div}_{\omega}{\cal W}.
\]
Equation \eqref{eq:v91-observable-transport} is the defining pushforward
identity \eqref{eq:v91-measured-action}.
\end{proof}

Thus \(DS[{\cal W}]\) alone is not a redundant owner.  The divergence is
part of the same measured coordinate tangent.

\subsection{Compatibility with an exact coarse map}

Let \({\cal B}:M_{\rm f}\to M_{\rm c}\) be a measurable block map and let
\(\nu_\epsilon\) be a differentiable family of fine measures.  Suppose
\[
 \rho_\epsilon d\omega_{\rm c}
 =
 {\cal B}_*\nu_\epsilon.
\]
For a coarse vector field \({\cal W}\) with flow \(\varphi_t\), define
\begin{equation}
 \widetilde\rho_\epsilon d\omega_{\rm c}
 :=
 (\varphi_{-\epsilon\lambda})_*
 (\rho_\epsilon d\omega_{\rm c}).
 \label{eq:v91-postcomposed-density}
\end{equation}

\begin{proposition}[Exact postcomposition]
\label{prop:v91-postcomposition}
The measure in \eqref{eq:v91-postcomposed-density} is the exact
pushforward of \(\nu_\epsilon\) under
\[
 \varphi_{-\epsilon\lambda}\circ{\cal B}.
\]
If the logarithmic response of the original block map is represented by
the measured current \(U\), the response after postcomposition is
represented by
\begin{equation}
 U+\lambda{\cal W}.
 \label{eq:v91-current-shift}
\end{equation}
All observables must be transported by
\eqref{eq:v91-observable-transport}.
\end{proposition}

\begin{proof}
Functoriality gives
\[
 (\varphi_{-\epsilon\lambda})_*{\cal B}_*\nu_\epsilon
 =
 (\varphi_{-\epsilon\lambda}\circ{\cal B})_*\nu_\epsilon.
\]
Differentiate at \(\epsilon=0\).  Lemma
\ref{lem:v91-measured-tangent} and linearity in the velocity give
\eqref{eq:v91-current-shift}.
\end{proof}

This proposition is compatible with the triangular pushforward of V90:
postcomposition occurs on the retained coarse variables after the exact
fine fibres have been integrated.

\subsection{The finite redundant response matrix}

Let \({\cal V}\) be a Banach space of local measured currents with norm
\(\|\cdot\|_{\cal V}\).  Let
\[
 {\mathfrak z}:{\cal V}\longrightarrow{\mathbb R}^m
\]
be a bounded finite-jet extractor for the candidate redundant rows.  Let
\[
 {\cal W}:{\mathbb R}^m\longrightarrow{\cal V},
 \qquad
 {\cal W}\lambda=\sum_{a=1}^m\lambda_a{\cal W}_a,
\]
be a bounded family of admissible local compensation fields.  Define the
response matrix
\begin{equation}
 C:={\mathfrak z}{\cal W}\in{\mathbb R}^{m\times m}.
 \label{eq:v91-response-matrix}
\end{equation}

\begin{theorem}[Multi-direction measured-coordinate section]
\label{thm:v91-multisection}
Assume \(C\) is invertible.  Define
\begin{align}
 \lambda(U)&:=-C^{-1}{\mathfrak z}U,
 \label{eq:v91-compensation-coeff}\\
 \Pi U&:=U+{\cal W}\lambda(U)
 =
 \left(I-{\cal W}C^{-1}{\mathfrak z}\right)U.
 \label{eq:v91-section-projector}
\end{align}
Then:
\begin{enumerate}
\item \({\mathfrak z}(\Pi U)=0\);
\item \(\Pi^2=\Pi\);
\item \(\operatorname{ran}\Pi=\ker{\mathfrak z}\) and
      \(\ker\Pi=\operatorname{ran}{\cal W}\);
\item
\begin{equation}
 \|\Pi\|
 \leq
 1+\|{\cal W}\|\,\|C^{-1}\|\,\|{\mathfrak z}\|;
 \label{eq:v91-projector-bound}
\end{equation}
and
\item every correction \({\cal W}\lambda(U)\) is realized by exact
      measured coarse-coordinate postcomposition as in Proposition
      \ref{prop:v91-postcomposition}.
\end{enumerate}
If every \({\cal W}_a\) has zero vacuum linearization, then \(\Pi U\) and
\(U\) have the same vacuum linearization.
\end{theorem}

\begin{proof}
Using \(C={\mathfrak z}{\cal W}\),
\[
 {\mathfrak z}\Pi U
 =
 {\mathfrak z}U-CC^{-1}{\mathfrak z}U=0.
\]
Also
\[
 \Pi{\cal W}
 =
 {\cal W}-{\cal W}C^{-1}C=0.
\]
Since \(\Pi U=U\) on \(\ker{\mathfrak z}\), the range, kernel, and
idempotence statements follow.  The triangle inequality gives
\eqref{eq:v91-projector-bound}.  Proposition
\ref{prop:v91-postcomposition} realizes each linear combination of
compensation fields by the flow of that combination for sufficiently
small postcomposition parameter.  Zero linearization is preserved
because the correction is a linear combination of fields with zero
linearization.
\end{proof}

The inverse bound on \(C\) is a finite-dimensional pivot.  It belongs in
the same interval and reserve ledger as the pivots of V63 and the
beta-Jacobian projectors of V78.

\subsection{Retained-parameter dependence}

Let \(r\) denote retained physical couplings and modes.  Suppose
\({\mathfrak z}(r)\), \({\cal W}(r)\), and
\[
 C(r)={\mathfrak z}(r){\cal W}(r)
\]
are differentiable and
\begin{equation}
 \sup_r\|C(r)^{-1}\|\leq K_C.
 \label{eq:v91-uniform-pivot}
\end{equation}

\begin{proposition}[Derivative of the section]
\label{prop:v91-section-derivative}
The projector
\[
 \Pi(r)=I-{\cal W}(r)C(r)^{-1}{\mathfrak z}(r)
\]
is differentiable and, in a retained direction \(h\),
\begin{align}
 D\Pi(r)[h]
 ={}&
 -D{\cal W}(r)[h]\,C(r)^{-1}{\mathfrak z}(r)
 \nonumber\\
 &+
 {\cal W}(r)C(r)^{-1}DC(r)[h]C(r)^{-1}{\mathfrak z}(r)
 \nonumber\\
 &-
 {\cal W}(r)C(r)^{-1}D{\mathfrak z}(r)[h].
 \label{eq:v91-projector-derivative}
\end{align}
Consequently,
\begin{align}
 \|D\Pi(r)\|
 \leq{}&
 \|D{\cal W}(r)\|K_C\|{\mathfrak z}(r)\|
 \nonumber\\
 &+
 \|{\cal W}(r)\|K_C^2\|DC(r)\|
       \|{\mathfrak z}(r)\|
 +
 \|{\cal W}(r)\|K_C\|D{\mathfrak z}(r)\|.
 \label{eq:v91-projector-derivative-bound}
\end{align}
\end{proposition}

\begin{proof}
Differentiate \eqref{eq:v91-section-projector} and use
\[
 D(C^{-1})[h]=-C^{-1}DC[h]C^{-1}.
\]
Taking operator norms gives the last estimate.
\end{proof}

This derivative is a genuine owner in a moving-reference RG step.  It
cannot be omitted when the retained variables move.

\subsection{Symmetry and locality}

Let a group \({\cal G}\) act continuously on \({\cal V}\) and on
\({\mathbb R}^m\).  Suppose
\[
 {\mathfrak z}(gU)=g{\mathfrak z}(U),
 \qquad
 {\cal W}(g\lambda)=g{\cal W}(\lambda).
\]
Then \(C\) and \(C^{-1}\) intertwine the finite-dimensional action and
\begin{equation}
 \Pi(gU)=g\Pi(U).
 \label{eq:v91-equivariant-section}
\end{equation}
If the extractors and compensation fields have fixed range, the section
is local in the owner norm.  If they are only exponentially quasi-local,
their tails must be included in the V83 leakage budget.

\begin{proof}
Intertwining gives
\[
 C(g\lambda)
 =
 {\mathfrak z}{\cal W}(g\lambda)
 =
 gC\lambda.
\]
Hence \(C^{-1}\) also intertwines.  Substitute in
\eqref{eq:v91-section-projector}.  The locality statement follows by
composition of the corresponding kernel bounds.
\end{proof}

\subsection{Quotient response and finite-jet stability}

Let \({\cal R}_r\) be one exact coarse response map on currents, and
define its sectioned representative by
\begin{equation}
 {\cal R}_r^\circ:=\Pi(r){\cal R}_r.
 \label{eq:v91-sectioned-RG}
\end{equation}
At fixed \(r\),
\begin{equation}
 D{\cal R}_r^\circ(U)
 =
 \Pi(r)D{\cal R}_r(U).
 \label{eq:v91-sectioned-Jacobian}
\end{equation}
If \(r\) changes with \(U\), the chain rule adds
\begin{equation}
 D\Pi(r(U))[Dr(U)]\,{\cal R}_{r(U)}(U),
 \label{eq:v91-moving-section-owner}
\end{equation}
controlled by Proposition \ref{prop:v91-section-derivative}.  Thus the
physical finite-jet stability matrix is the derivative after sectioning,
with the moving-section term included.  An eigenvalue seen only along
\(\operatorname{ran}{\cal W}\) is a coordinate eigenvalue, not a physical
instability.

\subsection{SM--Einstein admissibility}

For a Standard Model--Einstein regulator, a compensation field is
admissible only if:
\begin{enumerate}
\item its flow preserves the regulator domain, boundary conditions, and
      the declared constraint/gauge slice, or transports them with an
      exact coarea identity;
\item it is local or has a quantified quasi-local kernel;
\item it respects reflection and the real integration contour;
\item the full bosonic, ghost, and fermionic measure Jacobian is included;
\item its action on observables and sources is transported exactly; and
\item it does not change a quantized topological sector or a physical
      anomaly class.
\end{enumerate}

For chiral fermions, the divergence in
\eqref{eq:v91-measured-owner} is replaced by the logarithmic derivative
of the complete regulated Berezinian/determinant section.  A local formula
is not enough if the determinant line is globally nontrivial.

\begin{firewall}[Anomalous directions are not redundant]
\label{fire:v91-anomaly}
A field redefinition whose regulated Jacobian carries a nonzero local or
global anomaly cannot be placed in
\(\operatorname{ran}{\cal W}\) and quotiented by
\eqref{eq:v91-section-projector}.  Standard Model anomaly cancellation
must be proved in the common gravity-compatible regulator, including
boundary and determinant-line gluing.  V91 assumes admissibility; it does
not prove it.
\end{firewall}

\begin{definition}[V91 measured-section card]
\label{def:v91-section-card}
The redundant sector passes the card if:
\begin{enumerate}
\item the complete regulated density and its divergence are defined;
\item the selected compensation flows are exact admissible measured
      coordinate changes;
\item the response matrix \(C(r)\) has a uniform inverse;
\item its inverse, the extractors, and the flows pass the owner-locality
      norm;
\item the derivative reserve
      \eqref{eq:v91-projector-derivative-bound} fits the V78 Jacobian
      interval; and
\item anomaly, topology, boundary, and determinant-line directions are
      excluded unless their exact triviality is proved.
\end{enumerate}
\end{definition}

\begin{assessment}[Exact status after V91]
\label{ass:v91-status}
V91 proves a multi-direction measured-coordinate projection, its exact
realization by block-map postcomposition, norm and retained-parameter
derivative bounds, symmetry and locality inheritance, and the correct
sectioned response formula.  It supplies the abstract quotient required
before contracting finite SM--Einstein symmetry-breaking jets.  The
actual compensation matrix, its uniform pivot, chiral admissibility,
weighted exterior, retained-sector tightness, symmetry restoration, and
the full SM--Einstein continuum remain unproved.
\end{assessment}

\subsection{Next acquisition}

V92 will combine the V91 section with the one-distinguished-mark forest
method.  It will prove a finite-dimensional contraction theorem for a
general irrelevant symmetry-breaking jet with engineering dimension
\(-\delta<0\), including affine forcing, a moving section, and cofinal
restoration.  The theorem will state which numerical rows must be
computed for the SM--Einstein regulator.

\paragraph{Parent-version record.}
V91 was formed from
\texttt{Renewal\_SM\_Einstein\_Continuum\_Research\_Notes\_V90.tex}.
The V90 parent had \(31\,539\) counted source lines, \(1\,147\,719\) bytes,
and SHA--256
\[
 \texttt{9411AE565D2E10482FDFEE70C1B21E8622CE4386640A75735EA2D5B2C3B55EB1}.
\]

% =============================================================================
% END APPENDED TENTH-SERIES V91 MATERIAL
% =============================================================================
% =============================================================================
% BEGIN APPENDED TENTH-SERIES V92 MATERIAL
% =============================================================================

\section{Tenth-series V92: sectioned irrelevant-jet contraction}
\label{sec:v10v92}

V91 constructs the measured-coordinate quotient.  This note proves the
finite-dimensional contraction that becomes available after taking that
quotient.  It applies to genuinely irrelevant local symmetry-breaking
jets.  Relevant and marginal Ward, anomaly, and physical coupling rows
are not included.

\subsection{The sectioned jet space}

Let \(E_{\rm sb}\) be a finite-dimensional normed space of local
symmetry-breaking jets after:
\begin{enumerate}
\item constants and declared physical relevant or marginal couplings have
      been extracted;
\item measured redundant directions have been removed by the V91
      projector \(\Pi(r)\); and
\item the remaining stable owner complement has been fixed.
\end{enumerate}
Let \(r\) collect all retained parameters.  For blocking factor \(b>1\),
suppose the Gaussian principal derivative on \(E_{\rm sb}\) is
\(L_{0,r}\) and
\begin{equation}
 \sup_r\|L_{0,r}\|\leq\theta_0:=b^{-\delta},
 \qquad
 \delta>0.
 \label{eq:v92-engineering-gap}
\end{equation}
For an operator of canonical dimension \(\Delta>D\), the engineering
value is \(\delta=\Delta-D\).

Let \(w\geq0\) denote the complete nonreference rooted activity after the
triangular pushforward, including regular kernel defects, bad-component
witnesses, generated hit owners, chart-exit terms, and the stable
completion defect.  The physical exterior is part of \(w\) only if its
required bound has actually been proved.

\subsection{One distinguished mark}

Let \(M_{\rm sb}\) be the norm of one normalized
\(E_{\rm sb}\)-source insertion, let \(c_{\rm sb}\geq1\) contain the fixed
owner, source-radius, jet, and extractor conversions, and let
\(A_{\rm F}\) be the connected-forest exploration constant.  Put
\begin{equation}
 x:=A_{\rm F}c_{\rm sb}w.
 \label{eq:v92-marked-ratio}
\end{equation}

\begin{lemma}[Marked connected sum]
\label{lem:v92-marked-sum}
Assume \(x<1\) and that the total norm of connected words containing one
distinguished symmetry-breaking mark and \(n+1\) nonreference primitives
is at most
\[
 M_{\rm sb}c_{\rm sb}w\,(n+1)x^n
 \qquad(n\geq0).
\]
Then their normally convergent sum \({\cal E}_{\rm sb}\) obeys
\begin{equation}
 \|{\cal E}_{\rm sb}\|
 \leq
 \varepsilon_{\rm sb}(w)
 :=
 \frac{M_{\rm sb}c_{\rm sb}w}{(1-x)^2}.
 \label{eq:v92-marked-bound}
\end{equation}
In particular, if \(x\leq1/2\), then
\begin{equation}
 \varepsilon_{\rm sb}(w)
 \leq4M_{\rm sb}c_{\rm sb}w.
 \label{eq:v92-marked-half}
\end{equation}
\end{lemma}

\begin{proof}
Absolute convergence permits grouping by the number of nonreference
primitives.  Sum the assumed bounds and use
\[
 \sum_{n\geq0}(n+1)x^n=(1-x)^{-2}.
\]
The last assertion follows from \((1-x)^{-2}\leq4\).
\end{proof}

The hypothesis is exactly what a one-source derivative of the rooted
connected forest must prove.  The distinguished mark fixes a root, so
the constants do not grow with volume.

\subsection{Fixed-retained-parameter contraction}

Let \({\cal R}_{n,r}\) be one exact triangular coarse response followed
by all declared projections and the measured section.  Let
\({\mathfrak a}_r\) be the bounded extractor onto \(E_{\rm sb}\).  For
\(\alpha\) in a ball \(B_R\subset E_{\rm sb}\), define
\begin{equation}
 T_{n,r}(\alpha)
 :=
 {\mathfrak a}_r
 {\cal R}_{n,r}^{\circ}(S_r+\alpha),
 \label{eq:v92-sectioned-map}
\end{equation}
where the calibrated base \(S_r\) contains the Gaussian reference.  Write
\begin{equation}
 f_{n,r}:=T_{n,r}(0),
 \qquad
 F_{n,r}(\alpha):=T_{n,r}(\alpha)-f_{n,r}.
 \label{eq:v92-affine-split}
\end{equation}

\begin{theorem}[Irrelevant-jet contraction]
\label{thm:v92-jet-contraction}
Suppose:
\begin{enumerate}
\item the Gaussian principal derivative satisfies
      \eqref{eq:v92-engineering-gap};
\item the differentiated nonreference connected output satisfies Lemma
      \ref{lem:v92-marked-sum};
\item the extractor norm is at most \(C_{\mathfrak a}\); and
\item the retained parameter \(r\) is fixed during the one-step
      Lipschitz comparison.
\end{enumerate}
Then
\begin{equation}
 \|F_{n,r}(\alpha)-F_{n,r}(\alpha')\|
 \leq q_n
 \|\alpha-\alpha'\|,
 \label{eq:v92-Lipschitz}
\end{equation}
where
\begin{equation}
 q_n
 :=
 \theta_0+C_{\mathfrak a}\varepsilon_{\rm sb}(w_n).
 \label{eq:v92-contraction-factor}
\end{equation}
If \(q_n<1\) and
\begin{equation}
 \|f_{n,r}\|\leq(1-q_n)R,
 \label{eq:v92-invariant-ball}
\end{equation}
then \(T_{n,r}\) maps \(B_R\) to itself and has a unique fixed point
\(\alpha_{n,*}(r)\) there, with
\begin{equation}
 \|\alpha_{n,*}(r)\|
 \leq
 \frac{\|f_{n,r}\|}{1-q_n}.
 \label{eq:v92-fixed-graph-bound}
\end{equation}
\end{theorem}

\begin{proof}
Differentiate \eqref{eq:v92-affine-split} along the segment from
\(\alpha'\) to \(\alpha\).  The sectioned Gaussian principal term has
norm at most \(\theta_0\).  After subtracting that term, every derivative
contains one distinguished \(E_{\rm sb}\) mark and at least one
nonreference primitive.  Lemma \ref{lem:v92-marked-sum} and the extractor
norm bound give
\[
 \sup_{\alpha\in B_R}\|DF_{n,r}(\alpha)\|
 \leq q_n.
\]
Integration along the segment proves \eqref{eq:v92-Lipschitz}.  Also,
\[
 \|T_{n,r}(\alpha)\|
 \leq\|f_{n,r}\|+q_n\|\alpha\|\leq R
\]
under \eqref{eq:v92-invariant-ball}.  Banach's fixed-point theorem gives
existence and uniqueness, and the fixed-point equation gives
\eqref{eq:v92-fixed-graph-bound}.
\end{proof}

Because \(\theta_0<1\), there is a concrete activity threshold.  For
example, if
\begin{equation}
 A_{\rm F}c_{\rm sb}w_n\leq\frac12,
 \qquad
 4C_{\mathfrak a}M_{\rm sb}c_{\rm sb}w_n
 \leq\frac{1-\theta_0}{2},
 \label{eq:v92-explicit-gate}
\end{equation}
then
\begin{equation}
 q_n\leq\frac{1+\theta_0}{2}<1.
 \label{eq:v92-explicit-contraction}
\end{equation}

\subsection{Moving-section debit}

Suppose now that the retained coordinate used in the step depends on
\(\alpha\), so \(r=r(\alpha)\).  With
\[
 T_n(\alpha)
 =
 {\mathfrak a}_{r(\alpha)}
 \Pi(r(\alpha))
 {\cal R}_{n,r(\alpha)}(S_{r(\alpha)}+\alpha),
\]
the derivative has, in addition to the fixed-\(r\) term, contributions
from \(D{\mathfrak a}\), \(D\Pi\), \(D_r{\cal R}\), and \(DS_r\).

\begin{proposition}[Moving-section contraction reserve]
\label{prop:v92-moving-section}
Assume the total norm of all parameter-motion contributions is at most
\(\varepsilon_{{\rm mov},n}\) on \(B_R\).  A sufficient explicit bound
contains the term
\begin{equation}
 \|{\mathfrak a}_r\|\,
 \|D\Pi(r)\|\,
 \|Dr(\alpha)\|\,
 \|{\cal R}_{n,r}(S_r+\alpha)\|,
 \label{eq:v92-moving-projector-term}
\end{equation}
with \(\|D\Pi(r)\|\) bounded by
\eqref{eq:v91-projector-derivative-bound}, plus the analogous extractor,
response, and base-chart terms.  Then Theorem
\ref{thm:v92-jet-contraction} holds with
\begin{equation}
 q_n^{\rm mov}
 =
 \theta_0+C_{\mathfrak a}\varepsilon_{\rm sb}(w_n)
 +\varepsilon_{{\rm mov},n}.
 \label{eq:v92-moving-factor}
\end{equation}
\end{proposition}

\begin{proof}
Apply the product and chain rules.  Proposition
\ref{prop:v91-section-derivative} gives
\eqref{eq:v92-moving-projector-term}.  By hypothesis all parameter-motion
terms sum to \(\varepsilon_{{\rm mov},n}\).  Add this to the derivative
bound in the proof of Theorem \ref{thm:v92-jet-contraction}.
\end{proof}

This debit is absent when contraction is proved fibrewise at fixed
retained data and the fixed graph is only afterward shown analytic in
\(r\).  It is mandatory in a simultaneous coupled iteration.

\subsection{Cofinal restoration}

\begin{lemma}[Nonautonomous forced contraction]
\label{lem:v92-forced-contraction}
Let
\begin{equation}
 \|\alpha_{n+1}\|
 \leq q\|\alpha_n\|+g_n,
 \qquad
 0\leq q<1.
 \label{eq:v92-forced-recurrence}
\end{equation}
Then
\begin{equation}
 \|\alpha_N\|
 \leq
 q^N\|\alpha_0\|
 +
 \sum_{j=0}^{N-1}q^{N-1-j}g_j.
 \label{eq:v92-forced-solution}
\end{equation}
If \(g_n\to0\), then \(\alpha_n\to0\).
\end{lemma}

\begin{proof}
Induction gives \eqref{eq:v92-forced-solution}.  Given \(\epsilon>0\),
choose \(J\) so that \(g_j\leq\epsilon(1-q)/2\) for \(j\geq J\).
The tail of the sum is at most \(\epsilon/2\), and the finite initial part
and \(q^N\|\alpha_0\|\) tend to zero.
\end{proof}

\begin{corollary}[Conditional symmetry restoration]
\label{cor:v92-conditional-restoration}
Suppose along a cofinal regulator trajectory:
\begin{enumerate}
\item the section pivots and all owner radii remain in their uniform
      domains;
\item \(w_n\to0\) and
      \(\varepsilon_{{\rm mov},n}\to0\);
\item the factors in \eqref{eq:v92-moving-factor} are eventually bounded
      by one common \(q<1\); and
\item the affine symmetry-breaking forcing
      \(g_n:=\|T_n(0)\|\) tends to zero.
\end{enumerate}
Then the sectioned irrelevant symmetry-breaking coordinate tends to zero.
\end{corollary}

\begin{proof}
Apply Proposition \ref{prop:v92-moving-section} and Lemma
\ref{lem:v92-forced-contraction}.
\end{proof}

\subsection{Which SM--Einstein rows can use the theorem}

The engineering gap \(\delta>0\) is indispensable.  In four dimensions:
\begin{enumerate}
\item a local operator of dimension \(4+\delta\) can enter
      \(E_{\rm sb}\) if it is not physical, anomalous, topological, or
      redundant;
\item a dimension-four row has \(\theta_0=1\) and needs a genuine
      beta-function or Ward-identity calculation;
\item a relevant row has \(\theta_0>1\) and must be tuned or placed in
      the unstable coordinate set of V77; and
\item a measured redundant row must be removed by V91 before the
      extractor \({\mathfrak a}\) is defined.
\end{enumerate}

Candidate irrelevant rows include regulator-induced local Lorentz,
hypercubic, or diffeomorphism-breaking completions of dimension greater
than four.  Their actual independence, canonical dimension, source
normalization, and Gaussian principal matrix must be computed in the
chosen common regulator.  V92 does not declare any named row to pass
without that computation.

\begin{definition}[V92 finite-jet contraction card]
\label{def:v92-jet-card}
A proposed symmetry-restoration row passes the card if:
\begin{enumerate}
\item its class survives the physical, topological, anomaly, and measured
      redundant quotients;
\item its Gaussian principal derivative has a certified gap
      \(\theta_0<1\);
\item the one-mark connected forest satisfies
      \eqref{eq:v92-marked-bound};
\item the physical exterior is included in \(w_n\) with a proved
      activity;
\item the V91 pivot and moving-section reserve remain uniform;
\item affine forcing tends to zero along the cofinal trajectory; and
\item the entire trajectory stays inside the invariant owner domain.
\end{enumerate}
\end{definition}

\begin{firewall}[No restoration without a trajectory]
\label{fire:v92-no-trajectory}
The contraction theorem is conditional on the owner activity, the
measured-section pivot, and the affine forcing along an actual cofinal
trajectory.  It does not construct that trajectory.  It also does not
apply to marginal gauge couplings, the cosmological and Einstein terms,
Yukawa or Higgs marginal data, theta angles, anomaly classes, or retained
massless fields merely by renaming them as symmetry-breaking jets.
\end{firewall}

\begin{assessment}[Exact status after V92]
\label{ass:v92-status}
V92 proves a general sectioned irrelevant-jet contraction, an explicit
small-activity gate, the moving-section derivative reserve, a unique
calibrated fixed graph, and conditional cofinal restoration.  This
converts the V90 companion mechanism into a multi-field theorem
compatible with the V77 stable/centre/unstable ledger.  The named
SM--Einstein jet matrix, weighted physical exterior, cofinal trajectory,
anomaly and determinant-line data, retained-sector tightness, and the
full continuum remain unproved.
\end{assessment}

\subsection{Next acquisition}

V93 will return to the largest remaining analytic obstruction: the
weighted physical exterior.  It will formulate a contour-coercivity
theorem for a finite-dimensional conformal factor coupled to a positive
quartic Higgs direction, determine the exact Schur condition preventing
runaway mixed directions, and separate that toy coercive sector from the
unresolved derivative and gauge-field exterior.

\paragraph{Parent-version record.}
V92 was formed from
\texttt{Renewal\_SM\_Einstein\_Continuum\_Research\_Notes\_V91.tex}.
The V91 parent had \(31\,944\) counted source lines, \(1\,160\,662\) bytes,
and SHA--256
\[
 \texttt{63F0E1871976523CF5633568D8B79122ADA3F91165C8E8427E1C3699DDAA4677}.
\]

% =============================================================================
% END APPENDED TENTH-SERIES V92 MATERIAL
% =============================================================================
% =============================================================================
% BEGIN APPENDED TENTH-SERIES V93 MATERIAL
% =============================================================================

\section{Tenth-series V93: coercivity-to-tail theorem and the conformal--Higgs test}
\label{sec:v10v93}

The weighted physical exterior is now the largest common hypothesis in
V83, V84, V89, and V92.  This note proves a finite-dimensional
coercivity-to-tail theorem and an exact test for a mixed conformal--Higgs
quartic.  It also proves that a Higgs quartic cannot stabilize a pure
negative conformal direction when the conformal quartic vanishes.

\subsection{From a real-part lower bound to a weighted exterior}

Let \(S:{\mathbb R}^d\to{\mathbb C}\) be measurable.  The following
statement bounds the absolute density \(e^{-\operatorname{Re}S}\).  If
\(S\) is real, it is a probability-tail theorem after normalization.

\begin{theorem}[Coercivity-to-weighted-tail]
\label{thm:v93-coercive-tail}
Suppose there are constants \(a>0\), \(p>0\), \(b,B<\infty\), and
\(r_0>0\) such that
\begin{align}
 \operatorname{Re}S(x)&\geq a\|x\|^p-b
 \quad\text{for all }x,
 \label{eq:v93-global-coercivity}\\
 \operatorname{Re}S(x)&\leq B
 \quad\text{when }\|x\|\leq r_0.
 \label{eq:v93-local-upper}
\end{align}
Then
\[
 Z_{\rm abs}:=\int_{{\mathbb R}^d}e^{-\operatorname{Re}S(x)}\,dx
\]
is finite and satisfies
\begin{equation}
 Z_{\rm abs}
 \geq
 |B_d|r_0^d e^{-B},
 \label{eq:v93-partition-lower}
\end{equation}
where \(|B_d|\) is the unit-ball volume.  For every \(q\geq0\), there
are \(C<\infty\) and \(R_0<\infty\), depending only on the displayed
data and \(q,d\), such that
\begin{equation}
 \frac1{Z_{\rm abs}}
 \int_{\|x\|>R}
 (1+\|x\|^q)e^{-\operatorname{Re}S(x)}\,dx
 \leq
 C e^{-aR^p/2}
 \qquad(R\geq R_0).
 \label{eq:v93-weighted-tail}
\end{equation}
\end{theorem}

\begin{proof}
Equation \eqref{eq:v93-local-upper} gives
\eqref{eq:v93-partition-lower}.  The lower bound
\eqref{eq:v93-global-coercivity} gives
\[
 \int_{\|x\|>R}(1+\|x\|^q)e^{-\operatorname{Re}S(x)}\,dx
 \leq
 e^b|{\mathbb S}^{d-1}|
 \int_R^\infty(1+r^q)r^{d-1}e^{-ar^p}\,dr.
\]
For sufficiently large \(r\), the polynomial
\((1+r^q)r^{d-1}\) is at most a fixed constant times
\(e^{ar^p/4}\).  Increasing \(R_0\) if necessary, the remaining radial
integral of \(e^{-3ar^p/4}\) is bounded by a constant times
\(e^{-aR^p/2}\).  Divide by
\eqref{eq:v93-partition-lower}.  The same estimate without the
restriction \(\|x\|>R\) proves finiteness.
\end{proof}

The local upper bound is the normalization stock.  A coercive lower bound
alone does not prevent the normalized denominator from being
arbitrarily small in a parameter-dependent family.

\subsection{Exact mixed-quartic condition}

Consider two real amplitudes \(s,h\), representing one conformal and one
Higgs radial direction after a declared real contour has been chosen.
Let
\begin{equation}
 P_4(s,h)
 =
 a s^4+c s^2h^2+b h^4,
 \qquad a,b>0.
 \label{eq:v93-mixed-quartic}
\end{equation}
Put \(c_-:=\max\{0,-c\}\) and
\begin{equation}
 \lambda_*:=
 \frac{a+b-\sqrt{(a-b)^2+c_-^2}}2.
 \label{eq:v93-quartic-eigenvalue}
\end{equation}

\begin{theorem}[Mixed-quartic coercivity]
\label{thm:v93-mixed-quartic}
If
\begin{equation}
 c>-2\sqrt{ab},
 \label{eq:v93-copositive-gate}
\end{equation}
then \(\lambda_*>0\) and
\begin{equation}
 P_4(s,h)
 \geq
 \lambda_*(s^4+h^4)
 \geq
 \frac{\lambda_*}{2}(s^2+h^2)^2.
 \label{eq:v93-quartic-lower}
\end{equation}
Conversely, if \(c<-2\sqrt{ab}\), then \(P_4\) is negative along a ray.
At \(c=-2\sqrt{ab}\), it vanishes along a nonzero ray.
\end{theorem}

\begin{proof}
Set \(x=s^2\geq0\) and \(y=h^2\geq0\).  Since
\(cxy\geq-c_-xy\),
\[
 P_4(s,h)
 \geq
 \begin{pmatrix}x&y\end{pmatrix}
 \begin{pmatrix}
  a&-c_-/2\\
  -c_-/2&b
 \end{pmatrix}
 \begin{pmatrix}x\\y\end{pmatrix}.
\]
Condition \eqref{eq:v93-copositive-gate} says
\(c_-<2\sqrt{ab}\), so the matrix is positive definite.  Its smaller
eigenvalue is \(\lambda_*\), giving the first inequality.  The second
uses \(x^2+y^2\geq(x+y)^2/2\).

If \(c<0\), choose \(x/y=\sqrt{b/a}\).  Then
\[
 ax^2+cxy+by^2
 =
 y^2\left(2b+c\sqrt{\frac ba}\right).
\]
This is negative when \(c<-2\sqrt{ab}\) and zero at equality.  For
\(c\geq0\), the converse cases do not arise.
\end{proof}

Thus the mixed coupling is controlled by a genuine two-variable
condition.  Positivity of \(a\) and \(b\) separately is insufficient if
the mixed coefficient is too negative.

\subsection{Absorbing lower-degree terms}

\begin{lemma}[Quartic absorption]
\label{lem:v93-quartic-absorption}
Let \(z=(s,h)\in{\mathbb R}^2\), and suppose
\[
 |P_{\leq3}(z)|\leq C_0(1+\|z\|^3).
\]
For every \(\eta>0\), there is \(C_\eta<\infty\) such that
\begin{equation}
 |P_{\leq3}(z)|
 \leq\eta\|z\|^4+C_\eta.
 \label{eq:v93-lower-absorption}
\end{equation}
Consequently, under Theorem \ref{thm:v93-mixed-quartic},
\begin{equation}
 \operatorname{Re}\{P_4(z)+P_{\leq3}(z)\}
 \geq
 \frac{\lambda_*}{4}\|z\|^4-C
 \label{eq:v93-full-polynomial-coercivity}
\end{equation}
for some finite \(C\).
\end{lemma}

\begin{proof}
The scalar function
\(C_0(1+r^3)-\eta r^4\) has a finite maximum on \(r\geq0\), proving
\eqref{eq:v93-lower-absorption}.  Apply it with
\(\eta=\lambda_*/4\) and use
\eqref{eq:v93-quartic-lower}.
\end{proof}

Any finite-dimensional negative quadratic form is included in
\(P_{\leq3}\), so a strictly positive quartic can dominate it at large
amplitude.  Uniformity in the cutoff requires the quartic constant and
the lower-degree coefficient stock to remain uniform per local block.

\subsection{Pure conformal no-go}

\begin{proposition}[Higgs quartic alone does not cure a pure conformal
runaway]
\label{prop:v93-pure-conformal}
Suppose the leading polynomial has \(a=0\),
\[
 P_4(s,h)=c s^2h^2+b h^4,
 \qquad b>0,
\]
and the conformal direction contains a negative quadratic
\(-\mu s^2\) with \(\mu>0\).  Then
\begin{equation}
 -\mu s^2+P_4(s,h)\longrightarrow-\infty
 \quad\text{along }h=0,\ |s|\to\infty.
 \label{eq:v93-conformal-runaway}
\end{equation}
No value of \(b\) or \(c\) changes this conclusion.
\end{proposition}

\begin{proof}
On the ray \(h=0\), both Higgs-containing quartic terms vanish, leaving
\(-\mu s^2\).
\end{proof}

Therefore the Standard Model Higgs quartic cannot by itself solve the
Euclidean Einstein conformal-factor problem.  A valid construction needs
a contour on which the conformal quadratic real part is nonnegative, a
genuine positive conformal higher-order term, a constraint reduction
that removes the runaway direction, or another proved mechanism.  The
mechanism must remain compatible with reflection and the Lorentzian
reconstruction.

\subsection{Uniform conditional block exterior}

Let \(x_y\in{\mathbb R}^{d_*}\) be one local fast block and let \(W\)
denote its frozen exterior data.

\begin{corollary}[Regular-collar weighted exterior]
\label{cor:v93-conditional-exterior}
Suppose that on a regular exterior set \({\cal R}\), uniformly in \(W\),
\begin{align}
 \operatorname{Re}S_y(x_y;W)
 &\geq a_*\|x_y\|^p-b_*,
 \label{eq:v93-uniform-block-lower}\\
 \operatorname{Re}S_y(x_y;W)
 &\leq B_*
 \quad(\|x_y\|\leq r_*),
 \label{eq:v93-uniform-block-upper}
\end{align}
with \(a_*,r_*>0\) and finite \(b_*,B_*\).  Then for every \(q\geq0\)
there are uniform constants \(C,c>0\) such that the normalized
conditional absolute density obeys
\begin{equation}
 {\mathbb E}_{y,W}\left[
  (1+\|x_y\|^q){\bf1}_{\{\|x_y\|>R\}}
 \right]
 \leq Ce^{-cR^p}
 \label{eq:v93-uniform-block-tail}
\end{equation}
for all sufficiently large \(R\) and all \(W\in{\cal R}\).
\end{corollary}

\begin{proof}
Apply Theorem \ref{thm:v93-coercive-tail} with the same constants for
every \(W\).  The partition lower bound is uniform because of
\eqref{eq:v93-uniform-block-upper}.
\end{proof}

With \(R_n=\varepsilon_n^{-\alpha/2}\) and \(p=4\), the right side is
\[
 C\exp\{-c\varepsilon_n^{-2\alpha}\},
\]
which is summable for a geometric penalty schedule.  This would fill the
weighted exterior row of V84 on regular collars.  Exterior data outside
\({\cal R}\) still need a saturated witness and unique-payment schedule.

\subsection{What the theorem requires from gravity}

For a local conformal--Higgs block, one must derive the coefficients
\(a,b,c\) from the actual gauge-fixed, contour-chosen regulated action.
The following are separate checks:
\begin{enumerate}
\item the conformal contour makes the complete derivative quadratic form
      have a controlled real part;
\item the leading local polynomial passes
      \eqref{eq:v93-copositive-gate} or another explicit coercivity
      inequality;
\item boundary couplings and neighboring retained fields fit the uniform
      lower-degree stock;
\item the local upper bound
      \eqref{eq:v93-uniform-block-upper} supplies a nonvanishing
      conditional normalization;
\item the contour deformation crosses no singular determinant or gauge
      orbit and preserves the required reflection relation; and
\item irregular exterior fields possess a joint action witness.
\end{enumerate}

\begin{firewall}[Finite polynomial test is not Einstein coercivity]
\label{fire:v93-not-Einstein-coercivity}
Theorem \ref{thm:v93-mixed-quartic} is an exact two-amplitude result.
The Einstein conformal mode is a field with derivative couplings,
constraints, gauge Jacobians, and boundary terms.  V93 neither derives a
positive conformal quartic nor proves that a contour rotation is
reflection-compatible.  Proposition \ref{prop:v93-pure-conformal}
instead shows that the Higgs quartic leaves the pure conformal ray
uncontrolled when no independent conformal mechanism is present.
\end{firewall}

\begin{assessment}[Exact status after V93]
\label{ass:v93-status}
V93 proves a normalized weighted-tail theorem from uniform coercivity and
a local partition lower bound, gives the exact mixed-quartic gate
\(c>-2\sqrt{ab}\), absorbs all lower polynomial degrees, and proves the
pure conformal no-go for Higgs-only stabilization.  These results specify
a route to the V84 exterior row on regular blocks but do not establish
the required gravitational contour or irregular-collar witness.
Determinant-line gluing, anomaly control, retained-sector tightness,
symmetry restoration, regulator independence, and the full
SM--Einstein continuum remain open.
\end{assessment}

\subsection{Next acquisition}

V94 will examine the quadratic conformal contour itself.  It will prove a
finite-dimensional contour-rotation criterion with source and
determinant control, state the reflection-pair condition, and identify
which nonlinear terms obstruct a uniform deformation from the original
real metric cycle.

\paragraph{Parent-version record.}
V93 was formed from
\texttt{Renewal\_SM\_Einstein\_Continuum\_Research\_Notes\_V92.tex}.
The V92 parent had \(32\,320\) counted source lines, \(1\,173\,601\) bytes,
and SHA--256
\[
 \texttt{0C0FC4BED0BA393EBDF031CF64A820EC60351D525853589A8A5FF12FB40C8793}.
\]

% =============================================================================
% END APPENDED TENTH-SERIES V93 MATERIAL
% =============================================================================
% =============================================================================
% BEGIN APPENDED TENTH-SERIES V94 MATERIAL
% =============================================================================

\section{Tenth-series V94: conformal contour rotation and its obstruction ledger}
\label{sec:v10v94}

V93 showed that Higgs growth does not control the pure conformal runaway.
This note treats the quadratic contour rotation exactly.  It proves the
real-part Schur criterion, the Gaussian determinant and source formula,
and a contour-deformation lemma.  The lemma exposes an essential point:
the imaginary conformal contour generally defines an analytic
continuation; it is not homologous to the divergent real conformal cycle
through a uniformly decaying deformation.

\subsection{Quadratic conformal block}

Let \(s\in{\mathbb R}^{p}\) denote regulated conformal coordinates and
\(x\in{\mathbb R}^{q}\) the remaining real bosonic coordinates.  Consider
the real quadratic form
\begin{equation}
 S_2(s,x)
 =
 -\frac12\langle s,Ks\rangle
 +\frac12\langle x,Ax\rangle
 +\langle s,Bx\rangle,
 \label{eq:v94-indefinite-quadratic}
\end{equation}
where \(K=K^*>0\), \(A=A^*>0\), and
\(B:{\mathbb R}^q\to{\mathbb R}^p\).  Rotate the conformal line by
\begin{equation}
 s=e^{i\vartheta}t,\qquad t\in{\mathbb R}^p.
 \label{eq:v94-rotation}
\end{equation}
On the rotated real coordinates \(y=(t,x)\),
\begin{equation}
 \operatorname{Re}S_2(e^{i\vartheta}t,x)
 =
 \frac12
 \begin{pmatrix}t\\x\end{pmatrix}^{\!*}
 Q_\vartheta
 \begin{pmatrix}t\\x\end{pmatrix},
 \label{eq:v94-real-quadratic}
\end{equation}
with
\begin{equation}
 Q_\vartheta
 =
 \begin{pmatrix}
  -\cos(2\vartheta)K&\cos\vartheta B\\
  \cos\vartheta B^*&A
 \end{pmatrix}.
 \label{eq:v94-real-part-matrix}
\end{equation}

\begin{theorem}[Quadratic contour Schur criterion]
\label{thm:v94-Schur-contour}
Suppose \(-\cos(2\vartheta)>0\), and put
\(K_\vartheta=-\cos(2\vartheta)K\).  Then
\(Q_\vartheta>0\) if and only if
\begin{equation}
 A-\cos^2\vartheta\,B^*K_\vartheta^{-1}B>0.
 \label{eq:v94-contour-Schur}
\end{equation}
At the imaginary contour \(\vartheta=\pi/2\),
\begin{equation}
 Q_{\pi/2}
 =
 \begin{pmatrix}K&0\\0&A\end{pmatrix}>0.
 \label{eq:v94-imaginary-decoupling}
\end{equation}
More quantitatively, if
\[
 K_\vartheta\geq k_\vartheta I,
 \qquad
 A-\cos^2\vartheta B^*K_\vartheta^{-1}B\geq a_\vartheta I,
\]
then
\begin{equation}
 \operatorname{Re}S_2(e^{i\vartheta}t,x)
 \geq
 \frac12
 \min\left\{
  \frac{k_\vartheta a_\vartheta}
       {a_\vartheta+2\cos^2\vartheta\|B\|^2/k_\vartheta},
  \frac{a_\vartheta}{2}
 \right\}
 (\|t\|^2+\|x\|^2).
 \label{eq:v94-quantitative-floor}
\end{equation}
\end{theorem}

\begin{proof}
Complete the square:
\begin{align*}
 &\langle t,K_\vartheta t\rangle
 +2\cos\vartheta\langle t,Bx\rangle+\langle x,Ax\rangle\\
 &\quad=
 \left\langle
  t+\cos\vartheta K_\vartheta^{-1}Bx,\,
  K_\vartheta
  \{t+\cos\vartheta K_\vartheta^{-1}Bx\}
 \right\rangle\\
 &\qquad+
 \left\langle x,
  \{A-\cos^2\vartheta B^*K_\vartheta^{-1}B\}x
 \right\rangle.
\end{align*}
This proves the equivalence and
\eqref{eq:v94-imaginary-decoupling}.  For the quantitative bound, set
\(u=t+\cos\vartheta K_\vartheta^{-1}Bx\).  Then
\[
 \|t\|^2
 \leq2\|u\|^2+
 \frac{2\cos^2\vartheta\|B\|^2}{k_\vartheta^2}\|x\|^2.
\]
Combine this with
\(k_\vartheta\|u\|^2+a_\vartheta\|x\|^2\) and decrease the coefficient
to the displayed positive constant.
\end{proof}

The exact imaginary line removes the real quadratic mixing because the
mixed term becomes purely imaginary.  Away from that line, positivity
requires the full Schur condition; positivity of \(K\) and \(A\)
separately is insufficient.

\subsection{Gaussian determinant and sources}

Write
\[
 R_\vartheta=
 \operatorname{diag}(e^{i\vartheta}I_p,I_q),
 \qquad
 H_\vartheta=R_\vartheta^TH_0R_\vartheta,
\]
where \(H_0\) is the symmetric matrix of
\eqref{eq:v94-indefinite-quadratic}.  For a source \(J\), set
\(J_\vartheta=R_\vartheta^TJ\).

\begin{proposition}[Rotated Gaussian formula]
\label{prop:v94-rotated-Gaussian}
If \(\operatorname{Re}H_\vartheta\geq\lambda I\) for some
\(\lambda>0\), then
\begin{align}
 &\int_{R_\vartheta{\mathbb R}^{p+q}}
 \exp\left\{-\frac12z^TH_0z+J^Tz\right\}\,dz
 \nonumber\\
 &\quad=
 e^{ip\vartheta}(2\pi)^{(p+q)/2}
 \det(H_\vartheta)^{-1/2}
 \exp\left\{
  \frac12J_\vartheta^TH_\vartheta^{-1}J_\vartheta
 \right\}.
 \label{eq:v94-rotated-Gaussian}
\end{align}
The square-root branch is the analytic branch on the convex domain
\(\{H=H^T:\operatorname{Re}H>0\}\) that is positive for real
\(H>0\).  Moreover,
\begin{equation}
 \left|
 e^{-\frac12y^TH_\vartheta y+J_\vartheta^Ty}
 \right|
 \leq
 \exp\left\{
  -\frac{\lambda}{4}\|y\|^2+
  \lambda^{-1}\|J_\vartheta\|^2
 \right\}.
 \label{eq:v94-source-domination}
\end{equation}
\end{proposition}

\begin{proof}
The substitution \(z=R_\vartheta y\) gives
\(dz=e^{ip\vartheta}dy\).  The domain
\(\operatorname{Re}H>0\) is convex and contains no singular matrices, so
the standard real positive Gaussian formula extends analytically there
with the stated determinant branch.  Finally,
\[
 \operatorname{Re}(J_\vartheta^Ty)
 \leq
 \|J_\vartheta\|\|y\|
 \leq
 \frac{\lambda}{4}\|y\|^2+
 \lambda^{-1}\|J_\vartheta\|^2,
\]
which proves \eqref{eq:v94-source-domination}.
\end{proof}

The phase \(e^{ip\vartheta}\) and the determinant branch are one owner.
Discarding either changes the measure.

\subsection{When two contour rays are homologous}

\begin{lemma}[Sector deformation with decay]
\label{lem:v94-sector-deformation}
Let \(F\) be entire in one complex variable.  Suppose
\(\vartheta_0<\vartheta_1\) and there are \(C,c,\rho>0\) such that
\begin{equation}
 |F(re^{i\phi})|
 \leq Ce^{-cr^\rho}
 \quad
 (r\geq R_0,\ \vartheta_0\leq\phi\leq\vartheta_1).
 \label{eq:v94-sector-decay}
\end{equation}
Then the oriented ray integrals converge absolutely and
\begin{equation}
 \int_{e^{i\vartheta_0}{\mathbb R}_+}F(z)\,dz
 =
 \int_{e^{i\vartheta_1}{\mathbb R}_+}F(z)\,dz.
 \label{eq:v94-ray-equality}
\end{equation}
The same statement holds for a full line after applying it to both rays.
A finite-dimensional product contour can be rotated coordinate by
coordinate if the bound is uniform in the remaining coordinates and
provides an integrable majorant.
\end{lemma}

\begin{proof}
Integrate \(F\) around the boundary of the sector truncated at radius
\(R\).  Cauchy's theorem makes the boundary integral zero.  The circular
arc has length \(R(\vartheta_1-\vartheta_0)\), while
\eqref{eq:v94-sector-decay} bounds its contribution by
\(CR(\vartheta_1-\vartheta_0)e^{-cR^\rho}\), which tends to zero.
Let \(R\to\infty\).  The full-line and product statements follow by
repeating the argument and using dominated Fubini.
\end{proof}

\begin{proposition}[The real conformal line is not deformable by this
criterion]
\label{prop:v94-real-line-obstruction}
For the one-dimensional wrong-sign Gaussian
\[
 e^{-S(s)}=e^{ks^2/2},
 \qquad k>0,
\]
the real-line integral diverges.  On the ray \(s=re^{i\phi}\), its
absolute value is
\[
 \exp\left\{\frac{k}{2}r^2\cos(2\phi)\right\}.
\]
Hence no sector containing \(\phi=0\) and \(\phi=\pi/2\) satisfies
\eqref{eq:v94-sector-decay}.
\end{proposition}

\begin{proof}
At \(\phi=0\), the integrand grows as \(e^{kr^2/2}\).  In every
neighbourhood of that ray, \(\cos(2\phi)>0\) for sufficiently small
\(\phi\), so uniform decay fails.  The real integral itself is infinite.
\end{proof}

Thus the imaginary conformal contour can define a convergent analytic
continuation, but it is not justified as a deformation of the divergent
real Euclidean conformal integral.  A physical derivation must specify
the original Lorentzian cycle, a lapse or constraint prescription, or a
Picard--Lefschetz construction selecting the thimble.

\subsection{Nonlinear phase table}

For a real coefficient and a monomial \(s^kh^\ell\), the imaginary
rotation gives
\begin{equation}
 \operatorname{Re}\{(it)^kh^\ell\}
 =
 \cos\left(\frac{k\pi}{2}\right)t^kh^\ell.
 \label{eq:v94-monomial-phase}
\end{equation}
Consequently:
\[
 \begin{array}{c|c}
 k\bmod4&\text{real-part factor}\\ \hline
 0&+1\\
 1&0\\
 2&-1\\
 3&0
 \end{array}
\]
In particular, a positive \(s^4\) remains positive, while a positive
\(s^2h^2\) becomes negative.  If the rotated quartic is
\[
 a t^4-c t^2h^2+b h^4,
\]
Theorem \ref{thm:v93-mixed-quartic} requires
\begin{equation}
 c<2\sqrt{ab}
 \label{eq:v94-rotated-quartic-gate}
\end{equation}
when \(c\geq0\).  The quadratic rotation can therefore create a quartic
mixed-direction obstruction.

Derivative interactions and field-dependent determinants obey the same
phase accounting term by term, but an infinite Taylor series needs a
uniform real-part bound on the full chart, not merely checks of finitely
many coefficients.

\subsection{Perturbations of the rotated Gaussian}

\begin{corollary}[Quadratically subordinate interaction]
\label{cor:v94-subordinate-interaction}
Suppose
\(\operatorname{Re}H_\vartheta\geq\lambda I\), and let
\(N_\vartheta(y)\) satisfy
\begin{equation}
 \operatorname{Re}N_\vartheta(y)
 \geq-\eta\|y\|^2-C_N,
 \qquad
 0\leq\eta<\lambda/2.
 \label{eq:v94-subordinate-bound}
\end{equation}
Then
\begin{equation}
 \operatorname{Re}\left\{
  \frac12y^TH_\vartheta y+N_\vartheta(y)
 \right\}
 \geq
 \left(\frac{\lambda}{2}-\eta\right)\|y\|^2-C_N.
 \label{eq:v94-full-contour-coercivity}
\end{equation}
The absolute integral and all polynomial moments are finite.  Uniform
constants give the weighted exterior bound of V93 with \(p=2\).
\end{corollary}

\begin{proof}
Take real parts and use the spectral lower bound.  Theorem
\ref{thm:v93-coercive-tail} applies after a local upper bound is supplied.
\end{proof}

If \(N_\vartheta\) has a negative quartic direction, it is not
quadratically subordinate.  The mixed-quartic gate of V93 must then be
checked before this corollary can be used.

\subsection{Reflection-pair condition}

Let \(\Theta\) be the linear Euclidean-time reflection on regulated real
field coordinates and define the antilinear involution
\[
 {\mathscr R}z:=\overline{\Theta z}.
\]
A contour \(\Gamma\) used for a reflection form must at least satisfy
\begin{equation}
 {\mathscr R}\Gamma=\Gamma
 \label{eq:v94-reflection-contour}
\end{equation}
with compatible orientation and density.  For a single
\(\Theta\)-invariant real component, the line
\(e^{i\vartheta}{\mathbb R}\) is invariant as an unoriented set exactly
when
\begin{equation}
 \vartheta=0\quad\text{or}\quad\vartheta=\pi/2
 \pmod{\pi}.
 \label{eq:v94-reflection-lines}
\end{equation}

\begin{proof}
Complex conjugation maps
\(e^{i\vartheta}{\mathbb R}\) to
\(e^{-i\vartheta}{\mathbb R}\).  The two real lines coincide exactly
when \(e^{2i\vartheta}\) is real, which gives
\eqref{eq:v94-reflection-lines}.
\end{proof}

Contour invariance is necessary but not sufficient for
Osterwalder--Schrader positivity.  If \(s=it\), then its raw covariance
is the negative of the \(t\)-covariance.  Positivity can hold only after
the conformal variable is removed from the physical reflection algebra,
paired with constraints or ghosts by a proved quotient, or assigned a
modified involution justified by the physical reconstruction.

\subsection{Singularities and determinant lines}

The sector lemma assumes an entire integrand.  Gauge fixing and fermion
integration can introduce:
\begin{enumerate}
\item zeros of Faddeev--Popov or Dirac determinants;
\item poles from a chosen inverse or gauge chart;
\item logarithmic branch cuts used to express a determinant;
\item changes of orientation across Gribov or constraint-rank walls; and
\item nontrivial transition functions of the chiral determinant line.
\end{enumerate}
A deformation crossing any such locus requires residue, spectral-flow,
or line-bundle data.  A local analytic logarithm does not establish a
global contour equivalence.

\begin{definition}[V94 conformal-contour card]
\label{def:v94-contour-card}
A finite regulator passes the contour card if:
\begin{enumerate}
\item the conformal and physical quadratic blocks satisfy
      \eqref{eq:v94-contour-Schur} with a uniform floor;
\item all nonlinear terms satisfy a global real-part coercivity bound;
\item sources obey the uniform domination
      \eqref{eq:v94-source-domination};
\item the Jacobian phase and Gaussian determinant branch are owned once;
\item the selected cycle has a derivation from a convergent original
      cycle or an explicit analytic-continuation prescription;
\item the cycle is invariant under the physical reflection involution and
      the physical reflection algebra is identified;
\item no unaccounted gauge or determinant singularity is crossed; and
\item the bounds are uniform in volume, cutoff, retained data, and
      regulator removal.
\end{enumerate}
\end{definition}

\begin{firewall}[Convergence is not contour equivalence or OS positivity]
\label{fire:v94-contour-limit}
The imaginary conformal contour makes
\eqref{eq:v94-indefinite-quadratic} Gaussian-coercive, but the original
real Euclidean conformal integral diverges and cannot be connected to it
by Lemma \ref{lem:v94-sector-deformation}.  The rotation also changes the
real signs of nonlinear even-conformal terms.  Finally, invariance under
the reflection involution does not prove reflection positivity of the
physical quotient.
\end{firewall}

\begin{assessment}[Exact status after V94]
\label{ass:v94-status}
V94 proves the quadratic contour Schur criterion, determinant and source
formula, a sector deformation theorem, the obstruction to deforming the
divergent real conformal Gaussian, the nonlinear phase table, and the
reflection-line condition.  It establishes a conditional coercive
imaginary cycle but not its Lorentzian derivation, nonlinear global
coercivity, determinant-line compatibility, or OS positivity.  The full
SM--Einstein continuum remains open.
\end{assessment}

\subsection{Next acquisition}

V95 is the next compulsory companion audit.  It will inspect the newest
Yang--Mills and Navier--Stokes notes for a triangular exterior schedule,
measured-coordinate refinement, contour argument, or finite-jet
contraction relevant to the V93--V94 bottleneck.  V96 will continue with
the first transferable acquisition.

\paragraph{Parent-version record.}
V94 was formed from
\texttt{Renewal\_SM\_Einstein\_Continuum\_Research\_Notes\_V93.tex}.
The V93 parent had \(32\,649\) counted source lines, \(1\,184\,760\) bytes,
and SHA--256
\[
 \texttt{8F3CF81CD810677A5E6EAFD5DC9C05DE96149055E22096F9369976A946F4584B}.
\]

% =============================================================================
% END APPENDED TENTH-SERIES V94 MATERIAL
% =============================================================================
% =============================================================================
% BEGIN APPENDED TENTH-SERIES V95 MATERIAL
% =============================================================================

\section{Tenth-series V95: companion audit and retained topological sectors}
\label{sec:v10v95}

\subsection{Compulsory companion audit}

The newest active companion notes at this gate were inspected read-only:
\begin{align*}
 &\texttt{Renewal\_Yang\_Mills\_Mass\_Gap\_Research\_Notes\_V13.tex},
 \\
 &\qquad
 \operatorname{SHA256}
 =
 \texttt{C550C93F0488FBB5D364F9E512245EB3941AAB59130F94B4FB591FAF40A4EF06},
 \\
 &\texttt{Renewal\_NS\_Stability\_Research\_Notes\_V26.tex},
 \\
 &\qquad
 \operatorname{SHA256}
 =
 \texttt{82C887F8C2C69E4F28FAD7C3DC8DE5B9099CB5A4BF431EBA1FFF3E91935978AD}.
\end{align*}
The Yang--Mills companion advanced from V10 to V13 after the V90 audit.
The Navier--Stokes companion is unchanged.

\subsection{Sectorwise transport}

The Yang--Mills V11 note proves that continuous local integration on a
regular connected fibre cannot change an integer topological label.
Flux, toron, and integer-charge data are retained, and the exact
triangular pushforward is diagonal in those labels.  A theta factor
\(e^{i\theta n}\) therefore remains attached to the same integer sector;
it is not a local running vertex generated by conditional integration.

The SM--Einstein extension must retain a larger discrete ledger,
potentially including:
\begin{enumerate}
\item \(SU(3)\), \(SU(2)\), and Abelian flux or bundle data;
\item spin and spin-\(c\) structures;
\item gravitational Pontryagin, signature, or boundary eta data where
      defined;
\item global gauge components and boundary charges; and
\item determinant-line transition labels.
\end{enumerate}
Which of these occur depends on the fixed manifold, boundary conditions,
and regulator.  V96 will prove the abstract sectorwise projective
statement without asserting that every item in this list is present.

Transitions between sectors require a singular, nonadmissible, or
topology-changing configuration.  They cannot be hidden inside a regular
local polymer.  If the regulator allows them, their weights need a
separate bad-component and tightness estimate.

\subsection{Reflection-positive phases}

The V11 half-space phase identity is also transferable.  If a
reflection-positive functional has a real positive-time observable
\(q_+\), then
\[
 e^{i\theta(q_+-\Theta q_+)}
 =
 h_\theta\,\Theta h_\theta,
 \qquad
 h_\theta=e^{i\theta q_+}.
\]
Multiplication by this factor preserves the reflection-positive cone.
Therefore reflection positivity alone does not force a theta coordinate
to vanish.  Zero theta is a choice of physical section, though it is
exactly preserved by a sector-diagonal pushforward.

This refines V94.  Invariance of a complex contour under the antilinear
reflection is only a necessary geometric condition.  A phase produced by
the contour Jacobian, a topological term, or a determinant must also
factor through the reflection-positive cone, including any term supported
on the reflection plane.  Reflection invariance of a plane term is not
enough.

\subsection{Marginal trajectory input}

The Yang--Mills V12--V13 notes organize a calibrated product domain and
use the reciprocal coupling to shield a marginal response estimate.  The
method separates:
\begin{enumerate}
\item the physical marginal coefficient;
\item measured redundant directions;
\item irrelevant anisotropy and stable packet owners;
\item theta and integer labels; and
\item the remaining scalar reset mismatch.
\end{enumerate}
Its asymptotic trajectory theorem is conditional on a model-specific
scalar inverse-coupling increment.  The reciprocal-variable algebra is
relevant to asymptotically free \(SU(3)\) and possibly \(SU(2)\) rows,
but its numerical exponent and determinant coefficient are not imported
into the coupled SM--Einstein flow.  Gravity, Abelian, Yukawa, Higgs, and
cosmological coordinates form a genuinely coupled marginal or relevant
system.

The useful transfer at V95 is the dependency order: remove exact
topological and measured-coordinate directions before estimating
marginal leakage from stable graphs.  A multi-coupling trajectory theorem
can be attempted only after the sectorwise continuum measure has been
defined.

\subsection{No Navier--Stokes transfer}

The Navier--Stokes V26 rank-one kernel calculation has no discrete bundle
sector, determinant line, topological phase, or quantum reflection form.
It supplies no new input to the V94 contour or sector problem.

\subsection{Adjusted route}

The next steps are:
\begin{enumerate}
\item construct the projective limit sector by sector;
\item impose tightness in the discrete label as well as in the field
      distribution topology;
\item sum sectors only under an absolute-mass bound;
\item prove reflection-positive factorization of every phase, including
      reflection-plane contributions;
\item retain zero theta as an exact section when that theory is chosen;
      and
\item return afterward to the contour and weighted exterior within each
      fixed sector.
\end{enumerate}

\begin{definition}[V95 topological-sector card]
\label{def:v95-sector-card}
A regulator tower passes the sector card if:
\begin{enumerate}
\item every admissible configuration has a declared discrete sector label;
\item regular local fibres preserve that label;
\item coarse maps define compatible maps of label sets;
\item transitions or ambiguous labels belong to a separately weighted bad
      set;
\item the family is tight in both field variables and sector labels;
\item the absolute sum of sector masses is finite;
\item every phase has a proved half-space factorization, with a separate
      reflection-plane check; and
\item determinant-line transitions agree with the sector maps.
\end{enumerate}
\end{definition}

\begin{firewall}[No topological summation without label tightness]
\label{fire:v95-label-tightness}
Tightness within each fixed sector does not control escape to sectors of
unbounded charge or changing spin, flux, or boundary data.  Nor does
sectorwise reflection positivity imply positivity after an arbitrary
complex phase sum.  The discrete tail and phase factorization are
independent hypotheses.
\end{firewall}

\begin{assessment}[Status after the V95 audit]
\label{ass:v95-status}
The audit imports sector-diagonal triangular transport, the need for
discrete-label tightness, and the half-space phase criterion.  It rejects
reflection positivity as a reason to set theta to zero and does not
import the Yang--Mills scalar coupling exponent into the coupled
SM--Einstein trajectory.  V96 will prove the sectorwise projective and
reflection-positive reconstruction theorem.  The full continuum remains
open.
\end{assessment}

\paragraph{Parent-version record.}
V95 was formed from
\texttt{Renewal\_SM\_Einstein\_Continuum\_Research\_Notes\_V94.tex}.
The V94 parent had \(33\,093\) counted source lines, \(1\,199\,470\) bytes,
and SHA--256
\[
 \texttt{F9DE60334A72833F16FC236D45441D215AD5509C227EBA880FD6D6372F59D1AB}.
\]

% =============================================================================
% END APPENDED TENTH-SERIES V95 MATERIAL
% =============================================================================
% =============================================================================
% BEGIN APPENDED TENTH-SERIES V96 MATERIAL
% =============================================================================

\section{Tenth-series V96: sectorwise projective reconstruction and topological phases}
\label{sec:v10v96}

V95 requires topology to be part of the retained projective state.  This
note proves a cylinder-local repair theorem with a countable discrete
sector label, a tight Radon realization on the product of distribution
and sector spaces, and the exact bounded theta-phase limit.  Reflection
positivity requires a half-space factorization and is not implied by
assigning constant phases to sectors.

\subsection{Regular fibres preserve integer labels}

\begin{lemma}[Connected-fibre constancy]
\label{lem:v96-connected-label}
Let \(Y\) be connected and let \(Q:Y\to{\mathbb Z}^r\) be continuous.
Then \(Q\) is constant.
\end{lemma}

\begin{proof}
The continuous image of a connected set is connected, while every
connected subset of the discrete space \({\mathbb Z}^r\) is a singleton.
\end{proof}

Thus any continuous local integration over a regular connected fibre is
diagonal in an integer topological charge.  A change of charge requires
leaving that fibre, crossing a singular set, or changing the underlying
bundle or boundary data.  Such configurations belong to a separate
bad-component owner.

\subsection{Local repair in the sector sum}

Let \(\Sigma\) be a countable discrete sector space.  At scale \(n\), let
\(\mu_{n,\sigma}\) be finite positive measures on \(X_n\), with
\begin{equation}
 \sum_{\sigma\in\Sigma}\mu_{n,\sigma}(1)=1.
 \label{eq:v96-total-sector-mass}
\end{equation}
The associated probability on \(X_n\times\Sigma\) is
\[
 \mu_n^\Sigma(dx,d\sigma)
 =
 \sum_{\sigma\in\Sigma}\mu_{n,\sigma}(dx)\delta_\sigma(d\sigma).
\]
Assume the coarse maps preserve the label:
\[
 p_{m,n}^\Sigma(x,\sigma)=(p_{m,n}x,\sigma).
\]
A nontrivial compatible map between regulator-dependent label sets can be
included by replacing \(\Sigma\) with their countable inverse-limit
label.

For a fixed coarse scale \(k\) and finite cell set \(A\), define
\[
 \nu_{n;k,A,\sigma}
 :=
 (r_{k,A}\circ p_{n,k})_*\mu_{n,\sigma}.
\]

\begin{theorem}[Sector-summed cylinder repair]
\label{thm:v96-sector-repair}
Suppose that for every \((k,A)\) there is a summable sequence
\(e_n^\Sigma(k,A)\) such that
\begin{equation}
 \sum_{\sigma\in\Sigma}
 \|\nu_{n+1;k,A,\sigma}-\nu_{n;k,A,\sigma}\|_{\rm TV}
 \leq e_n^\Sigma(k,A).
 \label{eq:v96-sector-adjacent}
\end{equation}
Then the joint local probabilities
\[
 \nu_{n;k,A}^\Sigma
 :=
 \sum_{\sigma\in\Sigma}
 \nu_{n;k,A,\sigma}\otimes\delta_\sigma
\]
converge in total variation on \(X_{k,A}\times\Sigma\) to a probability
\(\bar\mu_{k,A}^\Sigma\), and
\begin{equation}
 \|\nu_{N;k,A}^\Sigma-\nu_{n;k,A}^\Sigma\|_{\rm TV}
 \leq
 \sum_{j=n}^{N-1}e_j^\Sigma(k,A).
 \label{eq:v96-sector-Cauchy}
\end{equation}
The limits are compatible under local restriction, coarse graining, and
the sector projection.  They define a unique probability on the joint
cylinder sigma-algebra.
\end{theorem}

\begin{proof}
Measures supported on distinct sector labels are mutually singular, so
\[
 \left\|
 \sum_\sigma\eta_\sigma\otimes\delta_\sigma
 \right\|_{\rm TV}
 =
 \sum_\sigma\|\eta_\sigma\|_{\rm TV}.
\]
Thus \eqref{eq:v96-sector-adjacent} is exactly the adjacent
total-variation estimate for the joint local probability.  Telescope as
in Theorem \ref{thm:v83-cylinder-repair} to obtain
\eqref{eq:v96-sector-Cauchy}.  Completeness, positivity, and preservation
of total mass give the limit probability.  Pushforward contracts total
variation, so all finite-scale compatibility identities pass to the
limit.
\end{proof}

If sector-changing bad configurations have total local variation
\(b_n(k,A)\), they can be included on the right side of
\eqref{eq:v96-sector-adjacent}.  The theorem applies when their sum is
finite.  Relabelling them as a local vertex without this debit is not
valid.

\subsection{A practical label-tightness criterion}

Sectorwise convergence for each fixed \(\sigma\) is weaker than
\eqref{eq:v96-sector-adjacent}; mass can escape to \(|\sigma|\to\infty\).

\begin{lemma}[Envelope criterion]
\label{lem:v96-sector-envelope}
Suppose there are numbers \(M_\sigma\geq0\) such that
\[
 \sup_n\mu_{n,\sigma}(1)\leq M_\sigma,
 \qquad
 \sum_{\sigma\in\Sigma}M_\sigma<\infty.
\]
Then the labels are uniformly tight:
for every \(\epsilon>0\), there is a finite
\(F\subset\Sigma\) such that
\begin{equation}
 \sup_n\sum_{\sigma\notin F}\mu_{n,\sigma}(1)<\epsilon.
 \label{eq:v96-label-tightness}
\end{equation}
\end{lemma}

\begin{proof}
Choose \(F\) so that
\(\sum_{\sigma\notin F}M_\sigma<\epsilon\), and use the envelope.
\end{proof}

An exponential charge action is one possible source of such an envelope,
but no such SM--Einstein estimate is asserted here.

\subsection{Radon realization with discrete labels}

Let \({\cal S}\) be the Polish distribution space and give
\(\Sigma\) the discrete topology.  Let
\(\iota_n:X_n\to{\cal S}\) be the reconstruction maps.

\begin{theorem}[Joint field-sector tightness]
\label{thm:v96-joint-tightness}
Suppose:
\begin{enumerate}
\item the joint cylinder repair of Theorem
      \ref{thm:v96-sector-repair} holds;
\item for every \(\epsilon>0\), there is a finite
      \(F\subset\Sigma\) satisfying
      \eqref{eq:v96-label-tightness}; and
\item for that finite \(F\), there are compact sets
      \(K_\sigma\subset{\cal S}\) such that
\begin{equation}
 \sup_n
 (\iota_n)_*\mu_{n,\sigma}(K_\sigma^c)
 <
 \frac{\epsilon}{\max\{1,|F|\}}
 \quad(\sigma\in F).
 \label{eq:v96-sector-field-tightness}
\end{equation}
\end{enumerate}
Then the reconstructed joint probabilities on
\({\cal S}\times\Sigma\) are tight.  If the joint cylinder projections
separate probabilities, they converge weakly to a unique Radon
probability \(P\) having the repaired local marginals.
\end{theorem}

\begin{proof}
The set
\[
 K:=\bigcup_{\sigma\in F}K_\sigma\times\{\sigma\}
\]
is compact because it is a finite union of compact sets.  The probability
of its complement is bounded by the label tail plus the sum of the field
tails in \eqref{eq:v96-sector-field-tightness}, which can be made
arbitrarily small after rescaling \(\epsilon\).  Prokhorov's theorem gives
subsequential Radon limits.  The repaired separating cylinder marginals
identify every subsequential limit, exactly as in Theorem
\ref{thm:v83-tight-realization}; hence the full sequence converges.
\end{proof}

Tightness in a fixed sector and tightness of the sector label are separate
rows.  Neither implies the other.

\subsection{Topological phase in the continuum}

Let \(Q:\Sigma\to{\mathbb Z}^r\) be a charge map and
\(\vartheta\in{\mathbb R}^r\).  Define the unit-modulus character
\begin{equation}
 \chi_\vartheta(\sigma)
 :=
 e^{i\vartheta\cdot Q(\sigma)}.
 \label{eq:v96-sector-character}
\end{equation}

\begin{theorem}[Bounded phase reconstruction]
\label{thm:v96-phase-reconstruction}
Let \(P_n\) converge weakly to \(P\) on
\({\cal S}\times\Sigma\) as in Theorem
\ref{thm:v96-joint-tightness}.  Define complex measures
\[
 d\nu_{n,\vartheta}
 :=
 \chi_\vartheta(\sigma)\,dP_n,
 \qquad
 d\nu_\vartheta
 :=
 \chi_\vartheta(\sigma)\,dP.
\]
Then:
\begin{enumerate}
\item \(\|\nu_{n,\vartheta}\|_{\rm TV}
      =\|\nu_\vartheta\|_{\rm TV}=1\);
\item \(\nu_{n,\vartheta}\) converges weakly to
      \(\nu_\vartheta\);
\item the sector coefficient \(\chi_\vartheta(\sigma)\) is unchanged by
      every sector-preserving coarse pushforward; and
\item if
\begin{equation}
 Z_\vartheta:=\nu_\vartheta(1)\ne0,
 \label{eq:v96-nonzero-partition}
\end{equation}
then \(Z_\vartheta^{-1}\nu_\vartheta\) is a normalized complex
continuum functional.
\end{enumerate}
At \(\vartheta=0\), \(Z_0=1\) and the limit is the positive probability
\(P\).
\end{theorem}

\begin{proof}
Multiplication by a unit-modulus function preserves total variation.
Because \(\Sigma\) is discrete,
\(\chi_\vartheta\) is bounded and continuous, so weak convergence passes
after multiplication.  Sector preservation leaves \(Q(\sigma)\), hence
its character, unchanged.  The last statements are immediate.
\end{proof}

The nonzero-partition condition is essential.  A phase sum may cancel
even when every sector measure is positive.

\subsection{Half-space phase factorization}

Let \({\cal A}\) be a commutative Euclidean field algebra with an
antilinear reflection involution \(\Theta\), and let
\({\cal A}_+\) be its positive-time subalgebra.  A functional \(L_0\) is
reflection positive when
\[
 L_0((\Theta F)F)\geq0
 \quad(F\in{\cal A}_+).
\]

\begin{theorem}[Half-space phase lemma]
\label{thm:v96-half-space-phase}
Let \(q_+\in{\cal A}_+\) be real and
\[
 h_\vartheta=e^{i\vartheta q_+}.
\]
Define
\begin{equation}
 w_\vartheta
 =
 h_\vartheta\,\Theta h_\vartheta
 =
 e^{i\vartheta(q_+-\Theta q_+)}.
 \label{eq:v96-half-space-factor}
\end{equation}
If \(L_0\) is reflection positive, then
\[
 L_\vartheta(F):=L_0(Fw_\vartheta)
\]
is reflection positive.  Moreover,
\(L_\vartheta(1)\geq0\); if it is positive, normalization preserves
reflection positivity.
\end{theorem}

\begin{proof}
Antilinearity gives
\(\Theta h_\vartheta=e^{-i\vartheta\Theta q_+}\).  For
\(F\in{\cal A}_+\), commutativity and
\eqref{eq:v96-half-space-factor} give
\[
 L_\vartheta((\Theta F)F)
 =
 L_0(\Theta(Fh_\vartheta)(Fh_\vartheta))
 \geq0.
\]
Take \(F=1\) for the partition value, and divide by it when positive.
\end{proof}

This proves that reflection positivity does not select
\(\vartheta=0\).  Zero theta is a physical choice of section, and
Theorem \ref{thm:v96-phase-reconstruction} shows that a
sector-preserving RG map preserves it exactly.

\subsection{Reflection-plane terms}

Let
\begin{equation}
 {\cal C}_{\rm RP}
 :=
 \overline{
  \left\{
   \sum_{j=1}^N(\Theta b_j)b_j:
   b_j\in{\cal A}_+,\ N<\infty
  \right\}}
 \label{eq:v96-RP-cone}
\end{equation}
in a topology in which \(L_0\) is continuous.

\begin{proposition}[Plane-factor criterion]
\label{prop:v96-plane-factor}
Suppose
\[
 Q=q_+-\Theta q_++q_0.
\]
The phase \(e^{i\vartheta Q}\) preserves reflection positivity if
\begin{equation}
 e^{i\vartheta q_0}\in{\cal C}_{\rm RP}.
 \label{eq:v96-plane-cone}
\end{equation}
Reflection invariance of \(q_0\) alone does not imply
\eqref{eq:v96-plane-cone}.
\end{proposition}

\begin{proof}
Approximate the plane factor by finite sums in
\eqref{eq:v96-RP-cone}.  For one such sum and
\(F\in{\cal A}_+\), the reflection form is a sum of nonnegative terms
\[
 L_0\left(
  \Theta(Fh_\vartheta b_j)(Fh_\vartheta b_j)
 \right).
\]
Pass to the closure.  The final statement follows because the fixed
subspace of \(\Theta\) is larger than the positive cone
\({\cal C}_{\rm RP}\).
\end{proof}

The conformal contour Jacobian and fermionic phase in V94 must pass this
factor test if they meet the reflection plane.

\subsection{Sectorwise identities and determinant lines}

Ward identities and constraints that are valid in every sector and are
uniformly integrable pass to \(P\) sectorwise and then to the zero-theta
mixture.  For a nonzero theta phase, the same algebraic identity passes
to the complex functional because the character is bounded.

A chiral determinant is not merely a scalar function if its determinant
line is nontrivial.  The sector sum
\[
 \sum_\sigma\chi_\vartheta(\sigma)\det D_\sigma\,d\mu_\sigma
\]
is globally defined as a scalar measure only after compatible frames and
transition functions have been chosen.  A phase assigned independently
inside each chart can violate projective compatibility on overlaps.

\begin{firewall}[Sector transport does not trivialize the determinant
line]
\label{fire:v96-no-line-trivialization}
V96 proves projective transport for scalar positive sector measures and
bounded scalar characters.  It does not prove that the Standard Model
chiral determinant is a globally compatible scalar over gauge,
gravitational, boundary, and spin-structure sectors.  Local anomaly
cancellation is necessary but does not by itself remove every global or
boundary cocycle.
\end{firewall}

\begin{definition}[V96 sector-continuum card]
\label{def:v96-sector-continuum-card}
The sector construction passes if:
\begin{enumerate}
\item regular fibres preserve the discrete labels;
\item sector-changing bad mass is summable in each cylinder seminorm;
\item the joint sector error satisfies
      \eqref{eq:v96-sector-adjacent};
\item field and discrete-label tightness satisfy Theorem
      \ref{thm:v96-joint-tightness};
\item topological phases are bounded characters of compatible labels;
\item every reflection phase has the half-space and plane-factor
      decomposition;
\item the chosen normalized phase partition is nonzero; and
\item determinant-line frames satisfy the same coarse and overlap
      cocycles.
\end{enumerate}
\end{definition}

\begin{assessment}[Exact status after V96]
\label{ass:v96-status}
V96 proves sector-summed cylinder repair, a joint field-sector Radon
limit, bounded theta-phase reconstruction, exact preservation of the
zero-theta section, the half-space phase lemma, and the reflection-plane
cone criterion.  It leaves the actual sector-tail estimate, singular
transition weight, determinant-line trivialization, gravitational
contour, weighted exterior, and retained-field tightness unproved.  The
full SM--Einstein continuum remains open.
\end{assessment}

\subsection{Next acquisition}

V97 will formulate the determinant-line projective cocycle.  It will
prove a finite-cover gluing theorem, identify local and global anomaly
obstructions as curvature and holonomy of the line, and state the exact
compatibility required for the fermionic coefficients in V79 and the
sector sum above.

\paragraph{Parent-version record.}
V96 was formed from
\texttt{Renewal\_SM\_Einstein\_Continuum\_Research\_Notes\_V95.tex}.
The V95 parent had \(33\,270\) counted source lines, \(1\,206\,780\) bytes,
and SHA--256
\[
 \texttt{F38AE6C577B8A31B907AB1611173B5E1D024301E0415C370E1731696FBBA43BD}.
\]

% =============================================================================
% END APPENDED TENTH-SERIES V96 MATERIAL
% =============================================================================
% =============================================================================
% BEGIN APPENDED TENTH-SERIES V97 MATERIAL
% =============================================================================

\section{Tenth-series V97: determinant-line cocycle and projective gluing}
\label{sec:v10v97}

V79 and V96 require fermionic coefficients to be scalar in one globally
compatible frame.  A chiral determinant is naturally a section of a
complex line bundle, so local determinant functions cannot simply be
identified across gauge, metric, boundary, and regulator charts.  This
note proves the exact finite-cover and projective gluing criteria and
separates local curvature from global holonomy.

\subsection{Finite-cover cocycle}

Let \(B\) be the regulated bosonic parameter space after the declared
gauge, boundary, and topological sector choices.  Let
\(\{U_i\}_{i=1}^N\) be a finite open cover.  A complex line bundle \(L\)
is described by nonvanishing transition functions
\begin{equation}
 g_{ij}:U_i\cap U_j\longrightarrow{\mathbb C}^{\times}
 \label{eq:v97-transition}
\end{equation}
satisfying
\begin{equation}
 g_{ii}=1,\qquad
 g_{ij}g_{ji}=1,\qquad
 g_{ij}g_{jk}g_{ki}=1
 \label{eq:v97-Cech-cocycle}
\end{equation}
on the corresponding overlaps.  A local determinant section is a family
\(D_i:U_i\to{\mathbb C}\) satisfying
\begin{equation}
 D_i=g_{ij}D_j.
 \label{eq:v97-local-determinant}
\end{equation}

\begin{theorem}[Scalarization criterion]
\label{thm:v97-scalarization}
The line bundle admits a global nowhere-vanishing frame if and only if
there are nonvanishing functions \(h_i:U_i\to{\mathbb C}^{\times}\) such
that
\begin{equation}
 g_{ij}=h_i h_j^{-1}.
 \label{eq:v97-coboundary}
\end{equation}
In that case
\begin{equation}
 d|_{U_i}:=h_i^{-1}D_i
 \label{eq:v97-global-scalar}
\end{equation}
defines one global scalar coefficient.  Conversely, any choice of global
frame gives functions \(h_i\) satisfying
\eqref{eq:v97-coboundary}.
\end{theorem}

\begin{proof}
If \eqref{eq:v97-coboundary} holds, then on an overlap
\[
 h_i^{-1}D_i
 =
 h_i^{-1}g_{ij}D_j
 =
 h_j^{-1}D_j,
\]
so the local functions glue.  The local frame
\(h_i^{-1}e_i\), where \(e_i\) is the defining frame on \(U_i\), also
glues and never vanishes.

Conversely, write a global nowhere-zero frame \(e\) as
\(e=h_i^{-1}e_i\) on \(U_i\).  Equality on overlaps gives
\(g_{ij}=h_i h_j^{-1}\).
\end{proof}

A partition of unity constructs a Hermitian metric on \(L\), but it does
not produce the functions \(h_i\).  Metric existence is not
trivialization.

The determinant section \(D\) itself may vanish.  A vanishing determinant
is not a frame and does not prove that \(L\) is trivial.  Anomaly
cancellation concerns the combined determinant line, not the absence of
zeros of one Dirac operator.

\subsection{Connection, curvature, and global holonomy}

Let \(\nabla\) be a connection on \(L\).  In a local frame, write
\[
 \nabla=d+A_i.
\]
On overlaps,
\begin{equation}
 A_i=A_j-g_{ij}^{-1}dg_{ij},
 \label{eq:v97-connection-change}
\end{equation}
and the curvature \(F_\nabla=dA_i\) is global.

\begin{theorem}[Flat-holonomy criterion]
\label{thm:v97-flat-holonomy}
Let \(B\) be connected and locally path connected.  A flat line bundle
with connection \((L,\nabla)\) admits a global nonzero parallel frame if
and only if its holonomy around every loop is the identity.
\end{theorem}

\begin{proof}
A global parallel frame returns to itself around every loop, so the
holonomy is trivial.  Conversely, choose a base point \(b_0\) and a
nonzero vector \(\ell_0\in L_{b_0}\).  For \(b\in B\), parallel transport
\(\ell_0\) along any path from \(b_0\) to \(b\).  Flatness makes transport
invariant under endpoint-fixed homotopy, and trivial loop holonomy makes
the result independent of the chosen path.  The resulting section is
nonzero, parallel, and continuous; standard local parallel frames make it
smooth.
\end{proof}

Thus two distinct obstructions remain:
\begin{enumerate}
\item nonzero curvature, detected infinitesimally and associated with a
      local anomaly; and
\item nontrivial flat holonomy, invisible to local curvature and
      associated with a global anomaly.
\end{enumerate}
Vanishing of the first does not imply vanishing of the second.

For a gauge group action, the required frame is equivariant.  Even a
topologically trivial line can have a nontrivial action on its fibre.
That equivariant character is another form of the same holonomy
obstruction on the gauge-orbit groupoid.

\subsection{Regulator-scale line maps}

At regulator scale \(n\), let \(B_n\) be the bosonic parameter space and
\(L_n\to B_n\) the combined fermionic determinant line.  Let
\[
 p_{n+1,n}:B_{n+1}\longrightarrow B_n
\]
be the coarse map.  A determinant-line projective structure is a family
of line-bundle isomorphisms
\begin{equation}
 \Phi_{n+1,n}:
 L_{n+1}\longrightarrow p_{n+1,n}^*L_n
 \label{eq:v97-scale-line-map}
\end{equation}
covering the identity on \(B_{n+1}\).  For three scales it must obey
\begin{equation}
 \Phi_{n+2,n}
 =
 p_{n+2,n+1}^*\Phi_{n+1,n}
 \circ\Phi_{n+2,n+1},
 \label{eq:v97-scale-coherence}
\end{equation}
after the canonical identification of iterated pullbacks.

\begin{theorem}[Coherent projective scalarization]
\label{thm:v97-projective-scalarization}
Suppose:
\begin{enumerate}
\item the line maps \(\Phi_{n+1,n}\) are isomorphisms satisfying
      \eqref{eq:v97-scale-coherence};
\item \(L_0\) has a global frame \(\tau_0\); and
\item the scale maps identify the complete combined determinant lines,
      with no omitted ultraviolet factor.
\end{enumerate}
Then there is a unique sequence of frames \(\tau_n\), once \(\tau_0\) is
fixed, satisfying
\begin{equation}
 \tau_{n+1}
 =
 p_{n+1,n}^*\tau_n\circ\Phi_{n+1,n}.
 \label{eq:v97-compatible-frame}
\end{equation}
In these frames, every compatible determinant section
\(D_n\in\Gamma(L_n)\) becomes a scalar \(d_n=\tau_n(D_n)\), and exact
section compatibility
\begin{equation}
 \Phi_{n+1,n}(D_{n+1})=p_{n+1,n}^*D_n
 \label{eq:v97-compatible-section}
\end{equation}
becomes exact scalar projective consistency
\[
 d_{n+1}=d_n\circ p_{n+1,n}.
\]
\end{theorem}

\begin{proof}
Define \(\tau_{n+1}\) recursively by
\eqref{eq:v97-compatible-frame}.  Since every \(\Phi_{n+1,n}\) is an
isomorphism, each \(\tau_{n+1}\) is a frame.  Coherence
\eqref{eq:v97-scale-coherence} makes the result independent of grouping
over several scale steps.  Uniqueness follows from the recursion.
Applying the compatible frames to
\eqref{eq:v97-compatible-section} gives the scalar identity.
\end{proof}

The third hypothesis prevents a common bookkeeping error.  Integrating a
fermionic shell can contribute a determinant line or phase that is not
part of the pullback of the coarse line.  Such a factor must be included
in \(\Phi_{n+1,n}\), not silently treated as a scalar normalization.

\subsection{Local scale cocycles}

Choose covers \(U_{n,i}\) of \(B_n\), with spatial transition functions
\(g^n_{ij}\).  On a fine chart contained in
\(p_{n+1,n}^{-1}(U_{n,i})\), represent
\(\Phi_{n+1,n}\) by a nonzero scale function \(k^n_{a i}\).  Compatibility
on a fine overlap requires the commuting-square identity
\begin{equation}
 g^{n+1}_{ab}\,k^n_{b j}
 =
 k^n_{a i}\,(g^n_{ij}\circ p_{n+1,n}).
 \label{eq:v97-space-scale-square}
\end{equation}
The three-scale condition gives a second cocycle equation for the
\(k^n\)'s.  Together, the spatial and scale functions form one cocycle on
the regulator-chart groupoid.

\begin{proposition}[Combined cocycle criterion]
\label{prop:v97-combined-cocycle}
Compatible scalar frames on all charts and scales exist if and only if
the combined spatial-scale cocycle is a coboundary.  Failure of either
\eqref{eq:v97-Cech-cocycle},
\eqref{eq:v97-space-scale-square}, or
\eqref{eq:v97-scale-coherence} prevents a scalar projective determinant.
\end{proposition}

\begin{proof}
A compatible collection of local frame multipliers trivializes every
spatial transition and every scale transition, so the combined cocycle is
a coboundary.  Conversely, a coboundary supplies those multipliers;
Theorem \ref{thm:v97-scalarization} glues within each scale and
Theorem \ref{thm:v97-projective-scalarization} glues across scales.
\end{proof}

\subsection{Summably approximate coefficient transport}

Suppose compatible frames have been constructed.  Let \(d_n\) be the
resulting scalar fermionic coefficient, possibly a finite Grassmann
coefficient rather than a positive density.  For a fixed cylinder
\((k,A)\), let \(\|\cdot\|_{k,A}^{\rm coeff}\) be the coefficient norm
used in V79.

\begin{theorem}[Coefficient repair in a compatible frame]
\label{thm:v97-coefficient-repair}
If
\begin{equation}
 \|d_{n+1}-d_n\circ p_{n+1,n}\|_{k,A}^{\rm coeff}
 \leq e_n^{\rm det}(k,A),
 \qquad
 \sum_{n\geq k}e_n^{\rm det}(k,A)<\infty,
 \label{eq:v97-coefficient-error}
\end{equation}
then the descended coefficient on every fixed cylinder is Cauchy and has
a unique limit.  The limit is independent of the compatible frame up to
one global nowhere-zero scalar function pulled back consistently from the
base scale.
\end{theorem}

\begin{proof}
Telescope \eqref{eq:v97-coefficient-error} in the complete coefficient
Banach space.  If \(\tau_n'\) is another compatible frame, the ratio
\(\tau_n'/\tau_n\) pulls back exactly under every scale map by
\eqref{eq:v97-compatible-frame}; hence it is determined by its base-scale
ratio.  Multiplication by that fixed compatible ratio relates the two
limits.
\end{proof}

This theorem cannot even be stated invariantly in unrelated local
frames; phase jumps on overlaps can masquerade as adjacent-scale errors.

\subsection{Reflection-compatible real structure}

Let \(\Theta:B\to B\) be reflection.  A determinant line has a compatible
antilinear lift when there are fibre-antilinear maps
\begin{equation}
 {\cal J}_b:L_b\longrightarrow L_{\Theta b},
 \qquad
 {\cal J}_{\Theta b}{\cal J}_b=I.
 \label{eq:v97-real-structure}
\end{equation}
A scalarizing frame \(\tau\) is reflection compatible when
\begin{equation}
 \tau_{\Theta b}({\cal J}_b\ell)
 =
 \overline{\tau_b(\ell)}.
 \label{eq:v97-real-frame}
\end{equation}

\begin{proposition}[Reflection frame obstruction]
\label{prop:v97-reflection-frame}
A complex trivialization of \(L\) need not admit a rescaling satisfying
\eqref{eq:v97-real-frame}.  Such a frame exists if and only if the
induced Real line cocycle is a coboundary.  If it exists and the scale
maps commute with \({\cal J}\), the recursive construction in Theorem
\ref{thm:v97-projective-scalarization} preserves it.
\end{proposition}

\begin{proof}
Writing \({\cal J}\) in local frames produces an antilinear transition
cocycle.  A rescaling satisfies \eqref{eq:v97-real-frame} exactly when it
trivializes that cocycle.  If \(\Phi_{n+1,n}{\cal J}={\cal J}
\Phi_{n+1,n}\), apply the recursion
\eqref{eq:v97-compatible-frame} and complex conjugate; the same identity
holds at the next scale.
\end{proof}

This Real cocycle is part of the V94--V96 reflection card.  A complex
scalar determinant with an arbitrary phase convention is insufficient
for OS reconstruction.

\subsection{Application boundary for the Standard Model}

For one Standard Model generation, perturbative representation-theoretic
cancellation is expected to remove the local gauge and
gauge-gravitational curvature of the combined determinant line.
That algebraic statement must still be matched to:
\begin{enumerate}
\item the regulator's exact chiral representation and hypercharge
      normalization;
\item the gauge and gravitational orbit space, including large
      transformations;
\item spin, boundary, and topological sectors;
\item any mod-two or eta-invariant holonomy;
\item the coarse determinant-line maps \(\Phi_{n+1,n}\); and
\item the reflection Real structure.
\end{enumerate}
V98 will compute the perturbative Standard Model anomaly sums and the
\(SU(2)\) doublet parity.  Those calculations can close only the
finite-dimensional representation rows, not the projective gluing above.

\begin{firewall}[Local anomaly cancellation is not global gluing]
\label{fire:v97-local-not-global}
Vanishing of the local anomaly polynomial corresponds to cancellation of
the determinant-line curvature in gauge directions.  Theorem
\ref{thm:v97-flat-holonomy} shows that a flat line can still have
nontrivial holonomy.  Boundary eta phases, large gauge transformations,
gravitational loops, scale-transition cocycles, and the reflection Real
structure remain independent checks.
\end{firewall}

\begin{definition}[V97 determinant-line card]
\label{def:v97-line-card}
The fermionic sector passes the card if:
\begin{enumerate}
\item the combined determinant line is defined on every regulator chart;
\item its spatial transitions satisfy the Cech cocycle;
\item local anomaly curvature cancels;
\item all global and equivariant holonomies are trivial in the chosen
      physical theory;
\item scale maps satisfy the spatial-scale square and three-scale
      coherence;
\item a compatible scalarizing frame exists;
\item coefficient errors are summable in that frame; and
\item the frame and scale maps preserve the reflection Real structure.
\end{enumerate}
\end{definition}

\begin{assessment}[Exact status after V97]
\label{ass:v97-status}
V97 proves the scalarization, flat-holonomy, combined
spatial-scale-cocycle, compatible-frame, coefficient-repair, and
reflection-frame criteria.  It identifies exactly what determinant-line
data V79 and V96 require.  It does not prove the Standard Model anomaly
sums, global holonomy, scale cocycle, or Real trivialization.  The
gravitational contour, weighted exterior, retained-field tightness, and
the full continuum remain open.
\end{assessment}

\subsection{Next acquisition}

V98 will compute the local four-dimensional Standard Model gauge anomaly
coefficients per generation, including
\(SU(3)^2U(1)\), \(SU(2)^2U(1)\), \(U(1)^3\), and
gravity-squared--\(U(1)\), and will check the parity of left-handed
\(SU(2)\) doublets.  It will state precisely which determinant-line rows
these algebraic cancellations do and do not close.

\paragraph{Parent-version record.}
V97 was formed from
\texttt{Renewal\_SM\_Einstein\_Continuum\_Research\_Notes\_V96.tex}.
The V96 parent had \(33\,699\) counted source lines, \(1\,221\,052\) bytes,
and SHA--256
\[
 \texttt{1CEF60893BC38492A43BF0A63C3503472B2DC9A6EB4A35A5E6B8F0CC7E2BD049}.
\]

% =============================================================================
% END APPENDED TENTH-SERIES V97 MATERIAL
% =============================================================================
% =============================================================================
% BEGIN APPENDED TENTH-SERIES V98 MATERIAL
% =============================================================================

\section{Tenth-series V98: Standard Model local anomaly sums}
\label{sec:v10v98}

V97 distinguishes local determinant-line curvature from global holonomy.
This note computes the four-dimensional perturbative gauge and mixed
gauge--gravitational anomaly coefficients of one Standard Model
generation.  The calculation closes the finite representation-theoretic
rows only.  It does not trivialize the global or projective determinant
line.

\subsection{Left-handed convention}

Use the convention
\[
 Q_{\rm electric}=T_3+Y
\]
and express every fermion as a left-handed Weyl field.  One generation is
\[
 \begin{array}{c|c|c|c|c}
 \text{field}&SU(3)&SU(2)&Y&
 \dim(SU(3))\dim(SU(2))\\ \hline
 Q_L&{\bf3}&{\bf2}&1/6&6\\
 u_R^c&\overline{\bf3}&{\bf1}&-2/3&3\\
 d_R^c&\overline{\bf3}&{\bf1}&1/3&3\\
 L_L&{\bf1}&{\bf2}&-1/2&2\\
 e_R^c&{\bf1}&{\bf1}&1&1
 \end{array}
 \label{tab:v98-generation}
\]
A sterile \(\nu_R^c\sim({\bf1},{\bf1})_0\), if included, contributes zero
to every coefficient below.  Scalars, gauge bosons, and
Faddeev--Popov ghosts do not add a chiral matter anomaly.

Normalize the quadratic index by
\[
 T({\bf3})=T(\overline{\bf3})=
 T({\bf2})=\frac12
\]
and the cubic \(SU(3)\) index by
\[
 A({\bf3})=1,\qquad A(\overline{\bf3})=-1.
\]
Only cancellation, not the overall anomaly-polynomial normalization, is
at issue.

\subsection{Pure non-Abelian anomalies}

\begin{proposition}[Pure non-Abelian cancellation]
\label{prop:v98-pure-nonabelian}
For one generation,
\begin{equation}
 {\cal A}_{SU(3)^3}
 =
 2A({\bf3})+A(\overline{\bf3})
 +A(\overline{\bf3})=0.
 \label{eq:v98-SU3-cubic}
\end{equation}
The perturbative \(SU(2)^3\) anomaly vanishes representationwise because
the symmetric cubic invariant of \(SU(2)\) is zero.
\end{proposition}

\begin{proof}
The two components of the weak doublet \(Q_L\) give two colour
fundamentals.  The left-handed conjugates \(u_R^c,d_R^c\) are colour
antifundamentals, giving
\(2-1-1=0\).  For \(SU(2)\), every generator trace
\(\operatorname{tr}(T^a\{T^b,T^c\})\) is proportional to a symmetric
invariant tensor \(d^{abc}\), which vanishes for \(\mathfrak{su}(2)\).
\end{proof}

Mixed anomalies with one non-Abelian generator and two Abelian insertions
vanish within each irreducible non-Abelian multiplet because
\(\operatorname{tr}T^a=0\).  Mixed non-Abelian--gravitational traces
vanish for the same reason.

\subsection{Two non-Abelian and one hypercharge insertion}

\begin{proposition}[Mixed non-Abelian--hypercharge cancellation]
\label{prop:v98-mixed-nonabelian}
For one generation,
\begin{align}
 {\cal A}_{SU(3)^2U(1)}
 &=
 2\left(\frac16\right)\frac12
 +\left(-\frac23\right)\frac12
 +\left(\frac13\right)\frac12
 =0,
 \label{eq:v98-SU3SU3Y}\\
 {\cal A}_{SU(2)^2U(1)}
 &=
 3\left(\frac16\right)\frac12
 +\left(-\frac12\right)\frac12
 =0.
 \label{eq:v98-SU2SU2Y}
\end{align}
\end{proposition}

\begin{proof}
For \(SU(3)^2U(1)\), multiply the hypercharge and colour quadratic index
by the weak multiplicity.  For \(SU(2)^2U(1)\), multiply by the colour
multiplicity.  The displayed arithmetic gives zero.
\end{proof}

\subsection{Cubic hypercharge and mixed gravitational anomaly}

\begin{theorem}[Abelian anomaly cancellation]
\label{thm:v98-Abelian-cancellation}
For one generation,
\begin{align}
 {\cal A}_{U(1)^3}
 &=
 6\left(\frac16\right)^3
 +3\left(-\frac23\right)^3
 +3\left(\frac13\right)^3
 +2\left(-\frac12\right)^3
 +1^3
 \nonumber\\
 &=
 \frac1{36}-\frac89+\frac19-\frac14+1
 =0,
 \label{eq:v98-Y-cubic}\\
 {\cal A}_{\operatorname{grav}^2U(1)}
 &=
 6\left(\frac16\right)
 +3\left(-\frac23\right)
 +3\left(\frac13\right)
 +2\left(-\frac12\right)+1
 \nonumber\\
 &=1-2+1-1+1=0.
 \label{eq:v98-grav-grav-Y}
\end{align}
\end{theorem}

\begin{proof}
The multiplicity of each Weyl field is the dimension of its non-Abelian
representation.  Bringing the cubic sum to denominator \(36\) gives
\[
 1-32+4-9+36=0.
\]
The linear sum is the displayed integer cancellation.
\end{proof}

In four spacetime dimensions there is no perturbative pure
diffeomorphism anomaly for an ordinary Weyl fermion without a gauge
insertion: the degree-six index polynomial has no purely gravitational
term because the \(\widehat A\)-class has components in degrees divisible
by four.  The mixed term is proportional to
\(\sum Y\) and vanishes by
\eqref{eq:v98-grav-grav-Y}.  This statement does not exclude trace or
axial anomalies, which are different Ward identities.

\subsection{Perturbative anomaly theorem}

\begin{theorem}[One-generation local anomaly polynomial vanishes]
\label{thm:v98-local-polynomial}
For the left-handed representation in
Table \ref{tab:v98-generation}, every perturbative four-dimensional gauge
and mixed gauge--gravitational anomaly coefficient vanishes.  The result
remains true for any number of identical generations and after adding
hypercharge-zero sterile right-handed neutrinos.
\end{theorem}

\begin{proof}
The possible coefficients are the pure non-Abelian cubic traces, the
two-non-Abelian--one-\(U(1)\) traces, the cubic \(U(1)\) trace, and the
linear \(U(1)\) trace multiplying the gravitational class.  Propositions
\ref{prop:v98-pure-nonabelian} and
\ref{prop:v98-mixed-nonabelian}, together with Theorem
\ref{thm:v98-Abelian-cancellation}, make all of them zero.  Coefficients
add over generations, and a completely neutral singlet adds zero.
\end{proof}

\subsection{The \(SU(2)\) mod-two row}

Perturbative \(SU(2)^3\) cancellation does not address the global
mod-two anomaly.

\begin{proposition}[Even doublet parity]
\label{prop:v98-Witten-parity}
One Standard Model generation contains an even number of left-handed
\(SU(2)\) fundamental doublets:
\begin{equation}
 N_{\bf2}=3\ \text{colour copies of }Q_L
 +1\ \text{copy of }L_L=4.
 \label{eq:v98-doublet-count}
\end{equation}
Hence the elementary \(SU(2)\) doublet-parity obstruction is absent per
generation and for three generations.
\end{proposition}

\begin{proof}
Colour provides three independent copies of the quark weak doublet, and
the lepton doublet provides the fourth.  The parity of \(4\) and of
\(12\) is even.
\end{proof}

This checks the usual fundamental-doublet mod-two row.  It does not
classify every possible global anomaly of the actual global gauge group,
especially after quotienting
\(SU(3)\times SU(2)\times U(1)\) by a common discrete subgroup or coupling
to nontrivial spin and gravitational backgrounds.

\subsection{What this implies for the determinant line}

For a family of chiral Dirac operators, the curvature of the combined
determinant-line connection in infinitesimal gauge directions is given by
the corresponding local anomaly polynomial.  Therefore Theorem
\ref{thm:v98-local-polynomial} supplies the representation-theoretic
cancellation needed for the curvature row of the V97 card, provided the
common regulator:
\begin{enumerate}
\item realizes exactly the stated chiral multiplets and hypercharges;
\item treats all members of a generation with compatible gauge and
      gravitational regularization;
\item includes the same local counterterm convention in every chart; and
\item identifies its determinant-line curvature with the regulated index
      density without an additional boundary term.
\end{enumerate}

\begin{corollary}[Conditional flatness row]
\label{cor:v98-conditional-flatness}
Under the four regulator-identification hypotheses above, the combined
Standard Model determinant-line curvature in infinitesimal gauge
directions vanishes generation by generation.
\end{corollary}

\begin{proof}
The curvature is linear in the anomaly polynomial, whose group
coefficients vanish by Theorem \ref{thm:v98-local-polynomial}.
\end{proof}

The qualifier \emph{conditional} is essential.  A boundary family can
contribute an eta-form, and an inconsistent regulator can violate the
common multiplet cancellation even when the continuum representation sum
is zero.

\subsection{Rows not closed by the arithmetic}

The following V97 requirements remain:
\begin{enumerate}
\item holonomy under large gauge transformations and diffeomorphisms;
\item the effect of the chosen global form of the Standard Model gauge
      group;
\item spin, spin-\(c\), bundle, and boundary sectors;
\item eta-invariant and mod-two phases beyond the elementary doublet
      count;
\item coherent determinant-line maps between regulator scales;
\item a global scalarizing frame;
\item reflection-compatible Real trivialization; and
\item summable coefficient errors in that frame.
\end{enumerate}

\begin{firewall}[Anomaly arithmetic is not a fermion measure]
\label{fire:v98-not-fermion-measure}
The equalities
\eqref{eq:v98-SU3-cubic}--\eqref{eq:v98-grav-grav-Y} prove cancellation
of the local representation traces.  They do not define a chiral lattice
measure, prove positivity, construct the determinant-line cocycle, or
remove global and boundary anomalies.  Proposition
\ref{prop:v98-Witten-parity} checks one mod-two obstruction only.
\end{firewall}

\begin{definition}[V98 anomaly card]
\label{def:v98-anomaly-card}
The local algebraic part of the fermionic card consists of:
\begin{enumerate}
\item \(SU(3)^3=0\);
\item perturbative \(SU(2)^3=0\);
\item \(SU(3)^2U(1)=0\);
\item \(SU(2)^2U(1)=0\);
\item \(U(1)^3=0\);
\item \(\operatorname{grav}^2U(1)=0\);
\item vanishing one-generator mixed traces; and
\item even \(SU(2)\) fundamental-doublet parity.
\end{enumerate}
All eight rows pass for the representation in
Table \ref{tab:v98-generation}.  They enter V97 as inputs, not as a
replacement for its global gluing rows.
\end{definition}

\begin{assessment}[Exact status after V98]
\label{ass:v98-status}
V98 proves all perturbative gauge and mixed gauge--gravitational anomaly
sums for one Standard Model generation and checks even \(SU(2)\) doublet
parity.  Conditional on a regulator realizing the common index-density
connection, the local determinant-line curvature cancels.  Global
holonomy, boundary eta data, scale coherence, reflection Real structure,
the gravitational contour and exterior, retained-field tightness, and
the full SM--Einstein continuum remain open.
\end{assessment}

\subsection{Next acquisition}

V99 will combine V83, V84, V92, V96, V97, and V98 into one
continuum-readiness theorem with a minimal nonduplicating list of
remaining hypotheses.  It will distinguish proved implication theorems
from model-specific acquisition rows before the V100 audit and rollover.

\paragraph{Parent-version record.}
V98 was formed from
\texttt{Renewal\_SM\_Einstein\_Continuum\_Research\_Notes\_V97.tex}.
The V97 parent had \(34\,087\) counted source lines, \(1\,235\,651\) bytes,
and SHA--256
\[
 \texttt{08490DEFE407C307F76673FDADFEE80D51D1DCFF56D614021549C74EC55CDA36}.
\]

% =============================================================================
% END APPENDED TENTH-SERIES V98 MATERIAL
% =============================================================================
% =============================================================================
% BEGIN APPENDED TENTH-SERIES V99 MATERIAL
% =============================================================================

\section{Tenth-series V99: continuum-readiness theorem and acquisition ledger}
\label{sec:v10v99}

The preceding notes prove many implication theorems but do not yet supply
all model-specific hypotheses.  This note assembles them once, without
turning a conditional implication into a construction.  It distinguishes
the positive bosonic body, the fermionic coefficient system, topological
phases, and physical reconstruction.

\subsection{Regulator-tower data}

A candidate SM--Einstein shot tower consists of:
\begin{enumerate}
\item finite regulator state spaces \(X_n\), retained sector labels
      \(\Sigma\), and local coarse maps \(p_{n+1,n}\);
\item positive bosonic probabilities \(\mu_n^\Sigma\) on
      \(X_n\times\Sigma\);
\item reconstruction maps \(\iota_n:X_n\to{\cal S}\) into a Polish
      distribution space;
\item exact local constraint, gauge, boundary, and interface residual
      charts;
\item exact triangular fine-fibre integrations and their normalized
      conditional owner factors;
\item a combined chiral determinant line \(L_n\to B_n\), its section
      \(D_n\), and coherent scale maps;
\item finite-dimensional physical, redundant, irrelevant, marginal, and
      unstable coupling coordinates; and
\item reflection, source, and observable algebras compatible with every
      coarse map.
\end{enumerate}

The qualifier \emph{positive} applies to the bosonic body after the
physical contour and gauge quotient.  Fermionic observables are recovered
as coefficient functionals in the exterior algebra.  A nonzero theta
phase can turn the scalar functional complex even when the zero-theta
body is positive.

\subsection{The readiness hypotheses}

For later reference, group the remaining requirements as follows.

\paragraph{(R1) Finite-regulator definition and coercivity.}
Every finite regulator integral is well defined.  The chosen conformal
and gauge contour has a uniform real-part lower bound, a local
normalization lower bound, and a weighted physical exterior estimate.
The contour has an admissible Lorentzian or thimble derivation and a
reflection-compatible physical quotient.

\paragraph{(R2) Exact local geometry.}
The independent constraint and gauge residual has constant rank after
Bianchi, Killing, zero, and boundary dependencies are retained.  The
V63 nested inverse and V73 analytic residual chart have uniform locality,
pivot, and radius bounds.

\paragraph{(R3) Exact triangular pushforward.}
The fine fibres are genuinely disjoint coordinates, their colour
integrations obey Fubini or Tonelli, and pass, hit, jet, determinant, and
bad-component owners follow a unique-payment history.  The connected
forest passes its numerical gate.

\paragraph{(R4) Summable local comparison.}
For every fixed cylinder \((k,A)\), the sector-summed adjacent error in
V96 is summable.  Quasi-local leakage satisfies the V83 logarithmic
budget.  Weighted errors for the Sobolev shell observable are summable.

\paragraph{(R5) Tightness.}
Ultraviolet shell moments pass V84 for some
\(s_0>(D-2)/2\); retained massless and infrared variables are tight;
zero, boundary, charge, and moduli variables are tight in their physical
quotients; and the discrete sector label is tight.  On a noncompact base,
spatial escape is controlled separately.

\paragraph{(R6) Fermionic line gluing.}
The V97 spatial-scale determinant cocycle is a coboundary, local anomaly
curvature cancels, global holonomy is trivial in the chosen physical
theory, scale maps are coherent, and the scalarizing frame respects the
reflection Real structure.  Coefficient errors are summable in that
frame.

\paragraph{(R7) Symmetry and coupling flow.}
The measured-coordinate pivot of V91 is uniform.  Irrelevant
symmetry-breaking jets pass V92.  All marginal and relevant physical
couplings satisfy their own calibrated beta and tuning equations, and
the resulting cofinal orbit remains in the common owner domain.

\paragraph{(R8) Identities and support.}
Constraint residual moments vanish; gauge, diffeomorphism, BRST or
physical quotient Ward identities hold with uniformly integrable
insertions; and topological labels and phases obey V96.

\paragraph{(R9) Euclidean axioms.}
The limiting cylinder functional is Euclidean covariant in the restored
continuum sense, reflection positive on the physical algebra, regular,
symmetric or graded-symmetric as appropriate, and satisfies the
clustering or spectral hypotheses required for the intended
reconstruction.

\paragraph{(R10) Cofinal uniqueness.}
Any two admissible regulator towers possess a common refinement whose
local cylinder and coefficient comparison tails vanish.  Physical
renormalization conditions select the same retained coupling orbit.

\subsection{Conditional continuum theorem}

\begin{theorem}[SM--Einstein continuum readiness]
\label{thm:v99-continuum-readiness}
If a regulator tower satisfies \emph{(R1)--(R10)}, then:
\begin{enumerate}
\item the zero-theta bosonic bodies converge to a unique Radon probability
      \(P\) on \({\cal S}\times\Sigma\);
\item the limit has the Sobolev support and sector tightness specified in
      \emph{(R5)};
\item compatible fermionic coefficients define a unique continuum
      Grassmann-valued cylinder functional in the V97 frame;
\item bounded topological characters define the complex theta
      functionals whenever their partition values are nonzero;
\item the constraints and uniformly integrable Ward identities in
      \emph{(R8)} pass to the limit;
\item the physical zero-theta functional satisfies the Euclidean axioms
      in \emph{(R9)} and therefore admits the corresponding OS
      reconstruction; and
\item the result is independent of the admissible cofinal regulator tower
      in the sense of all physical cylinder observables.
\end{enumerate}
\end{theorem}

\begin{proof}
Hypotheses (R3)--(R4) give the sector-summed cylinder estimate
\eqref{eq:v96-sector-adjacent}.  Theorem
\ref{thm:v96-sector-repair} constructs compatible joint local marginals.
Hypothesis (R5) and Theorem \ref{thm:v96-joint-tightness} realize them as
the unique Radon probability \(P\).  The weighted shell estimates and
Theorem \ref{thm:v84-sector-tightness} give the declared Sobolev support.

Hypothesis (R6), Theorem
\ref{thm:v97-projective-scalarization}, and Theorem
\ref{thm:v97-coefficient-repair} turn the local fermionic sections into
compatible continuum coefficients.  Theorem
\ref{thm:v96-phase-reconstruction} constructs the theta functionals.
Hypothesis (R8), Corollary \ref{cor:v83-local-identities}, and uniform
integrability pass the constraint and Ward rows.

Hypothesis (R9) is exactly the remaining input of the
Osterwalder--Schrader reconstruction theorem on the physical algebra, so
the stated reconstruction follows.  Finally, (R10) and the
common-refinement argument of Corollary
\ref{cor:v79-common-refinement}, applied in the cylinder-local and
coefficient norms of V83 and V97, identify the physical observables from
any two towers.
\end{proof}

The theorem is an implication.  It does not say that the current
programme has verified (R1)--(R10).

\subsection{Rows already proved as abstract mathematics}

The following implication mechanisms no longer need to be rediscovered:
\begin{enumerate}
\item exact residual scaling and anisotropic analytic gain
      (V72--V73);
\item exact two-jet gluing, Taylor control, and joint Gaussian chart exit
      (V81--V82);
\item cylinder-local and weighted projective repair (V83--V84);
\item normalized local kernel products and invariant resolvents
      (V86);
\item Gaussian influence, its massless obstruction, and projected-shell
      OU contraction (V87--V88);
\item exact local Metropolis correction and rejection ownership (V89);
\item multi-direction measured-coordinate section and irrelevant-jet
      contraction (V91--V92);
\item coercivity-to-tail, mixed-quartic, and conformal-contour criteria
      (V93--V94);
\item sectorwise projective reconstruction and reflection phases
      (V96);
\item determinant-line cocycle and coherent-frame criteria (V97); and
\item the Standard Model local anomaly sums and doublet parity (V98).
\end{enumerate}

These are reusable theorems.  Their hypotheses must be checked on the
actual SM--Einstein regulator.

\subsection{Model-specific acquisition ledger}

The unresolved rows, in dependency order, are:
\begin{enumerate}
\item \textbf{Physical conformal cycle.}
      Derive the integration cycle from the Lorentzian or constrained
      theory and prove its nonlinear, weighted, reflection-compatible
      coercivity.
\item \textbf{Gravity-compatible chiral regulator.}
      Construct the combined gauge--gravity fermion measure and its
      determinant line at finite cutoff.
\item \textbf{Uniform local residual chart.}
      Verify the constant-rank selection, nested inverse locality, coarea
      density, and boundary compatibility on the actual shot domain.
\item \textbf{Bad-component witness.}
      Prove a joint action lower bound for irregular metric, conformal,
      gauge, Higgs, and interface collars, with unique payment.
\item \textbf{Marginal and relevant flow.}
      Compute and control the coupled Newton, cosmological, gauge, Higgs,
      Yukawa, and permitted topological coordinates with physical
      renormalization conditions.
\item \textbf{Retained infrared tightness.}
      Bound graviton, photon, global gauge, boundary, moduli, and charge
      variables outside the contracted ultraviolet shell.
\item \textbf{Global determinant holonomy.}
      Close large gauge, diffeomorphism, spin, boundary eta, scale, and
      reflection cocycles after the V98 local cancellation.
\item \textbf{Continuum symmetry and OS rows.}
      Prove restoration of the physical diffeomorphism, local Lorentz,
      gauge, and Euclidean symmetries and reflection positivity on the
      quotient algebra.
\item \textbf{Common-refinement estimates.}
      Compare distinct meshes, gauges, contours, and block maps in the
      local weighted norms.
\end{enumerate}

The first four are upstream of any honest global density domination.
The fifth and sixth determine the cofinal orbit and tightness.  The last
three identify the physical continuum rather than one regulator-dependent
limit.

\subsection{Finite audit matrix}

\[
 \begin{array}{p{0.29\linewidth}|p{0.28\linewidth}|p{0.31\linewidth}}
 \text{row}&\text{proved mechanism}&\text{missing acquisition}\\ \hline
 local analytic shell&
 V72--V73, V81--V82&
 gravity-specific uniform chart and exterior\\
 local-to-continuum measure&
 V83--V84, V96&
 actual summable errors and all-sector tightness\\
 symmetry quotient&
 V91--V92&
 compensation pivot and named SM--Einstein jet matrix\\
 fermions&
 V97--V98&
 regulator line, global holonomy, scale and Real cocycles\\
 conformal field&
 V93--V94&
 physical contour derivation and nonlinear field coercivity\\
 physical reconstruction&
 V79, V83, V96&
 complete OS, Ward, clustering, and common-refinement rows
 \end{array}
 \label{tab:v99-readiness}
\]

\begin{firewall}[Readiness is not construction]
\label{fire:v99-readiness-only}
Theorem \ref{thm:v99-continuum-readiness} proves that the listed
quantitative and geometric acquisitions would suffice.  It does not
supply them.  In particular, V98's anomaly arithmetic cannot replace a
chiral regulator, and V94's convergent imaginary Gaussian cannot replace
a physical derivation of the conformal cycle.
\end{firewall}

\begin{assessment}[Exact status after V99]
\label{ass:v99-status}
V99 gives a nonduplicating continuum-readiness theorem and reduces the
open programme to nine model-specific acquisition families.  The
abstract passage from those hypotheses to a sectorwise, fermionic,
reflection-positive, regulator-independent continuum is proved.  Several
central hypotheses, beginning with the physical conformal cycle and
gravity-compatible chiral regulator, remain unproved.  The full
SM--Einstein continuum is not yet constructed.
\end{assessment}

\subsection{Next acquisition}

V100 is the mandatory companion audit and the terminal note of this
compact cycle.  It will inspect the newest Yang--Mills and Navier--Stokes
notes, update the acquisition order, and prepare the required rollover.
The completed V100 file will then be frozen as the next
\(\texttt{\_fN}\) archive, and a new compact V1 will summarize the major
mathematical results and resume at the first upstream acquisition.

\paragraph{Parent-version record.}
V99 was formed from
\texttt{Renewal\_SM\_Einstein\_Continuum\_Research\_Notes\_V98.tex}.
The V98 parent had \(34\,401\) counted source lines, \(1\,246\,744\) bytes,
and SHA--256
\[
 \texttt{95F2AA2AD6F209330551E84293EC26912206957313C130CD1B7A95AFD9E6B2B5}.
\]

% =============================================================================
% END APPENDED TENTH-SERIES V99 MATERIAL
% =============================================================================
% =============================================================================
% BEGIN APPENDED TENTH-SERIES V100 MATERIAL
% =============================================================================

\section{Tenth-series V100: terminal companion audit and rollover direction}
\label{sec:v10v100}

\subsection{Compulsory companion audit}

The newest active companion notes at the V100 gate were inspected
read-only:
\begin{align*}
 &\texttt{Renewal\_Yang\_Mills\_Mass\_Gap\_Research\_Notes\_V17.tex},
 \\
 &\qquad
 \operatorname{SHA256}
 =
 \texttt{E8ED68538DE49A4F657A2DE6041F32FAB1406ACDD4A6FAE43D166D752848AA0C},
 \\
 &\texttt{Renewal\_NS\_Stability\_Research\_Notes\_V26.tex},
 \\
 &\qquad
 \operatorname{SHA256}
 =
 \texttt{82C887F8C2C69E4F28FAD7C3DC8DE5B9099CB5A4BF431EBA1FFF3E91935978AD}.
\end{align*}
The Yang--Mills companion advanced from V13 to V17 after the V95 audit.
The Navier--Stokes companion remains unchanged.

\subsection{Averaged marginal matching versus projective convergence}

The Yang--Mills V14 note proves that absolute summability of every
same-scale reset mismatch is stronger than necessary for one averaged
marginal beta coefficient.  An exact shell determinant telescopes before
reference resets.  A reset contribution of coboundary form
\begin{equation}
 r_j=B_{j+1}-B_j
 \label{eq:v100-reset-coboundary}
\end{equation}
has cumulative memory
\begin{equation}
 \sum_{j=0}^{N-1}r_j=B_N-B_0.
 \label{eq:v100-reset-telescope}
\end{equation}
Therefore
\begin{equation}
 B_N=o(N)
 \quad\Longrightarrow\quad
 \frac1N\sum_{j=0}^{N-1}r_j\longrightarrow0,
 \label{eq:v100-sublinear-average}
\end{equation}
even when \(\sum_j|r_j|=\infty\).

This observation transfers only to averaged marginal coefficients.  The
V83 and V97 construction of an actual projective cylinder coefficient
requires convergence of the coefficient itself, for which a merely
Cesaro-small error is insufficient.  Thus two norms must remain distinct:
\begin{enumerate}
\item summable local coefficient errors for the continuum functional; and
\item sublinear cumulative coboundaries for a universal averaged beta
      coefficient.
\end{enumerate}
Using \eqref{eq:v100-sublinear-average} to claim projective determinant
convergence would be invalid.

\subsection{Exact intertwiners produce scale coboundaries}

The Yang--Mills V15 note shows that an exact measured conjugacy between
two regulator descriptions changes the one-loop scale cochain by an
endpoint coboundary.  This fits the V97 scale-square architecture.  In
abstract form, if invertible maps \(T_n\) intertwine two exact scale
operators,
\begin{equation}
 \widetilde R_n
 =
 T_{n+1}R_nT_n^{-1},
 \label{eq:v100-scale-intertwiner}
\end{equation}
then their logarithmic determinant increments differ by
\begin{equation}
 \log\det\widetilde R_n-\log\det R_n
 =
 \log\det T_{n+1}-\log\det T_n,
 \label{eq:v100-determinant-coboundary}
\end{equation}
with one coherent determinant branch.  Summation leaves endpoints only.

For the SM--Einstein determinant line, the scalar formula
\eqref{eq:v100-determinant-coboundary} is available only after the V97
line is trivialized.  Before that, the statement is an isomorphism of
line cocycles, not an equality of arbitrary local logarithms.

\subsection{Ward recombination must precede spectral splitting}

The Yang--Mills V16--V17 notes make the most important transferable
correction at this gate.  Low-mode, edge, alias, Haar, routing, quotient,
and overlap response words must first be placed in the complete Ward sum.
Only the recombined density may then be split into low and edge regions.
Splitting individual response words first can manufacture boundary terms
that cancel in the full Ward identity and can therefore produce a false
anomaly.

For the Standard Model this means that the following contributions must
be assembled before a spectral or position-space cutoff is estimated:
\begin{enumerate}
\item all left-handed fermion multiplets of one generation;
\item conjugate right-handed fields in the same hypercharge convention;
\item gauge-fixing and ghost Jacobians where they enter the same Ward
      identity;
\item coarea, routing, block-map, and overlap determinants;
\item boundary eta or reflection-plane terms; and
\item local counterterms fixed by one common regulator convention.
\end{enumerate}
The V98 representation sums show that the continuum local polynomial of
the complete fermion multiplet vanishes.  They do not justify applying
different cutoffs to the separate summands and canceling only afterward.

The Yang--Mills V17 reduction of its residual edge problem to one scalar
alias class is model-specific.  Its Fourier decimation constants and
finite-grid certificate do not transfer to the coupled chiral
SM--Einstein operator.  The general transfer is the Ward-first order of
operations and the cohomological distinction between a coboundary and a
nonzero mean class.

\subsection{No Navier--Stokes transfer}

The Navier--Stokes V26 calculation contains neither a determinant
cocycle nor a chiral Ward sum, topological sector, conformal contour, or
quantum reflection algebra.  It adds no input at V100.

\subsection{Direction after rollover}

The first new-cycle acquisition should construct a Ward-complete
regulator comparison before returning to component estimates.  Its
minimum content is:
\begin{enumerate}
\item a finite-dimensional combined chiral determinant response written
      without choosing unrelated logarithm branches;
\item an exact Ward sum over a complete Standard Model generation and all
      measure contributions;
\item a spectral split applied only to that sum;
\item a proof that changes of local regulator coordinates contribute a
      line-cocycle coboundary;
\item a separate edge or boundary class whose vanishing is a named
      acquisition; and
\item compatibility with the V97 scale and reflection cocycles.
\end{enumerate}
This direction attacks the gravity-compatible chiral regulator row of
V99.  The physical conformal cycle remains an equally upstream,
independent acquisition and is not declared solved.

\begin{definition}[V100 Ward-complete regulator card]
\label{def:v100-Ward-card}
A chiral regulator comparison passes the card if:
\begin{enumerate}
\item every determinant factor belongs to the combined V97 line;
\item the full regulated Ward response is formed before any cutoff split;
\item the V98 local representation traces cancel in that common response;
\item exact rechartings act by coherent scale coboundaries;
\item the residual edge class is computed rather than bounded away by a
      componentwise estimate;
\item boundary and reflection-plane terms are explicit; and
\item projective coefficient errors remain summable even if an averaged
      marginal calculation uses the weaker sublinear criterion.
\end{enumerate}
\end{definition}

\begin{firewall}[Averaged anomaly matching is not a continuum
determinant]
\label{fire:v100-average-not-projective}
A sublinear endpoint discrepancy proves that its Cesaro contribution to
one beta coefficient vanishes.  It does not produce a convergent
determinant coefficient on every cylinder, a global scalarizing frame, or
a reflection-compatible chiral measure.  Those V97 requirements remain.
\end{firewall}

\begin{assessment}[Terminal status of the tenth compact cycle]
\label{ass:v100-status}
The V100 audit imports the sublinear coboundary criterion for averaged
marginal data and, more importantly, the rule that Ward recombination
precedes spectral splitting.  Across V1--V100 this compact cycle has
built the local analytic, projective, weighted-tightness,
measured-coordinate, sector, contour, determinant-line, and local anomaly
implication machinery summarized in V99.  The physical conformal cycle,
gravity-compatible chiral regulator, uniform bad-component witness,
marginal cofinal orbit, retained infrared tightness, global determinant
holonomy, complete OS identities, and common-refinement estimates remain
unproved.  The full SM--Einstein continuum is not constructed.
\end{assessment}

\subsection{Rollover record}

This is the terminal active version of the tenth compact cycle.  Following
the standing version rule, the validated file will be frozen as
\[
 \texttt{Renewal\_SM\_Einstein\_Continuum\_Research\_Notes\_V100\_f10.tex}.
\]
A new compact V1 will give context, summarize the major acquired results,
retain the V99 acquisition ledger, and continue with the Ward-complete
chiral regulator comparison.  The frozen file will not be modified.

\paragraph{Parent-version record.}
V100 was formed from
\texttt{Renewal\_SM\_Einstein\_Continuum\_Research\_Notes\_V99.tex}.
The V99 parent had \(34\,695\) counted source lines, \(1\,259\,873\) bytes,
and SHA--256
\[
 \texttt{D6F51278D3F278EECD2EE1DBF721931833D30EACB9176E576B96A0D5EF60C94A}.
\]

% =============================================================================
% END APPENDED TENTH-SERIES V100 MATERIAL
% =============================================================================
\end{document}
